%!TEX program = lualatex
\documentclass[letterpaper, 11pt]{report}
\usepackage{fontspec}
\usepackage{fancyhdr}

%%% MARGINS %%%
% use this for normal letter documents
% \usepackage[
%     letterpaper,
%     margin=1in,
%     includehead
% ]{geometry}

%use this for bookletized pdfs
\usepackage[
    paperwidth=6.875in,
    paperheight=10.625in,
    margin=0.75in,     % Now your margins will be exactly what you type
    includehead
]{geometry}

%%% FONT %%%
\setmainfont{Liberation Serif}

%%% SPEAKER \& TALK VARIABLES %%%
\newcommand{\speaker}{}
\newcommand{\talk}{}
\newcommand{\talkdatestr}{}
\newcommand{\setspeaker}[1]{\renewcommand{\speaker}{#1}}
\newcommand{\settalk}[1]{\renewcommand{\talk}{#1}}
\newcommand{\settalkdate}[1]{\renewcommand{\talkdatestr}{#1}}

%%% HEADER \& FOOTER SETUP %%%
\pagestyle{fancy}
\fancyhf{}

\fancyhead[L]{\textbf{\talk}}
\fancyhead[R]{\textbf{\speaker}}

\renewcommand{\headrulewidth}{0.4pt}

%%% CUSTOM HEADER CLASS FOR TALK TITLE %%%
\newcommand{\talkheading}[2]{%
    \vspace{1em}
    \noindent{\huge\bfseries #1}
    \par\nopagebreak\vspace{0.25em}
    \noindent{\normalsize\bfseries #2\ifx\talkdatestr\empty\else,\ \talkdatestr\fi}
    \par\nopagebreak\vspace{0.5em}
}

%%% IN-TALK HEADERS
\newcommand{\intalkheader}[1]{
    \vspace{0.75em}
    \noindent{\normalsize\bfseries #1}
    \par\nopagebreak % Prevents a break immediately after the text
    \vspace{0.25em}
    \nopagebreak     % Prevents a break after the whitespace
}

\begin{document}
\newpage
\settalk{Call to the Holy Apostleship}
\setspeaker{Russell M. Nelson}
\settalkdate{April 1984}
\talkheading{Call to the Holy Apostleship}{Russell M. Nelson}

Saturday of April conference of 1984 has been circled on our calendar for many years, for that date was targeted as the first time in my life that our only son would be old enough to attend general priesthood meeting with me. Last night, that long awaited goal became a reality. Brothers and sisters, little did we know that on that day my name would be presented as a member of the Council of the Twelve.

Since we didn’t know, our children didn’t know either. Our married daughters telephoned us between sessions. One who was expecting a baby, said, “Daddy, I was so shocked by that announcement—I think I am going into labor.”

That she did.

So, President Hinckley, your announcement from the First Presidency should get credit at least for “an assist.” Our twenty-second grandchild arrived safely last evening!

A wide array of feelings has flashed through my heart since I heard the call that will change my life. The first feeling is that of personal inadequacy. That feeling is intensified as I think of the incomparable power of Elders LeGrand Richards and Mark E. Petersen, whose absence we keenly sense. They were, to me, dear friends as well as esteemed leaders. Then, as I look about and see the strength of those more qualified and able than I, I truly am humbled by this calling.

Fortunately, these feelings are blanketed by feelings of faith, for I know the words of Nephi are true: “I will go and do the things which the Lord hath commanded, for I know that the Lord giveth no commandments unto the children of men, save he shall prepare a way for them that they may accomplish the thing which he commandeth them.” (1 Ne. 3:7.) I have implicit faith in the Lord and in His prophets. I have learned not to put question marks but to use exclamation points when calls are issued through inspired channels of priesthood government.

About twenty years ago, two members of the Quorum of the Twelve were assigned to select a new president of the Bonneville stake in which we lived. They were Elders Spencer W. Kimball and LeGrand Richards. That call they extended to me. Now, I am to be ordained an Apostle under the hand of President Spencer W. Kimball, to fill a chair vacated by the passing of Elder LeGrand Richards.

Feelings of commitment well up from the depths of my soul. My sweetheart, Dantzel, and I first made those covenants in the temple of the Lord over thirty-eight years ago, to consecrate our lives to the service of the Lord. Today, I reaffirm that promise, to give all I have to the building of the kingdom of God on the earth. In accepting this call, knowing that challenges, charges, and keys will be conferred and that buffetings will likewise come, I commit my effort, my energy, and my all.

Feelings of gratitude emerge as I acknowledge the goodness of my parents, and of all eight of my pioneer great-grandparents who were converted to the Church in the populous nations of Europe. All eight of them ultimately immigrated to the little town of Ephraim, Utah. I know they are watching the events of this day from their celestial windows.

To my dear Dantzel, I acknowledge my debt, as well as my eternal love. She is the fountain from whom flows the nourishing love in our home. Her sacrifices to bring our ten wonderful children into this world, teaching and training them, while always supporting me without a murmur through my responsibilities in the Church and in my profession, are monumental. We love our eight sons-in-law as our own, and are so grateful for an unbroken string of temple marriages that link them to us forever. For those precious grandchildren already here and for the others on the way, we express our gratitude.

While nominally I come to you from the science of surgery and its mother of medicine, in a truer sense, I have been forged from the stern discipline of law—not the laws of men, as mastered by our brethren of the legal profession, but the eternal and unchanging laws of our Divine Creator. The surgeon soon learns the incontrovertibility of divine law. He knows that hopes and wishes are sometimes simply powerless sham. Desired blessings come only by obedience to divine law, and in no other way. My lifetime thus far has been focused on learning those laws. Only as the laws are known, and then obeyed, can the blessings we desire be earned. To this extent, there will be little difference for me in the activities of the past and those of the future. The endless laws of the Lord are the doctrines taught by His Apostles.

I am grateful to my colleagues in the medical profession who have permitted me to assist them in the care of their patients and who now will continue to care for them. Their selfless service satisfies the two great commandments, to love God and neighbor. They have learned that one loves best who serves best.

I sustain the servants of the Lord, from Isaiah and Ezekiel to Joseph Smith and Spencer W. Kimball. For their confidence in me, I am deeply appreciative. To them and all mankind they serve, I express my love.

Now, I understand fully that the call to the Holy Apostleship is one of witness to the world of the divinity of the Lord Jesus the Christ. I know that salvation is centered in Him! Said the angel to King Benjamin: “Salvation was, and is, and is to come, in and through the atoning blood of Christ, the Lord Omnipotent.” (Mosiah 3:18.)

The Prophet Joseph Smith was asked, “What are the fundamental principles of your religion?” He answered: “The fundamental principles of our religion are the testimony of the Apostles and Prophets, concerning Jesus Christ, that He died, was buried, and rose again the third day, and ascended into heaven; and all other things which pertain to our religion are only appendages to it.” (History of the Church, 3:30.)

The creation of the earth and all that dwell therein, the necessary fall that enabled man to be, and the atonement of the Lord are three fundamental components of God’s eternal plan. Without that atoning sacrifice, there would be neither immortality nor eternal life.

His message is the restored gospel of Jesus Christ; and it is administered by the Church that bears His name. He directs the affairs of His church by the power of the priesthood and by revelation through prophets who proclaim His doctrine to all people of the world.

The testimony that I bear is but an echo of the resounding testimonies of the eighty-four who previously have received this call to the Twelve since the spring of 1820. I know that God the Father and the Son appeared to usher in this last great dispensation of eternal truth. I testify that we are of the house of Israel, specifically of the lineage of Joseph, bearing the birthright and charged with the irrevocable responsibility to prepare the world for the second coming of the Savior. Then, numberless multitudes among every nation, kindred, tongue, and people will eventually join in proclaiming that Jesus is the Christ, the Son of the living God. So I solemnly testify in the name of Jesus Christ, amen.

\newpage
\settalk{Protect the Spiritual Power Line}
\setspeaker{Russell M. Nelson}
\settalkdate{October 1984}
\talkheading{Protect the Spiritual Power Line}{Russell M. Nelson}

One day while working around our home, trimming the hedges and vines, I had an interesting experience. I was at work with my electric clippers and long extension cord. I had done this often, each time reminding myself of the need to use these clippers with great care in order to avoid cutting things that I shouldn’t.

Suddenly the blades became jammed. Caught between them was the power cord itself. Because I had not seen it in the thicket I was trimming, I had cut into the very line that was providing the power to work.

“Isn’t that one of life’s great lessons?” I thought. “Power, if misused, can cut into the very source of that power.”

Just as the careless use of electrical power can sever the source of that power, so is it possible to misuse spiritual power to sever our spiritual power line. We would then lose that which enables us to generate success in our lives. Proper use of our spiritual power line allows us to learn, to labor, to be obedient to law, and to love. While these capabilities lead to fulfillment, at the same time they also carry risk.

\intalkheader{Power to Learn}

Consider the power to learn. How essential it is for progress, whether one’s career is professional or that of a salesman, a farmer, or a homemaker.

But learning can be misused! A sharp mind, misdirected, can cut into that line of spiritual power. Some “learned” souls delight in leading others astray, all in the so-called name of learning. Years later their victims may realize that they have climbed their ladder of learning, only to find it leaning against the wrong wall. A prophet of the Lord has counseled us on this point:

“O that cunning plan of the evil one! O the vainness, and the frailties, and the foolishness of men! When they are learned they think they are wise, and they hearken not unto the counsel of God, for they set it aside, supposing they know of themselves, wherefore, their wisdom is foolishness and it profiteth them not. And they shall perish.

“But to be learned is good if they hearken unto the counsels of God.” (2 Ne. 9:28–29; italics added.)

What happens without spiritual learning? What happens to the pilot of a glider when he is cut loose from the power of his tow plane? There may be banks and turns, but ultimately there is only one direction he can go, and that is down!

We must gain learning, but we must apply it wisely. Otherwise, we have politics without principle, industry without morality, knowledge without wisdom, science without humanity!

\intalkheader{Power to Labor}

Consider the power of labor. Labor well spent increases our capacity to do. President Heber J. Grant often said, “‘That which we persist in doing becomes easier for us to do; not that the nature of the thing itself is changed, but that our power to do is increased.’” (In Conference Report, Apr. 1901, p. 63.)

The Lord, through his prophet Lehi, said, “It must needs be, that there is an opposition in all things.” (2 Ne. 2:11.) So, in reality, competition forces us to improve. It deserves our praise and our gratitude. Without it we could not reach the heights that are otherwise ours to achieve.

But our labor can be misdirected. One can be anxiously engaged in a meaningless cause. Or one can do something wrong a hundred times and call it experience! Some would avoid labor while pursuing the goals of wealth without working for it, or an honored position without preparing for it. The converted Lamanites taught this lesson: “Rather than shed the blood of their brethren they would give up their own lives; and rather than take away from a brother they would give unto him; and rather than spend their days in idleness they would labor abundantly with their hands.” (Alma 24:18.)

\intalkheader{Power of Obedience to Law}

Next, consider the power of obedience to law. One of the great applications of spiritual power is obeying the laws of both God and man. Freedom to act and mastery of our actions both emanate from law. “When we obtain any blessing from God, it is by obedience to that law upon which it is predicated.” (D\&C 130:21.)

I relearned that so well from President Kimball. On one occasion when he needed an operation that I was to perform, he first asked me for a priesthood blessing. Following that, he said, “Now you may proceed to do that which must be done in order to make that blessing possible.”

He knew, and I knew, that not even for God’s prophet can anyone be exempted from law. Not even for God’s Son could divine law be broken!

\intalkheader{Power to Love}

Consider the power to love. I remember a mother I met once as I made a professional house call. This woman was confined in an iron lung. The ravages of polio had effectively destroyed all the breathing muscles so that her life was completely dependent upon this large metal tank and the electrical motor that powered its noisy bellows.

While there, I watched her three children as they related to their mother. The oldest interrupted our work to ask permission to go to a friend’s house for an hour. Later the second child asked her mother for help with arithmetic. Finally the youngest child, so small that she couldn’t see her mother’s face directly, looked up at her image in a mirror that had been placed over the mother’s head and asked, “Mommy, may I have a cookie?” I’ve never forgotten that lesson on the power of love. This woman, virtually disabled and certainly incapable of any degree of physical enforcement of parental authority, sweetly influenced that home solely with the power to love!

The power of love between a man and a woman is special. The love shared with my beloved companion, Dantzel, has increased the power of love for both of us. That love brought us to the altar in the temple of the Lord. Her love for me motivated her to teach school during the early years of our marriage. When things were tight, she held a second job at night. Once when things were exceptionally tight, she even sold her blood in between her two jobs to keep us solvent. (Her dear parents may have wondered what kind of a son-in-law they had on that occasion!) I thought of that many years later when she needed a transfusion urgently and her blood couldn’t be matched readily with donor blood from the blood bank. What a privilege it was for me to donate mine directly to her.

Can the power to love be misused? Sadly, yes. The illegitimate union of the sexes has, in my observation, been one of the greatest causes of grief. In some instances, the improper use of this sacred physical endowment has even destroyed its use in later years.

Misused, the power of love can cut off spiritual power. The abuse of the power to love can result in no love at all. Only its cheap facsimiles of lewdness and lust remain in the wake of pleasure without conscience. Instead of feasting at the banquet table of bounteous love with his own posterity, one is left with scraps from the table—only the refuse from what might have been.

The ashes of burned “love” smolder with the smoke of sadness. Yet the embers of evil still burn. But as loud voices argue on, let us remember that those who advocate abortion have already been born! Those who freely deny God with their amoral and agnostic practices will one day find that He may just as freely deny them! (See 3 Ne. 28:34.)

\intalkheader{Source of Spiritual Power}

The source of our spiritual power is the Lord! The ultimate source of spiritual power is God our Father. The messenger of this power is the Holy Ghost. This power differs from electrical power. An electrical appliance consumes power. The use of His spiritual power replenishes our power. While electrical power can be used only for measured periods of time, spiritual power can be used for time and eternity!

Our spiritual power line is strengthened through prayer. As we counsel with God in all our doings, he will direct us for good. (See Alma 37:37.) Prayer is available whenever we ask for it!

But the Lord places the initiative upon us. He expects us to reach for his power, just as we must insert the plug in the outlet for electricity. He said, “If thou shalt ask, thou shalt receive revelation upon revelation, knowledge upon knowledge, that thou mayest know the mysteries and peaceable things—that which bringeth joy, that which bringeth life eternal.” (D\&C 42:61; italics added.) Personal worthiness and scripture study enable us to do more with this power.

Rewards result from the righteous use of the spiritual power belonging to the priesthood! And they are so great that they are almost beyond human comprehension. To those couples who bear and share that priesthood worthily and remain faithful to the law of the everlasting covenant of eternal marriage, enduring the congested years and trials of diapers and dishes, crowded kitchen and thin pocketbook, service in the Church, education and the burning of the midnight oil, the Lord makes this promise: “Ye shall come forth in the first resurrection; … and shall inherit thrones, kingdoms, principalities, and powers, dominions, … [and there] shall be a fulness and a continuation of the seeds forever and ever.” (D\&C 132:19; italics added.)

The spiritual power line conveys revelation. Revelation is given to the President for the Church, and to leaders and teachers in their respective callings. Personal revelation is provided for the glorification of individual members and families in the Church. These lines are well insulated, and they are not scrambled! Our Father is a God of order! No one else receives revelation to govern the Church; that will be delivered only to the President of the Church. The father in the home next door will not receive revelation for his neighbor’s family.

It is possible to disregard or even misuse spiritual power. Some have misused the power of prayer by making that sacred communication trivial. Some well-meaning Saints even do the right things for the wrong reasons, if they narrowly center on the percentages they report rather than on the precious people they serve.

Like cutting the cord with clippers, it is possible to use spiritual power so carelessly as to destroy one’s very connection to that power. I know a husband who dominates his wife as though she were his possession. He seems to regard her about as he does his automobile or his suitcase, which he uses for his own purposes. And I know a wife who dominates her husband to the point that he has lost all feelings of worth.

Remember, “The rights of the priesthood are inseparably connected with the powers of heaven, and … the powers of heaven cannot be controlled nor handled only upon the principles of righteousness.” (D\&C 121:36; italics added.)

The unrighteous use of priesthood authority surely severs the connection to the Source of that authority. (See D\&C 121:37.)

Worship strengthens our power line to deity. There can be no true worship without sacrifice, and there can be no true sacrifice without a cause. The cause that earns our love and priority is the cause of Jesus Christ.

Speaking of his own atonement, the Savior proclaimed, “To this end was I born, and for this cause came I into the world.” (John 18:37.) His example of worship, sacrifice, and commitment to cause becomes ours! His is the ultimate source of all our power to do good.

As we carefully and prayerfully protect the spiritual power line that links us to the Savior, we become more like him!

Of his sacred name and mission I solemnly bear special witness. Of his divine sonship I proclaim sacred testimony. That this is his church I fervently declare, in the name of Jesus Christ, amen.

\newpage
\settalk{Reverence for Life}
\setspeaker{Russell M. Nelson}
\settalkdate{April 1985}
\talkheading{Reverence for Life}{Russell M. Nelson}

Unitedly we thank the Almighty for the wondrous prolongation of the life of Elder Bruce R. McConkie, enabling him to preach that powerful sermon. Our gratitude is profound!

I pray for the Spirit of the Lord to help me communicate his mind and will on a very vital and sensitive subject. I apologize for the use of words repugnant to me and ill-suited to this hallowed pulpit. I do so only for clarity of communication regarding reverence for human life.

As sons and daughters of God, we cherish life as a gift from him.

A heavy toll on life is included among the evils of war. Data from all nations are appalling. For the United States of America, one hundred thousand were killed in World War I; over four hundred thousand died in World War II. In the first two hundred years as a nation, the lives of over one million Americans were lost due to war.

Regrettable as is the loss of loved ones from war, these figures are dwarfed by the toll of a new war that annually claims more casualties than the total number of fatalities from all the wars of this nation.

It is a war on the defenseless—and the voiceless. It is a war on the unborn.

This war, labeled “abortion,” is of epidemic proportion and is waged globally. Over fifty-five million abortions were reported worldwide in the year 1974 alone. Sixty-four percent of the world’s population now live in countries that legally sanction this practice. In the United States of America, over 1.5 million abortions are performed annually. About 25–30 percent of all pregnancies now end in abortion. In some metropolitan areas, there are more abortions performed than live births. Comparable data also come from other nations.

Yet society professes reverence for human life. We weep for those who die, pray and work for those whose lives are in jeopardy. For years I have labored with other doctors here and abroad, struggling to prolong life. It is impossible to describe the grief a physician feels when the life of a patient is lost. Can anyone imagine how we feel when life is destroyed at its roots, as though it were a thing of naught?

What sense of inconsistency can allow people to grieve for their dead, yet be calloused to this baleful war being waged on life at the time of its silent development? What logic would encourage efforts to preserve the life of a critically ill twelve-week-old infant, but countenance the termination of another life twelve weeks after inception? More attention is seemingly focused on the fate of a life at some penitentiary’s death row than on the millions totally deprived of life’s opportunity through such odious carnage before birth.

The Lord has repeatedly declared this divine imperative: “Thou shalt not kill.” Recently he added, “Nor do anything like unto it.” (D\&C 59:6.) Even before the fulness of the gospel was restored, the enlightened understood the sanctity of life. John Calvin, the sixteenth-century reformer, wrote: “If it seems more horrible to kill a man in his own house than in a field, because a man’s house is his place of most secure refuge, it ought surely to be deemed more atrocious to destroy a fetus in the womb before it has come to light.”

But what impropriety could now legalize that which has been forbidden by the laws of God from the dawn of time? What twisted reasoning has transformed mythical concepts into contorted slogans assenting to a practice which is consummately wrong?

These slogans begin with proper concern for the health of the mother. Infrequently, instances may occur in which the continuation of pregnancy could be life-threatening to the mother. When deemed by competent medical authorities that the life of one must be terminated in order to save the life of the other, many agree that it is better to spare the mother. But these circumstances are rare, particularly where modern medical care is available.

Another sympathetic concern applies to pregnancies resulting from rape or incest. The tragedy of this despoilment is compounded because, in such relationships, freedom of choice is denied the woman who is innocently involved.

But less than 3 percent of all abortions are performed for these two reasons. The other 97 percent are performed for what may be termed “reasons of convenience.”

Some argue for abortion because a malformed child may result. The harmful effects of certain infectious or toxic agents in the first trimester of pregnancy are real.

The experience of a couple whom I shall identify as Brother and Sister Brown (fictitious names) is instructive. Sister Brown was only twenty-one years old at the time, a beautiful woman and a devoted wife. In her first trimester, she contracted the dreaded German measles.

Abortion was advised because the developing baby would almost surely be damaged. Some members of her family, out of loving concern, applied additional pressure for an abortion. “Don’t burden yourself financially with a handicapped child,” they argued. “You are too young and too poor.”

Devotedly Brother and Sister Brown consulted their bishop. He referred them to their stake president, who listened to their serious concern and counseled them not to terminate the life of this baby, even though the child might have a problem. He quoted this scripture:

“Trust in the Lord with all thine heart; and lean not unto thine own understanding. In all thy ways acknowledge him, and he shall direct thy paths.” (Prov. 3:5–6.)

They chose to follow that counsel and permit their child to be born—a beautiful little girl, normal in every respect, except for a hearing loss, that became evident later. After an evaluation at a school for the deaf, Brother and Sister Brown were advised that this child had the intellect of a genius. Now, some twenty years later, she attends a major university on a scholarship.

When recently asked how they felt about their once-weighty decision, the mother quickly responded, “She is one of the great joys of my life! She is such a choice spirit! Though she lost the sense of hearing, she has compensated with augmented ability otherwise. Her eyes are alive with constant attention. She excels in dancing, even though she perceives the sounds of music from vibrations. She has served as an officer in school. But most significant is her guileless spirit, her unconditional love. She has taught us to serve and to share. Her spiritual insights have helped us to know God and his purposes. My husband and I are so grateful that she is one of our children.”

Consider another individual weighing the consequences of her pregnancy. She was beyond the normal age for bearing children. She announced to her doctor that her husband was an alcoholic with a syphilitic infection. One of her children had been born dead. Another child was blind. Another had tuberculosis. Her family had a history of deafness. Finally she confessed that she was living in abject poverty. If this true historical situation were posed today, many would recommend abortion. The child born from that pregnancy became the renowned composer Ludwig van Beethoven.

But the principle involved extends beyond those who may become great. If one is to be deprived of life because of potential for developing physical problems, consistency would dictate that those who already have such deficiencies should likewise be terminated. Continuing, then, those who are either infirm, incompetent, or inconvenient should be eliminated by those in power. Such irreverence for life is unthinkable!

Another contention raised is that a woman is free to choose what she does with her own body. To a certain extent this is true for all of us. We are free to think. We are free to plan. And then we are free to do. But once an action has been taken, we are never free from its consequences. Those considering abortion have already exercised certain choices.

To clarify this concept, we can learn from the astronaut. Any time during the selection process, planning, and preparation, he is free to withdraw. But once the powerful rocket fuel is ignited, he is no longer free to choose. Now he is bound by the consequences of his choice. Even if difficulties develop and he might wish otherwise, the choice made was sealed by action.

So it is with those who would tamper with the God-given power of procreation. They are free to think and plan otherwise, but their choice is sealed by action.

The woman’s choice for her own body does not validate choice for the body of another. The expression “terminate the pregnancy” applies literally only to the woman. The consequence of terminating the fetus therein involves the body and very life of another. These two individuals have separate brains, separate hearts, and separate circulatory systems. To pretend that there is no child and no life there is to deny reality.

It is not a question of when “meaningful life” begins or when the spirit “quickens” the body. In the biological sciences, it is known that life begins when two germ cells unite to become one cell, bringing together twenty-three chromosomes from both the father and from the mother. These chromosomes contain thousands of genes. In a marvelous process involving a combination of genetic coding by which all the basic human characteristics of the unborn person are established, a new DNA complex is formed. A continuum of growth results in a new human being. The onset of life is not a debatable issue, but a fact of science.

Approximately twenty-two days after the two cells have united, a little heart begins to beat. At twenty-six days the circulation of blood begins.

Scripture declares that the “life of the flesh is in the blood.” (Lev. 17:11.) Abortion sheds that innocent blood.

Another excuse some use to justify abortion relates to population control. Many in developing nations unknowingly ascribe their lack of prosperity to overpopulation. While they grovel in ignorance of God and his commandments, they may worship objects of their own creation (or nothing at all), while unsuccessfully attempting to limit their population by the rampant practice of abortion. They live in squalor, oblivious to the divine teaching—stated in the scriptures not once, but thirty-four times—that people will prosper in the land only if they obey the commandments of God.

How can God fulfill his promise to prosper his children in obedience if they worship idols or destroy life created by him—destined to be in his very image?

They will prosper only when their education includes faith in and obedience to the God of this world, who said,

“I, the Lord, … built the earth, my very handiwork; and all things therein are mine. And it is my purpose to provide. … But it must needs be done in mine own way. … For the earth is full, and there is enough and to spare.” (D\&C 104:14–17; italics added.)

Now, as a servant of the Lord, I dutifully warn those who advocate and practice abortion that they incur the wrath of Almighty God, who declared, “If men … hurt a woman with child, so that her fruit depart from her, … he shall be surely punished.” (Ex. 21:22.)

Of those who shed innocent blood, a prophet declared: “The judgments which [God] shall exercise … in his wrath [shall] be just; and the blood of the innocent shall stand as a witness against them, yea, and cry mightily against them at the last day.” (Alma 14:11.)

The Church of Jesus Christ of Latter-day Saints has consistently opposed the practice of abortion. One hundred years ago the First Presidency wrote: “And we again take this opportunity of warning the Latter-day Saints against those … practices of foeticide and infanticide.”

Early in his presidency, our beloved President Spencer W. Kimball said, “We decry abortions and ask our people to refrain from this serious transgression.”

Why destroy a life that could bring such joy to others?

Now, is there hope for those who have so sinned without full understanding, who now suffer heartbreak? Yes. So far as is known, the Lord does not regard this transgression as murder. And “as far as has been revealed, a person may repent and be forgiven for the sin of abortion.” Gratefully, we know the Lord will help all who are truly repentant.

Yes, life is precious! No one can cuddle a cherished newborn baby, look into those beautiful eyes, feel the little fingers, and caress that miraculous creation without deepening reverence for life and for our Creator.

Life comes from life. It is a gift from our Heavenly Father. It is eternal, as he is eternal. Innocent life is not sent by him to be destroyed! This doctrine is not of me, but is that of the living God and of his divine Son, which I testify in the name of Jesus Christ, amen.

\newpage
\settalk{Self-Mastery}
\setspeaker{Russell M. Nelson}
\settalkdate{October 1985}
\talkheading{Self-Mastery}{Russell M. Nelson}

If our faith be united in prayer that we may be edified together, I should like to speak about our quest for self-mastery. In so doing, I would converse as a loving father counseling one of my own children.

Before you can master yourself, my precious one, you need to know who you are. You consist of two parts—your physical body, and your spirit which lives within your body. You may have heard the expression “mind over matter.” That’s what I would like to talk about—but phrase it a little differently: “spirit over body.” That is self-mastery.

When you arrived as a newborn baby, your little body was master. You had what I call the “I-want-what-I-want-when-I-want-it” philosophy. No amount of discussion could postpone your impatient demands when you wanted to be fed—and now! Like all parents, we anxiously anticipated the first smile, a word, a glimpse at the potential of the spirit within your tiny body. Is there a mother who has not cradled her baby as your sweet mother did, in wistful wonder of the destiny of her dear little one? Even Mary, the mother of Jesus, might have asked such questions:

\begin{verse}
Baby, lyin’ in a manger, slumberin’ so sweetly,\\
What you goin’ to be?\\
Baby, all the world is watchin’, all the world awaits to see.\\
What will you be?\\
Baby, sleepin’ in a stable, underneath the heaven,\\
What you goin’ to say? …\\
Baby, lyin’ in a manger,\\
Will you save the world one day?
\end{verse}

(Natalie Sleeth, “Baby, What You Goin’ to Be?” New York: Carl Fisher, Inc.)

Through those early years, we parents are properly concerned with physical needs of our children, such as food, clothing, and shelter.

But as you grow older, our concerns shift more toward your spiritual growth, in order that you might achieve your full potential. “For the natural man is an enemy to God, and has been from the fall of Adam, and will be, forever and ever, unless he yields to the enticings of the Holy Spirit … and becometh a saint.” (Mosiah 3:19.)

That requires self-mastery. Remember, “The spirit and the body are the soul of man.” (D\&C 88:15.) Both are of great importance. Your physical body is a magnificent creation of God. It is his temple as well as yours, and must be treated with reverence. Scripture declares: “Ye are the temple of God. … If any man defile [it], him shall God destroy; for the temple of God is holy, which temple ye are.” (1 Cor. 3:16–17.)

Remarkable as your body is, its prime purpose is even of greater importance—to serve as tenement for your spirit. Abraham taught that “these … spirits … existed before, they shall have no end … for they are … eternal.” (Abr. 3:18.)

Your spirit acquired a body at birth and became a soul to live in mortality through periods of trial and testing. Part of each test is to determine if your body can become mastered by the spirit that dwells within it.

Although your spirit had a veil of forgetfulness placed over it at the time of your birth into mortality, it retained its power to remember all that happens—precisely recording each event of life. Indeed, scriptures warn “that every idle word that men shall speak, they shall give account thereof in the day of judgment.” (Matt. 12:36.) Prophets refer to our “bright recollection” (Alma 11:43) and “perfect remembrance” (Alma 5:18) at that day of decision.

Since thoughts precede deeds, you must first learn to control your thoughts. “As [a man] thinketh in his heart, so is he.” (Prov. 23:7.)

In your quest for self-mastery, full participation in the activities of the Church will help. I’ll mention but a few. A first step comes as we learn together to keep the Sabbath day holy. This is one of the Ten Commandments. (See Ex. 20:8; Deut. 5:15.) We honor the Sabbath “to pay [our] devotions unto the Most High” (D\&C 59:10), and because the Lord declared: “It is a sign between me and you … that ye may know that I am the Lord that doth sanctify you.” (Ex. 31:13; see also Ezek. 20:20.)

Another step toward self-mastery comes when you are old enough to observe the law of the fast. As funds are contributed from meals missed, the needs of the poor may be met. But meanwhile, through your spirit, you develop personal power over your body’s drives of hunger and thirst. Fasting gives you confidence to know that your spirit can master appetite.

Some time ago your mother and I visited a third-world country where sanitary conditions were much poorer than ours. We joined with a delegation of other doctors from all over the world. The president of our group, an experienced traveler, warned of risks. In order to avoid water that might be contaminated, we were even counseled to brush our teeth with an alcoholic beverage. We chose not to follow that counsel, but simply did what we had learned to do once a month. We fasted that first day, thinking we could introduce simple food and fluids gradually thereafter. Later, we were the only ones in our group without disabling illness.

Fasting fortifies discipline over appetite and helps to protect against later uncontrolled cravings and gnawing habits.

Another step toward self-mastery comes from obedience to the Word of Wisdom. Remember, it contains a “promise, adapted to the capacity of … the weakest of all saints.” (D\&C 89:3.) It was given “in consequence of evils and designs which do and will exist in the hearts of conspiring men in the last days.” (D\&C 89:4.) Indeed, as you develop courage to say no to alcohol, tobacco, and other stimulants, you gain additional strength. You can then refuse conspiring men—those seditious solicitors of harmful substances or smut. You can reject their evil enticements to your body.

If you yield to anything that can addict, and thus defy the Word of Wisdom, your spirit surrenders to the body. The flesh then enslaves the spirit. This is contrary to the purpose of your mortal existence. And in the process of such addiction, your life span is likely to be shortened, thereby reducing the time available for repentance by which your spirit might attain self-mastery over your body.

Other physical appeals come during your courtship period. In your youth, you may be challenged by restraints of parents hoping to guide you through this wonderful period of life.

Because the adversary is keenly aware of the power of physical temptation, Alma instructed his son and all of us: “See that ye bridle all your passions.” (Alma 38:12.)

When you marry, you and your eternal companion may then invoke the power of procreation, that you may have joy and rejoicing in your posterity. This divine endowment is guarded by your Creator’s law of chastity. All through the years, remember: chastity is the powerful protector of virile manhood and the crown of beautiful womanhood.

In courtship and marriage, virtue seems to come under attack first. Mental turmoil that trails in the wake of weakness from lust has evoked many a tear from innocent loved ones. Without repentance, tumult within self does not quit either.

Shakespeare expressed such self-conflict as one of his characters contemplating conquest in lust spoke these lines:

\begin{verse}
What win I, if I gain the thing I seek?\\
A dream, a breath, a froth of fleeting joy.\\
Who buys a minute’s mirth to wail a week?\\
Or sells eternity to get a toy?\\
For one sweet grape who will the vine destroy?
\end{verse}

(“Lucrece,” lines 211–15.)

Prophets have repeatedly cautioned about moral sin. One, for example, warned: “O, my beloved brethren, remember the awfulness in transgressing against that Holy God, and also the awfulness of yielding to the enticings of that cunning one. Remember, to be carnally-minded is death, and to be spiritually-minded is life eternal.” (2 Ne. 9:39; see also Rom. 8:6; Alma 36:4; D\&C 29:35; 67:10.)

Now don’t misunderstand me. I would not want you to neglect your body. It deserves daily care. Physical conditioning through regular exercise requires self-mastery too. I marvel at Elder Joseph Anderson, now in his ninety-sixth year. For decades, the strength of his spirit over his body has induced him to swim regularly. But his motivation has never been to attain physical longevity. That has come only incidentally. His desire has been to serve God and His anointed. Elder Anderson has followed what I label as the Lord’s prescription for a long and useful life. Those faithful in “magnifying their calling, are sanctified by the Spirit unto the renewing of their bodies. They become … the elect of God.” (D\&C 84:33–34.)

Elder Anderson’s exercise program agrees with the perspective of Paul, who said: “Bodily exercise profiteth little: but godliness is profitable unto all things, having promise of the life that now is, and of that which is to come.” (1 Tim. 4:8.)

Handsome and fit, Elder Anderson personifies this scripture: “Glorify God in your body, and in your spirit, which are God’s.” (1 Cor. 6:20.)

As you work during the productive years of life, whether at home or in the field, in the factory or at a workbench, reputation is built and character is forged as you develop self-mastery. Faithful payment of tithing is part of that process. It defends you against dishonesty or shabby temptations. Courageous accountability for your own actions becomes a cherished prize.

It really matters what you listen to, what you look at, what you think, say, and do. Select music that will strengthen your spirit. Control your speech; keep it free from profanity and vulgarity. Follow the teachings of this proverb: “My mouth shall speak truth; and wickedness is an abomination to my lips.

“All the words of my mouth are in righteousness; there is nothing … perverse in them.” (Prov. 8:7–8.)

As you approach old age, you will face new challenges to self-mastery. Symptoms of the deteriorating body can be painful, even disabling. Deep aches of sadness are caused by the departing of loved ones. For some, these deepening trials come early in life. But when yours are thrust upon you, remember a concept expressed by my father some time after my mother had passed away. Your grandparents had been married for sixty-four years. When someone asked how he was doing, my father simply stated, “I’m lonely, but I’m not lonesome.” Do you know what he meant? Though he was now without his sweetheart, he was so busy assisting family and friends, he had replaced sorrow with service and had displaced self-pity with selfless love. He had found joy in following the timeless example of the Master.

Jesus, our Savior, was born in the lowliest of circumstances. For his baptism he was immersed in the lowest body of fresh water upon the planet. In service and suffering, he also “descended below” all things (D\&C 122:8), that he could rise above all things. Near the end of his life, he triumphantly declared, “I have overcome the world.” (John 16:33.) “Look unto me, and endure to the end, and ye shall live; for unto him that endureth to the end will I give eternal life.” (3 Ne. 15:9.) Scriptures tutor us at least twenty-six times* to endure to the end to attain eternal life. Then we will obtain a resurrected body—one that is incorruptible, glorified, and prepared to live in the presence of God.

To reach your highest destiny, emulate the Savior. He proclaimed, “What manner of men ought ye to be? … Even as I am.” (3 Ne. 27:27.) Our loftiest hope is to grow in spirit and attain “the stature of the fulness of Christ: That we henceforth be no more children.” (Eph. 4:13–14.)

You will then be well prepared for that pending day of judgment when, as taught by President Spencer W. Kimball, “the soul, composed of the resurrected body and the eternal spirit, … will come before the great judge to receive its final assignment for the eternity” (The Teachings of Spencer W. Kimball, ed. Edward L. Kimball, Salt Lake City: Bookcraft, 1982, p. 46.)

Remember, my dear one, not an age in life passes without temptation, trial, or torment experienced through your physical body. But as you prayerfully develop self-mastery, desires of the flesh may be subdued. And when that has been achieved, you may have the strength to submit to your Heavenly Father, as did Jesus, who said, “Not my will, but thine, be done.” (Luke 22:42.)

When deepening trials come your way, remember this glorious promise of the Savior: “To him that overcometh will I grant to sit with me in my throne, even as I also overcame, and am set down with my Father in his throne.” (Rev. 3:21.)

Christ is our great Exemplar. I declare, as a special witness, that he is the Son of God and “is the life and the light of the world.” (Alma 38:9; see also D\&C 11:28.) We develop self-mastery as we become like him, I testify in the name of Jesus Christ, amen.

\newpage
\settalk{In the Lord’s Own Way}
\setspeaker{Russell M. Nelson}
\settalkdate{April 1986}
\talkheading{In the Lord’s Own Way}{Russell M. Nelson}

Since my previous medical work took Sister Nelson and me to a number of developing nations, we have been exposed to many challenging scenes. In one country, so many people were sleeping in the streets and on sidewalks that we literally had to step over them as we walked. In another nation, our compassion was stretched almost to the breaking point as we yearned to help countless people in need. Young mothers with babies bundled on their backs begged for money while paddling their little sampan boats, which served both as their shelter and as their mode of transportation. And, oh, how our hearts ached for young men and women of another country who, one by one, were strapped as beasts of burden to wooden-wheeled carts heavily laden with weighty cargo. As far as our eyes could see, the endless caravan of vehicles continued, pulled by dint of human toil.

Although reasons vary according to time and place, the poor and the needy have nearly always been present. Regardless of cause, our Heavenly Father is concerned for them. They are all his children. He loves and cares for them.

Lessons from the Old Testament remind us that when the Lord sent prophets to call Israel back from apostasy, in almost every instance, one of the first charges made was that the poor had been neglected.

Scriptures teach us that the poor—especially widows, orphans, and strangers—have long been the concern of God and the godly. The poor have been especially favored by the law. Old Testament teachings authorized poor persons at harvest time to glean after the reapers. At fruit-picking time, what was left hanging on branches belonged to the poor. In the sabbatical seventh year, and in the jubilee fiftieth year, land was not planted nor tilled, and what grew of itself was free for the hungry.

To those who cared for the poor, blessings were promised. The Lord would deliver them in time of trouble. (See Ps. 41:1.) Truths were taught by these proverbs: “He that hath mercy on the poor, happy is he.” (Prov. 14:21.) “The righteous considereth the cause of the poor: but the wicked regardeth not to know it.” (Prov. 29:7.)

During the Savior’s earthly ministry, he reemphasized his timeless concern for the poor. Remember the reply the Lord gave to the question of the rich man: “If thou wilt be perfect, go and sell [all] that thou hast, and give to the poor, and thou shalt have treasure in heaven: and come and follow me.” (Matt. 19:21; see also Luke 18:22.)

In one of the Master’s precious parables, he illustrated this doctrine with the story of one who was hungry and was given meat, another who was thirsty and was given drink, and a stranger who was welcomed. The Lord related those as favors to him when he taught, “Verily I say unto you, Inasmuch as ye have done it unto one of the least of these my brethren, ye have done it unto me.” (Matt. 25:40.) And when they were not ministered unto, he admonished, “Verily I say unto you, Inasmuch as ye did it not to one of the least of these, ye did it not to me.” (Matt. 25:45; italics added.)

Indeed, the Church in New Testament times also had a binding obligation to care for the poor. (See also Luke 14:12–14; 16:19–31; Acts 2:44–45; 4:32; 6:1–6; 11:27–30; Rom. 15:26; Gal. 2:10; James 2:1–9.)

The Book of Mormon repeatedly declares this doctrine. From it we learn that care of the poor is an obligation that we take upon ourselves at the time of baptism. The prophet Alma so taught:

“Ye are desirous to come into the fold of God, and to be called his people, and are willing to bear one another’s burdens, that they may be light;

“Yea, and [ye] are willing to mourn with those that mourn; yea, and comfort those that stand in need of comfort, and to stand as witnesses of God at all times and in all things, … even until death, …

“Serve him and keep his commandments.” (Mosiah 18:8–10.)

Every person so baptized who receives the gift of the Holy Ghost, which seals the ordinance, is under solemn covenant with the Lord to obey his commandments. Caring for the poor is one of those commandments. Surely, in Book of Mormon times, members of the Church had a sacred obligation to care for the poor. (See 2 Ne. 9:30; Mosiah 4:16–27; Alma 1:29–31; 32:2–6; 34:28; Hel. 4:11–12; Morm. 8:35–39.)

Few, if any, of the Lord’s instructions are stated more often, or given greater emphasis, than the commandment to care for the poor and the needy. Our dispensation is no exception.

In December 1830, the very year in which The Church of Jesus Christ of Latter-day Saints was organized, the Lord declared that “the poor and the meek shall have the gospel preached unto them, and they shall be looking forth for the time of my coming, for it is nigh at hand.” (D\&C 35:15.)

Bishops were designated and their duties defined: “They shall look to the poor and the needy, and administer to their relief that they shall not suffer.” (D\&C 38:35.)

In 1831, the Lord said: “Remember the poor. … Inasmuch as ye impart of your substance unto the poor, ye will do it unto me.” (D\&C 42:30–31.) A little later, he again declared, “Visit the poor and the needy and administer to their relief.” (D\&C 44:6.) Later the same year, he warned: “Wo unto you rich men, that will not give your substance to the poor, for your riches will canker your souls.” (D\&C 56:16.)

With these teachings throbbing in our ears, stated and restated in accounts to all people in all days of recorded scripture, let our thoughts return to the homeless, beggars in boats, human beasts of burden, and to multitudes stricken with poverty.

Is it possible to be faithful to our solemn obligation to care for the poor and needy, to lift them and to love them—worldwide?

Where shall we begin?

When?

How?

Hear the answer of Almighty God:

“I, the Lord, stretched out the heavens, and built the earth, my very handiwork; and all things therein are mine.

“And it is my purpose to provide for my saints, for all things are mine.

“But it must needs be done in mine own way; and behold this is the way that I, the Lord, have decreed to provide for my saints, that the poor shall be exalted, in that the rich are made low.

“For the earth is full, and there is enough and to spare; yea, I prepared all things, and have given unto the children of men to be agents unto themselves.

“Therefore, if any man shall take of the abundance which I have made, and impart not his portion, according to the law of my gospel, unto the poor and the needy, he shall, with the wicked, lift up his eyes in hell, being in torment.” (D\&C 104:14–18.)

I repeat the Lord’s prescription: “But it must needs be done in mine own way”! We begin where we are, now, and work according to his plan. His “own way” includes these principles:

“Women have claim on their husbands for their maintenance. … [And] all children have claim upon their parents. … And after that, they have claim upon the church, or … upon the Lord’s storehouse, if their parents have not. … And the storehouse shall be kept by the consecrations of the church; and widows and orphans shall be provided for, as also the poor.” (D\&C 83:2, 4–6.)

An important part of the Lord’s storehouse is maintained as a year’s supply, stored, where possible, in the homes of faithful families of the Church.

Now, some may ask, “What about those who are poor because they are idle and unwilling to work?” They should heed these words of warning:

“Thou shalt not be idle; for he that is idle shall not eat the bread nor wear the garments of the laborer.” (D\&C 42:42.)

“Wo unto you poor men … who will not labor with your own hands!” (D\&C 56:17.)

Judgment of worthiness is made by the bishop, and ultimately by the Lord, as taught by Nephi: “With righteousness shall the Lord God judge the poor, and reprove with equity for the meek of the earth.” (2 Ne. 30:9.)

Ours is not to judge; ours is a covenantal obligation to care for the poor and the needy, to prepare for their rejoicing when the Messiah shall come again. (See D\&C 56:18–19.)

The Lord’s “own way” includes, first, reliance on self, then on the family. As parents care for their children, they, in turn, may reciprocate when parents become less able. Family pride promotes solicitude for each member, taking priority over other assistance.

If one’s family can’t help, the Lord’s “own way” includes the Church organization. The bishop is assisted by priesthood quorums and good sisters of the Relief Society, organized to look “to the wants of the poor, searching after objects of charity and … administering to their wants.” (Handbook of the Relief Society, 1931, p. 22.)

Members of priesthood quorums and groups have a duty to rehabilitate, spiritually and temporally, their erring or unfortunate brethren. While a bishop extends aid to one temporarily out of work, the quorum arranges for his employment until fully self-supporting again.

As individual members of the Church, you and I participate in the Lord’s “own way.” At least once a month, we fast and pray and contribute generous offerings to funds that enable bishops to disperse aid. This is part of the law of the gospel. Each of us truly can help the poor and the needy, now, and wherever they are. And we, too, will be blessed and protected from apostasy by so doing.

Limitations do exist. Measures of relief are at best temporary. Storehouses can provide only for some temporal needs. All people cannot be brought to the same living standards. And all needed things cannot be achieved by goods or gold.

To care fully for the poor, we must help the poor to change. As they are taught and abide doctrines of Deity, spiritual strength will come that enlightens the mind and liberates the soul from the yoke of bondage. When people of the earth accept the gospel of Christ, their attitudes change. Their understanding and capabilities increase.

A poet sensed the great power of the Spirit of the Lord to lift an individual when he wrote:

\begin{verse}
The chief of all thy wondrous works,\\
Supreme of all thy plan;\\
Thou hast put an upward reach\\
Within the heart of man.
\end{verse}

(Author unknown)

That upward reach, drawn from a knowledge of divine doctrines, transforms souls! May I share an illustration with you? Once, Sister Nelson and I were invited to the humble home of Polynesian Saints who had relatively recently joined the Church. By walking carefully on wooden planks, we approached their house, built on wooden piles emerging from the floor of the sea. We climbed a ladder to enter their little one-room dwelling. As we were invited to be seated on freshly woven grass mats, we could peek through holes in the floor and view sea water below. That home was starkly devoid of furniture, except for a used sewing machine provided by sisters of the Relief Society. But the love and warmth of this special family were apparent as our visit continued.

“We would like to sing for you,” the father said through an interpreter. He put one arm about his wife and the other about the children, as did his wife. Five little ones, dressed in newly sewn clothing, joined their parents in singing songs the father had composed.

Concluding, he said: “These songs express our feelings of deep gratitude. Before we joined the Church, we had so little. Now we have so much!”

While wiping tears from our moistened cheeks, Sister Nelson and I looked at each other, comprehending that the gospel brings spiritual wealth which may bear little relation, at first, to tangible abundance. Conversely, people with plenty can be spiritually poor. Yet the Lord is concerned for them all!

Missionary work throughout the world is part of his plan. It brings the light of the gospel to those who embrace the truth. Then, as Saints learn and obey the commandments of God, they will prosper. This promise has been recorded by prophets throughout time and in diverse places. (See Josh. 1:7; 1 Kgs. 2:3; 2 Chr. 24:20; 31:21; Ezra 6:14; 1 Ne. 2:20; 4:14; 2 Ne. 1:9, 20; 4:4; Jarom 1:9; Omni 1:6; Mosiah 1:7; 2:22, 31; Alma 9:13; 36:1, 30; 37:13; 38:1; 48:15, 25.)

Working with a will, Saints gain a new appreciation of who they are and of their eternal worth. Righteousness, independence, thrift, industry, and self-reliance become personal goals. These qualities transform lives. In time, in the Lord’s “own way,” the poor will no longer be poor.

The Church of Jesus Christ has been restored to the earth. The power of God is among men. Angelic ministrants communicate to legal administrators once again. A living prophet, the everlasting priesthood, covenant people, and the Lord’s “own way” are upon the earth to bless mankind, all races, in every clime—yes, all people, young and old, wealthy and poor, now and forevermore. This I testify, in the name of Jesus Christ, amen.

\newpage
\settalk{“Joy Cometh in the Morning”}
\setspeaker{Russell M. Nelson}
\settalkdate{October 1986}
\talkheading{“Joy Cometh in the Morning”}{Russell M. Nelson}

The title of my message is taken from the thirtieth Psalm, verse 5 [Ps. 30:5]: “Joy cometh in the morning.” As I discussed this scripture with members of our family, they recalled that “men are, that they might have joy” (2 Ne. 2:25), but they had not pondered the intriguing concept that “joy cometh in the morning.”

One of our family said, “News reports appear almost daily concerning people who have problems with drugs, drinking, and emotional distress. How can they, and we, attain the joy spoken of in the scriptures?”

“The gospel of Jesus Christ offers hope,” I answered. “It declares joy to be part of our divine destiny. And to experience joy in the morning becomes our special challenge. The true test,” I continued, “is to be able to look in the mirror, first thing in the morning, and feel real joy.”

One of our daughters, who had recently announced that she was expecting a new addition to the family, said, “But Dad, that’s the hardest time of the day for me!”

“My dear ones,” I replied, “in order to experience true joy in the morning, or at any time, at least three factors are needed. You need to feel good about the people with whom you live and work—your companions in life. You must feel good about yourself—not in any sense of conceit, but simply a proper esteem for yourself, well deserved. And possibly most important, you must feel good about your relation to God and sincerely love him.”

As I so counseled my family in that conversation, we all might consider those three steps to achieve real joy in life.

\intalkheader{Courtesy to Companions}

Joy in the morning begins with courtesy to companions. When shades of slumber first admit the light of dawn, I reach gently for my beloved companion nearby. I gain sweet reassurance that all is well with her even before my eyes are fully opened. That reminds me, incidentally, of advice given by President David O. McKay, who said, “During courtship we should keep our eyes wide open, but after marriage keep them half-shut” (in Conference Report, Apr. 1956, p. 9).

My sweetheart has done that. Through our many long years of postgraduate study, professional responsibilities, and a growing family, she did not complain. Recently I overheard a conversation she had with young mothers enduring similar stress. They asked her how she had managed with ten children and a husband whose time to help was so limited. Kindness was reflected in her reply: “Through our struggling years I didn’t expect much, so I was rarely disappointed.”

She is special. With her, it is easy to obey the scriptural injunction to “live joyfully with the wife whom thou lovest all the days of [thy] life” (Eccl. 9:9).

Not all of us are blessed with such wonderful eternal companions—not yet anyway. Many of us who are married cannot be together as much as we would like. Thankfully, we all have companionship of family and friends.

Recently another General Authority was my partner for mission tours to dusty places. On occasion, when I returned from a morning shower, I found to my surprise that this considerate companion had shined my shoes. Gratefully I wondered if each of the thirty thousand missionaries now laboring in the Lord’s service would have, and be, as kind a friend as he was to me, thoughtfully rendering simple acts of courtesy to a companion.

Joy cometh in the morning to those who have earned the night’s rest of a laborer. One of life’s sweetest returns is the privilege of rendering significant service of worth to others. To be able to do for fellow human beings something they could not do for themselves brings matchless satisfaction. Years of preparation are worth it.

And joy is derived in Church service. Alma so expressed this thought: “That perhaps I may be an instrument in the hands of God to bring some soul to repentance … is my joy” (Alma 29:9).

Through service in the temples, the concept of courtesy to companions can be nobly extended to those who have passed beyond the veil. The gospel brings glad tidings for the dead and a voice of gladness for the living and the dead—for all, glad tidings of great joy (see D\&C 128:19).

Even when death’s veil separates us from parents who gave so much that we might be, their righteous influence continues with us. And as they may watch from windows of heaven, their mornings will be brighter if they can truly say, as did the Apostle, “I have no greater joy than to hear that my children walk in truth” (3 Jn. 1:4).

Above all, courtesy to companions cannot be defiled by disobedience to the law of chastity. That sin is joy’s deadly poison. The first morning’s glance in the mirror cannot reflect joy if there is any recollection of misdeeds the night before. The surest step toward joy in the morning is virtue in the evening! Virtue includes courtesy to companions all day long.

\intalkheader{Self-esteem, Well Deserved}

The next prerequisite to joy is to feel good about yourself. The second of our Lord’s two great commandments carries a double charge: “Thou shalt love thy neighbour as thyself” (Matt. 22:39). Therefore, love of companion is governed, in part, by esteem of self, and so is joy in the morning.

Each individual should understand the nature of his or her own soul. Profound insight is provided by this revelation:

“For man is spirit. The elements are eternal, and spirit and element, inseparably connected, receive a fulness of joy;

“And when separated, man cannot receive a fulness of joy” (D\&C 93:33–34).

Therefore, spiritual and physical elements each must be nurtured if we are to earn proper self-esteem.

Spiritual self-esteem begins with the realization that each new morning is a gift from God. Even the air we breathe is a loving loan from him. He preserves us from day to day and supports us from one moment to another (see Mosiah 2:21).

Therefore, our first noble deed of the morning should be a humble prayer of gratitude. Scripture so counsels: “Pray unto God, and he will be favourable unto [you]: and [you] shall see his face with joy” (Job 33:26; see also Alma 34:21; 37:37).

I did not fully appreciate the significance of prayerful greetings until I became a father myself. I am so grateful that our children never gave their mother or dad the silent treatment. Now I sense how our Heavenly Father may appreciate our prayers, morning and night. But I can imagine the pangs of his sorrow because of silence from any of his children. To me, such ingratitude seems comparable to sullen goldfish oblivious to kind providers who sprinkle food in their bowl. Indeed, those who pray can “worship God with exceedingly great joy” (Alma 45:1).

I learned long ago that a period of uninterrupted scriptural study in the morning brings enduring enrichment. I feel as did Jeremiah: “Thy word was unto me the joy and rejoicing of mine heart” (Jer. 15:16). Sacred scriptures have been repeatedly described as “glad tidings of great joy” (Hel. 16:14; Mosiah 3:3; Alma 13:22; see also Luke 2:10). As we learn and abide their teachings, that joy becomes part of our lives.

Joy cometh in the morning when personal talents are developed. Each of us is blessed with different potential. I don’t think I could get up early enough to become a portrait painter. But I have appreciated teachings since my earliest childhood from parents who knew the joy that good music brings. And some of the sweetest sounds in our home have been those from songs and instruments of children improving their talents.

Even in Old Testament days, “David spake to the chief of the Levites to appoint their brethren to be the singers with instruments of musick, … sounding, by lifting up the voice with joy” (1 Chr. 15:16).

Confidence to begin each morning ready to meet the challenges of the day comes from spiritual self-esteem.

Physical self-esteem also requires nurturing. Our bodies deserve thoughtful care. I echo this declaration of Paul:

“Know ye not that ye are the temple of God, and that the Spirit of God dwelleth in you?

“If any man defile the temple of God, him shall God destroy; for the temple of God is holy, which temple ye are” (1 Cor. 3:16–17).

Physical conditioning from regular exercise is important. And we can do so much more to keep our bodies strong.

In 1833 the Prophet Joseph Smith received the Word of Wisdom by revelation. It includes these simple directives: We are not to drink alcohol, tea, or coffee, and we are not to use tobacco. Prophets in our generation and in this conference have told us also to avoid harmful drugs. Now medical science increasingly confirms the physical benefits of compliance with these teachings.

The damaging effects of alcohol are so widely known, additional comment is hardly needed. Harm inflicted by alcohol has been demonstrated, for example, by a study of the relationship between alcohol consumption among expectant mothers and the birth weight of their newborn infants. Findings published from the U. S. National Institutes of Health disclosed that consumption of one to two alcoholic drinks a day was associated with a substantially increased risk of producing a growth-retarded infant (see James L. Mills, et. al., “Maternal Alcohol consumption and Birth Weight,” Journal of the American Medical Association, 12 Oct. 1984, pp. 1875–79).

Scientists now know that smoking tobacco is the number one preventable cause of death in all the world. It is the leading preventable cause of heart disease, lung disease, artery disease, and cancer (see William W. Pollin and R. T. Ravenholt, “Tobacco Addiction and Tobacco Mortality,” Journal of the American Medical Association, 23 Nov. 1984, pp. 2849–54; 1986 Heart Facts, American Heart Association, 1986, p. 16; “The Health Consequences of Smoking: A Report of the Surgeon General,” publication DHHS (PHS) 84-50204, U.S. Department of Health and Human Services, Public Health Services, Office of Smoking and Health, Rockville, Md., 1983).

Still another report indicates that more than one-fourth of all deaths in the United States are now caused by conditions which physicians classify as addictive disorders (see Pollin and Ravenholt, p. 2849).

Obedience to the Word of Wisdom keeps one free from all such addictions. This protection is pronounced by covenant in the last verse of the eighty-ninth section of the Doctrine and Covenants:

“And I, the Lord, give unto them a promise, that the destroying angel shall pass by them, as the children of Israel, and not slay them” (D\&C 89:21).

This reference to the first Passover reminds us that, in faith, ancient Israel was obedient to the commandment to take blood and “strike it on the two side posts and on the upper door post of the houses” (Ex. 12:7).

“And the blood shall be to you for a token upon the houses where ye are: and when I see the blood, I will pass over you, and the plague shall not … destroy you” (Ex. 12:13).

So, in faith, modern Israel is commanded to obey the Word of Wisdom. It becomes our token of a covenant with the Lord—a spiritual separator of covenant Israel from the rest of the world.

Joy cometh in the morning—to those who can stand before the mirror and feel clean, to those whose mouths are free from the taste of flavors forbidden by the Lord, to those whose spirits and bodies are free from feelings of self-remorse.

\intalkheader{Love of God}

The crowning attribute that leads to joy is love of God. Even that first look in the mirror can be more enjoyable knowing we are created in his image. Each of us can say, as did the Apostle, “Thou hast made known to me the ways of life; thou shalt make me full of joy with thy countenance” (Acts 2:28; see also Ps. 16:11).

God, who gave us life, also gave us commandments to live by, that we might have joy. They have been revealed periodically by prophets from Adam to President Benson. One wrote this expression: “Consider … the blessed and happy state of those that keep the commandments of God. For behold, they are blessed in all things, both temporal and spiritual” (Mosiah 2:41).

But for those who have not known his ways or who have strayed from them, remember, it is not too late to change. Blessings from faith and repentance still can be yours.

To those who feel defeated and downtrodden, look to the early hours of the day for your rescue. The Lord tells us, “Cease to sleep longer than is needful; retire to thy bed early, that ye may not be weary; arise early, that your bodies and your minds may be invigorated” (D\&C 88:124).

The dawning of a brighter day heralds a time of forgiveness. Shadows of yesterday’s grief melt in the rays of early morn’s opportunity.

Joy comes from our posterity, and we rejoice as they are blessed by the ordinances of salvation and exaltation.

Our family experienced that in a special way recently as our youngest daughter was sealed to her eternal companion in the holy temple. To witness this event, along with other family members, were her parents and all eight of her older sisters and their husbands. For us, there was truly joy in the morning on that day. Then we really felt the scriptural truth, “Men are, that they might have joy” (2 Ne. 2:25).

These experiences, glorious as they are, become but prelude to that great day ahead, when the faithful will stand at the latter day upon the earth. They shall abide the Second Coming of the Lord and shall stand with him when he appears (see Mal. 3:2–12; 3 Ne. 24:2–12). On that joyous morning, the mirror will reflect the miracle of the first resurrection. The faithful shall be crowned with glory, immortality, and eternal life (see D\&C 75:5; 138:51).

Once again “morning stars [will] sing together, and … all the sons [and daughters] of God [will] shout for joy!” (D\&C 128:23; see also Job 38:7). For on that morning, “the glory of the Lord shall be revealed, and all flesh shall see it together” (Isa. 40:5; see also Ezek. 20:48; Luke 3:6; D\&C 101:23).

Then, “there’ll be peace and contentment evermore, ev’ry heart, ev’ry voice on that day will rejoice. … There’ll be joy in the morning on that day” (Natalie Sleeth, “Joy in the Morning,” Carol Stream, Illinois: Hope Publishing Co., 1977, pp. 4–5, 9–10). It shall be ours—through our faithfulness—I testify in the name of Jesus Christ, amen.

\newpage
\settalk{Life after Life}
\setspeaker{Russell M. Nelson}
\settalkdate{April 1987}
\talkheading{Life after Life}{Russell M. Nelson}

At this pulpit today, I am truly honored to follow the President of the Church, President Ezra Taft Benson. I sustain him as a living prophet of God! He was ordained an Apostle in 1943, before I entered medical school. Now he is the senior Apostle on earth, and I am privileged to add my gratitude for a prophet and for his message of instruction and inspiration.

He has taught us of the ministry of the resurrected Lord to the inhabitants of ancient America. This precious knowledge recorded in the Book of Mormon is of transcendent importance to all people. Truly, the Book of Mormon is another testament of Jesus Christ!

As President Benson spoke, my thoughts momentarily flashed back some years ago to a conversation I had with a book publisher. He was interested in the topic of possible continuation of life after what we know as death. The publisher asked if I could contribute stories from patients who had come close enough to death to experience the other side and yet had survived to share those accounts. Sensing public interest in that subject, he would entitle the book Life after Life.

When I considered that request, I remembered many such incidents which had been whispered in confidence to me over the years. But those seemed too sacred to share in a worldly way, especially to the benefit of a commercial venture. Besides, what would be the validity of isolated stories of life after life without supporting testimonies of witnesses?

To me, much more logical and convincing would be a study of well-documented and carefully witnessed evidences of life after life.

President Benson has spoken of one of these priceless records. The activities of the living Christ in America followed his own resurrection from the dead. Many witnesses in many places have seen the risen Lord, before, during, and after the account to which President Benson has referred.

\intalkheader{To Associates in the Holy Land}

\intalkheader{To the Nephites}

\intalkheader{To the Dead}

\intalkheader{To the Lost Tribes}

\intalkheader{To Those of This Dispensation}

After almost two thousand years, new witnesses to the resurrection of Jesus have added their testimonies of this transcendent truth.

Yes, the resurrection of Jesus Christ is one of the most carefully documented events in history. I have mentioned many of those appearances, but additional instances have been recorded.

Even more remarkable is the fact that his mission among men—the Atonement, the Resurrection—extends privileges of redemption from sin and a glorious resurrection to each and every one of us! In some marvelous way, fully comprehended only by Deity, this is his work and his glory—“to bring to pass the immortality and eternal life of man” (Moses 1:39).

Teachers in the Church so instruct the old and the young. Sometimes the results are humorous. One leader shared this story with me:

As a little boy came home from Primary one day, his mother asked him what he had learned, and he said, “My teacher told me that I used to be dust and I would be dust again once more. Is that true, Mommy?”

“Yes,” the mother replied. “A scripture tells us so: ‘For dust thou art, and unto dust shalt thou return’ (Gen. 3:19).”

The little boy was amazed at this! The next morning, he was scurrying around getting ready for school, looking for his shoes. He crawled under the bed. Lo and behold, there he saw balls of dust. He ran to his mother in wonder, saying, “Oh, Mommy, somebody’s under my bed, and they’re either coming or going.”

\intalkheader{Nature of the Resurrection}

Yes, compounds derived from dust—elements of the earth—are combined to make each living cell in our bodies. The miracle of the Resurrection is matched only by the miracle of our creation in the first place.

No one knows precisely how two germ cells unite to make one. Nor do we know how that resulting cell multiplies and divides to make others—some to become eyes that see, ears that hear, or fingers that feel glorious things about us. Each cell contains chromosomes with thousands of genes, chemically insuring identity and independence of each individual. Our bodies undergo constant rebuilding according to genetic recipes that are uniquely ours. Each time we take a bath, we lose not only dirt, but cells dead and dying, as they are replaced by a newer crop. This process of regeneration and renewal is but prelude to the promised phenomenon and future fact of our resurrection.

“If a man die, shall he live again?” asked Job (Job 14:14). In faith, he answered his own question:

“For I know that my redeemer liveth, and that he shall stand at the latter day upon the earth:

“And though after my skin worms destroy this body, yet in my flesh shall I see God” (Job 19:25–26).

At the time of our resurrection, we shall take up our immortal tabernacles. Bodies that now age, deteriorate, and decay will no longer be subject to processes of degeneration: “This mortal must put on immortality” (1 Cor. 15:53).

This great priesthood power of resurrection is vested in the Lord of this world. He taught that “all power is given unto me in heaven and in earth” (Matt. 28:18). Though he supplicated his Father for aid at the eleventh hour, the final victory over death was earned by the Son. These are his words:

“Therefore doth my Father love me, because I lay down my life, that I might take it again.

“No man taketh it from me, but I lay it down of myself. I have power to lay it down, and I have power to take it again. This commandment have I received of my Father” (John 10:17–18).

This power he subtly proclaimed when he said unto the Jews:

“Destroy this temple, and in three days I will raise it up. … But he spake of the temple of his body” (John 2:19–21).

The keys of the Resurrection repose securely with our Lord and Master. He said:

“I am the resurrection, and the life: he that believeth in me, though he were dead, yet shall he live:

“And whosoever liveth and believeth in me shall never die” (John 11:25–26).

But obedience to the commandments of God is requisite if one is to be resurrected with a celestial body. Laws leading to celestial glory will be discussed at this conference. Our challenge is to learn them and abide by them.

I thank God for his Son, Jesus Christ, for his mission in mortality, and for his ministry as the resurrected Lord. He brought about his own resurrection. Testimonies of thousands, from ancient and modern times, attest to the truth that the resurrected Jesus is the Savior of mankind. He brought about a universal resurrection: “For as in Adam all die, even so in Christ shall all be made alive” (1 Cor. 15:22; see also Mosiah 3:16).

His sacrifice and his glory assure that “the spirit and the body shall be reunited again in its perfect form; both limb and joint shall be restored to its proper frame, even as we now are at this time” (Alma 11:43).

Gratefully and positively, I affirm that there is life after life, first in the spirit world and then in the Resurrection, for each and every one of us. I know that God lives and that Jesus the Christ is his Son. He is “the resurrection, and the life” (John 11:25). He lives. He is my Master. I am his servant. With all my heart, I love him, and I testify of him, in his holy name—the name of Jesus Christ, amen.

\newpage
\settalk{Keys of the Priesthood}
\setspeaker{Russell M. Nelson}
\settalkdate{October 1987}
\talkheading{Keys of the Priesthood}{Russell M. Nelson}

You all know how valuable keys can be. Many, if not most, of us carry keys in our pockets at this very moment.

But the keys about which I shall speak are far more important. They are precious, powerful, and invisible! Some can lock and unlock in heaven as well as on earth. I speak of keys of the priesthood.

You boys hold what is sometimes referred to as the preparatory priesthood. Preparation, priesthood service, and keys are all related, but different. Service of any type requires preparation. But proper authorization to give that service requires keys. May I illustrate?

Prior to my call to the Twelve, I served as a medical doctor and surgeon. I had earned two doctor’s degrees. I had been certified by two specialty boards. That long preparation had consumed many years, yet it carried no legal permission. Keys were required. They were held by authorities of the state government and the hospitals in which I desired to work. Once those holding proper authority exercised those keys by granting me a license and permission, then I could perform operations. In return, I was obligated to obey the law, to be loyal, and to understand and not abuse the power of a surgeon’s knife. The important steps of preparation, permission, and obligation likewise pertain to other occupations.

Why is the power to act in the name of God more important? Because it is of eternal significance. We should understand the source of our authority and something of the keys that control its power. They may benefit every man, woman, and child who now lives, who has lived, and who yet will live upon the earth.

Let us consider keys of the priesthood through three scenes of history: in ancient days, during the Lord’s mortal ministry, and in modern times.

\intalkheader{Scene One—Ancient Days}

Scene one pertains to keys of the priesthood in ancient days, even before creation of the earth. We then existed as spirit children of God. Abraham was one of us. He was told that he was among rulers chosen—even before he was born (see Abr. 3:23). Scriptures also relate that the Lord God foreordained priests prepared from the foundation of the world according to his foreknowledge. Thus, our calling to the holy priesthood was foreseen before we were born (see Alma 13:1–5).

We know that Adam received priesthood authority before the world was formed (see Joseph Smith, Teachings of the Prophet Joseph Smith, sel. Joseph Fielding Smith, Salt Lake City: Deseret Book Co., 1938, pp. 157, 167–68; see also D\&C 78:16).

The potential of the priesthood is so vast that our comprehension of it is a challenge. The Prophet Joseph Smith declared, “The Priesthood is an everlasting principle, and existed with God from eternity, and will to eternity” (Teachings, p. 157).

President Brigham Young added that the priesthood is “the law by which the worlds are, were, and will continue for ever and ever” (Discourses of Brigham Young, sel. John A. Widtsoe, Salt Lake City: Deseret Book Co., 1941, p. 130).

Priesthood is the authority of God delegated to man to minister for the salvation of men. “The power of directing these labors constitutes the keys of the Priesthood” (Joseph F. Smith, “Editor’s Table,” Improvement Era, Jan. 1901, p. 230).

Many in scene one held those keys, such as Abraham, Isaac, Jacob, Joseph, Moses, Elias, and Elijah.

\intalkheader{Scene Two—The Mortal Ministry of the Lord}

Scene two pertains to keys of the priesthood during the mortal ministry of the Lord. Jesus revealed the extent of priesthood authority. To his Apostles, the Savior said, “I will give unto thee the keys of the kingdom of heaven: and whatsoever thou shalt bind on earth shall be bound in heaven: and whatsoever thou shalt loose on earth shall be loosed in heaven” (Matt. 16:19; see also D\&C 128:10).

Within a week of that promise, Jesus took Peter, James, and John to a high mountain. Keys of the priesthood were bestowed upon them by Moses and Elijah (Elias) under divine direction (see Matt. 17:1–5).

The Master then reminded his disciples of their sacred calls to the holy apostleship, “Ye have not chosen me, but I have chosen you, and ordained you” (John 15:16; see also D\&C 7:7).

Jesus not only called and ordained men, but he taught them their duties.

After Christ was crucified, and even before the early Apostles completed their labors, the Apostasy began. This occurred as prophesied when priesthood authority was abused and sacred ordinances were defiled.

\intalkheader{Scene Three—Modern Times}

With scenes one and two as background, let us consider scene three—the modern times in which we live. After centuries of apostasy, keys of the priesthood have again been restored. Under that authority, we have had hands laid upon our heads.

Obligations pertain to those who give and to those who receive ordinations or calls.

Perhaps that can be explained by example. I hold a set of keys to an automobile. In your mind, let them represent keys to something of value in your life—a tractor, an implement, or a powerful instrument. If I give keys to you, I have certain obligations, and you have certain obligations.

For me as the giver, I have a duty toward your success. Should you fail, in a measure I have failed. So I must teach and train adequately to ensure your personal safety and, at the same time, safeguard precious property you are to use.

For you as the receiver, obligations accompany the keys. You must know applicable laws and obey them. Loyalty is expected. And you should understand the power of your instrument. Obedience, loyalty, and understanding are implicit with your acceptance of those keys.

Now apply the same principles to keys of the priesthood. Your stake president, quorum president, and bishop hold keys of presidency. Their keys control the power of their unit of the Church. Those leaders not only may call and release, but they must train and bear sacred responsibility that the mission of the Church be accomplished.

They who receive ordinations or callings have obligations of obedience, loyalty, and understanding.

Obedience to law first and foremost means keeping the commandments of God. By so doing, one becomes worthy to receive personal revelation! Those who receive the Melchizedek Priesthood are under solemn oath and covenant to “live by every word that proceedeth forth from the mouth of God” (D\&C 84:44).

Loyalty is vital. Loyalty to him who has keys to call and to release you, even though he is an imperfect human being, will develop unity essential to success (see D\&C 124:45–46). The Lord defined this reality when he said, “Israel shall be saved in mine own due time; and by the keys which I have given shall they be led” (D\&C 35:25).

To understand the power of the priesthood, we must know its limitations. If an automobile is used recklessly, future permission from parents is likely to be denied. So, if priesthood power be abused, “the Spirit of the Lord is grieved; and when it is withdrawn, Amen to the priesthood or the authority of that man” (D\&C 121:37).

To understand the power of the priesthood, we must know the divinity of its restoration in these latter days. In 1820, our Heavenly Father and his Son Jesus Christ appeared to the Prophet Joseph Smith. In 1829, John the Baptist conferred the Aaronic Priesthood upon Joseph Smith and Oliver Cowdery (see D\&C 13; 27:8; JS—H 1:69, 72). Shortly thereafter they received the Melchizedek Priesthood under the hands of Peter, James, and John (see D\&C 27:11–12).

Later, the Lord spoke to Joseph and Oliver of others to whom specific keys of the priesthood had been committed. Each, in turn, did confer these keys:

Moroni: Keys of the Book of Mormon (see D\&C 27:5).

Moses: Keys of the gathering of Israel and the leading of the ten tribes (see D\&C 110:11).

Elias: Keys of the restoration of all things (see D\&C 27:6), including the Abrahamic covenant (see D\&C 110:12; Joseph Fielding Smith, Doctrines of Salvation, comp. Bruce R. McConkie, 3 vols., Salt Lake City: Bookcraft, 1954–56, 3:126–27).

Elijah: Keys of the sealing power (see D\&C 27:9; 110:13–16; see also 128:21).

Joseph Smith conferred all keys on all of the Twelve (see Doctrines of Salvation, 3:154–56). In turn, keys have been transferred to present leaders.

Today, President Ezra Taft Benson actively holds every restored key held by “all those who have received a dispensation at any time from the beginning of the creation” (D\&C 112:31; italics added; see also D\&C 128:18).

Surely a sacred moment of my life occurred April 12, 1984, when the First Presidency and members of the Quorum of the Twelve Apostles laid their hands upon my head. As had been done for others before me, all the keys of the priesthood were conferred. As it is with each member of the Quorum of the Twelve, some keys are not used until called upon by the Lord, or as directed by his senior Apostle.

I feel the weight of responsibility and the burden of timeless trust. I know those keys have been restored “for the last days and for the last time” (D\&C 112:30).

I am deeply grateful that we bear the priesthood—each of us foreordained from the foundation of the world for that responsibility (see Alma 13:1–5).

As a symbol of gratitude, I have penned a few lines to be sung as the concluding portion of my message. A verse for each of three scenes of history may summarize my remarks. This song will be new to you. The words have been written to music from Wales. For the gracious support of Brother Jerold D. Ottley and our male chorus I express my appreciation. To this song I have assigned an ancient title, “Hosanna,” the prayerful shout of fervent praise:

\begin{verse}
Through time’s immortal endless stay\\
In love he guides our way.\\
Beyond the realms of heaven’s beam,\\
Our great God, Elohim.\\
Hosanna to his holy name—\\
Our fathers’ God is still the same.
\end{verse}

May we be true to the trust he has given to us who bear the holy priesthood and hold its sacred keys, I pray in the name of Jesus Christ, amen.

\newpage
\settalk{Lessons from Eve}
\setspeaker{Russell M. Nelson}
\settalkdate{October 1987}
\talkheading{Lessons from Eve}{Russell M. Nelson}

President and Sister Benson, I hope you can feel the love and spirit emanating from the thousands of sisters at this General Women’s Meeting. President Hinckley and President Monson, we extend our love to you and your companions. The attendance of members of the First Presidency symbolizes their support for each of us. On their behalf, I extend their gratitude and love to all of you.

I deeply appreciate the invocation and the messages delivered by the General Presidency of the Relief Society and by this wonderful women’s chorus.

Sister Doxey has recounted the privilege that each of us has to strengthen the family—the basic component of society, and the fundamental unit of the Church. The family is provided in God’s plan to nurture his children and prepare them to return to him—as families.

Sister Evans has helped us understand that the human race doesn’t run on a smooth track. The route more closely resembles an obstacle course with hurdles, puddles, or snares lurking around every turn we make. Life isn’t meant to be easy. The race isn’t to be won by the swift or the flashy. Victory comes only to those who muster the faith to stay on the track—the strait and narrow way.

Sister Winder has put our purpose in perspective. She has taught an eternal truth—“Men are, that they might have joy” (2 Ne. 2:25). And she has emphasized one of the means by which true joy is realized: rendering selfless service of worth to others.

As the male participant tonight, I bear a great responsibility to convey deep feelings of gratitude. For the men of the Church, I say thank you!

Not only do I express gratitude, but our affection as well. Represented in this vast audience are our dear companions, our mothers, our sisters, our daughters, and our granddaughters. We respect and honor you. We support your wonderful efforts. We are grateful to receive the benefit of your counsel and insights in challenges we share as partners. And we thank you as mothers or potential mothers for your partnership with God in providing life. Without women, the whole purpose of the creation of this world would be in vain.

This truth we learn from scriptures about the priesthood, the Creation, Adam and Eve.

Before the world was formed, the Lord Jesus Christ was Jehovah, “the Great I Am … the beginning and the end, … [who] looked upon the wide expanse of eternity, and all the seraphic hosts of heaven” (D\&C 38:1; see also D\&C 29:1; D\&C 39:1).

He told Abraham, having first shown him “the intelligences that were organized before the world was; among all these there were many of the noble and great ones;

“And God saw these souls that they were good, and he stood in the midst of them, and he said: These I will make my rulers” (Abr. 3:22–23).

The Lord then revealed that Abraham was one of them, chosen and foreordained before he was born.

Scriptural verse then continues:

“There stood one among them that was like unto God, and he said unto those who were with him: We will go down, … and we will make an earth whereon these may dwell” (Abr. 3:24).

Then “the Gods, organized and formed the heavens and the earth” (Abr. 4:1).

You will recall that after the earth had been created, divided, beautified, and inhabited with plant and animal life, the crowning achievement of the Creation was to be man—the human being. “So the Gods went down to organize man in their own image, … male and female to form they them” (Abr. 4:27; see also Gen. 1:26).

The very purpose of creation was to provide bodies, to enable these eagerly awaiting spirits to enjoy mortal life and experiences.

Question: What was the role of the priesthood in the process of creation?

The Prophet Joseph Smith taught:

“The Priesthood is an everlasting principle, and existed with God from eternity. … Christ is the Great High Priest; Adam next” (Teachings of the Prophet Joseph Smith, sel. Joseph Fielding Smith, Salt Lake City: Deseret Book Co., 1938, pp. 157–58).

“The priesthood was first given to Adam; he obtained the First Presidency, and held the keys of it from generation to generation. He obtained it in the Creation, before the world was formed” (Teachings, p. 157).

President Brigham Young said, “Priesthood … is the law by which the worlds are, were, and will continue for ever and ever” (Discourses of Brigham Young, sel. John A. Widtsoe, Salt Lake City: Deseret Book Co., 1941, p. 130).

Thus, priesthood is the power of God. Its ordinances and covenants are to bless men and women alike. By that power, the earth was created. Under the direction of the Father, Jehovah was the creator. As Michael, Adam did his part. He became the first man. But, in spite of the power and glory of creation to that point, the final link in the chain of creation was still missing. All the purposes of the world and all that was in the world would be brought to naught without woman—a keystone in the priesthood arch of creation.

When Eve was created—when her body was made by God—Adam exclaimed, “Bone of my bones, and flesh of my flesh; she shall be called Woman, because she was taken out of man” (Moses 3:23).

From the rib of Adam, Eve was formed (see Gen. 2:22; Moses 3:22; Abr. 5:16). Interesting to me is the fact that animals fashioned by our Creator, such as dogs and cats, have thirteen pairs of ribs, but the human being has one less with only twelve. I presume another bone could have been used, but the rib, coming as it does from the side, seems to denote partnership. The rib signifies neither dominion nor subservience, but a lateral relationship as partners, to work and to live, side by side.

Adam and Eve were joined together in marriage for time and for all eternity by the power of that everlasting priesthood (see Gen. 2:24–25; Moses 3:25; Abr. 5:18–19). Eve came as a partner, to build and to organize the bodies of mortal men. She was designed by Deity to cocreate and nurture life, that the great plan of the Father might achieve fruition. Eve “was the mother of all living” (Moses 4:26). She was the first of all women.

From our study of Eve, we may learn five fundamental lessons of everlasting importance:

From these five fundamental lessons, we can study patterns which apply to present-day circumstances. Let us review them, lesson by lesson.

We can’t all do all things. Circumstances, available time, and talents vary widely among us. Your diversities take you to numerous arenas of activity. There let your presence be felt.

Your foes in a sordid society demean the sacredness of women and the sanctity of motherhood. Your world, sickened by unchastity and plagued with sexually transmitted disease, needs your righteous example. For the wrath of God is provoked by governments that sponsor gambling, condone pornography, or legalize abortion. These forces serve to denigrate women now, just as they did in the days of Sodom and Gomorrah.

You can—you must—make a difference. You are vital to the Lord’s team—one team with one purpose. Through your diversity, build strength in unity. Bind yourselves together in all holiness. Anchor yourselves to “the foundation of the apostles and prophets, Jesus Christ himself being the chief corner stone” (Eph. 2:20).

As Sister Winder mentioned, the funeral service for Sister Camilla Kimball was held earlier today. Magnificent and well-deserved tributes were paid to her. We can learn from her great example of courage. Let us emulate her faith and likewise heed this teaching from the Book of Mormon, which gave her such strength and security:

“Remember that it is upon the rock of our Redeemer, who is Christ, the Son of God, that ye must build your foundation; that when the devil shall send forth his mighty winds, yea, his shafts in the whirlwind, yea, when all his hail and his mighty storm shall beat upon you, it shall have no power over you to drag you down to the gulf of misery and endless wo, because of the rock upon which ye are built, which is a sure foundation, a foundation whereon if men build they cannot fall” (Hel. 5:12).

Beloved sisters, let your lives be committed to your Father in Heaven, to his Only Begotten Son, and to his church, restored by them in this latter day. I testify that it is true and that President Benson is the prophet on earth today. A dispensation of the gospel has been entrusted to our care. Everlasting priesthood principles, laws, and powers depend upon our partnership.

I invoke the blessings of Almighty God upon you that you may be successful in fulfilling your divine destiny, in the name of Jesus Christ, amen.

\newpage
\settalk{“With God Nothing Shall Be Impossible”}
\setspeaker{Russell M. Nelson}
\settalkdate{April 1988}
\talkheading{“With God Nothing Shall Be Impossible”}{Russell M. Nelson}

I echo the testimony of our beloved associate, Elder Dallin H. Oaks. With him and the other Brethren, I applaud the efforts of Latter-day Saints throughout the world who willingly serve in building the kingdom of God. Likewise, I respect those who quietly do their duty though deepening trials come their way. And I admire those who strive to be more worthy by overcoming a personal fault or who work to achieve a difficult goal.

I feel impressed to counsel those engaged in personal challenges to do right. In particular, my heart reaches out to those who feel discouraged by the magnitude of their struggle. Many shoulder heavy burdens of righteous responsibility which, on occasion, seem so difficult to bear. I have heard those challenges termed impossible.

As a medical doctor, I have known the face of adversity. I have seen much of death and dying, suffering and sorrow. I also remember the plight of students overwhelmed by their studies and of those striving to learn a foreign language. And I recall the fatigue and frustration felt by young parents with children in need. Amidst circumstances seemingly impossible, I have also experienced the joyous relief that comes when one’s understanding is deepened by scriptural insight.

The Lord has often chosen to instruct His people in their times of trial. Scriptures show that some of His lasting lessons have been taught with examples terrible as war, commonplace as childbearing, or obvious as hazards of deep water. His teachings are frequently based on common understanding, but with uncommon results. Indeed, one might say that to teach His people, the Lord employs the unlikely.

Warfare, for example, has been known since time began. Even in that ugly circumstance, the Lord has helped those obedient to His counsel. Going into battle, all would assume the obvious advantage of outnumbering an enemy. But when God’s disciple Gideon was leading an army against the Midianites, “the Lord said unto Gideon, The people that are with thee are too many … , lest Israel vaunt themselves … , saying, Mine own hand hath saved me” (Judg. 7:2).

So the Lord directed Gideon to reduce his numbers. He first decreased the troops from twenty-two thousand to ten thousand.

Then the Lord said to Gideon, “The people are yet too many” (Judg. 7:4). So another reduction was made. Finally, only three hundred remained. Then the Lord delivered the victory to the outnumbered few (see Judg. 7:5–25).

Even more widely known than war is an understanding of childbearing. Everyone “knows” that old women do not bear children. So upon whom did the Lord call to bear Abraham’s birthright son? Sarah, at age ninety! When told this was to be, she asked a logical question: “Shall I [which am old] of a surety bear a child?” (Gen. 18:13). From heaven came this reply: “Is any thing too hard for the Lord?” (Gen. 18:14).

So decreed, she gave birth to Isaac, to carry the crucial Abrahamic covenant into the second generation (see Gen. 26:1–4, 24).

Later, for one of the most important events ever to occur, the other extreme was chosen. As all knew that an elderly woman could not bear children, it was just as obvious that a virgin could not have children. But Isaiah had made this prophetic utterance:

“The Lord himself shall give you a sign; Behold, a virgin shall conceive, and bear a son, and shall call his name Immanuel” (Isa. 7:14).

When Mary was notified of her sacred responsibility, the announcing angel reassured, “For with God nothing shall be impossible” (Luke 1:37).

The expression deep water means danger! That very hazard challenged the Israelites led by Moses at the Red Sea (see Ex. 14). Later, they were led by Joshua to the river Jordan at flood time (see Josh. 3). In each instance, deep water was divinely divided to allow the faithful to reach their destination safely. To teach His people, the Lord employs the unlikely.

Turning to our day, have you ever wondered why the Master waited so long to inaugurate the promised “restitution of all things” (Acts 3:21)? Any competitor knows the disadvantage of allowing an opponent to get too far ahead. Wouldn’t the work of the restoration of the Church have been easier if begun earlier?

Suppose for a moment you are a member of a team. The coach beckons you from the bench and says: “You are to enter this contest. I not only want you to win; you shall win. But the going will be tough. The score at this moment is 1,143,000,000 to six, and you are to play on the team with the six points!”

That large number was the approximate population of the earth in the year 1830 when the restored church of Jesus Christ was officially organized with six members (see James Avery Joyce, sel., World Population Basic Documents, 4 vols., Dobbs Ferry, New York: Oceana Publications, Inc., 1976, 4:2214). The setting was remote and rural. By standards of the world, its leaders were deemed to be unlearned. Their followers seemed so ordinary. But with them, the work was begun. Assignments had been revealed:

The great dispensation of the latter days had commenced, and they were the ones to usher it forth!

Furthermore, the Prophet Joseph Smith was unjustly held in the unspeakable isolation of a distant prison. In such obscurity, then and there, he was told by the Lord that “the ends of the earth shall inquire after thy name” (D\&C 122:1).

If any tasks ever deserved the label impossible, those would seem to qualify. But, in fact, our Lord had spoken: “With men this is impossible; but with God all things are possible” (Matt. 19:26; see also Mark 10:27; Luke 18:27). To teach His people, the Lord employs the unlikely.

A century and a half later, the burdening baton of that opportunity has now been passed to us. We are children of the noble birthright, who must carry on in spite of our foredetermined status to be broadly outnumbered and widely opposed. Challenges lie ahead for the Church and for each member divinely charged toward self-improvement and service.

How is it possible to achieve the “impossible”? Learn and obey the teachings of God. From the holy scriptures, heaven-sent lift will be found for heaven-sent duties. To so achieve, at least three basic scriptural themes loom repeatedly as requirements.

\intalkheader{Faith}

The foremost requisite is faith. It is the first principle of the gospel (see A of F 1:4). In his epistle to the Hebrews, Paul so taught. He concluded that by faith the great deeds of Noah, Abraham, Sarah, Isaac, Jacob, Joseph, Moses, Joshua, and others were accomplished (see Heb. 11:4–34).

Prophets on the American hemisphere similarly taught the fundamental importance of faith. Moroni said it included things “hoped for and not seen” and then warned his skeptics, “Dispute not because ye see not, for ye receive no witness until after the trial of your faith” (Ether 12:6). Then he spoke of leaders whose faith preceded their miraculous deeds, including Alma, Amulek, Nephi, Lehi, Ammon, the brother of Jared, and the three who were promised that they should not taste of death (see Ether 12:13–20).

The Lord personally taught this truth to his disciples: “If ye have faith,” he said, “nothing shall be impossible unto you” (Matt. 17:20).

Faith is nurtured through knowledge of God. It comes from prayer and feasting upon the words of Christ through diligent study of the scriptures.

\intalkheader{Focus}

The second requisite I have classified as focus. Imagine, if you will, a pair of powerful binoculars. Two separate optical systems are joined together with a gear to focus two independent images into one three-dimensional view. To apply this analogy, let the scene on the left side of your binoculars represent your perception of your task. Let the picture on the right side represent the Lord’s perspective of your task—the portion of His plan He has entrusted to you. Now, connect your system to His. By mental adjustment, fuse your focus. Something wonderful happens. Your vision and His are now the same. You have developed an “eye single to the glory of God” (D\&C 4:5; see also Morm. 8:15). With that perspective, look upward—above and beyond mundane things about you. The Lord said, “Look unto me in every thought” (D\&C 6:36). That special vision will also help clarify your wishes when they may be a bit fuzzy and out of focus with God’s hopes for your divine destiny. Indeed, the precise challenge you regard now as “impossible” may be the very refinement you need, in His eye.

Recently I visited the home of a man terminally ill. The stake president introduced me to the man’s family. His wife demonstrated such focus when she asked for a blessing for her dying husband—not for healing, but for peace, not for a miracle, but for ability to abide to the end. She could see from an eternal viewpoint, not merely from the perspective of one weighted with the responsibilities of her husband’s day-to-day care.

Elsewhere, a mother with focus nurtures her son, crippled for the whole of this life. Daily she thanks her Heavenly Father for the privilege of laboring in love with a child for whom mortality’s vale of tears will be mercifully brief. Her focus is fixed on eternity. With celestial sight, trials impossible to change become possible to endure.

\intalkheader{Strength and Courage}

A third theme in the scriptures requisite for significant accomplishment is difficult to summarize in one word, so I shall link two to describe it—strength and courage. Repeatedly, scriptures yoke these attributes of character together, especially when difficult challenges are to be conquered (see Deut. 31:6, 7, 23; Josh. 1:6, 7, 9, 18; 10:25; 1 Chr. 22:13; 28:20; 2 Chr. 32:7; Ps. 27:14; 31:24; Alma 43:43; 53:20).

Perhaps this is more easily illustrated than defined. Our pioneer forefathers are good examples. They sang, “Gird up your loins; fresh courage take” (“Come, Come, Ye Saints,” Hymns, 1985, no. 30). They feared no toil and no labor. Among them were Johan Andreas Jensen and his wife, Petra, who left their native Norway in 1863. Their family included six-week-old tiny twin daughters. As handcarts were pulled in their rugged journey, one of those little girls died along the way. The child who survived grew up to become my Grandmother Nelson!

There are pioneers in the Church today just as strong and courageous. Recently, I interviewed a married couple three days after their release as full-time missionaries in a large metropolis. “We are converts,” they said. “We joined the Church ten years ago. Even though we just completed a mission, we want to go again! But this time, we would like to volunteer for a more difficult assignment. We want to teach and serve children of God who live in remote areas of the world!”

As I countered with the grim realities of their request, they continued their expression of commitment. “Our three children and their spouses will assist with our expenses. Two of those couples have joined the Church already, and the third is equally supportive. Please send us among humble people who love the Lord and desire to know that His Church has again been restored to the earth.” Needless to say, their petition was gratefully heard, and now they have received their second call to missionary service.

Strength and courage also characterize another couple. As faithful members of the Church, they had always upheld its doctrines, including the twelfth article of faith. When their country went to war, military conscription called the dutiful husband away from his wife before either had learned she was to bear their child. He was captured by enemy troops and taken as a prisoner of war. Months elapsed. Their baby came. Still no word to know whether the new father was alive. A year after his capture, he was permitted to write to his wife.

Meanwhile, though countries apart, they each remained faithful to covenants made at baptism. Even though clothed in prisoner’s stripes and able to speak the language of his captors’ country only in a limited way, he became Sunday School superintendent of the branch. He baptized four fellow prisoners during their confinement. Three years after the war ended, he returned home to his wife and a son he had never seen. Later, he served for ten years as the first stake president of his country. Now he is a member of the presidency of one of our temples! His wife stands faithfully beside him in the privilege of that sacred assignment.

You who may be momentarily disheartened, remember, life is not meant to be easy. Trials must be borne and grief endured along the way. As you remember that “with God nothing shall be impossible” (Luke 1:37), know that He is your Father. You are a son or daughter created in His image, entitled through your worthiness to receive revelation to help with your righteous endeavors. You may take upon you the holy name of the Lord. You can qualify to speak in the sacred name of God (see D\&C 1:20). It matters not that giants of tribulation torment you. Your prayerful access to help is just as real as when David battled his Goliath (see 1 Sam. 17).

Foster your faith. Fuse your focus with an eye single to the glory of God. “Be strong and courageous” (2 Chr. 32:7), and you will be given power and protection from on high. “For I will go before your face,” the Lord declared. “I will be on your right hand and on your left, and my Spirit shall be in your hearts, and mine angels round about you, to bear you up” (D\&C 84:88).

The great latter-day work of which we are a part shall be accomplished. Prophecies of the ages shall be fulfilled. “For with God all things are possible” (Mark 10:27), I testify in the name of Jesus Christ, amen.

\newpage
\settalk{Addiction or Freedom}
\setspeaker{Russell M. Nelson}
\settalkdate{October 1988}
\talkheading{Addiction or Freedom}{Russell M. Nelson}

I am impressed to speak out on a problem of deep concern—the worldwide epidemic of drug addiction. As a medical doctor, my study of drugs began early in medical school. Each doctor spends months in specialized courses learning potential benefits and risks of medicinal agents. Proper prescription of drugs is the forte of skilled physicians. Generally, when their advice is carefully followed, results are remarkably successful. In addressing this topic, I specifically exclude such application of modern knowledge by educated professionals.

But I raise my voice with others throughout the world who warn against abuse of drugs beyond prescribed limits, and the recreational or social use of chemical substances so often begun naively by the ill-informed.

From an initial experiment thought to be trivial, a vicious cycle may follow. From trial comes a habit. From habit comes dependence. From dependence comes addiction. Its grasp is so gradual. Enslaving shackles of habit are too small to be sensed until they are too strong to be broken. Indeed, drugs are the modern “mess of pottage” for which souls are sold. No families are free from risk.

But this problem is broader than hard drugs. Their use most often begins with cigarette smoking. Tobacco and alcoholic beverages contain addicting drugs. They lead the list in incidence and cost to society.

As I speak with governmental and medical leaders of many nations, they voice grave concern over the consumption of alcohol and other substances by their citizens. Though the extent of the challenge is international, data from the United States of America will be cited solely to indicate the monstrous scope of this worldwide problem.

\intalkheader{Tobacco}

Consider the magnitude of tobacco’s harm. Cigarette smoking is the most frequent preventable cause of heart disease, artery disease, lung disease, and cancer. In the USA in 1982, 16 percent of all deaths (314,000) were attributed to the smoking of tobacco.

For the year 1985, the estimated cost of both smoking-related health care and lost productivity amounted to \$65 billion. That calculates to an average of \$2.17 per pack of cigarettes sold. Social consequences of smoking far exceed the price paid to purchase cigarettes.

An insurance company recently reported that one-fifth of all its claims were for afflictions that could have been prevented by simply not smoking. We all bear this financial burden of illness that need not be.

Surgeon General Dr. C. Everett Koop and his team of more than fifty scientists recently published a landmark report. For the USA alone, they attributed 320,000 deaths annually to tobacco, 125,000 to alcohol, and lesser mortality to cocaine (2,000) and other opioids (4,000). They declared nicotine to be a powerfully addicting drug in the same sense as are drugs such as heroin and cocaine. Comparable views have been recorded by medical authorities in many other nations. Yet many of our good friends who use tobacco may not believe it to be addicting. Some are reluctant to admit that their behavior is substantially controlled by a drug. We understand those feelings.

\intalkheader{Alcohol}

There is mounting concern worldwide over the consumption of alcohol. The U.S. government estimates that 10.6 million adults are alcoholics and that one family in four is troubled by alcohol. It is a factor in half of all the nation’s traffic deaths.

Last year, a tragic milestone was reached. More Americans had been killed from alcohol-related motor vehicle accidents (1,350,000) than had been killed in all the wars America has ever fought (1,156,000).

\intalkheader{Other Drugs}

Drugs such as LSD, marijuana, heroin, and cocaine are also endangering people throughout the earth. The noble attributes of reason, integrity, and dignity, which distinguish men and women from all other forms of life, are often the first to be attacked by these drugs and alcohol.

\intalkheader{Reaching Help}

We reach out in love to family, friends, and neighbors, regardless of nationality or creed, who suffer addiction. The Church of Jesus Christ of Latter-day Saints continues to help relieve this international plague.

The solution to this problem ultimately is neither governmental nor institutional. Nor is it a question of legality. It is a matter of individual choice and commitment. Agency must be understood. The importance of the will in making crucial choices must be known. Then steps toward relief can follow.

\intalkheader{Agency}

Agency, or the power to choose, was ours as spirit children of our Creator before the world was. (See Alma 13:3; Moses 4:4.) It is a gift from God, nearly as precious as life itself.

Often, however, agency is misunderstood. While we are free to choose, once we have made those choices, we are tied to the consequences of those choices.

We are free to take drugs or not. But once we choose to use a habit-forming drug, we are bound to the consequences of that choice. Addiction surrenders later freedom to choose. Through chemical means, one can literally become disconnected from his or her own will!

\intalkheader{Road to Recovery}

For relief of an ailment, as a doctor of medicine I might write a prescription. As an ordained Apostle, I would invoke the spiritual blessing of eternal perspective. Combined, my spiritual prescription would return the gift of agency to its rightful owner.

Each one who resolves to climb that steep road to recovery must gird up for the fight of a lifetime. But a lifetime is a prize well worth the price.

This challenge uniquely involves the will, and the will can prevail. Healing doesn’t come after the first dose of any medicine. So the prescription must be followed firmly, bearing in mind that it often takes as long to recover as it did to become ill. But if made consistently and persistently, correct choices can cure.

\intalkheader{Spiritual Prescription}

My spiritual prescription includes six choices which I shall list alphabetically, A through F, and then comment about each:

The Lord has revealed His sacred standard to guide people in a troubled world. You and I were born free to follow His divine guidance. We may choose for ourselves. Those choices may bring addiction or freedom. For freedom and joy, choose to “be faithful in Christ.” He will lift you up. May “the hope of his glory and of eternal life, rest in your mind forever” (Moro. 9:25), I pray, in the name of Jesus Christ, amen.

\newpage
\settalk{The Canker of Contention}
\setspeaker{Russell M. Nelson}
\settalkdate{April 1989}
\talkheading{The Canker of Contention}{Russell M. Nelson}

Some months ago my esteemed colleague Elder Carlos E. Asay and I stood atop Mount Nebo, where Moses once stood. (See Deut. 34:1–4.) We saw what he saw. In the distance to our right was the Sea of Galilee. The river Jordan flowed from there to the Dead Sea on our left. Ahead was the promised land into which Joshua led the Israelite faithful so long ago.

Later we were permitted to do what Moses could not. We were escorted from the Hashemite kingdom of Jordan to its western border with Israel. From there, we and our associates walked over the Allenby Bridge. We felt the tension as armed soldiers nearby guarded both sides of the international boundary.

After safely enduring this experience, I thought of the irony of it all. Here in the land made holy by the Prince of Peace, contention has existed almost continuously from that day to this.

Prior to His ascension from the Holy Land, the Savior pronounced a unique blessing: “Peace I leave with you, my peace I give unto you: not as the world giveth, give I unto you.” (John 14:27.)

His peace is not necessarily political; His peace is personal. But that spirit of inner peace is driven away by contention. Contention does not usually begin as strife between countries. More often, it starts with an individual, for we can contend within ourselves over simple matters of right and wrong. From there, contention can infect neighbors and nations like a spreading sore.

As we dread any disease that undermines the health of the body, so should we deplore contention, which is a corroding canker of the spirit. I appreciate the counsel of Abraham Lincoln, who said:

“Quarrel not at all. No man resolved to make the most of himself can spare time for personal contention. … Better give your path to a dog than be bitten by him.” (Letter to J. M. Cutts, 26 Oct. 1863, in Concise Lincoln Dictionary of Thoughts and Statements, comp. and arr. Ralph B. Winn, New York: New York Philosophical Library, 1959, p. 107.)

President Ezra Taft Benson in his keynote address yesterday described contention as “another face of pride.”

My concern is that contention is becoming accepted as a way of life. From what we see and hear in the media, the classroom, and the workplace, all are now infected to some degree with contention. How easy it is, yet how wrong it is, to allow habits of contention to pervade matters of spiritual significance, because contention is forbidden by divine decree:

“The Lord God hath commanded that men should not murder; that they should not lie; that they should not steal; that they should not take the name of the Lord their God in vain; that they should not envy; that they should not have malice; that they should not contend one with another.” (2 Ne. 26:32.)

\intalkheader{Creator of Contention}

To understand why the Lord has commanded us not to “contend one with another,” we must know the true source of contention. A Book of Mormon prophet revealed this important knowledge even before the birth of Christ:

“Satan did stir them up to do iniquity continually; yea, he did go about spreading rumors and contentions upon all the face of the land, that he might harden the hearts of the people against that which was good and against that which should come.” (Hel. 16:22.)

When Christ did come to the Nephites, He confirmed that prophecy:

“He that hath the spirit of contention is not of me [saith the Lord], but is of the devil, who is the father of contention, and he stirreth up the hearts of men to contend with anger, one with another.

“Behold, this is not my doctrine, to stir up the hearts of men with anger, one against another; but this is my doctrine, that such things should be done away.” (3 Ne. 11:29–30.)

\intalkheader{Origin of Contention}

Contention existed before the earth was formed. When God’s plan for creation and mortal life on the earth was first announced, sons and daughters of God shouted for joy. The plan was dependent on man’s agency, his subsequent fall from the presence of God, and the merciful provision of a Savior to redeem mankind. Scriptures reveal that Lucifer sought vigorously to amend the plan by destroying the agency of man. Satan’s cunning motive was unmasked in his statement:

“Behold, here am I, send me, I will be thy son, and I will redeem all mankind, that one soul shall not be lost, and surely I will do it; wherefore give me thine honor.” (Moses 4:1.)

Satan’s selfish efforts to alter the plan of God resulted in great contention in heaven. The Prophet Joseph Smith explained:

“Jesus said there would be certain souls that would not be saved; and the devil said he could save them all, and laid his plans before the grand council, who gave their vote in favor of Jesus Christ. So the devil rose up in rebellion against God, and was cast down.” (Teachings of the Prophet Joseph Smith, sel. Joseph Fielding Smith, Salt Lake City: Deseret Book Co., 1938, p. 357.)

This war in heaven was not a war of bloodshed. It was a war of conflicting ideas—the beginning of contention.

Scriptures repeatedly warn that the father of contention opposes the plan of our Heavenly Father. Satan’s method relies on the infectious canker of contention. Satan’s motive: to gain personal acclaim even over God Himself.

\intalkheader{Targets of the Adversary}

The work of the adversary may be likened to loading guns in opposition to the work of God. Salvos containing germs of contention are aimed and fired at strategic targets essential to that holy work. These vital targets include—in addition to the individual—the family, leaders of the Church, and divine doctrine.

\intalkheader{The Family}

The family has been under attack ever since Satan first taunted Adam and Eve. (See Gen. 3; Moses 4.) So today, each must guard against the hazard of contention in the family. It usually begins innocently. Years ago when our daughters were little girls who wanted to be big girls, the style of the day was to wear multiple petticoats. A little contention could have crept in as the girls soon learned that the one to get dressed first was the one best dressed.

In a large family of boys, those with the longest reach were the best fed. In order to avoid obvious contention, they adopted a rule that required them at mealtime to leave at least one foot on the floor.

The home is the great laboratory of learning and love. Here parents help children overcome these natural tendencies to be selfish. In rearing our own family, Sister Nelson and I have been very grateful for this counsel from the Book of Mormon:

“Ye will not suffer your children that they go hungry, or naked; neither will ye suffer that they transgress the laws of God, and fight and quarrel one with another, …

“But ye will teach them to walk in the ways of truth and soberness; ye will teach them to love one another, and to serve one another.” (Mosiah 4:14–15.)

And I might add, please be patient while children learn those lessons.

Parents should be partners to cherish and protect one another, knowing that the aim of the adversary is to destroy the integrity of the family.

\intalkheader{Leaders of the Church}

Leaders of the Church are targets for attack by those who stir contention. This is true even though not a single leader has called himself or herself to a position of responsibility. Each General Authority, for instance, chose another path to pursue as his life’s occupation. But the reality is, as with Peter or Paul, each was surely “called of God, by prophecy, and by the laying on of hands by those who are in authority.” (A of F 1:5.) With that call comes the commitment to emulate the patterns of the Prince of Peace.

That goal is shared by worthy servants of the Master, who would not speak ill of the Lord’s anointed nor provoke contention over teachings declared by ancient or living prophets.

Certainly no faithful follower of God would promote any cause even remotely related to religion if rooted in controversy, because contention is not of the Lord.

Surely a stalwart would not lend his or her good name to periodicals, programs, or forums that feature offenders who do sow “discord among brethren.” (Prov. 6:19; see also Prov. 6:14.)

Such agitators unfortunately fulfill long-foretold prophecy: they “take counsel together, against the Lord, and against his anointed.” (Ps. 2:2.)

Yet, mercifully, the anointed pray for those who attack them, knowing the sad fate prophesied for their attackers. (See D\&C 121:16–22.)

Throughout the world, Saints of the Lord follow Him and His anointed leaders. They have learned that the path of dissent leads to real dangers. The Book of Mormon carries this warning:

“Now these dissenters, having the same instruction and the same information … , having been instructed in the same knowledge of the Lord, nevertheless, it is strange to relate, not long after their dissensions they became more hardened and impenitent, and more wild, wicked and ferocious … ; giving way to indolence, and all manner of lasciviousness; yea, entirely forgetting the Lord their God.” (Alma 47:36.)

How divisive is the force of dissension! Small acts can lead to such great consequences. Regardless of position or situation, no one can safely assume immunity to contention’s terrible toll.

Thomas B. Marsh, once one of the Twelve, left the Church. His spiritual slide to apostasy started because his wife and another woman had quarreled over a little cream! After an absence from the Church of nearly nineteen years, he came back. To a congregation of Saints, he then said:

“If there are any among this people who should ever apostatize and do as I have done, prepare your backs for a good whipping, if you are such as the Lord loves. But if you will take my advice, you will stand by the authorities.” (In Journal of Discourses, 5:206; see also Gordon B. Hinckley, in Conference Report, Apr. 1984, pp. 109–12; or Ensign, May 1984, pp. 81–83.)

Of course the authorities are human. But to them God has entrusted the keys to His divine work. And He holds us accountable for our responses to the teachings of His servants. These are the words of the Lord:

“If my people will hearken unto my voice, and unto the voice of my servants whom I have appointed to lead my people, behold, verily I say unto you, they shall not be moved out of their place.

“But if they will not hearken to my voice, nor unto the voice of these men whom I have appointed, they shall not be blest.” (D\&C 124:45–46.)

\intalkheader{Divine Doctrine}

Divine doctrine of the Church is the prime target of attack by the spiritually contentious. Well do I remember a friend who would routinely sow seeds of contention in Church classes. His assaults would invariably be preceded by this predictable comment: “Let me play the role of devil’s advocate.” Recently he passed away. One day he will stand before the Lord in judgment. Then, I wonder, will my friend’s predictable comment again be repeated?

Such contentious spirits are not new. In an epistle to Timothy, the Apostle Paul gave this warning, “that the name of God and his doctrine be not blasphemed.” (1 Tim. 6:1.)

“If any man teach otherwise, and consent not to wholesome words, even the words of our Lord Jesus Christ, and to [his] doctrine … doting about questions and strifes of words, … supposing that gain is godliness: from such withdraw thyself.” (1 Tim. 6:3–5; see also Isa. 29:21; 2 Ne. 27:32; D\&C 19:30; 38:41; 60:14.)

Dissecting doctrine in a controversial way in order to draw attention to oneself is not pleasing to the Lord. He declared:

“Bring to light the true points of my doctrine, yea, and the only doctrine which is in me.

“And this I do that I may establish my gospel, that there may not be so much contention; yea, Satan doth stir up the hearts of the people to contention concerning the points of my doctrine; and in these things they do err, for they do wrest the scriptures and do not understand them.” (D\&C 10:62–63.)

Contention fosters disunity. The Book of Mormon teaches the better way:

“Alma, having authority from God, … commanded them that there should be no contention one with another, but that they should look forward with one eye, having one faith and one baptism, having their hearts knit together in unity and in love one towards another.” (Mosiah 18:18, 21; see also Mosiah 23:15.)

\intalkheader{Steps to Supplant Contention}

What can we do to combat this canker of contention? What steps may each of us take to supplant the spirit of contention with a spirit of personal peace?

To begin, show compassionate concern for others. Control the tongue, the pen, and the word processor. Whenever tempted to dispute, remember this proverb: “He that is void of wisdom despiseth his neighbour: but a man of understanding holdeth his peace.” (Prov. 11:12; see also Prov. 17:28.)

Bridle the passion to speak or write contentiously for personal gain or glory. The Apostle Paul thus counseled the Philippians, “Let nothing be done through strife or vainglory; but in lowliness of mind let each esteem other better than themselves.” (Philip. 2:3.)

Such high mutual regard would then let us respectfully disagree without being disagreeable.

But the ultimate step lies beyond beginning control of expression. Personal peace is reached when one, in humble submissiveness, truly loves God. Heed carefully this scripture:

“There was no contention in the land, because of the love of God which did dwell in the hearts of the people.” (4 Ne. 1:15; italics added; see also 4 Ne. 1:2.)

Thus, love of God should be our aim. It is the first commandment—the foundation of faith. As we develop love of God and Christ, love of family and neighbor will naturally follow. Then will we eagerly emulate Jesus. He healed. He comforted. He taught, “Blessed are the peacemakers: for they shall be called the children of God.” (Matt. 5:9; see also 3 Ne. 12:9.)

Through love of God, the pain caused by the fiery canker of contention will be extinguished from the soul. This healing begins with a personal vow: “Let there be peace on earth, and let it begin with me.” (Sy Miller and Jill Jackson, “Let There Be Peace on Earth,” Beverly Hills, Calif.: Jan-Lee Music, 1972.) This commitment will then spread to family and friends and will bring peace to neighborhoods and nations.

Shun contention. Seek godliness. Be enlightened by eternal truth. Be like-minded with the Lord in love and united with Him in faith. Then shall “the peace of God, which passeth all understanding” (Philip. 4:7), be yours, to bless you and your posterity through generations yet to come. I so testify in the name of Jesus Christ, amen.

\newpage
\settalk{Woman—Of Infinite Worth}
\setspeaker{Russell M. Nelson}
\settalkdate{October 1989}
\talkheading{Woman—Of Infinite Worth}{Russell M. Nelson}

Words are inadequate to express the gratitude we feel for these brethren who have just been released from active assignment as General Authorities of the Church. We appreciate their marvelous ministry and powerful influence for good in the world.

Our privilege of sustaining leaders is granted by the Lord. (See D\&C 20:63–66; 26:2; 28:13; 38:34; 93:51; 104:21; 124:144.) Sustaining makes known to the Church who has authority (see D\&C 42:11) and enables each of us to show support. We honor all our leaders, both men and women, and are grateful for brothers and sisters so united in this kingdom of God on earth.

At a recent news conference in an Eastern European country, I was asked about the potential for women in The Church of Jesus Christ of Latter-day Saints. I replied that perhaps the Church does more to enlighten understanding about and to lift the cause of women than any other institution on earth. It provides the path to her eternal destiny.

A worthy woman personifies the truly noble and worthwhile attributes of life. A faithful woman can become a devoted daughter of God—more concerned with being righteous than with being selfish, more anxious to exercise compassion than to exercise dominion, more committed to integrity than to notoriety. And she knows of her own infinite worth.

Each faithful young woman in the Church proclaims that individual worth is one of her most cherished values. She declares, “I am of infinite worth with my own divine mission which I will strive to fulfill.” (Personal Progress, Salt Lake City: The Church of Jesus Christ of Latter-day Saints, 1989, p. 7.) Each daughter of God is of infinite worth because of her divine mission.

Important lessons about her divine mission may be learned from women of the scriptures. Mother Eve was a great example. She labored beside her husband as a partner. They both knew the plan of salvation. They both heeded commandments of obedience to God. Likewise, she prayed for divine guidance. She bore children. She taught the gospel to them. (See Moses 5:1–12; D\&C 138:39.)

Sarah, in becoming the mother of Isaac in her advanced years, verified that nothing is “too hard for the Lord.” (Gen. 18:14.)

Mary, mother of our Redeemer, was the perfect example of complete submission to the will of God. (See Luke 1:38.) She kept confidences. (See Luke 2:19.) In faith, she endured grief. (See John 20:11.)

The stories of these and other scriptural heroines show that women are essential in God’s plan for His children. Likewise, men have important but different assignments. We learn from the scriptures that men are to hold priesthood responsibilities and, as husbands and fathers, provide for (see Rom. 12:17; 1 Tim. 5:8; D\&C 75:28; 83:2, 4) and protect their families (see D\&C 134:6, 11). From the beginning, the priesthood has been conferred only upon men, coming through the “lineage of the fathers.” (See D\&C 84:6–16; 86:8; 107:40–41; Abr. 1:3–4.)

Blessings of the priesthood are shared by men and women. All may qualify for baptism and the gift of the Holy Ghost. All may take upon themselves the name of the Lord and partake of the sacrament. All may pray and receive answers to their prayers. Gifts of the Spirit and testimonies of the truth are bestowed regardless of gender. Men and women receive the highest ordinance in the house of the Lord together and equally, or not at all (see D\&C 131:1–3).

Opportunities for development of spiritual and intellectual potential are equal. Masculinity has no monopoly on the mind, and femininity has no exclusive dominion over the heart. The highest titles of human achievement—teacher, educated professional, loyal employee, faithful friend, student of the scriptures, child of God, disciple of Christ, trusted companion, loving parent—are earned under a uniform requirement of worthiness.

Loving parent. What a noble title! There are no greater roles in life for a man than those of husband and father. Likewise, there are no greater roles for a woman than those of wife and mother.

As I have watched my beloved companion and our dear daughters stretch to meet the challenging demands of these sacred roles, I have truly been inspired.

I have marveled at Mother’s forecast that a child would break out with measles or chicken pox well before my trained medical eye would have so predicted. I have watched the incredible calm of Mother when her child experienced a convulsive seizure.

So much is expected of a woman. Often she is a detective; she must search for lost articles and solve on a daily basis baffling “who done it” mysteries.

Frequently, her eager audience requests her talent as a musician by calling upon her to sing—any time, any place. She is an artist, using crayons and coloring books, needles and thread, or other means to create works of art by her own hands. (See Ps. 90:17; D\&C 42:40.)

During the early life of most children, a woman is the principal disciplinarian. She treads the tightrope of judgment between being too strict or too permissive.

Mother is “secretary of labor” for her home. She teaches the work ethic with its responsibilities and rewards. Father, too, shares that duty. I remember a hot Saturday afternoon years ago when one of our little daughters heard the chimes of an approaching ice-cream wagon. She asked me for some money. A bit heartlessly, I replied, “Sweetheart, why don’t you earn your money like everyone else does?” I’ll never forget her reply.

“But, Daddy,” she said, “I don’t like to work!” (Things are different now with four children of her own.)

A woman is a master communicator. And she communicates best in humble prayer. How many of us first learned to pray beside the bended knees of our mothers? Surely she knows that her children can walk alone only when they have found their pathway to Father in Heaven through prayer.

Certainly, a woman is a teacher. Someone said, “When you teach a boy, you teach an individual, but when you teach a girl, you teach a whole generation.”

J. Edgar Hoover said that “the cure of crime is not the electric chair but the high chair” (in Emerson Roy West, comp., Vital Quotations, Salt Lake City: Bookcraft, 1968, p. 78.)

I honor women who are not mothers. They know that motherhood is but one of the realms of womankind. The virtue and intelligence of women are uniquely applicable to other realms as well, such as compassionate service and teaching.

I am indebted to so many wonderful men and women who were my teachers. From grade school, I remember Miss Crow, Miss McLean, Miss Starr, and others. Later, Miss Bradford, Miss Cunningham, and Miss Snow were among my favorites. They were modest, gracious, and moral examples. They were not concerned with what I was to acquire, but with what I was to become. These marvelous unmarried teachers exerted an influence distinct from that of my angel mother. Their gleaming hopes, vicarious ambitions, and exacting demands were vitally important to my preparation for life.

A wise woman renews herself. In proper season, she develops her talents and continues her education. She musters the discipline to reach her goals. She dispels darkness and opens windows of truth to light her way.

A woman teaches priorities by precept and example. Recently I watched a television program in which a female lawyer was being interviewed. She was at home with her child on a full-time basis. When asked of her decision, she replied, “Oh, I may go back to the law sometime, but not now. For me, the issue is simple. Any lawyer could take care of my clients, but only I should be the mother of this child.”

Such a decision is made not in terms of rights but in terms of obligations and responsibilities. She knows that as she rises to meet responsibilities, rights will take care of themselves.

The Prophet Joseph Smith was taught this lesson while unjustly incarcerated in Liberty Jail—scene of so few rights, so little freedom, and so much abusive authority. As the Master tutored His Prophet regarding rights, those instructions were surrounded by schooling about obligations and responsibilities. (See D\&C 121.)

A righteous woman is a student of the scriptures. Many apply uniquely to her life. (See Gen. 27:46; Ps. 113:9; Prov. 31:10–31; Eph. 5:22–33; Col. 3:18; Titus 2:3–5; Jacob 3:7; Mosiah 4:14—16; D\&C 25.) In the scriptures she finds “great treasures of knowledge, even hidden treasures.” (D\&C 89:19.)

She need not have majored in physics to know divine truths, as “there is no such thing as immaterial matter.” (D\&C 131:7.) She need not have graduated in astronomy to learn lessons taught by God to Abraham—the relation of the earth to the sun, the sun to the planets, the planets to the center of the universe, and more. (See Abr. 3.) When she sings “Twinkle, twinkle, little star, how I wonder what you are,” she knows scriptural answers.

Well do I remember attending an important international symposium. The speaker was one of the world’s leading scholars, addressing a learned society with faculty representatives from major universities near and far. He said, in effect, that libraries of the world are being filled with case histories and anecdotal research. He pled for a significant change in direction. “What we need,” he said, “is research in things that really matter. We need to learn more about why we are here, where we have come from, and where we are going!”

That unforgettable address comes to mind each time I hear teachers at Primary and mothers at home enlighten children with revealed answers to these same questions. While so doing, they bear firm testimonies of the divinity of the Lord Jesus Christ.

Of course, there are times when a woman’s ability to endure is taxed to the limit. A teacher may have had enough of childish pranks, or a mother might be heard to say she’s “ready to resign.” She could become discouraged, especially if comparing herself unrealistically to others or focusing on what she is to do instead of on what she is to be.

Her self-esteem cannot be based on physical features, possession or lack of a particular talent, or comparative quantities of anything. Her self-esteem is earned by individual righteousness and a close relationship with God. Her outward glow is generated by goodness within. And her patience is much more apparent than any imperfection. (See D\&C 67:13.)

Sweet serenity is found in fervent prayer. Then, we forget ourselves and remember the reaching hands of the Savior, who said, “Come unto me, all ye that labour and are heavy laden, and I will give you rest.” (Matt. 11:28.) As our burdens are shared with Him, they do become lighter.

Feelings of worth come when a woman follows the example of the Master. Her sense of infinite worth comes from her own Christlike yearning to reach out with love, as He does.

When her husband, children, grandchildren, nieces, or nephews return from a day marred by the world’s rude realities, a loving woman can say, “Come unto me. I will give you rest.” Wherever she is can become a sanctified place, safe from the storms of life. Refuge is there because of her ability to nurture and to love unconditionally.

Sometimes this true love necessarily takes the tone of tough love. Her lessons of obedience and accountability must resemble those of her Master, who said, “If ye love me, keep my commandments.” (John 14:15; see also Ex. 20:6; Deut. 5:10; Mosiah 13:14; D\&C 46:9; 124:87.)

The Good Shepherd said, “Feed my lambs.” (John 21:15.) So a woman feeds her loved ones, providing succor and sustenance just as the Savior would do. Her divine gift is to nurture, to help the young, to care for the poor, to lift the brokenhearted.

The Lord said, “My work and my glory [is] to bring to pass the immortality and eternal life of man.” (Moses 1:39.) So His devoted daughter-disciple may truly say, “My work and my glory is to help my loved ones reach that heavenly goal.”

To help another human being reach one’s celestial potential is part of the divine mission of woman. As mother, teacher, or nurturing saint, she molds living clay to the shape of her hopes. In partnership with God, her divine mission is to help spirits live and souls be lifted. This is the measure of her creation. It is ennobling, edifying, and exalting.

Her saintly calling is opposed by Satan. He would shatter the family unit and demean the worth of woman. He would triumph if one man would offend or fail to honor her, or if one woman would deny her infinite worth and behave beneath her dignity. The vulgar portrayal of her beauty as an object of lust, the vile invasion of her private purity, should provoke righteous indignation from all caring people.

The gospel has been restored in these latter days so that the light of the Lord can prevail over efforts of the adversary. This day has long been prophesied. The Lord has promised worthy Saints of our time, “Upon the servants and upon the handmaids in those days will I pour out my spirit.” (Joel 2:29.)

A woman’s richest rewards will come as she rises to fulfill her destiny as a devoted daughter of God. To all faithful Saints He has promised thrones, kingdoms, principalities, glory, immortality, and eternal lives. (See Rom. 2:7; D\&C 75:5; 128:12, 23; 132:19.) That is the potential for women in The Church of Jesus Christ of Latter-day Saints. It is exalting, everlasting, and divine.

God bless us to honor each woman in her divine mission as a woman of infinite worth, I pray in the name of Jesus Christ, amen.

\newpage
\settalk{“Thus Shall My Church Be Called”}
\setspeaker{Russell M. Nelson}
\settalkdate{April 1990}
\talkheading{“Thus Shall My Church Be Called”}{Russell M. Nelson}

Today I would like to speak about a name. We are all pleased when our names are pronounced and spelled correctly. Sometimes a nickname is used instead of the real name. But a nickname may offend either the one named or the parents who gave the name.

The name of which I shall speak is not a personal name, yet the same principles apply. I refer to a name given by the Lord:

“Thus shall my church be called in the last days, even The Church of Jesus Christ of Latter-day Saints” (D\&C 115:4).

Note carefully the language of the Lord. He did not say, “Thus shall my church be named.” He said, “Thus shall my church be called.” Years ago, members were cautioned by the Brethren, who wrote, “We feel that some may be misled by the too frequent use of the term ‘Mormon Church’” (Member-Missionary Class, Instructor’s Guide [Salt Lake City: The Church of Jesus Christ of Latter-day Saints, 1982], p. 2). Before any other name is considered to be a legitimate substitute, the thoughtful person might reverently consider the feelings of the Heavenly Parent who bestowed that name.

Surely every word that proceeds from the mouth of the Lord is precious. So each word in this name must be important—divinely designated for a reason. If we study the key words in that name, we can better understand the name’s full significance.

\intalkheader{Saints}

The last word in the title is Saints. I smile when I remember a comment made after my call to the Quorum of the Twelve. A doctor friend relayed a report made at a professional meeting that “Dr. Nelson was no longer practicing cardiac surgery because his church had made him ‘a saint.’”

Such a comment was not only amusing but revealing. It evidenced unfamiliarity with the language of the Bible, in which the word saint is used much more frequently than is the term Christian.

The word Christian appears in only three verses of the King James Version of the Bible. One verse describes the historical fact that “disciples were called Christians first in Antioch” (Acts 11:26); another quotes a sarcastic nonbeliever, King Agrippa (see Acts 26:28); and the third indicates that one known as “a Christian” must be prepared to suffer (1 Pet. 4:16).

In contrast, the term saint (or saints) appears in thirty-six verses of the Old Testament and in sixty-two verses of the New Testament.

Paul addressed an epistle “to the saints which are at Ephesus, and to the faithful in Christ Jesus” (Eph. 1:1).

To recent converts there he said, “Ye are no more strangers and foreigners, but fellowcitizens with the saints, and of the household of God” (Eph. 2:19; see also Eph. 3:17–19).

In his epistle to the Ephesians, Paul used the word saint at least once in every chapter!

Despite its use in ninety-eight verses of the Bible, the term saint is still not well understood. Some mistakenly think that it implies beatification or perfection. Not so! A saint is a believer in Christ and knows of His perfect love. The giving saint shares in a true spirit of that love, and the receiving saint accepts in a true spirit of gratitude. A saint serves others, knowing that the more one serves, the greater the opportunity for the Spirit to sanctify and purify.

A saint is tolerant and is attentive to the pleadings of other human beings, not only to spoken messages but to unspoken messages as well. A saint is different from an individual whose response to a concern might be a selfish “What do I care” attitude. A real saint responds, “What? I do care!” Do is an action verb, and it becomes the driving force in the reply of one who will care for another in need (see 1 Cor. 12:25–27; 2 Cor. 7:12).

A saint “refrain[s] from idleness” (Alma 38:12) and seeks learning by study and also by faith. Education not only helps in communication with others, it enables one to discern truth from error, particularly through studying the scriptures (see D\&C 88:118).

A saint is honest and kind, paying financial obligations promptly and fully, treating others as she or he would want to be treated (see Matt. 7:12; 3 Ne. 14:12; D\&C 112:11).

A saint is an honorable citizen, knowing that the very country which provides opportunity and protection deserves support, including prompt payment of taxes and personal participation in its legal political process (see D\&C 134:5).

A saint resolves any differences with others honorably and peacefully and is constant in courtesy—even in traffic at the rush hour.

A saint shuns that which is unclean or degrading and avoids excess even of that which is good.

Perhaps above all, a saint is reverent. Reverence for the Lord, for the earth He created, for leaders, for the dignity of others, for the law, for the sanctity of life, and for chapels and other buildings are all evidences of saintly attitudes (see Lev. 19:30; Alma 47:22; D\&C 107:4; 134:7).

A reverent saint loves the Lord and gives highest priority to keeping His commandments. Daily prayer, periodic fasting, and payment of tithes and offerings are privileges important to a faithful saint.

Finally, a saint is one who receives the gifts of the Spirit that God has promised to all His faithful sons and daughters (see Joel 2:28–29; Acts 2:17–18).

\intalkheader{Latter-day}

The term latter-day is an expression especially difficult for translators who labor in languages in which there is not a good equivalent term. Some translations may suggest last day.

It is true that scriptures foretell the final days of the earth’s temporal existence as a telestial sphere. The earth will then be renewed and receive its paradisiacal, or terrestrial, glory (see A of F 1:10). Ultimately, the earth will become celestialized (see Rev. 21:1; D\&C 77:1; 88:25–26). But its last days must be preceded by its latter days!

We live in those latter days, and they are really remarkable. The Lord’s Spirit is being poured out upon all inhabitants of the earth, precisely as the Prophet Joel foretold. His prophecy was of such significance that the angel Moroni reaffirmed it to the Prophet Joseph Smith (see Joel 2:28–32; JS—H 1:41).

For millennia, methods of farming, travel, and communication were largely unchanged from ancient techniques. Developments since the birth of Joseph Smith, however, have risen in remarkable contrast.

Joseph Smith had long been foreordained as God’s prophet for the restoration of the gospel in the fulness of times (see 2 Ne. 3:7–15). Twenty-five years after his birth, The Church of Jesus Christ of Latter-day Saints was officially organized.

Later in that same century, the telegraph was developed, the Atlantic Ocean was first crossed by a steamship, and the telephone, the automobile, and motion pictures were invented.

The twentieth century has been even more extraordinary. Farming has become mechanized.

Modern transportation allows travel to nearly any destination in the world within a day or two.

Computers have been developed that allow the Church to serve living members and to organize information relative to progenitors who live on the other side of the veil. People throughout the world, once little concerned with family history, now search for roots of their ancestral heritage using technologies unavailable a century ago.

Long-distance telephone, telefax, radio, television, and satellite communications have become routine. In these latter days it is possible for the word of the Lord to be broadcast from world headquarters of His Church and heard in the most remote areas of the globe.

The divine promise is being fulfilled that this restored “gospel shall be preached unto every nation, and kindred, and tongue, and people” (D\&C 133:37).

Political changes have occurred recently in many countries. Previous restrictions of personal liberties have been relieved. The shell of spiritual confinement has been shattered. Swelling shouts of freedom fill the air. Surely the hand of the Lord is apparent. He said, “I will hasten my work in its time” (D\&C 88:73), and that time of hastening is now.

\intalkheader{Jesus Christ}

By divine directive, the title of the Church bears the sacred name of Jesus Christ, whose church this is (see D\&C 115:3–4). He so decreed more than once. Nearly two thousand years ago, the Lord said, “Ye shall call the church in my name; …

“And how be it my church save it be called in my name?” (3 Ne. 27:7–8; italics added).

We worship God the Eternal Father in the name of His Son by the power of the Holy Ghost. We know the premortal Jesus to be Jehovah, God of the Old Testament. We know Him to be “the chief corner stone” upon which the organization of His Church is based (Eph. 2:20). We know Him to be the Rock from whom revelation comes to His authorized agents (see 1 Cor. 10:4; Hel. 5:12) and to all who worthily seek Him (see D\&C 88:63).

We know that He came into the world to do the will of His Father, who sent Him (see 3 Ne. 27:13). His divine mission was to effect the Atonement, which was to break the bands of death and enable us to receive immortality and eternal life.

The living Lord’s divine mission still continues. One day we will stand before Him in judgment. He has foretold that event:

“Whoso repenteth and is baptized in my name shall be filled; and if he endureth to the end, behold, him will I hold guiltless before my Father at that day when I shall stand to judge the world” (3 Ne. 27:16).

We revere the name of Jesus Christ. He is our risen Redeemer.

\intalkheader{The Church}

The first two words of the name the Lord chose for His earthly organization are The Church.

Note that the article The begins with a capital letter. This is an important part of the title, for the Church is the official organization of baptized believers who have taken upon themselves the name of Christ (see D\&C 10:67–69; 18:21–25).

The foundation of the Church is the reality that God is our Father and that His Only Begotten Son, Jesus Christ, is the Savior of the world. The witness and inspiration of the Holy Ghost confirm those realities.

The Church is the way by which the Master accomplishes His work and bestows His glory. Its ordinances and related covenants are the crowning rewards of our membership. While many organizations can offer fellowship and fine instruction, only His church can provide baptism, confirmation, ordination, the sacrament, patriarchal blessings, and the ordinances of the temple—all bestowed by authorized priesthood power. That power is destined to bless all children of our Heavenly Father, regardless of their nationality:

“The keys of the kingdom of God are committed unto man on the earth, and from thence shall the gospel roll forth unto the ends of the earth” (D\&C 65:2; see also Dan. 2:37–45; D\&C 109:72).

Admission into His church is by baptism. This sacred ordinance is reserved only for children after they reach the age of accountability and for adults who are truly converted, prepared, and worthy to pass this scriptural test:

“Ye are desirous to come into the fold of God, and to be called his people, and are willing to bear one another’s burdens, that they may be light;

“Yea, and are willing to mourn with those that mourn; yea, and comfort those that stand in need of comfort, and to stand as witnesses of God at all times and in all things, and in all places that ye may be in” (Mosiah 18:8–9).

Through the ordinance of baptism, we take upon ourselves the name of the Lord and covenant to be saints in these latter days. We covenant to live by the doctrines of the Church as recorded in sacred scriptures and as revealed to prophets, ancient and modern.

“We believe all that God has revealed, all that He does now reveal, and we believe that He will yet reveal many great and important things pertaining to the Kingdom of God” (A of F 1:9).

These revelations include fundamental truths essential to our everlasting happiness and joy. They teach of priorities with eternal potential, such as love of God, family, mother, father, children, and home; self-mastery; care of the poor and needy; service; and thoughtful consideration for others.

This church, established under the direction of Almighty God, fulfills promises made in biblical times. It is part of the “restitution of all things, which God hath spoken by the mouth of all his holy prophets since the world began” (Acts 3:21). It has been restored and given a name by the Lord Himself.

He issued this solemn warning: “Let all men beware how they take my name in their lips.” “Remember,” He added, “that which cometh from above is sacred, and must be spoken with care” (D\&C 63:61, 64). Therefore, just as we revere His holy name, we likewise revere the name that He decreed for His church.

As members of His church, we are privileged to participate in its divine destiny. May we so honor Him who declared, “Thus shall my church be called … The Church of Jesus Christ of Latter-day Saints,” I pray in the sacred name of Jesus Christ, amen.

\newpage
\settalk{Choices}
\setspeaker{Russell M. Nelson}
\settalkdate{October 1990}
\talkheading{Choices}{Russell M. Nelson}

Not long ago a beautiful young mother asked me for guidance with a very difficult decision she was facing. It pertained to an important surgical operation that was being considered. Consequences of her choice would affect her husband and her family as well. She said, “Decisions are really hard for me. I even have trouble choosing what to wear each morning.”

“You are not so different,” I replied. “Each of us must make choices. That is one of life’s great privileges.”

I told this lovely mother that my fellow physicians are regularly asked questions about the human body. Some questions relate to surgical intervention to save a life or to save a part of the body. Other questions relate to elective procedures to alter the body’s structure or function. In recent years, many questions relate to the “choice” to abort the life of a newly forming human being. Ironically, such “choice” would deny that developing individual both life and choice.

I reminded her that questions regarding our bodies represent only an important fraction of life’s most challenging choices. Others include “Where shall I live?” “What shall I do with my life?” “To which cause should I commit my effort and my good name?” These are but a few of the many choices that we must make each day.

I will not disclose the name of the sister, nor the specific operation she was contemplating. To do so might divert our attention to a specific topic and away from those fundamental principles that pertain to important decisions generally.

Because challenging choices face all of us from time to time, I invite others to join with us as I extend my conversation with this young mother.

I would suggest three questions you might ask yourself as you consider your options. Whether they are once-in-a-lifetime or routine daily decisions, serious reflection on these three questions will help clarify your thinking. You might wish to review these questions first alone and then with your husband. They are:

Truthful answers to these three questions will remind you of important anchors and unchanging principles.

As you consider these fundamental questions, it will become clear that decisions you first thought to be purely personal virtually always impact the lives of others. In answering these questions, then, you must be mindful of the broader circle of family and friends who will be affected by the consequences of your choice. This self-evaluation will be a silent examination. No one else will hear your replies. Though I will suggest some answers, the ultimate responses must be uniquely yours.

\intalkheader{“Who Am I?”}

Remember, you are a daughter of God, just as your husband is a son of God. Our Heavenly Father loves you. He has created you to be successful and to have joy.

“He created man, male and female, after his own image and in his own likeness” (D\&C 20:18; see also Gen. 1:26–27; Mosiah 7:27; Alma 18:34; 22:12; Ether 3:15; Moses 2:27).

These bodies, created in God’s image, are to be preserved, protected, and well cared for. I feel as did the Apostle Paul, who likened the human body to a temple:

“Know ye not that ye are the temple of God, and that the Spirit of God dwelleth in you?

“If any man defile the temple of God, him shall God destroy; for the temple of God is holy, which temple ye are” (1 Cor. 3:16–17).

You are one of God’s noble and great spirits, held in reserve to come to earth at this time (see D\&C 86:8–11). In your premortal life you were appointed to help prepare the world for the great gathering of souls that will precede the Lord’s second coming. You are one of a covenant people. You are an heir to the promise that all the earth will be blessed by the seed of Abraham and that God’s covenant with Abraham will be fulfilled through his lineage in these latter days (see 1 Ne. 15:18; 3 Ne. 20:25).

As a member of the Church, you have made sacred covenants with the Lord. You have taken upon yourself the name of Christ (see D\&C 18:28; 20:29, 37). You have promised to always remember Him and to keep His commandments. In return, He has agreed to grant His Spirit to be with you (see Moro. 4:3; 5:2; D\&C 20:77).

Having briefly considered some answers to question number one, let us turn our attention to question number two.

\intalkheader{“Why Am I Here?”}

This question is one I have often asked myself. Well do I remember doing so many years ago while in military service, separated from family and friends, surrounded by the horrible devastation of war. On another unforgettable occasion, I was stranded in a cold, remote area, far from transportation, food, or shelter. No doubt you have had similar anxious moments. But those experiences are the exceptions. I’d like to discuss the greater question.

Why are you here on planet earth?

One of the most important reasons is to receive a mortal body. Another is to be tested—to experience mortality—to determine what you will do with life’s challenging opportunities. Those opportunities require you to make choices, and choices depend on agency. A major reason for your mortal existence, therefore, is to test how you will exercise your agency (see 2 Ne. 2:15, 25).

Agency is a divine gift to you. You are free to choose what you will be and what you will do. And you are not without help. Counsel with your parents is a privilege at any age. Prayer provides communication with your Heavenly Father and invites the promptings of personal revelation. And in certain circumstances, consultation with professional advisers and with your local leaders in the Church may be highly advisable, especially when very difficult decisions must be made.

That is precisely the pattern chosen by President Spencer W. Kimball. In 1972, Elder Kimball, then a member of the Council of the Twelve, knew that his mortal life was slipping away because of heart disease. He obtained competent medical counsel and prayerfully consulted with the Lord and with his file leaders in the Church. Elder and Sister Kimball and the First Presidency carefully weighed available alternatives. Then President Harold B. Lee, speaking for the First Presidency, counseled Elder Kimball. With great conviction, President Lee said: “Spencer, you have been called! You are not to die! You must do everything you need to do to care for yourself and continue to live” (“Spencer W. Kimball: Man of Faith,” Ensign, Dec. 1985, p. 40).

President Kimball chose to have an operation performed upon his heart that was known to carry a high risk. He was blessed with a successful result. He lived thirteen more years, eventually to succeed President Lee as President of the Church.

That precious privilege of choice—man’s agency—was decreed before the world was created (see D\&C 93:29–31). It is a moral agency (see D\&C 101:78). Thus, it was opposed by Satan (see Moses 4:3) but affirmed by the Lord (see Moses 4:2) and reaffirmed through prophets in ancient and in modern times (see D\&C 58:26–28; Moses 6:56; 7:32).

The proper exercise of moral agency requires faith. Faith in the Lord Jesus Christ is the first principle of the gospel (see A of F 1:4). Because of Him, you have your agency. He must be the very foundation of your faith, and the testing of that faith is a fundamental reason for your freedom to choose.

You are free to develop and exercise faith in God and in His divine Son, faith in His word, faith in His Church, faith in His servants, and faith in His commandments.

Facing difficult challenges is neither new nor unique. Centuries ago, Joshua spoke of a choice his family faced. He declared:

“Choose you this day whom ye will serve; … as for me and my house, we will serve the Lord” (Josh. 24:15; see also Moses 6:33).

Cultivation of that faith will entitle you to the companionship of the Holy Ghost, who will help you make wise decisions (see 2 Ne. 2:27–28; D\&C 14:8).

Many may profess a measure of faith in God, but without sincere repentance, faith cannot be fully operative. This concept was made known to the Nephites:

“Many of them … are brought to the knowledge of the truth, … and are led to believe the holy scriptures, … which leadeth them to faith on the Lord, and unto repentance, which faith and repentance bringeth a change of heart unto them” (Hel. 15:7).

Faith, repentance, and obedience will qualify you for sublime gifts of justice and mercy, which are bestowed upon those worthy of the blessings of the Atonement (see Alma 34:16–17).

Yes, every test, every trial, every challenge and hardship you endure is an opportunity to further develop your faith (see D\&C 63:11; 101:4).

Faith can be fortified through prayer. Prayer is the powerful key to making decisions, not only concerning your physical body, but concerning all other important aspects of your life. Humbly seek the Lord in prayer with a sincere heart and real intent, and He will help you (see Alma 33:23; Moro. 7:9; 10:4; D\&C 9:7–9).

Remember that faith and prayer alone are seldom sufficient. Personal effort is usually necessary to accomplish your heart’s desire. “Faith, if it hath not works, is dead, being alone” (James 2:17; see also vs. 18, 20, 26; Alma 26:22).

The answers to question number two emphasize that you are here to exercise faith, to pray, and to work hard.

Now let’s turn our attention to question number three.

\intalkheader{“Where Am I Going?”}

This question reminds us that eventually you (and I) are going to die, be resurrected, be judged, and be awarded a place in eternal realms (see 1 Cor. 15:22; Alma 12:24; 21:9; Hel. 14:16–17; D\&C 138:19). With each passing sunset, you are closer to that inevitable day of judgment. Then you will be asked to account for your faith, your hopes, and your works. The Lord said:

“Every man may act in doctrine and principle … according to the moral agency which I have given unto him, that every man may be accountable for his own sins in the day of judgment” (D\&C 101:78; see also Mosiah 3:24).

As all will be resurrected, your physical body will then be restored to its proper and perfect frame (see Alma 11:43; 40:23). The day of your resurrection will be a day of judgment that will determine the kind of life you shall have hereafter.

That judgment will consider not only your actions, but also your innermost intent and heartfelt desires. Your everyday thoughts have not been lost. Scriptures speak of the “bright recollection” (Alma 11:43) and “perfect remembrance” (Alma 5:18) that your mind will provide in times of divine judgment.

The Lord knows the desires of our hearts. At the time of judgment, surely the special yearnings of single sisters and childless couples, for example, will be given compassionate consideration by Him who said:

“I, the Lord, will judge all … according to their works, according to the desire of their hearts” (D\&C 137:9; see also Heb. 4:12; Alma 18:32; D\&C 6:16; 33:1; 88:109).

He will know of your longings as a wife and mother who tried diligently to serve your family and society properly.

As I listen to those who argue for causes contrary to the commandments of God and observe individuals who revel in the pleasures of the world with apparent disregard for eventual judgment, I think of this divine description of their folly:

“They despised my judgments, and walked not in my statutes, … for their heart went after their idols” (Ezek. 20:16).

Interviews, as for temple recommends, with your bishop and members of your stake presidency are precious experiences. And, in a way, they could be considered meaningful “dress rehearsals” for that grand colloquy when you will stand before the Great Judge.

After the Resurrection and Judgment, you will be assigned to your everlasting home on high. The revelations liken the glory of those dwelling places to the differing lights of heavenly bodies. Paul said:

“There is one glory of the sun, and another glory of the moon, and another glory of the stars” (1 Cor. 15:41).

The Lord revealed more to the Prophet Joseph Smith, who wrote of the telestial glory, where those will ultimately abide “who received not the gospel of Christ, neither the testimony of Jesus” while in this life (D\&C 76:82).

The Prophet taught of the terrestrial glory as the abode for the “honorable … of the earth, who were blinded by the craftiness of men,” who rejected the gospel while on the earth (D\&C 76:75).

And then he wrote of the celestial glory, which “glory is that of the sun, even the glory of God, the highest of all” (D\&C 76:70). There the faithful will dwell together with their families, enjoying exaltation with our Heavenly Father and His Beloved Son. With them will be those who have been obedient to ordinances and covenants made in holy temples, where they were sealed to predecessors and posterity.

As you continue to face many challenging choices in life, remember, there is great protection when you know who you are, why you are here, and where you are going. Let your unique identity shape each decision you make on the path toward your eternal destiny. Accountability for your choices now will bear on all that lies ahead.

May each of us choose wisely and with faith in Him who created us, I pray in the name of Jesus Christ, amen.

\newpage
\settalk{Listen to Learn}
\setspeaker{Russell M. Nelson}
\settalkdate{April 1991}
\talkheading{Listen to Learn}{Russell M. Nelson}

In his invocation for this session of conference, Elder Hugh W. Pinnock prayed that we might listen carefully. Many articles in Church literature have dealt with the important art of listening. They support a proverb that teaches this vital lesson: “Hear counsel, and receive instruction, that thou mayest be wise” (Prov. 19:20). Surely wisdom will come as we listen to learn from children, parents, partners, neighbors, Church leaders, and the Lord.

\intalkheader{Children}

Parents and teachers, learn to listen, then listen to learn from children. A wise father once said, “I do a greater amount of good when I listen to my children than when I talk to them.”

When our youngest daughter was about four years of age, I came home from hospital duties quite late one evening. I found my dear wife to be very weary. I don’t know why. She only had nine children underfoot all day. So I offered to get our four-year-old ready for bed. I began to give the orders: “Take off your clothes; hang them up; put on your pajamas; brush your teeth; say your prayers” and so on, commanding in a manner befitting a tough sergeant in the army. Suddenly she cocked her head to one side, looked at me with a wistful eye, and said, “Daddy, do you own me?”

She taught me an important lesson. I was using coercive methods on this sweet soul. To rule children by force is the technique of Satan, not of the Savior. No, we don’t own our children. Our parental privilege is to love them, to lead them, and to let them go.

The time to listen is when someone needs to be heard. Children are naturally eager to share their experiences, which range from triumphs of delight to trials of distress. Are we as eager to listen? If they try to express their anguish, is it possible for us to listen openly to a shocking experience without going into a state of shock ourselves? Can we listen without interrupting and without making snap judgments that slam shut the door of dialogue? It can remain open with the soothing reassurance that we believe in them and understand their feelings. Adults should not pretend an experience did not happen just because they might wish otherwise.

Even silence can be misinterpreted. A story was written of “a little boy [who] looked up at his mother and said, ‘Why are you mad at me?’ She answered, ‘I’m not angry at you. What makes you say that?’ ‘Well, your hands are on your hips, and you are not saying anything.’”

Parents with teenage youth may find that time for listening is often less convenient but more important when young people feel lonely or troubled. And when they seem to deserve favor least, they may need it most.

Wise parents and teachers, listen to learn from children.

\intalkheader{Parents}

Children of all ages, learn to listen, and listen to learn from parents, as Elder Oaks taught us this morning. Spiritually or physically, it can be a matter of life and death.

Several years ago, I was invited to give an important lecture at a medical school in New York City. The night before the lecture, Sister Nelson and I were invited to dinner at the home of our host professor. There he proudly introduced us to an honor medical student—his beautiful daughter.

Some weeks later, that professor telephoned me in an obvious state of grief. I asked, “What is the matter?”

“Remember our daughter whom you met at our home?”

“Of course,” I replied. “I’ll never forget such a stunning young lady.”

Then her father sobbed and said, “Last night she was killed in an automobile accident!” Trying to gain composure, he continued: “She asked permission to go to a dance with a certain young man. I didn’t have a good feeling about it. I told her so and asked her not to go. She asked, ‘Why?’ I simply told her that I was uneasy. She had always been an obedient daughter, but she said that if I could not give her a good reason to decline, she wanted to go. And so she did. At the dance, alcoholic beverages were served. Her escort drank a bit—we don’t know how much. While returning home, he was driving too fast, missed a turn, and careened through a guardrail into a reservoir below. They were both submerged and taken to their death.”

As I shared my feeling of sadness, he concluded: “My grief is made worse because I had the distinct feeling that trouble lay ahead. Why couldn’t I have been more persuasive?”

This experience will not have been in vain if others can listen and learn from it. Children, honor your parents, even when they cannot give a satisfactory explanation for their feelings. Please have faith in this scripture, which applies to all age-groups: “Hear the instruction of thy father, and forsake not the law of thy mother” (Prov. 1:8).

Parents have a divine duty to teach their children to love the Lord. Children have an equal obligation to “obey [their] parents in the Lord” (Eph. 6:1).

Wise children, listen to learn from parents.

\intalkheader{Partners}

Husbands and wives, learn to listen, and listen to learn from one another. I was amused to read of an experience recorded by Elder F. Burton Howard in his biography of President Marion G. Romney: “His good-humored love for Ida was manifested in many ways. He delighted in telling of her hearing loss. ‘I once went to see a doctor about her hearing,’ he would say. ‘He asked me how bad it was, and I said I didn’t know. He told me to go home and find out. The doctor instructed me to go into a far room and speak to her. Then I should move nearer and nearer until she does hear. Following the doctor’s instructions, I spoke to her from the bedroom while she was in the kitchen—no answer. I moved nearer and spoke again—no answer. So I went right up to the door of the kitchen and said, “Ida, can you hear me?” She responded, “What is it, Marion—I’ve answered you three times.”’”

Even with normal hearing, some couples seem not to listen to one another. Taking time to talk is essential to keep lines of communication intact. If marriage is a prime relationship in life, it deserves prime time! Yet less important appointments are often given priority, leaving only leftover moments for listening to precious partners.

Keeping the garden of marriage well cultivated and free from weeds of neglect requires the time and commitment of love. It is not only a pleasant privilege, it is a scriptural requirement with promise of eternal glory.

Wise partners, listen to learn from one another.

\intalkheader{Neighbors}

Learn to listen, and listen to learn from neighbors. Repeatedly the Lord has said, “Thou shalt love thy neighbour” (Lev. 19:18; Matt. 19:19). Opportunities to listen to those of diverse religious or political persuasion can promote tolerance and learning. And a good listener will listen to a person’s sentiments as well. I learned much from Brother David M. Kennedy as we met with many dignitaries in nations abroad. When one of them spoke, Brother Kennedy not only looked eye to eye and listened with real intent, but he even removed his reading glasses, as if to show that he wanted nothing in the way of his total concentration.

The wise listen to learn from neighbors.

\intalkheader{Church Leaders}

Members, learn to listen, and listen to learn from Church leaders. Faithful members love the Savior and honor His servants, having faith in the Lord’s declaration that “whether by mine own voice or by the voice of my servants, it is the same” (D\&C 1:38).

One day in Italy I met a wonderful priesthood leader and his wife. In him I saw a man with great potential. But my language was foreign to them. Through an interpreter, I challenged them to study the English language. They listened obediently and studied diligently. Now six years later, ably sustained by his wife, Carolina, Vincenzo Conforte is faithfully serving his second assignment as a mission president, interviewing missionaries well in Italian or in English.

President Ezra Taft Benson has proclaimed the importance of studying the Book of Mormon. People throughout the earth are being blessed as they follow this and other counsel he has given.

Gratefully we thank God for a prophet to guide us in these latter days. But many turn a deaf ear to his teachings, oblivious to his prophetic position. They do so at great risk, for scriptures contain this warning:

“A prophet shall the Lord your God raise up unto you … ; him shall ye hear in all things whatsoever he shall say unto you. …

“Every soul, which will not hear that prophet, shall be destroyed from among the people” (Acts 3:22–23).

President J. Reuben Clark, Jr., said, “We do not lack a prophet; what we lack is a listening ear.” Words of the Lord are taught by His disciples (see D\&C 1:4). Wise members listen to learn from Church leaders.

\intalkheader{The Lord}

Above all, God’s children should learn to listen, then listen to learn from the Lord. On several sacred occasions in the world’s history, our Heavenly Father has personally appeared to introduce His divine Son with a specific charge to “hear him.”

Jesus taught this first and great commandment: “Thou shalt love the Lord thy God with all thy heart, and with all thy soul, and with all thy mind” (Matt. 22:37).

Scriptures recorded in all dispensations teach that we show our love of God as we hearken to His commandments and obey them. These actions are closely connected. In fact, the Hebrew language of the Old Testament in most instances uses the same term for both hearkening (to the Lord) and obedience (to His word).

In addition to hearing the word of the Lord obediently, we manifest our love for God through prayer. And listening is an essential part of prayer. Answers from the Lord come ever so quietly. Hence He has counseled us to “be still and know that I am God” (D\&C 101:16).

President Spencer W. Kimball said, “It would not hurt us, either, if we paused at the end of our prayers to do some intense listening—even for a moment or two—always praying, as the Savior did, ‘not my will, but thine, be done’ (Luke 22:42).”

In a world scarred by scourges of tyranny and war, many of its inhabitants earnestly pray for inner peace. For example, not long ago a beautiful young mother named Svetlana developed an intense desire to obtain a Bible. But in her city of Leningrad, a Bible was very rare and expensive. Frequently and fervently she prayed for a Bible. Ultimately, she and her husband were impressed to travel with their small child to Helsinki, Finland, with that hope in mind. There one day while walking in a park, she stumbled across an object buried under the cover of autumn leaves. She picked it up and found it to be a Bible written in the Russian language! Excitedly she recounted the story of this great discovery to another mother who was also in the park with her youngster. The second mother rejoiced with Svetlana and added, “Would you like to have another book about Jesus Christ?” Svetlana, of course, answered in the affirmative. The other mother provided Svetlana with a copy of a Russian-language edition of the Book of Mormon and invited the family to church. She eagerly embraced the teachings of the missionaries and shortly thereafter joined The Church of Jesus Christ of Latter-day Saints. Then they returned to their home, where they have helped pioneer the work in the Leningrad Branch of the Church.

Her experience typifies this promise of the Savior to those who seek Him: “Thus saith the Lord God: I will give unto the children of men line upon line, precept upon precept, here a little and there a little; and blessed are those who hearken unto my precepts, and lend an ear unto my counsel, for they shall learn wisdom; for unto him that receiveth I will give more” (2 Ne. 28:30; italics added).

While stressing the importance of listening well, I am mindful of those who cannot hear. Many labeled as deaf have received the Spirit “by the hearing of faith” (Gal. 3:2). The example of Rachel Ivins Grant is inspiring to me. She never complained about her own deafness. Though most women in their seventies would be completely worn out while rearing six growing children of another mother, she undertook that task. Rachel’s deafness seemed to save her from the wear and tear of noise. Sometimes, when two were arguing, Rachel would burst out laughing. She said they had no idea how funny it was to see their angry faces and hear none of their words.

Before her son, Heber J. Grant, became the seventh President of the Church, she declared, “Of course the greatest trial I have is that I cannot hear, but I have so many blessings I cannot complain, but if we only will live so that we may receive the instructions of God, there is nothing we are called to pass through but will be for our good.”

The Redeemer loves such faithful souls: “For the eyes of the Lord are over the righteous, and his ears are open unto their prayers” (1 Pet. 3:12).

They qualify for this prophetic promise: “Before they call, I [the Lord] will answer; and while they are yet speaking, I will hear” (Isa. 65:24).

To all of God’s children, either able to hear or deaf to mortal sound, He offers this reward: “Incline your ear, and come unto me: hear, and your soul shall live” (Isa. 55:3).

Your soul will be blessed as you learn to listen, then listen to learn from children, parents, partners, neighbors, and Church leaders, all of which will heighten capacity to hear counsel from on high.

Carefully listen to learn from the Lord through the still small voice—the Holy Spirit—which leads to truth. Listen to learn by studying scriptures that record His holy mind and will. Listen to learn in prayer, for He will answer the humble who truly seek Him.

The wise listen to learn from the Lord. I testify of Him and certify that as we “hearken and … hear the voice of the Lord,” we will be blessed, “for the hour of his coming is nigh” (D\&C 133:16–17), in the name of Jesus Christ, amen.

\newpage
\settalk{“These … Were Our Examples”}
\setspeaker{Russell M. Nelson}
\settalkdate{October 1991}
\talkheading{“These … Were Our Examples”}{Russell M. Nelson}

In June of this year, Sister Nelson and I had the great privilege of accompanying the Mormon Tabernacle Choir on its historic concert tour in Europe. We are grateful to the First Presidency for that assignment. Much has been written regarding the success of the choir and of its favorable influence that will yet continue. Members and friends of the Church worldwide join with me in expressing appreciation to the officers, staff, directors, accompanists, and all vocalists for their wonderful service. I won’t mention anyone by name; I will simply refer to them all as members of the choir.

I’ll not comment as a music critic would. While musical experts of the world focus on what choir members can do, I would like to focus on what choir members can be. This I do because I witnessed in choir members great examples that can inspire and improve the lives of each of you who honestly strives to emulate the Lord who said, “I am the light; I have set an example for you” (3 Ne. 18:16; see also John 13:15). So we should strive to learn from His example—and from the good example of those who love and follow Him.

Members of the Mormon Tabernacle Choir are not superhuman. They are ordinary people with ordinary frailties. But therein lies the power of their example. They believe this promise from the Lord: “Because thou hast seen thy weakness, thou shalt be made strong” (D\&C 135:5).

Before attending their first rehearsal, choir members brought with them not only musical education and talent but qualities of personal righteousness. Before singing to their first audience, they were blessed by another promise from our Savior:

“Ye may show forth good examples unto them in me, and I will make an instrument of thee in my hands unto the salvation of many souls” (Alma 17:11).

Have you not learned that strength comes to an ordinary soul when given an extraordinary calling? The choir has! Indeed, each member seemed to be imbued with a real sense of mission, striving for those ten traits that missionaries are expected to possess and practice:

“Faith, virtue, knowledge, temperance, patience, brotherly kindness, godliness, charity, humility, diligence” (D\&C 4:6).

Those are attributes the Lord deserves from His disciples. Because each of us is to be “an example of the believers” (1 Tim. 4:12), I would like to address those ten topics as members of the Mormon Tabernacle Choir “were our examples” (1 Cor. 10:6) in many memorable ways.

\intalkheader{Faith}

Their great faith was strengthened by the faith of our leaders. I pay tribute to the First Presidency and to leaders of the choir who had the foresight to plan as they did and when they did. How bold and inspired they were to conceive this tour many months—even years—before Europe’s unwelcoming walls began to crumble! The Brethren had the faith to believe that the choir could sing in cities such as Warsaw, Budapest, Prague, Leningrad, and Moscow long before such dreams seemed plausible. Then in January 1991, hopeful plans were seriously threatened when war erupted in the Persian Gulf. Even then, our leaders decided against canceling the tour. They knew of its potential for good and had faith that countless obstacles could be overcome. Often they prayed that the choir’s tour might be successfully accomplished.

Those prayers were answered!

Think of the timing. In one thousand years of Russia’s existence, its first popular national election ever to be held occurred in June 1991. Six days later, the Mormon Tabernacle Choir performed in Moscow! That very night, after the strains of “Come, Come, Ye Saints” (see Hymns, 1985, no. 30) had resounded from the Bolshoi Theater, the vice president of the republic announced that The Church of Jesus Christ of Latter-day Saints had been granted recognition in the Republic of Russia. On the eve of a supreme crisis that was yet ahead, Russian people heard songs of faith, courage, hope, and love.

That faith of our Church leaders filled the hearts of members, families, and friends of the choir as well. Real were the risks of separation of husbands and wives, of parents and children. Hundreds of dependents were temporarily deprived of moms or dads for almost a month. Thanks be to all who cared for those families—in faith. If each of us could muster that same faith in the service we are called to render, we would also be blessed.

\intalkheader{Virtue}

Virtue radiated from the choir. Each member seemed to exemplify President Brigham Young’s counsel: “Learn the will of God, keep his commandments and do his will, and you will be a virtuous person” (Discourses of Brigham Young, sel. John A. Widtsoe [Salt Lake City: Deseret Book Co., 1941], p. 194).

They applied the word of God (see Alma 31:5) not only in song but in sermons of example. After one concert, I was greeted by an individual who expressed gratitude in an unusual way. He said, “I am thankful for the choir’s message.”

I asked, “What message did you receive?”

His answer: “The choir was trying to teach me a better way to live.”

That comment seemed profound to me. What inspired him to feel such a force for good? Was it complicated classical music or the excellence of its rendition? I doubt it. I believe it was more likely communication by the Spirit, which allowed both the giver and receiver to be edified (see D\&C 50:21–22). And very possibly, the spirit of the listener was moved most by melodies soft and simple, sung with sweet sincerity. When the choir sang “The Lord’s Prayer,” for example, audiences were hushed and attentive. They seemed to understand, without fully knowing either the language or the history of the song.

The virtue of choir members brought blessings to their own lives. When the tour was completed and each was safely home, I thought of this promise from the Lord: “If thou art faithful and walk in the paths of virtue before me, I will preserve thy life” (D\&C 25:2). That same promise applies to me and to you.

\intalkheader{Knowledge}

Knowledge is essential to competent service for any missionary. So it was with members of the choir in this tour to eight nations. To communicate more effectively, the choir sang in ten languages! In these times of changing political views, much study was also required to determine which songs should be sung and which should not be sung.

But their quest for knowledge did not begin or end with music. They eagerly studied the culture, history, and ways of the people they had come to serve. On a moment’s notice, their knowledge of the gospel had to be retrieved when questioned by the media or interested individuals.

In these lands that have known so much of hardship and strife, choir members echoed the scriptural hope “that perhaps they might bring [others] to the knowledge of the Lord their God, … that they might also be brought to rejoice in the Lord their God, that they might become friendly to one another, and that there should be no more contentions in all the land” (Mosiah 28:2).

Members of the choir knew “that after a people have been once enlightened by the Spirit of God, and have had great knowledge of things pertaining to righteousness” (Alma 24:30), lives would never again be the same.

Just think of the good you can do if you accept a difficult challenge and pursue knowledge—then use it to bless others as did the choir!

\intalkheader{Temperance}

“To knowledge,” the Apostle Peter wrote, add “temperance” (2 Pet. 1:6). Temperance suggests sobriety and self-restraint in action. It reminds one of covenants made.

Members of the choir were always temperate and well disciplined, not from without but from within. Completely obedient to the Word of Wisdom, they were blessed with health and strength. Keeping a schedule that left little time for leisure, their grueling pace allowed them to perform for audiences that otherwise might have been excluded.

Repeatedly, scriptures teach that we be “temperate in all things” (1 Cor. 9:25; Alma 7:23; 38:10; D\&C 12:8). Temperance can protect each of us from the aftermath of excess.

\intalkheader{Patience}

Patience is one of the most practiced attributes of choir members. Checking into a hotel with a group of five hundred travelers and more than a thousand pieces of luggage provided practice in patience nearly every day. One dear sister never did receive her baggage. Her patience flowered into ingenuity as she attempted to feel fresh with the same clothing day after day.

Patience is a divine attribute. The Book of Mormon invites us to “come to a knowledge of the goodness of God, and his matchless power, and his wisdom, and his patience, and his long-suffering towards the children of men” (Mosiah 4:6).

Although choir members are not perfect, each one seemed to epitomize scriptural counsel to “continue in patience until ye are perfected” (D\&C 67:13). If they can do it, each of us can also develop that precious talent of patience.

\intalkheader{Brotherly Kindness}

Brotherly kindness was a hallmark of this tour. Never did I hear a disparaging remark. Especially did I watch those heroic members of the choir with serious physical disabilities. Some have severe visual handicaps and can read music only with their fingers. Others walk solely with the aid of crutches or appliances. Their great courage was superbly matched by the courtesy of companions who gave much that all might triumph together. Indeed, they typified this scriptural analogy:

“The whole body fitly joined together … according to the effectual working … of every part, maketh increase of the body unto the edifying of itself in love” (Eph. 4:16).

Brotherly kindness overcomes the rudeness of selfish intent. Each of us can develop brotherly kindness at home, at school, at work, or at play.

\intalkheader{Charity}

The Book of Mormon defines charity as the pure love of Christ (see Moro. 7:47). It further teaches that the “Lord God hath given a commandment that all men should have charity, which charity is love” (2 Ne. 26:30; see also Ether 12:34).

I saw choir members extend that love to countless souls. For those in need, members quietly contributed money, food, and goods. They shared freely of their precious time and talent without any thought of personal acclaim or recompense. Truly, “charity never faileth” any of us (1 Cor. 13:8; Moro. 7:46).

\intalkheader{Humility}

The choir’s humility seemed to increase throughout their journey. Though their mounting successes gave them much to be proud of, members grew into this scriptural pattern:

“They did fast and pray oft, and did wax stronger and stronger in their humility, and firmer and firmer in the faith of Christ, unto the filling their souls with joy and consolation, yea, even to the purifying and the sanctification of their hearts” (Hel. 3:35).

Choir members became “subject one to another” (1 Pet. 5:5) and imparted “the word of God, one with another” (Alma 1:20) as true disciples do. They were humble and teachable, as submissive Saints should be (see Mosiah 3:19; D\&C 105:12). Their examples of humility should ennoble our souls.

\intalkheader{Diligence}

Diligence was demonstrated by all members of the choir, but extraordinarily so by one precious mother whom I shall never forget. Just five days after the choir had left America, I was asked to inform her that her beautiful 37-year-old daughter had died after a long illness. The leaders of this sorrowing mother offered means for her to return home from Europe. She declined the offer. She and her family had already anticipated this possibility. Their decision had been made. It was not to be interpreted as a pattern for anyone else to follow, but for her alone. Her children and grandchildren had pleaded that she remain on assignment. So she continued in all diligence, never missing a single concert! She nobly fulfilled this scriptural counsel: “Therefore be diligent … in whatsoever difficult circumstances [you] may be” (D\&C 6:18).

Her example can bless each of you, just as it did her family and me.

\intalkheader{Godliness}

Godliness is an attribute that seems as difficult to define as it is to attain. Scriptures refer to “the mystery of godliness” (1 Tim. 3:16; D\&C 19:10). Because it is so special, I have chosen to speak of it last.

Simon Peter counseled us “to be in all holy conversation and godliness, looking for and hasting unto the coming of the day of God” (2 Pet. 3:11–12).

“The power of godliness is manifest” in the ordinances of the priesthood (D\&C 84:20). Godliness is not a product of perfection; it comes of concentration and consecration.

Godliness characterizes each of you who truly loves the Lord. You are constantly mindful of the Savior’s atonement and rejoice in His unconditional love. Meanwhile you vanquish personal pride and vain ambition. You consider your accomplishments important only if they help establish His kingdom on earth.

The Mormon Tabernacle Choir’s songs helped to convey their uncommon spirit of godliness. From their hearts the choir sang one number which bore testimony of love so amazing, so divine. Tears moistened the faces of more than a few as they expressed personal feelings of conversion and commitment to godliness. This song which the choir will soon sing includes these verses penned by Isaac Watts:

\begin{verse}
When I survey the wondrous cross\\
On which the Prince of glory died,\\
My richest gain I count but loss,\\
And pour contempt on all my pride.
\end{verse}

[“Love So Amazing, So Divine”]

Such are lessons taught by those who “serve unto the example and shadow of heavenly things” (Heb. 8:5). Unitedly members of the Mormon Tabernacle Choir testify of the living Lord and of His Church restored in these latter days, as do I. May God bless us to lift our lives by their example I pray in the name of Jesus Christ, amen.

\newpage
\settalk{Doors of Death}
\setspeaker{Russell M. Nelson}
\settalkdate{April 1992}
\talkheading{Doors of Death}{Russell M. Nelson}

This morning Elder Boyd K. Packer spoke of the “fountains of life.” This afternoon, I should like to speak of the “doors of death.” Both are crucial components of life.

Recently at the funeral of a friend, I visited with two distinguished brothers—former surgical colleagues of mine—whose lovely companions had both passed away. They said they were going through the most difficult period of their lives, adjusting to the almost unbearable loss of their partners. These wonderful men then told of their cooking breakfast for each other once a week, sharing that rotation with their sister, trying to lessen their loneliness imposed by the doors of death.

Death separates “the spirit and the body [which] are the soul of man” (D\&C 88:15). That separation evokes pangs of sorrow and shock among those left behind. The hurt is real. Only its intensity varies. Some doors are heavier than others. The sense of tragedy may be related to age. Generally the younger the victim, the greater the grief. Yet even when the elderly or infirm have been afforded merciful relief, their loved ones are rarely ready to let go. The only length of life that seems to satisfy the longings of the human heart is life everlasting.

\intalkheader{Mourning}

Irrespective of age, we mourn for those loved and lost. Mourning is one of the deepest expressions of pure love. It is a natural response in complete accord with divine commandment: “Thou shalt live together in love, insomuch that thou shalt weep for the loss of them that die” (D\&C 42:45).

Moreover, we can’t fully appreciate joyful reunions later without tearful separations now. The only way to take sorrow out of death is to take love out of life.

\intalkheader{Eternal Perspective}

Eternal perspective provides peace “which passeth all understanding” (Philip. 4:7). In speaking at a funeral of a loved one, the Prophet Joseph Smith offered this admonition: “When we lose a near and dear friend, upon whom we have set our hearts, it should be a caution unto us. … Our affections should be placed upon God and His work, more intensely than upon our fellow beings” (Teachings of the Prophet Joseph Smith, sel. Joseph Fielding Smith [Salt Lake City: Deseret Book Co., 1938], p. 216).

Life does not begin with birth, nor does it end with death. Prior to our birth, we dwelled as spirit children with our Father in Heaven. There we eagerly anticipated the possibility of coming to earth and obtaining a physical body. Knowingly we wanted the risks of mortality, which would allow the exercise of agency and accountability. “This life [was to become] a probationary state; a time to prepare to meet God” (Alma 12:24). But we regarded the returning home as the best part of that long-awaited trip, just as we do now. Before embarking on any journey, we like to have some assurance of a round-trip ticket. Returning from earth to life in our heavenly home requires passage through—and not around—the doors of death. We were born to die, and we die to live (see 2 Cor. 6:9). As seedlings of God, we barely blossom on earth; we fully flower in heaven.

\intalkheader{Physical Death}

The writer of Ecclesiastes said, “To every thing there is a season, and a time to every purpose under the heaven:

“A time to be born, and a time to die” (Eccl. 3:1–2; see also Alma 12:27).

Think of the alternative. If all 69 billion people who have ever lived on earth were still here, imagine the traffic jam! And we could own virtually nothing and scarcely make any responsible decisions.

\intalkheader{Plan of Happiness}

Scriptures teach that death is essential to happiness: “Now behold, it was not expedient that man should be reclaimed from this temporal death, for that would destroy the great plan of happiness” (Alma 42:8; italics added; see also 2 Ne. 9:6).

Our limited perspective would be enlarged if we could witness the reunion on the other side of the veil, when doors of death open to those returning home. Such was the vision of the psalmist who wrote, “Precious in the sight of the Lord is the death of his saints” (Ps. 116:15).

\intalkheader{Spiritual Death}

But there is another type of separation known in scripture as spiritual death (see 2 Ne. 9:12; Alma 12:16; 42:9; Hel. 14:16, 18). It “is defined as a state of spiritual alienation from God” (Joseph Fielding Smith, Doctrines of Salvation, comp. Bruce R. McConkie, 3 vols. [Salt Lake City: Bookcraft, 1954–56], 2:217). Thus, one can be very much alive physically but dead spiritually.

Spiritual death is more likely when goals are unbalanced toward things physical. Paul explained this concept to the Romans: “If ye live after the flesh, ye shall die: but if ye through the Spirit do mortify the deeds of the body, ye shall live” (Rom. 8:13).

If physical death should strike before moral wrongs have been made right, opportunity for repentance will have been forfeited. Thus, “the [real] sting of death is sin” (1 Cor. 15:56).

Even the Savior cannot save us in our sins. He will redeem us from our sins, but only upon condition of our repentance. We are responsible for our own spiritual survival or death (see Rom. 8:13–14; Hel. 14:18; D\&C 29:41–45).

\intalkheader{Coping with Trials}

Physical and spiritual trials provide continuing challenges in life. Each of you could provide illustrations from personal experience. Many of you, for instance, are at the twilight of life and endure long and difficult days. You know well the meaning of that divine injunction to “endure to the end” (see Matt. 24:13; Mark 13:13; 1 Ne. 13:37; 22:31; 2 Ne. 31:16; 33:4; Omni 1:26; 3 Ne. 15:9; D\&C 14:7; 18:22; 24:8).

The Savior of the world repeatedly asked that we pattern our lives after His (see John 13:15; 14:6; 1 Pet. 2:21; 2 Ne. 31:9, 16; 3 Ne. 18:16; 27:27). So we must endure trials—as did He. “Though he were a Son, yet learned he obedience by the things which he suffered” (Heb. 5:8).

When hardship heaps its heavy load upon us, good may yet be gleaned. Shakespeare so wrote:

\begin{verse}
Sweet are the uses of adversity,\\
Which, like the toad, ugly and venomous,\\
Wears yet a precious jewel in his head.
\end{verse}

[As You Like It, act 2, scene 1, lines 12–14]

The Lord’s expression is even more explicit: “After much tribulation come the blessings” (D\&C 58:4).

\intalkheader{Postmortal Life}

Mortality, temporary as it is, is terminated by the doors of death. Questions then come to searching minds of those left behind: “Where is my loved one now?” “What happens after death?” While many questions cannot be fully answered with available knowledge, much is known.

\intalkheader{Paradise}

The first station in postmortal life is named paradise. Alma wrote: “Concerning the state of the soul between death and the resurrection—Behold, it has been made known unto me … that the spirits of all men, as soon as they are departed from this mortal body, … are taken home to that God who gave them life.

“… The spirits of those who are righteous are received into a state of happiness, which is called paradise, a state of rest, a state of peace” (Alma 40:11–12).

\intalkheader{Resurrection and Immortality}

Some facetiously state that nothing is as permanent as death. Not so! The grip of physical death is temporary. It began with the fall of Adam; it ended with the atonement of Jesus the Christ. The waiting period in paradise is temporary too. It ends with the Resurrection. From the Book of Mormon we learn that the “paradise of God must deliver up the spirits of the righteous, and the grave deliver up the body of the righteous; and the spirit and the body is restored to itself again, and all men become incorruptible, and immortal, and they are living souls” (2 Ne. 9:13).

A few years ago, our stake president and his wife had a wonderful son taken in his youthful prime because of an automobile accident. We are consoled by the knowledge that the very laws that could not allow his broken body to survive here are the same eternal laws which the Lord will employ at the time of the Resurrection, when that body “shall be restored to [its] proper and perfect frame” (Alma 40:23; see also Alma 11:42–45).

The Lord who created us in the first place surely has power to do it again. The same necessary elements now in our bodies will still be available—at His command. The same unique genetic code now embedded in each of our living cells will still be available to format new ones then. The miracle of the Resurrection, wondrous as it will be, is marvelously matched by the miracle of our creation in the first place.

\intalkheader{Judgment}

Our resurrection will not be an end but a new beginning. It will prepare us for judgment by the Lord, who said, “As I have been lifted up [upon the cross] by men even so should men be lifted up by the Father, to stand before me, to be judged of their works” (3 Ne. 27:14).

Even before we approach that threshold of the eternal court of justice, we know who will personally preside: “The keeper of the gate is the Holy One of Israel; and he employeth no servant there; and there is none other way save it be by the gate; for he cannot be deceived, for the Lord God is his name.

“And whoso knocketh, to him will he open” (2 Ne. 9:41–42).

\intalkheader{Family Ties}

Loving relationships continue beyond the doors of death and judgment. Family ties endure because of sealings in the temple. Their importance cannot be overstated.

I remember vividly an experience I had as a passenger in a small two-propeller airplane. One of its engines suddenly burst open and caught on fire. The propeller of the flaming engine was starkly stilled. As we plummeted in a steep spiral dive toward the earth, I expected to die. Some of the passengers screamed in hysterical panic. Miraculously, the precipitous dive extinguished the flames. Then, by starting up the other engine, the pilot was able to stabilize the plane and bring us down safely.

Throughout that ordeal, though I “knew” death was coming, my paramount feeling was that I was not afraid to die. I remember a sense of returning home to meet ancestors for whom I had done temple work. I remember my deep sense of gratitude that my sweetheart and I had been sealed eternally to each other and to our children, born and reared in the covenant. I realized that our marriage in the temple was my most important accomplishment. Honors bestowed upon me by men could not approach the inner peace provided by sealings performed in the house of the Lord.

That harrowing experience consumed but a few minutes, yet my entire life flashed before my mind. Having had such rapid recall when facing death, I do not doubt the scriptural promise of “perfect remembrance” when facing judgment (Alma 5:18; see also Alma 11:43).

\intalkheader{Eternal Life}

After judgment comes the possibility of eternal life—the kind of life that our Heavenly Father lives. His celestial realm has been compared with the glory of the sun (see 1 Cor. 15:41; D\&C 76:96). It is available to all who prepare for it, the requirements of which have been clearly revealed: “Ye must press forward with a steadfastness in Christ, having a perfect brightness of hope, and a love of God and of all men. Wherefore, if ye shall press forward, feasting upon the word of Christ, and endure to the end, behold, thus saith the Father: Ye shall have eternal life” (2 Ne. 31:20; see also John 17:3).

\intalkheader{Time to Prepare}

Meanwhile, we who tarry here have a few precious moments remaining “to prepare to meet God” (Alma 34:32). Unfinished business is our worst business. Perpetual procrastination must yield to perceptive preparation. Today we have a little more time to bless others—time to be kinder, more compassionate, quicker to thank and slower to scold, more generous in sharing, more gracious in caring.

Then when our turn comes to pass through the doors of death, we can say as did Paul: “The time of my departure is at hand. I have fought a good fight, I have finished my course, I have kept the faith” (2 Tim. 4:6–7).

We need not look upon death as an enemy. With our full understanding and preparation, faith supplants fear. Hope displaces despair. The Lord said, “Fear not even unto death; for in this world your joy is not full, but in me your joy is full” (D\&C 101:36). He bestowed this gift: “Peace I leave with you, my peace I give unto you: not as the world giveth, give I unto you. Let not your heart be troubled, neither let it be afraid” (John 14:27).

As a special witness of Jesus Christ, I testify that He lives! I also testify that the veil of death is very thin. I know by experiences too sacred to relate that those who have gone before are not strangers to leaders of this Church. To us and to you, our loved ones may be just as close as the next room—separated only by the doors of death.

With that assurance, brothers and sisters, love life! Cherish each moment as a blessing from God (see Mosiah 2:21). Live it well—even to your loftiest potential. Then the anticipation of death shall not hold you hostage. With the help of the Lord, your deeds and desires will qualify you to receive everlasting joy, glory, immortality, and eternal lives. For this I pray in the name of Jesus Christ, amen.

\newpage
\settalk{Where Is Wisdom?}
\setspeaker{Russell M. Nelson}
\settalkdate{October 1992}
\talkheading{Where Is Wisdom?}{Russell M. Nelson}

Today I would like to pose a question asked long ago by Job: “Where shall wisdom be found?” (Job 28:12).

Leaders of this Church have repeatedly emphasized the importance of education. It is a vital component of wisdom. Not long after the pioneers began construction of their temple in Illinois, they established the University of the City of Nauvoo. The First Presidency proclaimed that this university “will enable us to teach our children wisdom, to instruct them in all the knowledge and learning, in the arts, sciences, and learned professions.”

A similar scene followed after the persecuted pioneers entered the valley of the Great Salt Lake. Less than three years later, on 28 February 1850, they instituted the University of the State of Deseret. Later several academies of learning were established.

Now as Church membership worldwide exceeds eight million, it is evident that a direct role of the Church in secular education is no longer feasible. Yet our commitment to education remains constant.

Scriptures teach that “the glory of God is intelligence” (D\&C 93:36). They also teach that individual “intelligences … were organized before the world was” (Abr. 3:22). “Man was also in the beginning with God. Intelligence, or the light of truth, was not created or made, neither indeed can be” (D\&C 93:29).

Our personal intelligence is everlasting and divine. I believe Thomas Jefferson felt that dignity of the human spirit when he wrote, “I have sworn upon the Altar of God, eternal hostility against every form of tyranny over the mind of man.”

\intalkheader{Seek Education}

Because of our sacred regard for each human intellect, we consider the obtaining of an education to be a religious responsibility. Yet opportunities and abilities differ. I believe that in the pursuit of education, individual desire is more influential than institution, and personal faith more forceful than faculty.

Our Creator expects His children everywhere to educate themselves. He issued a commandment: “Seek ye diligently and teach one another words of wisdom; yea, seek ye out of the best books words of wisdom; seek learning, even by study and also by faith” (D\&C 88:118). And He assures us that knowledge acquired here will be ours forever (see D\&C 130:18–19).

Measured by this celestial standard, it is apparent that those who impulsively “drop out” and cut short their education not only disregard divine decree but frustrate the realization of their own potential.

I remember my moment of resolution many years ago when, as an untrained teenager, I secured temporary employment at Christmastime. The work was monotonous. Each hour and each day passed slowly. I resolved then and there that I must obtain an education that would qualify me better for life. I determined to stay in school and work for an education as though my very life depended upon it.

Later as stake president I was questioned by many young people about their own educational pursuits. Some asked me how long it took to become a doctor of medicine. “The general pattern would be four years at a university, followed by four years in medical school,” I replied. “And if you choose to become a specialist, that could take another five years or more, depending upon your desire.”

That occasionally evoked a reaction: “That adds up to thirteen years—and maybe more? That’s too long for me!”

“It all depends,” I would respond. “Preparation for your career is not too long if you know what you want to do with your life. How old will you be thirteen years from now if you don’t pursue your education? Just as old, whether or not you become what you want to be!”

So my counsel then—and now—is to continue your education wherever you are, whatever your interest and opportunity, however you determine you can best serve your family and society.

\intalkheader{Beware of Unbalance}

Choose what you will learn and whose purposes you will serve. But don’t place all your intellectual eggs in one basket of secular learning. Remember this warning from the Book of Mormon:

“O the vainness, and the frailties, and the foolishness of men! When they are learned they think they are wise, and they hearken not unto the counsel of God, for they set it aside, supposing they know of themselves, wherefore, their wisdom is foolishness and it profiteth them not. And they shall perish.

“But to be learned is good if they hearken unto the counsels of God” (2 Ne. 9:28–29).

That scripture reminds me of a friend who proudly boasted that his climb toward wealth had come from tireless work and lessons learned in the “school of hard knocks.” But his fortune had come at the expense of his spiritual development. When it was too late, he regretfully discovered that his ladder of success had been leaning against the wrong wall. He had never read this instruction from his Maker:

“Seek not for riches but for wisdom, and behold, the mysteries of God shall be unfolded unto you, and then shall you be made rich. Behold, he that hath eternal life is rich” (D\&C 6:7; see also D\&C 11:7).

Mankind’s unfamiliarity with the scriptures has sometimes brought sorrow to great numbers of people over long periods of time. The suffering that has resulted from such ignorance is truly tragic. May I illustrate with excerpts from history that pertain to the spread of infection.

In the nineteenth century, health officials and others were concerned about pollution of the air, not by visible smoggy hydrocarbons of today, but by an invisible miasma that was blamed for almost any infection. In 1867, for example, Lord Lister indicted bad air as the chief cause of infection. Because of that, in 1869 Simpson from Edinburgh urged that hospitals be taken down and rebuilt every few years. Such an extravagant practice was also advocated by other experts.

Even Florence Nightingale, a living legend following her heroic efforts in the Crimean War, failed to recognize the transmission of infection from one patient to another—this despite her careful notations that wound infection accounted for 40 percent of postoperative mortality.

But others missed the connection too. For centuries, lives of innumerable mothers and children were claimed by “childbirth fever”—infections unknowingly transmitted among the innocent by unwashed hands of attendants.

It was only a short century ago that the great work of Koch, Pasteur, and others proved that infection could be caused by bacteria in contaminated body fluids—or infected issues—passed from one individual to another.

With these highlights of history in mind, may I quote the word of the Lord recorded long ago in Leviticus, chapter 15:

“The Lord spake unto Moses and to Aaron, saying,

“Speak unto the children of Israel, and say unto them, When any man hath a running issue out of his flesh, because of his issue he is unclean.

“And this shall be his uncleanness in his issue. …

“Every bed, whereon he lieth that hath the issue, is unclean: and every thing, whereon he sitteth, shall be unclean.

“And whosoever toucheth his bed shall wash his clothes, and bathe himself in water. …

“And he that toucheth the flesh of him that hath the issue shall wash his clothes, and bathe himself in water” (Lev. 15:1–5, 7; emphasis added).

Several verses follow which re-emphasize and illustrate those important principles. Then we read this conclusion:

“And when he that hath an issue is cleansed of his issue; then he shall … wash his clothes, and bathe his flesh in running water, and shall be clean” (Lev. 15:13).

Thus, our loving Heavenly Father had clearly revealed principles of clean technique in the handling of infected patients more than three thousand years ago! These scriptures are in complete harmony with modern medical guidelines. But during those many millennia, how many mothers needlessly perished? How many children suffered because man’s quest for knowledge had failed to incorporate the word of the Lord?

\intalkheader{Contemporary Challenges}

In our day, many challenges face us. Some are new, some are old—simply clothed in modern attire. The epistles of Paul include prophecies pertaining to our day. Do these descriptions sound familiar?

“In the last days perilous times shall come.

“For men shall be lovers of their own selves, covetous, boasters, proud, [and the list of insidious qualities goes on] …

“Without natural affection, …

“Lovers of pleasures more than lovers of God;

“Having a form of godliness, but denying the power thereof: …

“Ever learning, and never able to come to the knowledge of the truth” (2 Tim. 3:1–5, 7).

Paul’s warnings describe apostasy and other dangers of our day. Some of these perils are contrary to God’s purposes and are championed by persuasive people possessing more ability than morality, more knowledge than wisdom. Their rationalization breeds justification. The Bible affirms that “the way of a fool is right in his own eyes” (Prov. 12:15). Indeed, individuals with malignity of purpose often wear the mask of honesty. So we must constantly be on guard.

To build a house straight and strong, you do not choose crooked boards. So to build your eternal destiny, you cannot—you must not—limit lessons only to those warped to exclude revelation from God. The Book of Mormon offers this note of caution and hope:

“Seek not to counsel the Lord, but to take counsel from his hand. For behold, ye yourselves know that he counseleth in wisdom, and in justice, and in great mercy, over all his works” (Jacob 4:10).

Remember the terrible price paid for ignorance of divine instruction. Until the turn of this century, infection was spread as if no one had ever read or taken seriously the fifteenth chapter of Leviticus. Where is wisdom?

Today we are seriously concerned with the increasing incidence of human infection with HIV (human immunosuppressive virus) and variant viruses and the associated outbreak of AIDS (acquired immune deficiency syndrome). An epidemic has been forecast—a plague fueled by a vocal few who exhibit greater concern for civil rights than for public health, a plague abetted by the immoral. Some live in lust as though God’s commandment to be chaste was written with an asterisk, exempting them from obeying. And regrettably, as in previous plagues, many innocent victims are doomed to suffer. Where is wisdom?

Avoidable deaths and mounting financial burdens are also incurred worldwide because of indifference to or ignorance of God’s declaration that tobacco “is not good for man” (D\&C 89:8). Many other societal problems could be listed, such as alcohol and drug abuse, gambling, civil strife, and erosion of family stability. We may know so much, yet learn so little. “Where is the wisdom we have lost in knowledge? Where is the knowledge we have lost in information?” Again I ask, Where is wisdom?

\intalkheader{Wisdom Found}

Wisdom is to be found in pure intelligence—in that divine light which can guide people in all countries, all climes, and all continents. The Lord has promised that “a light shall break forth among them that sit in darkness, and it shall be the fulness of my gospel” (D\&C 45:28). Then He lamented:

“But they receive it not; for they perceive not the light, and they turn their hearts from me because of the precepts of men. …

“And there shall be men standing in that generation, that shall not pass until they shall see an overflowing scourge; for a desolating sickness shall cover the land.

“But my disciples shall stand in holy places, and shall not be moved; but among the wicked, men shall lift up their voices and curse God and die.

“And there shall be earthquakes also in divers places, and many desolations; yet men will harden their hearts against me” (D\&C 45:29, 31–33; see also D\&C 87:6).

In bright contrast to such bitter chaos, the light of the gospel of Jesus Christ beams as the hope of the world. Missionaries and members courageously proclaim its brilliance. Wise students throughout the world heed its light and enrich their education by adding the curriculum of Church seminaries and institutes. The Lord hides His wisdom from no one: “If any of you lack wisdom, let him ask of God” (James 1:5).

Again I pose the question once asked by Job: “Where shall wisdom be found?” (Job 28:12). Answer: It emanates from the Lord. He Himself said:

“I will give unto the children of men line upon line, precept upon precept, here a little and there a little; and blessed are those who hearken unto my precepts, and lend an ear unto my counsel, for they shall learn wisdom; for unto him that receiveth I will give more” (2 Ne. 28:30).

Divine light and wisdom continue to increase when love for Deity grows:

“That which is of God is light; and he that receiveth light, and continueth in God, receiveth more light; and that light groweth brighter and brighter until the perfect day” (D\&C 50:24; see also D\&C 88:67).

“He that keepeth his commandments receiveth truth and light, until he is glorified in truth and knoweth all things” (D\&C 93:28).

Where is wisdom? It pulses and surges with the Lord’s light of truth! With that light He lifts us toward eternal life, I testify in the name of Jesus Christ, amen.

\newpage
\settalk{Honoring the Priesthood}
\setspeaker{Russell M. Nelson}
\settalkdate{April 1993}
\talkheading{Honoring the Priesthood}{Russell M. Nelson}

Brethren, relatively little is written on my subject. Yet we are all expected to know about it. I speak of honoring the priesthood.

This is The Church of Jesus Christ of Latter-day Saints. He who stands at the head of His restored Church so ordered His priesthood “that every man might speak in the name of God the Lord, even the Savior of the world” (D\&C 1:20).

Remarkable! He chose to honor us with His priesthood. So we honor Him by honoring His priesthood—both its power and those who bear it. By so doing, men, women, and children throughout the world will be blessed. Honoring the priesthood fosters respect, respect promotes reverence, and reverence invites revelation.

President Ezra Taft Benson has specifically asked us to follow proper priesthood protocol—principles, he noted, “that many of us have learned by observation while listening to senior brethren.” He said, “Protocol is a long-established practice prescribing complete deference to … an order of correct procedure.” I will quote from President Benson and other leaders because, as you will note, much of my message pertains to that protocol.

\intalkheader{Types of Organizations}

Differences exist in practice and organization between the Lord’s Church and man-made institutions. Men and women may form associations for and among themselves and be governed by stipulations that are mutually acceptable.

The Church of Jesus Christ of Latter-day Saints, however, is neither a democracy nor a republic. His is a kingdom—the kingdom of God on earth. His is a hierarchical church, with ultimate authority at the top. The Lord directs His anointed servants. They testify to all the world that God has again spoken. The heavens have been opened. A living linkage has been formed between heaven and earth in our day.

That supreme authority is supported by a firm foundation following an organizational pattern established anciently. Jesus Christ is the chief cornerstone, with Apostles and prophets and all the gifts, powers, and blessings that characterized the Church in earlier days (see 1 Cor. 12:28).

\intalkheader{Leaders and Titles}

Secular and spiritual institutions have differing patterns of leadership. Man-made organizations are governed by officers with titles that designate rank or accomplishment. A military officer, judge, senator, doctor, or professor is properly addressed by title. We appropriately honor individuals who have attained such positions.

In contrast, the kingdom of God is governed by the authority of the priesthood. It is not conferred for honor, but for a ministry of service. Priesthood titles are not created by man; neither are they for adornment, nor do they express mastership. They denote appointment to service in the work of the Lord. We are called, sustained, and ordained—not by ourselves, but “by prophecy, and by the laying on of hands by those who are in authority, to preach the Gospel and administer in the ordinances thereof” (A of F 1:5; see also Heb. 5:4).

Titles pertaining to the holy priesthood deserve our utmost care and respect. Each member of the First Presidency is addressed and spoken of as “President” (see D\&C 107:22). The title President is also used when referring to the presidency of a stake or mission, and in reference to a quorum or branch president. The title Apostle is sacred. It has been given of God and belongs only to those who have been called and ordained as “special witnesses of the name of Christ in all the world” (D\&C 107:23). An Apostle speaks in the name of Him whose special witness he is. This hallowed title is not used in ordinary forms of address. The preferred title for one of the Twelve is Elder or Brother.

The title bishop is also expressive of presidency; the bishop is the president of the Aaronic Priesthood in his ward and the presiding high priest of the ward organization. Reverently we refer to him as “the bishop.”

Elder is a sacred title shared by all who bear the Melchizedek Priesthood.

\intalkheader{General Counsel}

May I offer counsel of a general nature, first with comments about General Authorities. We recognize them as instruments in the hand of the Lord, yet realize that they are ordinary human beings. They require haircuts, laundry services, and occasional reminders just like anyone else. President Benson once shared with us a story to illustrate. He said:

“Orson F. Whitney … was a great man to concentrate. One day when he was traveling by train, he was so preoccupied that he did not notice the train pass the station where he was to get off. So he had to [be driven] back to where he should have been. Meanwhile the stake president waited and waited. … Finally when he decided that something had more than likely happened to Brother Whitney and he was not going to make it, they commenced the meeting. As Elder Whitney approached, he was greeted by the opening hymn, which was ‘Ye Simple Souls Who Stray.’”

We honor such a man because of his extraordinary calling. His official acts are valid on earth and in heaven. Well do I remember the first time I met one of the General Authorities. It was a feeling beyond description. Though I was but a boy, immediately—almost instinctively—I rose to my feet. Even now I feel that same way when one of the Brethren enters the room. A General Authority is an oracle of God.

Often we speak of keys of priesthood authority. Fifteen living men—the First Presidency and the Twelve—have been ordained as Apostles and have had all keys of priesthood authority conferred upon them. President Gordon B. Hinckley recently explained:

“Only the President of the Church has the right to exercise [those keys] in their fulness. He may delegate the exercise of various of them to one or more of his Brethren. …

“Such agency has been given by President Benson to his Counselors and to the Twelve according to various responsibilities delegated to them.”

Under assignment from the First Presidency and the Twelve, General Authorities confer the appropriate keys upon presidents of stakes and of missions, who in turn confer the needed keys upon bishops and upon quorum and branch presidents.

Assigned to each one who bears the priesthood is a loving leader because “mine house is a house of order, saith the Lord God, and not a house of confusion” (D\&C 132:8).

That order also defines bounds of revelation. The Prophet Joseph Smith taught that “it is contrary to the economy of God for any member of the Church, or any one, to receive instruction for those in authority, higher than themselves.” That same principle precludes receiving revelation for anyone outside one’s defined circle of responsibility.

Honoring the priesthood also means to honor your personal call to serve. A few do’s and don’ts may be helpful:

The one who extends and the one who receives a call are both under obligation of accountability. I quote from Elder James E. Talmage:

“Those through whom the call came to him … are as surely held answerable for their acts as is he for his; and of every one shall be demanded a strict and personal accounting for his stewardship, a report in full of service or of neglect, of use or abuse in the administration of the trust to him committed.”

Some aspects of the priesthood are not related to position or title. Authority to administer a priesthood blessing, for example, is dependent only upon ordination and worthiness. The Lord would not withhold blessings from any of His children for want of one with a particular calling. Every elder in the Church holds the same priesthood as the President of the Church.

Brethren, please remember: the highest degree of glory is available to you only through that order of the priesthood linked to the new and everlasting covenant of marriage (see D\&C 131:1–4). Therefore, your first priority in honoring the priesthood is to honor your eternal companion.

\intalkheader{Specific Counsel}

Now for counsel more specific. Husbands and fathers: with your dear partner, shape attitudes at home. Establish a pattern of prayer. Pray regularly and vocally for your priesthood and auxiliary leaders, both local and general. Your manners of courtesy at home and of reverence in the chapel will be copied by members of your family. Help your loved ones follow proper channels when they seek guidance. Teach that counsel should be obtained from trusted parents and leaders on a local level, not from General Authorities. In the past two decades, the First Presidency has sent out essentially the same letter six times to reaffirm that policy.

Fathers, you understand the principle of temporal self-reliance and try to provide for a year’s supply, stored at home. Please also consider the need for spiritual food and self-reliance—not just for a year, but for a lifetime—also stored at home. A worthy father should have first opportunity to administer blessings to members of his family. As time moves on, his sons may then draw from that spiritual reservoir, worthy to administer to their own families and to their parents.

Now to young men who bear the Aaronic (or preparatory) Priesthood: if you honor it and prepare for and are worthy of a call to be a missionary, I promise that you will then “speak in the name of God the Lord” and bring His light to searching souls. To them you will be as a ministering angel, remembered with love forever. (See D\&C 13).

Though I next speak to our beloved presidents and bishops, the principles apply to all. When one who presides over you comes into a meeting where you have been presiding, please consult with him immediately for instruction. Determine his desires. Be certain to allow adequate time for a message from him. A poignant illustration was once related by Elder James E. Faust:

“I learned some time ago of the distress felt by members of a stake in this valley when their stake presidency was reorganized. The presiding officer was one of the most venerated and unique Apostles in all the history of the Church. [Elder] LeGrand Richards was then in his nineties but was sharp and alert. During the conference, the local people who were called on to speak took most of the time. As a result, Elder Richards had only ten or fifteen minutes remaining in the meeting. What did he do? Go overtime? No. He bore a brief testimony and closed the meeting on time.

“The members of the stake did not necessarily want to go overtime. … They were upset, however, because the local membership, who would have other opportunity to hear from their local leaders, would never again, and in fact never did again, have an opportunity to hear from this venerable Apostle. In short, the speakers did not respect the presiding officer.”

When a presiding General Authority has spoken, no one speaks following him. After the meeting has concluded, presidents and bishops, remain at the side of your file leader until excused. He may be impressed to give additional teaching or direction. And you may also prevent problems. For example, if a member asks a question of your leader that should not be directed to him, you are there to respond.

Now for comments about the stake high council. It has no president. It has no autonomy and meets, even when divided into committees, only upon call from the stake presidency. Although high councilors may be seated in the order of their call to the council, no one member has seniority over another.

In contrast, seniority is honored among ordained Apostles—even when entering or leaving a room. President Benson related to us this account:

“Some [years] ago Elder Haight extended a special courtesy to President Romney while they were in the upper room in the temple. President Romney was lingering behind for some reason, and [Elder Haight] did not want to precede him out the door. When President Romney signaled [for him] to go first, Elder Haight replied, ‘No, President, you go first.’

“President Romney replied with his humor, ‘What’s the matter, David? Are you afraid I’m going to steal something?’”

Such deference from a junior to a senior Apostle is recorded in the New Testament. When Simon Peter and John the Beloved ran to investigate the report that the body of their crucified Lord had been taken from the sepulcher, John, being younger and swifter, arrived first, yet he did not enter. He deferred to the senior Apostle, who entered the sepulcher first. (See John 20:2–6.) Seniority in the apostleship has long been a means by which the Lord selects His presiding high priest.

\intalkheader{Rebuke and Repentance}

Brethren, these matters are important. More than a century and a half ago, the Lord issued a sharp rebuke to His people. These are His words:

“Verily, condemnation resteth upon you, who are appointed to lead my Church, … and also upon the Church; and there must needs be a repentance and a reformation among you, in all things, in your examples before the Church and before the world, in all your manners, habits and customs, and salutations one toward another; rendering unto every man the respect due the office, calling, and priesthood whereunto I, the Lord, have appointed and ordained you.”

If any among us are also guilty of treating as trivial such things that are sacred, we may repent and resolve to honor the priesthood and those to whom the Lord has entrusted its keys.

Brethren, to all mankind we proclaim these everlasting truths: “The Melchizedek Priesthood holds the right of presidency, and has power and authority over all the offices in the church in all ages of the world” (D\&C 107:8). This power holds “the keys of all the spiritual blessings of the church” (D\&C 107:18). May we fully honor that priesthood, I pray in the name of Jesus Christ, amen.

\newpage
\settalk{Combatting Spiritual Drift—Our Global Pandemic}
\setspeaker{Russell M. Nelson}
\settalkdate{October 1993}
\talkheading{Combatting Spiritual Drift—Our Global Pandemic}{Russell M. Nelson}

Dear colleagues and friends, it is a privilege for me to represent The Church of Jesus Christ of Latter-day Saints at this historic parliament. I pay tribute to those who have organized it and to all who have participated.

Prior to my call to full-time service as a Church leader, I worked as a medical doctor specializing in the teaching, research, and practice of thoracic and cardiovascular surgery. I performed about seven thousand operations. During nearly four decades as a student of medical sciences, I learned to distinguish between the physical and spiritual components of the human soul. A patient could be physically weak but spiritually strong, and vice versa. For example, on many occasions I watched the critically ill mobilize great spiritual strength in order to undergo a high-risk operation to correct a physical problem within the heart. In contrast, I saw individuals without physical abnormality who were seriously ill because of spiritual loss of heart.

Duties have taken me to approximately one hundred nations of the earth, where I have observed the physical and spiritual dualism of life on a broader scale. There I have noted examples of the same physical and spiritual disparity that I previously encountered in my profession. I have observed that if a great physical disaster should strike, such as a devastating earthquake or flood, people are motivated spiritually by an uncommon desire to help one another. After the great earthquake in Armenia in December 1988, for example, I saw spontaneous outpourings of help from all over the world. Donors from many nations contributed aid to hundreds of thousands of homeless victims. Generosity seemed to transcend political or racial differences.

Fortunately, such major disasters are rare. But unfortunately, when normal life resumes, the pendulum seems to swing from spiritual vitality to laxity. It is ironic that as affluence and physical comforts increase, spiritual strength declines. This observation has prompted the title for my message: “Combatting Spiritual Drift—Our Global Pandemic.” Reversing this crisis in the health of the human spirit is an enormous challenge.

\intalkheader{Changes Since the Last Parliament of the World’s Religions}

In the century that has elapsed since the last meeting of the Parliament of the World’s Religions, there have been notable advances in virtually every field of human endeavor. Think of the progress made in transportation, communication, commerce, agriculture, medicine, science, and electronics.

But spiritual progress has lagged behind. We see evidences of increasing ethnic strife and hatred. Nationalism seems to be taking priority over brotherly love. Violence and civil wars are raging. Divorce and diminishing regard for the sanctity of human life have eroded the strength of the family—the basic unit of society. Immorality, infidelity, and promiscuity—once shunned—are now tolerated and even condoned. We have witnessed the insidious intrusion of pornography, with its attendant denigration of the human soul. And gambling, which preys upon the poor and the compulsive, has crept from the realm of the illegal into the arena of governmental sponsorship.

\intalkheader{Religion’s Role in the Past Century}

We may well ask, “What has been the role of religion during this past century?” The lofty objectives and high ideals promoted by the religions of the world are truly commendable. Much good has been done. We should all feel a deep sense of gratitude for noble predecessors in every faith group.

At the same time, we should be objective enough to analyze past activity and note any efforts that might have been inadequate, misdirected, or even counterproductive. For instance, worldly trends may have been assimilated into some religious ideology. On occasion, there may have been misguided incursions of religion into the political arena. And in spite of the exemplary behavior of the majority of our religious leaders and members, inappropriate conduct by any individual identified with a specific religious movement tends to tarnish the image of all. Sadly, we have also witnessed instances of interfaith contention, maligning, and intolerance.

Moreover, honest people searching for answers to life’s greatest questions may have turned to religion, asking, “Where did I come from? Why am I here? Where am I going?” Some may have come away unsatisfied.

Given these realities, religion must shoulder greater responsibility in healing the spiritual sickness that engulfs our globe.

\intalkheader{The Role of Religion in the Future}

The dismal dusk of today’s spiritual drift provides a rare opportunity for the radiance of religion to light the way to a new tomorrow. This can happen only as we proclaim eternal truths that have the power to engender spiritual strength. Human nature cannot be changed by reforming public policy; that kind of change comes by exposing the human mind and heart to the transforming teachings of the Lord Jesus Christ. I have learned that when we teach His correct principles, people govern themselves appropriately.

We at this world parliament represent many religious persuasions. Because there is much that is praiseworthy in each of our faiths, it is important for us to maintain the integrity of our religious institutions and to preserve tolerance of each other’s sacred beliefs. These factors are essential to the strength of a pluralistic society. Tolerance and understanding are enhanced as we teach clearly and courteously the tenets of our religions.

In that spirit, and as done so well by previous speakers, I would like to speak of the organization I represent. I would like to speak of the institution and of the doctrine of The Church of Jesus Christ of Latter-day Saints, sometimes known as the “Mormon” church. That is not its correct name, as many of you may know; it is only a nickname—derived from our belief in the Book of Mormon, to which I will refer later.

\intalkheader{The Church of Jesus Christ of Latter-day Saints}

The Church of Jesus Christ of Latter-day Saints has endured much persecution, oppression, and hardship. From these experiences has come our genuine passion for tolerance and a sincere appreciation for feelings held sacred by others. Early members of this church were some of the great pioneers in the development of the United States of America. All eight of my great-grandparents, for example, were early converts to this church from populous nations of Europe. All eight of them immigrated to the U.S.A. and became part of the nineteenth century’s historic trailblazing trek westward.

Now the Church has a worldwide membership in excess of eight million people. Although its world headquarters are located in Salt Lake City, Utah, its congregations flourish throughout the world in approximately one hundred fifty nations on six continents and upon many islands of the sea.

Ours is a lay ministry. For example, I previously served as general president of the Sunday School of the Church for eight years, while continuing my professional work as a cardiovascular surgeon. Our religious activities throughout the world are led by male and female volunteers. We believe that everyone can be happy because everyone can serve. Congregations in Germany are led by Germans, in Peru by Peruvians, in Japan by Japanese, and so on. Government by consent of the governed is a rule of the Church.

Members of the Church in biblical times were addressed as saints. Likewise, members of the Church today are known as Latter-day Saints, often abbreviated as “LDS.” We regard a saint as a believer in Christ, one who knows of His perfect love and who strives to live in accord with His commandments. The term connotes a high level of commitment to following the Savior.

A summary of the tenets of The Church of Jesus Christ of Latter-day Saints has long been known as the Articles of Faith. Copies are available for any who are interested. I will quote liberally from this statement in discussing the doctrines of our church.

First, “we believe in God, the Eternal Father, and in His Son, Jesus Christ, and in the Holy Ghost” (A of F 1:1). We worship our Heavenly Father and pray to Him in the name of His Beloved Son, Jesus Christ, through the power of the Holy Ghost. We revere the Son of God as our Savior, our Redeemer, and our Master. He came into the world to do the will of His Father (see 3 Ne. 27:13). He died for us and was resurrected from the tomb. His triumph over death brings the blessings of resurrection and immortality to all mankind (see 1 Cor. 15:20–22). Our adoration for Him is typified in this passage from the Book of Mormon:

“We talk of Christ, we rejoice in Christ, we preach of Christ, we prophesy of Christ, … that our children may know to what source they may look for a remission of their sins” (2 Ne. 25:26).

Nothing in human history equals the wonder, the magnitude, or the fruits of the matchless life of the Son of God. He is our exemplar and the author of our faith. And one day He will again come to the earth to begin His promised millennial reign.

We also believe in prophets. God has always taught His children upon the earth through prophets. In ancient times, Abraham, Isaac, Jacob, Isaiah, Ezekiel, Jeremiah, and many others were among those chosen by God as His prophets. Likewise, we believe that Joseph Smith was called by God as the prophet to inaugurate this present era, known as “the dispensation of the fulness of times” (see Eph. 1:10; D\&C 112:30; D\&C 121:31; D\&C 124:41; D\&C 128:18, 20; D\&C 138:48). He was called and chosen to bring together all former “dispensations, and keys, and powers … from the days of Adam even to the present time” (D\&C 128:18). Previous dispensations include those identified with Adam, with Noah, with Moses, and with Jesus and His Apostles in the meridian of time.

This modern dispensation of which I speak fulfills the biblical promise of a “restitution of all things” (Acts 3:21; see also Rev. 14:6, 7). It also fulfills another scriptural promise that “the Lord doth grant unto all nations, of their own nation and tongue, to teach his word” (Alma 29:8). We believe in the restitution of the same organization that once existed in the early Christian Church, with Apostles, prophets, and teachers (see A of F 1:6). Our origins did not spring from protest of any kind; hence, we are not a Protestant denomination. Rather, we understand that the church of Jesus Christ has been restored to the earth, established, as the Apostle Paul once declared, “upon the foundation of the apostles and prophets, Jesus Christ himself being the chief corner stone” (Eph. 2:20).

God the Father and His Son, Jesus Christ, literally appeared to the Prophet Joseph Smith in April 1820. While this may be relatively unknown in the annals of history, yet it constituted an epochal event in the history of humanity—a profound demonstration of the love of God for all of His children. Because this milestone is of such importance, I would like to quote excerpts from the Prophet’s written account:

“I saw two Personages, whose brightness and glory defy all description, standing above me in the air. One of them spake unto me, calling me by name and said, pointing to the other—This is My Beloved Son. Hear Him!” JS—H 1:17).

Joseph Smith encountered intense persecution after this experience had become public knowledge. He wrote:

“I had actually seen a light, and in the midst of that light I saw two Personages, and they did in reality speak to me; and though I was hated and persecuted for saying that I had seen a vision, yet it was true; and while [abusers] were persecuting me, reviling me, and speaking all manner of evil against me falsely for so saying, I was led to say in my heart: Why persecute me for telling the truth? I have actually seen a vision; and who am I that I can withstand God, or why does the world think to make me deny what I have actually seen? For I had seen a vision; I knew it, and I knew that God knew it, and I could not deny it, neither dared I do it; at least I knew that by so doing I would offend God, and come under condemnation” (JS—H 1:25).

We regard that theophany—which we designate as the First Vision—and subsequent communications of heavenly beings with the Prophet Joseph Smith to be as real as my conversation with you now. These and other visitations relating to the promised restitution of the Church are well documented.

For his singular service as the prophet called to inaugurate this dispensation, Joseph Smith gave his life as a martyr at the young age of thirty-eight years. As we honor prophets of former days, so we honor the Prophet Joseph Smith. May I quote from an account written shortly after his martyrdom in 1844, penned by one of his contemporaries:

“In the short space of twenty years, he [the Prophet Joseph Smith] has brought forth the Book of Mormon, which he translated by the gift and power of God, and has been the means of publishing it on two continents; has sent the fulness of the everlasting gospel, which it contained, to the four quarters of the earth; has brought forth the revelations and commandments which compose this book of Doctrine and Covenants, and many other wise documents and instructions for the benefit of the children of men; gathered many thousands of the Latter-day Saints, founded a great city, and left a fame and name that cannot be slain. He lived great, and he died great in the eyes of God and his people; and like most of the Lord’s anointed in ancient times, has sealed his mission and his works with his own blood” (D\&C 135:3).

We regard the Book of Mormon as companion scripture to the Holy Bible (see A of F 1:8). While the Bible bears record of teachings the Lord gave to His prophets who lived in the Old World, the Book of Mormon contains teachings the Lord gave to His prophets who lived in the American hemisphere. This record covers a period of about a thousand years, from approximately 600 b.c. to a.d. 400 (although there is reference to other migrations and peoples preceding that time). The Book of Mormon bears record of the personal ministry of Jesus Christ, as the resurrected Lord, to inhabitants of ancient America. Prophets recorded His teachings among them. Hence, the Book of Mormon is “Another Testament of Jesus Christ” (Book of Mormon Title Page). It came in fulfillment of biblical prophecy: “In the mouth of two or three witnesses shall every word be established” (2 Cor. 13:1). It came as a heavenly gift to the entire human family. Originally inscribed on metallic plates and written in a modified Egyptian language, that record was preserved for many centuries, then made available for translation into the English language. That translation was done by the Prophet Joseph Smith through the gift and power of God. Subsequently, the book has been translated into many other languages. Since its first publication in 1830, millions of copies of the Book of Mormon have been printed. In the year 1992 alone, nearly five million copies were distributed.

We have observed that sincere study of the Book of Mormon transforms the spirit and engenders a testimony of Jesus Christ that heals the spiritual hearts of men and women throughout the world. The more widely it is read and understood, the greater is its influence for good. We find that as the Book of Mormon spreads to the four corners of the globe, men and women exhibit greater devotion and become promulgators of peace within their own families and society.

Other divine revelations to prophets, ancient and modern, have been canonized into additional scriptures—books known as the Doctrine and Covenants and as the Pearl of Great Price. Thus, we believe the Holy Bible, the Book of Mormon, the Doctrine and Covenants, and the Pearl of Great Price to be four standard works of sacred scripture.

These books of revelation, recorded by prophets from age to age, provide broad eternal perspective and clarify the purpose of man’s existence. These scriptures teach that all of us lived as spirit children of God before we came into this life. We came to earth to obtain a body, to gain earthly experiences, and to prove our faith as we confront life’s many challenges. And because there are differences in degree of obedience to God’s eternal laws, there will be different gradations, or degrees of glory, in our eternal destinies. Jesus taught that “in my Father’s house are many mansions: if it were not so, I would have told you” (John 14:2; see also 1 Cor. 15:40–41; D\&C 76:70–98; D\&C 98:18).

This plan of salvation reveals the supernal potential that each of us has as a unique son or daughter of our Heavenly Father. It is literally possible for each of us to return to His presence, to become heir to all that He has in store for the faithful. This plan informs us where we came from, why we are here, and where we are going. We cherish this knowledge and are grateful for the joyful satisfaction it produces.

We believe that our priesthood leaders are ordained by authority restored by heavenly messengers in these latter days and that leaders may receive direct inspiration and revelation (see A of F 1:5, 7, 9). Both a lesser priesthood—such as was operative in the days of Moses—and a greater priesthood—such as was operative in the days of Jesus and the Apostles—have been restored, along with many additional specific priesthood offices and keys. During His mortal ministry, Jesus “chose twelve, whom also he named apostles” (Luke 6:13). That holy apostleship has been restored. Living Apostles again bear witness of the name of Christ to all the world. Restoration of all priesthood offices and keys has been accomplished under the direction of the Lord by heavenly messengers, including John the Baptist, Peter, James, John, Moses, Elias, Elijah, Moroni, and others. Priesthood power extends beyond life through the veil of death into the eternities.

The Apostle Paul once wrote concerning the priesthood: “No man taketh this honour unto himself, but he that is called of God, as was Aaron” (Heb. 5:4). As in former days, it is again conferred by ordination, by the laying on of hands by those in authority. The influence of the priesthood is controlled according to conditions established by the Lord, who said:

“The rights of the priesthood are inseparably connected with the powers of heaven, and … the powers of heaven cannot be controlled nor handled only upon the principles of righteousness.

“That they may be conferred upon [men] it is true; but when [they] undertake to cover [their] sins, or to gratify [their] pride, [their] vain ambition, or to exercise control or dominion or compulsion upon the souls of the children of men, in any degree of unrighteousness, … the Spirit of the Lord is grieved; and when it is withdrawn, Amen to the priesthood or the authority of that man …

“No power or influence can or ought to be maintained by virtue of the priesthood, only by persuasion, by long-suffering, by gentleness and meekness, and by love unfeigned;

“By kindness, and pure knowledge, which shall greatly enlarge the soul without hypocrisy, and without guile” (D\&C 121:36–37, 41–42).

This priesthood authority has divine power to confer sacred ordinances that are binding for all eternity. To illustrate, a husband and wife may be eternally married so that even after death, they and their family can be together forever in the presence of Deity. This enabling ordinance is performed in our holy temples. Furthermore, these and other necessary ordinances of salvation, such as baptism, may be performed by living individuals as proxy for deceased family members.

Hence, as part of our religious responsibility, we identify our ancestors, for as the Apostle Paul said to the Hebrew Saints, “they without us should not be made perfect” (Heb. 11:40). For this purpose, Church family history (genealogical) libraries dot the earth; they are patronized extensively and without fee, both by members of our Church and by those who may have other religious affiliation.

Members of this church practice principles of Christian living by assisting some of the homeless, the hungry, and the needy throughout the world. We strive to practice the admonition of James: “Pure religion and undefiled before God and the Father is this, To visit the fatherless and widows in their affliction” (James 1:27).

Once a month, if members are physically able and desirous of doing so, they voluntarily go without food for two meals and contribute, as a fast offering, the money that would have been spent. These funds are then used to assist the poor and the needy as determined by leaders of local congregations. Contributions in excess of local need are pooled and made available for aid elsewhere. In addition, our members electively donate millions of hours of compassionate service each year.

During 1992, for example, voluntary humanitarian donations enabled the Church to help care for the poor through more than three hundred fifty projects in Asia, Eastern Europe, Africa, Latin America, the Caribbean, Canada, and the United States of America. This help was in the form of food, clothing, shelter, textbooks, medical supplies, volunteer labor, and skilled advisers. In these and other activities, the Church has participated with more than two hundred private volunteer organizations, other churches, and community agencies such as the International Red Cross and Red Crescent Societies, Catholic Relief Services, Food for the Poor, and the Salvation Army, to name but a few.

Most recently, in response to the devastating drought in Eastern Africa, the Church provided approximately \$1 million worth of powdered milk, oil, flour, beans, and clothing to help care for expectant or nursing mothers, young children, and others. Twenty thousand pounds of medical equipment were sent to Albania, where modern medical care is urgently needed. Hundreds of hygiene kits and essential clothing were sent to Croatia to assist the suffering people in that war-torn country.

Each of these humanitarian endeavors is an application of the biblical command “Love thy neighbour as thyself” (Lev. 19:18; see also Luke 10:27).

As we repent and live in accordance with the commandments of God, the Atonement of Jesus Christ becomes a force for permanent change in our lives. We become more pure in heart and true disciples of the Master. Personal prayer, contributions of tithes and offerings, honoring the Sabbath day, and partaking of the sacrament of the Lord’s Supper are all regular practices that protect us from the bondage of sin and bring us true freedom and joy.

The daily practice of our faith is not confined solely to religious rites. We reach out to aid the communities in which we live. For example, parent-teacher associations, public schools, and similar worthy civic activities receive our support.

We value good teachers. We affirm the scriptural admonition, “Teach ye diligently … that you may be instructed more perfectly” (D\&C 88:78). We assert that “the glory of God is intelligence” (D\&C 93:36), and we seek wisdom out of the best books (see D\&C 88:118; D\&C 109:7, 14). We regard one’s education as a religious responsibility with everlasting reward. Scriptures state that “whatever principle of intelligence we attain unto in this life, it will rise with us in the resurrection. And if a person gains more knowledge and intelligence in this life through his diligence and obedience than another, he will have so much the advantage in the world to come” (D\&C 130:18–19).

We neither endorse games of chance nor rely on random favor. Scripture attests that predictable results come only from obedience to divine law:

“There is a law, irrevocably decreed in heaven before the foundations of this world, upon which all blessings are predicated—

“And when we obtain any blessing from God, it is by obedience to that law upon which it is predicated” (D\&C 130:20–21).

As a cardiac surgeon, I can certify that this understanding helped me a great deal. In the operating room, hopes were never enough. I knew that I must remain there and perform with precision all that was required to put things into proper order. Only then could the desired blessings of healing and health be expected to follow.

We know that alcoholic beverages and tobacco are not good for human beings. Our understanding came by way of divine revelation recorded more than a century and a half ago in the year 1833. Moreover, we teach that stimulants and other harmful substances are contrary to the will of God for our bodies. Scientific studies document the benefits of obedience to this code of health. Research from the University of California at Los Angeles published in 1989, for example, demonstrated that the length of life of faithful members of this church was significantly longer than the norm. Latter-day Saint men and women were reported to have lower mortality rates from cancer and heart disease when compared with control population groups.

Such data are of interest because we maintain, as did the Apostle Paul, that our bodies are temples for our spirits, divinely created by our loving Father in Heaven (see 1 Cor. 3:16–17).

We also believe in “being honest, true, chaste, benevolent, virtuous, and in doing good to all men … If there is anything virtuous, lovely, or of good report or praiseworthy, we seek after these things” (A of F 1:13). As a result, we cherish happy and wholesome relationships with our families and others wherever we live. We honor the family as the basis of a virtuous, moral life.

Because these doctrines are so ennobling and fulfilling, we have a genuine desire to share the good news of the gospel with others (see Prov. 25:25). We have many missionaries who serve at the call of their Church leaders. They go wherever called to live and labor among the people. Most of our missionaries are young men and women who serve for one and a half to two years, and do so without financial remuneration, supporting their service from personal savings and contributions from family members and loved ones.

Our missionaries are now teaching the gospel throughout the world in sixty-nine languages. When released, they return home to resume their education and the normal pursuits of life. But they also continue as advocates of the people whose culture, language, and ways they have come to love. They serve because the Lord issued a commandment: “Go ye therefore, and teach all nations, baptizing them in the name of the Father, and of the Son, and of the Holy Ghost” (Matt. 28:19; see also Mark 16:15). On another occasion, the Master stated: “Except a man be born of water and of the Spirit, he cannot enter into the kingdom of God” John 3:5). We comply by performing baptisms by immersion and then conferring the gift of the Holy Ghost by the imposition of hands by those duly authorized. Thus, our missionary efforts are responses to divine decrees.

“We believe in being subject to kings, presidents, rulers, and magistrates, in obeying, honoring, and sustaining the law” (A of F 1:12). We encourage our members to be upright and loyal citizens in the countries that give them citizenship and to uphold governmental leaders and obey the civil laws wherever they live.

We hold to a single standard of morality. Both men and women are required to observe the law of chastity. The commandment “Thou shalt not commit adultery” applies to all. Promiscuity and all extramarital sexual relationships are sinful in the sight of God, and members of the Church who engage in such practices are subject to disciplinary measures.

Before proceeding further, it might be helpful to pause for a brief synopsis of this discussion of doctrines of The Church of Jesus Christ of Latter-day Saints. May I offer ten summary statements. (I apologize that time won’t allow for a more comprehensive list.) We believe:

While implementing these and other important precepts, the Church is tolerant of all faith groups, claiming for itself no right or privilege that it would deny to others. It affirms itself to be Christ’s church of old, reestablished anew. Its message to the world is one of peace and goodwill, with a sincere invitation to come and partake of all the blessings incident to the new and everlasting covenant between God and His children.

\intalkheader{Interfaith Cooperation}

Moving now from the distinctive doctrines and tenets of The Church of Jesus Christ of Latter-day Saints, I would like to discuss areas in which faith groups may cooperate. Members of our church often join with other like-minded citizens, regardless of religious persuasion, in support of worthy causes and humanitarian projects. This can be done without losing independent identity and strength. We are mindful of history’s lesson that attempts of interfaith groups to unify through theological blending have not always been successful. Indeed, when divine doctrines are compromised to accommodate social pressures, religious institutions become no more relevant in their mission than other agencies of good intent. At the same time, we should be models of tolerance of others whose sacred beliefs may differ from our own. We recognize that if one religion is persecuted, all are attacked.

We need to safeguard the proper relationships of religion to governments and to the people we serve. We are under obligation to take a strong stand on moral issues in order to preserve a wholesome and ethical environment. As religious leaders, our participation in the political process may be needed to protect precious religious liberty. This is particularly important in countries of the earth whose citizens, until recently, have been deprived of freedom of religious expression and practice. Many of those nations are now shaping laws to allow freedoms that citizens of other countries have long enjoyed. These developing democracies deserve support.

Selected humanitarian projects can be pursued independently or cooperatively. We can promote education regarding moral values, the arts, service to humanity, and the value of education itself. We can give encouragement for family life, knowing that patience, righteousness, and love will be rewarded with joy.

\intalkheader{Conclusion}

Returning to the medical metaphor with which I began, I firmly feel that the people of this world can be healed as we do our best in combatting spiritual drift—our global pandemic. Healing requires a concurrent development of the will to live, the faith to become well, and the motivation to improve our world for future generations.

May I conclude on a personal note. My expressions represent more than mere belief. I know from personal experience that God can and does communicate with individuals. Through that personal revelation, one may surely know that God lives, that His Beloved Son, Jesus, is the Christ, and that the doctrines I have discussed which are pertinent to man’s salvation and eternal progress are true. Similar strong convictions in the hearts of millions of individual members comprise the real vitality of The Church of Jesus Christ of Latter-day Saints.

After this exposure to the roots of our faith, you deserve to experience some of the good fruits by which you shall also know us. So I have asked several children from the Chicago area, all members of our church, to sing to you. The title of their song is “I Am a Child of God.” As they acknowledge the Fatherhood of God, they also demonstrate the brotherhood of all mankind. Their song will comprise the concluding portion of my presentation. Together we echo a statement previously published by the First Presidency of our church:

“Our message is one of special love and concern for the eternal welfare of all men and women, regardless of religious belief, race, or nationality, knowing that we are truly brothers and sisters because we are the sons and daughters of the same eternal Father.”

To each of you I express heartfelt appreciation along with sincere wishes for your success in all of your worthy endeavors. May God bless you and your loved ones and sustain our mutual quest for spiritual strength, goodwill, and peace on earth, I pray in the name of Jesus Christ, amen.

\newpage
\settalk{Constancy amid Change}
\setspeaker{Russell M. Nelson}
\settalkdate{October 1993}
\talkheading{Constancy amid Change}{Russell M. Nelson}

I echo President Monson’s compliments to Brother Peterson, Brother Komatsu, and Brother de Jager. They have earned our greatest commendation. And I, too, join with others who feel a debt of gratitude to this wonderful chorus of youth from Brigham Young University.

Our youth are wonderful and especially able to ask thoughtful questions. Recently I had a conversation with “Ruth” and “John.” Ruth opened the discussion. With a sigh, she lamented, “Our world is constantly changing, isn’t it?”

“Yes,” I replied, “ever since its creation—geologically and geographically. And its populations are changing—politically and spiritually. You might ask your grandparents about life when they were your age and discover their thoughts.”

“Oh, I already have,” Ruth continued. “My grandpa summarized his opinion with a clever quip: ‘Give me the good old days—plus penicillin.’”

Then John expressed deep concern. “Continually changing conditions make the future shaky for us,” he said. “It’s kind of scary. We seem to be standing on shifting sand.”

Together they asked, “What can we trust? Is anything constant that will not change as we grow older?”

To that question I responded with an emphatic, “Yes! Many things!” Because Ruth and John are typical of many today who seek for unchanging constants in a changing world, I would like to address that subject, titling my remarks “Constancy amid Change.” Through the years, prophets and Apostles have spoken of many unchanging constants. To facilitate this discussion, I will group some of these constants into three categories: heavenly personages, plans, and principles.

\intalkheader{I. Personages}

Our Heavenly Father has a glorified body of flesh and bone, inseparably connected with His spirit. Scriptures state that He is “infinite and eternal, from everlasting to everlasting the same unchangeable God” (D\&C 20:17).

His Beloved Son, Jesus Christ, is our Savior and the chief cornerstone of our religion. “He is the life and the light of the world” (Alma 38:9). “There shall be no other name … nor any other way … whereby salvation can come unto the children of men, only in and through the name of Christ, the Lord Omnipotent” (Mosiah 3:17).

Another personage is the Holy Ghost, whose enduring influence transcends time. Scripture assures that “the Holy Ghost shall be thy constant companion, and thy scepter an unchanging scepter of righteousness and truth; and thy dominion shall be an everlasting dominion, and without compulsory means it shall flow unto thee forever and ever” (D\&C 121:46; emphasis added).

Brothers and sisters, these Heavenly Beings love you. Their love is as constant as is the greatest love of earthly parents.

But there is another personage about whom you should be reminded. Satan also exists and seeks “that all men might be miserable like unto himself” (2 Ne. 2:27).

\intalkheader{II. Plans}

I speak now of category two—unchanging plans. A great council in heaven was once convened, in which it seems that all of us participated. There our Heavenly Father announced His plan. Scriptures refer to this plan of God by many names. Perhaps out of deference to the sacred name of Deity, or to depict its broad scope, it is also called the plan of happiness, the plan of salvation, the plan of redemption, the plan of restoration, the plan of mercy, the plan of deliverance, and the everlasting gospel. Prophets have used these terms interchangeably. Regardless of designation, the enabling essence of the plan is the atonement of Jesus Christ. As it is central to the plan, we should try to comprehend the meaning of the Atonement. Before we can comprehend it, though, we must understand the fall of Adam. And before we can fully appreciate the Fall, we must first comprehend the Creation. These three events—the Creation, the Fall, and the Atonement—are three preeminent pillars of God’s plan, and they are doctrinally interrelated.

The creation of the earth was a preparatory part of our Father’s plan. Then “the Gods went down to organize man in their own image, … male and female to form they them.

“And the Gods said: We will bless them” (Abr. 4:27–28). And bless us they did, with a plan that would give us physical bodies of our very own.

Adam and Eve were the first people to live upon the earth. They were different from the plant and animal life that had been created previously. Adam and Eve were children of God. Their bodies of flesh and bone were made in the express image of God’s. In that state of innocence, they were not yet mortal. They could have had no children, were not subject to death, and could have lived in Eden’s garden forever. Thus, we might speak of the Creation in terms of a paradisiacal creation.

If that state had persisted, you and I would still be stranded among the heavenly host as unborn sons and daughters of God. “The great plan of [happiness] would have been frustrated” (Alma 42:5).

That leads us to the fall of Adam. To bring the plan of happiness to fruition, God issued to Adam and Eve the first commandment ever given to mankind. It was a commandment to beget children. A law was explained to them. Should they eat from “the tree of the knowledge of good and evil” (Gen. 2:17), their bodies would change; mortality and eventual death would come upon them. But partaking of that fruit was prerequisite to their parenthood.

While I do not fully understand all the biochemistry involved, I do know that their physical bodies did change; blood began to circulate in their bodies. Adam and Eve thereby became mortal. Happily for us, they could also beget children and fulfill the purposes for which the world was created. Happily for them, “the Lord said unto Adam [and Eve]: Behold I have forgiven thee thy transgression in the Garden of Eden” (Moses 6:53). We and all mankind are forever blessed because of Eve’s great courage and wisdom. By partaking of the fruit first, she did what needed to be done. Adam was wise enough to do likewise. Accordingly, we could speak of the fall of Adam in terms of a mortal creation, because “Adam fell that men might be” (2 Ne. 2:25).

Other blessings came to us through the Fall. It activated two closely coupled additional gifts from God, nearly as precious as life itself—agency and accountability. We became “free to choose liberty and eternal life … or to choose captivity and death” (2 Ne. 2:27). Freedom of choice cannot be exercised without accountability for choices made.

Now we come to the third pillar of God’s plan—the Atonement. Just as Adam and Eve were not to live forever in the Garden of Eden, so our final destination was not to be planet earth. We were to return to our heavenly home.

Given that reality, still another change was necessary. An infinite atonement was required to redeem Adam, Eve, and all of their posterity. That atonement must enable our physical bodies to be resurrected and changed to a bloodless form, no longer liable to disease, deterioration, or death.

According to eternal law, that atonement required a personal sacrifice by an immortal being not subject to death. Yet He must die and take up His own body again. The Savior was the only one who could accomplish this. From His mother He inherited power to die. From His Father He obtained power over death. The Redeemer so explained:

“I lay down my life, that I might take it again.

“No man taketh it from me, but I lay it down of myself. I have power to lay it down, and I have power to take it again” (John 10:17–18).

The Lord declared that “this is my work and my glory—to bring to pass the immortality and eternal life of man” (Moses 1:39). He who had created the earth came into mortality to fulfill the will of His Father and all prophecies of His atonement. And His atonement redeems every soul from penalties of personal transgression, on the condition of repentance.

Thus, we might speak of the Atonement in terms of the immortal creation. “For as in Adam all die, even so in Christ shall all be made alive” (1 Cor. 15:22).

I have recounted the importance of the Creation, the Fall, and the Atonement, knowing that parents are accountable to teach these precepts of God’s plan to their children.

Before leaving our discussion of unchanging plans, however, we need to remember that the adversary sponsors a cunning plan of his own. It invariably attacks God’s first commandment for husband and wife to beget children. It tempts with tactics that include infidelity, unchastity, and other abuses of procreative power. Satan’s band would trumpet choice but mute accountability. Nevertheless, his capacity has long been limited, “for he knew not the mind of God” (Moses 4:6).

I speak now of category three—unchanging principles.

\intalkheader{III. Principles}

Unchanging principles are so because they come from our unchanging Heavenly Father. Try as they might, no parliament or congress could ever repeal the law of earth’s gravity or amend the Ten Commandments. Those laws are constant. All laws of nature and of God are part of the everlasting gospel. Thus, there are many unchanging principles. Time will permit consideration of only a few.

One of them is that of the priesthood. The Prophet Joseph Smith taught that “the Priesthood is an everlasting principle, and existed with God from eternity, and will to eternity, without beginning of days or end of years” (Teachings of the Prophet Joseph Smith, sel. Joseph Fielding Smith [Salt Lake City: Deseret Book Co., 1938], p. 157).

We know that “the Priesthood was first given to Adam; he obtained the First Presidency, and held the keys of it from generation to generation. He obtained it in the Creation, before the world was formed” (Teachings of the Prophet Joseph Smith, p. 157).

Scriptures certify that the priesthood has continued and will continue “through the lineage of [the] fathers” (D\&C 86:8). Ordination to its offices has timeless implication as well. Tenure in priesthood office may extend into postmortal realms. For example, scriptures declare that one ordained as a high priest may be a high priest forever. Promised blessings of the priesthood extend to men, women, and children throughout the world and may endure forever.

The use of the priesthood is carefully controlled according to conditions established by the Lord, who said:

“No power or influence can or ought to be maintained by virtue of the priesthood, only by persuasion, by long-suffering, by gentleness and meekness, and by love unfeigned” (D\&C 121:41).

“That [the rights of the priesthood] may be conferred upon [men], it is true; but when [they] undertake to cover [their] sins, or to gratify [their] pride, [their] vain ambition, or to exercise control or dominion or compulsion upon the souls of the children of men, in any degree of unrighteousness, … the Spirit of the Lord is grieved; and when it is withdrawn, Amen to the priesthood or the authority of that man” (D\&C 121:37).

While the priesthood is an everlasting principle, those privileged to exercise its authority must maintain themselves daily as worthy vessels.

Another unchanging principle is that of divine or moral law. Transgression of moral law brings retribution; obedience to it brings blessings “immutable and unchangeable” (D\&C 104:2). Blessings are always predicated upon obedience to law. So the Church teaches us to embrace the right and to renounce the wrong—that we might have joy.

The Savior and His servants do not speak words of complacency but teach what people need to know. Through the ages, history attests that contemporary critics have pressed Church leaders to modify a decree of the Lord. But such is eternal law, and it cannot be altered. Not even for His Beloved Son could God change the law that required the Atonement. Divine doctrines cannot be squeezed into compact molds to make them fit fashionable patterns of the day. Nor can they be fully expressed on a bumper sticker.

Another unchanging principle, brothers and sisters, is that of your eventual judgment. Each of you will be judged according to your individual works and the desires of your hearts. You will not be required to pay the debt of any other. Your eventual placement in the celestial, terrestrial, or telestial kingdom will not be determined by chance. The Lord has prescribed unchanging requirements for each. You can know what the scriptures teach, and pattern your lives accordingly.

Other unchanging principles include divine commandments—even those that seem to be temporal. Tithing, for example, is not temporal (or temporary); it is an everlasting principle. The Lord said:

“Those who have thus been tithed shall pay one-tenth of all their interest annually; and this shall be a standing law unto them forever” (D\&C 119:4; emphasis added).

We know that tithe payers shall not be burned at the Second Coming.

Another unchanging principle is that of truth. Scripture reminds us that “the truth abideth forever and ever” (D\&C 1:39). Even though one’s understanding of the truth may be fragmentary, truth itself does not change. Everlasting truth and wisdom come from the Lord. The first truth ever taught to man came directly from Deity. From generation to generation, God has given additional light. Whether truth comes from a laboratory of science or directly by revelation, truth is embraced by the gospel.

May I mention one more everlasting principle—the family. A family can be together forever. Though each of us will pass through the doors of death, the timing of that departure is less important than is the preparation for eternal life. Part of that preparation includes service in the Church. It is not to be a burden but a blessing to a family. The Lord said, “Thy duty is unto the church forever, and this because of thy family” (D\&C 23:3).

Ruth, John, and each of you will more fully understand that concept in light of this scriptural promise:

“If a man marry a wife by my word, which is my law, and by the new and everlasting covenant, and it is sealed unto them … [they] shall inherit thrones, kingdoms, principalities, and powers, dominions, … exaltation and glory in all things, … which glory shall be a fulness and a continuation of the seeds forever and ever” (D\&C 132:19).

A promise like that is worth your personal effort and endurance.

Constancy amid change is assured by heavenly personages, plans, and principles. Our trust can be safely anchored to them. They provide peace, eternal progression, hope, freedom, love, and joy to all who will be guided by them. They are true—now and forever—I testify in the name of Jesus Christ, amen.

\newpage
\settalk{“Teach Us Tolerance and Love”}
\setspeaker{Russell M. Nelson}
\settalkdate{April 1994}
\talkheading{“Teach Us Tolerance and Love”}{Russell M. Nelson}

Dear brothers and sisters, I join my brethren in extending Easter greetings to each of you, while expressing personal gratitude for the atonement of Jesus Christ, for His example, and for His teachings that have motivated my message today.

I have been impressed to speak on the subject of tolerance—a virtue much needed in our turbulent world. But in discussing this topic, we must recognize at the outset that there is a difference between tolerance and tolerate. Your gracious tolerance for an individual does not grant him or her license to do wrong, nor does your tolerance obligate you to tolerate his or her misdeed. That distinction is fundamental to an understanding of this vital virtue.

I attended a “laboratory of tolerance” some months ago when I had the privilege of participating in the Parliament of the World’s Religions. There I conversed with good men and women representing many religious groups. Again I sensed the advantages of ethnic and cultural diversity and reflected once more on the importance of religious freedom and tolerance.

I marvel at the inspiration of the Prophet Joseph Smith when he penned the eleventh article of faith: “We claim the privilege of worshiping Almighty God according to the dictates of our own conscience, and allow all men the same privilege, let them worship how, where, or what they may.”

That noble expression of religious tolerance is particularly poignant in light of the Prophet’s personal persecution. On one occasion he wrote, “I am at this time persecuted the worst of any man on the earth, as well as this people, … and all our sacred rights are trampled under the feet of the mob.”

Joseph Smith endured incessant persecution and finally heartless martyrdom—at the hands of the intolerant. His brutal fate stands as a stark reminder that we must never be guilty of any sin sown by the seed of intolerance.

\intalkheader{Two Great Commandments to Love}

Revealed to that revered prophet was the fulness of the gospel. He was tutored by the resurrected Christ, whom Joseph adored. He taught doctrines declared by the Lord, including these He gave in response to the question of an exacting lawyer:

“Master, which is the great commandment in the law?

“Jesus said unto him, Thou shalt love the Lord thy God with all thy heart, and with all thy soul, and with all thy mind.

“This is the first and great commandment.

“And the second is like unto it, Thou shalt love thy neighbour as thyself.

“On these two commandments hang all the law and the prophets.”

Hence, our highest priorities in life are to love God and to love our neighbors. That broadly includes neighbors in our own family, our community, our nation, and our world. Obedience to the second commandment facilitates obedience to the first commandment. “When ye are in the service of your fellow beings ye are only in the service of your God.”

\intalkheader{Parental Love}

That concept is easy for mothers and fathers to understand. Parental love includes gratitude for service extended to any of their children, especially in their time of need.

I was amused recently when one of our grown children confided that she had always thought that she was her daddy’s favorite daughter. She was surprised to discover later that each of her eight sisters harbored that same feeling. Only when they had become mothers themselves did they realize that parents hardly have favorites. (Incidentally, our only son never had to wonder who was our favorite son.)

Our Father in Heaven loves all of His children, too. Peter taught that “God is no respecter of persons:

“But in every nation he that feareth him, and worketh righteousness, is accepted with him.”

Yet His children can be so intolerant with one another. Neighboring factions, whether they be identified as groups or gangs, schools or states, counties or countries, often develop animosity. Such tendencies make me wonder: Cannot boundary lines exist without becoming battle lines? Could not people unite in waging war against the evils that beset mankind instead of waging war on each other? Sadly, answers to these questions are often no. Through the years, discrimination based on ethnic or religious identity has led to senseless slaughter, vicious pogroms, and countless acts of cruelty. The face of history is pocked by the ugly scars of intolerance.

How different our world would be if all parents would apply this inspired instruction from the Book of Mormon: “Ye will not suffer your children … that they transgress the laws of God, and fight and quarrel one with another. …

“But ye will teach them to walk in the ways of truth and soberness; ye will teach them to love one another, and to serve one another.”

If such training occurred, children and parents around this globe would join in singing, “Fill our hearts with sweet forgiving; Teach us tolerance and love.” Men and women would respect their neighbors and the beliefs held sacred by them. No longer would ethnic jokes and cultural slurs be acceptable. The tongue of the tolerant speaks no guile.

\intalkheader{Independence and Cooperation}

While we strive for the virtue of tolerance, other commendable qualities need not be lost. Tolerance does not require the surrender of noble purpose or of individual identity. The Lord gave instruction to leaders of His restored church to establish and maintain institutional integrity—“that the church may stand independent.”

Meanwhile, its members are encouraged to join with like-minded citizens in doing good. We are grateful for the many examples of heroic service rendered in times of earthquakes, floods, hurricanes, or other disasters. Such cooperative efforts to help neighbors in distress transcend any barriers posed by religion, race, or culture. Those good deeds are latter-day love in action!

Humanitarian relief rendered by members of this church is extensive, multinational, and generally unpublicized. Even so, there are doubtless many who wonder why we don’t do more to assist the innumerable worthy causes to which our hearts respond.

Of course we are concerned with the need for ambulances in the valley below. But at the same time, we cannot ignore the greater need for protective guardrails on the cliffs above. Limited resources needed for the accomplishment of the higher work cannot be depleted in rescue efforts that provide only temporary relief.

The biblical prophet Nehemiah must have felt that same commitment to his important calling. When he was asked to divert attention away from his primary purpose, he replied, “I am doing a great work, so that I cannot come down: why should the work cease, whilst I leave it, and come down to you?”

Fortunately, we in the Church rarely have to make such a decision. We consider love of neighbor an integral part of our mission. And while we serve one another, we continue to build a spiritual house of refuge on the cliffs above. Such a sanctuary becomes a blessing for all mankind. We are but the builders; the architect is almighty God.

\intalkheader{Missionary Responsibilities}

Latter-day Saints throughout the world work side by side with others—regardless of race, color, or creed—hoping to be good examples worthy of emulation. The Savior said: “I give unto you a commandment, that every man, both elder, priest, teacher, and also member, … prepare and accomplish the things which I have commanded.

“And let your preaching be the warning voice, every man to his neighbor, in mildness and in meekness.”

This we are to do with tolerance. While in Moscow in June 1991, in that spirit of preparation and with sincere respect for leaders of other religious denominations, Elder Dallin H. Oaks and I had the privilege of meeting with the presiding official of the Russian Orthodox Church. We were accompanied by Elder Hans B. Ringger and the mission president, Gary L. Browning. Patriarch Aleksei was most gracious in sharing a memorable hour with us. We perceived the great difficulties endured for so many years by this kind man and his fellow believers. We thanked him for his perseverance and for his faith. Then we assured him of our good intentions and of the importance of the message that missionaries of The Church of Jesus Christ of Latter-day Saints would be teaching among his countrymen. We affirmed that ours is a global church and that we honor and obey the laws of each land in which we labor.

To those with an interest in the fulness of the restored gospel—regardless of nationality or religious background—we say as did Elder Bruce R. McConkie: “Keep all the truth and all the good that you have. Do not abandon any sound or proper principle. Do not forsake any standard of the past which is good, righteous, and true. Every truth found in every church in all the world we believe. But we also say this to all men—Come and take the added light and truth that God has restored in our day. The more truth we have, the greater is our joy here and now; the more truth we receive, the greater is our reward in eternity. This is our invitation to men [and women] of good will everywhere.”

Each of you with a testimony of the truth of the restored gospel has opportunity to share that precious gift. The Lord expects you to “be ready always to give an answer to every man that asketh you a reason of the hope that is in you with meekness.”

\intalkheader{Baptism Transcends Background}

On every continent and across isles of the sea, the faithful are being gathered into The Church of Jesus Christ of Latter-day Saints. Differences in cultural background, language, gender, and facial features fade into insignificance as members lose themselves in service to their beloved Savior. Paul’s declaration is being fulfilled: “As many of you as have been baptized into Christ have put on Christ.

“There is neither Jew nor Greek, there is neither bond nor free, there is neither male nor female: for ye are all one in Christ Jesus.”

Only the comprehension of the true Fatherhood of God can bring full appreciation of the true brotherhood of man. That understanding inspires desire to build bridges of cooperation instead of walls of segregation.

Our Creator decreed “that there should be no contention one with another, but that they should look forward with one eye, having one faith and one baptism, having their hearts knit together in unity and in love one towards another.”

Intolerance seeds contention; tolerance supersedes contention. Tolerance is the key that opens the door to mutual understanding and love.

\intalkheader{Risks of Boundless Tolerance}

Now may I offer an important note of caution. An erroneous assumption could be made that if a little of something is good, a lot must be better. Not so! Overdoses of needed medication can be toxic. Boundless mercy could oppose justice. So tolerance, without limit, could lead to spineless permissiveness.

The Lord drew boundary lines to define acceptable limits of tolerance. Danger rises when those divine limits are disobeyed. Just as parents teach little children not to run and play in the street, the Savior taught us that we need not tolerate evil. “Jesus went into the temple of God … and overthrew the tables of the moneychangers.” Though He loved the sinner, the Lord said that He “cannot look upon sin with the least degree of allowance.” His Apostle Paul specified some of those sins in a letter to the Galatians. The list included “adultery, fornication, uncleanness, lasciviousness,

“Idolatry, witchcraft, hatred, … wrath, strife, seditions, heresies,

“Envyings, murders, drunkenness, revellings, and such like.”

To Paul’s list I might add the regrettable attitudes of bigotry, hypocrisy, and prejudice. These were also decried in 1834 by early Church leaders who foresaw the eventual rise of this church “amid the frowns of bigots and the calumny of hypocrites.” The Prophet Joseph Smith prayed that “prejudices may give way before the truth.” Hatred stirs up strife and digs beneath the dignity of mature men and women in our enlightened era.

Paul’s list included “uncleanness.” As members of the Church entrusted with its holy temples, we are commanded that “no unclean thing shall be permitted to come into [His] house to pollute it.”

That assignment requires great fortitude as well as love. In former days, disciples of the Lord “were firm, and would suffer even unto death rather than commit sin.” In latter days, devoted disciples of the Lord are just as firm. Real love for the sinner may compel courageous confrontation—not acquiescence! Real love does not support self-destructive behavior.

\intalkheader{Tolerance and Mutual Respect}

Our commitment to the Savior causes us to scorn sin yet heed His commandment to love our neighbors. Together we live on this earth, which is to be tended, subdued, and shared with gratitude. Each of us can help to make life in this world a more pleasant experience. Not long ago the First Presidency and the Twelve approved a public statement from which I quote: “It is morally wrong for any person or group to deny anyone his or her inalienable dignity on the tragic and abhorrent theory of racial or cultural superiority.

“We call upon all people everywhere to recommit themselves to the time-honored ideals of tolerance and mutual respect. We sincerely believe that as we acknowledge one another with consideration and compassion we will discover that we can all peacefully coexist despite our deepest differences.”

That pronouncement is a contemporary confirmation of the Prophet Joseph’s earlier entreaty for tolerance. Unitedly we may respond. Together we may stand, intolerant of transgression but tolerant of neighbors with differences they hold sacred. Our beloved brothers and sisters throughout the world are all children of God. He is our Father. His Son, Jesus, is the Christ. His church has been restored to the earth in these latter days to bless all of God’s children. I so testify in the name of Jesus Christ, amen.

\newpage
\settalk{The Spirit of Elijah}
\setspeaker{Russell M. Nelson}
\settalkdate{October 1994}
\talkheading{The Spirit of Elijah}{Russell M. Nelson}

Within the past year, four of our beloved General Authorities have concluded their mission in mortality. The passing of President Ezra Taft Benson and Elders Marvin J. Ashton, Sterling W. Sill, and Clinton L. Cutler has prompted a great outpouring of love to their families. Their departure has also underscored a scriptural explanation of activities on the other side of the veil: “The faithful elders of this dispensation, when they depart from mortal life, continue their labors in the preaching of the gospel … in the great world of the spirits of the dead.”

President Spencer W. Kimball taught that “our great part in this aspect of missionary work is to perform on this earth the ordinances required for those who accept the gospel over there.”

These thoughts, coupled with the fact that next month marks the one hundredth anniversary of the Genealogical Society of Utah, highlight the great importance and influence of “the spirit of Elijah.” Happily, the date of this centennial closely coincides with the birthday of President Howard W. Hunter, who once presided over that society and who now beckons us to the house of the Lord.

\intalkheader{Baptism Is Essential to Enter the Kingdom of God}

Basic to all Christian understanding is the timeless statement made by Jesus: “Except a man be born of water and of the Spirit, he cannot enter into the kingdom of God.”

Following His crucifixion, Jesus ministered in the spirit world, setting in motion missionary work among those who had died without hearing the gospel. Baptism for these souls would logically be expected. Yet only one verse in the New Testament refers directly to that need: “Else what shall they do which are baptized for the dead, if the dead rise not at all? why are they then baptized for the dead?”

Those questions of the Apostle Paul, without latter-day revelation, would remain an enigma. With latter-day revelation, they become clear. Clarification began when the Prophet Joseph Smith was tutored by the angel Moroni, who said:

“I will reveal unto you the Priesthood, by the hand of Elijah the prophet, before the coming of the great and dreadful day of the Lord.

“And he shall plant in the hearts of the children the promises made to the fathers, and the hearts of the children shall turn to their fathers.

“If it were not so, the whole earth would be utterly wasted.”

Joseph compared Moroni’s teaching to a similar prophecy by Malachi—that Elijah would come again. We know that Elijah did return—at least twice—after Malachi’s promise. At Christ’s transfiguration, Elijah appeared on the mount to Peter, James, and John. At the Kirtland Temple, April 3, 1836, Elijah appeared to the Prophet and Oliver Cowdery and said, “The keys of this dispensation are committed into your hands.”

The Prophet Joseph Smith was chosen to reestablish the Church, to receive and administer priesthood authority, and to restore plain and precious truths lost to human knowledge. One such doctrine was that of salvation for the dead—a pivotal part of the promised “restitution of all things.”

Several years ago Elder Howard W. Hunter said: “Does it seem reasonable that persons who have lived upon the earth and died without the opportunity of baptism should be deprived throughout eternity? Is there anything unreasonable about the living performing the baptisms for the dead? Perhaps the greatest example of vicarious work for the dead is the Master himself. He gave his life as a vicarious atonement, that all who die shall live again and have life everlasting. He did for us what we could not do for ourselves. In a similar way we can perform ordinances for those who did not have the opportunity to do them in lifetime.”

Elder Hunter added: “Not only may baptisms be performed for the dead, but endowments; also sealings, by which wives become eternal companions to husbands and their children sealed to them as a family. The sealing of family units can be continued until the family of God is made perfect. This is the great work of the dispensation of the fulness of times. … The uniting and redemption of the family of God was the divine plan before the foundations of the earth were laid.”

President Hunter’s classic statement emphasizes the importance of temple work for our own families and helps us to understand the Old Testament prophecy that “saviours shall come up on mount Zion.” This exalting service for others unseen is one of the most noble acts of human kindness.

\intalkheader{Temple Work—Ancient and Modern}

From the days of Adam to the meridian of time, temple ordinances were performed for the living only. Ordinances for the dead had to await the Atonement and postmortal ministry of the Savior.

There was no provision for baptism for the dead when the Kirtland Temple was designed. Yet it served an important preparatory purpose. One week after its dedication, the Lord came personally to the temple to accept it. Then, under His direction, Moses, Elias, and Elijah restored specific keys of priesthood authority.

Five years later, the Saints were in Nauvoo, Illinois. There the Lord again commanded them to build a temple—this time with additional facilities—because, He said, “a baptismal font there is not upon the earth, that they, my saints, may be baptized for those who are dead—

“For this ordinance belongeth to my house.”

Then, to make certain that there could be no misunderstanding, He issued a solemn word of warning: “If you do not these [baptisms for the dead] ye shall be rejected as a church, with your dead, saith the Lord.”

Though the Nauvoo Temple was destroyed by fire, it served its sacred purposes.

\intalkheader{Genealogical Records for Use in Temples}

Throughout the world, members of the Church faithfully prepare family records for use in our many temples. When ordinances are performed there, further documentation is required because the Lord said:

“When any of you are baptized for your dead, let there be a recorder, and let him be eye-witness of your baptisms; …

“That in all your recordings it may be recorded in heaven; whatsoever you bind on earth, may be bound in heaven; whatsoever you loose on earth, may be loosed in heaven.”

This weighty doctrine pressed itself upon the Prophet’s mind. His thoughts resonated with those of previous prophets. Joseph wrote, “John the Revelator was contemplating this very subject in relation to the dead, when he declared, … I saw the dead … stand before God; and the books were opened; … and the dead were judged out of those things which were written in the books.”

Then the Prophet Joseph added, “Whatsoever you do not record on earth shall not be recorded in heaven; for out of the books shall your dead be judged.”

\intalkheader{Elijah and Keys of Priesthood Authority}

In 1844, Joseph Smith asked, “What is this office and work of Elijah?” The Prophet promptly answered his own question: “It is one of the greatest and most important subjects that God has revealed. …

“… This is the spirit of Elijah, that we redeem our dead, and connect ourselves with our fathers which are in heaven. … This is the power of Elijah and the keys of the kingdom of Jehovah.”

Some among us still have neither perceived the Spirit of Elijah nor its power. Yet we are bound by this warning:

“These are principles in relation to the dead and the living that cannot be lightly passed over. … For their salvation is necessary and essential to our salvation, [for] they without us cannot be made perfect—neither can we without our dead be made perfect.”

Joseph Smith’s responsibility was to “lay the foundation” for this great work. Important details were to be revealed later. At April conference 1894, President Wilford Woodruff announced this revelation: “We want the Latter-day Saints from this time to trace their genealogies as far as they can, and to be sealed to their fathers and mothers. Have children sealed to their parents, and run this chain through as far as you can get it. … This is the will of the Lord to his people.”

\intalkheader{One Hundredth Anniversary of the Genealogical Society of Utah}

That revelation in April led to the organization of the Genealogical Society of Utah the following November 13, 1894. Its objectives were “to establish and maintain a Genealogical Library … ; to teach members how to compile acceptable family records and to trace their pedigrees; and to foster temple ordinances.” Events of that historic year established family history research and temple service as one work in the Church.

In the century since then, much has been accomplished. More and more people are becoming excited about discovering their roots, and the Church is doing its best to help them. The Church adopted the term family history to encourage this activity among all its members, especially those who might be intimidated by the word genealogy. In addition, 2,150 busy and productive family history centers have been established throughout the world. For example, the FamilySearch® Center in the Joseph Smith Memorial Building has served hundreds of thousands of visitors, at least two-thirds of whom have found something in the computer file about their ancestors.

More than three hundred thousand copies of the Church’s computer program Personal Ancestral File® are used in libraries and homes by many hundreds of thousands of people. Our FamilySearch program is used by millions of genealogical researchers throughout the world, most of whom are not members of our church. TempleReady® allows convenient and almost immediate clearance by personal computer of names for temple ordinances that formerly required much time and labor.

Sister Nelson, our family, and I have submitted our own ancestral names to the temple and have performed ordinances for them. Because we are fortunate to live near a temple, we like to meet there early in the morning. Usually in less than an hour, the initiatory work is accomplished, our youth are taken directly to school, their mothers return home, and their fathers get to work—on time! When we do endowments or sealings, available adults prefer to meet early in the evening to share that choice experience. Following that, we gather at home to update our records and enjoy some of Sister Nelson’s homemade goodies.

We are also doing temple work for ancestors of a Russian convert to the Church who is not able to travel to a temple. While our son was serving on his mission in Russia, this devoted convert entrusted records of his relatives to our son, along with a plea that their temple work be done. When our children and grandchildren go to the temple to perform those ordinances, our son’s help is needed for pronunciation of names, but not for the perception of joy among all participants.

Service in the temple together is a sublime activity for a family. It provides its own sustaining motivation and verification of the truth of this unique work.

\intalkheader{Rights of Personal Choice, Accountability, and Privacy}

For whom will such temple work be efficacious? Principles of agency pertain on both sides of the veil. There, in postmortal realms, personal choice and accountability will be of paramount importance. Not all will accept these ordinances. Not all who would choose to do so may be worthy to receive them. Scriptures indicate that individual faith, repentance, and obedience will be required to consummate this vicarious work.

Here, on this side of the veil, there are limitations of available time and temples. This means that choosing to identify and perform ordinances for our own kindred should receive our highest priority. The Spirit of Elijah will inspire individual members of the Church to link their generations rather than submit lists of people or popular personalities to whom they are unrelated.

Now, we are mindful of those not of our faith who are concerned about or even offended by the practice of temple ordinances for the dead. To them we say, our Heavenly Father directed the restoration of keys of priesthood authority and surely intended no offense to any of His children. Quite to the contrary. He intended to bless them. This doctrine and its ordinances are laden with love and are intended to perpetuate the sweetest of all relationships—in families forever.

Nevertheless, the Church is sensitive to these concerns. The First Presidency has asked that, as far as possible, individual rights of privacy be protected. In 1972, they wrote, “Persons submitting names for other than direct ancestors [should] have obtained approval from the closest living relative of the deceased before submitting records of persons born within the last ninety-five years.” In addition, reminders of rights of precedence and privacy appear each time our computer programs are used.

Meanwhile, as a gesture of generosity and goodwill, leaders of the Church continue to make its family history facilities available to interested individuals, regardless of religious affiliation and without admission fees! All patrons, in turn, are invited to make valuable additions to the world’s ever-expanding pool of genealogical information.

\intalkheader{Opportunities to Serve}

In a recent statement, President Howard W. Hunter included these remarks: “Let us be a temple-attending and a temple-loving people. Let us hasten to the temple … not only for our kindred dead, but let us also go for the personal blessing of temple worship.”

President Hunter’s invitation reminds us that we can provide names and ordinances for ancestors for whom information is readily available, and, where possible, we can regularly attend the temple. What and how much we do should depend upon personal circumstances and abilities, direction from Church leaders, and guidance from the Spirit. Throughout our lives, each of us can do something significant.

I would add that the daily building of happy memories in our families is an important part of making family history pleasant. Each day on earth can bring a little bit of heaven.

Many travel the highways of life without a companion. They too are needed by their families on both sides of the veil. Others may never be able to attend a temple during their mortal lifetime. To the faithful, comfort comes from the knowledge that no blessings will be withheld from any who love the Lord and strive earnestly to keep His commandments. We will be judged by our deeds and the desires of our hearts—in the Lord’s merciful way and time.

No mortal mind could have conceived this divine work. It is evidence of the restoration of the gospel in its fulness and is sparked by the Spirit of Elijah. “Let us, therefore, as a church and a people … offer unto the Lord an offering in righteousness; and let us present in his holy temple … a book containing the records of our dead … worthy of all acceptation.” Then we shall bless and be blessed as saviors upon mount Zion, I testify in the name of Jesus Christ, amen.

\newpage
\settalk{Children of the Covenant}
\setspeaker{Russell M. Nelson}
\settalkdate{April 1995}
\talkheading{Children of the Covenant}{Russell M. Nelson}

The title of my message is the scriptural phrase “children of the covenant.” In introducing this topic I will reflect on recent events as a colleague of President Howard W. Hunter and as a father and upon earlier experiences as a doctor of medicine.

These past weeks have been challenging for Sister Nelson and me. Not only have we bid farewell to our beloved President Hunter, but thirty-three days earlier, we suffered the demise of our precious daughter Emily. A mother of five young children, Emily had just celebrated her thirty-seventh birthday when called to the other side.

President Hunter influenced Emily’s life in a real way. She welcomed his invitation for all adult members of the Church to hold a temple recommend. She and her husband, Bradley Wittwer, regarded their regular time in the temple as a sacred privilege. They viewed “the temple of the Lord as the great symbol of their membership and the supernal setting for their most sacred covenants.” She strived to emulate the example of the Lord Jesus Christ.

Even though illness brought intense suffering to President Hunter and Emily, an angry word never fell from their lips. Instead, they chose to endure with loving faith. When well-meaning friends and family expressed concern for Emily, she cheerfully replied, “Don’t worry, I’ll be OK!” Even when she concluded a telephone call, she did not close with the customary “good-bye.” She would say, “I love you!”

When President Boyd K. Packer and I last visited President Hunter, he beckoned for Sister Hunter, reached for her hand, and said with a smile, “I feel better when you are near me.”

My tears of sorrow have flowed along with wishes that I could have done more for our daughter and for our President. If I had the power of resurrection, I would have been tempted to bring them back. Though one of the ordained Apostles, each of whom is entrusted with all the keys of the kingdom of God, I do not hold keys of the Resurrection. Jesus Christ holds those keys and will use them for Emily, for President Hunter, and for all people in the Lord’s own time.

Emily and President Hunter had no fear of death. They had made and honored sacred covenants with the Lord, and they knew that his covenants to them will be kept with equal fidelity. They lived nobly as “children of the covenant.”

Years ago as a young medical student I saw many patients afflicted with diseases that are now preventable. Today it is possible to immunize individuals against conditions that once were disabling—even deadly. One medical method by which acquired immunity is conferred is inoculation. The term inoculate is fascinating. It comes from two Latin roots: in, meaning “within”; and oculus, meaning “an eye.” The verb to inoculate, therefore, literally means “to put an eye within”—to monitor against harm.

An affliction like polio can cripple or destroy the body. An affliction like sin can cripple or destroy the spirit. The ravages of polio can now be prevented by immunization, but the ravages of sin require other means of prevention. Doctors cannot immunize against iniquity. Spiritual protection comes only from the Lord—and in his own way. Jesus chooses not to inoculate, but to indoctrinate. His method employs no vaccine; it utilizes the teaching of divine doctrine—a governing “eye within”—to protect the eternal spirits of his children.

\intalkheader{Identification and Indoctrination}

In so teaching, Jesus often established his own identity, then the identity of his followers. I quote his words to the people of ancient America. He said, “I am Jesus Christ, the Son of God.”

“All the prophets from Samuel and those that follow after … have testified of me.

“And behold, ye are the children of the prophets; and ye are of the house of Israel; and ye are of the covenant which the Father made with your fathers, saying unto Abraham: And in thy seed shall all the kindreds of the earth be blessed.

“The Father having raised me up unto you first, and sent me to bless you in turning away every one of you from his iniquities; and this because ye are the children of the covenant.”

A giant step toward spiritual immunity is taken when we understand the expression “children of the covenant.” To what covenant did the Savior refer? “The covenant which he made with Abraham.” The Lord added, “I will remember the covenant which I have made with my people; and I have covenanted with them that I would gather them together in mine own due time.”

\intalkheader{The Abrahamic Covenant}

The covenant that the Lord first made with Abraham and reaffirmed with Isaac and Jacob is of transcendent significance. It contained several promises:

Some of these promises have been accomplished; others have yet to be. I quote from a prophecy given nearly 600 years b.c.: “Our father hath not spoken of our seed alone, but also of all the house of Israel, pointing to the covenant which should be fulfilled in the latter days; which covenant the Lord made to our father Abraham.”

Precisely as promised, the Master appeared in these latter days to renew the Abrahamic covenant. To the Prophet Joseph Smith the Lord declared: “Abraham received promises concerning his seed, and of the fruit of his loins—from whose loins ye are, … my servant Joseph. … This promise is yours also, because ye are of Abraham.”

We are also children of the covenant. We have received, as did they of old, the holy priesthood and the everlasting gospel. Abraham, Isaac, and Jacob are our ancestors. We are of Israel. We have the right to receive the gospel, blessings of the priesthood, and eternal life. Nations of the earth will be blessed by our efforts and by the labors of our posterity. The literal seed of Abraham and those who are gathered into his family by adoption receive these promised blessings—predicated upon acceptance of the Lord and obedience to his commandments.

Elijah the prophet came to plant a knowledge of these promises made to the fathers. Later the Book of Mormon came forth as a sign that the Lord had commenced to gather children of the covenant. This book, written for our day, states: “Then ye may know that the covenant which the Father hath made with the children of Israel … is already beginning to be fulfilled. …

“… For behold, the Lord will remember his covenant which he hath made unto his people of the house of Israel.”

\intalkheader{The New and Everlasting Covenant}

Indeed, the Lord has not forgotten us. And to ensure that we do not forget him, children of the covenant receive his doctrine and claim it by covenant. Brigham Young said: “All Latter-day Saints enter the new and everlasting covenant when they enter this Church. … They enter the new and everlasting covenant to sustain the Kingdom of God and no other kingdom.”

At baptism, we covenant to serve the Lord and keep his commandments. When we partake of the sacrament, we renew those covenants. We may receive covenants of the priesthood and the crowning blessings of the endowment, the doctrine, and the covenants unique to the holy temple.

The new and everlasting covenant of the gospel allows us to qualify for marriage in the temple and be blessed to “come forth in the first resurrection” and “inherit thrones, kingdoms, principalities, and powers, dominions, … to [our] exaltation and glory in all things.”

Children born to parents thus married are natural heirs to the blessings of the priesthood. They are born in the covenant. Hence, “they require no rite of adoption or sealing to insure them place in the posterity of promise.”

Rewards for obedience to the commandments are almost beyond mortal comprehension. Here, children of the covenant become a strain of sin-resistant souls. And hereafter, President Hunter, Emily, other children of the covenant, and “each generation would be linked to the one which went on before [in] the divine family of God.” Great comfort comes from the knowledge that our loved ones are secured to us through the covenants.

\intalkheader{Unity among Children of the Covenant}

Latter-day Saints understand the word of the Lord, who declared, “I say unto you, be one; and if ye are not one ye are not mine.”

“This great unity is the hallmark of the true church of Christ,” said President Gordon B. Hinckley. “It is felt among our people throughout the world.” President Hinckley continued, “We pray for one another that we may go on in unity and strength.”

Throughout the world, however, strident voices are engaged in divisive disputation and name-calling. Often demeaning nicknames are added to—or even substituted for—given names. Unfortunately, terms of derision obscure the true identity of children of the covenant.

In contrast, God employs names that unify and sanctify. When we embrace the gospel and are baptized, we are born again and take upon ourselves the sacred name of Jesus Christ. We are adopted as his sons and daughters and are known as brothers and sisters. He is the Father of our new life. We become joint heirs to promises given by the Lord to Abraham, Isaac, Jacob, and their posterity.

Peter used uplifting terms in a prophecy regarding our day. He identified members of the Church as “a chosen generation, a royal priesthood, an holy nation, a peculiar people.” The adjectives chosen, royal, and holy we recognize as elevating. But what about peculiar? A modern dictionary defines peculiar as “unusual,” “eccentric,” or “strange.” What kind of compliment is that?

But the term peculiar as used in the scriptures is quite different. In the Old Testament, the Hebrew term from which peculiar was translated is segullah, which means “valued property” or “treasure.” In the New Testament, the Greek term from which peculiar was translated is peripoiesis, which means “possession” or “an obtaining.”

Thus we see that the scriptural term peculiar signifies “valued treasure,” “made” or “selected by God.” For us to be identified by servants of the Lord as his peculiar people is a compliment of the highest order.

When we know who we are and what God expects of us—when his “law [is] written in [our] hearts”—we are spiritually protected. We become better people. When the Nephites were truly righteous, they avoided divisive nicknames and “there was no contention in the land, because of the love of God which did dwell in the hearts of the people.”

“There were no … Lamanites, nor any manner of -ites; but they were in one, the children of Christ, and heirs to the kingdom of God.”

That lesson from history suggests that we also delete from our personal vocabularies names that segregate and hyphens that separate. Paul taught that “there is neither Jew nor Greek, there is neither bond nor free, there is neither male nor female: for ye are all one in Christ Jesus.”

He invites us “to come unto him and partake of his goodness; and he denieth none that come unto him, black and white, bond and free, male and female; … all are alike unto God.”

The Church of Jesus Christ of Latter-day Saints has been restored in these latter days to fulfill ancient promises of the Lord. It is part of the “restitution of all things.” Committed children of the covenant remain steadfast, even in the midst of adversity. We shall “be chastened and tried, even as Abraham, who was commanded to offer up his only son.” Yet we are strengthened by this promise of the Lord:

“Ye are lawful heirs, according to the flesh, and have been hid from the world with Christ in God—

“Therefore your life and the priesthood have remained, and must needs remain through you and your lineage until the restoration of all things. …

“Therefore, blessed are ye if ye continue in my goodness, a light unto the Gentiles, and through this priesthood, a savior unto my people Israel.”

With that doctrine implanted deeply within our souls, the sting of death is soothed and spiritual protection is provided. Children of the covenant will be blessed—here and hereafter—I testify in the name of Jesus Christ, amen.

\newpage
\settalk{Perfection Pending}
\setspeaker{Russell M. Nelson}
\settalkdate{October 1995}
\talkheading{Perfection Pending}{Russell M. Nelson}

If I were to ask which of the Lord’s commandments is most difficult to keep, many of us might cite Matthew 5:48: “Be ye therefore perfect, even as your Father which is in heaven is perfect.”

Keeping this commandment can be a concern because each of us is far from perfect, both spiritually and temporally. Reminders come repeatedly. We may lock keys inside the car or even forget where the car is parked. And not infrequently we walk intently from one part of the house to another, only to forget the reason for the errand.

When comparing one’s personal performance with the supreme standard of the Lord’s expectation, the reality of imperfection can at times be depressing. My heart goes out to conscientious Saints who, because of their shortcomings, allow feelings of depression to rob them of happiness in life.

We all need to remember: men are that they might have joy—not guilt trips! We also need to remember that the Lord gives no commandments that are impossible to obey. But sometimes we fail to comprehend them fully.

Our understanding of perfection might be aided if we classify it into two categories. The first could pertain uniquely to this life—mortal perfection. The second category could pertain uniquely to the next life—immortal or eternal perfection.

\intalkheader{Mortal Perfection}

In this life, certain actions can be perfected. A baseball pitcher can throw a no-hit, no-run ball game. A surgeon can perform an operation without an error. A musician can render a selection without a mistake. One can likewise achieve perfection in being punctual, paying tithing, keeping the Word of Wisdom, and so on. The enormous effort required to attain such self-mastery is rewarded with a deep sense of satisfaction. More importantly, spiritual attainments in mortality accompany us into eternity.

James gave a practical standard by which mortal perfection could be measured. He said, “If any man offend not in word, the same is a perfect man.”

Scriptures have described Noah, Seth, and Job as perfect men. No doubt the same term might apply to a large number of faithful disciples in various dispensations. Alma said that “there were many, exceedingly great many,” who were pure before the Lord.

This does not mean that these people never made mistakes or never had need of correction. The process of perfection includes challenges to overcome and steps to repentance that may be very painful. There is a proper place for chastisement in the molding of character, for we know that “whom the Lord loveth he chasteneth.”

Mortal perfection can be achieved as we try to perform every duty, keep every law, and strive to be as perfect in our sphere as our Heavenly Father is in His. If we do the best we can, the Lord will bless us according to our deeds and the desires of our hearts.

\intalkheader{Eternal Perfection}

But Jesus asked for more than mortal perfection. The moment He uttered the words “even as your Father which is in heaven is perfect,” He raised our sights beyond the bounds of mortality. Our Heavenly Father has eternal perfection. This very fact merits a much broader perspective.

Recently I studied the English and Greek editions of the New Testament, concentrating on each use of the term perfect and its derivatives. Studying both languages together provided some interesting insights, since Greek was the original language of the New Testament.

In Matthew 5:48, the term perfect was translated from the Greek teleios, which means “complete.” Teleios is an adjective derived from the noun telos, which means “end.” The infinitive form of the verb is teleiono, which means “to reach a distant end, to be fully developed, to consummate, or to finish.” Please note that the word does not imply freedom from error; it implies achieving a distant objective. In fact, when writers of the Greek New Testament wished to describe perfection of behavior—precision or excellence of human effort—they did not employ a form of teleios; instead, they chose different words.

Teleios is not a total stranger to us. From it comes the prefix tele- that we use every day. Telephone literally means “distant talk.” Television means “to see distantly.” Telephoto means “distant light,” and so on.

With that background in mind, let us consider another highly significant statement made by the Lord. Just prior to His crucifixion, He said that on “the third day I shall be perfected.” Think of that! The sinless, errorless Lord—already perfect by our mortal standards—proclaimed His own state of perfection yet to be in the future. His eternal perfection would follow His resurrection and receipt of “all power … in heaven and in earth.”

The perfection that the Savior envisions for us is much more than errorless performance. It is the eternal expectation as expressed by the Lord in His great intercessory prayer to His Father—that we might be made perfect and be able to dwell with them in the eternities ahead.

The Lord’s entire work and glory pertains to the immortality and eternal life of each human being. He came into the world to do the will of His Father, who sent Him. His sacred responsibility was foreseen before the Creation and was foretold by all His holy prophets since the world began.

The atonement of Christ fulfilled the long-awaited purpose for which He had come to the earth. His concluding words upon Calvary’s cross referred to the culmination of His assignment—to atone for all humankind. Then He said, “It is finished.” Not surprisingly, the Greek word from which finished was derived is teleios.

That Jesus attained eternal perfection following His resurrection is confirmed in the Book of Mormon. It records the visit of the resurrected Lord to the people of ancient America. There He repeated the important injunction previously cited but with one very significant addition. He said, “I would that ye should be perfect even as I, or your Father who is in heaven is perfect.” This time He listed Himself along with His Father as a perfected personage. Previously He had not.

Resurrection is requisite for eternal perfection. Thanks to the atonement of Jesus Christ, our bodies, corruptible in mortality, will become incorruptible. Our physical frames, now subject to disease, death, and decay, will acquire immortal glory. Presently sustained by the blood of life and ever aging, our bodies will be sustained by spirit and become changeless and beyond the bounds of death.

Eternal perfection is reserved for those who overcome all things and inherit the fulness of the Father in His heavenly mansions. Perfection consists in gaining eternal life—the kind of life that God lives.

\intalkheader{Ordinances and Covenants of the Temple}

Scriptures identify other important prerequisites to eternal perfection. They relate to the ordinances and covenants of the temple. No accountable individual can receive exaltation in the celestial kingdom without the ordinances of the temple. Endowments and sealings are for our personal perfection and are secured through our faithfulness.

This requirement also pertains to our ancestors. Paul taught “that they without us should not be made perfect.” Again, in that verse, the Greek term from which perfect was translated was a form of teleios.

In latter-day revelation, the Lord was even more explicit. His prophet wrote: “My dearly beloved brethren and sisters, let me assure you that these are principles in relation to the dead and the living that cannot be lightly passed over, as pertaining to our salvation. For their salvation is necessary and essential to our salvation. … They without us cannot be made perfect—neither can we without our dead be made perfect.”

\intalkheader{Encouragement from the Savior’s Example}

Our climb up the path to perfection is aided by encouragement from the scriptures. They hold the promise that we shall, if faithful in all things, become like Deity. John the beloved Apostle wrote:

“We should be called the sons [and daughters] of God. …

“… When he shall appear, we shall be like him; for we shall see him as he is.

“And every man that hath this hope in him purifieth himself, even as he is pure.”

Continuing encouragement comes as we follow the example of Jesus, who taught, “Be ye holy; for I am holy.” His hope for us is crystal clear! He declared: “What manner of men ought ye to be? Verily I say unto you, even as I am.” Thus, our adoration of Jesus is best expressed by our emulation of Jesus.

People have never failed to follow Jesus because His standards were imprecise or insufficiently high. Quite to the contrary. Some have disregarded His teachings because they were viewed as being too precise or impractically high! Yet such lofty standards, when earnestly pursued, produce great inner peace and incomparable joy.

There is no other individual to compare with Jesus Christ, nor is there any other exhortation equal to His sublime expression of hope: “I would that ye should be perfect even as I, or your Father who is in heaven is perfect.”

This divine entreaty is consistent with the fact that as begotten children of heavenly parents, we are endowed with the potential to become like them, just as mortal children may become like their mortal parents.

The Lord restored His church to help us prepare for perfection. Paul said that the Savior placed in the Church apostles, prophets, and teachers “for the perfecting of the saints, … for the edifying of the body of Christ:

“Till we all come in the unity of the faith, and of the knowledge of the Son of God, unto a perfect man, unto the measure of the stature of the fulness of Christ.”

The perfect man described in Paul’s quotation is the completed person—teleios—the glorified soul!

Moroni taught how to gain this glorious objective. His instruction stands in any age as an antidote for depression and a prescription for joy. I echo his plea: “Come unto Christ, and be perfected in him, and deny yourselves of all ungodliness; … love God with all your might, mind and strength. … [Then] ye may be perfect in Christ, … holy, [and] without spot.”

Meanwhile, brothers and sisters, let us do the best we can and try to improve each day. When our imperfections appear, we can keep trying to correct them. We can be more forgiving of flaws in ourselves and among those we love. We can be comforted and forbearing. The Lord taught, “Ye are not able to abide the presence of God now … ; wherefore, continue in patience until ye are perfected.”

We need not be dismayed if our earnest efforts toward perfection now seem so arduous and endless. Perfection is pending. It can come in full only after the Resurrection and only through the Lord. It awaits all who love Him and keep His commandments. It includes thrones, kingdoms, principalities, powers, and dominions. It is the end for which we are to endure. It is the eternal perfection that God has in store for each of us. I so testify in the name of Jesus Christ, amen.

\newpage
\settalk{“Thou Shalt Have No Other Gods”}
\setspeaker{Russell M. Nelson}
\settalkdate{April 1996}
\talkheading{“Thou Shalt Have No Other Gods”}{Russell M. Nelson}

As I face this congregation, I sense the responsibility entrusted to me. As you face the General Authorities and officers of the Church, you also bear responsibility. We are all accountable for acting upon the truths that are taught.

Facing you, I am reminded of military days long ago when our platoon heard shouts from a sergeant: “Attention!” “Right face!” “Left face!” “About face!” We learned to respond to those orders with instant precision. In retrospect, I don’t recall ever having heard his command to “face upward.” Yet scriptures tell us to “look to God and live.”

My topic today relates to the first of the Lord’s Ten Commandments: “Thou shalt have no other gods before me.” This commandment may be better known than obeyed. May I share a suggestion that I have found useful in testing my own allegiance to this commandment? When confronted with a challenging choice, I ask myself, “Which way do you face?”

\intalkheader{Life without Looking to God}

Sadly, many individuals don’t know where to find God, and exclude Him from their lives. When spiritual needs arise, they may look to the left, the right, or round about. But looking to other people on the same level cannot satisfy spiritual shortages. When the immortal spirit is starved, hunger persists for something more filling. Even when material success comes, there is a hollow ache if living well falls short of living worthily. Inner peace cannot be found in affluence accompanied by spiritual privation.

\intalkheader{Invitation to Come to the Lord}

Members of The Church of Jesus Christ of Latter-day Saints invite all to come unto Christ and enjoy the spiritual feast that His gospel provides. The Saints savor a sweet spiritual nourishment that sustains them through life. This sustenance comes because they have made covenants to take upon themselves the name of the Lord and strive to obey his precepts. Strength comes in recognizing and in being grateful for the Lord’s gifts of immortality and the opportunity for eternal life.

\intalkheader{Loyal Citizens}

These gifts are available to all. Citizens of many countries claim membership in the Church. Regardless of their flag or form of government, they find that allegiance to the Lord does not preclude their being loyal citizens of their nations. Fidelity to God enables one to develop a more profound patriotic allegiance and become a better citizen.

In addition to their national citizenship, members of the Church are also citizens of God’s kingdom. Their commitment to it, however, may be variable. The great majority seek “first to build up the kingdom of God, and to establish his righteousness.” Some allow their allegiance to God and His kingdom to slip below that of other interests in life. They have not yet determined which way they face.

\intalkheader{Representatives of the Lord}

I perceived such confusion in the mind of a newspaper reporter who asked one of our leaders when a representative of such-and-such a country would become a General Authority. While that question was being answered, I thought about our beloved General Authorities born in the countries of Asia; of Europe; of North, Central, and South America; and of the islands of the sea. Though these Brethren come from many nations and speak several tongues, not one of them was called to represent his native country. Presiding quorums of the Church are not representative assemblies. Each leader has been called to face the people as a representative of the Lord, not the other way around.

General Authorities are “called of God, by prophecy, and by the laying on of hands by those who are in authority.” They are called as “especial witnesses” unto all the world, to teach and testify of the Lord Jesus Christ.

\intalkheader{Obeying Commandments of the Lord}

No matter where we live or in what position we serve, all of us need to determine which way we face. God’s commandments serve as a standard against which priorities can be measured. Our respect for the first commandment fashions our feelings for all the others. Consider the commandment to keep the Sabbath day holy, for example. We live in a time when many people throughout the world have transferred their allegiance on the Sabbath from places of worship to places of amusement. Again I ask, “Which way do you face?”

Scriptures give us encouragement to do right: “If thou turn away … from doing thy pleasure on my holy day; and call the sabbath a delight, the holy of the Lord, … and shalt honour him, not doing thine own ways, nor finding thine own pleasure, …

“Then shalt thou delight thyself in the Lord.”

Self-esteem is also earned by obedience to God’s commandments regarding chastity. Yet in our day those commands have been attacked and trivialized. The morality of self-discipline with appropriate “denial or restraint has been popularly depicted as unhealthy and de-humanizing.” The truth is, “it is dehumanizing to define ourselves by our desires alone.” Each human being is a child of God—created in His image—with natural appetites to control.

If we break God’s first commandment, we cannot escape retribution. If we allow any other person or cause to come before allegiance to Him, we will reap a bitter harvest. Paul foresaw “destruction” for those “whose God is their belly.” (I might include all forms of anatomical affection.) Any who choose to serve “the creature more than the Creator” deprive themselves of spiritual reward.

Thus, our priorities should be honestly evaluated in terms of that first commandment. If change in direction is needed, we may want to issue a self-command to “about face!” Doing so would please the Lord, who said, “Repent, and turn yourselves from your idols; and turn away your faces from all your abominations.”

Trees reach up for the light and grow in the process. So do we as sons and daughters of heavenly parents. Facing upward provides a loftier perspective than facing right or facing left. Looking up in search of holiness builds strength and dignity as disciples of Deity.

\intalkheader{Facing Our Families}

Facing upward is crucial for successful parenting. Families deserve guidance from heaven. Parents cannot counsel children adequately from personal experience, fear, or sympathy. But when parents face children as would the Creator who gave them life, parents will be endowed with wisdom beyond that of their own. Wise mothers and fathers will teach members of their family how to make personal decisions based upon divine law. They will teach them that “this life is the time … to prepare to meet God.” They will teach them that decisions of a moral and spiritual character cannot be based on freedom to choose without accountability to God for those choices. With that understanding, parents and children will be rewarded with strength of character, peace of mind, joy, and rejoicing in their posterity.

\intalkheader{Facing Our Neighbors}

Similarly, relationships with neighbors, friends, and associates will be enhanced as we approach them with “the pure love of Christ.” A desire to emulate the Lord provides powerful motivation for good. Our craving for compassion will cause us to act in accord with the Golden Rule. By so doing, we will find joy in feeding the poor, clothing the naked, or doing volunteer work of worth.

Service to neighbors takes on new stature when we first look to God. In the Church, when priesthood and auxiliary leaders face their congregations, quorums, and classes as would the Lord, they learn that it does not matter where they serve, but how. Position in the Church does not exalt anyone, but faithfulness does. On the other hand, aspiring to a visible position—striving to become a master rather than a servant—can destroy the spirit of the worker and the work.

Occasionally confusion exists regarding servants and masters. The Bible reports that a group of men “had disputed among themselves, who should be the greatest” among them. Jesus said, “If any man desire to be first, the same shall be last of all, and servant of all.”

Was Jesus asking His disciples to respond to random requests from the crowd or to serve tables? No! He was asking them to serve in His way. The people were not to be masters of his disciples. The Lord is their Master.

In rendering service to others, which way do we face? From the right or the left, we can only push or pull. We can lift only from a higher plane. To reach it we don’t look sideways; we look up to our Master. Just as we must look to God to live well, so we must look to God to serve well.

\intalkheader{Attitudes of Effective Disciples}

If we are called to positions of leadership, we are accountable to the Savior for the acts we perform in that office. Those actions are shaped by attitudes, and attitudes are elevated while lowering our heads in humble prayer. So state words in the hymn “Before Thee, Lord, I Bow My Head”:

\begin{verse}
Look up, my soul; be not cast down.\\
Keep not thine eyes upon the ground.\\
Break off the shackles of the earth.\\
Receive, my soul, the spirit’s birth.\\
And now as I go forth again\\
To mingle with my fellowmen,\\
Stay thou nearby, my steps to guide,\\
That I may in thy love abide.
\end{verse}

Praying helps us to face trials in life. Prayer centers our attitudes precisely. With that focus, we do not wander to the right or left through land mined with traps of temptation. Disciples do not flirt with danger at the jagged edge of disaster. Experienced mountain climbers do not lean toward the dangerous edge but toward safety, with ropes and other safeguards to secure them to those they trust. So it is with us. When we climb mountainous challenges of life, we should lean toward our Master and be yoked with him, clinging tightly to the iron rod of the gospel, to family, and to trusted friends.

President David O. McKay taught this about edges: “Many of us through selfishness are lingering near the edge of the animal jungle where Nature’s law demands us to do everything with self in view.”

The Lord said, “Look unto me in every thought; doubt not, fear not.” I have learned that such faith gives emancipating power. Facing God first lets us decide firmly what we shall not do; then we are free to pursue what we ought to do.

Recently President Gordon B. Hinckley declared: “Love of God is the root of all virtue, of all goodness, of all strength of character, of all fidelity to do right. Love the Lord your God, and love his Son, and be ever grateful for their love for us. Whenever other love fades, there will be that shining, transcendent, everlasting love of God for each of us and the love of his Son, who gave his life for us.”

Brothers and sisters, race, nationality, occupation, or other interests need not stand in the way. All can look to the Lord. All can place him first in their lives. Those who do so and remain faithful will qualify for his sublime promise:” “Every soul who forsaketh his sins and cometh unto me, and calleth on my name, and obeyeth my voice, and keepeth my commandments, shall see my face and know that I am.” This glorious destiny can be ours, I so testify in the name of Jesus Christ, amen.

\newpage
\settalk{The Atonement}
\setspeaker{Russell M. Nelson}
\settalkdate{October 1996}
\talkheading{The Atonement}{Russell M. Nelson}

Humbly I join the Book of Mormon prophet Jacob, who asked, “Why not speak of the atonement of Christ?” This topic comprises our third article of faith: “We believe that through the Atonement of Christ, all mankind may be saved, by obedience to the laws and ordinances of the Gospel.”

Before we can comprehend the Atonement of Christ, however, we must first understand the Fall of Adam. And before we can understand the Fall of Adam, we must first understand the Creation. These three crucial components of the plan of salvation relate to each other.

\intalkheader{The Creation}

The Creation culminated with Adam and Eve in the Garden of Eden. They were created in the image of God, with bodies of flesh and bone. Created in the image of God and not yet mortal, they could not grow old and die. “And they would have had no children” nor experienced the trials of life. (Please forgive me for mentioning children and the trials of life in the same breath.) The creation of Adam and Eve was a paradisiacal creation, one that required a significant change before they could fulfill the commandment to have children and thus provide earthly bodies for premortal spirit sons and daughters of God.

\intalkheader{The Fall}

That brings us to the Fall. Scripture teaches that “Adam fell that men might be; and men are, that they might have joy.” The Fall of Adam (and Eve) constituted the mortal creation and brought about the required changes in their bodies, including the circulation of blood and other modifications as well. They were now able to have children. They and their posterity also became subject to injury, disease, and death. And a loving Creator blessed them with healing power by which the life and function of precious physical bodies could be preserved. For example, bones, if broken, could become solid again. Lacerations of the flesh could heal themselves. And miraculously, leaks in the circulation could be sealed off by components activated from the very blood being lost.

Think of the wonder of that power to heal! If you could create anything that could repair itself, you would have created life in perpetuity. For example, if you could create a chair that could fix its own broken leg, there would be no limit to the life of that chair. Many of you walk on legs that were once broken and do so because of your remarkable gift of healing.

Even though our Creator endowed us with this incredible power, He consigned a counterbalancing gift to our bodies. It is the blessing of aging, with visible reminders that we are mortal beings destined one day to leave this “frail existence.” Our bodies change every day. As we grow older, our broad chests and narrow waists have a tendency to trade places. We get wrinkles, lose color in our hair—even the hair itself—to remind us that we are mortal children of God, with a “manufacturer’s guarantee” that we shall not be stranded upon the earth forever. Were it not for the Fall, our physicians, beauticians, and morticians would all be unemployed.

Adam and Eve, as mortal beings, were instructed to “worship the Lord their God, and … offer the firstlings of their flocks, for an offering unto the Lord.” They were further instructed that “the life of the flesh is in the blood: … for it is the blood that maketh an atonement for the soul.” Probation, procreation, and aging were all components of—and physical death was essential to—God’s “great plan of happiness.”

But mortal life, glorious as it is, was never the ultimate objective of God’s plan. Life and death here on planet Earth were merely means to an end—not the end for which we were sent.

\intalkheader{The Atonement}

That brings us to the Atonement. Paul said, “As in Adam all die, even so in Christ shall all be made alive.” The Atonement of Jesus Christ became the immortal creation. He volunteered to answer the ends of a law previously transgressed. And by the shedding of His blood, His and our physical bodies could become perfected. They could again function without blood, just as Adam’s and Eve’s did in their paradisiacal form. Paul taught that “flesh and blood cannot inherit the kingdom of God; … this mortal must put on immortality.”

\intalkheader{Meaning of Atonement}

With this background in mind, let us now ponder the deep meaning of the word atonement. In the English language, the components are at-one-ment, suggesting that a person is at one with another. Other languages employ words that connote either expiation or reconciliation. Expiation means “to atone for.” Reconciliation comes from Latin roots re, meaning “again”; con, meaning “with”; and sella, meaning “seat.” Reconciliation, therefore, literally means “to sit again with.”

Rich meaning is found in study of the word atonement in the Semitic languages of Old Testament times. In Hebrew, the basic word for atonement is kaphar, a verb that means “to cover” or “to forgive.” Closely related is the Aramaic and Arabic word kafat, meaning “a close embrace”—no doubt related to the Egyptian ritual embrace. References to that embrace are evident in the Book of Mormon. One states that “the Lord hath redeemed my soul … ; I have beheld his glory, and I am encircled about eternally in the arms of his love.” Another proffers the glorious hope of our being “clasped in the arms of Jesus.”

I weep for joy when I contemplate the significance of it all. To be redeemed is to be atoned—received in the close embrace of God with an expression not only of His forgiveness but of our oneness of heart and mind. What a privilege! And what a comfort to those of us with loved ones who have already passed from our family circle through the gateway we call death!

Scriptures teach us more about the word atonement. The Old Testament has many references to atonement, which called for animal sacrifice. Not any animal would do. Special considerations included:

The Atonement of Christ fulfilled these prototypes of the Old Testament. He was the firstborn Lamb of God, without blemish. His sacrifice occurred by the shedding of blood. No bones of His body were broken—noteworthy in that both malefactors crucified with the Lord had their legs broken. And His was a vicarious sacrifice for others.

While the words atone or atonement, in any of their forms, appear only once in the King James translation of the New Testament, they appear 35 times in the Book of Mormon. As another testament of Jesus Christ, it sheds precious light on His Atonement, as do the Doctrine and Covenants and the Pearl of Great Price. Latter-day revelation has added much to our biblical base of understanding.

\intalkheader{Infinite Atonement}

In preparatory times of the Old Testament, the practice of atonement was finite—meaning it had an end. It was a symbolic forecast of the definitive Atonement of Jesus the Christ. His Atonement is infinite—without an end. It was also infinite in that all humankind would be saved from never-ending death. It was infinite in terms of His immense suffering. It was infinite in time, putting an end to the preceding prototype of animal sacrifice. It was infinite in scope—it was to be done once for all. And the mercy of the Atonement extends not only to an infinite number of people, but also to an infinite number of worlds created by Him. It was infinite beyond any human scale of measurement or mortal comprehension.

Jesus was the only one who could offer such an infinite atonement, since He was born of a mortal mother and an immortal Father. Because of that unique birthright, Jesus was an infinite Being.

\intalkheader{The Ordeal of the Atonement}

The ordeal of the Atonement centered about the city of Jerusalem. There the greatest single act of love of all recorded history took place. Leaving the upper room, Jesus and His friends crossed the deep ravine east of the city and came to a garden of olive trees on the lower slopes of the Mount of Olives. There in the garden bearing the Hebrew name of Gethsemane—meaning “oil press”—olives had been beaten and pressed to provide oil and food. There at Gethsemane, the Lord “suffered the pain of all men, that all … might repent and come unto him.” He took upon Himself the weight of the sins of all mankind, bearing its massive load that caused Him to bleed from every pore.

Later He was beaten and scourged. A crown of sharp thorns was thrust upon His head as an additional form of torture. He was mocked and jeered. He suffered every indignity at the hands of His own people. “I came unto my own,” He said, “and my own received me not.” Instead of their warm embrace, He received their cruel rejection. Then He was required to carry His own cross to the hill of Calvary, where He was nailed to that cross and made to suffer excruciating pain.

Later He said, “I thirst.” To a doctor of medicine, this is a very meaningful expression. Doctors know that when a patient goes into shock because of blood loss, invariably that patient—if still conscious—with parched and shriveled lips cries for water.

Even though the Father and the Son knew well in advance what was to be experienced, the actuality of it brought indescribable agony. “And [Jesus] said, Abba, Father, all things are possible unto thee; take away this cup from me: nevertheless not what I will, but what thou wilt.” Jesus then complied with the will of His Father. Three days later, precisely as prophesied, He rose from the grave. He became the firstfruits of the Resurrection. He had accomplished the Atonement, which could give immortality and eternal life to all obedient human beings. All that the Fall allowed to go awry, the Atonement allowed to go aright.

The Savior’s gift of immortality comes to all who have ever lived. But His gift of eternal life requires repentance and obedience to specific ordinances and covenants. Essential ordinances of the gospel symbolize the Atonement. Baptism by immersion is symbolic of the death, burial, and Resurrection of the Redeemer. Partaking of the sacrament renews baptismal covenants and also renews our memory of the Savior’s broken flesh and of the blood He shed for us. Ordinances of the temple symbolize our reconciliation with the Lord and seal families together forever. Obedience to the sacred covenants made in temples qualifies us for eternal life—the greatest gift of God to man—the “object and end of our existence.”

\intalkheader{The Atonement Enabled the Purpose of the Creation to Be Accomplished}

The Creation required the Fall. The Fall required the Atonement. The Atonement enabled the purpose of the Creation to be accomplished. Eternal life, made possible by the Atonement, is the supreme purpose of the Creation. To phrase that statement in its negative form, if families were not sealed in holy temples, the whole earth would be utterly wasted.

The purposes of the Creation, the Fall, and the Atonement all converge on the sacred work done in temples of The Church of Jesus Christ of Latter-day Saints. The earth was created and the Church was restored to make possible the sealing of wife to husband, children to parents, families to progenitors, worlds without end.

This is the great latter-day work of which we are a part. That is why we have missionaries; that is why we have temples—to bring the fullest blessings of the Atonement to faithful children of God. That is why we respond to our own calls from the Lord. When we comprehend His voluntary Atonement, any sense of sacrifice on our part becomes completely overshadowed by a profound sense of gratitude for the privilege of serving Him.

As one of the “special witnesses of the name of Christ in all the world,” I testify that He is the Son of the living God. Jesus is the Christ—our atoning Savior and Redeemer. This is His Church, restored to bless God’s children and to prepare the world for the Second Coming of the Lord. I so testify in the sacred name of Jesus Christ, amen.

\newpage
\settalk{Endure and Be Lifted Up}
\setspeaker{Russell M. Nelson}
\settalkdate{April 1997}
\talkheading{Endure and Be Lifted Up}{Russell M. Nelson}

Early in our married life when Sister Nelson and I lived in Minneapolis, we decided to enjoy a free afternoon with our two-year-old daughter. We went to one of Minnesota’s many beautiful lakes and rented a small boat. After rowing far from shore, we stopped to relax and enjoy the tranquil scene. Suddenly, our little toddler lifted one leg out of the boat and started to go overboard, exclaiming, “Time to get out, Daddy!”

Quickly we caught her and explained, “No, dear, it’s not time to get out; we must stay in the boat until it brings us safely back to land.” Only with considerable persuasion did we succeed in convincing her that leaving the boat early would have led to disaster.

Children are prone to do such dangerous things simply because they have not acquired the wisdom their parents have. Similarly, we as children of our Heavenly Father may foolishly want to get “out of the boat” before we arrive at destinations He would like us to reach. The Lord teaches over and over that we are to endure to the end. This is a dominant theme of the scriptures. One example may serve to represent many passages that convey a similar message:

“Blessed are they who shall seek to bring forth my Zion … , for they shall have the gift and the power of the Holy Ghost; and if they endure unto the end they shall be lifted up at the last day, and shall be saved in the everlasting kingdom of the Lamb.”

Blessings bestowed by God are always predicated upon obedience to law. Applied to my analogy, we are first to get “on the boat” with Him. Then we are to stay with Him. And if we don’t get “out of the boat” before we should, we shall reach His kingdom, where we will be lifted up to eternal life.

The term lifted up relates to a physical law that can be illustrated by a simple demonstration. I will use a spool of thread and blow into the axial hole of the spool. The force of my breath will move a piece of tissue paper away from me. Next I will take an ordinary card and a straight pin. I will place the pin through the card. With the pin in the hole of the spool, I will hold the card close to the spool. I will again blow into the hole of the spool. As I blow, I will let go of the card so that it can respond to physical forces. Before I proceed, would you like to predict what will happen? Will I blow the card away from me, or will the card be lifted up toward me? Are you ready? [Elder Nelson demonstrates that blowing down the axial hole of the spool lifts the card up toward the spool.]

Did you notice? As long as I had sufficient breath, the card was lifted up. But when I could endure no longer, the card fell. When my breath gave out, the opposing force of gravity prevailed. If my energy could have endured, the card would have been lifted up indefinitely.

Energy is always required to provide lift over opposing forces. These same laws apply in our personal lives. Whenever an undertaking is begun, both the energy and the will to endure are essential. The winner of a five-kilometer race is declared at the end of five kilometers, not at one or two. If you board a bus to Boston, you don’t get off at Burlington. If you want to gain an education, you don’t drop out along the way—just as you don’t pay to dine at an elegant restaurant only to walk away after sampling the salad.

Whatever your work may be, endure at the beginning, endure through opposing forces along the way, and endure to the end. Any job must be completed before you can enjoy the result for which you are working. So wrote the poet:

\begin{verse}
Stick to your task till it sticks to you;\\
Beginners are many, but enders are few.\\
Honor, power, place, and praise\\
Will [always] come … to the one who stays.
\end{verse}

Sometimes the need to endure comes when facing a physical challenge. Anyone afflicted with a serious illness or with the infirmities of age hopes to be able to endure to the end of such trials. Most often, intense physical challenges are accompanied by spiritual challenges as well.

Think of the early pioneers. What if they had not endured the hardships of their westward migration? There would be no sesquicentennial celebration this year. Steadfastly they endured—through persecution, expulsion, a governmental order of extermination, expropriation of property, and much more. Their enduring faith in the Lord provided lift for them as it will for you and for me.

The Lord’s ultimate concern is for the salvation and exaltation of each individual soul. What if the Apostle Paul’s conversion had not been enduring? He never would have testified as he did at the end of his ministry: “I have fought a good fight, I have finished my course, I have kept the faith.”

What if Jesus had wavered in His commitment to do His Father’s will? His Atonement would not have been accomplished. The dead would not be resurrected. The blessings of immortality and eternal life would not be. But Jesus did endure. During His final hour, Jesus prayed to His Father, saying, “I have glorified thee on the earth: I have finished the work which thou gavest me to do.”

Early in His mortal ministry, Jesus became concerned about the commitment of His followers. He had just fed the 5,000, then had taught them the doctrines of the kingdom. But some had murmured, “This is an hard saying; who can hear it?” Even after He had fed them, many lacked the faith to endure with Him. He turned to the Twelve and said, “Will ye also go away?

“Then Simon Peter answered him, Lord, … thou hast the words of eternal life.

“And we believe and are sure that thou art that Christ, the Son of the living God.”

Peter’s answer defines the real core of commitment. When we know without a doubt that Jesus is the Christ, we will want to stay with Him. When we are surely converted, the power to endure is ours.

This power to endure is critical in those two most important relationships we enter into in life. One is marriage; the other is membership in the Lord’s Church. These are also unique in that they are both covenant—not contractual—relationships.

Marriage—especially temple marriage—and family ties involve covenant relationships. They cannot be regarded casually. With divorce rates escalating throughout the world today, it is apparent that many spouses are failing to endure to the end of their commitments to each other. And some temple marriages fail because a husband forgets that his highest and most important priesthood duty is to honor and sustain his wife. The best thing that a father can do for his children is to “love their mother.”

President Gordon B. Hinckley made a statement recently that each Latter-day Saint husband should heed: “Magnify your [wife],” he said, “and in so doing you will magnify your priesthood.” To his profound advice we might couple the timeless counsel of Paul, who said, “Let every one of you … love his wife even as himself; and the wife see that she reverence her husband.” Enduring love provides enduring lift through life’s trials. An enduring marriage results when both husband and wife regard their union as one of the two most important commitments they will ever make.

The other commitment of everlasting consequence is to the Lord. Unfortunately, some souls make a covenant with God—signified by the sacred ordinance of baptism—without a heartfelt commitment to endure with Him. Baptism is an extremely important ordinance. But it is only initiatory. The supreme benefits of membership in the Church can be realized only through the exalting ordinances of the temple. These blessings qualify us for “thrones, kingdoms, principalities, and powers” in the celestial kingdom.

The Lord can readily discern between those with superficial signs of activity and those who are deeply rooted in His Church. This Jesus taught in the parable of the sower. He observed that some “have no root in themselves, and so endure but for a time: afterward, when affliction or persecution ariseth for the word’s sake, immediately they are offended.”

Loyalty to the Lord carries an obligation of loyalty to those called by the Lord to lead His Church. He has empowered that men be ordained to speak in His holy name. As they guide His unsinkable boat safely toward the shore of salvation, we would do well to stay on board with them. “No waters can swallow the ship where lies / The Master of ocean and earth and skies.”

Nevertheless, some individuals want to jump “out of the boat” before reaching land. And others, sadly, are persuaded out by companions who insist that they know more about life’s perilous journey than do prophets of the Lord. Problems often arise that are not of your own making. Some of you may innocently find yourselves abandoned by one you trusted. But you will never be forsaken by your Redeemer, who said, “I, the Lord, am bound when ye do what I say.”

Without a strong commitment to the Lord, an individual is more prone to have a low level of commitment to a spouse. Weak commitments to eternal covenants lead to losses of eternal consequence. Laments later in life are laced with remorse, as expressed in these lines:

\begin{verse}
For of all sad words of tongue or pen,\\
The saddest are these: “It might have been!”
\end{verse}

We are speaking of the most important of all blessings. The Lord said, “If you keep my commandments and endure to the end you shall have eternal life, which gift is the greatest of all the gifts of God.”

Each of you who really wants to endure to the glorious end that our Heavenly Father has foreseen should firmly establish some personal priorities. With many interests competing for your loyalty, you need to be careful first to stay safely “on the boat.” No one can serve two masters. If Satan can get you to love anything—fun, flirtation, fame, or fortune—more than a spouse or the Lord with whom you have made sacred covenants to endure, the adversary begins to triumph. When faced with such temptations, you will find that strength comes from commitments made well in advance. The Lord said, “Settle this in your hearts, that ye will do the things which I shall teach, and command you.” He declared through His prophet Jeremiah, “I will put my law in their inward parts, and write it in their hearts; and will be their God, and they shall be my people.”

When priorities are proper, the power to endure is increased. And when internalized, those priorities will help keep you from “going overboard.” They will protect you from cheating—in marriage, in the Church, and in life.

If you really want to be like the Lord—more than anything or anyone else—you will remember that your adoration of Jesus is best shown by your emulation of Him. Then you will not allow any other love to become more important than love for your companion, your family, and your Creator. You will govern yourself not by someone else’s set of rules but by revealed principles of truth.

Your responsibility to endure is uniquely yours. But you are never alone. I testify that the lifting power of the Lord can be yours if you will “come unto Christ” and “be perfected in him.” You will “deny yourselves of all ungodliness.” And you will “love God with all your might, mind and strength.”

The living prophet of the Lord has issued a clarion call: “I invite every one of you,” said President Hinckley, “to stand on your feet and with a song in your heart move forward, living the gospel, loving the Lord, and building the kingdom. Together we shall stay the course and keep the faith.”

I pray that each of us may so endure and be lifted up at the last day, in the name of Jesus Christ, amen.

\newpage
\settalk{Spiritual Capacity}
\setspeaker{Russell M. Nelson}
\settalkdate{October 1997}
\talkheading{Spiritual Capacity}{Russell M. Nelson}

Those invited by the First Presidency to speak at general conference are not assigned subjects. Speakers pray for inspiration and prepare according to promptings they receive. I have been impressed to speak on “spiritual capacity.”

A verse of scripture opens a door of opportunity for each of us: “There is a spirit in man,” said Job, “and the inspiration of the Almighty giveth them understanding.” To take advantage of such an opportunity, we need more than a verbal incentive. We need an example—someone to show us how spiritual capacity can be developed. I have selected as a model for my message President Gordon B. Hinckley. I hope he will pardon me. My motive in doing so is not one of adulation, but of emulation. We can draw upon his example in order to improve our own spiritual attributes.

This year, Sister Nelson and I have had the privilege of accompanying President and Sister Hinckley to 11 countries for which I have had some responsibilities. That has given us a rare opportunity to observe him closely under a variety of conditions. His teachings are always inspiring and relevant. They should be studied carefully and applied individually. They represent the word of the Lord for His people.

But my purpose is not to review the content of President Hinckley’s messages. Instead, I would like to focus upon his spiritual capacities. He has developed many, including “faith, virtue, knowledge, temperance, patience, brotherly kindness, godliness, charity, humility, [and] diligence.”

His humility, for example, is so sincere that he would have me point only to the Lord Jesus Christ as our great exemplar. Of course, He is! The Master said, “I have given you an example, that ye should do as I have done.” We must never lose sight of the Savior’s enduring standard as the ultimate for each of us.

But we can also learn much from a man who has spent his entire lifetime in striving to be more like the Lord. More than 87 years ago, Gordon B. Hinckley was an infant cradled in the arms of his loving parents. That newborn babe looked much as any other, I presume. An infant’s body is tiny, and its spiritual capacities are undeveloped. While the body may reach the peak of its maturation in a few years, the development of the spirit may never reach the limit of its capacity, because there is no end to progression.

President Hinckley’s personality, manner, and native intelligence have always been uniquely his. To these inborn attributes, however, he has added spiritual capacities, and they are continuing to increase.

Both his parents and he understood the importance of education and a mission. After his graduation from the university, he faced a major decision in 1933, when he was called to go on a mission. At that time, most young men in the Church were not able to serve because of a global economic depression that deprived nearly everyone of available cash. Earlier, his wonderful mother, with foresight and faith, had established a small savings account for his mission. Though she died before his call, her fund sent him on his way.

Shortly after Elder Hinckley’s labors began in England, he became discouraged and wrote to his father. After reading that letter, his father’s wise reply closed with these words: “Forget yourself and go to work.” Thanks to noble parents and a crucial decision to remain, Elder Hinckley completed his mission with honor. Now he often states that the good things that have happened to him since have all hinged upon that decision to stay. On his mission, he developed good habits of study, work, communication, budgeting, time management, and more. There he learned that nothing is too hard for the Lord.

Long ago, President Hinckley harnessed the power of prayer. I have watched him pray over many weighty matters and receive inspired answers. Prayer invites those ennobling attributes of the spirit that are ultimately “bestowed upon all who are true followers of … Jesus Christ.”

Hobbies can aid in spiritual development. Worthy music, dance, art, and writing are among the creative activities that can enrich the soul. A good hobby can dispel heartache and give zest to life. Through the years, one of President Hinckley’s hobbies has been his home. As a young father, he learned how to build. He acquired the skills necessary to remodel a house and make needed repairs. And more important, he has built and maintained the trust of his wife and their children. Together they have established—and are still adding to—wonderful memories with their children and grandchildren, who know that they are part of “a chosen generation … called … out of darkness into [the] marvellous light” of the Lord. From the Hinckleys’ parental example, we can learn a great lesson. Love at home comes when companions cultivate their commitment to keep the commandments of God.

President Hinckley’s love of learning is catalyzed by curiosity. He grasps every opportunity to learn from others. On one occasion, I heard him quiz a local security officer for nearly an hour regarding crime control in a major city. I have heard him converse with building contractors, reporters, and those who specialize in the arts, architecture, business, government, law, medicine, and other disciplines. He knows their vocabularies, their challenges, and their strengths.

His remarkable ability as a writer has been gained by his living close to the Spirit. Similar skills can come to others too, for scriptures state that such has been “given unto as many as called upon God to write by the spirit of inspiration.”

Through the years, President Hinckley has developed a remarkable sense of humor. You have heard his quip that “Sister Hinckley and I are learning that the so-called golden years are laced with lead.” I might add a pun, President. We are grateful to be led by that kind of lead. It gives ballast to balance one who might lean too far in any direction. And it gives stability to character.

While I focus upon President Hinckley, Sister Hinckley should also be included. They have been married for 60 years and have long been one in spirit, while maintaining their individuality. They do not waste time pondering the past or fretting about the future. And they persevere in spite of adversity.

While going from a chapel to an airport in Central America, their vehicle was involved in an accident. Sister Nelson and I were traveling behind them and saw it occur. A truck loaded on top with unsecured metal rods approached them at an intersection. To avoid a collision, its driver suddenly stopped the truck, launching those iron rods like javelins to pierce the Hinckleys’ car. Windows were smashed; fenders and doors were dented. The accident could have been very serious. While shattered glass was being removed from their clothing and skin, President Hinckley said, “Thank the Lord for His blessing; now let’s continue on in another car.”

Among President Hinckley’s spiritual attributes is that of compassion. He is sympathetic to people and feels a strong urge to help them. I have watched him weep with those who mourn and rejoice when Saints are blessed. Such compassion can come to anyone whose heart has truly been touched by the Spirit of the Lord.

President and Sister Hinckley have demonstrated that the capacity to understand increases as one learns and then teaches with diligence. Unless illness interferes, age does not diminish—it augments—the capacity for spiritual development.

Each President of the Church, armed with the Holy Ghost as a constant companion, inherits an enormous workload at an age when most men would be retired. President Hinckley sets a pace that is unprecedented. In 1996, he visited missionaries, members, and friends of the Church in 23 nations on four continents. During that year, he gave more than 200 major discourses. His stride in 1997 continues to follow that same pattern. His strenuous schedule is driven by his determination to be “anxiously engaged” in building the kingdom of God. Often I have heard him say, “I don’t know how to get anything done except getting on my knees and pleading for help and then getting on my feet and going to work.” Unshakable faith, hard work, and contagious optimism epitomize our prophet.

I have watched President Hinckley, in speaking before great congregations, depend upon the Holy Ghost, who serves “to enlighten and ennoble the mind, to purify and sanctify the soul, to incite to good works, and to reveal the things of God.”

President Hinckley has achieved spiritual supremacy over physical feelings. Even when entitled to normal complaints of “jet lag” or “burnout,” he is attentive. I believe that his personal antidote for fatigue is enthusiasm for the work. He is energized by the Lord, who said, “I will impart unto you of my Spirit, which shall enlighten your mind … [and] fill your soul with joy.”

One of our most memorable experiences occurred when we visited the temple construction site in Guayaquil, Ecuador. There President Hinckley recounted to us how that property was selected. On a prior visit, he had been shown several possible locations, but none seemed to satisfy him. While prayerfully searching, he asked about ground on a hill not far from the airport. But it was said to be not for sale. President Hinckley directed that they visit that property anyway. There he received inspiration from the Almighty that this was the right place for the temple. Now we were privileged to stand on that spot reserved by the Lord and then procured for this sacred purpose. Our joy was indescribable.

The prophet makes major decisions on a daily basis. This he does with great capacity. Meanwhile, he encourages each of us to make choices that will “give us growth and joy in this life and eternal life in the world to come.”

This President of the Church calls many people to serve, knowing that much is required of them. He is keenly aware of their opportunities and risks. “Yes, this work requires sacrifice,” he said. “It requires effort, it means courage to speak out and faith to try. … It needs men and women of solemn purpose.” “We know that there are some limits on what you can do, but we know also that there need be no limits on enthusiasm, planning, thoughtful consideration, and effort.”

Brothers and sisters, the spirit that dwells within each of us can be enriched with enthusiasm and enlightened by the Almighty. The process of spiritual growth is revealed in the scriptures: “Intelligence cleaveth unto intelligence; wisdom receiveth wisdom; truth embraceth truth; … [and] light cleaveth unto light.” “That which is of God is light; and he that receiveth light, and continueth in God, receiveth more light; and that light groweth brighter and brighter until the perfect day.”

Gratefully, we follow prophets who have been given a divine commission: “Whatsoever they shall speak when moved upon by the Holy Ghost shall be scripture, shall be the will of the Lord, shall be the mind of the Lord, shall be the word of the Lord, shall be the voice of the Lord, and the power of God unto salvation.”

While we follow prophetic teachings, we can develop our spiritual capacities by emulating one such as President Gordon B. Hinckley. I thank God for this prophet. He is the Lord’s anointed. Willingly I follow him. I love him and sustain him. I so testify in the name of Jesus Christ, amen.

\newpage
\settalk{A New Harvest Time}
\setspeaker{Russell M. Nelson}
\settalkdate{April 1998}
\talkheading{A New Harvest Time}{Russell M. Nelson}

Family love is wonderful. Nothing is as specific as the love of a baby for its mother. Nothing is as predictable as the love of children for their parents or the love of parents for their children.

Recently I was tenderly hugging one of our precious little five-year-old granddaughters and said to her, “I love you, sweetheart.”

She responded rather blandly, “I know.”

I asked, “How do you know that I love you?”

“Because! You’re my grandfather!”

That was reason enough for her. Indeed, we do love our grandchildren. We also love our grandparents. I cherish the memories of life with three of my four grandparents. I never met my Grandfather Nelson. He died when my father was only 16 years old. At the time of Grandfather’s passing, he was superintendent of public instruction for the state of Utah. He owned a handsome pocket watch, which my father later gave to me. Now that watch is a tangible link between us.

I think of my Grandfather Nelson with deep feelings of gratitude. I received much of my early education in schools he helped to develop. And I cherish my membership in this Church, to which both of his parents were converted in Denmark about a century and a half ago. In fact, all eight of my great-grandparents were converts to the Church in Europe. Of the others, one joined the Church in Sweden, two in England, and three in Norway. How grateful I am to these pioneer predecessors! My debt to them is reflected in these biblical verses: “One soweth, and another reapeth” that “both he that soweth and he that reapeth may rejoice together.”

Today we are reaping a harvest of family love from seeds sown years ago. Preparations to strengthen family ties came in 1823, when the angel Moroni first appeared to the Prophet Joseph Smith. Moroni announced the coming of Elijah, who would cause the hearts of children to be turned to their fathers.

Elijah’s return to earth occurred at the first temple built in this dispensation, where he and other heavenly messengers, under direction of the Lord, entrusted special keys of priesthood authority to the restored Church:

With that, natural affection between generations began to be enriched. This restoration was accompanied by what is sometimes called the Spirit of Elijah—a manifestation of the Holy Ghost bearing witness of the divine nature of the family. Hence, people throughout the world, regardless of religious affiliation, are gathering records of deceased relatives at an ever-increasing rate.

Elijah came not only to stimulate research for ancestors. He also enabled families to be eternally linked beyond the bounds of mortality. Indeed, the opportunity for families to be sealed forever is the real reason for our research. The Lord declared through the Prophet Joseph Smith: “These are principles in relation to the dead and the living that cannot be lightly passed over, as pertaining to our salvation. For their salvation is necessary and essential to our salvation, … they without us cannot be made perfect—neither can we without our dead be made perfect.”

Among the first in this dispensation to sow seeds of interest in family history were the brothers Orson and Parley P. Pratt, members of the Quorum of the Twelve Apostles. Their efforts resulted in a Pratt family genealogy and the performance of temple ordinances for about 3,000 of their ancestors.

Yet there were many Church members who did not fully understand the responsibility for their own kindred. President Wilford Woodruff was so concerned that he made the issue a matter of fervent prayer. Then, at April 1894 general conference, he presented a revelation to the membership of the Church. From it I quote: “We want the Latter-day Saints from this time to trace their genealogies as far as they can, and to be sealed to their fathers and mothers. Have children sealed to their parents, and run this chain through as far as you can get it. … This is the will of the Lord to his people.”

Later that year, the First Presidency and the Twelve established the Genealogical Society of Utah. From modest beginnings in an upstairs room of the Church historian’s office, its collection and facilities have grown. Today the Family History Library™ occupies a modern five-story building with access to 280,000 books, 700,000 microfiches, and more than 2 million rolls of microfilm, making it the largest library of its kind in the world.

In 1964 the department began to establish branch libraries. Today more than 3,000 Family History Centers™ dot the globe.

Technology used to support this important work has changed greatly over the years. In 1927 a card file was instituted to index all endowments performed. The index was maintained through 1969, when new endowments were recorded in the first major computer system, identified by the acronym GIANT. It was used for more than two decades.

The Society’s extensive microfilming has permitted the gathering of records at their sources, with copies made available later at the Family History Library and Family History Centers. Microfilming has been done in 110 countries, accumulating more than 2 billion exposures with approximately 13 billion names. Microfilming has enabled the Family History Library to expand its collections dramatically and provide resources for an explosive growth of genealogical research worldwide. These microfilms comprise the core of information contained in our present automated systems.

By the 1980s, the personal computer had revolutionized the management of information. The Family History Department employed this technology in developing Personal Ancestral File® to help members organize data regarding their ancestors. In 1990, FamilySearch® was announced. At October conference that year, Elder Richard G. Scott described components of FamilySearch: Ancestral File™, Family History Library Catalog™, International Genealogical Index®, and more. His message stimulated Sister Nelson and me to use these tools to organize information that we and our relatives had gathered over many years.

Meanwhile, objectives of decentralization and simplification led to record extraction programs, in which thousands of Church members have participated. Extraction projects have now produced records for more than 300 million individuals.

Many people have joined with members of the Church in efforts to index the burgeoning bank of genealogical information. An example is the 1881 British census. For this project, more than 8,000 volunteers from family history societies throughout the British Isles have transcribed 30 million names. Gratefully, we announce that fruits of this labor are now on fiche and will soon be available on compact disc from the Church’s distribution centers.

We are also pleased to announce that data from the 1880 census of the United States will soon be released on compact disc. Meanwhile, volunteers are working on other projects, such as arrival records for immigrants to the USA through Ellis Island.

May I express our deep appreciation to all valiant volunteers—past, present, and future—for their diligent work on these and other projects.

In describing these achievements, I realize that for some who are less involved in this work, I may have intensified feelings of guilt. I apologize for that. I know that fear and unfamiliarity may stand in your way. For others, even the mention of a computer may be an additional intimidator. Some secretly hope that they can slip through their remaining days on earth without ever having to touch a computer. To those with access to computers, I say: “Reach out! Have hope! Try! I have exciting news for you!”

“The time of harvest is come.” A new era of family history work has arrived. As President Gordon B. Hinckley recently noted, “The Lord has inspired skilled men and women in developing new technologies which we can use to our great advantage in moving forward this sacred work.” Previously, efforts have focused on gathering names and dates and organizing that information. Now, computer products are available that can actually guide you to find your kindred.

May I introduce you to the new Family History SourceGuide™. This compact disc is now available at the Church’s distribution centers. It can lead you to genealogical records in countries, states, and provinces around the world and shows how you can use these records to identify your ancestors. It includes other aids, such as maps, letter-writing guides, translations of words for several non-English-speaking countries, definitions, and terms often found in genealogical records. Family History SourceGuide puts at your fingertips much of the collected knowledge and experience of hundreds of genealogical experts. It can all be yours—at the touch of a button. Use it, and rejoice!

A new Vital Records Index™ will make available on compact disc the results of extraction programs prepared from many civil and ecclesiastical records. Some overlap will exist between this resource and records in the International Genealogical Index, but most of the names in the Vital Records Index have not yet had temple ordinance work performed. The entire index will include approximately 25 million records. During the next few months, it will be released in segments by geographic area, such as the British Isles (5 million records) and North America (4.5 million records). This file represents years of work of many extraction workers.

I am excited about these and other developments. Tasks that once seemed beyond reach are now within our grasp. “With God nothing shall be impossible.” A new harvest time has come. The way is opening by which we can obey His will and provide welding links between all dispensations and generations.

To get started, you do not need equipment. Begin with a pedigree chart and a family group record. List the names of those you know. Add information learned from living relatives. This simple start at home will prepare you to receive additional help. And when you are baptized for a deceased ancestor, you will sense a feeling of validation of this divine work that will bring great joy.

As we ponder the importance of our ancestral responsibilities, we also need to be reminded of the Lord’s vast ministry. I quote from President Joseph F. Smith: “Jesus had not finished his work when his body was slain, neither did he finish it after his resurrection from the dead; although he had accomplished the purpose for which he then came to the earth, he had not fulfilled all his work. And when will he? Not until he has redeemed and saved every son and daughter of our father Adam that have been or ever will be born upon this earth to the end of time. … That is his mission. We will not finish our work until we have saved ourselves, and then not until we shall have saved all depending upon us; for we are to become saviors upon Mount Zion, as well as Christ. We are called to this mission. The dead are not perfect without us, neither are we without them.”

To this end, the will of the Lord has been impressed upon President Hinckley to build more temples. The Latter-day Saints are to be an endowed people, and they are to be sealed to their posterity and progenitors.

My grandfather’s watch reminds me that our grandparents watch—and wait—for us to identify them, be linked to them, and provide temple ordinances for them. May God bless us all with success in this sacred service, I pray in the name of Jesus Christ, amen.

\newpage
\settalk{We Are Children of God}
\setspeaker{Russell M. Nelson}
\settalkdate{October 1998}
\talkheading{We Are Children of God}{Russell M. Nelson}

Recently I observed young men who attracted attention to themselves because of their extreme styles of dress and grooming. One made a revealing remark when he said, “I’m trying to find out who I really am.” This occurred after I had come from a Church meeting where Primary children had sung “I Am a Child of God.” Such contrasting experiences emphasize the importance of knowing that we are literally children of God.

We are dual beings. Each soul is comprised of body and spirit, both of which emanate from God. A firm understanding of body and spirit will shape our thoughts and deeds for good.

\intalkheader{The Body}

The marvel of our physical bodies is often overlooked. Who has not encountered feelings of low self-esteem because of physique or appearance? Many people wish their bodies could be more to their liking. Some with naturally straight hair want it curly. Others with curly hair want it straight. Occasionally some ladies, believing that “gentlemen prefer blonds,” become “decided blonds.”

Your body, whatever its natural gifts, is a magnificent creation of God. It is a tabernacle of flesh—a temple for your spirit. A study of your body attests to its divine design.

Its formation begins with the union of two reproductive cells—one from the mother and one from the father. Together, these two cells contain all of the new individual’s hereditary information, stored in a space so small it cannot be seen by the naked eye. Twenty-three chromosomes from each parent unite in one new cell. These chromosomes contain thousands of genes which determine all of the physical characteristics of the unborn person. Approximately 22 days after these two cells unite, a little heart begins to beat. At 26 days, blood begins to circulate. Cells multiply and divide. Some become eyes that see; others become ears that hear.

Each organ is a wondrous gift from God. The eye has a self-focusing lens. Nerves and muscles control two separate eyes to make a single three-dimensional image. The eyes are connected to the brain, which records the sights seen. No cords or batteries are needed.

Each ear is connected to compact equipment designed to convert sound waves into audible tones. An eardrum serves as a diaphragm. Minute ossicles amplify sound vibrations and transmit a signal via nerves to the brain, which senses and remembers the sounds.

The heart is an incredible pump. It has four delicate valves that control the direction of blood flow. These valves open and close more than 100,000 times a day—36 million times a year. Yet, unless altered by disease, they are able to withstand this stress almost indefinitely. No man-made material developed to date can be flexed so frequently and so long without breaking.

Each day an adult heart pumps enough fluid to fill a 2,000-gallon tank. This work is equivalent to lifting a grown man to the top of the Empire State Building while expending only about four watts of energy. At the crest of the heart is an electrical generator that transmits energy down special lines, causing myriads of muscle fibers to work together.

Much could be said about each of the other precious organs in the body. They function in a marvelous manner, beyond my time or ability to describe.

Other attributes of the body are equally amazing, though less evident. For example, backup is provided. Each paired organ has instant backup available from the other of the pair. Single organs, such as the brain, the heart, and the liver, are nourished by two routes of blood supply. This design protects the organ if harm should come to any one channel.

Think of the body’s system of self-defense. To protect it from harm, the body perceives pain. In response to infection, it generates antibodies. They not only help to combat the immediate problem, but they persist to strengthen resistance to infection in the future. One day my attention was directed to some three-year-old children who had lapped up water from a street gutter. The number of germs they ingested must have been incalculable, but not one of those youngsters became ill. As soon as that dirty drink reached each little stomach, its hydrochloric acid went to work to treat the water and protect the life of the child.

The skin provides protection. It also warns against injuries that excessive heat or cold might cause. It even sends signals that indicate trouble elsewhere. With fever, the skin perspires. When one is frightened, it pales. When one is embarrassed, it blushes.

The body repairs itself. Broken bones mend and become strong once again. Skin lacerations heal themselves. A leak in the circulation can seal itself.

The body renews its own outdated cells. The average red blood cell, for instance, lives about 120 days. Then it is replaced by a newly regenerated cell.

The body regulates its own vital ingredients. Essential elements and chemical constituents are adjusted continuously. And regardless of wide fluctuations in temperature of the environment, the temperature of the body is carefully controlled within narrow bounds.

If these qualities of normal function, defense, repair, regeneration, and regulation were to prevail in perpetuity, life here would continue without limit. Mercifully, our Creator provided for aging and other processes which ultimately result in physical death. We often think of death as untimely or tragic. But death, like birth, is part of life. Scripture tells us that “it was not expedient that man should be reclaimed from this temporal death, for that would destroy the great plan of happiness.” To return to God through the gateway of death is a joy for those who love Him.

When death claims an individual in the prime of life, we take comfort in knowing that the very laws which do not allow life to persist here are the same laws that will be implemented at the time of the Resurrection, when the body will be endowed with immortality.

\intalkheader{The Spirit}

Next I speak of the spirit. Prior to our mortal existence here, each spirit son and daughter lived with God. The spirit is eternal; it existed in innocence in the premortal realm and will exist after the body dies. The spirit provides the body with animation and personality. “All spirit is matter, but it is more fine or pure.”

“The spirit of man [is] in the likeness of his person.” Jesus so explained as the brother of Jared viewed the Lord’s premortal body:

“Seest thou that ye are created after mine own image? Yea, even all men were created in the beginning after mine own image.

“… This body, which ye now behold, is the body of my spirit; and man have I created after the body of my spirit; and even as I appear unto thee to be in the spirit will I appear unto my people in the flesh.”

Development of the spirit is of eternal consequence. The attributes by which we shall be judged one day are those of the spirit. These include the virtues of integrity, compassion, love, and more. Your spirit, by being housed in your body, is able to develop and express these attributes in ways that are vital to your eternal progression.

Spirit and body, when joined together, become a living soul of supernal worth. Indeed, we are children of God—physically and spiritually.

\intalkheader{Physical Limitations}

For reasons usually unknown, some people are born with physical limitations. Specific parts of the body may be abnormal. Regulatory systems may be out of balance. And all of our bodies are subject to disease and death. Nevertheless, the gift of a physical body is priceless. Without it, we cannot attain a fulness of joy.

A perfect body is not required to achieve a divine destiny. In fact, some of the sweetest spirits are housed in frail frames. Great spiritual strength is often developed by those with physical challenges precisely because they are challenged. Such individuals are entitled to all the blessings that God has in store for His faithful and obedient children.

Eventually the time will come when each “spirit and … body shall be reunited again in … perfect form; both limb and joint shall be restored to its proper frame.” Then, thanks to the Atonement of Jesus Christ, we can become perfected in Him.

\intalkheader{Personal Behavior}

How should these truths influence our personal behavior? We should gratefully acknowledge God as our Creator. Otherwise, we would be as guilty as goldfish swimming in a bowl, oblivious to the goodness of their provider. “Ye must give thanks unto God,” said the Lord, “for whatsoever blessing ye are blessed with.” And we can practice virtue and holiness before Him continually.

We will regard our body as a temple of our very own. We will not let it be desecrated or defaced in any way. We will control our diet and exercise for physical fitness.

Should not equal attention be paid to spiritual fitness? Just as physical strength requires exercise, so spiritual strength requires effort. Among the most important of spiritual exercises is prayer. It engenders harmony with God and a desire to keep His commandments. Prayer is a key to wisdom, virtue, and humility.

We will be careful about which counsel we heed. Many so-called experts give advice for the body—without thought for the spirit. Anyone who accepts direction contrary to the Word of Wisdom, for example, forsakes a law revealed to bring both physical and spiritual blessings. Some recommendations regarding use of our reproductive organs are based solely—and inadequately—upon physical considerations. Beware of such one-sided views! Paul taught that “if ye live after the flesh, ye shall die: but if ye through the Spirit do mortify the deeds of the body, ye shall live.”

That caution pertains to pornography, which is highly addictive. Scriptural warning is clear: “Let not sin therefore reign in your mortal body, that ye should obey it in the lusts thereof.” In time, addictions enslave both the body and the spirit. Full repentance from addiction is best accomplished in this life, while we still have a mortal body to help us.

As children of God, we should not let anything enter the body that might defile it. To allow sensors of sight, touch, or hearing to supply the brain with unclean memories is a sacrilege. We will cherish our chastity and avoid “foolish and hurtful lusts, which drown [us] in destruction and perdition.” We will “flee these things; and follow after righteousness, godliness, faith, love, patience, [and] meekness”—traits that edify the whole soul.

Who are we? We are children of God. Our potential is unlimited. Our inheritance is sacred. May we always honor that heritage—in every thought and deed—I pray in the name of Jesus Christ, amen.

\newpage
\settalk{Our Sacred Duty to Honor Women}
\setspeaker{Russell M. Nelson}
\settalkdate{April 1999}
\talkheading{Our Sacred Duty to Honor Women}{Russell M. Nelson}

It is a joy to be with you tonight, brethren, and wonderful to see so many young men with their fathers. We are assembled because of our desire to hearken to the leaders of the Church. But this congregation is unique. I don’t see any mothers. Not one of us could be here without a mother, yet we are all here—without our mothers.

Tonight I am attending with a son, sons-in-law, and grandsons. Where are their mothers? Gathered in the kitchen of our home! What are they doing? Making large batches of homemade doughnuts! And when we return home, we will feast on those doughnuts. While we enjoy them, these mothers, sisters, and daughters will listen intently as each of us speaks of things he learned here tonight. It’s a nice family tradition, symbolic of the fact that everything we learn and do as priesthood bearers should bless our families.

Let us speak about our worthy and wonderful sisters, particularly our mothers, and consider our sacred duty to honor them.

When I was a young university student, one of my classmates urgently pleaded with a group of us—his Latter-day Saint friends—to donate blood for his mother, who was bleeding profusely. We went directly to the hospital to have our blood typed and tested. I’ll never forget our shock when told that one of the prospective donors was unfit because of a positive blood test for a venereal disease. That infected blood was his own! Fortunately, his mother survived, but I’ll never forget his lingering sorrow. He bore the burden of knowing that his personal immorality had disqualified him from giving needed aid to his mother, and he had added to her grief. I learned a great lesson: if one dishonors the commandments of God, one dishonors mother, and if one dishonors mother, one dishonors the commandments of God.

\intalkheader{Honor Motherhood}

During my professional career as a doctor of medicine, I was occasionally asked why I chose to do that difficult work. I responded with my opinion that the highest and noblest work in this life is that of a mother. Since that option was not available to me, I thought that caring for the sick might come close. I tried to care for my patients as compassionately and competently as Mother cared for me.

Many years ago the First Presidency issued a statement that has had a profound and lasting influence upon me. “Motherhood,” they wrote, “is near to divinity. It is the highest, holiest service to be assumed by mankind. It places her who honors its holy calling and service next to the angels.”

Because mothers are essential to God’s great plan of happiness, their sacred work is opposed by Satan, who would destroy the family and demean the worth of women.

You young men need to know that you can hardly achieve your highest potential without the influence of good women, particularly your mother and, in a few years, a good wife. Learn now to show respect and gratitude. Remember that your mother is your mother. She should not need to issue orders. Her wish, her hope, her hint should provide direction that you would honor. Thank her and express your love for her. And if she is struggling to rear you without your father, you have a double duty to honor her.

The influence of your mother will bless you throughout life, especially when you serve as a missionary. Long years ago, Elder Frank Croft was serving in the state of Alabama. While preaching to the people, he was forcefully abducted by a vicious gang, to be whipped and lashed across his bare back. Elder Croft was ordered to remove his coat and shirt before he was tied to a tree. As he did so, a letter he had recently received from his mother fell to the ground. The vile leader of the gang picked up the letter. Elder Croft closed his eyes and uttered a silent prayer. The attacker read the letter from Elder Croft’s mother. From a copy of that letter, I quote:

“My beloved son, … remember the words of the Savior when He said, … ‘Blessed are ye when men shall revile you and persecute you and say all manner of evil against you falsely for my name’s sake. Rejoice and be exceedingly glad for you will have your reward in Heaven for so persecuted they the prophets which were before you.’ Also remember the Savior upon the cross suffering from the sins of the world when He had uttered these immortal words, ‘Father forgive them, for they know not what they do.’ Surely, my boy, they who are mistreating you … know not what they do or they would not do it. Sometime, somewhere, they will understand and then they will regret their action and they will honor you for the glorious work you are doing. So be patient, my son, love those who mistreat you and say all manner of evil against you and the Lord will bless you and magnify you. … Remember also, my son, that day and night, your mother is praying for you.”

Elder Croft watched the hateful man as he studied the letter. He would read a line or two, then sit and ponder. He arose to approach his captive. The man said: “Feller, you must have a wonderful mother. You see, I once had one, too.” Then addressing the mob, he said: “Men, after reading this Mormon’s mother’s letter, I just can’t go ahead with the job. Maybe we had better let him go.” Elder Croft was released without harm.

We are deeply grateful for the faithful mothers and fathers of our wonderful missionaries. The love they bear for their children is sublime.

\intalkheader{Honor the Sisters}

We who bear the holy priesthood have a sacred duty to honor our sisters. We are old enough and wise enough to know that teasing is wrong. We respect sisters—not only in our immediate families but all the wonderful sisters in our lives. As daughters of God, their potential is divine. Without them, eternal life would be impossible. Our high regard for them should spring from our love of God and from an awareness of their lofty purpose in His great eternal plan.

Hence, I warn against pornography. It is degrading of women. It is evil. It is infectious, destructive, and addictive. The body has means by which it can cleanse itself from harmful effects of contaminated food or drink. But it cannot vomit back the poison of pornography. Once recorded, it always remains subject to recall, flashing its perverted images across your mind, with power to draw you away from the wholesome things in life. Avoid it like the plague!

\intalkheader{Honor Your Wife}

To you who are not yet married, think about your future marriage. Choose your companion well. Remember the scriptures that teach the importance of marriage in the temple:

“In the celestial glory there are three heavens or degrees;

“And in order to obtain the highest, a man must enter into this order of the priesthood [meaning the new and everlasting covenant of marriage];

“And if he does not, he cannot obtain it.”

The highest ordinances in the house of the Lord are received by husband and wife together and equally—or not at all!

In retrospect I see that the most important day of my life was the day when my sweetheart, Dantzel, and I were married in the holy temple. Without her I could not have the highest and most enduring blessings of the priesthood. Without her I would not be the father to our wonderful children or grandfather to our precious grandchildren.

As fathers we should have love unbounded for the mothers of our children. We should accord to them the gratitude, respect, and praise that they deserve. Husbands, to keep alive the spirit of romance in your marriage, be considerate and kind in the tender intimacies of your married life. Let your thoughts and actions inspire confidence and trust. Let your words be wholesome and your time together be uplifting. Let nothing in life take priority over your wife—neither work, recreation, nor hobby.

An ideal marriage is a true partnership between two imperfect people, each striving to complement the other, to keep the commandments, and to do the will of the Lord.

\intalkheader{Fathers Preside over the Family in Love}

The family is the most important unit of society and of the Church. The family is ordained of God. It is central to His plan for the eternal destiny of His children. “God has established families to bring happiness to his children, to allow them to learn correct principles in a loving atmosphere, and to prepare them for eternal life.”

Parents have the primary responsibility for the welfare of their children. The Church does not replace that parental responsibility. Ideally, the Latter-day Saint family is presided over by a worthy man who holds the priesthood. This patriarchal authority has been honored among the people of God in all dispensations. It is of divine origin, and that union, if sealed by proper authority, will continue throughout eternity. He who is the Father of us all and the source of this authority demands that governance in the home be in love and righteousness.

You fathers can help with the dishes, care for a crying baby, and change a diaper. And perhaps some Sunday you could get the children ready for Church, and your wife could sit in the car and honk.

“Husbands, love your wives, even as Christ also loved the church, and gave himself for it.” With that kind of love, brethren, we will be better husbands and fathers, more loving and spiritual leaders. Happiness at home is most likely to be achieved when practices there are founded upon the teachings of Jesus Christ. Ours is the responsibility to ensure that we have family prayer, scripture study, and family home evening. Ours is the responsibility to prepare our children to receive the ordinances of salvation and exaltation and the blessings promised to tithe payers. Ours is the privilege to bestow priesthood blessings of healing, comfort, and direction.

The home is the great laboratory of love. There the raw chemicals of selfishness and greed are melded in the crucible of cooperation to yield compassionate concern and love one for another.

Honor the special sisters in your lives, brethren. Express your love to your wife, to your mother, and to the sisters. Praise them for their forbearance with you even when you are not at your best. Thank the Lord for these sisters who—like our Heavenly Father—love us not only for what we are but for what we may become. Humbly I thank God for my mother, my sisters, my daughters, granddaughters, and for my special sweetheart, companion, and friend—my wife!

May God bless us to honor each virtuous woman, I pray in the name of Jesus Christ, amen.

\newpage
\settalk{A Testimony of the Book of Mormon}
\setspeaker{Russell M. Nelson}
\settalkdate{October 1999}
\talkheading{A Testimony of the Book of Mormon}{Russell M. Nelson}

Not long after my call to serve as one of the Twelve Apostles, I was summoned to the office of the President of our Quorum, President Ezra Taft Benson. He expressed deep concern that members of the Church did not fully appreciate the value of the Book of Mormon. With emotion in his voice, he read to me from the 84th section of the Doctrine and Covenants:

“Your minds in times past have been darkened because of unbelief, and because you have treated lightly the things you have received—

“Which vanity and unbelief have brought the whole church under condemnation.”

By that time, President Benson had completely captured my attention. He then concluded his admonition:

“And they shall remain under this condemnation until they repent and remember the new covenant, even the Book of Mormon.”

I shall never forget that lesson. Since then, President Howard W. Hunter, President Gordon B. Hinckley, and many other leaders of the Church have continued to extol the Book of Mormon to people throughout the world.

I would like to add my testimony of the divinity of this book. I have read it many times. I have also read much that has been written about it. Some authors have focused upon its stories, its people, or its vignettes of history. Others have been intrigued by its language structure or its records of weapons, geography, animal life, techniques of building, or systems of weights and measures.

Interesting as these matters may be, study of the Book of Mormon is most rewarding when one focuses on its primary purpose—to testify of Jesus Christ. By comparison, all other issues are incidental.

When you read the Book of Mormon, concentrate on the principal figure in the book—from its first chapter to the last—the Lord Jesus Christ, Son of the Living God. And look for a second undergirding theme: God will keep His covenants with the remnants of the house of Israel.

The Book of Mormon is a crucial component of that covenant. It is holy scripture that encompasses sacred writings from the small and large plates of Nephi, the plates of Mormon, the plates of Ether, and the plates of brass, which contained “the five books of Moses, … a record of the Jews … ; and … prophecies of the holy prophets.”

When Mormon abridged these records, he noted that he could not write a “hundredth part” of their proceedings. Thus, historical aspects of the book assume secondary significance.

The Holy Bible has 66 individual books; the Book of Mormon contains 15. Its first book of Nephi—written some six centuries before the birth of Jesus—records that the prophet Lehi received a vision of the tree of life. His son Nephi prayed to know its meaning. In answer, he was given a remarkable vision. He beheld a virgin bearing a Child in her arms. He envisioned the Redeemer of the world, His earthly ministry, and His Crucifixion. He saw 12 others who would follow the Holy One. And he foresaw the ongoing opposition to the work of God and of His Apostles.

Other great prophets of the Book of Mormon—in their own way and time—testified of the divinity of the Lord Jesus Christ. Among them were the brother of Jared, Zenock, Neum, and Zenos. Testimonies of Jesus Christ that predated His birth in Bethlehem were also recorded from King Benjamin, Abinadi, Alma the Elder, Alma the Younger, Amulek, the sons of Mosiah, Captain Moroni, the brothers Nephi and Lehi, and Samuel the Lamanite. In a seemingly endless sequence of prophetic proclamations—testimonies of “all the holy prophets” for “a great many thousand years before his coming”—the Book of Mormon makes the solemn declaration that Jesus is the Christ, our Savior and Redeemer.

\intalkheader{Authors}

Most books contained in libraries of the world were authored for contemporary readers. And they were generally written for profit, with royalties accruing from successful sales.

Not so with the Book of Mormon. It was written anciently for our day. It reveals the endless Lordship of Jesus Christ in accounts of two ancient American dispensations, preserved for the benefit of us who live in this dispensation of the fulness of times. Certainly no royalties came to its authors. In fact, they paid dearly for their privilege of participation. What motivated them? Their devotion to God! The book’s four major writers—Nephi, Jacob, Mormon, and Moroni—were all eyewitnesses of the Lord, as was its martyred translator, the Prophet Joseph Smith.

\intalkheader{Content}

Their writings centered upon the Lord, His mission, and His ministry. Jacob, for example, repeatedly referred to the Atonement and Resurrection of Christ. “Beloved brethren,” wrote Jacob, “be reconciled unto [God] through the atonement of Christ, his Only Begotten Son, and ye may obtain a resurrection … and be presented as the first-fruits of Christ unto God. …

“And now, … why not speak of the atonement of Christ, and attain to a perfect knowledge of him,” and a “knowledge of a resurrection and the world to come?”

Jacob’s advice is priceless and timeless.

The Savior declared that the Book of Mormon contains “the fulness of [His] everlasting gospel.” How did He define the gospel? The resurrected Lord taught, “This is the gospel which I have given unto you—that I came into the world to do the will of my Father, because my Father sent me.”

Then He amplified that one-sentence definition: “My Father sent me that I might be lifted up upon the cross; and after that I had been lifted up upon the cross, that I might draw all men unto me.”

This unique mortal mission of the Lord—the gospel as He defined it—we know as the Atonement. The fulness of the gospel, therefore, connotes a fuller comprehension of the Atonement. This we do not obtain from the Bible alone. The word atonement, in any of its forms, is mentioned only once in the King James Version of the New Testament. In the Book of Mormon, it appears 39 times! The Book of Mormon also contains more references to the Resurrection than does the Bible.

The Savior referred to the Book of Mormon as His “new covenant” with the house of Israel. It is a tangible sign of Christ’s culminating covenant with mankind. The divine teachings of this book, as a third testament, clarify doctrine and unify the Old Testament with the New Testament. Scriptural covenants, testaments, and witnesses since the beginning of time relate to the Atonement of Jesus Christ, the central act of all human history.

The Book of Mormon is the most important religious text to be revealed from God to man “since the writings of the New Testament were compiled nearly two millennia ago.” Joseph Smith declared the Book of Mormon to be “the most correct of any book on earth, and the keystone of our religion.” It is the only book that the Lord Himself has testified to be true.

The crowning event of this sacred record is the personal ministry of the resurrected Lord to people of ancient America. To them He made this revelatory announcement:

“Behold, I am Jesus Christ the Son of God. I created the heavens and the earth, and all things that in them are. I was with the Father from the beginning. …

“… The scriptures concerning my coming are fulfilled. …

“I am the light and the life of the world. …

“… Whoso cometh unto me with a broken heart and a contrite spirit, him will I baptize with fire and with the Holy Ghost. …

“… I have come unto the world to bring redemption unto the world, to save the world from sin.

“Therefore, whoso repenteth and cometh unto me as a little child, him will I receive, for of such is the kingdom of God. … I have laid down my life, and have taken it up again; therefore repent, and come unto me … , and be saved.”

After that supernal introduction, the Master reinforced His identity by allowing the multitudes to thrust their hands into His side and feel the prints of the nails in His hands and in His feet. Then they knew that the very God of Israel stood in their presence—He who had been slain for the sins of the world.

He instructed the people. He taught them to pray, to repent, to be baptized, to partake of the sacrament, to know of His doctrine, to understand the importance of sacred ordinances and covenants, and to endure to the end.

The Book of Mormon is a gift from God to all humankind, and He has “commanded his people that they should persuade all men to repentance.” He invites all “to come unto him and partake of his goodness,” and He denies “none that come unto him, black and white, bond and free, male and female.”

\intalkheader{Translation}

This appeal to all people must involve many languages and the work of skilled translators. The King James Version of the Bible, for example, was produced by 50 English scholars who accomplished their work in seven years, translating at the rate of one page per day. Expert translators today do well if they can also translate scripture at the rate of one page per day.

In contrast, Joseph Smith translated the Book of Mormon at the rate of about 10 pages per day, completing the task in about 85 days! (Many of us feel good if we can read the book in that time.)

Such a pace is even more remarkable considering the circumstances under which the Prophet labored. In that same period, while enduring constant distractions and incessant hostility, Joseph Smith moved more than 100 miles from Harmony, Pennsylvania, to Fayette, New York. He applied for a copyright. He received revelations comprising 12 sections of the Doctrine and Covenants. Heavenly beings restored the holy priesthood. Yet he completed the translation in less than three months.

The First Presidency provided opportunity for the Twelve Apostles to see portions of the original manuscript and the printer’s manuscript of the Book of Mormon. Words cannot describe the deep emotions we felt as we examined these precious documents and observed that marks of editing were remarkably rare.

\intalkheader{Personal Testimony and Blessings}

Each individual who prayerfully studies the Book of Mormon can also receive a testimony of its divinity. In addition, this book can help with personal problems in a very real way. Do you want to get rid of a bad habit? Do you want to improve relationships in your family? Do you want to increase your spiritual capacity? Read the Book of Mormon! It will bring you closer to the Lord and His loving power. He who fed a multitude with five loaves and two fishes—He who helped the blind to see and the lame to walk—can also bless you! He has promised that those who live by the precepts of this book “shall receive a crown of eternal life.”

The Book of Mormon is true! I so testify in the name of Jesus Christ, amen.

\newpage
\settalk{The Creation}
\setspeaker{Russell M. Nelson}
\settalkdate{April 2000}
\talkheading{The Creation}{Russell M. Nelson}

We will long remember this inspiring conference in the new Conference Center. Not long ago, there was only a deep hole in the ground where this building now stands. We have watched its construction with interest and awe.

The process of construction is truly inspiring to me. From conception to completion, any major building project reflects upon the work of the Master Creator. In fact, the Creation—of planet Earth and of life upon it—undergirds all other creative capability. Any manmade creation is possible only because of our divine Creator. The people who design and build are given life and capacity by that Creator. And all materials used in the construction of an edifice are ultimately derived from the rich resources of the earth. The Lord declared, “The earth is full, and there is enough and to spare; yea, I prepared all things.”

It is difficult for mortal minds to comprehend the majesty of the Creation. It is much easier for us to think about good things to eat or fun things to do. But I would like to stretch our minds to think of things beyond our easy grasp. The creation of man and woman was wondrous and great. So was the creation of the earth as their mortal dwelling place.

The entire Creation was planned by God. A council in heaven was once convened in which we participated. There our Heavenly Father announced His divine plan. It is also called the plan of happiness, the plan of salvation, the plan of redemption, the plan of restoration, the plan of mercy, the plan of deliverance, and the everlasting gospel. The purpose of the plan is to provide opportunity for the spirit children of God to progress toward an eternal exaltation.

\intalkheader{Components of the Plan}

The plan required the Creation, and that in turn required both the Fall and the Atonement. These are the three fundamental components of the plan. The creation of a paradisiacal planet came from God. Mortality and death came into the world through the Fall of Adam. Immortality and the possibility of eternal life were provided by the Atonement of Jesus Christ. The Creation, the Fall, and the Atonement were planned long before the actual work of the Creation began.

While visiting the British Museum in London one day, I read a most unusual book. It is not scripture. It is an English translation of an ancient Egyptian manuscript. From it I quote a dialogue between the Father and the Son. Referring to His Father, Jehovah—the premortal Lord—says:

“He took the clay from the hand of the angel, and made Adam according to Our image and likeness, and He left him lying for forty days and forty nights without putting breath into him. And He heaved sighs over him daily, saying, ‘If I put breath into this [man], he must suffer many pains.’ And I said unto My Father, ‘Put breath into him; I will be an advocate for him.’ And My Father said unto Me, ‘If I put breath into him, My beloved Son, Thou wilt be obliged to go down into the world, and to suffer many pains for him before Thou shalt have redeemed him, and made him to come back to his primal state.’ And I said unto My Father, ‘Put breath into him; I will be his advocate, and I will go down into the world, and will fulfil Thy command.’”

Although this text is not scripture, it reaffirms scriptures that teach of the deep and compassionate love of the Father for the Son, and of the Son for us—attesting that Jesus volunteered willingly to be our Savior and Redeemer.

The Lord God declared, “This is my work and my glory—to bring to pass the immortality and eternal life of man.” He who, under direction of the Father, had created the earth, subsequently came into mortality to do the will of His Father and to fulfill all prophecies of the Atonement. His Atonement would redeem every soul from the penalties of personal transgression, on conditions that He set.

\intalkheader{Phases of the Creation}

Each phase of the Creation was well planned before it was accomplished. Scripture tells us that “the Lord God, created all things … spiritually, before they were naturally upon the face of the earth.”

The physical Creation itself was staged through ordered periods of time. In Genesis and Moses, those periods are called days. But in the book of Abraham, each period is referred to as a time. Whether termed a day, a time, or an age, each phase was a period between two identifiable events—a division of eternity.

Period one included the creation of atmospheric heavens and physical earth, culminating in the emergence of light from darkness.

In period two, the waters were divided between the surface of the earth and its atmospheric heavens. Provision was made for clouds and rain to give life to all that would later dwell upon the earth.

In period three, plant life began. The earth was organized to bring forth grass, herbs, trees, and vegetation—each growing from its own seed.

Period four was a time of further development. Lights in the expanse of the heaven were organized so there could be seasons and other means of measuring time. During this period, the sun, the moon, the stars, and the earth were placed in proper relationship to one another. The sun, with its vast stores of hydrogen, was to serve as a giant furnace to provide light and heat for the earth and life upon it.

In period five, fish, fowl, and “every living creature” were added. They were made fruitful and able to multiply—in the sea and on the earth—each after its own kind.

In the sixth period, creation of life continued. The beasts of the earth were made after their kind, cattle after their kind, and everything which “creepeth upon the earth”—again, after its own kind. Then the Gods counseled together and said: “Let us go down and form man in our image, after our likeness. …

“So the Gods went down to organize man in their own image, in the image of the Gods to form they him, male and female to form they them.” Thus, Adam and Eve were formed. And they were blessed to “be fruitful, and multiply, and replenish the earth, and subdue it: and have dominion over the fish of the sea, and over the fowl of the air, and over every living thing that moveth upon the earth.”

The seventh period was designated as a time of rest.

\intalkheader{The Creation Testifies of a Creator}

I testify that the earth and all life upon it are of divine origin. The Creation did not happen by chance. It did not come ex nihilo (out of nothing). And human minds and hands able to build buildings or create computers are not accidental. It is God who made us and not we ourselves. We are His people! The Creation itself testifies of a Creator. We cannot disregard the divine in the Creation. Without our grateful awareness of God’s hand in the Creation, we would be just as oblivious to our provider as are goldfish swimming in a bowl. With deep gratitude, we echo the words of the Psalmist, who said, “O Lord, how manifold are thy works! in wisdom hast thou made them all: the earth is full of thy riches.”

\intalkheader{Purpose and Destiny of the Earth}

This earth is but one of many creations over which God presides. “Worlds without number have I created,” He said. “And I also created them for mine own purpose; and by the Son I created them, which is mine Only Begotten.” Grand as it is, planet Earth is part of something even grander—that great plan of God. Simply summarized, the earth was created that families might be. Scripture explains that a husband and wife “shall be one flesh, and all this that the earth might answer the end of its creation.”

And as part of the planned destiny of the earth and its inhabitants, here our kindred dead are also to be redeemed. Families are to be sealed together for all eternity. A welding link is to be forged between the fathers and the children. In our time, a whole, complete, and perfect union of all dispensations, keys, and powers is to be welded together. For these sacred purposes, holy temples now dot the earth.

Though our understanding of the Creation is limited, we know enough to appreciate its supernal significance. And that store of knowledge will be augmented in the future. Scripture declares: “In that day when the Lord shall come [again], he shall reveal all things—

“Things which have passed, and hidden things which no man knew, things of the earth, by which it was made, and the purpose and the end thereof—

“Things most precious, things that are above, and things that are beneath, things that are in the earth, and upon the earth, and in heaven.”

Yes, further light and knowledge will come. The Lord said: “If there be bounds set to the heavens or to the seas, or to the dry land, or to the sun, moon, or stars—

“All the times of their revolutions, all the appointed days, months, and years, … and all their glories, laws, and set times, shall be revealed in the days of the dispensation of the fulness of times.”

Eventually, “the earth will be renewed and receive its paradisiacal glory.” At the Second Coming of the Lord, the earth will be changed once again. It will be returned to its paradisiacal state and be made new. There will be a new heaven and a new earth.

\intalkheader{Our Responsibilities}

Meanwhile, brothers and sisters, we should understand our significant responsibilities. Both the creations of God and the creations of man teach us the importance of each component. Do you think that the absence of one piece of granite from the face of this building would be noticed? Of course it would!

So it is with each son or daughter of God. We cannot let “the head say unto the feet it [has] no need of the feet; for without the feet how shall the body be able to stand?” Just as “the body [has] need of every member,” so the family has need of every member. All members of a family are to be linked, sealed, and “edified together, that the system may be kept perfect.”

The Creation, great as it is, is not an end in itself but a means to an end. We come to the earth for a brief period of time, endure our tests and trials, and prepare to move onward and upward to a glorious homecoming. Our thoughts and deeds while here will surely be more purposeful if we understand God’s plan and are thankful for and obedient to His commandments.

As beneficiaries of the divine Creation, what shall we do? We should care for the earth, be wise stewards over it, and preserve it for future generations. And we are to love and care for one another.

We are to be creators in our own right—builders of an individual faith in God, faith in the Lord Jesus Christ, and faith in His Church. We are to build families and be sealed in holy temples. We are to build the Church and kingdom of God upon the earth. We are to prepare for our own divine destiny—glory, immortality, and eternal lives. These supernal blessings can all be ours through our faithfulness.

I testify that God lives! Jesus is the Christ and Creator! He is Lord over all the earth. He has established His Church in these latter days to accomplish His divine purposes. Joseph Smith is the great prophet of the Restoration. President Gordon B. Hinckley is His prophet today, whom I sustain with all my heart, in the name of Jesus Christ, amen.

\newpage
\settalk{Living by Scriptural Guidance}
\setspeaker{Russell M. Nelson}
\settalkdate{October 2000}
\talkheading{Living by Scriptural Guidance}{Russell M. Nelson}

Recently Sister Nelson and I were in Denmark during the commemoration of the 150th anniversary of the Church in Scandinavia. Between meetings, we took a few hours to search for villages where two of my father’s grandparents were born. They were among the early converts to the Church in Denmark. Father’s paternal grandmother’s family lived in the western part of the country. His paternal grandfather’s family lived in northern Denmark. Thanks to a good driver and a superb map, we found each town on our list and obtained treasured information. During the entire journey, my hands were riveted to that valuable map so essential to achieve our goals.

In contrast, many people travel through life without good guidance, lacking knowledge of a desired destination or how to get there. But if rapt attention is paid to a road map for a day’s journey, isn’t it also wise to pay attention to authoritative guidance on our journey through life? To this end I would like to speak—on why we need guidance, where we obtain it, and how we can achieve it.

\intalkheader{Why We Need Guidance}

The question why focuses on the purpose of life. The ultimate objective in our mortal journey has been revealed by our Creator, who said, “If you keep my commandments and endure to the end you shall have eternal life, which gift is the greatest of all the gifts of God.”

His gift of eternal life is subject to conditions established by Him. Those conditions constitute a plan or, to use my analogy, a spiritual road map. And when trouble comes, guidance is needed most. In our journey in Denmark, we met an unexpected detour that led us astray. In order to get back on course, we stopped the car. We studied the map with great care. Then we made the necessary course correction.

What if you are lost and have no map? Suppose you are alone. You do not know where you are. What can you do? You call for help! You call home! Call the Church! Pray! When connected with your help line, you learn that you need to make a climb here or a turn there to get back on course. Or you may have to go back to the beginning in order to be certain that you can get where you want to go.

\intalkheader{Where We Obtain Guidance}

That brings us to the question of where do we obtain the guidance we need. We turn to Him who knows us best—our Creator. He allowed us to come to earth with freedom to choose our own course. In His great love, He did not leave us alone. He provided a guide—a spiritual road map—to help us achieve success in our journey. We call that guide the standard works, so named because they—the Holy Bible, the Book of Mormon, the Doctrine and Covenants, and the Pearl of Great Price—constitute the standard by which we should live. They serve as a standard of reference, as are standards of time, weights, and measures that are kept in national bureaus of standards.

To reach our objective of eternal life, we need to follow teachings in the standard works and other revelations received from prophets of God. Our loving Lord foresaw our need for guidance: “For strait is the gate,” He said, “and narrow the way that leadeth unto the exaltation and continuation of the lives, and few there be that find it.”

Few find the way, because they ignore the divine road map provided by the Lord. An even more serious mistake is to ignore the Maker of the map. God declared in the first of His Ten Commandments, “Thou shalt have no other gods before me.” Yet carnal man tends to let his loyalty drift toward idols.

For example, we marvel at computers and the Internet that enable transmission of data with remarkable speed. We are truly grateful for these electronic servants. But if we let them take over our time, pervert our potential, or poison our minds with pornography, they cease being servants and become instead false gods.

The Master warned of those who “seek not the Lord to establish his righteousness, but every man walketh in his own way, and after the image of his own god, whose image is in the likeness of the world, and whose substance is that of an idol.”

False gods can only lead to dead ends. If our journey through life is to be successful, we need to follow divine direction. The Lord said, “Look unto me in every thought; doubt not, fear not.” And the Psalmist wrote, “Thy word is a lamp unto my feet, and a light unto my path.”

Following such counsel demands not only conviction but conversion and often repentance. That would please the Lord, who said, “Repent, and turn yourselves from your idols; and turn … from all your abominations.”

In your journey through life, you meet many obstacles and make some mistakes. Scriptural guidance helps you to recognize error and make the necessary correction. You stop going in the wrong direction. You carefully study the scriptural road map. Then you proceed with repentance and restitution required to get on the “strait and narrow path which leads to eternal life.”

Brothers and sisters, our busy lives force us to focus on things we do from day to day. But the development of character comes only as we focus on who we really are. To establish and accomplish those greater goals, we do need heavenly help.

\intalkheader{How We Can Achieve Scriptural Guidance}

Once we understand why we need guidance and where we obtain it, we then ask, how can we achieve it? How can we truly live, not “by bread alone, but by every word that proceedeth out of the mouth of God”?

We begin with a determination to “liken all scriptures unto us … for our profit and learning.” If we “press forward, feasting upon the word of Christ, and endure to the end, … [we] shall have eternal life.”

To feast means more than to taste. To feast means to savor. We savor the scriptures by studying them in a spirit of delightful discovery and faithful obedience. When we feast upon the words of Christ, they are embedded “in fleshy tables of the heart.” They become an integral part of our nature.

Many years ago a medical colleague chastised me for failing to separate my professional knowledge from my religious convictions. That startled me because I did not feel that truth should be fractionalized. Truth is indivisible.

Danger lurks when we divide ourselves with expressions such as “my private life,” “my professional life,” or even “my best behavior.” Living life in separate compartments can lead to internal conflict and exhausting tension. To escape that tension, many people unwisely resort to addicting substances, pleasure seeking, or self-indulgence, which in turn produce more tension, thus creating a vicious cycle.

Inner peace comes only as we maintain the integrity of truth in all aspects of our lives. When we covenant to follow the Lord and obey His commandments, we accept His standards in every thought, action, and deed.

Living the Lord’s standards requires that we cultivate the gift of the Holy Ghost. That gift helps us understand doctrine and apply it personally. Because truth given by revelation can only be understood by revelation, our studies need to be prayerful. Scriptures attest to the efficacy of prayer in daily life. One is in Proverbs: “In all thy ways acknowledge [God], and he shall direct thy paths.” Another comes from the Book of Mormon: “Counsel with the Lord in all thy doings, and he will direct thee for good.”

As you ponder and pray about doctrinal principles, the Holy Ghost will speak to your mind and your heart. From events portrayed in the scriptures, new insights will come and principles relevant to your situation will distill upon your heart.

You cultivate such revelatory experiences by living according to the light already given you and by searching the scriptures with pure motives—with real intent to “come unto Christ.” As you do so, your confidence will “wax strong in the presence of God,” and the Holy Ghost will be your constant companion.

Achieving scriptural guidance is aided by posing pertinent questions. You might ask, “What principle can be learned from these teachings of the Lord?” For example, scriptures teach that the Creation was accomplished in six periods of time. Principles learned from that study show that any great attainment requires proper planning, timing, patience, labor, and no shortcuts.

Next, I suggest that you shape the style of your study to fit you. One way is to read a book of scripture from the first page to the last. This method gives good overall perspective. But other approaches also have merit. Attention to a particular topic or a specific theme, supplemented by use of cross-referencing footnotes and study guides, can help to switch on the light of doctrinal understanding.

Guidance can come when grappling with a serious challenge in life. Years ago, in the days of my early scientific research in a field that was then new to medical practice, a scriptural standard of truth gave me the courage needed to persevere. I leaned heavily upon these verses in the Doctrine and Covenants:

“All kingdoms have a law given;

“And there are many kingdoms; for there is no space in the which there is no kingdom; and there is no kingdom in which there is no space. …

“And unto every kingdom is given a law; and unto every law there are certain bounds also and conditions.” We learned laws that pertained to the “kingdom” of our concern and mastered control that had previously been relegated by ignorance to chance alone.

Motivation for scriptural guidance comes when important choices must be made—even between options that are equally right. The Brethren are often faced with these kinds of decisions. On such an occasion, we turn to the scriptures. We may read all of the standard works afresh, looking for insights relative to a specific issue.

Time for scripture study requires a schedule that will be honored. Otherwise, blessings that matter most will be at the mercy of things that matter least. Time for family scripture study may be difficult to establish. Years ago when our children were at home, they attended different grades in several schools. Their daddy had to be at the hospital no later than 7:00 in the morning. In family council we determined that our best time for scripture study was 6:00 a.m. At that hour our little ones were very sleepy but supportive. Occasionally we had to awaken one when a turn came to read. I would be less than honest with you if I conveyed the impression that our family scripture time was a howling success. Occasionally it was more howling than successful. But we did not give up.

Now, a generation later, our children are all married with families of their own. Sister Nelson and I have watched them enjoy family scripture study in their own homes. Their efforts are much more successful than were ours. We shudder to think what might have happened if we had quit trying.

We all need guidance through life. We obtain it best from the standard works and teachings of the prophets of God. With diligent effort, we can achieve that guidance and thus qualify for all of the blessings that God has in store for His faithful children. I so testify in the name of Jesus Christ, amen.

\newpage
\settalk{Personal Preparation for Temple Blessings}
\setspeaker{Russell M. Nelson}
\settalkdate{April 2001}
\talkheading{Personal Preparation for Temple Blessings}{Russell M. Nelson}

Under President Gordon B. Hinckley’s inspired direction, temples are now more accessible than ever. Inscribed on each temple are the words “Holiness to the Lord.” That statement designates both the temple and its purposes as holy. Those who enter the temple are also to bear the attribute of holiness. It may be easier to ascribe holiness to a building than it is to a people. We can acquire holiness only by enduring and persistent personal effort. Through the ages, servants of the Lord have warned against unholiness. Jacob, brother of Nephi, wrote, “I would speak unto you of holiness; but as ye are not holy, and ye look upon me as a teacher, [I] must … teach you the consequences of sin.”

Now I feel that same sense of responsibility to teach. As temples are prepared for our members, our members need to prepare for the temple.

\intalkheader{The Temple}

The temple is the house of the Lord. The basis for every temple ordinance and covenant—the heart of the plan of salvation—is the Atonement of Jesus Christ. Every activity, every lesson, all we do in the Church points to the Lord and His holy house. Our efforts to proclaim the gospel, perfect the Saints, and redeem the dead all lead to the temple. Each holy temple stands as a symbol of our membership in the Church, as a sign of our faith in life after death, and as a sacred step toward eternal glory for us and our families.

President Hinckley said that “these unique and wonderful buildings, and the ordinances administered therein, represent the ultimate in our worship. These ordinances become the most profound expressions of our theology.”

To enter the temple is a tremendous blessing. But first we must be worthy. We should not be rushed. We cannot cut corners of preparation and risk the breaking of covenants we were not prepared to make. That would be worse than not making them at all.

\intalkheader{The Endowment}

In the temple we receive an endowment, which is, literally speaking, a gift. In receiving this gift, we should understand its significance and the importance of keeping sacred covenants. Each temple ordinance “is not just a ritual to go through, it is an act of solemn promising.”

The temple endowment was given by revelation. Thus, it is best understood by revelation, prayerfully sought with a sincere heart. President Brigham Young said, “Your endowment is, to receive all those ordinances in the house of the Lord, which are necessary for you, after you have departed this life, to enable you to walk back to the presence of the Father … and gain your eternal exaltation.”

\intalkheader{Sealing Authority}

In preparing to receive the endowment and other ordinances of the temple, we should understand the sealing authority of the priesthood. Jesus referred to this authority long ago when He taught His Apostles, “Whatsoever thou shalt bind on earth shall be bound in heaven.” That same authority has been restored in these latter days. Just as priesthood is eternal—without beginning or end—so is the effect of priesthood ordinances that bind families together forever.

Temple ordinances, covenants, endowments, and sealings enable individuals to be reconciled with the Lord and families to be sealed beyond the veil of death. Obedience to temple covenants qualifies us for eternal life, the greatest gift of God to man. Eternal life is more than immortality. Eternal life is exaltation in the highest heaven—the kind of life that God lives.

\intalkheader{Temple Recommend}

Preparation also includes qualification for a temple recommend. Our Redeemer requires that His temples be protected from desecration. No unclean thing may enter His hallowed house. Yet anyone is welcome who prepares well. Each person applying for a recommend will be interviewed by a judge in Israel—the bishop—and by a stake president. They hold keys of priesthood authority and the responsibility to help us know when our preparation and timing are appropriate to enter the temple. Their interviews will assess several vital issues. They will ask if we obey the law of tithing, if we keep the Word of Wisdom, and if we sustain the authorities of the Church. They will ask if we are honest, if we are morally clean, and if we honor the power of procreation as a sacred trust from our Creator.

Why are these issues so crucial? Because they are spiritual separators. They help to determine if we truly live as children of the covenant, able to resist temptation from servants of sin. These interviews help to discern if we are willing to live in accord with the will of the true and living God or if our hearts are still set “upon riches and … vain things of the world.”

Such requirements are not difficult to understand. Because the temple is the house of the Lord, standards for admission are set by Him. One enters as His guest. To hold a temple recommend is a priceless privilege and a tangible sign of obedience to God and His prophets.

\intalkheader{Physical Preparation for the Temple}

One prepares physically for the temple by dressing properly. It is not a place for casual attire. “We should dress in such a way that we might comfortably attend a sacrament meeting or a gathering that is proper and dignified.”

Within the temple, all are dressed in spotless white to remind us that God is to have a pure people. Nationality, language, and position in the Church are of secondary significance. In that democracy of dress, all sit side by side and are considered equal in the eyes of our Maker.

Brides and grooms enter the temple to be married for time and all eternity. There brides wear white dresses—long sleeved, modest in design and fabric, and free of elaborate ornamentation. Grooms also dress in white. And brethren who come to witness weddings do not wear tuxedos.

Wearing the temple garment has deep symbolic significance. It represents a continuing commitment. Just as the Savior exemplified the need to endure to the end, we wear the garment faithfully as part of the enduring armor of God. Thus we demonstrate our faith in Him and in His eternal covenants with us.

\intalkheader{Spiritual Preparation for the Temple}

In addition to physical preparation, we prepare spiritually. Because the ordinances and covenants of the temple are sacred, we are under solemn obligation not to speak outside the temple of that which occurs in the temple. There are, however, some principles we can discuss.

Each temple is a house of learning. There we are taught in the Master’s way. His way differs from modes of others. His way is ancient and rich with symbolism. We can learn much by pondering the reality for which each symbol stands. Teachings of the temple are beautifully simple and simply beautiful. They are understood by the humble, yet they can excite the intellect of the brightest minds.

Spiritual preparation is enhanced by study. I like to recommend that members going to the temple for the first time read short explanatory paragraphs in the Bible Dictionary, listed under seven topics: “Anoint,” “Atonement,” “Christ,” “Covenant,” “Fall of Adam,” “Sacrifices,” and “Temple.” Doing so will provide a firm foundation.

One may also read in the Old Testament and the books of Moses and Abraham in the Pearl of Great Price. Such a review of ancient scripture is even more enlightening after one is familiar with the temple endowment. Those books underscore the antiquity of temple work.

With each ordinance is a covenant—a promise. A covenant made with God is not restrictive, but protective. Such a concept is not new. For example, if our water supply is not clean, we filter the water to screen out harmful ingredients. Divine covenants help us to filter out of our minds impurities that could harm us. When we choose to deny ourselves of all ungodliness, we lose nothing of value and gain the glory of eternal life. Covenants do not hold us down; they elevate us beyond the limits of our own power and perspective.

\intalkheader{Eternal Perspective}

President Hinckley has explained that lofty perspective: “There is a goal beyond the Resurrection,” he said. “That is exaltation in our Father’s kingdom. … It will begin with acceptance of him as our Eternal Father and of his Son as our living Redeemer. It will involve participation in various ordinances, each one important and necessary. The first of these is baptism by immersion in water, without which, according to the Savior, a man cannot enter into the kingdom of God. There must follow the birth of the Spirit, the gift of the Holy Ghost. Then in succession through the years will come, for men, ordination to the priesthood, followed by the blessings of the temple for both men and women who are worthy to enter therein. These temple blessings include our washings and anointings that we may be clean before the Lord. They include the … endowment of obligations and blessings that motivate us to behavior compatible with the principles of the gospel. They include the sealing ordinances by which that which is bound on earth is bound in heaven, providing for the continuity of the family.”

I have learned that temple blessings are most meaningful when death takes a loved one away from the family circle. To know that the pain of separation is only temporary provides peace that passes ordinary understanding. Death cannot sever families sealed in the temple. They understand death as a necessary part of God’s great plan of happiness.

Such perspective helps us to maintain fidelity to covenants made. President Boyd K. Packer emphasized that “ordinances and covenants become our credentials for admission into [God’s] presence. To worthily receive them is the quest of a lifetime; to keep them thereafter is the challenge of mortality.”

Ordinances of the temple relate to personal progress and to the redemption of departed ancestors as well. “For their salvation is necessary and essential to our salvation, … they without us cannot be made perfect—neither can we without our dead be made perfect.” Service in their behalf provides repeated opportunities for temple worship. And that service deserves commitment to a planned schedule. By doing for others what they cannot do for themselves, we emulate the pattern of the Savior, who wrought the Atonement to bless the lives of other people.

One day we will meet our Maker and stand before Him at Judgment. We will be judged according to our ordinances, covenants, deeds, and the desires of our hearts.

Meanwhile, in this world smitten with spiritual decay, can individuals prepared for temple blessings make a difference? Yes! Those Saints are “the covenant people of the Lord, … armed with righteousness and with the power of God in great glory.” Their example can lift the lives of all humankind. I so testify in the name of Jesus Christ, amen.

\newpage
\settalk{“Set in Order Thy House”}
\setspeaker{Russell M. Nelson}
\settalkdate{October 2001}
\talkheading{“Set in Order Thy House”}{Russell M. Nelson}

Years ago when Sister Nelson and I had several teenaged daughters, we took our family on a vacation far away from telephones and boyfriends. We went on a raft trip down the Colorado River through the Grand Canyon. As we started our journey, we had no idea how dangerous this trip could be.

The first day was beautiful. But on the second day, when we approached Horn Creek rapids and saw that precipitous drop ahead, I was terrified. Floating on a rubber raft, our precious family was about to plunge over a waterfall! Instinctively I put one arm around my wife and the other around our youngest daughter. To protect them, I tried to hold them close to me. But as we reached the precipice, the bended raft became a giant sling and shot me into the air. I landed into the roiling rapids of the river. I had a hard time coming up. Each time I tried to find air, I hit the underside of the raft. My family couldn’t see me, but I could hear them shouting, “Daddy! Where’s Daddy?”

I finally found the side of the raft and rose to the surface. The family pulled my nearly drowned body out of the water. We were thankful to be safely reunited.

The next several days were pleasant and delightful. Then came the last day, when we were to go over Lava Falls, known as the most dangerous drop of the journey. When I saw what was ahead, I immediately asked to beach the raft and hold an emergency family council meeting, knowing that if we were to survive this experience, we needed to plan carefully. I reasoned with our family: “No matter what happens, the rubber raft will remain on top of the water. If we cling with all our might to ropes secured to the raft, we can make it. Even if the raft should capsize, we will be all right if we hang tightly to the ropes.”

I turned to our little seven-year-old daughter and said, “All of the others will cling to a rope. But you will need to hold on to your daddy. Sit behind me. Put your arms around me and hold me tightly while I hold the rope.”

That we did. We crossed those steep, rough rapids—hanging on for dear life—and all of us made it safely.

\intalkheader{The Lesson}

Brothers and sisters, I nearly lost my life learning a lesson that I now give to you. As we go through life, even through very rough waters, a father’s instinctive impulse to cling tightly to his wife or to his children may not be the best way to accomplish his objective. Instead, if he will lovingly cling to the Savior and the iron rod of the gospel, his family will want to cling to him and to the Savior.

This lesson is surely not limited to fathers. Regardless of gender, marital status, or age, individuals can choose to link themselves directly to the Savior, hold fast to the rod of His truth, and lead by the light of that truth. By so doing, they become examples of righteousness to whom others will want to cling.

\intalkheader{The Commandment}

With the Lord, families are essential. He created the earth that we could gain physical bodies and form families. He established His Church to exalt families. He provides temples so that families can be together forever.

Of course, He expects fathers to preside over, provide for, and protect their families. But the Master has asked for much more. Etched in sacred scripture is a commandment to “set in order thy house.” Once we as parents understand the importance and meaning of that commandment, we need to learn how to do it.

\intalkheader{How to Set Your House in Order}

To set our house in an order pleasing to the Lord, we need to do it His way. We are to employ His attributes of “righteousness, godliness, faith, love, patience, [and] meekness.” Each father should remember that “no power or influence can or ought to be maintained by virtue of the priesthood, only by persuasion, by long-suffering, by gentleness and meekness, and by love unfeigned.”

Parents are to be living examples of “kindness, and pure knowledge, which … greatly enlarge the soul.” Each mother and father should lay aside selfish interests and avoid any thought of hypocrisy, physical force, or evil speaking. Parents soon learn that each child has an inborn yearning to be free. Each individual wants to make his or her own way. No one wants to be restrained, even by a well-intentioned parent. But all of us can cling to the Lord.

Ages ago, Job taught that concept. He said, “My righteousness I hold fast, and will not let it go.” Nephi also taught, “Whoso would hearken unto the word of God, and … hold fast unto it, … would never perish.”

These tenets are timeless as the gospel and endless as eternity. Ponder these additional scriptural admonitions:

From the Old Testament Proverbs we read, “Take fast hold of instruction; let her not go: keep her; for she is thy life.”

From the New Testament: “Brethren, stand fast, and hold the traditions which ye have been taught.”

From the Book of Mormon we learn about multitudes who were “continually holding fast to the rod of iron,” which Nephi likened to “the word of God.” Anchored in truth, that iron rod is immovable and immutable.

\intalkheader{Other Divine Mandates}

Not only are parents to cling to the word of the Lord, but they have a divine mandate to teach it to their children. Scriptural direction is very clear: “Inasmuch as parents have children in Zion … that teach them not to understand the doctrine of repentance, faith in Christ the Son of the living God, and of baptism and the gift of the Holy Ghost by the laying on of the hands, when eight years old, the sin be upon the heads of the parents.”

That commandment places responsibility and accountability for the teaching of children squarely upon the shoulders of the parents. The proclamation to the world regarding the family warns that individuals “who fail to fulfill family responsibilities will one day stand accountable before God.” Today I solemnly reaffirm that reality.

In discharging these duties, we need both the Church and the family. They work hand in hand to strengthen each other. The Church exists to exalt the family. And the family is the fundamental unit of the Church.

These interrelationships are evident as we study the early history of the Church. In 1833 the Lord rebuked young leaders of His Church because of parental shortcomings. The Lord said:

“I have commanded you to bring up your children in light and truth.

“But verily I say unto you, …

“You have not taught your children light and truth, according to the commandments. …

“And now a commandment I give unto you [that] you shall set in order your own house, for there are many things that are not right in your house. … First set in order thy house.”

This revelation represents one of the many powerful validations of the integrity of the Prophet Joseph Smith. He did not delete from scripture words of stinging rebuke, even though some were directed to himself.

In our day, the First Presidency has again stressed parental priority. From their recent letter to the Saints, I quote: “We call upon parents to devote their best efforts to the teaching and rearing of their children in gospel principles which will keep them close to the Church. The home is the basis of a righteous life, and no other instrumentality can take its place or fulfill its essential functions in carrying forward this God-given responsibility.”

\intalkheader{What Should Parents Teach?}

With this sacred charge in mind, let us consider what we should teach. Scriptures direct parents to teach faith in Jesus Christ, repentance, baptism, and the gift of the Holy Ghost. Parents are to teach the plan of salvation and the importance of living in complete accord with the commandments of God. Otherwise, their children will surely suffer in ignorance of God’s redeeming and liberating law. Parents should also teach by example how to consecrate their lives—using their time, talents, tithing, and substance to establish the Church and kingdom of God upon the earth. Living in that manner will literally bless their posterity. A scripture states, “Thy duty is unto the church forever, and this because of thy family.”

\intalkheader{Opposition to the Family}

Parents and children should realize that strong opposition will always come against the work and will of the Lord. Because the work (and glory) of God is to bring to pass our immortality and eternal life as a family, it logically follows that the work of the adversary will strike directly at the heart of the home—the family. Relentlessly Lucifer attacks the sanctity of life and the joy of parenthood.

Because the evil one is ever at work, our vigilance cannot be relaxed—not even for a moment. A small and seemingly innocent invitation can turn into a tall temptation which can lead to tragic transgression. Night and day, at home or away, we must shun sin and “hold fast that which is good.”

The seditious evils of pornography, abortion, and addiction to harmful substances serve as termites to erode the undergirding strength of a happy home and a faithful family. We cannot yield to any iniquity without putting our families at risk.

Satan wants us to be miserable just as he is. He would animate our carnal appetites and entice us to live in spiritual darkness and doubt the reality of life after death. The Apostle Paul observed, “If in this life only we have hope in Christ, we are of all men most miserable.”

\intalkheader{Perpetuation of Family Blessings}

An understanding of God’s great plan of happiness, however, fortifies our faith in the future. His plan provides answers to ageless questions: Are all our sympathies and love for each other only temporary—to be lost in death? No! Can family life endure beyond this period of mortal probation? Yes! God has revealed the eternal nature of celestial marriage and the family as the source of our greatest joy.

Brethren and sisters, material possessions and honors of the world do not endure. But your union as wife, husband, and family can. The only duration of family life that satisfies the loftiest longings of the human soul is forever. No sacrifice is too great to have the blessings of an eternal marriage. To qualify, one needs only to deny oneself of ungodliness and honor the ordinances of the temple. By making and keeping sacred temple covenants, we evidence our love for God, for our companion, and our real regard for our posterity—even those yet unborn. Our family is the focus of our greatest work and joy in this life; so will it be throughout all eternity, when we can “inherit thrones, kingdoms, principalities, … powers, dominions, … exaltation and glory.”

These priceless blessings can be ours if we set our houses in order now and faithfully cling to the gospel. God lives. Jesus is the Christ. This is His Church. President Gordon B. Hinckley is His prophet. I so testify in the name of Jesus Christ, amen.

\newpage
\settalk{How Firm Our Foundation}
\setspeaker{Russell M. Nelson}
\settalkdate{April 2002}
\talkheading{How Firm Our Foundation}{Russell M. Nelson}

About two decades ago when a temple was to be built in Mexico City, architects faced a great challenge. Because Mexico City is situated on a basin over water, some of its buildings settle and become tilted with the passage of time. Construction of a temple there required a special foundation. Two hundred twenty-one large, reinforced concrete piles were driven more than 100 feet deep into the ground. Onto those pilings, steel straps were anchored and fastened to a unit that can be adjusted if need be to keep the building level. With this unseen but sure foundation, that temple today stands steady and straight.

A firm foundation is necessary for any building, institution, or individual to endure. With that in mind, let us consider the foundation of The Church of Jesus Christ of Latter-day Saints. Then let us see how the solid foundation of the Church supports our foundation of faith as individual members of the Church.

\intalkheader{I. Foundation of The Church of Jesus Christ of Latter-day Saints}

This Church stands on a unique foundation, anchored to a bedrock of timeless truth. Brothers and sisters, the holy cause in which we are engaged did not begin in 1820 in the state of New York. It did not begin in Bethlehem. It did not begin in the Garden of Eden. The underpinnings of the everlasting gospel were in place even before the world was.

This reality is affirmed repeatedly in holy writ. I have studied the scriptures pertaining to that part of eternity before the earth was formed. Don’t worry. I won’t quote them all. But I will append some of those citations to the published text of this message. These unseen and timeless truths constitute “premortal pillars” that undergird the foundation of this Church.

The Plan of Redemption and Possibility of an Inheritance in the Kingdom of God

Before the foundation of the earth, the plan of salvation was prepared. It included the glorious possibility of a divine inheritance in the kingdom of God.

The Atonement of Jesus Christ

Central to that plan was the Atonement of Jesus Christ. In premortal councils, He was foreordained by His Father to atone for our sins and break the bands of physical and spiritual death. Jesus declared: “I … was prepared from the foundation of the world to redeem my people. … In me shall all mankind have life, and that eternally, even they who shall believe on my name.” Later, Paul added that the Church is “built upon the foundation of the apostles and prophets, Jesus Christ himself being the chief corner stone.”

Opposition from the Adversary

Even before the world was formed, the Lord was opposed by Satan. He and his minions have always fought against the Savior’s sacred work and will continue to do so.

Foreordination for and Conferral of the Priesthood

Before the foundation of the world, each of us existed as spirit children with our Father in Heaven. Among us were noble and great men foreordained to bear the priesthood. Abraham, Jeremiah, Joseph Smith, and others were destined to become prophets of God. Also foreseen was that the blood of prophets would be shed in this holy cause. The manner of conferring the priesthood was to come down from the fathers “from the beginning of time, … even … before the foundation of the earth.”

Sisters Empowered to Glorify God

Sisters received special gifts. They, according to the Lord, were empowered “to multiply and replenish the earth, according to my commandment, and to fulfil the promise which was given by my Father before the foundation of the world, … for their exaltation in the eternal worlds, that they may bear the souls of men; … herein is the work of my Father continued, that he may be glorified.” Think of it: When a mother bears and cares for a child, she not only helps the earth answer the end of its creation, but she glorifies God!

Children of the Covenant

Children of the covenant were designated in the premortal realm. Paul taught that the Lord “hath chosen us in him before the foundation of the world.”

Blessings Predicated upon Obedience to Law

Then conditions were established that allowed us to receive blessings from God—upon obedience to laws upon which those blessings were predicated.

Sacred Things to Be Revealed in the Latter Days

Another pillar pertains to revelation reserved for the latter days. The Lord had long planned to reveal “things which have been kept secret from the foundation of the world.” Those things included revelations recorded in the Book of Mormon. They also included ordinances and covenants of the holy temple. To the Prophet Joseph Smith the Lord declared that “I deign to reveal unto my church things which have been kept hid from before the foundation of the world, things that pertain to the dispensation of the fulness of times.”

Redemption of the Dead

The Lord revealed that “the ordinance of baptizing for the dead … was instituted from before the foundation of the world.” Salvation was thus made available for those “who should die without a knowledge of the gospel.” A welding link between generations was provided, that a whole, complete, and perfect union of dispensations, keys, powers, and glories should take place.

Brothers and sisters, these unseen but sure pillars were in place before the world was. They undergird the everlasting gospel—now restored in its fulness. With such a foundation, this Church will not be moved from its place, even through the Millennium.

\intalkheader{II. Foundation of Individual Members of the Church}

Just as buildings and institutions have foundations, so do we as individuals have foundations that support our faith. Some are weak; some are strong. We can waffle “like a wave of the sea driven with the wind and tossed,” or we can stand on a firm foundation and anchor ourselves with straps of spiritual steel, rooted and grounded to the timeless pillars of the gospel.

President Gordon B. Hinckley has asked that our new converts be sustained in the Church by a friend, a call to serve, and be nourished by the good word of God. Often we sing, “How firm a foundation, ye Saints of the Lord, / Is laid for your faith in his excellent word!” When we receive that nurturing word of God, we feast upon His love.

Remember: “It is upon the rock of our Redeemer, who is Christ, the Son of God, that ye must build your foundation; that when the devil shall send forth his mighty winds, … when all his hail and his mighty storm shall beat upon you, it shall have no power … to drag you down to the gulf of misery and endless wo, because of the rock upon which ye are built, which is a sure foundation.”

Even firm foundations cannot prevent life’s problems. Wayward children cause parents to grieve. Some broken families don’t get fixed. Gender disorientation is poorly understood. Married couples, for whatever reason, may not be blessed with children. Even in our day, “the guilty and the wicked go unpunished because of their money.” Some things just don’t seem fair.

With strong underpinnings, however, we are better able to reach upward for help, even when faced with questions without easy answers. The poet so wrote:

\begin{verse}
Who Thou art I know not,\\
But this much I know:\\
Thou hast set the Pleiades\\
In a silver row;
\end{verse}

Though we don’t know all things, we know that God lives and that He loves us. Standing on that firm foundation, we can reach up and find strength to endure the heavy burdens of life.

For example, I honor those special souls who face challenges of parenthood with unwavering faith in their Maker. To a family dear to Sister Nelson and me, a son was recently born. This child was afflicted with multiple anomalies affecting virtually every system of his little body. Two operations were required in his first week of life. More will be necessary. When I spoke with the child’s parents, they did not ask, “Why did this happen to us?” Instead, they declared: “We know that this child is meant for us. God has entrusted this special baby to us. We will love him and care for him to the best of our ability.” Thank the Lord for such parents!

Not long ago, death came to the husband of a friend of ours, suddenly and without warning. Of her, these words were penned by her own sister: “Searching the years and days just past, she is awestruck, recognizing specific skills and experiences put into her life by a loving Father, things that might have seemed circumstantial at the time but that have specifically prepared her to cope successfully with this heartbreaking loss. Rather than feeling abandoned and bitter, she feels cradled and cared for. … She said to me, ‘When I see how carefully Heavenly Father has prepared and planned for my present circumstance, how can I be frightened about my future? Surely He is putting into place today all that I will need to face the unknown times ahead.’”

From a missionary I received an unforgettable letter. He wrote: “I still am not certain why it was that I was [afflicted] with [cancer], particularly during the time I was serving the Lord on a mission, but I can say with all honesty and sincerity that I am eternally grateful to our gracious Heavenly Father for allowing me to have that experience. … Not a day goes by,” he continued, “that I don’t think about the days I spent lying in the hospital suffering through chemotherapy or grimacing with the pain from another operation. … Not a day … passes when I don’t think of the days I spent studying the scriptures, particularly the Book of Mormon, and remembering the overwhelming feelings of comfort and peace which I felt. I often think of the nights when I would retire to bed and pour out my soul to my Heavenly Father and thank Him for preserving my life.” Then the elder shared this wonderful news: “I returned to the doctor this week … and … he found no evidence of any disease in my body.” I love such faithful missionaries!

How firm our foundation? The foundation of this Church was forged long before the world was. It is strong. It is true. It is eternal. The foundation of one’s individual faith, if anchored firmly to eternal truth, allows each of us to reach upward with an eternal perspective. That faith will give hope when there is no hope. It will give us joy here and eternal life hereafter. I so testify in the name of Jesus Christ, amen.

\newpage
\settalk{“Blessed Are the Peacemakers”}
\setspeaker{Russell M. Nelson}
\settalkdate{October 2002}
\talkheading{“Blessed Are the Peacemakers”}{Russell M. Nelson}

Among the most memorable of all our family activities have been trips to the Holy Land. For us, visits to that part of the world have been life changing. But now the Holy Land is a cauldron of turmoil and off-limits to those who would like to go there for spiritual enrichment. Virtually all parts of the world are plagued by acts of terror previously unknown. Confusion comes to many who pray for peace but fearfully face this foe of terror.

\intalkheader{Peace versus Contention}

The perilous times in which we live have been prophesied in the scriptures. Our day has been foreseen as one with “fires, and tempests, and vapors of smoke in foreign lands; … wars, rumors of wars, and earthquakes in divers places … great pollutions upon the face of the earth; … and all manner of abominations.”

That prophecy echoes an earlier scriptural account of the second generation of human life upon the earth: “In those days Satan had great dominion among men, and raged in their hearts; and from thenceforth came wars and bloodshed; and a man’s hand was against his own brother, in … seeking for power.” From the days of Cain and Abel, Esau and Jacob, and Joseph who was sold into Egypt, the flames of hostility have been fueled by family feuding.

Hatred among brothers and neighbors has now reduced sacred cities to sites of sorrow. When I think of the plight of such places, I am reminded of an ancient proverb: “Scornful men bring a city into a snare: but wise men turn away wrath.”

\intalkheader{Doctrinal Direction}

Scripture sheds light on both the cause of and the cure for the sickness of human hatred: “The natural man is an enemy to God, and has been from the fall of Adam, and will be, forever and ever, unless he yields to the enticings of the Holy Spirit, and putteth off the natural man and becometh a saint through the atonement of Christ.”

Peace can prevail only when that natural inclination to fight is superseded by self-determination to live on a loftier level. Coming unto Jesus Christ as the “Prince of Peace” is the pathway to peace on earth and goodwill among men. He made a promise to us: “Blessed are the peacemakers: for they shall be called the children of God.”

Jesus taught people how to live with one another. He declared the two great commandments: first, to “love the Lord thy God with all thy heart, and with all thy soul, and with all thy mind,” and the second, to “love thy neighbour as thyself.”

Then He added, “Love your enemies, [and] bless them that curse you.”

He taught the Golden Rule: “All things whatsoever ye would that men should do to you, do ye even so to them.” This principle is found in nearly every major religion. Others such as Confucius and Aristotle have also taught it. After all, the gospel did not begin with the birth of the Babe in Bethlehem. It is everlasting. It was proclaimed in the beginning to Adam and Eve. Portions of the gospel have been preserved in many cultures. Even heathen mythologies have been enriched by fragments of truth from earlier dispensations.

Wherever it is found and however it is expressed, the Golden Rule encompasses the moral code of the kingdom of God. It forbids interference by one with the rights of another. It is equally binding upon nations, associations, and individuals. With compassion and forbearance, it replaces the retaliatory reactions of “an eye for an eye, and a tooth for a tooth.” If we were to stay on that old and unproductive path, we would be but blind and toothless.

This concept of treating others as one would like to be treated is easy to understand. And it acknowledges the precious nature of each of God’s sons and daughters. Scripture asks parents to teach their children not to “fight and quarrel one with another, and serve the devil, who is the master of sin.” Instead, we “teach them to love one another, and to serve one another.”

Jesus taught the importance of reconciliation and resolution of dispute on a personal basis. He said:

“Whosoever is angry with his brother … shall be in danger of the judgment. …

“Therefore if thou bring thy gift to the altar, and there rememberest that thy brother hath ought against thee;

“Leave there thy gift before the altar, and go thy way; first be reconciled to thy brother, and then come and offer thy gift.”

The Master Teacher taught us to “forgive, if ye have ought against any: that your Father also which is in heaven may forgive you your trespasses.

“But if ye do not forgive, neither will your Father which is in heaven forgive your trespasses.”

Jesus declared that a day of judgment would come. All individuals will give an account of their mortal lives and of how they have treated other people.

\intalkheader{Civic Obligations}

The commandments to love God and neighbor are interrelated. We cannot fully love God without loving our neighbor. We cannot fully love our neighbor without loving God. Men really are brothers because God really is our Father. Nevertheless, scriptures are studded with stories of contention and combat. They strongly condemn wars of aggression but sustain obligations of citizens to defend their families and their freedoms. Because “we believe in being subject to kings, presidents, rulers, and magistrates, in obeying, honoring, and sustaining the law,” members of this Church will be called into military service of many nations. “We believe that governments were instituted of God for the benefit of man; and that he holds men accountable for their acts in relation to them, both in making laws and administering them, for the good and safety of society.”

During the Second World War, when members of the Church were forced to fight on opposing sides, the First Presidency affirmed that “the state is responsible for the civil control of its citizens or subjects, for their political welfare, and for the carrying forward of political policies, domestic and foreign. … But the Church itself, as such, has no responsibility for these policies, [other] than urging its members fully to render … loyalty to their country.”

\intalkheader{Peace Is Possible}

Because of the long history of hostility upon the earth, many feel that peace is beyond hope. I disagree. Peace is possible. We can learn to love our fellow human beings throughout the world. Whether they be Jewish, Islamic, or fellow Christians, whether Hindu, Buddhist, or other, we can live together with mutual admiration and respect, without forsaking our religious convictions. Things we have in common are greater than are our differences. Peace is a prime priority that pleads for our pursuit. Old Testament prophets held out hope and so should we. The Psalmist said, “God is our refuge and strength, a very present help in trouble.” “He maketh wars to cease unto the end of the earth.”

Father Abraham was uniquely called a “Friend of God.” Peace was one of Abraham’s highest priorities. He sought to be a “prince of peace.” His influence could loom large in our present pursuit of peace. His sons Ishmael and Isaac, though born of different mothers, overcame their differences when engaged in a common cause. After their father died, they worked together to bury the mortal remains of their exalted father. Their descendants could well follow that pattern.

Abraham’s posterity has a divinely decreed potential. The Lord declared that Ishmael would become a great nation and that the seed of Abraham, Isaac, and Jacob would bless all the nations of the earth.

So the descendants of Abraham, entrusted with great promises of infinite influence, are in a pivotal position to emerge as peacemakers. Chosen by the Almighty, they can direct their powerful potential toward peace.

Resolution of present political problems will require much patience and negotiation. The process would be enhanced greatly if pursued prayerfully.

Isaiah prophesied of hope for our day. Speaking of the gathering of Israel and the Restoration of the Church through the Prophet Joseph Smith, Isaiah wrote:

“It shall come to pass in that day, that the Lord shall set his hand again the second time to recover the remnant of his people. …

“And he shall set up an ensign for the nations, and shall assemble the outcasts of Israel, and gather together the dispersed of Judah from the four corners of the earth.”

These prophecies of hope could materialize if leaders and citizens of nations would apply the teachings of Jesus Christ. Ours could then be an age of unparalleled peace and progress. Barbarism of the past would be buried. War with its horrors would be relegated to the realm of maudlin memory. Aims of nations would be mutually supportive. Peacemakers could lead in the art of arbitration, give relief to the needy, and bring hope to those who fear. Of such patriots, future generations would shout praises, and our Eternal God would pass judgments of glory.

The hope of the world is the Prince of Peace—our Creator, Savior, Jehovah, and Judge. He offers us the good life, the abundant life, and eternal life. Peaceful—even prosperous—living can come to those who abide His precepts and follow His pathway to peace. This I declare to all the world.

\intalkheader{Church Members}

Now, as members of The Church of Jesus Christ of Latter-day Saints, what does the Lord expect of us? As a Church, we must “renounce war and proclaim peace.” As individuals, we should “follow after the things which make for peace.” We should be personal peacemakers. We should live peacefully—as couples, families, and neighbors. We should live by the Golden Rule. We have writings of the descendants of Judah as now merged with writings of the descendants of Ephraim. We should employ them and expand our circle of love to embrace the whole human family. We should bring divine love and revealed doctrines of restored religion to our neighbors and friends. We should serve them according to our abilities and opportunities. We should keep our principles on a high level and stand for the right. We should continue to gather scattered Israel from the four corners of the earth and offer the ordinances and covenants that seal families together forever. These blessings we are to bring to people of all nations.

As we so live, our Master will bless us. He provided this promise: “Fear thou not; for I am with thee: be not dismayed; for I am thy God: I will strengthen thee; yea, I will help thee; yea, I will uphold thee with the right hand of my righteousness.”

Ours is His supernal cause. Ours is the cause of eternal glory for all humankind. And as peacemakers, we shall be called the children of God. I so testify in the name of Jesus Christ, amen.

\newpage
\settalk{Sweet Power of Prayer}
\setspeaker{Russell M. Nelson}
\settalkdate{April 2003}
\talkheading{Sweet Power of Prayer}{Russell M. Nelson}

In this day of computers, phones, and pagers, people communicate with one another better than ever before. Even so, good communication is often lacking. Recently while visiting a nursing care facility, I spoke with a woman about her family. She told me that she had three sons, two of whom visit her regularly.

“What about your third son?” I asked.

“I don’t know where he is,” she replied tearfully. “I haven’t heard from him for years. I don’t even know how many grandchildren I have.”

\intalkheader{Why We Pray}

If such a mother yearns to hear from her sons, it is easy to see why a loving Father in Heaven wants to hear from His children. Through prayer, we can show our love for God. And He has made it so easy. We may pray to Him any time. No special equipment is needed. We don’t even need to charge batteries or pay a monthly service fee.

Some people pray only when confronted with personal problems. Others don’t pray at all. A scripture makes this observation: “Ye do not remember the Lord your God in the things with which he hath blessed you, but ye do always remember your riches, not to thank the Lord … for them.”

Prophets have long told us to pray humbly and frequently.

\intalkheader{How to Pray}

Jesus taught us how. We pray to our Heavenly Father, in the name of Jesus Christ, by the power of the Holy Ghost. This is the “true order of prayer,” in contrast to “vain repetitions” or recitations given to “be seen of men.”

Jesus revealed that we pray to a wise Father who knows what things we have need of, before we ask Him.

Mormon taught his son Moroni that we should pray “with all the energy of heart.” Nephi exclaimed, “I pray continually for [my people] by day, and mine eyes water my pillow by night, … and I cry unto my God in faith, and I know that he will hear my cry.”

The sweet power of prayer can be intensified by fasting, on occasion, when appropriate to a particular need.

Prayers can be offered even in silence. One can think a prayer, especially when words would interfere. We often kneel to pray; we may stand or be seated. Physical position is less important than is spiritual submission to God.

We close our prayers “in the name of Jesus Christ, amen.” When we hear another’s prayer, we audibly add our “amen,” meaning, “That is my prayer too.”

\intalkheader{When to Pray}

When should we pray? The Lord said, “Search diligently, pray always, and be believing, and all things shall work together for your good.”

Alma said, “Counsel with the Lord in all thy doings, and he will direct thee for good; yea, when thou liest down at night lie down unto the Lord, that he may watch over you in your sleep; and when thou risest in the morning let thy heart be full of thanks unto God.”

We pray privately, with our families regularly, at mealtime, and in daily activities. Simply summarized, we are a praying people.

\intalkheader{Personal Experience with Prayer}

Many of us have had experiences with the sweet power of prayer. One of mine was shared with a stake patriarch from southern Utah. I first met him in my medical office more than 40 years ago, during the early pioneering days of surgery of the heart. This saintly soul suffered much because of a failing heart. He pleaded for help, thinking that his condition resulted from a damaged but repairable valve in his heart.

Extensive evaluation revealed that he had two faulty valves. While one could be helped surgically, the other could not. Thus, an operation was not advised. He received this news with deep disappointment.

Subsequent visits ended with the same advice. Finally, in desperation, he spoke to me with considerable emotion: “Dr. Nelson, I have prayed for help and have been directed to you. The Lord will not reveal to me how to repair that second valve, but He can reveal it to you. Your mind is so prepared. If you will operate upon me, the Lord will make it known to you what to do. Please perform the operation that I need, and pray for the help that you need.”

His great faith had a profound effect upon me. How could I turn him away again? Following a fervent prayer together, I agreed to try. In preparing for that fateful day, I prayed over and over again, but still did not know what to do for his leaking tricuspid valve. Even as the operation commenced, my assistant asked, “What are you going to do for that?”

I said, “I do not know.”

We began the operation. After relieving the obstruction of the first valve, we exposed the second valve. We found it to be intact but so badly dilated that it could no longer function as it should. While examining this valve, a message was distinctly impressed upon my mind: Reduce the circumference of the ring. I announced that message to my assistant. “The valve tissue will be sufficient if we can effectively reduce the ring toward its normal size.”

But how? We could not apply a belt as one would use to tighten the waist of oversized trousers. We could not squeeze with a strap as one would cinch a saddle on a horse. Then a picture came vividly to my mind, showing how stitches could be placed—to make a pleat here and a tuck there—to accomplish the desired objective. I still remember that mental image—complete with dotted lines where sutures should be placed. The repair was completed as diagrammed in my mind. We tested the valve and found the leak to be reduced remarkably. My assistant said, “It’s a miracle.”

I responded, “It’s an answer to prayer.”

The patient’s recovery was rapid and his relief gratifying. Not only was he helped in a marvelous way, but surgical help for other people with similar problems had become a possibility. I take no credit. Praise goes to this faithful patriarch and to God, who answered our prayers. This faithful man lived for many more years and has since gone to his eternal glory.

\intalkheader{Inquire of the Lord}

When we pray, we should not presume to give counsel but should inquire of the Lord and hearken to His counsel. Joseph Smith’s first prayer ushered in the Restoration of the gospel. In 1833 he received the Word of Wisdom after asking the Lord for counsel. The revelation on the priesthood received in 1978 by President Spencer W. Kimball came after intense inquiry. Inspiration regarding the construction of smaller temples came after the pondering of President Gordon B. Hinckley.

\intalkheader{Answers to Prayers}

Not all of our prayers will be answered as we might wish. Occasionally the answer will be no. We should not be surprised. Loving mortal parents do not say yes to every request of their children.

At a recent extended family home evening, our grandchildren were having a wonderful time. A six-year-old grandson became very upset when his father said it was time to go home. So what did this dear boy do? He came to me and said, “Grandfather, may I have your permission to disobey my father?”

I said, “No, sweetheart. One of life’s great lessons is to learn that happiness comes through obedience. Go home with your family, and you will be happy.” Though disappointed, he dutifully obeyed.

We should pray in accord with the will of our Heavenly Father. He wants to test us, to strengthen us, and to help us achieve our full potential. When the Prophet Joseph Smith was held in Liberty Jail, he pled for relief. His prayers were answered with an explanation: “All these things shall give thee experience, and shall be for thy good.”

\intalkheader{Song of Prayer}

I have felt impressed to conclude this message on prayer with a prayer—presented as a hymn. The Lord has said that “the song of the righteous is a prayer unto me.” The music comes from our book of Hymns, for which I have written new words. With thanks to Craig Jessop, Mack Wilberg, and other dear friends in the Tabernacle Choir, we may hear that song of prayer. Brother Jessop, please:

\begin{verse}
We pray to Thee, our Heavenly Father,\\
With grateful hearts and fond emotion.\\
We thank Thee for our great Exemplar—\\
Thy beloved and atoning Son—\\
Who gave Himself as sacred ransom,\\
That we could live again with Thee!\\
Our joy is full, our song so gladsome;\\
Renew our faith and hope in Thee.
\end{verse}

In the name of Jesus Christ, amen.

For an arrangement of Elder Russell M. Nelson’s hymn text, see the inside back cover of this magazine.

\newpage
\settalk{Personal Priesthood Responsibility}
\setspeaker{Russell M. Nelson}
\settalkdate{October 2003}
\talkheading{Personal Priesthood Responsibility}{Russell M. Nelson}

My dear brethren of the priesthood, though we are from many nations, we are, as Paul said, of “one Lord, one faith, one baptism.” But the strength of the faith within each of us is developed individually, not as a group.

For example, think of the faith of a boy, about eight years of age, who was facing an emergency operation for acute appendicitis. As he lay on the operating table, he looked up at the surgeon and said, “Doctor, before you begin to operate, will you pray for me?”

The surgeon looked at the boy in amazement and said, “Why, I can’t pray for you.”

Then the little fellow said, “If you won’t pray for me, please wait while I pray for myself.” There on the operating table, the boy got on his knees, folded his hands, and began to pray. He said: “Heavenly Father, I am only a little orphan boy. I am awful sick, and these doctors are going to operate. Will you please help them that they will do it right? Heavenly Father, if you will make me well, I will be a good boy. Thank you for making me well.” He then lay on his back, looked up at the tear-filled eyes of the doctors and nurses, and said, “Now I am ready.”

His physical recovery was complete, and his spiritual power was developing. You brethren are older and have had the priesthood conferred upon you. Your priesthood quorums provide opportunities for friendship, service, and learning. But the responsibility to develop power in the priesthood is personal. Only as an individual can you develop a firm faith in God and a passion for personal prayer. Only as an individual can you keep the commandments of God. Only as an individual can you repent. Only as an individual can you qualify for the ordinances of salvation and exaltation. And when your wife is sealed to you, her power and potential will increase yours.

I belong to a wonderful priesthood quorum. We enjoy a precious brotherhood. We pray together; we serve together. We teach, love, and sustain one another. The Twelve come from different backgrounds—business, education, law, and science. But not one was called to serve because of that background. In fact, all men called to positions of priesthood responsibility are chosen because of who they are and who they can become.

Throughout life you will have a wide variety of duties and responsibilities. Many of these are temporary and will be relinquished upon your release. (You probably won’t object to your release from a call to pull weeds at the welfare farm.) But you never will be released from responsibilities related to your personal and family development.

When ordained to an office in the priesthood, you are granted authority. But power comes from exercising that authority in righteousness.

\intalkheader{Responsibility to the Lord}

From the President of the Church to the newest deacon, we are responsible to the Lord. We are to be true and faithful and live by every principle and doctrine that He has given to us. We cannot compromise a revelation or a commandment committed to our charge. He trusts us to “build up the kingdom of God, and to establish his righteousness.”

One day each of us will give an account to the Lord. This awareness was evident in a serious conversation I had years ago with a dear friend facing the end of his mortal life. I asked him if he was ready to die. I’ll never forget his answer. With courage and conviction, he said, “My life is ready for inspection.”

When the Prophet Joseph Smith faced death, he said, “I am going like a lamb to the slaughter; but I am calm as a summer’s morning; I have a conscience void of offense towards God, and towards all men.”

Now is the time to prepare for your own ultimate interview. You might ask yourself: “Do I pay tithing with a willing heart? Do I obey the Word of Wisdom? Is my language free from obscenities and swearing? Am I morally righteous? Am I truly grateful for the Atonement, which makes my resurrection a reality and eternal life a possibility? Do I honor temple covenants that seal loved ones to me forever?” If you can honestly say yes, you are developing power in the priesthood.

The gift of the Holy Ghost can add to that power. Scriptures tell of people who had received the Holy Ghost but did not know it. Don’t let that happen to you. Cultivate that gift and qualify for this promise from God: “Speak the thoughts that I shall put into your hearts, and you shall not be confounded before men; For it shall be given you in the very hour, yea, in the very moment, what ye shall say.”

\intalkheader{Personal Responsibility and Priesthood Power}

Priesthood authority has existed in many dispensations, such as those of Adam, Noah, Enoch, Abraham, Moses, the meridian of time, the Jaredites, the Nephites, and others. All previous dispensations were limited in time, as each ended in apostasy. They were also limited to small segments of planet Earth. In contrast, our dispensation—the dispensation of the fulness of times—will not be limited in time or place. Globally, it will host a whole, complete, and perfect union, welding together dispensations, keys, powers, and glories from the days of Adam even to the present time.

The Aaronic Priesthood was restored on May 15, 1829, by John the Baptist; the Melchizedek Priesthood was restored shortly thereafter by Peter, James, and John. Other heavenly messengers conveyed specific keys of the priesthood. Moroni held keys of the Book of Mormon. Moses brought keys of the gathering of Israel and the leading of the ten tribes. Elias conveyed keys of the restoration of all things, including the Abrahamic covenant. And Elijah conferred keys of the sealing authority.

You know something about keys. In your pocket there might be a key to your home or car. Priesthood keys, on the other hand, are intangible and invisible. They “switch on” the authority of the priesthood. Some keys even convey power to bind in heaven as well as on earth.

Joseph Smith conferred priesthood keys upon all of the Twelve. Those keys have been transferred to successive leaders. Today President Gordon B. Hinckley holds authority for every restored key held by “all those who have received a dispensation at any time from the beginning of the creation.”

With this doctrinal history in mind, it is clear that one cannot buy the priesthood. Scripture declares that “no man taketh this honour unto himself, but he that is called of God, as was Aaron.”

To bear the priesthood means you have a personal responsibility to magnify your calling. Let each opportunity to serve help to develop your power in the priesthood. In your personal grooming, follow the example of the living prophets. Doing so gives silent expression that you truly comprehend the importance of “the Holy Priesthood, after the Order of the Son of God.”

When you brethren have an opportunity to exercise the Melchizedek Priesthood, ponder what you are to do. When you lay hands upon the head of another, you are not offering a prayer, which of course requires no authority. You are authorized to set apart, to ordain, to bless, and to speak in the name of the Lord. Remember His promises: “Whomsoever you bless I will bless,” and “I will impart unto you of my Spirit, … and then shall ye know … all things … pertaining unto things of righteousness, in faith believing in me that you shall receive.”

To magnify your callings in the Aaronic Priesthood, you young men should shape your personal efforts toward five personal objectives to:

How can you remember those five objectives? It’s easy. Look at your hand. Let your pointer finger point to the scriptures. From them gain a better knowledge of the gospel of Jesus Christ, and then live in accord with His teachings. Let your middle finger remind you to be worthy of missionary service. Let your ring finger remind you of marriage, endowment, sealing, and blessings of the temple. Let your end finger remind you that pursuit of an education is a religious responsibility. Let your thumb go up, reminding you to uphold the standards of the Church and be worthy of your eternal companion. The realization of these five objectives will bless your lives.

You bearers of the Melchizedek Priesthood should qualify for the highest degree of celestial glory. “In order to obtain [it], a man must enter into this order of the priesthood [meaning the new and everlasting covenant of marriage]; And if he does not, he cannot obtain it.”

That covenant is honored when you honor your wife. A husband’s foremost priority should be the care of his wife. Be true to her. Don’t ever allow your eyes to gaze upon pornography or let your language be lewd. The very choices made by reason of agency limit one’s agency in the future. You cannot exercise agency and escape accountability and responsibility for each choice.

Never forget that “the rights of the priesthood are inseparably connected with the powers of heaven. … [This power] cannot be controlled nor handled [except] upon the principles of righteousness.” If we abuse that power to cover our sins, to gratify our pride, to pursue vain ambition, or to control others in any degree of unrighteousness, we lose both the authority and the power of the priesthood.

Brethren, serve with gentleness, long-suffering, kindness, meekness, love unfeigned, pure knowledge, and charity toward all. Then the “doctrine of the priesthood shall distil upon [your] soul as the dews from heaven.”

Please know of our love and gratitude for each of you. We thank you for your faith, your service, and your sustaining strength. May you, your loved ones, and your posterity be blessed by your righteous pursuit of power in the priesthood.

God lives. Jesus is the Christ. He directs His Church through His prophets and apostles. I so testify, in the name of Jesus Christ, amen.

\newpage
\settalk{Roots and Branches}
\setspeaker{Russell M. Nelson}
\settalkdate{April 2004}
\talkheading{Roots and Branches}{Russell M. Nelson}

Each time we experience security checks at an airport, we are asked to show photographic identification. We understand the need and comply, knowing that it is necessary and helpful. But I submit my photograph as evidence of my true identity somewhat apologetically. If someone were to examine my passport photo and say that it’s a good likeness, I would know it’s time to go home. But I feel apologetic for another reason. The photo shows nothing about my roots and branches. They are important parts of my identity. Could you tell much about a tree by looking at a photograph of only its trunk? No! Roots and branches of trees provide much more information. So it is with us both personally and with our religion.

\intalkheader{Personal Roots}

Personal roots are really important. Sister Nelson and I know a family that proudly display evidence of their ancestral roots with large paintings portrayed on the outside walls of their home. Beautiful artwork there depicts distinctive patterns of identity for both of their family lines.

When relatives gather around a new baby, one inevitably hears comments such as “She has red hair, just like her mother” or “He has a dimple in his chin, just like his father.”

Each of us has ancestral roots. Each man has received some genetic markers that are just like those of his father. Each woman has received some genetic markers that are just like those of her mother. In addition, each of us has received other genetic gifts that make us unique.

Because we have a spirit as well as a physical body, we also have spiritual roots that go way back. They shape our values, our beliefs, and our faith. Spiritual roots guide our commitment to the ideals and teachings of the Lord.

Children have a natural desire to emulate the examples of their parents. Generally, boys incline toward the attitudes and work of their fathers; girls aspire to live as their mothers do. And you parents, don’t be too surprised if, sometime along the way, your children become better than you.

Personal roots, physical and spiritual, merit gratitude. For my life, I am grateful to my Creator as well as to my dear parents and progenitors. I try to honor them by learning of them and serving them in the temple. Parents have a responsibility to share knowledge of their personal roots with their children and grandchildren. Learning their history together unifies a family.

\intalkheader{Religious Roots}

We also need to know the roots of our religion. The Church of Jesus Christ of Latter-day Saints, though officially organized in 1830, has been restored from roots that also go way back. Truths from previous dispensations have now been gathered, amplified, and clarified. For us as parents and teachers, we have an excellent teaching resource in the Articles of Faith. Written by the Prophet Joseph Smith, this document refers to many doctrines that undergird our religion. It mentions the Godhead, moral agency, the Fall of Adam, and the Atonement of Jesus Christ. It spells out the foundational principles and ordinances of faith, repentance, baptism, and the laying on of hands for the gift of the Holy Ghost. It addresses matters of priesthood authority and organization. It notes as sacred scripture the Holy Bible, the Book of Mormon, and an open canon of continuing revelation from God. And it proclaims the actuality of the gathering of Israel. What a treasure-house of truth is this precious document as we teach of our religious roots.

Other revealed doctrines at the root of our religion include the Creation, the Resurrection, the law of tithing, prayer, and the consummate blessings of the temple. As we teach of these doctrines, we realize how very firm is our foundation. As we apply these doctrines in our lives, the roots of our religion become part of our own spiritual strength.

Converts need to strengthen their religious roots. President Gordon B. Hinckley has taught that each convert needs a friend, a responsibility, and nurturing by the good word of God. With such roots to support them and their children, precious converts become pioneers for their own families to follow.

Unfortunately, some members of faithful families drift away because their own roots are weak. My heart aches when I learn of those who turn from the faith of their pioneer predecessors. One professionally acclaimed friend and gifted son of faithful ancestors has allowed one doctrinal doubt to dim his view of the fulness of the gospel and drive an ever-widening wedge between him and the temple. Another acquaintance, a sweet sister with illustrious pioneer progenitors, now politely states that she is not a “practicing member” of the Church.

Have these dear people become so fashionable that they have forgotten their roots? Have they forgotten what the Restoration really means and what it cost? Have they forgotten their pioneer heritage and their lineage as declared in patriarchal blessings? For a few fleeting favors now, would they forget and forfeit eternal life? Oblivious to the roots that have blessed them, they no longer enjoy the spiritual sparkle of Saints engaged in the work of Almighty God.

Their noble ancestors “were brought to the knowledge of the truth, … according to the spirit of revelation and of prophecy, and the power of God.” Their forebears “were converted unto the Lord [and] never did fall away.” How will those progenitors feel about the drift of their descendants? Their disappointment will likely turn to sorrow, for fruit detached from roots cannot long survive.

The Lord issued this solemn warning:

“After ye have been nourished by the good word of God … , will ye reject these words … of the prophets; and will ye reject all the words which have been spoken concerning Christ, … and the power of God, and the gift of the Holy Ghost, … and make a mock of the great plan of redemption, which [has] been laid for you?

“… The resurrection … will bring you to stand with shame and awful guilt before the bar of God.”

I plead with each of us to heed that sacred warning.

\intalkheader{Personal Branches}

Just as our roots determine to a significant degree who we are, our branches are also an important extension of our identity. Personal branches bear the fruit of our loins. Scriptures teach, “By their fruits ye shall know them.” Earlier in life Sister Nelson and I often met young people who said they felt like they knew us because they knew our children. Now we are greeted fondly by those who know us because they know our grandchildren.

\intalkheader{Religious Branches}

In much the same way, our religion is known by the fruit of its branches. Recently I met with government officials from a land far from here who were deeply impressed with the Church and its efforts throughout the world. They liked our teachings about the family and wanted copies of our proclamation to the world and guidebooks for family home evening. They wanted to know more about our welfare program and humanitarian help. We complied as we could and then shifted attention from what we do to why we do it. I explained with an analogy to a tree. “You are attracted by various fruits of our faith,” I said. “They are plentiful and powerful. But you cannot savor this fruit unless you know the tree that produces it. And you cannot understand the tree unless you comprehend its roots. With our religion, you cannot have the fruits without the roots.” This they understood.

Fruits from the branching tree of the gospel include “love, joy, peace, longsuffering, gentleness, goodness, [and] faith.” President Harold B. Lee once said: “Beautiful, luscious fruit does not grow unless the roots of the … tree have been planted in rich, fertile soil and unless due care is given to proper pruning, cultivation, and irrigation. So likewise the luscious fruits of virtue and chastity, honesty, temperance, integrity, and fidelity are not to be found growing in that individual whose life is not founded on a firm testimony of the truths of the gospel and of the life and the mission of the Lord Jesus Christ.”

The fruits of the gospel are delicious to those who obey the Lord. We pursue an education knowing that “the glory of God is intelligence.” The blessing of tithing comes by paying tithing. Rewards are reaped from the Word of Wisdom by obedience to it. We learn from experience borne of gospel living that prayer, honoring the Sabbath day, and partaking of the sacrament protect us from the bondage of sin. We shun pornography and immorality, knowing that the peace of personal purity can be ours only as we live according to the laws of the gospel.

The Lord gave this promise and commandment: “Now ye are clean through the word which I have spoken unto you. Abide in me, and I in you. … I am the vine, ye are the branches.” Simply summarized, life’s greatest blessings will come to us if our love of Jesus Christ is rooted deeply in our hearts.

\intalkheader{Testimony}

Personal identity is much more than a passport photograph. We also have roots and branches. Divinity is rooted in each of us. “We all are the work of [our Creator’s] hand.” We are eternal beings. In premortal realms, we brethren were foreordained for our priesthood responsibilities. Before the foundation of the world, women were prepared that they may bear children and glorify God.

We came to this mortal experience to acquire a body, to be tried and tested. We are to form families and be sealed in holy temples, with joy and loving relationships that endure eternally. To these everlasting truths we are personally rooted.

Branches of our families and of the gospel bear fruit to enrich our lives. God’s work and His glory—“to bring to pass the immortality and eternal life of man”—can become ours. We can dwell with Him and with our families forever. Those blessings will be granted to the faithful in His own way and time.

God lives. Jesus is the Christ. Joseph Smith is the revelator and prophet of this last dispensation. The Book of Mormon is true. The Church of Jesus Christ of Latter-day Saints is the Lord’s kingdom established once again upon the earth. President Gordon B. Hinckley is His living prophet. If rooted to these truths, the fruit of our branches will remain. I so testify in the name of Jesus Christ, amen.

\newpage
\settalk{Senior Missionaries and the Gospel}
\setspeaker{Russell M. Nelson}
\settalkdate{October 2004}
\talkheading{Senior Missionaries and the Gospel}{Russell M. Nelson}

Gratefully we welcome Elder Dieter F. Uchtdorf and Elder David A. Bednar to the Quorum of the Twelve Apostles. Prayerfully and unitedly we will serve the Lord Jesus Christ.

Assignments this year have taken me to many nations of the earth. In some of those countries the Church is relatively new. No matter where I go, I meet our missionaries. They are remarkably resilient and ever effective. They give visible and tangible evidence that the Church of Jesus Christ has been restored in its fulness. It was He who said, “Go ye into all the world, and preach the gospel to every creature.” This commandment throbs in the heart of every missionary who testifies of Jesus Christ and teaches His message.

When we think of missionaries, we generally picture in our minds young men with shirts and ties and young women dressed modestly. But along with them are marvelous senior missionaries who have answered the pleadings of prophets and apostles for more missionary couples.

I express gratitude for our senior missionaries. They are young in spirit, wise, and willing to work. They even tolerate remarks from their fun-filled children who might change President Spencer W. Kimball’s plea “Lengthen your stride” to “Hasten your shuffle.” These dear members are willing to serve and strengthen the lives of others. Even if these seniors don’t know the local language, their accomplishments are great and their spirit of sacrifice is precious.

\intalkheader{Examples of Senior Missionary Service}

For example, I think of Elder Lloyd Poelman and his wife, Sister Catherine Poelman. Parents of 9 grown children and grandparents of 20 grandchildren, they now serve in a remote part of Chile, working in a small branch. They make frequent visits among less-active members and with families recently converted to the Church. These visits provide opportunity for the Poelmans to read with those families and bear testimony of temple blessings. In their mission branches they have also taught people how to conduct music and play simplified versions of the hymns on small electronic keyboards. Elder and Sister Poelman recently wrote: “Baptism is only the first step in conversion. When the initial excitement subsides and the new converts continue facing the need to work long hours just to put bread on the table, they need others to help them who share the joy of the gospel. That is our specialty. Part of our work is preventive—staying close to new converts. Yet others who rarely attend meetings have not lost conviction and receive our messages gratefully. As we watch the changes brought about in the lives of those we visit, we feel blessed to be receiving constant tutoring and help from the Lord in this work and, at the same time, to know that our family members back home are vicariously sharing our calling and those special blessings.”

Such marvelous couples are engaged in the work of reclaiming souls who have previously made covenants to take upon themselves the name of Jesus Christ.

Other missionary couples render service in sacred temples of the Church. Elder Kenneth and Sister Barbara Willits, for instance, serve in the Accra Ghana Temple. They developed a special love for the people of Ghana while serving there as missionaries more than two decades earlier. They are energetic and enthusiastic converts of 50 years, with 3 children, 16 grandchildren, and 12 great-grandchildren. They perform the exalting ordinances of the temple. Brother Willits serves as a sealer. On several occasions they have been pleasantly surprised to meet members whom they had previously encountered during their first mission. Recently Elder Willits performed the sealing of a husband and wife whom they had taught in 1982, and to that couple Elder Willits sealed four of their deceased children. Elder and Sister Willits write: “Our willingness to leave our family and home is motivated by the temple covenants we have made, and our deepest desire is to become an eternal family. Our family is fully supportive as we serve, and they share in many blessings we have received. We are humbly grateful for the privilege of assisting others to receive their temple blessings.”

Courageous and caring couples like Elder and Sister Willits enable and enrich the work done in many of our temples across the earth. Some, such as the Accra Ghana Temple, are located where most local members had not had previous opportunities to attend a temple. Ordinances for those members are now enhanced by experienced couples who serve as temple missionaries. To them, we also express our heartfelt gratitude.

Earlier this year Elder Douglas L. Callister and I were in Kiev, capital city of Ukraine. We were there to create the first stake in the former Union of Soviet Socialist Republics. We were pleased to note that the Kiev Ukraine District was well prepared to become a stake—fully organized and ready to take its place among the stakes of Zion. There we also met with the missionaries, among whom were several stalwart senior couples. We listened attentively to their expressions.

We remember the account of Elder Rudi and Sister Eva Hegewald, who grew up in what was then known as East Germany. Speaking with a slight and sweet German accent, they recounted the difficult days of World War II and the subsequent Soviet occupation. They spoke of their many deprivations. Finding the Lord’s true Church and later immigrating to America were counted as treasured blessings. The ensuing years brought them five healthy children, along with spiritual and financial increase. They felt that serving a mission would be a good way for them to show gratitude to the Lord. They expressed a deep desire to serve in Eastern Europe. Their call came to serve in the Ukraine Kiev Mission. Elder and Sister Hegewald write: “Now, close to the end of our mission in the land of our former enemy, we are thankful for the opportunity to teach and love the Ukrainian people. As we have served the Lord, our souls have been healed and our family has become more united. We have had a truly remarkable and satisfying experience and have seen many small miracles.”

Notice that all three couples wrote of their blessings. Another couple tells of blessings that come from missionary service. They wrote: “Good people replaced our parenting functions better than we. … If a family problem has not yielded to prayer and fasting, a mission might be considered.”

No senior missionary finds it convenient to leave. Neither did Joseph or Brigham or John or Wilford. They had children and grandchildren too. They loved their families not one whit less, but they also loved the Lord and wanted to serve Him. Someday we may meet these stalwarts who helped to establish this dispensation. Then will we rejoice that we did not seek the shadows when a call to missionary service came from the prophet, even in the autumn years of our lives.

At general conference in October 1925, President Heber J. Grant issued a clarion call for “men of mature years and sound judgment, who have had experience in the preaching of the gospel, … to go forth and labor in the mission fields.”

That need persists. At the most recent training broadcast to priesthood leaders throughout the world, President Gordon B. Hinckley issued a similar call: “There is a constant need for more couple missionaries,” he said. “They perform wonderful service throughout the world. You [leaders] need not wait for the couples to volunteer. The sacrifices associated with serving the Lord full-time will abundantly bless the couples, their families, and the people they serve.”

\intalkheader{Qualifications for the Work}

Bishops also need to heed that prophetic call and ask such members if they could serve. Opportunities for senior missionaries are varied and vast. Their calls to serve are officially made after prayerful consideration has been given to their occupational background, language experience, and personal capabilities. Of all qualifications to serve, a desire to serve may be the most important. The Lord has declared:

“O ye that embark in the service of God, see that ye serve him with all your heart, might, mind and strength, that ye may stand blameless before God at the last day.

“Therefore, if ye have desires to serve God ye are called to the work.”

Many humble Latter-day Saints fear that they are not qualified for missionary labors. But to such a prospective missionary, the Lord has given this assurance: “Faith, hope, charity and love, with an eye single to the glory of God, qualify him for the work.”

\intalkheader{Limitations Imposed by Age and Health}

As I extol the work of senior missionaries, I realize that there are many more who would like to serve but are not able to do so. Limitations imposed by age or by poor health deserve realistic appraisal, as do the important needs of family members. When desire burns within yet such limitations exist, you can extend your service through others. They can be your arms and legs, and you can provide needed funds. Still others can contribute time and talents as live-at-home missionaries. Each will be pleasing to the Lord, and each will receive His praise.

\intalkheader{The Gospel}

All of us may preach the gospel by precept and example. The word gospel means “good news.” The good news is the Lord Jesus Christ and His message of salvation. Jesus equated the gospel with both His mission and with His ministry in mortality. In His mission statement, Jesus said:

“This is the gospel which I have given unto you—that I came into the world to do the will of my Father, because my Father sent me.

“And my Father sent me that I might be lifted up upon the cross.”

The Savior’s mortal mission we know as the Atonement.

The Savior’s mortal ministry includes everything else that He did—His teachings, expressions of love, attention to ordinances, patterns of prayer, perseverance, and more. He lived to be our Exemplar, which He also equated to the gospel in His ministerial statement. “This is my gospel,” He said, “for the works which ye have seen me do that shall ye also do.” Thus, faith; repentance; baptism by water, fire, and of the Holy Ghost; the gathering of the elect; and enduring to the end are all part of the gospel. All of us can emulate the Lord’s example, regardless of age, status, or location.

As one among the “special witnesses of the name of Christ in all the world,” I declare that He is the Son of the living God, our atoning Savior and Redeemer. This is His Church, restored in these latter days to fulfill its divine destiny. His prophet today is President Gordon B. Hinckley. I so testify in the name of Jesus Christ, amen.

\newpage
\settalk{Now Is the Time to Prepare}
\setspeaker{Russell M. Nelson}
\settalkdate{April 2005}
\talkheading{Now Is the Time to Prepare}{Russell M. Nelson}

My dear brethren and sisters, since our last general conference, my sweetheart—my beloved wife for 59 years—passed away. While I was at home on a rare Saturday with no assignment, we had worked together. She had washed our clothing. I had helped to carry it, fold it, and put it in place. Then while we were sitting on the sofa, holding hands, enjoying a program on television, my precious Dantzel slipped peacefully into eternity. Her passing came suddenly and unexpectedly. Just four days earlier, our doctor’s report at a routine checkup indicated that her laboratory tests were good. After my efforts to revive her proved fruitless, feelings of shock and sorrow overwhelmed me. My closest friend, angel mother of our 10 children, grandmother of our 56 grandchildren, had been taken from us.

Dantzel was not only a loved and loving companion. She was a teacher: by her noble example, she taught faith, virtue, obedience, and mercy. She taught me how to listen and to love. Because of her, I know all the blessings that can come to a husband, father, and grandfather.

With deep gratitude I acknowledge the tremendous outpouring of love from dear friends across the world. Countless letters, calls, cards, and other messages have been sent. All tributes expressed loving admiration for her and sympathy for us whom she left behind. Those messages came in such large numbers that we, regretfully, were unable to respond to all of them individually. May I thank each and all for your great kindness toward us. Thank you so very, very much. Your expressions have brought much comfort through this time of heartache for our family. We really love dear Dantzel! We miss her!

From her sudden departure we can learn a very important lesson: now is the time to prepare to meet God. Tomorrow may be too late. Prophets through the ages have so declared: “This life is the time for men to prepare to meet God. … Do not procrastinate the day of your repentance.”

\intalkheader{The Need to Prepare Now}

Yet many do procrastinate. A prophet warns us: “Ye cannot say … that I will repent, that I will return to my God. Nay, ye cannot say this; for that same spirit which doth possess your bodies at the time that ye go out of this life, that same spirit will … possess your body in that eternal world.” Another prophet adds, “He that is filthy shall be filthy still; and he that is righteous shall be righteous still.”

Great is the knowledge that “whatever principle of intelligence we attain … in this life, it will rise with us in the resurrection.” From the Prophet Joseph Smith we also learn that “God has … a time … appointed … when He will bring all His subjects, who have obeyed His voice and kept His commandments, into His celestial rest. This rest is of such perfection and glory, that man has need of a preparation before he can, according to the laws of that kingdom, enter it and enjoy its blessings. … God has given certain laws to the human family, which, if observed, are sufficient to prepare them to inherit this rest.” Sister Nelson was so prepared!

That glorious goal seems mighty distant if one is discouraged by worldly trouble and gloom. I remember when a friend having a difficult day exclaimed, “Oh, why was I ever born?” God’s plan answers his question. We came into this life to acquire a physical body. We may fall in love and be married. We may have children and experience the trials of mortal life. (Please forgive me for mentioning children and the trials of life in the same breath. I sense that they are both part of our growing process.) The Church was restored—the earth was created—so that those families could be sealed in holy temples. Otherwise, the whole earth would be “utterly wasted.”

We came to be tried, to be tested, and to choose. Our decisions determine our destiny. We are “free to choose liberty and eternal life, through the great Mediator … , or to choose captivity and death.” Those who choose the Lord’s way will likely endure persecution. But their reward is certain. Those who prove faithful “shall inherit the kingdom of God, … and their joy shall be full forever.” Sister Nelson has earned that reward. What comfort that brings to me and our family!

Trials and tests apply to rich and poor alike. Years ago I was asked to perform an operation upon a very wealthy man. A surgical biopsy confirmed that he had an advanced cancer that had spread throughout his body. As I reported this news, his immediate response was to rely upon his wealth. He would go anywhere or do anything to treat his condition. He thought he could buy his way back to health. But he soon passed away. Someone asked, “How much wealth did he leave?” The answer, of course, was, “All of it!”

His priorities were set upon things of the world. His ladder of success had been leaning against the wrong wall. I think of him when I read this scripture: “Behold, your days of probation are past; ye have procrastinated the day of your salvation until it is … too late.”

In radiant contrast, Sister Nelson prepared throughout her life for the time when she would return to God. She lived each day as though it were her last. She cherished every hour, knowing that time on earth is precious.

Some people live as if there were no day of reckoning. Others waste today’s time with a disabling fear of tomorrow or a paralyzing preoccupation over mistakes of yesterday. Each of us might well heed the words of a poet, as posted on a sundial:

\begin{verse}
The shadow by my finger cast\\
Divides the future from the past:\\
Before it, sleeps the unborn hour,\\
In darkness, and beyond thy power:\\
Behind its unreturning line,\\
The vanished hour, no longer thine:\\
One hour alone is in thy hands,—\\
The NOW on which the shadow stands.
\end{verse}

\intalkheader{How Do We Prepare?}

Now is the time. But how do we prepare? Begin with repentance! Scripture declares, “If ye have sought to do wickedly in the days of your probation, then ye are found unclean before the judgment-seat of God; … no unclean thing can dwell with God.” He gave this simple rule: “Except ye abide my law ye cannot attain to this glory.”

Now is the time to show reverent respect for one’s physical body. It serves as the tabernacle for one’s spirit throughout all eternity. Physical appetites are to be controlled by the will of one’s spirit. We are to “deny [our]selves of all ungodliness.” We are to “forsake all evil and cleave unto all good, [and] live by every word which proceedeth forth out of the mouth of God.”

Because of frequent and frightening calamities in the world, some people doubt the existence of God. But, in fact, He is trying to help us. He revealed these words: “How oft have I called upon you by the mouth of my servants, and by the ministering of angels, and by mine own voice, and by the voice of thunderings, and … tempests, … earthquakes, … great hailstorms, … famines and pestilences of every kind, … and would have saved you with an everlasting salvation, but ye would not!”

His hope for us is eternal life. We qualify for it by obedience to covenants and ordinances of the temple—for ourselves, our families, and our ancestors. We cannot be made perfect without them. We cannot wish our way into the presence of God. We are to obey the laws upon which those blessings are predicated.

God’s plan is fair. Even those “who have died without a knowledge of this gospel, who would have received it if they had been permitted to tarry, shall be heirs of the celestial kingdom of God.” His plan is also merciful. He “will judge all men according to their works, according to the desire of their hearts.”

Now is the time to enroll our names among the people of God. This we do by paying tithing. He tithes His people to bless them. Sister Nelson taught that lesson to our family over and over again.

Now is the time to align our goals with God’s goals. His work and His glory—“to bring to pass the immortality and eternal life of man”—can become ours. Of temple marriage the Savior declared, “If a man marry a wife by my word, which is my law, and by the new and everlasting covenant, … [they] shall inherit thrones, kingdoms, principalities, and powers, dominions, … exaltation and glory in all things.” We are to emulate the example of the Lord, to love as He did, to pray as He did, and to endure to the end as He did.

\intalkheader{Importance of Death in God’s Eternal Plan}

Death is a necessary component of our eternal existence. No one knows when it will come, but it is essential to God’s great plan of happiness. Thanks to the Atonement of the Lord, eventual resurrection is a reality and eternal life is a possibility for all humankind. That possibility becomes a reality as we obey God’s law. He said, “Except ye shall keep my commandments, … ye shall in no case enter into the kingdom of heaven.” One day we will be judged by the Lord and go to our own mansion prepared in our Father’s heavenly house. Celestial glory awaits those who have been faithful to God’s gentle commands.

Brothers and sisters, we live to die and we die to live—in another realm. If we are well prepared, death brings no terror. From an eternal perspective, death is premature only for those who are not prepared to meet God.

Now is the time to prepare. Then, when death comes, we can move toward the celestial glory that Heavenly Father has prepared for His faithful children. Meanwhile, for sorrowing loved ones left behind—such as our family and me—the sting of death is soothed by a steadfast faith in Christ, a perfect brightness of hope, a love of God and of all men, and a deep desire to serve them. That faith, that hope, that love will qualify us to come into God’s holy presence and, with our eternal companions and families, dwell with Him forever. Of this I testify in the name of Jesus Christ, amen.

\newpage
\settalk{Jesus Christ—the Master Healer}
\setspeaker{Russell M. Nelson}
\settalkdate{October 2005}
\talkheading{Jesus Christ—the Master Healer}{Russell M. Nelson}

My beloved brothers and sisters, I extend love and greetings to each of you. From the Brethren, I express gratitude for your goodness, for your many generous acts of kindness, for your prayers and sustaining influence in our lives. Our challenges are like yours. We are all subject to sorrow and suffering, to disease and death. Through times good and bad, the Lord expects each of us to endure to the end. As we all go forward together in His sacred work, the Brethren realize the importance of your thoughtful consideration, so lovingly offered and gratefully received. We love you and pray for you, as you pray for us.

I express special gratitude to the Lord Jesus Christ. I am thankful for His loving-kindness and for His open invitation to come unto Him. I marvel at His matchless power to heal. I testify of Jesus Christ as the Master Healer. It is but one of many attributes that characterize His incomparable life.

Jesus is the Christ, the Messiah, the Son of God, the Creator, the great Jehovah, the promised Immanuel, our atoning Savior and Redeemer, our Advocate with the Father, our great Exemplar. And one day we will stand before Him as our just and merciful Judge.

\intalkheader{Miracles of Healing}

As the Master Healer, Jesus directed His friends to “go … and tell … what things ye have seen and heard; how that the blind see, the lame walk, the lepers are cleansed, the deaf hear, [and] the dead are raised.”

The books of Matthew, Mark, Luke, and John repeatedly report that Jesus went about preaching the gospel and healing all types of sickness.

When the risen Redeemer appeared to the people of ancient America, He mercifully invited those “afflicted in any manner” to come unto Him and be healed.

Marvelously, His divine authority to heal the sick was conferred upon worthy priesthood bearers in earlier dispensations and again in these latter days, when His gospel has been restored in its fulness.

\intalkheader{Influence of Prayer on Healing}

We can also access His healing power through prayer. I’ll never forget an experience that Sister Nelson and I had about three decades ago with President Spencer W. Kimball and his beloved Camilla. We were in Hamilton, New Zealand, for a large conference with the Saints. I was not a General Authority at that time. I had been invited to participate in this and similar meetings in other Pacific Islands while serving as general president of the Sunday School. And as a doctor of medicine, I had attended President and Sister Kimball for many years. I knew each of them very well—inside and out.

A Saturday evening cultural program had been prepared for this conference by local youth of the Church. Unfortunately, President and Sister Kimball both became very ill, each with a high fever. After receiving priesthood blessings, they rested at the nearby home of the president of the New Zealand Temple. President Kimball asked his counselor, President N. Eldon Tanner, to preside at the cultural event and to excuse President and Sister Kimball.

Sister Nelson went with President and Sister Tanner and other leaders to the event while President Kimball’s secretary, Brother D. Arthur Haycock, and I watched over our feverish friends.

While President Kimball was sleeping, I was quietly reading in his room. Suddenly President Kimball was awakened. He asked, “Brother Nelson, what time was this evening’s program to begin?”

“At seven o’clock, President Kimball.”

“What time is it now?”

“It’s almost seven,” I replied.

President Kimball quickly said, “Tell Sister Kimball we are going!”

I checked President Kimball’s temperature. It was normal! I took Sister Kimball’s temperature. It was also normal!

They quickly dressed and got into an automobile. We were driven to the stadium of the Church College of New Zealand. As the car entered the arena, there was a very loud shout that erupted spontaneously. It was most unusual! After we took our seats, I asked Sister Nelson about that sudden sound. She said that when President Tanner began the meeting, he dutifully excused President and Sister Kimball because of illness. Then one of the young New Zealanders was called upon to pray.

With great faith, he gave what Sister Nelson described as a rather lengthy but powerful prayer. He so prayed: “We are 3,000 New Zealand youth. We are assembled here, having prepared for six months to sing and dance for Thy prophet. Wilt Thou heal him and deliver him here!” After the “amen” was pronounced, the car carrying President and Sister Kimball entered the stadium. They were identified immediately, and instantly everyone shouted for joy!

I had witnessed the healing power of the Lord! I had also witnessed revelation as received and responded to by His living prophet!

I recognize that on occasion some of our most fervent prayers may seem to go unanswered. We wonder, “Why?” I know that feeling! I know the fears and tears of such moments. But I also know that our prayers are never ignored. Our faith is never unappreciated. I know that an all-wise Heavenly Father’s perspective is much broader than is ours. While we know of our mortal problems and pain, He knows of our immortal progress and potential. If we pray to know His will and submit ourselves to it with patience and courage, heavenly healing can take place in His own way and time.

\intalkheader{Steps toward Healing}

Afflictions can come from spiritual as well as physical causes. Alma the Younger remembered that his sin was so painful that he wished to “become extinct both soul and body, that [he] might not be brought to stand in the presence of … God, to be judged of [his] deeds.” At such times, how can we be healed by Him?

We can more fully repent! We can become more fully converted! Then the “Son of Righteousness” can more fully bless us by His healing hand.

Early in His mortal ministry, Jesus announced that He had been sent “to heal the brokenhearted.” Wherever He taught them, His pattern was consistent. As I quote His words spoken at four different times and locations, note the pattern.

The sequence of His pattern is significant. Faith, repentance, baptism, a testimony, and enduring conversion lead to the healing power of the Lord. Baptism is a covenant act—a sign of a commitment and a promise. Testimony develops when the Holy Ghost gives conviction to the earnest seeker of the truth. True testimony fosters faith; it promotes repentance and obedience to God’s commandments. Testimony engenders enthusiasm to serve God and fellow human beings. Conversion means “to turn with.” Conversion is a turning from the ways of the world to, and staying with, the ways of the Lord. Conversion includes repentance and obedience. Conversion brings a mighty change of heart. Thus, a true convert is “born again,” walking with a newness of life.

As true converts, we are motivated to do what the Lord wants us to do and to be who He wants us to be. The remission of sins, which brings divine forgiveness, heals the spirit.

How do we know if we are truly converted? Self-examination tests are available in the scriptures. One measures the degree of conversion prerequisite to baptism. Another measures our willingness to serve others. To His disciple Peter, the Lord said, “I have prayed for thee, that thy faith fail not: and when thou art converted, strengthen thy brethren.” Willingness to serve and strengthen others stands as a symbol of one’s readiness to be healed.

\intalkheader{Magnitude of His Healing}

John the Beloved declared, “Behold the Lamb of God, who taketh away the sin of the world!” What power! Only the Master Healer could take away the sin of the world. Our debt to Him is incalculably great.

Well do I remember an experience while speaking to a group of missionaries. After I had invited questions, one elder stood. With tears in his eyes he asked, “Why did Jesus have to suffer so much?” I asked the elder to open his book of hymns and recite words from “How Great Thou Art.” He read:

\begin{verse}
And when I think that God, his Son not sparing,\\
Sent him to die, I scarce can take it in,\\
That on the cross, my burden gladly bearing,\\
He bled and died to take away my sin.
\end{verse}

Then I asked this elder to read from “Reverently and Meekly Now.” These words are particularly poignant because they are written as the Lord would express His own answer to the very question that had been asked:

\begin{verse}
Think of me, thou ransomed one;\\
Think what I for thee have done.\\
With my blood that dripped like rain,\\
Sweat in agony of pain,\\
With my body on the tree\\
I have ransomed even thee. …
\end{verse}

Jesus suffered deeply because He loves us deeply! He wants us to repent and be converted so that He can fully heal us.

When sore trials come upon us, it’s time to deepen our faith in God, to work hard, and to serve others. Then He will heal our broken hearts. He will bestow upon us personal peace and comfort. Those great gifts will not be destroyed, even by death.

\intalkheader{Resurrection—the Consummate Act of Healing}

The gift of resurrection is the Lord’s consummate act of healing. Thanks to Him, each body will be restored to its proper and perfect frame. Thanks to Him, no condition is hopeless. Thanks to Him, brighter days are ahead, both here and hereafter. Real joy awaits each of us—on the other side of sorrow.

I testify that God lives, that Jesus is the Christ—the Master Healer—in the sacred name of Jesus Christ, amen.

\newpage
\settalk{Nurturing Marriage}
\setspeaker{Russell M. Nelson}
\settalkdate{April 2006}
\talkheading{Nurturing Marriage}{Russell M. Nelson}

My beloved brethren and sisters, thank you for your love of the Lord and His gospel. Wherever you live, your righteous lives provide good examples in these days of decaying morals and disintegrating marriages.

As we Brethren travel about the world, sometimes we see worrisome scenes. On a recent flight, I sat behind a husband and wife. She obviously loved her husband. As she stroked the back of his neck I could see her wedding ring. She would nestle close to him and rest her head upon his shoulder, seeking his companionship.

In contrast, he seemed totally oblivious to her presence. He was focused solely upon an electronic game player. During the entire flight, his attention was riveted upon that device. Not once did he look at her, speak to her, or acknowledge her yearning for affection.

His inattention made me feel like shouting: “Open your eyes, man! Can’t you see? Pay attention! Your wife loves you! She needs you!”

I don’t know more about them. I haven’t seen them since. Perhaps I was alarmed unduly. And very possibly, if this man knew of my concern for them, he might feel sorry for me in not knowing how to use such an exciting toy.

But these things I do know: I know “that marriage between a man and a woman is ordained of God and that the family is central to the Creator’s plan for the eternal destiny of His children.” I know that the earth was created and that the Lord’s Church was restored so that families could be sealed and exalted as eternal entities. And I know that one of Satan’s cunning methods of undermining the work of the Lord is to attack the sacred institutions of marriage and the family.

Marriage brings greater possibilities for happiness than does any other human relationship. Yet some married couples fall short of their full potential. They let their romance become rusty, take each other for granted, allow other interests or clouds of neglect to obscure the vision of what their marriage really could be. Marriages would be happier if nurtured more carefully.

I realize that many mature members of the Church are not married. Through no failing of their own, they deal with the trials of life alone. Be we all reminded that in the Lord’s own way and time, no blessings will be withheld from His faithful Saints. For those who are now or will be married, I suggest two steps you can take to have a more joyful marriage.

\intalkheader{I. Doctrinal Foundation}

The first step is to comprehend the doctrinal foundation for marriage. The Lord declared that marriage is the legal wedding of one man and one woman: “Marriage is ordained of God unto man.

“Wherefore, it is lawful that he should have one wife, and they twain shall be one flesh, and all this that the earth might answer the end of its creation.”

Worldly trends to define marriage in some other way would sadly serve to destroy the institution of marriage. Such schemes are contrary to the plan of God.

It was He who said, “For this cause shall a man leave father and mother, and shall cleave to his wife: and they twain shall be one flesh.”

Scripture further reaffirms that “the man [is not] without the woman, neither the woman without the man, in the Lord.”

Marriage is the foundry for social order, the fountain of virtue, and the foundation for eternal exaltation. Marriage has been divinely designated as an eternal and everlasting covenant. Marriage is sanctified when it is cherished and honored in holiness. That union is not merely between husband and wife; it embraces a partnership with God. “Husband and wife have a solemn responsibility to love and care for each other.” Children born of that marital union are “an heritage of the Lord.” Marriage is but the beginning bud of family life; parenthood is its flower. And that bouquet becomes even more beautiful when graced with grandchildren. Families may become as eternal as the kingdom of God itself.

Marriage is both a commandment and an exalting principle of the gospel. Because it is ordained of God, the intimate physical expressions of married love are sacred. Yet all too commonly, these divine gifts are desecrated. If a couple allows lewd language or pornography to corrupt their intimacy, they offend their Creator while they degrade and diminish their own divine gifts. True happiness is predicated upon personal purity. Scripture commands: “Be ye clean.” Marriage should ever be a covenant to lift husbands and wives to exaltation in celestial glory.

Marriage was intended by the Lord to endure beyond physical death. His plan offers eternal perpetuation of the family in the kingdom of God. His plan provides temples and opportunities to officiate therein for the living and the dead. A marriage sealed there launches a husband and wife into that grand order of unity so necessary to the perfection of God’s work.

Doctrines related to marriage include individual agency and accountability. All of us are accountable for our choices. Couples blessed with children are accountable to God for the care they give to their children.

As I meet with priesthood leaders, I often ask about the priorities of their various responsibilities. Usually they mention their important Church duties to which they have been called. Too few remember their responsibilities at home. Yet priesthood offices, keys, callings, and quorums are meant to exalt families. Priesthood authority has been restored so that families can be sealed eternally. So brethren, your foremost priesthood duty is to nurture your marriage—to care for, respect, honor, and love your wife. Be a blessing to her and your children.

\intalkheader{II. Strengthening Marriage}

With these doctrinal underpinnings in mind, let us consider the second step—specific actions that would strengthen a marriage. I will offer sample suggestions and invite each couple privately to ponder them and adapt them as needed to their own particular circumstances.

My suggestions use three action verbs: to appreciate, to communicate, and to contemplate.

To appreciate—to say “I love you” and “thank you”—is not difficult. But these expressions of love and appreciation do more than acknowledge a kind thought or deed. They are signs of sweet civility. As grateful partners look for the good in each other and sincerely pay compliments to one another, wives and husbands will strive to become the persons described in those compliments.

Suggestion number two—to communicate well with your spouse—is also important. Good communication includes taking time to plan together. Couples need private time to observe, to talk, and really listen to each other. They need to cooperate—helping each other as equal partners. They need to nurture their spiritual as well as physical intimacy. They should strive to elevate and motivate each other. Marital unity is sustained when goals are mutually understood. Good communication is also enhanced by prayer. To pray with specific mention of a spouse’s good deed (or need) nurtures a marriage.

My third suggestion is to contemplate. This word has deep meaning. It comes from Latin roots: con, meaning “with,” and templum, meaning “a space or place to meditate.” It is the root from which the word temple comes. If couples contemplate often—with each other in the temple—sacred covenants will be better remembered and kept. Frequent participation in temple service and regular family scripture study nourish a marriage and strengthen faith within a family. Contemplation allows one to anticipate and to resonate (or be in tune) with each other and with the Lord. Contemplation will nurture both a marriage and God’s kingdom. The Master said, “Seek not the things of this world but seek ye first to build up the kingdom of God, and to establish his righteousness, and all these things shall be added unto you.”

I invite each marital partner to consider these suggestions and then determine specific goals to nurture your own relationship. Begin with sincere desire. Identify those actions needed to bless your spiritual unity and purpose. Above all, do not be selfish! Generate a spirit of selflessness and generosity. Celebrate and commemorate each day together as a treasured gift from heaven.

President Harold B. Lee said “that the most important of the Lord’s work you and I will ever do will be within the walls of our own homes.” And President David O. McKay declared, “No other success can compensate for failure in the home.”

When you as husband and wife recognize the divine design in your union—when you feel deeply that God has brought you to each other—your vision will be expanded and your understanding enhanced. Such feelings are expressed in words of a song that has long been a favorite of mine:

\begin{verse}
Because you come to me with naught save love,\\
And hold my hand and lift mine eyes above,\\
A wider world of hope and joy I see,\\
Because you come to me.
\end{verse}

That each marriage may be so nurtured is my prayer, in the name of Jesus Christ, amen.

\newpage
\settalk{The Gathering of Scattered Israel}
\setspeaker{Russell M. Nelson}
\settalkdate{October 2006}
\talkheading{The Gathering of Scattered Israel}{Russell M. Nelson}

My beloved brothers and sisters, thank you for your faith, your devotion, and your love. We share an enormous responsibility to be who the Lord wants us to be and to do what He wants us to do. We are part of a great movement—the gathering of scattered Israel. I speak of this doctrine today because of its unique importance in God’s eternal plan.

\intalkheader{Abrahamic Covenant}

Anciently, the Lord blessed Father Abraham with a promise to make his posterity a chosen people. References to this covenant occur throughout the scriptures. Included were promises that the Son of God would come through Abraham’s lineage, that certain lands would be inherited, that nations and kindreds of the earth would be blessed through his seed, and more. While some aspects of that covenant have already been fulfilled, the Book of Mormon teaches that this Abrahamic covenant will be fulfilled only in these latter days! It also emphasizes that we are among the covenant people of the Lord. Ours is the privilege to participate personally in the fulfillment of these promises. What an exciting time to live!

\intalkheader{Israel Became Scattered}

As descendants of Abraham, the tribes of ancient Israel had access to priesthood authority and blessings of the gospel, but eventually the people rebelled. They killed the prophets and were punished by the Lord. Ten tribes were carried captive into Assyria. From there they became lost to the records of mankind. (Obviously, the ten tribes are not lost to the Lord.) Two remaining tribes continued a short time and then, because of their rebellion, were taken captive into Babylon. When they returned, they were favored of the Lord, but again they honored Him not. They rejected and vilified Him. A loving but grieving Father vowed, “I will scatter you among the heathen,” and that He did—into all nations.

\intalkheader{Israel to Be Gathered}

God’s promise for the gathering of scattered Israel was equally emphatic. Isaiah, for example, foresaw that in the latter days the Lord would send “swift messengers” to these people who were so “scattered and peeled.”

This promise of the gathering, woven all through the fabric of the scriptures, will be fulfilled just as surely as were the prophecies of the scattering of Israel.

\intalkheader{The Church of Jesus Christ in the Meridian of Time and the Apostasy}

Prior to His Crucifixion, the Lord Jesus Christ had established His Church. It included apostles, prophets, seventies, teachers, and so forth. And the Master sent His disciples into the world to preach His gospel.

After a time the Church as established by the Lord fell into spiritual decay. His teachings were altered; His ordinances were changed. The Great Apostasy came as had been foretold by Paul, who knew that the Lord would not come again “except there come a falling away first.”

This Great Apostasy followed the pattern that had ended each previous dispensation. The very first was in the time of Adam. Then came dispensations of Enoch, Noah, Abraham, Moses, and others. Each prophet had a divine commission to teach of the divinity and the doctrine of the Lord Jesus Christ. In each age these teachings were meant to help the people. But their disobedience resulted in apostasy. Thus, all previous dispensations were limited in time and location. They were limited in time because each ended in apostasy. They were limited in location to a relatively small segment of planet earth.

\intalkheader{The Restoration of All Things}

Thus a complete restoration was required. God the Father and Jesus Christ called upon the Prophet Joseph Smith to be the prophet of this dispensation. All divine powers of previous dispensations were to be restored through him. This dispensation of the fulness of times would not be limited in time or in location. It would not end in apostasy, and it would fill the world.

\intalkheader{The Gathering of Israel—an Integral Part of the Restoration of All Things}

As prophesied by Peter and Paul, all things were to be restored in this dispensation. Therefore, there must come, as part of that restoration, the long-awaited gathering of scattered Israel. It is a necessary prelude to the Second Coming of the Lord.

This doctrine of the gathering is one of the important teachings of The Church of Jesus Christ of Latter-day Saints. The Lord has declared, “I give unto you a sign … that I shall gather in, from their long dispersion, my people, O house of Israel, and shall establish again among them my Zion.” The coming forth of the Book of Mormon is a sign to the entire world that the Lord has commenced to gather Israel and fulfill covenants He made to Abraham, Isaac, and Jacob. We not only teach this doctrine, but we participate in it. We do so as we help to gather the elect of the Lord on both sides of the veil.

The Book of Mormon is central to this work. It declares the doctrine of the gathering. It causes people to learn about Jesus Christ, to believe His gospel, and to join His Church. In fact, if there were no Book of Mormon, the promised gathering of Israel would not occur.

To us the honored name of Abraham is important. It is mentioned in more verses of scriptures of the Restoration than in all verses of the Bible. Abraham is linked to all members of The Church of Jesus Christ of Latter-day Saints. The Lord reaffirmed the Abrahamic covenant in our day through the Prophet Joseph Smith. In the temple we receive our ultimate blessings, as the seed of Abraham, Isaac, and Jacob.

\intalkheader{The Dispensation of the Fulness of Times}

This dispensation of the fulness of times was foreseen by God as the time to gather, both in heaven and on earth. Peter knew that after a period of apostasy, a restoration would come. He, who had been with the Lord on the Mount of Transfiguration, declared:

“Repent ye therefore, and be converted, that your sins may be blotted out, when the times of refreshing shall come from the presence of the Lord; …

“Whom the heaven must receive until the times of restitution of all things, which God hath spoken by the mouth of all his holy prophets since the world began.”

In modern times the Apostles Peter, James, and John were sent by the Lord with “the keys of [His] kingdom, and a dispensation of the gospel for the last times; and for the fulness of times,” in which He would “gather together in one all things, both which are in heaven, and which are on earth.”

In the year 1830 the Prophet Joseph Smith learned of a heavenly messenger named Elias, who possessed keys to bring to pass “the restoration of all things.”

Six years later the Kirtland Temple was dedicated. After the Lord accepted that holy house, heavenly messengers came with priesthood keys. Moses appeared “and committed … the keys of the gathering of Israel from the four parts of the earth, and the leading of the ten tribes from the land of the north.

“After this, Elias appeared, and committed the dispensation of the gospel of Abraham, saying that in us and our seed all generations after us should be blessed.”

Then Elijah the prophet came and proclaimed, “Behold, the time has fully come, which was spoken of by the mouth of Malachi—testifying that he [Elijah] should be sent, before the great and dreadful day of the Lord come—to turn the hearts of the fathers to the children, and the children to the fathers, lest the whole earth be smitten with a curse.”

These events occurred on April 3, 1836, and thus fulfilled Malachi’s prophecy. Sacred keys of this dispensation were restored.

\intalkheader{Gathering of Souls on the Other Side of the Veil}

Mercifully, the invitation to “come unto Christ” can also be extended to those who died without a knowledge of the gospel. Part of their preparation requires earthly efforts of others. We gather pedigree charts, create family group sheets, and do temple work vicariously to gather individuals unto the Lord and into their families.

\intalkheader{To Participate in the Gathering: A Commitment by Covenant}

Here on earth, missionary work is crucial to the gathering of Israel. The gospel was to be taken first to the “lost sheep of the house of Israel.” Consequently, servants of the Lord have gone forth proclaiming the Restoration. In many nations our missionaries have searched for those of scattered Israel; they have hunted for them “out of the holes of the rocks”; and they have fished for them as in ancient days.

The choice to come unto Christ is not a matter of physical location; it is a matter of individual commitment. People can be “brought to the knowledge of the Lord” without leaving their homelands. True, in the early days of the Church, conversion often meant emigration as well. But now the gathering takes place in each nation. The Lord has decreed the establishment of Zion in each realm where He has given His Saints their birth and nationality. Scripture foretells that the people “shall be gathered home to the lands of their inheritance, and shall be established in all their lands of promise.” “Every nation is the gathering place for its own people.” The place of gathering for Brazilian Saints is in Brazil; the place of gathering for Nigerian Saints is in Nigeria; the place of gathering for Korean Saints is in Korea; and so forth. Zion is “the pure in heart.” Zion is wherever righteous Saints are. Publications, communications, and congregations are now such that nearly all members have access to the doctrines, keys, ordinances, and blessings of the gospel, regardless of their location.

Spiritual security will always depend upon how one lives, not where one lives. Saints in every land have equal claim upon the blessings of the Lord.

This work of Almighty God is true. He lives. Jesus is the Christ. This is His Church, restored to accomplish its divine destiny, including the promised gathering of Israel. President Gordon B. Hinckley is God’s prophet today. I so testify in the name of Jesus Christ, amen.

\newpage
\settalk{Repentance and Conversion}
\setspeaker{Russell M. Nelson}
\settalkdate{April 2007}
\talkheading{Repentance and Conversion}{Russell M. Nelson}

Last year while Elder David S. Baxter and I were driving to a stake conference, we stopped at a restaurant. Later when returning to our car, we were approached by a woman who called out to us. We were startled by her appearance. Her grooming (or lack of it) was what I might politely call “extreme.” She asked if we were elders in the Church. We said yes. Almost unrestrained, she told the story of her tragic life, swamped in sin. Now, only 28 years old, she was miserable. She felt worthless, with nothing to live for. As she spoke, the sweetness of her soul began to emerge. Pleading tearfully, she asked if there was any hope for her, any way up and out of her hopelessness.

“Yes,” we responded, “there is hope. Hope is linked to repentance. You can change. You can ‘come unto Christ, and be perfected in him.’” We urged her not to procrastinate. She sobbed humbly and thanked us sincerely.

As Elder Baxter and I continued our journey, we pondered that experience. We recalled the counsel given to a hopeless soul by Aaron, who said, “If thou wilt repent of all thy sins, and will bow down before God, and call on his name in faith, … then shalt thou receive the hope which thou desirest.”

Now, at this closing session of general conference, I too speak on repentance. I do so because the Lord has commanded His servants to cry repentance unto all people. The Master has restored His gospel to bring joy to His children, and repentance is a crucial component of that gospel.

The doctrine of repentance is as old as the gospel itself. Biblical teachings from the books of Genesis to Revelation teach repentance. Lessons from Jesus Christ during His mortal ministry include these warnings: “The kingdom of God is at hand: repent ye, and believe the gospel” and “Except ye repent, ye shall all likewise perish.”

References to repentance are even more frequent in the Book of Mormon. To the people of ancient America, the Lord gave this commandment: “Again I say unto you, ye must repent, and be baptized in my name, and become as a little child, or ye can in nowise inherit the kingdom of God.”

With the Restoration of the gospel, our Savior has again stressed this doctrine. The word repent in any of its forms appears in 47 of the 138 sections of the Doctrine and Covenants!

\intalkheader{Repent from Sin}

What does it mean to repent? We begin with a dictionary’s definition that to repent is “to turn from sin … to feel sorrow [and] regret.” To repent from sin is not easy. But the prize is worth the price. Repentance needs to be done one step at a time. Humble prayer will facilitate each essential step. As prerequisites to forgiveness, there must first be recognition, remorse, then confession. “By this ye may know if a man repenteth of his sins—behold, he will confess them and forsake them.” Confession is to be made to the person who has been wronged. Confession should be sincere and not merely an admission of guilt after proof is evident. If many persons have been offended, confession should be made to all offended parties. Acts that may affect one’s standing in the Church or the right to its privileges should be confessed promptly to the bishop, whom the Lord has called as a common judge in Israel.

The next step is restitution—to repair damage done—if possible. Then come steps to resolve to do better and refrain from relapse—to repent “with full purpose of heart.” Thanks to the ransom paid by the Atonement of Jesus Christ, full forgiveness is given to the sinner who repents and remains free from sin. To the repentant soul, Isaiah said, “Though your sins be as scarlet, they shall be as white as snow; though they be red like crimson, they shall be as wool.”

The Lord’s imperative emphasis on repentance is evident as we read from section 19 of the Doctrine and Covenants: “I command you to repent—repent, lest I smite you by the rod of my mouth, and by my wrath, and by my anger, and your sufferings be sore—how sore you know not, how exquisite you know not, yea, how hard to bear you know not.

“For behold, I, God, have suffered these things for all, that they might not suffer if they would repent;

“But if they would not repent they must suffer even as I.”

While the Lord insists on our repentance, most people don’t feel such a compelling need. They include themselves among those who try to be good. They have no evil intent. Yet the Lord is clear in His message that all need to repent—not only from sins of commission but from sins of omission as well. Such is the case in His warning to parents: “Inasmuch as parents have children in Zion … that teach them not to understand the doctrine of repentance, faith in Christ the Son of the living God, and of baptism and the gift of the Holy Ghost … , the sin be upon the heads of the parents.”

\intalkheader{Broader Meaning of the Word Repent}

The doctrine of repentance is much broader than a dictionary’s definition. When Jesus said “repent,” His disciples recorded that command in the Greek language with the verb metanoeo. This powerful word has great significance. In this word, the prefix meta means “change.” The suffix relates to four important Greek terms: nous, meaning “the mind”; gnosis, meaning “knowledge”; pneuma, meaning “spirit”; and pnoe, meaning “breath.”

Thus, when Jesus said “repent,” He asked us to change—to change our mind, knowledge, and spirit—even our breath. A prophet explained that such a change in one’s breath is to breathe with grateful acknowledgment of Him who grants each breath. King Benjamin said, “If ye should serve him who has created you … and is preserving you from day to day, by lending you breath … from one moment to another—I say, if ye should serve him with all your whole souls yet ye would be unprofitable servants.”

Yes, the Lord has commanded us to repent, to change our ways, to come unto Him, and be more like Him. This requires a total change. Alma so taught his son: “Learn wisdom in thy youth,” he said. “Learn in thy youth to keep the commandments of God. … Let all thy thoughts be directed unto the Lord; yea, let the affections of thy heart be placed upon the Lord forever.”

To repent fully is to convert completely to the Lord Jesus Christ and His holy work. Alma taught that concept when he posed these questions: “I ask of you, my brethren of the church, have ye spiritually been born of God? Have ye received his image in your countenances? Have ye experienced this mighty change in your hearts?” That change comes when we are “born again,” converted and focused upon our journey to the kingdom of God.

\intalkheader{Fruits of Repentance}

The fruits of repentance are sweet. Repentant converts find that the truths of the restored gospel govern their thoughts and deeds, shape their habits, and forge their character. They are more resilient and able to deny themselves of all ungodliness. Moreover, uncontrolled appetite, addiction to pornography or harmful drugs, unbridled passion, carnal desire, and unrighteous pride are diminished with complete conversion to the Lord and a determination to serve Him and to emulate His example. Virtue garnishes their thoughts, and self-confidence grows. Tithing is seen as a joyful and protective blessing, not as a duty or a sacrifice. Truth becomes more attractive, and things praiseworthy become more engaging.

Repentance is the Lord’s regimen for spiritual growth. King Benjamin explained that “the natural man is an enemy to God, and has been from the fall of Adam, and will be, forever and ever, unless he yields to the enticings of the Holy Spirit, and putteth off the natural man and becometh a saint through the atonement of Christ the Lord, and becometh as a child, submissive, meek, humble, patient, full of love, willing to submit to all things which the Lord seeth fit to inflict upon him, even as a child doth submit to his father.” Brothers and sisters, that means conversion! Repentance is conversion! A repentant soul is a converted soul, and a converted soul is a repentant soul.

\intalkheader{Repentance for Those Who Are Dead}

Each living person can repent. But what about those who have died? They also have opportunities to repent. Scripture declares that “the faithful elders of this dispensation, when they depart from mortal life, continue their labors in the preaching of the gospel of repentance … among those who are … under the bondage of sin in the great world of the spirits of the dead.

“The dead who repent will be redeemed, through obedience to the ordinances of the house of God,

“And after they have paid the penalty of their transgressions, and are washed clean, [they] shall receive a reward according to their works.”

The Prophet Joseph Smith further revealed that “the earth will be smitten with a curse unless there is a welding link of some kind or other between the fathers and the children. … We without [our dead] cannot be made perfect; neither can they without us be made perfect. … [This] dispensation is now beginning to usher in, that a whole and complete and perfect union, and welding together of dispensations, and keys, and powers, and glories should take place.”

“Jesus wants me for a sunbeam”? Yes! And you too! He also wants us as bonding blacksmiths—creating celestial welding links—to curb the curse of family fragmentation. The earth was created and temples provided so that families can be together forever. Many, if not most, of us could repent and be converted to more temple and family history work for our ancestors. Thus, our repentance is necessary and essential for their repentance.

For all our kindred dead, to the 28-year-old woman mired in the swamp of sin, and to each one of us, I declare that the sweet blessing of repentance is possible. It comes through complete conversion to the Lord and His holy work.

I know that God lives. Jesus is the Christ. This is His Church. His prophet today is President Gordon B. Hinckley. I so testify in the name of Jesus Christ, amen.

\newpage
\settalk{Scriptural Witnesses}
\setspeaker{Russell M. Nelson}
\settalkdate{October 2007}
\talkheading{Scriptural Witnesses}{Russell M. Nelson}

We express love and admiration for President Henry B. Eyring, Elder Quentin L. Cook, and Elder Walter F. González and pray for the Lord’s blessings to attend them in their new callings.

We extend heartfelt feelings of gratitude to each of you, brothers and sisters. Throughout the world, your examples of service and compassion are receiving a great deal of attention. At the same time, many people are wondering about this Church’s history and doctrines. Among those questioners are some who choose to cast aspersions at the Book of Mormon.

Disregard for the Book of Mormon or for any other sacred scripture concerns me deeply. In addressing that concern, I have titled my remarks “Scriptural Witnesses.”

\intalkheader{Definitions}

I will define the term scriptural as pertaining to the Bible and the scriptures of the Restoration. Members of the Church “believe the Bible to be the word of God as far as it is translated correctly; we also believe the Book of Mormon to be the word of God.” Scriptures of the Restoration also include the Doctrine and Covenants as well as the Pearl of Great Price.

A dictionary defines the noun witness as an “attestation of a fact or event: [a] testimony.” The term witness bears special significance when applied to the word of God. In the Bible we read this important declaration: “In the mouth of two or three witnesses shall every word be established.” This assures God’s children that divine doctrines are confirmed by more than one scriptural witness.

\intalkheader{Scriptures Testify of Jesus Christ}

The Bible and the Book of Mormon are both witnesses of Jesus Christ. They teach that He is the Son of God, that He lived an exemplary life, that He atoned for all mankind, that He died upon the cross and rose again as the resurrected Lord. They teach that He is the Savior of the world.

Scriptural witnesses authenticate each other. This concept was explained long ago when a prophet wrote that the Book of Mormon was “written for the intent that ye may believe [the Bible]; and if ye believe [the Bible] ye will believe [the Book of Mormon] also.” Each book refers to the other. Each book stands as evidence that God lives and speaks to His children by revelation to His prophets.

Love for the Book of Mormon expands one’s love for the Bible and vice versa. Scriptures of the Restoration do not compete with the Bible; they complement the Bible. We are indebted to martyrs who gave their lives so that we could have the Bible. It establishes the everlasting nature of the gospel and of the plan of happiness. The Book of Mormon restores and underscores biblical doctrines such as tithing, the temple, the Sabbath day, and the priesthood.

An angel proclaimed that the Book of Mormon shall establish the truth of the Bible. He also revealed that writings in the Bible available in our day are not as complete as they were when originally written by prophets and apostles. He declared that the Book of Mormon shall restore plain and precious things taken away from the Bible.

A prophecy in the Book of Mormon warned that some people would object to the notion of additional scriptures. To those who think they “need no more Bible,” consider this God-given counsel:

“Know ye not that there are more nations than one? Know ye not that I, the Lord your God, … created all men, … and that I rule in the heavens above and in the earth beneath; and I bring forth my word unto the children of men, yea, even upon all the nations of the earth?

“… Know ye not that the testimony of two nations is a witness unto you that I am God, that I remember one nation like unto another? Wherefore, I speak the same words unto one nation like unto another. And … the testimony of the two nations shall run together also.”

The scriptural story of Jesus Christ is indeed a two-hemisphere story. While Mary and Joseph in the East were making preparations for the holy child’s birth in Bethlehem, Nephi in the West was being taught by the premortal Messiah. To Nephi the Lord said, “Be of good cheer; … on the morrow come I into the world, to show unto the world that I will fulfil all that which I have caused to be spoken by the mouth of my holy prophets.”

To doubters of that second witness—the Book of Mormon—the Lord issued a warning: “Because you have treated lightly the things you have received … [you] shall remain under … condemnation until [you] repent and remember … the Book of Mormon and the former commandments which I have given … , [the Bible, and] … do according to that which I have written.”

The Lord provided other scriptures of the Restoration and declared that those words will also be fulfilled. With these scriptural witnesses, false doctrines will be confounded. With these scriptural witnesses, doctrines of the Bible are not only reaffirmed but clarified.

\intalkheader{Restoration Scriptures Clarify the Bible}

How do scriptures of the Restoration clarify the Bible? Many examples exist. I will cite but a few, beginning with the Old Testament.

Isaiah wrote, “Thou shalt … speak out of the ground, and thy speech shall be low out of the dust, and thy voice shall be, as of one that hath a familiar spirit, out of the ground, and thy speech shall whisper out of the dust.” Could any words be more descriptive of the Book of Mormon, coming as it did “out of the ground” to “whisper out of the dust” to people of our day?

But Isaiah was not the only Old Testament prophet who foretold the Book of Mormon. Ezekiel wrote:

“Take thee one stick, and write upon it, For Judah, and for the children of Israel … : then take another stick, and write upon it, For Joseph, the stick of Ephraim, and for all the house of Israel … :

“And join them one to another into one stick; and they shall become one in thine hand.”

Today, Saints living in many nations of the earth gratefully hold the Bible (the stick of Judah) and the Book of Mormon (the stick of Ephraim) bound as one in their hands.

What about the New Testament? The Book of Mormon also stands as a scriptural witness to its teachings. Examples include the miraculous birth of the babe in Bethlehem, His Sermon on the Mount, and the intense suffering of the Savior. The doctrine of the Resurrection is mentioned more frequently in the Book of Mormon than in the Bible.

The need for the Holy Ghost was mentioned by Paul. He asked, “Have ye received the Holy Ghost? … And they said unto him, We have not so much as heard whether there be any Holy Ghost.” That doctrine is clarified by another scriptural witness, conveyed through the Lord’s prophet of the Restoration. He taught us to “believe in the gift of the Holy Ghost by the laying on of hands.” That precious and powerful gift is again available to children of God.

Paul referred to the three degrees of postmortal glory when he taught that “there is one glory of the sun, … another glory of the moon, and another glory of the stars.” That glimpse into postmortal glory has been clarified by another scriptural witness. The Lord revealed that “the glory of the celestial is one, even as the glory of the sun is one.

“And the glory of the terrestrial is one, even as the glory of the moon is one.

“And the glory of the telestial is one, even as the glory of the stars is one.”

The highest of these kingdoms, the celestial, is reserved for those who obey the law of that kingdom:

“They who are not sanctified through the law … of Christ, must inherit another kingdom, even that of a terrestrial kingdom, or that of a telestial kingdom.

“For he who is not able to abide the law of a celestial kingdom cannot abide a celestial glory.”

These three degrees of glory relate to postmortal life. They relate to the immortality of the human soul. That gift of immortality became a reality because of the Atonement of Jesus Christ. This important word—atonement—in any of its forms, is mentioned only once in the King James Version of the New Testament! In the Book of Mormon it appears 39 times!

John, writing in the New Testament’s book of Revelation, foresaw “another angel fly in the midst of heaven, having the everlasting gospel to preach unto them that dwell on the earth, and to every nation, and kindred, and tongue, and people.” One specific angel held keys of responsibility for the Book of Mormon. That was the angel Moroni! These examples are but a few of many biblical doctrines that are clarified by scriptures of the Restoration.

\intalkheader{The Book of Mormon: Another Testament of Jesus Christ}

Willingly we share scriptures of the Restoration with people across the world. The Book of Mormon records the personal ministry of the resurrected Lord to people of ancient America. Ponder these everlasting truths that He proclaimed:

“Behold, I am Jesus Christ the Son of God. I created the heavens and the earth, and all things that in them are. I was with the Father from the beginning. …

“… The scriptures concerning my coming are fulfilled. …

“I am the light and the life of the world.”

The Savior continued:

“Whoso cometh unto me with a broken heart and a contrite spirit, him will I baptize with fire and with the Holy Ghost. …

“… I have come unto the world to bring redemption unto the world, to save the world from sin.

“Therefore, whoso repenteth and cometh unto me as a little child, him will I receive, for of such is the kingdom of God. … I have laid down my life, and have taken it up again; therefore repent, and come unto me … and be saved.”

Those declarations of the Lord summarize who He really is and who He really wants us to be. He wants us to come unto Him and, in due course, be embraced gloriously in His loving arms.

I express my profound gratitude for scriptural witnesses. I have seen the mighty change that comes to people who apply the Lord’s teachings in their lives. Such transformation leads to the blessing of eternal life.

I know that God lives. Jesus is the Christ. His gospel has been restored in these latter days. President Gordon B. Hinckley is His prophet at this time. I so testify in the name of Jesus Christ, amen.

\newpage
\settalk{Salvation and Exaltation}
\setspeaker{Russell M. Nelson}
\settalkdate{April 2008}
\talkheading{Salvation and Exaltation}{Russell M. Nelson}

Gratefully we welcome Elder D. Todd Christofferson to the Quorum of the Twelve Apostles. Wholeheartedly we sustain this wonderful First Presidency and all who have been called.

Brethren and sisters, when we received the news that President Gordon B. Hinckley had passed away, each of us immediately felt a deep sense of loss. Knowing that his destiny was in the hands of the Lord, however, we have felt our mood shift from grief to gratitude. We are very grateful for what we have learned from this great prophet of God.

Today, at this solemn assembly, we have complied with the will of the Lord, who said that “it shall not be given to any one to go forth to preach my gospel, or to build up my church, except he be ordained by … one who has authority, and it is known to the church that he has authority and has been regularly ordained by the heads of the church.” This law of common consent has been invoked, and the Church will move forward on its prescribed course.

Members throughout the world sustain President Thomas S. Monson and his able counselors. We are “no more strangers and foreigners, but fellowcitizens with the saints, and of the household of God;

“And are built upon the foundation of the apostles and prophets, Jesus Christ himself being the chief corner stone.”

The Lord revealed why “he gave some, apostles; and some, prophets.” It is “for the perfecting of the saints, for the work of the ministry, for the edifying of the body of Christ:

“Till we all come in the unity of the faith, and of the knowledge of the Son of God.”

Thus the ministry of the Apostles—the First Presidency and the Twelve—is to bring about that unity of the faith and to proclaim our knowledge of the Master. Our ministry is to bless the lives of all who will learn and follow the “more excellent way” of the Lord. And we are to help people prepare for their potential salvation and exaltation.

The third article of faith declares that “through the Atonement of Christ, all mankind may be saved, by obedience to the laws and ordinances of the Gospel.”

To be saved—or to gain salvation—means to be saved from physical and spiritual death. Because of the Resurrection of Jesus Christ, all people will be resurrected and saved from physical death. People may also be saved from individual spiritual death through the Atonement of Jesus Christ, by their faith in Him, by living in obedience to the laws and ordinances of His gospel, and by serving Him.

To be exalted—or to gain exaltation—refers to the highest state of happiness and glory in the celestial realm. These blessings can come to us after we leave this frail and mortal existence. The time to prepare for our eventual salvation and exaltation is now.

As part of that preparation, one must first hear and understand the gospel. For this reason the gospel of Jesus Christ is being taken to “every nation, and kindred, and tongue, and people.”

\intalkheader{Individual Responsibility}

Some years ago I met with a tribal king in Africa. When he realized that he was being taught by an Apostle of the Lord, he was deeply moved. He said that throngs of his people would be baptized if he were to give them such a mandate. I thanked him for his kindness but explained that the Lord does not work in that way.

The development of faith in the Lord is an individual matter. Repentance is also an individual matter. Only as an individual can one be baptized and receive the Holy Ghost. Each of us is born individually; likewise, each of us is “born again” individually. Salvation is an individual matter.

\intalkheader{Family Responsibilities}

Individual progression is fostered in the family, which is “central to the Creator’s plan for the eternal destiny of His children.” The home is to be God’s laboratory of love and service. There a husband is to love his wife, a wife is to love her husband, and parents and children are to love one another.

Throughout the world, the family is increasingly under attack. If families fail, many of our political, economic, and social systems will also fail. And if families fail, their glorious eternal potential cannot be realized.

Our Heavenly Father wants husbands and wives to be faithful to each other and to esteem and treat their children as an heritage from the Lord. In such a family we study the scriptures and pray together. And we fix our focus on the temple. There we receive the highest blessings that God has in store for His faithful children.

Thanks to God’s great plan of happiness, families can be together forever—as exalted beings. Our Heavenly Father declared, “This is my work and my glory—to bring to pass the immortality and eternal life of man.” Both of His objectives were enabled by the Atonement of His Beloved Son, Jesus Christ. His Atonement made resurrection a reality and eternal life a possibility for all who would ever live.

Resurrection, or immortality, comes to every man and every woman as an unconditional gift.

Eternal life, or celestial glory or exaltation, is a conditional gift. Conditions of this gift have been established by the Lord, who said, “If you keep my commandments and endure to the end you shall have eternal life, which gift is the greatest of all the gifts of God.” Those qualifying conditions include faith in the Lord, repentance, baptism, receiving the Holy Ghost, and remaining faithful to the ordinances and covenants of the temple.

No man in this Church can obtain the highest degree of celestial glory without a worthy woman who is sealed to him. This temple ordinance enables eventual exaltation for both of them.

In Church callings we are subject to release. But we cannot be released as parents. From the first days of human history, the Lord has commanded parents to teach the gospel to their children. Moses wrote, “Thou shalt teach … diligently … thy children, and shalt talk of [God’s words] when thou sittest in thine house, and when thou walkest by the way, and when thou liest down, and when thou risest up.”

In our day the Lord has added, “Bring up your children in light and truth.” The Church is to assist and not to replace parents in their responsibilities to teach their children.

In this day of rampant immorality and addictive pornography, parents have a sacred responsibility to teach their children the importance of God in their lives. Those evils, so highly destructive of divine potential, are to be strictly shunned by children of God.

We are also to teach our children to honor their parents. The fifth commandment states, “Honour thy father and thy mother: that thy days may be long upon the land which the Lord thy God giveth thee.”

How can we best teach our children? The Lord has given us specific instruction:

“No power or influence can or ought to be maintained by virtue of the priesthood, only by persuasion, by long-suffering, by gentleness and meekness, and by love unfeigned;

“By kindness, and pure knowledge, which shall greatly enlarge the soul without hypocrisy, and without guile—

“Reproving betimes with sharpness, when moved upon by the Holy Ghost; and then showing forth afterwards an increase of love toward him whom thou hast reproved, lest he esteem thee to be his enemy.”

When a child needs correction, you might ask yourself, “What can I say or do that would persuade him or her to choose a better way?” When giving necessary correction, do it quietly, privately, lovingly, and not publicly. If a rebuke is required, show an increase of love promptly so that seeds of resentment may not remain. To be persuasive, your love must be sincere and your teachings based on divine doctrine and correct principles.

Do not try to control your children. Instead, listen to them, help them to learn the gospel, inspire them, and lead them toward eternal life. You are God’s agents in the care of children He has entrusted to you. Let His divine influence remain in your hearts as you teach and persuade.

\intalkheader{Ancestors}

Any discussion of family responsibilities to prepare for exaltation would be incomplete if we included only mother, father, and children. What about grandparents and other ancestors? The Lord has revealed that we cannot become perfect without them; neither can they without us be made perfect. Sealing ordinances are essential to exaltation. A wife needs to be sealed to her husband; children need to be sealed to their parents; and we all need to be connected with our ancestors.

What about those who are not able to marry in this life or those who cannot be sealed to their parents in this life? We know that the Lord will judge each of us according to the desires of our hearts, as well as our works, and that the blessings of exaltation will be given to all who are worthy.

We as children of the covenant are highly favored. In our hearts have been planted the promises made to Fathers Abraham, Isaac, and Jacob. The Lord has said:

“Ye are lawful heirs, according to the flesh, and have been hid from the world with Christ in God. …

“Therefore, blessed are ye if ye continue in my goodness, a light unto the Gentiles, and through this priesthood, a savior unto my people Israel.”

This life is the time to prepare for salvation and exaltation. In God’s eternal plan, salvation is an individual matter; exaltation is a family matter.

As children of the covenant, we have met in this morning’s solemn assembly. Attention has been focused upon the sacred titles of prophets and apostles. But the final responsibility to prepare for salvation and exaltation rests upon each person, accountable for individual agency, acting in one’s own family, bearing another sacred title of mother, father, daughter, son, grandmother, or grandfather.

In those responsible roles, may we go forward in faith, led by Jesus Christ, whose Church this is, and by His prophet, through whom He speaks, I pray in the name of Jesus Christ, amen.

\newpage
\settalk{Celestial Marriage}
\setspeaker{Russell M. Nelson}
\settalkdate{October 2008}
\talkheading{Celestial Marriage}{Russell M. Nelson}

My beloved brethren and sisters, I am deeply grateful for each of you. Together we feel a profound sense of gratitude for the gospel of Jesus Christ. In this world abounding with misery, we are truly thankful for God’s “great plan of happiness.” His plan declares that men and women are “that they might have joy.” That joy comes when we choose to live in harmony with God’s eternal plan.

The importance of choice may be illustrated by a homespun concept that came to mind one day when I was shopping in a large retail store. I call it “patterns of the shopper.” As shopping is part of our daily life, these patterns may be familiar.

Wise shoppers study their options thoroughly before they make a selection. They focus primarily on the quality and durability of a desired product. They want the very best. In contrast, some shoppers look for bargains, and others may splurge, only to learn later—much to their dismay—that their choice did not endure well. And sadly, there are those rare individuals who cast aside their personal integrity and steal what they want. We call them shoplifters.

The patterns of the shopper may be applied to the topic of marriage. A couple in love can choose a marriage of the highest quality or a lesser type that will not endure. Or they can choose neither and brazenly steal what they want as “marital shoplifters.”

The subject of marriage is debated across the world, where various arrangements exist for conjugal living. My purpose in speaking out on this topic is to declare, as an Apostle of the Lord, that marriage between a man and a woman is sacred—it is ordained of God. I also assert the virtue of a temple marriage. It is the highest and most enduring type of marriage that our Creator can offer to His children.

While salvation is an individual matter, exaltation is a family matter. Only those who are married in the temple and whose marriage is sealed by the Holy Spirit of Promise will continue as spouses after death and receive the highest degree of celestial glory, or exaltation. A temple marriage is also called a celestial marriage. Within the celestial glory are three levels. To obtain the highest, a husband and wife must be sealed for time and all eternity and keep their covenants made in a holy temple.

The noblest yearning of the human heart is for a marriage that can endure beyond death. Fidelity to a temple marriage does that. It allows families to be together forever.

This goal is glorious. All Church activities, advancements, quorums, and classes are means to the end of an exalted family.

To make this goal possible, our Heavenly Father has restored priesthood keys in this dispensation so that essential ordinances in His plan can be performed by proper authority. Heavenly messengers—including John the Baptist; Peter, James, and John; Moses, Elias, and Elijah—have participated in that restoration.

Knowledge of this revealed truth is spreading across the earth. We, as the Lord’s prophets and apostles, again proclaim to the world that “the family is central to the Creator’s plan for the eternal destiny of His children.”

We further proclaim that “all human beings—male and female—are created in the image of God. Each is a beloved spirit son or daughter of heavenly parents, and, as such, each has a divine nature and destiny. Gender is an essential characteristic of individual premortal, mortal, and eternal identity and purpose.

“In the premortal realm, spirit sons and daughters knew and worshipped God as their Eternal Father and accepted His plan by which His children could obtain a physical body and gain earthly experience to progress toward perfection and ultimately realize their divine destiny as heirs of eternal life. [Heavenly Father’s great] plan of happiness enables family relationships to be perpetuated beyond the grave. Sacred ordinances and covenants available in holy temples make it possible for individuals to return to the presence of God and for families to be united eternally.”

That proclamation on the family helps us realize that celestial marriage brings greater possibilities for happiness than does any other relationship. The earth was created and this Church was restored so that families could be formed, sealed, and exalted eternally.

Scriptures declare that “it is lawful that [a man] should have one wife, and they twain shall be one flesh, and all this that the earth might answer the end of its creation.” Another affirms that “the man [is not] without the woman, neither the woman without the man, in the Lord.” Thus, marriage is not only an exalting principle of the gospel; it is a divine commandment.

Our Heavenly Father declared, “This is my work and my glory—to bring to pass the immortality and eternal life of man.” The Atonement of His Beloved Son enabled both of these objectives to be realized. Because of the Atonement, immortality—or resurrection from the dead—became a reality for all. And because of the Atonement, eternal life—which is living forever in God’s presence, the “greatest of all the gifts of God”—became a possibility. To qualify for eternal life, we must make an eternal and everlasting covenant with our Heavenly Father. This means that a temple marriage is not only between husband and wife; it embraces a partnership with God.

The family proclamation also reminds us that “husband and wife have a solemn responsibility to love and care for each other.” Children born of that union are “an heritage of the Lord.” When a family is sealed in the temple, that family may become as eternal as the kingdom of God itself.

Such a reward requires more than a hopeful wish. On occasion I read in a newspaper obituary of an expectation that a recent death has reunited that person with a deceased spouse, when, in fact, they did not choose the eternal option. Instead, they opted for a marriage that was valid only as long as they both should live. Heavenly Father had offered them a supernal gift, but they refused it. And in rejecting the gift, they rejected the Giver of the gift.

One strong sentence of scripture clearly distinguishes between a hopeful wish and eternal truth: “All covenants, contracts, … obligations, oaths, vows, … or expectations, that are not made and entered into and sealed by the Holy Spirit of promise, of him who is anointed, both as well for time and for all eternity, … are of no efficacy, virtue, or force in and after the resurrection from the dead; for all contracts that are not made unto this end have an end when men are dead.”

These truths are absolute. Members of this Church invite all people to learn them and to qualify for eternal life. We invite all to gain faith in God the Eternal Father and in His Son, Jesus Christ, to repent, to receive the Holy Ghost, to obtain the blessings of the temple, to make and keep sacred covenants, and to endure to the end.

Mercifully, God’s great plan of happiness and its eternal blessings can be extended to those who did not have the opportunity to hear the gospel in mortality. Temple ordinances can be done vicariously for them.

But what of the many mature members of the Church who are not married? Through no failing of their own, they deal with the trials of life alone. Be we all reminded that, in the Lord’s own way and time, no blessings will be withheld from His faithful Saints. The Lord will judge and reward each individual according to heartfelt desire as well as deed.

Meanwhile, mortal misunderstandings can make mischief in a marriage. In fact, each marriage starts with two built-in handicaps. It involves two imperfect people. Happiness can come to them only through their earnest effort. Just as harmony comes from an orchestra only when its members make a concerted effort, so harmony in marriage also requires a concerted effort. That effort will succeed if each partner will minimize personal demands and maximize actions of loving selflessness.

President Thomas S. Monson has said: “To find real happiness, we must seek for it in a focus outside ourselves. No one has learned the meaning of living until he has surrendered his ego to the service of his fellow man. Service to others is akin to duty—the fulfillment of which brings true joy.”

Harmony in marriage comes only when one esteems the welfare of his or her spouse among the highest of priorities. When that really happens, a celestial marriage becomes a reality, bringing great joy in this life and in the life to come.

God’s plan of happiness allows us to choose for ourselves. As with the patterns of the shopper, we may choose celestial marriage or lesser alternatives. Some marital options are cheap, some are costly, and some are cunningly crafted by the adversary. Beware of his options; they always breed misery!

The best choice is a celestial marriage. Thankfully, if a lesser choice has previously been made, a choice can now be made to upgrade it to the best choice. That requires a mighty change of heart and a permanent personal upgrade. Blessings so derived are worth all efforts made.

The full realization of the blessings of a temple marriage is almost beyond our mortal comprehension. Such a marriage will continue to grow in the celestial realm. There we can become perfected. As Jesus ultimately received the fulness of the glory of the Father, so we may “come unto the Father … and in due time receive of his fulness.”

Celestial marriage is a pivotal part of preparation for eternal life. It requires one to be married to the right person, in the right place, by the right authority, and to obey that sacred covenant faithfully. Then one may be assured of exaltation in the celestial kingdom of God. I so testify in the name of Jesus Christ, amen.

\newpage
\settalk{Lessons from the Lord’s Prayers}
\setspeaker{Russell M. Nelson}
\settalkdate{April 2009}
\talkheading{Lessons from the Lord’s Prayers}{Russell M. Nelson}

With you, my dear brothers and sisters, I express love and admiration for Elder Neil L. Andersen. His call to the holy apostleship has come from the Lord as revealed to His prophet, President Thomas S. Monson. Throughout his life, President Monson has refined his ability to hearken to the will of the Lord. As the Savior submitted His will to Heavenly Father, so the prophet submits his will to the Lord. Thank you, President Monson, for developing and using that power. We congratulate you, Elder Andersen, and we pray for you!

\intalkheader{The Lord’s Prayer}

Our prayers follow patterns and teachings of the Lord Jesus Christ. He taught us how to pray. From His prayers we can learn many important lessons. We can begin with the Lord’s Prayer and add lessons from other prayers He has given.

As I recite the Lord’s Prayer, listen for lessons:

“Our Father which art in heaven, Hallowed be thy name.

“Thy kingdom come. Thy will be done in earth, as it is in heaven.

“Give us this day our daily bread.

“And forgive us our debts, as we forgive our debtors.

“And lead us not into temptation, but deliver us from evil: For thine is the kingdom, and the power, and the glory, for ever. Amen.”

The Lord’s Prayer is recorded twice in the New Testament and once in the Book of Mormon. It is also included in the Joseph Smith Translation of the Bible, where clarification is provided by these two phrases:

The clarification on forgiveness is supported by other statements of the Master. He said to His servants, “Inasmuch as you have forgiven one another your trespasses, even so I, the Lord, forgive you.” In other words, if one is to be forgiven, one must first forgive. The clarification on temptation is helpful, for surely we would not be led into temptation by Deity. The Lord said, “Watch and pray, that ye enter not into temptation.”

Though the four versions of the Lord’s Prayer are not identical, they all open with a salutation to “Our Father,” signifying a close relationship between God and His children. The phrase “hallowed be thy name” reflects the respect and worshipful attitude that we should feel as we pray. “Thy will be done” expresses a concept that we will discuss later.

His request for “daily bread” includes a need for spiritual nourishment as well. Jesus, who called Himself “the bread of life,” gave a promise: “He that cometh to me shall never hunger.” And as we partake of sacramental emblems worthily, we are further promised that we may always have His Spirit to be with us. That is spiritual sustenance that cannot be obtained in any other way.

As the Lord closes His prayer, He acknowledges God’s great power and glory, ending with “Amen.” Our prayers also close with amen. Though it is pronounced differently in various languages, its meaning is the same. It means “truly” or “verily.” Adding amen solemnly affirms a sermon or a prayer. Those who concur should each add an audible amen to signify “that is my solemn declaration too.”

The Lord prefaced His prayer by first asking His followers to avoid “vain repetitions” and to pray “after this manner.” Thus, the Lord’s Prayer serves as a pattern to follow and not as a piece to memorize and recite repetitively. The Master simply wants us to pray for God’s help while we strive constantly to resist evil and live righteously.

\intalkheader{Intercessory Prayers}

Other prayers of the Lord are also instructive, especially His intercessory prayers. They are so named because the Lord prayerfully interceded with His Father for the benefit of His disciples. Picture in your mind the Savior of the world kneeling in prayer, as I quote from John chapter 17:

“These words spake Jesus, and lifted up his eyes to heaven, and said, Father, … glorify thy Son, that thy Son also may glorify thee. …

“… I have finished the work which thou gavest me to do. …

“For I have given unto them the words which thou gavest me; and they have received them, and have known surely that I came out from thee, and they have believed that thou didst send me.

“I pray for them.”

From this prayer of the Lord we learn how keenly He feels His responsibility as our Mediator and Advocate with the Father. Just as keenly, we should feel our responsibility to keep His commandments and endure to the end.

An intercessory prayer was also given by Jesus for the people of ancient America. The record states that “no one can conceive of the joy which filled our souls at the time we heard him pray for us unto the Father.” Then Jesus added: “Blessed are ye because of your faith. And now behold, my joy is full.”

In a later prayer, Jesus included a plea for unity. “Father,” He said, “I pray unto thee for them, … that they may believe in me, that I may be in them as thou, Father, art in me, that we may be one.” We too can pray for unity. We can pray to be of one heart and one mind with the Lord’s anointed and with our loved ones. We can pray for mutual understanding and respect between ourselves and our neighbors. If we really care for others, we should pray for them. “Pray one for another,” taught James, for “the effectual fervent prayer of a righteous man availeth much.”

\intalkheader{Lessons from Other Prayers}

Other lessons about prayer were taught by the Lord. He told His disciples that “ye must always pray unto the Father in my name.” The Savior further emphasized, “Pray in your families unto the Father, always in my name.” Obediently, we apply that lesson when we pray to our Heavenly Father in the name of Jesus Christ.

Another of the Lord’s prayers teaches a lesson repeated in three consecutive verses:

“Father, I thank thee that thou hast given the Holy Ghost unto these whom I have chosen. …

“Father, I pray thee that thou wilt give the Holy Ghost unto all them that shall believe in their words.

“Father, thou hast given them the Holy Ghost because they believe in me.”

If companionship of the Holy Ghost is that important, we should pray for it too. We should likewise help all converts and our children cultivate the gift of the Holy Ghost. As we so pray, the Holy Ghost can become a vital force for good in our lives.

\intalkheader{Enhancing Our Prayers}

The Lord has taught ways by which our prayers can be enhanced. For example, He said that “the song of the righteous is a prayer unto me, and it shall be answered with a blessing upon their heads.”

Prayer can also be enhanced by fasting. The Lord said, “I give unto you a commandment that ye shall continue in prayer and fasting from this time forth.” A plea for wisdom in fasting was offered by President Joseph F. Smith, who cautioned that “there is such a thing as overdoing. A man may fast and pray till he kills himself; and there isn’t any necessity for it; nor wisdom in it. … The Lord can hear a simple prayer, offered in faith, in half a dozen words, and he will recognize fasting that may not continue more than twenty-four hours, just as readily and as effectually as He will answer a prayer of a thousand words and fasting for a month. … The Lord will accept that which is enough, with a good deal more pleasure and satisfaction than that which is too much and unnecessary.”

The concept of “too much and unnecessary” could also apply to the length of our prayers. A closing prayer in a Church meeting need not include a summary of each message and should not become an unscheduled sermon. Private prayers can be as long as we want, but public prayers ought to be short supplications for the Spirit of the Lord to be with us or brief declarations of gratitude for what has transpired.

Our prayers can be enhanced in other ways. We can use “right words”—special pronouns—in reference to Deity. While worldly manners of daily dress and speech are becoming more casual, we have been asked to protect the formal, proper language of prayer. In our prayers we use the respectful pronouns Thee, Thou, Thy, and Thine instead of You, Your, and Yours. Doing so helps us to be humble. That can also enhance our prayers. Scripture so declares, “Be thou humble; and the Lord thy God shall lead thee by the hand, and give thee answer to thy prayers.”

Prayer begins with individual initiative. “Behold,” saith the Lord, “I stand at the door, and knock: if any man hear my voice, and open the door, I will come in to him, and will sup with him, and he with me.” That door is opened when we pray to our Heavenly Father in the name of Jesus Christ.

When should we pray? Whenever we desire! Alma taught, “Counsel with the Lord in all thy doings, and he will direct thee for good; yea, when thou liest down at night lie down unto the Lord, … and when thou risest in the morning let thy heart be full of thanks unto God; and if ye do these things, ye shall be lifted up at the last day.” Jesus reminded His disciples “that they should not cease to pray in their hearts.”

The practice of Church members is to kneel in family prayer each morning and evening, plus having daily personal prayers and blessings on our food. President Monson said, “As we offer unto the Lord our family and our personal prayers, let us do so with faith and trust in Him.” And so, in praying for temporal and spiritual blessings, we should all plead, as did Jesus in the Lord’s Prayer, “Thy will be done.”

Jesus Christ, the Savior of the world—He who ransomed us with His blood—is our Redeemer and our Exemplar. At the close of His mortal mission, He prayed that His will—as the Beloved Son—might be swallowed up in the will of the Father. In that crucial hour the Savior cried, “Father, … not as I will, but as thou wilt.” So we should pray to God, “Thy will be done.”

And let us ever pray “that [the Lord’s] kingdom may go forth upon the earth, that the inhabitants … may … be prepared for the days … [when] the Son of Man shall come down … in the brightness of his glory, to meet the kingdom of God which is set up on the earth.”

In our daily lives and in our own crucial hours, may we fervently apply these precious lessons from the Lord, I pray in the sacred name of Jesus Christ, amen.

\newpage
\settalk{Ask, Seek, Knock}
\setspeaker{Russell M. Nelson}
\settalkdate{October 2009}
\talkheading{Ask, Seek, Knock}{Russell M. Nelson}

My beloved brethren and sisters, I am very grateful for each of you. I am thankful too for the miracle of modern communication that allows this conference to reach millions of people throughout the world.

Today’s technology also allows us to use wireless telephones to exchange information rapidly. Recently Wendy and I were on assignment on another continent when we learned that a new baby had arrived in our family. We received the good news minutes after that birth had occurred half a world away.

Even more amazing than modern technology is our opportunity to access information directly from heaven, without hardware, software, or monthly service fees. It is one of the most marvelous gifts the Lord has offered to mortals. It is His generous invitation to “ask, and it shall be given you; seek, and ye shall find; knock, and it shall be opened unto you.”

This timeless offer to provide personal revelation is extended to all of His children. It almost sounds too good to be true. But it is true! I have received and responded to that heavenly help. And I have learned that I always need to be ready to receive it.

Years ago, while immersed in the task of preparing a talk for general conference, I was aroused from a sound sleep with an idea impressed strongly upon my mind. Immediately I reached for pencil and paper near my bed and wrote as rapidly as I could. I went back to sleep, knowing I had captured that great impression. The next morning I looked at that piece of paper and found, much to my dismay, that my writing was totally illegible! I still keep pencil and paper at my bedside, but I write more carefully now.

To access information from heaven, one must first have a firm faith and a deep desire. One needs to “ask with a sincere heart [and] real intent, having faith in [Jesus] Christ.” “Real intent” means that one really intends to follow the divine direction given.

The next requirement is to study the matter diligently. This concept was taught to leaders of this restored Church when they were first learning how to gain personal revelation. The Lord instructed them, “I say unto you, that you must study it out in your mind; then you must ask me if it be right, and if it is right I will cause that your bosom shall burn within you; therefore, you shall feel that it is right.”

Part of being prepared is to know and obey the relevant teachings of the Lord. Some of His timeless truths are applicable generally, such as the commandments not to steal, not to kill, and not to bear false witness. Other teachings or commandments are also general, such as those regarding the Sabbath, the sacrament, baptism, and confirmation.

Some revelations have been given for unique circumstances, such as Noah’s building of the ark or the necessity for prophets like Moses, Lehi, and Brigham to lead their followers in arduous travel. God’s long-established pattern of teaching His children through prophets assures us that He will bless each prophet and that He will bless those who heed prophetic counsel.

A desire to follow the prophet requires much effort because the natural man knows very little of God and even less of His prophet. Paul wrote that “the natural man [receives] not the things of the Spirit of God: for they are foolishness unto him: neither can he know them, because they are spiritually discerned.” The change from being a natural man to a devoted disciple is a mighty one.

Another prophet taught that “the natural man is an enemy to God, and has been from the fall of Adam, and will be, forever and ever, unless he yields to the enticings of the Holy Spirit, and putteth off the natural man and becometh a saint through the atonement of Christ the Lord, and becometh as a child, submissive, meek, humble, patient, full of love, willing to submit to all things which the Lord seeth fit to inflict upon him, even as a child doth submit to his father.”

Recently I observed such a mighty change in a man whom I first met about 10 years ago. He had come to a stake conference at which his son was sustained as a member of the new stake presidency. This father was not a member of the Church. After his son had been set apart, I put my arms around this father and praised him for having such a wonderful son. Then I boldly declared: “The day will come when you will want to have this son sealed to you and your wife in a holy temple. And when that day comes, I would be honored to perform that sealing for you.”

During the subsequent decade, I did not see this man. Six weeks ago he and his wife came to my office. He greeted me warmly and recounted how startled he was with my earlier invitation. He didn’t do much about it until later, when his hearing began to fail. Then he awakened to the realization that his body was changing and that his time on earth was indeed limited. In due course he ultimately lost his hearing. At the same time, he became converted and joined the Church.

During our visit he summarized his total transformation: “I had to lose my hearing before I could heed the great importance of your message. Then I realized how much I wanted my loved ones to be sealed to me. I am now worthy and prepared. Will you please perform that sealing?” This I did with a deep sense of gratitude to God.

After such a conversion takes place, even further spiritual refinement can come. Personal revelation can be honed to become spiritual discernment. To discern means to sift, to separate, or to distinguish. The gift of spiritual discernment is a supernal gift. It allows members of the Church to see things not visible and to feel things not tangible.

Bishops are entitled to that gift as they face the task of seeking out the poor and caring for the needy. With that gift, sisters may view trends in the world and detect those that, however popular, are shallow or even dangerous. Members can discern between schemes that are flashy and fleeting and those refinements that are uplifting and enduring.

Discernment was implicit in important instructions President John Taylor gave long ago. He taught stake presidents, bishops, and others: “It is the right of those holding [these positions] to obtain the word of God with regard to the duties of their presidencies that they may more effectually carry out His holy purposes. None of the callings or positions in the priesthood are intended for the personal benefits, emoluments and fame of those who hold them, but are expressly given to fulfil the purposes of our Heavenly Father and build up the Kingdom of God upon the earth. We … seek to understand the will of God, and then carry it out; and see that it is carried out by those over whom we have the charge.”

For each of you to receive revelation unique to your own needs and responsibilities, certain guidelines prevail. The Lord asks you to develop “faith, hope, charity and love, with an eye single to the glory of God.” Then with your firm “faith, virtue, knowledge, temperance, patience, brotherly kindness, godliness, charity, humility, [and] diligence,” you may ask, and you will receive; you may knock, and it will be opened unto you.

Revelation from God is always compatible with His eternal law. It never contradicts His doctrine. It is facilitated by proper reverence for Deity. The Master gave this instruction:

“I, the Lord, am merciful and gracious unto those who fear me, and delight to honor those who serve me in righteousness and in truth unto the end.

“Great shall be their reward and eternal shall be their glory.

“… To them will I reveal all mysteries [and] my will concerning all things pertaining to my kingdom.”

Revelation need not all come at once. It may be incremental. “Saith the Lord God: I will give unto the children of men line upon line, precept upon precept, here a little and there a little; and blessed are those who hearken unto my precepts, and lend an ear unto my counsel, for they shall learn wisdom; for unto him that receiveth I will give more.” Patience and perseverance are part of our eternal progression.

Prophets have described what they felt while receiving revelation. Joseph Smith and Oliver Cowdery reported that “the veil was taken from our minds, and the eyes of our understanding were opened.” President Joseph F. Smith wrote, “As I pondered over these things which are written, the eyes of my understanding were opened, and the Spirit of the Lord rested upon me.”

Every Latter-day Saint may merit personal revelation. The invitation to ask, seek, and knock for divine direction exists because God lives and Jesus is the living Christ. It exists because this is His living Church. And we are blessed today because President Thomas S. Monson is His living prophet. That we may hearken to and heed his prophetic counsel is my prayer in the name of Jesus Christ, amen.

\newpage
\settalk{Generations Linked in Love}
\setspeaker{Russell M. Nelson}
\settalkdate{April 2010}
\talkheading{Generations Linked in Love}{Russell M. Nelson}

Correspondence we receive at Easter or Christmastime provides refreshing memories of dear friends and relatives. Some of those messages are enriched with precious family photographs. Here is one that really caught my attention.

This is one of our great-granddaughters. I’ll call her “Dear Ruby.” This picture reminded me of her mother when she was about the same age. From my files I retrieved this photo of “Dear Ruby’s” mother—one of our granddaughters.

I took this picture of “Dear Ruby’s” mother some 29 years ago. Her eyes are still just as blue.

Fond memories surfaced from half a century ago, when “Dear Ruby’s” grandmother—one of our daughters—was then the newest member of our family. This is one of her baby pictures. Now she is a loving grandmother, and I am “Dear Ruby’s” great-grandfather. (I won’t show you my picture as a baby. That would not help.) These photographs suggest the love that links our four generations together.

When I think of the love I feel for each member of our family, I sense, to a slight degree, the love that our Heavenly Father bears for His children. While the family is under attack throughout the world, The Church of Jesus Christ of Latter-day Saints proclaims, promotes, and protects the truth that the family is central to the Creator’s plan for the eternal destiny of His children. “The Family: A Proclamation to the World” and our vast family history efforts are but two evidences of how this Church brings hope and help to the sacred institution of the family.

We teach that God’s love for His children is infinite. Regardless of race, nationality, or gender, He loves all of them. He has done so from the beginning and will continue to do so. He invites all to gain eternal exaltation for their family. His work and His glory is to bring to pass the immortality and eternal life—the exaltation—of His children. “For God so loved the world, that he gave his only begotten Son, that whosoever believeth in him should not perish, but have everlasting life.”

The Atonement of His Beloved Son enabled both of the Father’s objectives to be fulfilled. Without the Atonement, there would be no immortality. Without the Atonement, there would be no return to the presence of the Father and no continuation of the family beyond the grave.

Because of the Atonement, these consummate blessings can be realized by each of God’s children who obey His eternal laws. Through the ages, many of His children have had access to the blessings of the gospel, but many more have not. Before the foundation of the world, our Heavenly Father instituted the ordinance of baptism for those who die without a knowledge of the gospel. He loves those children too.

He also provided a way for them to be part of an eternal family. Every human being who comes to this earth is the product of generations of parents. We have a natural yearning to connect with our ancestors. This desire dwells in our hearts, regardless of age.

Consider the spiritual connections that are formed when a young woman helps her grandmother enter family information into a computer or when a young man sees the name of his great-grandfather on a census record. When our hearts turn to our ancestors, something changes inside us. We feel part of something greater than ourselves. Our inborn yearnings for family connections are fulfilled when we are linked to our ancestors through sacred ordinances of the temple.

Because of the importance of this work, the Church has built temples closer to the people, and family history research is being facilitated as never before. Methods to find and prepare names for temple ordinances are also improving. At the October 2005 conference, President Gordon B. Hinckley announced an exciting step forward in family history and temple work. He said: “One of the most troublesome aspects of our temple activity is that as we get more and more temples … across the earth there is duplication of effort in proxy work. … We, therefore, have been engaged for some time in a very difficult undertaking. … The solution lies in complex computer technology.”

Since then, not only has duplication been reduced, but procedures have been simplified so that virtually every member of the Church can participate in temple and family history work. Gone are the days when this sacred work was done only by specialists. No matter your situation, you can make family history a part of your life right now. Primary children can draw a family tree. Youth can participate in proxy baptisms. They can also help the older generation work with computers. Parents can relate stories of their lives to their posterity. Worthy adult members can hold a temple recommend and perform temple ordinances for their own kin.

The Prophet Joseph Smith said, “The greatest responsibility in this world that God has laid upon us is to seek after our dead.” New technology makes it easier than ever to fulfill that responsibility. Temple and family history work is now facilitated by a system known as the “new FamilySearch.” This Internet-based system helps members identify their ancestors, determine what ordinance work needs to be done for them, and prepare their names for the temple. It can be accessed from home, a family history center, or wherever the Internet is available. The steps are easy to follow.

You first identify individuals for whom you desire to do temple work.

Then you print out a Family Ordinance Request. This document provides the information needed at the temple and eliminates the need to take computer discs with you.

From the Family Ordinance Request, ordinance cards are printed at the temple. After an ordinance is performed, it is recorded and entered into new FamilySearch on that very day.

Now, what about those of you who have no access to a computer or prefer not to use this technology? Don’t worry! Take one step at a time. Start at home. Begin with an empty cardboard box, as suggested by President Boyd K. Packer. Put in that box important information about you and your family. Add data gathered from others of your family. Then avail yourself of assistance from your ward or branch family history consultant. The new FamilySearch system enables a consultant to perform all needed computer functions for you, including preparing names for the temple. About 60,000 consultants serve throughout the world. One in your ward or branch can be very helpful to you.

The new FamilySearch system changes the dynamics of family history work by facilitating the creation of one common pedigree. In the past, a person worked separately, keeping his or her own family records. One often worked without knowing what other family members were doing. Now each person can contribute information while coordinating with others in developing their family tree.

While the new FamilySearch is a giant step forward, it is still just a step. There is more work ahead. Because the system facilitates access to information submitted to the Church over many decades from many sources, new FamilySearch may expose duplicate entries or errors that had not previously been recognized. This feature is especially useful for those with early pioneer ancestry. Duplications and errors need correction, and no one can do it better than each individual for his or her own family.

You may be frustrated as you work through these challenges. Please be assured that your concerns are well understood. The Church, under the inspired leadership of President Thomas S. Monson, is working diligently to assist you in solving these problems. Together we are striving to organize the family tree for all of God’s children. This is an enormous endeavor with enormous rewards.

This is joyful work. Look at this photograph of new converts doing temple work for their own kin. These dear Saints are from the San Salvador El Salvador Ilopango Stake and are attending the Guatemala City Temple for the first time. They are holding their temple cards, each with the name of a deceased relative for whom they have performed a vicarious baptism.

In order for the Church’s family history efforts to succeed, priesthood direction and leadership are essential. Leaders teach and testify of the doctrine undergirding this sacred work. They issue callings and see that instruction is available. They view temple and family history activity as a way to elevate the spirit of their ward, strengthen the spiritual roots of new converts, and bless the lives of all members.

While temple and family history work has the power to bless those beyond the veil, it has an equal power to bless the living. It has a refining influence on those who are engaged in it. They are literally helping to exalt their families.

We are exalted when we can dwell together with our extended families in the presence of Almighty God. The Prophet Joseph Smith foresaw our duty: “The great day of the Lord is at hand,” he said. “Let us, therefore, as a church and a people, and as Latter-day Saints, offer unto the Lord an offering in righteousness; and let us present in his holy temple … a book containing the records of our dead, which shall be worthy of all acceptation.”

The preparation of that record is our individual and collective responsibility. As we work together, we can make it worthy of all acceptation by the Lord. That record enables ordinances to be performed for and accepted by our deceased ancestors, as they may choose. Those ordinances can bring liberty to captives on the other side of the veil.

Our children, grandchildren, “Dear Ruby,” and all our great-grandchildren are linked in love. They are also linked in love to ancestors. Those links, welded through sacred ordinances, lead to the exaltation of our families. That this sacred goal may be realized for each of us is my prayer in the name of Jesus Christ, amen.

\newpage
\settalk{Be Thou an Example of the Believers}
\setspeaker{Russell M. Nelson}
\settalkdate{October 2010}
\talkheading{Be Thou an Example of the Believers}{Russell M. Nelson}

My beloved brethren, this evening we are assembled in many locations throughout the world. Wonderful full-time missionaries are among us. I would like to invite all full-time missionaries to stand. Wherever you are, elders and mission presidencies, please stand. We are grateful for each one of you! We thank you! We love you! You may be seated.

From time to time we need to remind ourselves why we have missionaries. It is because of a commandment from the Lord, who said:

“Go ye therefore, and teach all nations, baptizing them in the name of the Father, and of the Son, and of the Holy Ghost:

“Teaching them to observe all things whatsoever I have commanded you: and, lo, I am with you alway, even unto the end of the world.”

This commandment is one of many that have been renewed because the gospel of Jesus Christ has been restored in its fulness. Missionaries serve now just as they did in New Testament times. The book of Acts describes early missionary labors of the Apostles and other disciples following the Lord’s mortal ministry. There we read of the remarkable conversion and baptism of Saul of Tarsus, who had previously been “breathing out threatenings and slaughter against the disciples of the Lord” and persecuting members of the fledgling Church. From such beginnings, Saul became the converted Paul, one of the Lord’s greatest missionaries. The final 15 chapters of the book of Acts report the missionary labors of Paul and his companions.

In a letter to one of his most trusted companions, Paul wrote to young Timothy, “Let no man despise thy youth; but be thou an example of the believers, in word, in conversation, in charity, in spirit, in faith, in purity.” That counsel is just as valid for us now as it was then. It applies to our full-time missionaries; it applies equally to each member of the Church. Whether full-time missionaries or members, we should all be good examples of the believers in Jesus Christ.

\intalkheader{Full-Time Missionaries}

Full-time missionaries, some 52,000 and more, serve in 340 missions around the world. They are believers and devoted servants of the Lord. Their purpose is to “invite others to come unto Christ by helping them receive the restored gospel through faith in Jesus Christ and His Atonement, repentance, baptism, receiving the gift of the Holy Ghost, and enduring to the end.”

Like Timothy, most full-time missionaries are young men. Some are sisters; some are senior missionaries. We love each one! Missionaries serve to make life better for God’s children. Heavenly Father loves every one of His children. After all, He is their Father. He wants to bless them with His greatest gift, that of eternal life. Missionaries so teach wherever they serve. They help people to develop faith in the Lord, repent, be baptized, receive the Holy Ghost, receive the ordinances of the temple, and endure faithfully to the end. God’s work and glory—“to bring to pass the immortality and eternal life of man”—is also the sacred work and glory of each missionary.

We need more missionaries—more worthy missionaries. During His earthly ministry, the Lord told His disciples, “The harvest truly is great, but the labourers are few: pray ye therefore the Lord of the harvest, that he would send forth labourers into his harvest.”

At this morning’s session of general conference, our beloved President Thomas S. Monson made an impassioned plea for each young man of this Church to prepare for missionary service. I hope his entire message will be heeded in every home of the Church.

To President Monson’s wise counsel, I add my witness. In my family, I have observed the blessings that come to each missionary. Thus far, the number of our children, grandchildren, and their spouses called to serve as full-time missionaries is 49, and that number will continue to increase. In each instance, I have seen the growth in wisdom, maturity in judgment, and flowering of faith that developed in each missionary. They, just as many generations before them, embarked in the service of God to “serve him with all [their] heart, might, mind and strength.” Missionary service has helped to shape their divine destiny.

\intalkheader{Member Missionaries}

Paul’s counsel, “Be thou an example of the believers,” applies equally to members. Most have not been and may never be full-time missionaries. But all can be member missionaries. That statement reminds me of a report of a humorous event. On a large playing field at a missionary training center, a sign was posted. It read, “Missionaries Only!” People who also wanted to play on that field posted a new sign of their own. Their sign read, “Every Member a Missionary!”

Each member can be an example of the believers. Brethren, as followers of Jesus Christ, each of you can live in accord with His teachings. You can have “a pure heart and clean hands”; you can have “the image of God engraven upon your [countenance].” Your good works will be evident to others. The light of the Lord can beam from your eyes. With that radiance, you had better prepare for questions. The Apostle Peter so counseled, “Be ready always to give an answer to every man that asketh you a reason of the hope that is in you.”

Let your response be warm and joyful. And let your response be relevant to that individual. Remember, he or she is also a child of God, that very God who dearly wants that person to qualify for eternal life and return to Him one day. You may be the very one to open the door to his or her salvation and understanding of the doctrine of Christ.

After your initial response, be ready to take the next step. You may invite your friend to attend church with you. Many of our friends do not know they are welcome in our Church buildings. “Come and see” was the Savior’s invitation to those who desired to learn more about Him. An invitation to attend a Sunday meeting with you or to participate in a Church social or service activity will help to dispel mistaken myths and make visitors feel more comfortable among us.

As a member of the Church, reach out to those you do not know and greet them warmly. Each Sunday extend a hand of fellowship to at least one person you did not know before. Each day of your life, strive to enlarge your own circle of friendship.

You can invite a friend to read the Book of Mormon. Explain that it is not a novel or a history book. It is another testament of Jesus Christ. Its very purpose is “to the convincing of the Jew and Gentile that Jesus is the Christ, the Eternal God, manifesting himself unto all nations.” There is a power in this book that can touch the hearts and lift the lives of honest seekers of truth. Invite your friend to read the book prayerfully.

The Prophet Joseph Smith said “that the Book of Mormon was the most correct of any book on earth, and the keystone of our religion, and a man would get nearer to God by abiding by its precepts, than by any other book.” The Book of Mormon teaches of the Atonement of Jesus Christ and is the instrument by which God will fulfill His ancient promise to gather scattered Israel in these latter days.

Many years ago two colleagues of mine—a nurse and her doctor husband—asked me why I lived the way I did. I answered, “Because I know the Book of Mormon is true.” I let them borrow my copy of the book, inviting them to read it. A week later they returned my book with a polite “Thanks a lot.”

I responded, “What do you mean, ‘Thanks a lot’? That’s a totally inappropriate response for one who has read this book. You didn’t read it, did you! Please take it back and read it; then I would like my book back.”

Admitting that they had only turned its pages, they accepted my invitation. When they returned, they said tearfully, “We have read the Book of Mormon. We know it is true! We want to know more.” They learned more, and it was my privilege to baptize both of them.

Another way that you can share the gospel is to invite friends to meet with full-time missionaries in your home. Those missionaries are called and prepared to teach the gospel. Your friends, in the comfort of your home and with your constant reassurance, can begin their journey toward salvation and exaltation. The Lord said, “Ye are called to bring to pass the gathering of mine elect; for mine elect hear my voice and harden not their hearts.”

Scripture tells us that “there are many yet on the earth … who are only kept from the truth because they know not where to find it.” Isn’t that your opportunity? You can become their own disciple of discovery!

Now in this day of the Internet, there are new and exciting ways you can do missionary work. You can invite friends and neighbors to visit the new mormon.org website. If you have blogs and online social networks, you could link your sites to mormon.org. And there you can create your own personal profile. Each profile includes an expression of belief, an experience, and a testimony. Because this is a new feature, most of these profiles are available in English. Profiles in other languages will follow.

These profiles can have a profound influence for good. Two months ago a young man named Zac—a freshman in college—saw an ad for mormon.org on television in Baton Rouge, Louisiana. He connected with the website and was intrigued by the profiles of Church members. At our website he found the link that informed him where he could attend church. The next Sunday, dressed in a white shirt and tie, he attended church, was introduced to members of the ward, and enjoyed all three hours of meetings. He was invited to a member’s home for dinner, followed by his first missionary lesson. In less than two weeks, he was baptized and confirmed as a member of the Church. Welcome, Zac! (He is listening.)

Each exemplary follower of Jesus Christ can become an effective member missionary. Members and full-time missionaries may walk arm in arm in bringing the blessings of the gospel to cherished friends and neighbors. Many of them are of Israel, now being gathered as promised. This is all part of the preparation for the Second Coming of the Lord. He wants each of us truly to be an example of the believers.

I know that God lives. Jesus is the Christ. This is His Church. The Book of Mormon is true. Joseph Smith is its translator and the prophet of this last dispensation. President Thomas S. Monson is God’s prophet today. I so testify in the sacred name of Jesus Christ, amen.

\newpage
\settalk{Face the Future with Faith}
\setspeaker{Russell M. Nelson}
\settalkdate{April 2011}
\talkheading{Face the Future with Faith}{Russell M. Nelson}

My beloved brethren and sisters, thank you for your sustaining influence, not only by your uplifted hand but also by your uplifting service at home, in the Church, and in your communities. We love to be with you and see you among your families and friends. Wherever you live, we observe your efforts to make this world a better place. We sustain you! We love you! As you pray for us, so we pray for you!

We envision your families gathered around the television or online to watch the proceedings of general conference at home. An alert mother and father sent me a copy of a picture they took at conference time. They observed the reaction of their then 18-month-old son, who recognized the features and voice of the speaker. The child started to throw kisses toward the TV. He wanted to come closer. So his thoughtful older sister quickly hoisted her little brother on her shoulders and brought him closer. Here is that photograph.

Yes, the image on the TV is mine, and those children are our grandchildren. In a few years this boy will be an elder endowed in the temple and ready for his mission. Later he will be sealed to an eternal companion of his choice. Can you see him one day as a husband and father, with children of his own? And one day he will say farewell to his grandfathers, with a sure knowledge that death is part of life.

It is true. We live to die, and we die to live again. From an eternal perspective, the only death that is truly premature is the death of one who is not prepared to meet God.

As apostles and prophets, we are concerned not only for our children and grandchildren but for yours as well—and for each of God’s children. All that the future holds in store for each sacred child of God will be shaped by his or her parents, family, friends, and teachers. Thus, our faith now becomes part of our posterity’s faith later.

Each individual will make his or her way in a constantly changing world—a world of competing ideologies. The forces of evil will ever be in opposition to the forces of good. Satan constantly strives to influence us to follow his ways and make us miserable, even as he is. And the normal risks of life, such as illness, injury, and accident, will ever be present.

We live in a time of turmoil. Earthquakes and tsunamis wreak devastation, governments collapse, economic stresses are severe, the family is under attack, and divorce rates are rising. We have great cause for concern. But we do not need to let our fears displace our faith. We can combat those fears by strengthening our faith.

Start with your children. You parents bear the primary responsibility to strengthen their faith. Let them feel your faith, even when sore trials come upon you. Let your faith be focused on our loving Heavenly Father and His Beloved Son, the Lord Jesus Christ. Teach that faith with deep conviction. Teach each precious boy or girl that he or she is a child of God, created in His image, with a sacred purpose and potential. Each is born with challenges to overcome and faith to be developed.

Teach of faith in God’s plan of salvation. Teach that our sojourn in mortality is a period of probation, a time of trial and testing to see if we will do whatever the Lord commands us to do.

Teach of faith to keep all the commandments of God, knowing that they are given to bless His children and bring them joy. Warn them that they will encounter people who pick which commandments they will keep and ignore others that they choose to break. I call this the cafeteria approach to obedience. This practice of picking and choosing will not work. It will lead to misery. To prepare to meet God, one keeps all of His commandments. It takes faith to obey them, and keeping His commandments will strengthen that faith.

Obedience allows God’s blessings to flow without constraint. He will bless His obedient children with freedom from bondage and misery. And He will bless them with more light. For example, one keeps the Word of Wisdom knowing that obedience will not only bring freedom from addiction, but it will also add blessings of wisdom and treasures of knowledge.

Teach of faith to know that obedience to the commandments of God will provide physical and spiritual protection. And remember, God’s holy angels are ever on call to help us. The Lord so declared: “I will go before your face. I will be on your right hand and on your left, and my Spirit shall be in your hearts, and mine angels round about you, to bear you up.” What a promise! When we are faithful, He and His angels will help us.

Unfailing faith is fortified through prayer. Your heartfelt pleadings are important to Him. Think of the intense and impassioned prayers of the Prophet Joseph Smith during his dreadful days of incarceration in Liberty Jail. The Lord responded by changing the Prophet’s perspective. He said, “Know thou, my son, that all these things shall give thee experience, and shall be for thy good.”

If we pray with an eternal perspective, we need not wonder if our most tearful and heartfelt pleadings are heard. This promise from the Lord is recorded in section 98 of the Doctrine and Covenants:

“Your prayers have entered into the ears of the Lord … and are recorded with this seal and testament—the Lord hath sworn and decreed that they shall be granted.

“Therefore, he giveth this promise unto you, with an immutable covenant that they shall be fulfilled; and all things wherewith you have been afflicted shall work together for your good, and to my name’s glory, saith the Lord.”

The Lord chose His strongest words to reassure us! Seal! Testament! Sworn! Decreed! Immutable covenant! Brothers and sisters, believe Him! God will heed your sincere and heartfelt prayers, and your faith will be strengthened.

To develop enduring faith, an enduring commitment to be a full-tithe payer is essential. Initially it takes faith to tithe. Then the tithe payer develops more faith to the point that tithing becomes a precious privilege. Tithing is an ancient law from God. He made a promise to His children that He would open “the windows of heaven, and pour … out a blessing, that there shall not be room enough to receive it.” Not only that, tithing will keep your name enrolled among the people of God and protect you in “the day of vengeance and burning.”

Why do we need such resilient faith? Because difficult days are ahead. Rarely in the future will it be easy or popular to be a faithful Latter-day Saint. Each of us will be tested. The Apostle Paul warned that in the latter days, those who diligently follow the Lord “shall suffer persecution.” That very persecution can either crush you into silent weakness or motivate you to be more exemplary and courageous in your daily lives.

How you deal with life’s trials is part of the development of your faith. Strength comes when you remember that you have a divine nature, an inheritance of infinite worth. The Lord has reminded you, your children, and your grandchildren that you are lawful heirs, that you have been reserved in heaven for your specific time and place to be born, to grow and become His standard bearers and covenant people. As you walk in the Lord’s path of righteousness, you will be blessed to continue in His goodness and be a light and a savior unto His people.

Available to each of you brethren and sisters are blessings obtained through the power of the holy Melchizedek Priesthood. These blessings can change the circumstances of your lives, in matters such as health, companionship of the Holy Ghost, personal relationships, and opportunities for the future. The power and authority of this priesthood holds the keys to all spiritual blessings of the Church. And most remarkably, the Lord has declared that He will sustain those blessings, according to His will.

The greatest of all the blessings of the priesthood are bestowed in holy temples of the Lord. Fidelity to covenants made there will qualify you and your family for the blessings of eternal life.

Your rewards come not only hereafter. Many blessings will be yours in this life, among your children and grandchildren. You faithful Saints do not have to fight life’s battles alone. Think of that! The Lord declared, “I will contend with him that contendeth with thee, and I will save thy children.” Later came this promise to His faithful people: “I, the Lord, would fight their battles, and their children’s battles, and their children’s children’s, … to the third and fourth generation.”

Our beloved President Thomas S. Monson has given us his prophetic witness. He said: “I testify to you that our promised blessings are beyond measure. Though the storm clouds may gather, though the rains may pour down upon us, our knowledge of the gospel and our love of our Heavenly Father and of our Savior will comfort and sustain us and bring joy to our hearts as we walk uprightly and keep the commandments.”

President Monson continued: “My beloved brothers and sisters, fear not. Be of good cheer. The future is as bright as your faith.”

To President Monson’s powerful declaration I add my own. I testify that God is our Father. Jesus is the Christ. His Church has been restored to the earth. His truth, covenants, and ordinances enable us to overcome fear and face the future with faith! I so testify in the sacred name of Jesus Christ, amen.

\newpage
\settalk{Covenants}
\setspeaker{Russell M. Nelson}
\settalkdate{October 2011}
\talkheading{Covenants}{Russell M. Nelson}

One week after a recent assignment to create the first stake in Moscow, Russia, I attended a district conference in St. Petersburg. While speaking about my gratitude for early missionaries and local leaders who brought strength to the Church in Russia, I mentioned the name of Vyacheslav Efimov. He was the first Russian convert to become a mission president. He and his wife did wonderfully well in that assignment. Not long after they had completed their mission, and much to our sorrow, President Efimov suddenly passed away. He was only 52 years of age.

While speaking of this pioneering couple, I felt impressed to ask the congregation if Sister Efimov might be present. Far in the rear of the room, a woman stood. I invited her to come to the microphone. Yes, it was Sister Galina Efimov. She spoke with conviction and bore a powerful testimony of the Lord, of His gospel, and of His restored Church. She and her husband had been sealed in the holy temple. She said they were united forever. They were still missionary companions, she on this side of the veil and he on the other side. With tears of joy, she thanked God for sacred temple covenants. I wept too, with full realization that the everlasting unity exemplified by this faithful couple was the righteous result of making, keeping, and honoring sacred covenants.

One of the most important concepts of revealed religion is that of a sacred covenant. In legal language, a covenant generally denotes an agreement between two or more parties. But in a religious context, a covenant is much more significant. It is a sacred promise with God. He fixes the terms. Each person may choose to accept those terms. If one accepts the terms of the covenant and obeys God’s law, he or she receives the blessings associated with the covenant. We know that “when we obtain any blessing from God, it is by obedience to that law upon which it is predicated.”

Through the ages, God has made covenants with His children. His covenants occur throughout the entire plan of salvation and are therefore part of the fulness of His gospel. For example, God promised to send a Savior for His children, asking in turn for their obedience to His law.

In the Bible we read of men and women in the Old World who were identified as children of the covenant. What covenant? “The covenant which God made with [their] fathers, saying unto Abraham, And in thy seed shall all the kindreds of the earth be blessed.”

In the Book of Mormon we read of people in the New World who were also identified as children of the covenant. The resurrected Lord so informed them: “Behold, ye are the children of the prophets; and ye are of the house of Israel; and ye are of the covenant which the Father made with your fathers, saying unto Abraham: And in thy seed shall all the kindreds of the earth be blessed.”

The Savior explained the importance of their identity as children of the covenant. He said, “The Father having raised me up unto you first, … sent me to bless you in turning away every one of you from his iniquities; and this because ye are the children of the covenant.”

The covenant God made with Abraham and later reaffirmed with Isaac and Jacob is of transcendent significance. It contained several promises, including:

Some of these promises have been fulfilled; others are still pending. I quote from an early Book of Mormon prophecy: “Our father [Lehi] hath not spoken of our seed alone, but also of all the house of Israel, pointing to the covenant which should be fulfilled in the latter days; which covenant the Lord made to our father Abraham.” Isn’t that amazing? Some 600 years before Jesus was born in Bethlehem, prophets knew that the Abrahamic covenant would be finally fulfilled only in the latter days.

To facilitate that promise, the Lord appeared in these latter days to renew that Abrahamic covenant. To the Prophet Joseph Smith, the Master declared:

“Abraham received promises concerning his seed, and of the fruit of his loins—from whose loins ye are, … my servant Joseph. …

“This promise is yours also, because ye are of Abraham.”

With this renewal, we have received, as did they of old, the holy priesthood and the everlasting gospel. We have the right to receive the fulness of the gospel, enjoy the blessings of the priesthood, and qualify for God’s greatest blessing—that of eternal life.

Some of us are the literal seed of Abraham; others are gathered into his family by adoption. The Lord makes no distinction. Together we receive these promised blessings—if we seek the Lord and obey His commandments. But if we don’t, we lose the blessings of the covenant. To assist us, His Church provides patriarchal blessings to give each recipient a vision for his or her future as well as a connection with the past, even a declaration of lineage back to Abraham, Isaac, and Jacob.

Brethren of the covenant have the right to qualify for the oath and covenant of the priesthood. If you are “faithful unto the obtaining these two priesthoods … and the magnifying [of your] calling, [you] are sanctified by the Spirit unto the renewing of [your] bodies.” That is not all. Men who worthily receive the priesthood receive the Lord Jesus Christ, and those who receive the Lord receive God the Father. And those who receive the Father receive all that He has. Incredible blessings flow from this oath and covenant to worthy men, women, and children in all the world.

Ours is the responsibility to help fulfill the Abrahamic covenant. Ours is the seed foreordained and prepared to bless all people of the world. That is why priesthood duty includes missionary work. After some 4,000 years of anticipation and preparation, this is the appointed day when the gospel is to be taken to the kindreds of the earth. This is the time of the promised gathering of Israel. And we get to participate! Isn’t that exciting? The Lord is counting on us and our sons—and He is profoundly grateful for our daughters—who worthily serve as missionaries in this great time of the gathering of Israel.

The Book of Mormon is a tangible sign that the Lord has commenced to gather His children of covenant Israel. This book, written for our day, states as one of its purposes that “ye may know that the covenant which the Father hath made with the children of Israel … is already beginning to be fulfilled. … For behold, the Lord will remember his covenant which he hath made unto his people of the house of Israel.”

Indeed, the Lord has not forgotten! He has blessed us and others throughout the world with the Book of Mormon. One of its purposes is for “the convincing of the Jew and Gentile that Jesus is the Christ.” It helps us to make covenants with God. It invites us to remember Him and to know His Beloved Son. It is another testament of Jesus Christ.

Children of the covenant have the right to receive His doctrine and to know the plan of salvation. They claim it by making covenants of sacred significance. Brigham Young said: “All Latter-day Saints enter the new and everlasting covenant when they enter this Church. … They enter the new and everlasting covenant to sustain the Kingdom of God.” They keep the covenant by obedience to His commandments.

At baptism we covenant to serve the Lord and keep His commandments. When we partake of the sacrament, we renew that covenant and declare our willingness to take upon ourselves the name of Jesus Christ. Thereby we are adopted as His sons and daughters and are known as brothers and sisters. He is the father of our new life. Ultimately, in the holy temple, we may become joint heirs to the blessings of an eternal family, as once promised to Abraham, Isaac, Jacob, and their posterity. Thus, celestial marriage is the covenant of exaltation.

When we realize that we are children of the covenant, we know who we are and what God expects of us. His law is written in our hearts. He is our God and we are His people. Committed children of the covenant remain steadfast, even in the midst of adversity. When that doctrine is deeply implanted in our hearts, even the sting of death is soothed and our spiritual stamina is strengthened.

The greatest compliment that can be earned here in this life is to be known as a covenant keeper. The rewards for a covenant keeper will be realized both here and hereafter. Scripture declares that “ye should consider on the blessed and happy state of those that keep the commandments of God. For behold, they are blessed in all things, … and if they hold out faithful to the end they are received into heaven … [and] dwell with God in a state of never-ending happiness.”

God lives. Jesus is the Christ. His Church has been restored to bless all people. President Thomas S. Monson is His prophet today. And we, as faithful children of the covenant, will be blessed now and forever. I so testify in the name of Jesus Christ, amen.

\newpage
\settalk{Thanks Be to God}
\setspeaker{Russell M. Nelson}
\settalkdate{April 2012}
\talkheading{Thanks Be to God}{Russell M. Nelson}

Dear brothers and sisters, we thank you for your sustaining support and your devotion. We express our gratitude and love to each of you.

Recently, Sister Nelson and I enjoyed the beauty of tropical fish in a small private aquarium. Fish with vivid colors and of a variety of shapes and sizes darted back and forth. I asked the attendant nearby, “Who provides food for these beautiful fish?”

She responded, “I do.”

Then I asked, “Have they ever thanked you?”

She replied, “Not yet!”

I thought of some people I know who are just as oblivious to their Creator and their true “bread of life.” They live from day to day without an awareness of God and His goodness unto them.

How much better it would be if all could be more aware of God’s providence and love and express that gratitude to Him. Ammon taught, “Let us give thanks to [God], for he doth work righteousness forever.” Our degree of gratitude is a measure of our love for Him.

God is the Father of our spirits. He has a glorified, perfected body of flesh and bone. We lived with Him in heaven before we were born. And when He created us physically, we were created in the image of God, each with a personal body.

Think of our physical sustenance. It is truly heaven-sent. The necessities of air, food, and water all come to us as gifts from a loving Heavenly Father. The earth was created to support our brief sojourn in mortality. We were born with a capacity to grow, love, marry, and form families.

Marriage and family are ordained of God. The family is the most important social unit in time and in eternity. Under God’s great plan of happiness, families can be sealed in temples and be prepared to return to dwell in His holy presence forever. That is eternal life! It fulfills the deepest longings of the human soul—the natural yearning for endless association with beloved members of one’s family.

We are part of His divine purpose: “My work and my glory,” He said, is “to bring to pass the immortality and eternal life of man.” In order to achieve those objectives, “God so loved the world, that he gave his only begotten Son, that whosoever believeth in him should not perish, but have everlasting life.” That act was a supernal manifestation of God’s love. “For [He] sent not his Son into the world to condemn the world; but that the world through him might be saved.”

Central to God’s eternal plan is the mission of His Son, Jesus Christ. He came to redeem God’s children. Because of the Lord’s Atonement, resurrection (or immortality) became a reality. Because of the Atonement, eternal life became a possibility for all who would qualify. Jesus so explained:

“I am the resurrection, and the life: he that believeth in me, though he were dead, yet shall he live:

“And whosoever liveth and believeth in me shall never die.”

For the Atonement of the Lord and His gift of resurrection—for this sublime message of Easter—thanks be to God!

\intalkheader{Physical Gifts}

Our Heavenly Father loves His children. He has blessed each with physical and spiritual gifts. Let me speak of each type. When you sing “I Am a Child of God,” think of His gift to you of your own physical body. The many amazing attributes of your body attest to your own “divine nature.”

Each organ of your body is a wondrous gift from God. Each eye has an autofocusing lens. Nerves and muscles control two eyes to make a single three-dimensional image. The eyes are connected to the brain, which records the sights seen.

Your heart is an incredible pump. It has four delicate valves that control the direction of blood flow. These valves open and close more than 100,000 times a day—36 million times a year. Yet, unless altered by disease, they are able to withstand such stress almost indefinitely.

Think of the body’s defense system. To protect it from harm, it perceives pain. In response to infection, it generates antibodies. The skin provides protection. It warns against injury that excessive heat or cold might cause.

The body renews its own outdated cells and regulates the levels of its own vital ingredients. The body heals its cuts, bruises, and broken bones. Its capacity for reproduction is another sacred gift from God.

Be we reminded that a perfect body is not required to achieve one’s divine destiny. In fact, some of the sweetest spirits are housed in frail or imperfect bodies. Great spiritual strength is often developed by people with physical challenges, precisely because they are so challenged.

Anyone who studies the workings of the human body has surely “seen God moving in his majesty and power.” Because the body is governed by divine law, any healing comes by obedience to the law upon which that blessing is predicated.

Yet some people erroneously think that these marvelous physical attributes happened by chance or resulted from a big bang somewhere. Ask yourself, “Could an explosion in a printing shop produce a dictionary?” The likelihood is most remote. But if so, it could never heal its own torn pages or reproduce its own newer editions!

If the body’s capacity for normal function, defense, repair, regulation, and regeneration were to prevail without limit, life here would continue in perpetuity. Yes, we would be stranded here on earth! Mercifully for us, our Creator provided for aging and other processes that would ultimately result in our physical death. Death, like birth, is part of life. Scripture teaches that “it was not expedient that man should be reclaimed from this temporal death, for that would destroy the great plan of happiness.” To return to God through the gateway we call death is a joy for those who love Him and are prepared to meet Him. Eventually the time will come when each “spirit and … body shall be reunited again in … perfect form; both limb and joint shall be restored to its proper frame,” never to be separated again. For these physical gifts, thanks be to God!

\intalkheader{Spiritual Gifts}

Important as is the body, it serves as a tabernacle for one’s eternal spirit. Our spirits existed in the premortal realm and will continue to live after the body dies. The spirit provides the body with animation and personality. In this life and in the next, spirit and body, when joined together, become a living soul of supernal worth.

Because one’s spirit is so important, its development is of eternal consequence. It is strengthened as we communicate in humble prayer with our loving Heavenly Father.

The attributes by which we shall be judged one day are all spiritual. These include love, virtue, integrity, compassion, and service to others. Your spirit, coupled with and housed in your body, is able to develop and manifest these attributes in ways that are vital to your eternal progression. Spiritual progress is attained through the steps of faith, repentance, baptism, the gift of the Holy Ghost, and enduring to the end, including the endowment and sealing ordinances of the holy temple.

Just as the body requires daily food for survival, the spirit needs nourishment as well. The spirit is nurtured by eternal truth. Last year we celebrated the 400th anniversary of the King James translation of the Holy Bible. And we have had the Book of Mormon for nearly 200 years. It has now been translated in its entirety or as selections into 107 languages. Because of these and other precious scriptures, we know that God is our Eternal Father and that His Son, Jesus Christ, is our Savior and Redeemer. For these spiritual gifts, thanks be to God!

\intalkheader{Gospel Gifts}

We know that prophets of many dispensations, such as Adam, Noah, Moses, and Abraham, all taught of the divinity of our Heavenly Father and of Jesus Christ. Our present dispensation was introduced by Heavenly Father and Jesus Christ when They appeared to the Prophet Joseph Smith in 1820. The Church was organized in 1830. Now, 182 years later, we remain under covenant to take the gospel to “every nation, kindred, tongue, and people.” As we do so, both givers and receivers will be blessed.

Ours is the responsibility to teach His children and awaken in them an awareness of God. Long ago, King Benjamin said:

“Believe in God; believe that he is, and that he created all things, both in heaven and in earth; believe that he has all wisdom, and all power, both in heaven and in earth; …

“… Believe that ye must repent of your sins and forsake them, and humble yourselves before God; and ask in sincerity of heart that he would forgive you; and now, if you believe all these things see that ye do them.”

God is the same yesterday, today, and forever, but we are not. Each day, ours is the challenge to access the power of the Atonement so that we can truly change, become more Christlike, and qualify for the gift of exaltation and live eternally with God, Jesus Christ, and our families. For these powers, privileges, and gospel gifts, thanks be to God!

I testify that He lives, that Jesus is the Christ, and that this is His Church, restored in these latter days to accomplish its divine destiny. We are led today by President Thomas S. Monson, whom we love and sustain with all our hearts, as we also sustain his counselors and the Twelve Apostles as prophets, seers, and revelators. I so testify in the sacred name of Jesus Christ, amen.

\newpage
\settalk{Ask the Missionaries! They Can Help You!}
\setspeaker{Russell M. Nelson}
\settalkdate{October 2012}
\talkheading{Ask the Missionaries! They Can Help You!}{Russell M. Nelson}

My beloved brothers, sisters, and friends, we extend our love and greetings to each of you. We are thrilled with President Thomas S. Monson’s announcement this morning, which adjusts the minimum age for missionary service to 18 for young men and 19 for young women. Through this option, more of our youth may enjoy the blessings of a mission.

Two years ago and powerfully reaffirmed again this morning, President Monson declared “that every worthy, able young man should prepare to serve a mission. Missionary service is a priesthood duty—an obligation the Lord expects of us who have been given so very much.” Again he explained that for young sisters, a mission is a welcome option but not a responsibility. And again he invited many more mature couples to serve.

Preparation for a mission is important. A mission is a voluntary act of service to God and humankind. Missionaries support that privilege with their personal savings. Parents, families, friends, and donors to the General Missionary Fund may also assist. All missionaries, younger and older, serve with the sole hope of making life better for other people.

The decision to serve a mission will shape the spiritual destiny of the missionary, his or her spouse, and their posterity for generations to come. A desire to serve is a natural outcome of one’s conversion, worthiness, and preparation.

In this great worldwide audience, many of you are not affiliated with The Church of Jesus Christ of Latter-day Saints and know very little about us and our missionaries. You are here or tuned in because you want to know more about the Mormons and what our missionaries teach. As you learn more about us, you will find that we share many of the same values. We encourage you to keep all that is good and true and then see if we can add more. In this world filled with challenges, we do need help from time to time. Religion, eternal truth, and our missionaries are vital parts of that help.

Our young missionaries set aside their education, occupation, dating, and whatever else young adults would typically be doing at this stage of life. For 18 to 24 months they put it all on hold because of their deep desire to serve the Lord. And some of our missionaries serve in their more mature years of life. I know their families are blessed. In our own family, eight are currently serving as full-time missionaries—three daughters, their husbands, one granddaughter, and one grandson.

Some of you may wonder about the name Mormon. It is a nickname for us. It is not our real name, though we are widely known as Mormons. The term is derived from a book of sacred scripture known as the Book of Mormon.

The true name of the Church is The Church of Jesus Christ of Latter-day Saints. It is the reestablished original Church of Jesus Christ. When He walked upon the earth, He organized His Church. He called Apostles, Seventies, and other leaders to whom He gave priesthood authority to act in His name. After Christ and His Apostles passed away, men changed the ordinances and doctrine. The original Church and the priesthood were lost. After the Dark Ages, and under the direction of Heavenly Father, Jesus Christ brought back His Church. Now it lives again, restored and functioning under His divine direction.

We follow the Lord Jesus Christ and teach of Him. We know that after His glorious triumph over death, the resurrected Lord appeared to His disciples on numerous occasions. He ate with them. He walked with them. Before His final Ascension, He commissioned them to “go … and teach all nations, baptizing them in the name of the Father, and of the Son, and of the Holy Ghost.” The Apostles heeded that instruction. They also called upon others to help them fulfill the Lord’s command.

Today, under the direction of modern apostles and prophets, that same charge has been extended to missionaries of The Church of Jesus Christ of Latter-day Saints. These missionaries serve in more than 150 nations. As representatives of the Lord Jesus Christ, they strive to fulfill that divine command—renewed in our day by the Lord Himself—to take the fulness of the gospel abroad and bless the lives of people everywhere.

Missionaries in their late teens or early 20s are young in ways of the world. But they are blessed with gifts—such as the power of the Holy Spirit, the love of God, and testimonies of the truth—that make them powerful ambassadors of the Lord. They share the good news of the gospel that will bring true joy and everlasting happiness to all who heed their message. And in many instances they do so in a country and a language foreign to them.

Missionaries strive to follow Jesus Christ in both word and deed. They preach of Jesus Christ and of His Atonement. They teach of the literal Restoration of Christ’s ancient Church through the Lord’s first latter-day prophet, Joseph Smith.

You may have previously encountered, or even ignored, our missionaries. My hope is that you will not fear them but learn from them. They can be a heaven-sent resource to you.

That happened to Jerry, a Protestant gentleman in his mid-60s who lives in Mesa, Arizona. Jerry’s father was a Baptist minister; his mother, a Methodist minister. One day Jerry’s close friend Pricilla shared with him the pain she felt from the death of her child during childbirth and a bitter divorce that occurred shortly thereafter. Struggling as a single mother, Pricilla has four children—three daughters and a son. As she opened her heart to Jerry, she confessed that she was thinking of taking her own life. With all the strength and love Jerry could muster, he tried to help her understand that her life had value. He invited her to attend his church, but Pricilla explained that she had given up on God.

Jerry did not know what to do. Later, while watering trees in his yard, this man of faith prayed to God for guidance. As he prayed, he heard a voice in his mind saying, “Stop the boys on the bikes.” Jerry, a little bewildered, wondered what this meant. As he reflected on this impression, he gazed up the street and saw two young men in white shirts and ties riding bicycles toward his home. Stunned by this “coincidence,” he watched them ride by. Then, realizing that the situation required him to act, he shouted out, “Hey, you, please stop! I need to talk to you!”

With a puzzled but excited look, the young men stopped. As they approached, Jerry noticed that they wore name tags identifying them as missionaries in The Church of Jesus Christ of Latter-day Saints. Jerry looked at them and said, “This may sound a little weird, but I was praying and was told to ‘stop the boys on the bikes.’ I looked up the street, and here you are. Can you help me?”

The missionaries smiled, and one said, “Yes, I am sure we can.”

Jerry explained the worrisome plight of Pricilla. Soon the missionaries were meeting with Pricilla, her children, and Jerry. They discussed the purpose of life and God’s eternal plan for them. Jerry, Pricilla, and her children grew in faith through sincere prayer, their study of the Book of Mormon, and the loving fellowship with members of the Church. Jerry’s already strong faith in Jesus Christ grew even stronger. Pricilla’s doubts and thoughts of suicide turned to hope and happiness. They were baptized and became members of Christ’s restored Church.

Yes, missionaries can help in many ways. For example, some of you might want to know more about your ancestors. You may know the names of your parents and your four grandparents, but what about your eight great-grandparents? Do you know their names? Would you like to know more about them? Ask the missionaries! They can help you! They have ready access to the vast family history records of The Church of Jesus Christ of Latter-day Saints.

Some of you are members but not presently participating. You love the Lord and often think of returning to His fold. But you don’t know how to start. I suggest that you ask the missionaries! They can help you! They can also help by teaching your loved ones. We and the missionaries love you and desire to bring joy and the light of the gospel back into your lives.

Some of you may want to know how to conquer an addiction or live longer and enjoy better health. Ask the missionaries! They can help you! Independent studies have shown that, as a group, members of The Church of Jesus Christ of Latter-day Saints are a healthy lot. Their death rates are among the lowest and their longevity greater than any yet reported in any well-defined group studied over a lengthy period of time in the United States.

Some of you may feel that life is busy and frenetic, yet down deep in your heart you feel a gnawing emptiness, without direction or purpose. Ask the missionaries! They can help you! They can help you to learn more about the true purpose of life—why you are here on earth and where you are going after death. You can learn how the restored gospel of Jesus Christ will bless your life beyond anything you can presently even imagine.

If you have concerns about your family, ask the missionaries! They can help you! Strengthening marriages and families is of utmost importance to Latter-day Saints. Families can be together forever. Ask the missionaries to teach you how this is possible for your family.

Missionaries can also help you with your desire for greater knowledge. The human spirit yearns for enlightenment. Whether truth comes from a scientific laboratory or by revelation from God, we seek it! The glory of God indeed is intelligence.

Increase in learning includes spiritual as well as temporal knowledge. We stress the importance of understanding sacred scriptures. An independent study recently found that Latter-day Saints were the most knowledgeable about Christianity and the Bible. If you want to understand the Bible better, to understand the Book of Mormon better, and gain a broader comprehension of the brotherhood of man and the fatherhood of God, ask the missionaries! They can help you!

Many of you have a deep desire to help people in need. Because we follow Jesus Christ, Latter-day Saints are also compelled by that insatiable urge. Anyone may join with us to help the needy and provide relief to victims of disaster anywhere in the world. If you want to participate, ask the missionaries! They can help you!

And if you want to know more about life after death, about heaven, about God’s plan for you; if you want to know more about the Lord Jesus Christ, His Atonement, and the Restoration of His Church as it was originally established, ask the missionaries! They can help you!

I know that God lives. Jesus is the Christ. His Church has been restored. Fervently I pray that God may bless each of you and each of our precious missionaries. In the name of Jesus Christ, amen.

\newpage
\settalk{Catch the Wave}
\setspeaker{Russell M. Nelson}
\settalkdate{April 2013}
\talkheading{Catch the Wave}{Russell M. Nelson}

Dear brothers and sisters, I add my voice to that of President Thomas S. Monson and others in praising those who have responded to a prophet’s call for more worthy missionaries. Now an unprecedented wave of enthusiasm for missionary work is sweeping the entire earth. Since President Monson’s historic announcement last October, thousands of elders, sisters, and couples have been called, and many more are preparing. Now we get questions like “What are you going to do with all these missionaries?” The answer is simple. They will do what missionaries have always done. They will preach the gospel! They will bless the children of Almighty God!

More of you young men and women will catch this wave as you strive to be worthy of mission calls. You see this as a wave of truth and righteousness. You see your opportunity to be on the crest of that wave.

You teenagers, embrace your new curriculum and teach one another the doctrine of Jesus Christ. Now is your time to prepare to teach others about the goodness of God.

Young men and women, your education is ever important—to us, to you, and to God. Where feasible, if you wish to attend a college or university after your mission, we encourage you to apply to your institution of choice before beginning your mission. Many institutions of higher learning will grant an 18- to 30-month deferral to prospective missionaries. This will enable you elders and sisters to serve without worrying about where you will begin your advanced education. We are very grateful to leaders of educational institutions who are making such planning possible!

You parents, teachers, and others, catch the wave as you prepare our rising generation to be worthy of missionary service. Meanwhile, your exemplary lives will attract the interest of your friends and neighbors. Be ready to give an answer to those who ask why you live as you do. Be ready to give a reason for the hope and joy that they see in you. When such questions come, you might respond by saying, “Let’s ask the missionaries! They can help us! And if you desire, I will be at your side as the missionaries respond and teach you.”

You adults, catch the wave with help for the spiritual, physical, and financial preparation of future missionaries. Pinching pennies for piggy banks becomes part of your practice. You senior couples, you plan for the day when you can go on your mission. We will be most grateful for your service. Until then, perhaps some of you could send your dollars on missions by contributing to the General Missionary Fund, as suggested by President Monson again this morning.

Increasing numbers of selected men and their dear companions catch the wave as they are called to preside over missions of the Church. In that service they will shape the destiny of generations born and unborn. Mission presidents hold keys of responsibility for the welfare, safety, and success of their missionaries. After consulting with stake and district presidents in his mission, each mission president assigns missionaries to serve in specific stakes, wards, and branches.

Stake presidents and bishops catch the wave as they spend more and more hours interviewing prospective missionaries. These priesthood leaders hold keys of responsibility for missionary work in their units, and they inspire members to participate.

Brothers and sisters on each ward council are beginning to catch the wave. On that council is the ward mission leader. I would like to speak specifically to each of you ward mission leaders. You have been called by your bishop to lead missionary work in the ward. And some of you are so successful that an assistant has also been called to help you. Along with others on the ward council, you identify less-active members, part-member families, and interested neighbors. You meet regularly with assigned full-time missionaries. You advise and assist the missionaries. Please help them to fill their daily planners with focused and meaningful teaching opportunities. This is your responsibility. Your role is crucial, truly crucial to the success of this work. If you catch the wave with faith and enthusiasm, others will also. You, as the ward mission leader, are the connecting link between members and missionaries in this sacred work of rescuing God’s children.

Our inquiring friends and neighbors not of our faith can also catch the wave. We encourage them to keep all that is good and true in their lives. And we invite them to receive more, especially the glorious truth that through God’s eternal plan, families can be together forever.

This wave of truth and righteousness is wondrous! It is not man-made! It comes from the Lord, who said, “I will hasten my work in its time.” This wave is empowered by a divine announcement made 193 years ago. It consisted of only seven words: “This is My Beloved Son. Hear Him!” Uttered by Almighty God, that announcement introduced a young Joseph Smith to the Lord Jesus Christ. Those seven words launched the Restoration of His gospel. Why? Because our living God is a loving God! He wants His children to know Him and Jesus Christ, whom He has sent! And He wants His children to gain immortality and eternal life!

For this glorious purpose, our missionaries teach of the Restoration. They know that some 2,000 years ago, the Lord established His Church. After His Crucifixion and the death of His Apostles, men changed the Church and its doctrine. Then, after generations of spiritual darkness, and as predicted by previous prophets, Heavenly Father and Jesus Christ restored the Church, its doctrine, and its priesthood authority. Because of that Restoration, knowledge and essential ordinances for salvation and exaltation are again available to all people. Ultimately, that exaltation allows each of us to dwell with our families in the presence of God and Jesus Christ forever!

I cannot speak of the Restoration in tempered tones. This fact of history is absolutely stunning! It is incredible! It is breathtaking! How amazing is it that messengers from heaven came to give authority and power to this work?

Our Eternal Father and Jesus Christ made multiple appearances to the Prophet Joseph Smith. Under Their direction, other heavenly messengers came, each with a specific purpose. For example:

In addition, the Restoration added knowledge that Saints had anciently. The Lord provided a new book of scripture. To the Holy Bible, He added the Book of Mormon: Another Testament of Jesus Christ. It is a record of prophecies and the ministry of the resurrected Lord to the people of ancient America. It explains God’s great plan of happiness—the plan of salvation. The Book of Mormon is in full harmony with the Bible. Both sacred records affirm the truth of the gospel of Jesus Christ and the importance of His Atonement.

The Restoration fulfills many biblical prophecies. For example, Isaiah prophesied that the Lord’s house shall be established in the top of the mountains. The exodus of Mormon pioneers to the mountains of western America is a fulfilling saga of sacrifice and faith. Isaiah also foretold that God will do “a marvellous work and a wonder.” That is now being fulfilled by the sacred work of our growing army of missionaries.

Old Testament teachings about tithing have been restored. As a result, more tithe payers are blessed because of their obedience. References to Melchizedek are clarified by scriptures of the Restoration. Prophecies that the stick of Joseph (the Book of Mormon) and the stick of Judah (the Bible) would be as one in God’s hand have now been fulfilled.

The Restoration also clarifies New Testament scriptures. Its reference to baptism for the dead is now better understood. Ordinances for our deceased ancestors are now performed vicariously in 141 temples throughout the world! There is no other way to offer salvation to our ancestors who died without a knowledge of the gospel! John the Revelator’s vision of “another angel [flying] in the midst of heaven, having the everlasting gospel to preach unto them that dwell on the earth,” foretold the mission of the angel Moroni and the Book of Mormon.

The Book of Mormon is the centerpiece of the Restoration. It was written, preserved, and transmitted under the Lord’s direction. It was translated “by the gift and power of God.” The book of Doctrine and Covenants contains many additional revelations given to the Prophet Joseph Smith. Through him we have received more pages of scripture than we have from any other prophet. In what must have been a wistful moment, he said to the Saints in Nauvoo, Illinois, “I never told you I was perfect; but there is no error in the revelations which I have taught.”

Together, members and missionaries invite all to learn of God, of Jesus Christ, and of His gospel. Each inquiring individual should seek earnestly and pray fervently for the assurance that these things are true. The truth will be manifest by the power of the Holy Ghost.

I thank God and His Son, Jesus Christ, for the Restoration and its power to propel a magnificent wave of truth and righteousness across the earth. May we catch this wave and fulfill the Lord’s command to take the gospel “to every nation, and kindred, and tongue, and people,” I pray in the name of Jesus Christ, amen.

\newpage
\settalk{Decisions for Eternity}
\setspeaker{Russell M. Nelson}
\settalkdate{October 2013}
\talkheading{Decisions for Eternity}{Russell M. Nelson}

My dear brothers and sisters, each day is a day of decision. President Thomas S. Monson has taught us that “decisions determine destiny.” The wise use of your freedom to make your own decisions is crucial to your spiritual growth, now and for eternity. You are never too young to learn, never too old to change. Your yearnings to learn and change come from a divinely instilled striving for eternal progression. Each day brings opportunity for decisions for eternity.

We are eternal beings—spirit children of heavenly parents. The Bible records that “God created man in his own image, … male and female created he them.” Recently I heard a chorus of children sing the beloved song “I Am a Child of God.” I wondered, “Why haven’t I heard that song rendered more often by singing mothers or faithful fathers?” Are we not all children of God? In truth, not one of us can ever stop being a child of God!

As children of God, we should love Him with all our heart and soul, even more than we love our earthly parents. We should love our neighbors as brothers and sisters. No other commandments are greater than these. And we should ever revere the worth of human life, through each of its many stages.

Scripture teaches that the body and the spirit are the soul of man. As a dual being, each of you can thank God for His priceless gifts of your body and your spirit.

\intalkheader{The Human Body}

My professional years as a medical doctor gave me a profound respect for the human body. Created by God as a gift to you, it is absolutely amazing! Think of your eyes that see, ears that hear, and fingers that feel all the wondrous things around you. Your brain lets you learn, think, and reason. Your heart pumps tirelessly day and night, almost without your awareness.

Your body protects itself. Pain comes as a warning that something is wrong and needs attention. Infectious illnesses strike from time to time, and when they do, antibodies are formed that increase your resistence to subsequent infection.

Your body repairs itself. Cuts and bruises heal. Broken bones can become strong once again. I have cited but a tiny sample of the many amazing God-given qualities of your body.

Even so, it seems that in every family, if not in every person, some physical conditions exist that require special care. A pattern for coping with such a challenge has been given by the Lord. He said, “I give unto men weakness that they may be humble; … for if they humble themselves … and have faith in me, then will I make weak things become strong unto them.”

Stellar spirits are often housed in imperfect bodies. The gift of such a body can actually strengthen a family as parents and siblings willingly build their lives around that child born with special needs.

The aging process is also a gift from God, as is death. The eventual death of your mortal body is essential to God’s great plan of happiness. Why? Because death will allow your spirit to return home to Him. From an eternal perspective, death is only premature for those who are not prepared to meet God.

With your body being such a vital part of God’s eternal plan, it is little wonder that the Apostle Paul described it as a “temple of God.” Each time you look in the mirror, see your body as your temple. That truth—refreshed gratefully each day—can positively influence your decisions about how you will care for your body and how you will use it. And those decisions will determine your destiny. How could this be? Because your body is the temple for your spirit. And how you use your body affects your spirit. Some of the decisions that will determine your eternal destiny include:

\intalkheader{The Human Spirit}

Your spirit is an eternal entity. The Lord said to His prophet Abraham: “Thou wast chosen before thou wast born.” The Lord said something similar about Jeremiah and many others. He even said it about you.

Your Heavenly Father has known you for a very long time. You, as His son or daughter, were chosen by Him to come to earth at this precise time, to be a leader in His great work on earth. You were chosen not for your bodily characteristics but for your spiritual attributes, such as bravery, courage, integrity of heart, a thirst for truth, a hunger for wisdom, and a desire to serve others.

You developed some of these attributes premortally. Others you can develop here on earth as you persistently seek them.

A pivotal spiritual attribute is that of self-mastery—the strength to place reason over appetite. Self-mastery builds a strong conscience. And your conscience determines your moral responses in difficult, tempting, and trying situations. Fasting helps your spirit to develop dominance over your physical appetites. Fasting also increases your access to heaven’s help, as it intensifies your prayers. Why the need for self-mastery? God implanted strong appetites within us for nourishment and love, vital for the human family to be perpetuated. When we master our appetites within the bounds of God’s laws, we can enjoy longer life, greater love, and consummate joy.

It is not surprising, then, that most temptations to stray from God’s plan of happiness come through the misuse of those essential, God-given appetites. Controlling our appetites is not always easy. Not one of us manages them perfectly. Mistakes happen. Errors are made. Sins are committed. What can we do then? We can learn from them. And we can truly repent.

We can change our behavior. Our very desires can change. How? There is only one way. True change—permanent change—can come only through the healing, cleansing, and enabling power of the Atonement of Jesus Christ. He loves you—each of you! He allows you to access His power as you keep His commandments, eagerly, earnestly, and exactly. It is that simple and certain. The gospel of Jesus Christ is a gospel of change!

A strong human spirit with control over appetites of the flesh is master over emotions and passions and not a slave to them. That kind of freedom is as vital to the spirit as oxygen is to the body! Freedom from self-slavery is true liberation!

We are “free to choose liberty and eternal life … or to choose captivity and death.” When we choose the loftier path toward liberty and eternal life, that path includes marriage. Latter-day Saints proclaim that “marriage between a man and a woman is ordained of God and that the family is central to the Creator’s plan for the eternal destiny of His children.” We also know that “gender is an essential characteristic of individual premortal, mortal, and eternal identity and purpose.”

Marriage between a man and a woman is fundamental to the Lord’s doctrine and crucial to God’s eternal plan. Marriage between a man and a woman is God’s pattern for a fulness of life on earth and in heaven. God’s marriage pattern cannot be abused, misunderstood, or misconstrued. Not if you want true joy. God’s marriage pattern protects the sacred power of procreation and the joy of true marital intimacy. We know that Adam and Eve were married by God before they ever experienced the joy of uniting as husband and wife.

In our day civil governments have a vested interest in protecting marriage because strong families constitute the best way of providing for the health, education, welfare, and prosperity of rising generations. But civil governments are heavily influenced by social trends and secular philosophies as they write, rewrite, and enforce laws. Regardless of what civil legislation may be enacted, the doctrine of the Lord regarding marriage and morality cannot be changed. Remember: sin, even if legalized by man, is still sin in the eyes of God!

While we are to emulate our Savior’s kindness and compassion, while we are to value the rights and feelings of all of God’s children, we cannot change His doctrine. It is not ours to change. His doctrine is ours to study, understand, and uphold.

The Savior’s way of life is good. His way includes chastity before marriage and total fidelity within marriage. The Lord’s way is the only way for us to experience enduring happiness. His way brings sustained comfort to our souls and perennial peace to our homes. And best of all, His way leads us home to Him and our Heavenly Father, to eternal life and exaltation. This is the very essence of God’s work and glory.

My dear brothers and sisters, each day is a day of decision, and our decisions determine our destiny. One day each of us will stand before the Lord in judgment. We will each have a personal interview with Jesus Christ. We will account for decisions that we made about our bodies, our spiritual attributes, and how we honored God’s pattern for marriage and family. That we may choose wisely each day’s decisions for eternity is my earnest prayer in the sacred name of Jesus Christ, amen.

\newpage
\settalk{Let Your Faith Show}
\setspeaker{Russell M. Nelson}
\settalkdate{April 2014}
\talkheading{Let Your Faith Show}{Russell M. Nelson}

Dear brothers and sisters, we express our deepest feelings of love and gratitude for you. We are grateful for our assignments among you.

On a recent flight, our pilot announced that we would encounter turbulence during our descent and that all passengers must fasten their seat belts securely. Sure enough, turbulence came. It was really rough. Across the aisle and a couple of rows behind me, a terrified woman panicked. With each frightening drop and jarring bump, she screamed loudly. Her husband tried to comfort her but to no avail. Her hysterical shouts persisted until we passed through that zone of turbulence to a safe landing. During her period of anxiety, I felt sorry for her. Because faith is the antidote for fear, I silently wished that I could have strengthened her faith.

Later, as passengers were leaving the aircraft, this woman’s husband spoke to me. He said, “I’m sorry my wife was so terrified. The only way I could comfort her was to tell her that ‘Elder Nelson is on this flight, so you don’t need to worry.’”

I’m not sure that my presence on that flight should have given her any comfort, but I will say that one of the realities of mortal life is that our faith will be tested and challenged. Sometimes those tests come as we face what appear to be life-and-death encounters. For this frightened woman, a violently rocking plane presented one of those moments when we come face-to-face with the strength of our faith.

When we speak of faith—the faith that can move mountains—we are not speaking of faith in general but of faith in the Lord Jesus Christ. Faith in the Lord Jesus Christ can be bolstered as we learn about Him and live our religion. The doctrine of Jesus Christ was designed by the Lord to help us increase our faith. In today’s vernacular, however, the word religion can mean different things to different people.

The word religion literally means “to ligate again” or “to tie back” to God. The question we might ask ourselves is, are we securely tied to God so that our faith shows, or are we actually tied to something else? For example, I have overheard conversations on Monday mornings about professional athletic games that took place on the preceding Sunday. For some of these avid fans, I have wondered if their “religion” would “tie them back” only to some kind of a bouncing ball.

We might each ask ourselves, where is our faith? Is it in a team? Is it in a brand? Is it in a celebrity? Even the best teams can fail. Celebrities can fade. There is only One in whom your faith is always safe, and that is in the Lord Jesus Christ. And you need to let your faith show!

God declared in the first of His Ten Commandments, “Thou shalt have no other gods before me.” He also said, “Look unto me in every thought; doubt not, fear not.” Yet so many people look only to their bank balance for peace or to fellow human beings for models to follow.

Clinicians, academicians, and politicians are often put to a test of faith. In pursuit of their goals, will their religion show or will it be hidden? Are they tied back to God or to man?

I had such a test decades ago when one of my medical faculty colleagues chastised me for failing to separate my professional knowledge from my religious convictions. He demanded that I not combine the two. How could I do that? Truth is truth! It is not divisible, and any part of it cannot be set aside.

Whether truth emerges from a scientific laboratory or through revelation, all truth emanates from God. All truth is part of the gospel of Jesus Christ. Yet I was being asked to hide my faith. I did not comply with my colleague’s request. I let my faith show!

In all professional endeavors, rigorous standards of accuracy are required. Scholars cherish their freedom of expression. But full freedom cannot be experienced if part of one’s knowledge is ruled “out-of-bounds” by edicts of men.

Spiritual truth cannot be ignored—especially divine commandments. Keeping divine commandments brings blessings, every time! Breaking divine commandments brings a loss of blessings, every time!

Problems abound in this world because it is populated by imperfect people. Their objectives and desires are heavily influenced by their faith or lack of it. Many put other priorities ahead of God. Some challenge the relevance of religion in modern life. As in every age, so today there are those who mock or decry the free exercise of religion. Some even blame religion for any number of the world’s ills. Admittedly, there have been times when atrocities have been committed in the name of religion. But living the Lord’s pure religion, which means striving to become a true disciple of Jesus Christ, is a way of life and a daily commitment that will provide divine guidance. As you practice your religion, you are exercising your faith. You are letting your faith show.

The Lord knew that His children would need to learn how to find Him. “For strait is the gate,” He said, “and narrow the way that leadeth unto … exaltation … , and few there be that find it.”

The scriptures provide one of the best ways to find our course and stay on it. Scriptural knowledge also provides precious protection. For example, throughout history, infections like “childbirth fever” claimed the lives of many innocent mothers and babies. Yet the Old Testament had the correct principles for the handling of infected patients, written more than 3,000 years ago! Many people perished because man’s quest for knowledge had failed to heed the word of the Lord!

My dear brothers and sisters, what are we missing in our lives if we are “ever learning, [but] never able to come to the knowledge of the truth”? We can gain great knowledge from the scriptures and obtain inspiration through prayers of faith.

Doing so will help us as we make daily decisions. Especially when the laws of man are created and enforced, God’s laws must ever be our standard. In dealing with controversial issues, we should first search for God’s guidance.

We should “liken all scriptures unto us … for our profit and learning.” Danger lurks when we try to divide ourselves with expressions such as “my private life” or even “my best behavior.” If one tries to segment his or her life into such separate compartments, one will never rise to the full stature of one’s personal integrity—never to become all that his or her true self could be.

The temptation to be popular may prioritize public opinion above the word of God. Political campaigns and marketing strategies widely employ public opinion polls to shape their plans. Results of those polls are informative. But they could hardly be used as grounds to justify disobedience to God’s commandments! Even if “everyone is doing it,” wrong is never right. Evil, error, and darkness will never be truth, even if popular. A scriptural warning so declares: “Woe unto them that call evil good, and good evil; that put darkness for light, and light for darkness.”

After World War I, a rather risqué song became popular. In promoting immorality, it vowed that 50 million people cannot be wrong. But in fact, 50 million people can be wrong—totally wrong. Immorality is still immorality in the eyes of God, who one day will judge all of our deeds and desires.

Contrast the fear and faithlessness so prevalent in the world today with the faith and courage of my dearly beloved daughter Emily, who now lives on the other side of the veil. As mortal life was leaving her cancer-ridden body, she could barely speak. But with a smile on her face, she said to me, “Daddy, don’t worry about me. I know I will be all right!” Emily’s faith was showing—showing brightly—in that tender moment, right when we needed it most.

This beautiful young mother of five had full faith in her Heavenly Father, in His plan, and in the eternal welfare of her family. She was securely tied back to God. She was totally faithful to covenants made with the Lord and with her husband. She loved her children but was at peace, despite her impending separation from them. She had faith in her future, and theirs too, because she had faith in our Heavenly Father and His Son.

In 1986, President Thomas S. Monson said: “Of course we will face fear, experience ridicule, and meet opposition. Let us have the courage to defy the consensus, the courage to stand for principle. Courage, not compromise, brings the smile of God’s approval. … Remember that all men have their fears, but those who face their fears with [faith] have courage as well.”

President Monson’s counsel is timeless! So I plead with you, my dear brothers and sisters: Day after day, on your path toward your eternal destiny, increase your faith. Proclaim your faith! Let your faith show!

I pray that you will be securely tied back to God, that His eternal truths will be etched on your heart forever. And I pray that, throughout your life, you will let your faith show! In the name of Jesus Christ, amen.

\newpage
\settalk{Sustaining the Prophets}
\setspeaker{Russell M. Nelson}
\settalkdate{October 2014}
\talkheading{Sustaining the Prophets}{Russell M. Nelson}

President Eyring, we thank you for your instructive and inspiring message. My dear brothers and sisters, we thank you for your faith and devotion. Yesterday, we were each invited to sustain Thomas S. Monson as the prophet of the Lord and President of the Lord’s Church. And often we sing, “We thank thee, O God, for a prophet.” Do you and I really understand what that means? Imagine the privilege the Lord has given us of sustaining His prophet, whose counsel will be untainted, unvarnished, unmotivated by any personal aspiration, and utterly true!

How do we really sustain a prophet? Long before he became President of the Church, President Joseph F. Smith explained, “It is an important duty resting upon the Saints who … sustain the authorities of the Church, to do so not only by the lifting of the hand, the mere form, but in deed and in truth.”

Well do I remember my most unique “deed” to sustain a prophet. As a medical doctor and cardiac surgeon, I had the responsibility of performing open-heart surgery on President Spencer W. Kimball in 1972, when he was Acting President of the Quorum of the Twelve Apostles. He needed a very complex operation. But I had no experience doing such a procedure on a 77-year-old patient in heart failure. I did not recommend the operation and so informed President Kimball and the First Presidency. But, in faith, President Kimball chose to have the operation, only because it was advised by the First Presidency. That shows how he sustained his leaders! And his decision made me tremble!

Thanks to the Lord, the operation was a success. When President Kimball’s heart resumed beating, it did so with great power! At that very moment, I had a clear witness of the Spirit that this man would one day become President of the Church!

You know the outcome. Only 20 months later, President Kimball became President of the Church. And he provided bold and courageous leadership for many years.

Since then we have sustained Presidents Ezra Taft Benson, Howard W. Hunter, Gordon B. Hinckley, and now Thomas S. Monson as Presidents of the Church—prophets in every sense of the word!

My dear brothers and sisters, if the Restoration did anything, it shattered the age-old myth that God had stopped talking to His children. Nothing could be further from the truth. A prophet has stood at the head of God’s Church in all dispensations, from Adam to the present day. Prophets testify of Jesus Christ—of His divinity and of His earthly mission and ministry. We honor the Prophet Joseph Smith as the prophet of this last dispensation. And we honor each man who has succeeded him as President of the Church.

When we sustain prophets and other leaders, we invoke the law of common consent, for the Lord said, “It shall not be given to any one to go forth to preach my gospel, or to build up my church, except he be ordained by some one who has authority, and it is known to the church that he has authority and has been regularly ordained by the heads of the church.”

This gives us, as members of the Lord’s Church, confidence and faith as we strive to keep the scriptural injunction to heed the Lord’s voice as it comes through the voice of His servants the prophets. All leaders in the Lord’s Church are called by proper authority. No prophet or any other leader in this Church, for that matter, has ever called himself or herself. No prophet has ever been elected. The Lord made that clear when He said, “Ye have not chosen me, but I have chosen you, and ordained you.” You and I do not “vote” on Church leaders at any level. We do, though, have the privilege of sustaining them.

The ways of the Lord are different from the ways of man. Man’s ways remove people from office or business when they grow old or become disabled. But man’s ways are not and never will be the Lord’s ways. Our sustaining of prophets is a personal commitment that we will do our utmost to uphold their prophetic priorities. Our sustaining is an oath-like indication that we recognize their calling as a prophet to be legitimate and binding upon us.

Twenty-six years before he became President of the Church, then-Elder George Albert Smith said: “The obligation that we make when we raise our hands … is a most sacred one. It does not mean that we will go quietly on our way and be willing that the prophet of the Lord shall direct this work, but it means … that we will stand behind him; we will pray for him; we will defend his good name, and we will strive to carry out his instructions as the Lord shall direct.”

The living Lord leads His living Church! The Lord reveals His will for the Church to His prophet. Yesterday, after we were invited to sustain Thomas S. Monson as President of the Church, we also had the privilege to sustain him, the counselors in the First Presidency, and members of the Quorum of the Twelve Apostles as prophets, seers, and revelators. Think of that! We sustain 15 men as prophets of God! They hold all the priesthood keys that have ever been conferred upon man in this dispensation.

The calling of 15 men to the holy apostleship provides great protection for us as members of the Church. Why? Because decisions of these leaders must be unanimous. Can you imagine how the Spirit needs to move upon 15 men to bring about unanimity? These 15 men have varied educational and professional backgrounds, with differing opinions about many things. Trust me! These 15 men—prophets, seers, and revelators—know what the will of the Lord is when unanimity is reached! They are committed to see that the Lord’s will truly will be done. The Lord’s Prayer provides the pattern for each of these 15 men when they pray: “Thy will be done on earth as it is in heaven.”

The Apostle with the longest seniority in the office of Apostle presides. That system of seniority will usually bring older men to the office of President of the Church. It provides continuity, seasoned maturity, experience, and extensive preparation, as guided by the Lord.

The Church today has been organized by the Lord Himself. He has put in place a remarkable system of governance that provides redundancy and backup. That system provides for prophetic leadership even when the inevitable illnesses and incapacities may come with advancing age. Counterbalances and safeguards abound so that no one can ever lead the Church astray. Senior leaders are constantly being tutored such that one day they are ready to sit in the highest councils. They learn how to hear the voice of the Lord through the whisperings of the Spirit.

While serving as First Counselor to President Ezra Taft Benson, who was then nearing the end of his mortal life, President Gordon B. Hinckley explained:

“The principles and procedures which the Lord has put in place for the governance of His church make provision for any … circumstance. It is important … that there be no doubts or concerns about the governance of the Church and the exercise of the prophetic gifts, including the right to inspiration and revelation in administering the affairs and programs of the Church, when the President may be ill or is not able to function fully.

“The First Presidency and the Council of the Twelve Apostles, called and ordained to hold the keys of the priesthood, have the authority and responsibility to govern the Church, to administer its ordinances, to expound its doctrine, and to establish and maintain its practices.”

President Hinckley continued:

“When the President is ill or not able to function fully in all of the duties of his office, his two Counselors together comprise a Quorum of the First Presidency. They carry on with the day-to-day work of the Presidency. …

“… But any major questions of policy, procedures, programs, or doctrine are considered deliberately and prayerfully by the First Presidency and the Twelve together.”

Last year, when President Monson reached the milestone of 5 years of service as President of the Church, he reflected on his 50 years of apostolic service and made this statement: “Age eventually takes its toll on all of us. However, we join our voices with King Benjamin, who said, … ‘I am like as yourselves, subject to all manner of infirmities in body and mind; yet I have been chosen … and consecrated by my father, … and have been kept and preserved by his matchless power, to serve you with all the might, mind and strength which the Lord hath granted unto me’ (Mosiah 2:11).”

President Monson continued: “Despite any health challenges that may come to us, despite any weakness in body or mind, we serve to the best of our ability. I assure you that the Church is in good hands. The system set up for the Council of the First Presidency and Quorum of the Twelve [Apostles] assures [us] that it will always be in good hands and that, come what may, there is no need to worry or to fear. Our Savior, Jesus Christ, whom we follow, whom we worship, and whom we serve, is ever at the helm.”

President Monson, we thank you for those truths! And we thank you for your lifetime of exemplary and dedicated service. May I presume to speak for the members of the Church throughout the world in our united and sincere expression of gratitude for you. We honor you! We love you! We sustain you, not only with uplifted hands but with all our hearts and consecrated efforts. Humbly and fervently, “we ever pray for thee, our prophet dear”! In the name of Jesus Christ, amen.

\newpage
\settalk{The Sabbath Is a Delight}
\setspeaker{Russell M. Nelson}
\settalkdate{April 2015}
\talkheading{The Sabbath Is a Delight}{Russell M. Nelson}

Dear brothers and sisters, these two days of conference have been glorious. We have been uplifted by inspiring music and eloquent prayers. Our spirits have been edified by messages of light and truth. On this Easter Sunday, we again unitedly and sincerely thank God for a prophet!

The question for each of us is: because of what I have heard and felt during this conference, how will I change? Whatever your answer might be, may I invite you also to examine your feelings about, and your behavior on, the Sabbath day.

I am intrigued by the words of Isaiah, who called the Sabbath “a delight.” Yet I wonder, is the Sabbath really a delight for you and for me?

I first found delight in the Sabbath many years ago when, as a busy surgeon, I knew that the Sabbath became a day for personal healing. By the end of each week, my hands were sore from repeatedly scrubbing them with soap, water, and a bristle brush. I also needed a breather from the burden of a demanding profession. Sunday provided much-needed relief.

What did the Savior mean when He said that “the sabbath was made for man, and not man for the sabbath”? I believe He wanted us to understand that the Sabbath was His gift to us, granting real respite from the rigors of daily life and an opportunity for spiritual and physical renewal. God gave us this special day, not for amusement or daily labor but for a rest from duty, with physical and spiritual relief.

In Hebrew, the word Sabbath means “rest.” The purpose of the Sabbath dates back to the Creation of the world, when after six days of labor the Lord rested from the work of creation. When He later revealed the Ten Commandments to Moses, God commanded that we “remember the sabbath day, to keep it holy.” Later, the Sabbath was observed as a reminder of the deliverance of Israel from their bondage in Egypt. Perhaps most important, the Sabbath was given as a perpetual covenant, a constant reminder that the Lord may sanctify His people.

In addition, we now partake of the sacrament on the Sabbath day in remembrance of the Atonement of Jesus Christ. Again, we covenant that we are willing to take upon us His holy name.

The Savior identified Himself as Lord of the Sabbath. It is His day! Repeatedly, He has asked us to keep the Sabbath or to hallow the Sabbath day. We are under covenant to do so.

How do we hallow the Sabbath day? In my much younger years, I studied the work of others who had compiled lists of things to do and things not to do on the Sabbath. It wasn’t until later that I learned from the scriptures that my conduct and my attitude on the Sabbath constituted a sign between me and my Heavenly Father. With that understanding, I no longer needed lists of dos and don’ts. When I had to make a decision whether or not an activity was appropriate for the Sabbath, I simply asked myself, “What sign do I want to give to God?” That question made my choices about the Sabbath day crystal clear.

Though the doctrine pertaining to the Sabbath day is of ancient origin, it has been renewed in these latter days as part of a new covenant with a promise. Listen to the power of this divine decree:

“That thou mayest more fully keep thyself unspotted from the world, thou shalt go to the house of prayer and offer up thy sacraments upon my holy day;

“For verily this is a day appointed unto you to rest from your labors, and to pay thy devotions unto the Most High. …

“And on this day … let thy food be prepared with singleness of heart that thy fasting may be perfect, … that thy joy may be full. …

“And inasmuch as ye do these things with thanksgiving, with cheerful hearts and countenances, … the fulness of the earth is yours.”

Imagine the scope of that statement! The fulness of the earth is promised to those who keep the Sabbath day holy. No wonder Isaiah called the Sabbath “a delight.”

How can you ensure that your behavior on the Sabbath will lead to joy and rejoicing? In addition to your going to church, partaking of the sacrament, and being diligent in your specific call to serve, what other activities would help to make the Sabbath a delight for you? What sign will you give to the Lord to show your love for Him?

The Sabbath provides a wonderful opportunity to strengthen family ties. After all, God wants each of us, as His children, to return to Him as endowed Saints, sealed in the temple as families, to our ancestors, and to our posterity.

We make the Sabbath a delight when we teach the gospel to our children. Our responsibility as parents is abundantly clear. The Lord said, “Inasmuch as parents have children in Zion … that teach them not to understand the doctrine of repentance, faith in Christ the Son of the living God, and of baptism and the gift of the Holy Ghost by the laying on of the hands, when eight years old, the sin be upon the heads of the parents.”

Years ago the First Presidency stressed the importance of quality family time. They wrote:

“We call upon parents to devote their best efforts to the teaching and rearing of their children in gospel principles which will keep them close to the Church. The home is the basis of a righteous life, and no other instrumentality can take its place or fulfill its essential functions in carrying forward this God-given responsibility.

“We counsel parents and children to give highest priority to family prayer, family home evening, gospel study and instruction, and wholesome family activities. However worthy and appropriate other demands or activities may be, they must not be permitted to displace the divinely-appointed duties that only parents and families can adequately perform.”

When I ponder this counsel, I almost wish I were a young father once again. Now parents have such wonderful resources available to help them make family time more meaningful, on the Sabbath and other days as well. They have LDS.org, Mormon.org, the Bible videos, the Mormon Channel, the Media Library, the Friend, the New Era, the Ensign, the Liahona, and more—much more. These resources are so very helpful to parents in discharging their sacred duty to teach their children. No other work transcends that of righteous, intentional parenting!

As you teach the gospel, you will learn more. This is the Lord’s way of helping you to comprehend His gospel. He said:

“I give unto you a commandment that you shall teach one another the doctrine of the kingdom.

“Teach ye diligently … , that you may be instructed more perfectly … in doctrine, in the law of the gospel, in all things that pertain unto the kingdom of God.”

Such study of the gospel makes the Sabbath a delight. This promise pertains regardless of family size, composition, or location.

In addition to time with family, you can experience true delight on the Sabbath from family history work. Searching for and finding family members who have preceded you on earth—those who did not have an opportunity to accept the gospel while here—can bring immense joy.

I have seen this firsthand. Several years ago, my dear wife Wendy determined to learn how to do family history research. Her progress at first was slow, but little by little she learned how easy it is to do this sacred work. And I have never seen her happier. You too need not travel to other countries or even to a family history center. At home, with the aid of a computer or mobile device, you can identify souls who are yearning for their ordinances. Make the Sabbath a delight by finding your ancestors and liberating them from spirit prison!

Make the Sabbath a delight by rendering service to others, especially those who are not feeling well or those who are lonely or in need. Lifting their spirits will lift yours as well.

When Isaiah described the Sabbath as “a delight,” he also taught us how to make it delightful. He said:

“If thou turn away … from doing thy pleasure on my holy day; and call the sabbath a delight, … and shalt honour [the Lord], not doing thine own ways, nor finding thine own pleasure, nor speaking thine own words:

“Then shalt thou delight thyself in the Lord.”

Not pursuing your “own pleasure” on the Sabbath requires self-discipline. You may have to deny yourself of something you might like. If you choose to delight yourself in the Lord, you will not permit yourself to treat it as any other day. Routine and recreational activities can be done some other time.

Think of this: In paying tithing, we return one-tenth of our increase to the Lord. In keeping the Sabbath holy, we reserve one day in seven as His. So it is our privilege to consecrate both money and time to Him who lends us life each day.

Faith in God engenders a love for the Sabbath; faith in the Sabbath engenders a love for God. A sacred Sabbath truly is a delight.

Now, as this conference comes to a close, we know that wherever we live we are to be examples of the believers among our families, neighbors, and friends. True believers keep the Sabbath day holy.

I conclude with the farewell plea of Moroni, as he closed the Book of Mormon. He wrote, “Come unto Christ, and be perfected in him, and deny yourselves of all ungodliness; and if ye shall deny yourselves of all ungodliness, and love God with all your might, mind and strength, then … are ye sanctified in Christ.”

With love in my heart, I leave this with you as my prayer, testimony, and blessing in the sacred name of Jesus Christ, amen.

\newpage
\settalk{A Plea to My Sisters}
\setspeaker{Russell M. Nelson}
\settalkdate{October 2015}
\talkheading{A Plea to My Sisters}{Russell M. Nelson}

Dear Elders Rasband, Stevenson, and Renlund, we, your Brethren, welcome you to the Quorum of the Twelve Apostles. We thank God for the revelations that He gives to His prophet, President Thomas S. Monson.

Brothers and sisters, when we met in general conference six months ago, none of us anticipated the coming changes that would tug at the heartstrings of the entire Church. Elder L. Tom Perry delivered a powerful message about the irreplaceable role that marriage and family occupy in the Lord’s plan. We were stunned when just a few days later, we learned of the cancer that would soon take him from us.

Though President Boyd K. Packer’s health had been declining, he continued to “soldier on” in the work of the Lord. He was frail last April, yet he was determined to declare his witness as long as he had breath. Then, just 34 days after Elder Perry’s passing, President Packer also stepped across the veil.

We missed Elder Richard G. Scott at our last general conference, but we’ve reflected upon the powerful witness of the Savior he had borne in many previous conferences. And just 12 days ago, Elder Scott was called home and reunited with his beloved Jeanene.

I had the privilege of being with all of these Brethren during their final days, including joining members of President Packer’s and Elder Scott’s immediate families just before their passing. It has been difficult for me to believe that these three treasured friends, these magnificent servants of the Lord, are gone. I miss them more than I can say.

As I’ve reflected on this unexpected turn of events, one of the impressions that has lingered with me is that which I observed in these surviving wives. Etched in my mind are the serene images of Sister Donna Smith Packer and Sister Barbara Dayton Perry at their husbands’ bedsides, both women filled with love, truth, and pure faith.

As Sister Packer sat next to her husband in his final hours, she radiated that peace that passes all understanding. Though she realized that her beloved companion of almost 70 years would soon depart, she showed the tranquility of a faith-filled woman. She seemed angelic, just as she was in this photo of them at the dedication of the Brigham City Utah Temple.

I saw that same kind of love and faith emanating from Sister Perry. Her devotion to both her husband and the Lord was obvious, and it moved me deeply.

Through their husbands’ final hours and continuing to the present day, these stalwart women have shown the strength and courage that covenant-keeping women always demonstrate. It would be impossible to measure the influence that such women have, not only on families but also on the Lord’s Church, as wives, mothers, and grandmothers; as sisters and aunts; as teachers and leaders; and especially as exemplars and devout defenders of the faith.

This has been true in every gospel dispensation since the days of Adam and Eve. Yet the women of this dispensation are distinct from the women of any other because this dispensation is distinct from any other. This distinction brings both privileges and responsibilities.

Thirty-six years ago, in 1979, President Spencer W. Kimball made a profound prophecy about the impact that covenant-keeping women would have on the future of the Lord’s Church. He prophesied: “Much of the major growth that is coming to the Church in the last days will come because many of the good women of the world … will be drawn to the Church in large numbers. This will happen to the degree that the women of the Church reflect righteousness and articulateness in their lives and to the degree that the women of the Church are seen as distinct and different—in happy ways—from the women of the world.”

My dear sisters, you who are our vital associates during this winding-up scene, the day that President Kimball foresaw is today. You are the women he foresaw! Your virtue, light, love, knowledge, courage, character, faith, and righteous lives will draw good women of the world, along with their families, to the Church in unprecedented numbers!

We, your brethren, need your strength, your conversion, your conviction, your ability to lead, your wisdom, and your voices. The kingdom of God is not and cannot be complete without women who make sacred covenants and then keep them, women who can speak with the power and authority of God!

President Packer declared:

“We need women who are organized and women who can organize. We need women with executive ability who can plan and direct and administer; women who can teach, women who can speak out. …

“We need women with the gift of discernment who can view the trends in the world and detect those that, however popular, are shallow or dangerous.”

Today, let me add that we need women who know how to make important things happen by their faith and who are courageous defenders of morality and families in a sin-sick world. We need women who are devoted to shepherding God’s children along the covenant path toward exaltation; women who know how to receive personal revelation, who understand the power and peace of the temple endowment; women who know how to call upon the powers of heaven to protect and strengthen children and families; women who teach fearlessly.

Throughout my life, I have been blessed by such women. My departed wife, Dantzel, was such a woman. I will always be grateful for the life-changing influence she had on me in all aspects of my life, including my pioneering efforts in open-heart surgery.

Fifty-eight years ago I was asked to operate upon a little girl, gravely ill from congenital heart disease. Her older brother had previously died of a similar condition. Her parents pleaded for help. I was not optimistic about the outcome but vowed to do all in my power to save her life. Despite my best efforts, the child died. Later, the same parents brought another daughter to me, then just 16 months old, also born with a malformed heart. Again, at their request, I performed an operation. This child also died. This third heartbreaking loss in one family literally undid me.

I went home grief stricken. I threw myself upon our living room floor and cried all night long. Dantzel stayed by my side, listening as I repeatedly declared that I would never perform another heart operation. Then, around 5:00 in the morning, Dantzel looked at me and lovingly asked, “Are you finished crying? Then get dressed. Go back to the lab. Go to work! You need to learn more. If you quit now, others will have to painfully learn what you already know.”

Oh, how I needed my wife’s vision, grit, and love! I went back to work and learned more. If it weren’t for Dantzel’s inspired prodding, I would not have pursued open-heart surgery and would not have been prepared to do the operation in 1972 that saved the life of President Spencer W. Kimball.

Sisters, do you realize the breadth and scope of your influence when you speak those things that come to your heart and mind as directed by the Spirit? A superb stake president told me of a stake council meeting in which they were wrestling with a difficult challenge. At one point, he realized that the stake Primary president had not spoken, so he asked if she had any impressions. “Well, actually I have,” she said and then proceeded to share a thought that changed the entire direction of the meeting. The stake president continued, “As she spoke, the Spirit testified to me that she had given voice to the revelation we had been seeking as a council.”

My dear sisters, whatever your calling, whatever your circumstances, we need your impressions, your insights, and your inspiration. We need you to speak up and speak out in ward and stake councils. We need each married sister to speak as “a contributing and full partner” as you unite with your husband in governing your family. Married or single, you sisters possess distinctive capabilities and special intuition you have received as gifts from God. We brethren cannot duplicate your unique influence.

We know that the culminating act of all creation was the creation of woman! We need your strength!

Attacks against the Church, its doctrine, and our way of life are going to increase. Because of this, we need women who have a bedrock understanding of the doctrine of Christ and who will use that understanding to teach and help raise a sin-resistant generation. We need women who can detect deception in all of its forms. We need women who know how to access the power that God makes available to covenant keepers and who express their beliefs with confidence and charity. We need women who have the courage and vision of our Mother Eve.

My dear sisters, nothing is more crucial to your eternal life than your own conversion. It is converted, covenant-keeping women—women like my dear wife Wendy—whose righteous lives will increasingly stand out in a deteriorating world and who will thus be seen as different and distinct in the happiest of ways.

So today I plead with my sisters of The Church of Jesus Christ of Latter-day Saints to step forward! Take your rightful and needful place in your home, in your community, and in the kingdom of God—more than you ever have before. I plead with you to fulfill President Kimball’s prophecy. And I promise you in the name of Jesus Christ that as you do so, the Holy Ghost will magnify your influence in an unprecedented way!

I bear witness of the reality of the Lord Jesus Christ and of His redeeming, atoning, and sanctifying power. And as one of His Apostles, I thank you, my dear sisters, and bless you to rise to your full stature, to fulfill the measure of your creation, as we walk arm in arm in this sacred work. Together we will help prepare the world for the Second Coming of the Lord. Of this I testify, as your brother, in the name of Jesus Christ, amen.

\newpage
\settalk{The Price of Priesthood Power}
\setspeaker{Russell M. Nelson}
\settalkdate{April 2016}
\talkheading{The Price of Priesthood Power}{Russell M. Nelson}

Six months ago in the October 2015 general conference, I spoke to the sisters of the Church about their divine role as women of God. Now I wish to speak to you brethren about your divine role as men of God. As I travel the world, I marvel at the strength and sheer goodness of the men and boys of this Church. There is simply no way to number the hearts you’ve healed and the lives you’ve lifted. Thank you!

In my last conference message, I related my devastating experience many years ago when, as a heart surgeon, I was not able to save the lives of two little sisters. With permission of their father, I would like to say more about that family.

Congenital heart disease afflicted three children born to Ruth and Jimmy Hatfield. Their first son, Jimmy Jr., died without a definitive diagnosis. I entered the picture when the parents sought help for their two daughters, Laural Ann and her younger sister, Gay Lynn. I was heartbroken when both girls died following their operations. Understandably, Ruth and Jimmy were spiritually shattered.

Over time, I learned that they harbored lingering resentment toward me and the Church. For almost six decades, I have been haunted by this situation and have grieved for the Hatfields. I tried several times to establish contact with them, without success.

Then one night last May, I was awakened by those two little girls from the other side of the veil. Though I did not see or hear them with my physical senses, I felt their presence. Spiritually, I heard their pleadings. Their message was brief and clear: “Brother Nelson, we are not sealed to anyone! Can you help us?” Soon thereafter, I learned that their mother had passed away, but their father and younger brother were still alive.

Emboldened by the pleadings of Laural Ann and Gay Lynn, I tried again to contact their father, who I learned was living with his son Shawn. This time they were willing to meet with me.

In June, I literally knelt in front of Jimmy, now 88 years old, and had a heart-to-heart talk with him. I spoke of his daughters’ pleadings and told him I would be honored to perform sealing ordinances for his family. I also explained that it would take time and much effort on his and Shawn’s part to be ready and worthy to enter the temple, as neither of them had ever been endowed.

The Spirit of the Lord was palpable throughout that meeting. And when Jimmy and Shawn each accepted my offer, I was overjoyed! They worked diligently with their stake president, bishop, home teachers, and ward mission leader, as well as with young missionaries and a senior missionary couple. And then, not long ago, in the Payson Utah Temple, I had the profound privilege of sealing Ruth to Jimmy and their four children to them. Wendy and I wept as we participated in that sublime experience. Many hearts were healed that day!

On reflection, I have marveled at Jimmy and Shawn and what they were willing to do. They have become heroes to me. If I could have the wish of my heart, it would be that each man and young man in this Church would demonstrate the courage, strength, and humility of this father and son. They were willing to forgive and let go of old hurts and habits. They were willing to submit to guidance from their priesthood leaders so that the Atonement of Jesus Christ could purify and magnify them. Each was willing to become a man who worthily bears the priesthood “after the holiest order of God.”

To bear means to support the weight of that which is held. It is a sacred trust to bear the priesthood, which is the mighty power and authority of God. Think of this: the priesthood conferred upon us is the very same power and authority through which God created this and numberless worlds, governs the heavens and the earth, and exalts His obedient children.

Recently, Wendy and I were in a meeting where the organist was poised and ready to play the opening hymn. His eyes were on the music, and his fingers were on the keys. He began pressing the keys, but there was no sound. I whispered to Wendy, “He has no power.” I reasoned that something had stopped the flow of electrical power to that organ.

Well, brethren, in like manner, I fear that there are too many men who have been given the authority of the priesthood but who lack priesthood power because the flow of power has been blocked by sins such as laziness, dishonesty, pride, immorality, or preoccupation with things of the world.

I fear that there are too many priesthood bearers who have done little or nothing to develop their ability to access the powers of heaven. I worry about all who are impure in their thoughts, feelings, or actions or who demean their wives or children, thereby cutting off priesthood power.

I fear that too many have sadly surrendered their agency to the adversary and are saying by their conduct, “I care more about satisfying my own desires than I do about bearing the Savior’s power to bless others.”

I fear, brethren, that some among us may one day wake up and realize what power in the priesthood really is and face the deep regret that they spent far more time seeking power over others or power at work than learning to exercise fully the power of God. President George Albert Smith taught that “we are not here to while away the hours of this life and then pass to a sphere of exaltation; but we are here to qualify ourselves day by day for the positions that our Father expects us to fill hereafter.”

Why would any man waste his days and settle for Esau’s mess of pottage when he has been entrusted with the possibility of receiving all of the blessings of Abraham?

I urgently plead with each one of us to live up to our privileges as bearers of the priesthood. In a coming day, only those men who have taken their priesthood seriously, by diligently seeking to be taught by the Lord Himself, will be able to bless, guide, protect, strengthen, and heal others. Only a man who has paid the price for priesthood power will be able to bring miracles to those he loves and keep his marriage and family safe, now and throughout eternity.

What is the price to develop such priesthood power? The Savior’s senior Apostle, Peter—that same Peter who with James and John conferred the Melchizedek Priesthood upon Joseph Smith and Oliver Cowdery—declared qualities we should seek to “be partakers of the divine nature.”

He named faith, virtue, knowledge, temperance, patience, godliness, brotherly kindness, charity, and diligence. And don’t forget humility! So I ask, how would our family members, friends, and coworkers say you and I are doing in developing these and other spiritual gifts? The more those attributes are developed, the greater will be our priesthood power.

How else can we increase our power in the priesthood? We need to pray from our hearts. Polite recitations of past and upcoming activities, punctuated with some requests for blessings, cannot constitute the kind of communing with God that brings enduring power. Are you willing to pray to know how to pray for more power? The Lord will teach you.

Are you willing to search the scriptures and feast on the words of Christ—to study earnestly in order to have more power? If you want to see your wife’s heart melt, let her find you on the Internet studying the doctrine of Christ or reading your scriptures!

Are you willing to worship in the temple regularly? The Lord loves to do His own teaching in His holy house. Imagine how pleased He would be if you asked Him to teach you about priesthood keys, authority, and power as you experience the ordinances of the Melchizedek Priesthood in the holy temple. Imagine the increase in priesthood power that could be yours.

Are you willing to follow President Thomas S. Monson’s example of serving others? For decades he has taken the long way home, following promptings of the Spirit to arrive on someone’s doorstep and then hear words such as, “How did you know it was the anniversary of our daughter’s death?” or “How did you know it was my birthday?” And if you truly want more priesthood power, you will cherish and care for your wife, embracing both her and her counsel.

Now, if all of this sounds excessive, please consider how different our relationships with our wife, children, and associates at work would be if we were as concerned about gaining priesthood power as we are in progressing at work or increasing the balance in our bank account. If we will humbly present ourselves before the Lord and ask Him to teach us, He will show us how to increase our access to His power.

In these latter days, we know there will be earthquakes in diverse places. Perhaps one of those diverse places will be in our own homes, where emotional, financial, or spiritual “earthquakes” may occur. Priesthood power can calm the seas and heal fractures in the earth. Priesthood power can also calm the minds and heal fractures in the hearts of those we love.

Are we willing to pray, fast, study, seek, worship, and serve as men of God so we can have that kind of priesthood power? Because two little girls were so eager to be sealed to their family, their father and brother were willing to pay the price to bear the holy Melchizedek Priesthood.

My dear brethren, we have been given a sacred trust—the authority of God to bless others. May each one of us rise up as the man God foreordained us to be—ready to bear the priesthood of God bravely, eager to pay whatever price is required to increase his power in the priesthood. With that power, we can help prepare the world for the Second Coming of our Savior, Jesus Christ. This is His Church, led today by His prophet, President Thomas S. Monson, whom I dearly love and sustain. I so testify in the name of Jesus Christ, amen.

\newpage
\settalk{Joy and Spiritual Survival}
\setspeaker{Russell M. Nelson}
\settalkdate{October 2016}
\talkheading{Joy and Spiritual Survival}{Russell M. Nelson}

My dear brothers and sisters, today I would like to discuss a principle that is key to our spiritual survival. It is a principle that will only become more important as the tragedies and travesties around us increase.

These are the latter days, so none of us should be surprised when we see prophecy fulfilled. A host of prophets, including Isaiah, Paul, Nephi, and Mormon, foresaw that perilous times would come, that in our day the whole world would be in commotion, that men would “be lovers of their own selves, … without natural affection, … lovers of pleasures more than lovers of God,” and that many would become servants of Satan who uphold the adversary’s work. Indeed, you and I “wrestle … against the rulers of the darkness of this world, [and] against spiritual wickedness in high places.”

As conflicts between nations escalate, as cowardly terrorists prey on the innocent, and as corruption in everything from business to government becomes increasingly commonplace, what can help us? What can help each of us with our personal struggles and with the rigorous challenge of living in these latter days?

The prophet Lehi taught a principle for spiritual survival. First, consider his circumstances: He had been persecuted for preaching truth in Jerusalem and had been commanded by the Lord to leave his possessions and flee with his family into the wilderness. He had lived in a tent and survived on what food could be found on the way to an unknown destination, and he had watched two of his sons, Laman and Lemuel, rebel against the teachings of the Lord and attack their brothers Nephi and Sam.

Clearly, Lehi knew opposition, anxiety, heartache, pain, disappointment, and sorrow. Yet he declared boldly and without reservation a principle as revealed by the Lord: “Men are, that they might have joy.” Imagine! Of all the words he could have used to describe the nature and purpose of our lives here in mortality, he chose the word joy!

Life is filled with detours and dead ends, trials and challenges of every kind. Each of us has likely had times when distress, anguish, and despair almost consumed us. Yet we are here to have joy?

Yes! The answer is a resounding yes! But how is that possible? And what must we do to claim the joy that Heavenly Father has in store for us?

Eliza R. Snow, second General President of the Relief Society, offered a riveting answer. Because of Missouri’s infamous extermination order, issued at the onset of the grueling winter of 1838, she and other Saints were forced to flee the state that very winter. One evening, Eliza’s family spent the night in a small log cabin used by refugee Saints. Much of the chinking between the logs had been extracted and burned for firewood by those who preceded them, so there were holes between the logs large enough for a cat to crawl through. It was bitter cold, and their food was frozen solid.

That night some 80 people huddled inside that small cabin, only 20 feet square (6.1 meters square). Most sat or stood all night trying to keep warm. Outside, a group of men spent the night gathered around a roaring fire, with some singing hymns and others roasting frozen potatoes. Eliza recorded: “Not a complaint was heard—all were cheerful, and judging from appearances, strangers would have taken us to be pleasure excursionists rather than a band of gubernatorial exiles.”

Eliza’s report of that exhausting, bone-chilling evening was strikingly optimistic. She declared: “That was a very merry night. None but saints can be happy under every circumstance.”

That’s it! Saints can be happy under every circumstance. We can feel joy even while having a bad day, a bad week, or even a bad year!

My dear brothers and sisters, the joy we feel has little to do with the circumstances of our lives and everything to do with the focus of our lives.

When the focus of our lives is on God’s plan of salvation, which President Thomas S. Monson just taught us, and Jesus Christ and His gospel, we can feel joy regardless of what is happening—or not happening—in our lives. Joy comes from and because of Him. He is the source of all joy. We feel it at Christmastime when we sing, “Joy to the world, the Lord is come.” And we can feel it all year round. For Latter-day Saints, Jesus Christ is joy!

That is why our missionaries leave their homes to preach His gospel. Their goal is not to increase the number of Church members. Rather, our missionaries teach and baptize to bring joy to the people of the world!

Just as the Savior offers peace that “passeth all understanding,” He also offers an intensity, depth, and breadth of joy that defy human logic or mortal comprehension. For example, it doesn’t seem possible to feel joy when your child suffers with an incurable illness or when you lose your job or when your spouse betrays you. Yet that is precisely the joy the Savior offers. His joy is constant, assuring us that our “afflictions shall be but a small moment” and be consecrated to our gain.

How, then, can we claim that joy? We can start by “looking unto Jesus the author and finisher of our faith” “in every thought.” We can give thanks for Him in our prayers and by keeping covenants we’ve made with Him and our Heavenly Father. As our Savior becomes more and more real to us and as we plead for His joy to be given to us, our joy will increase.

Joy is powerful, and focusing on joy brings God’s power into our lives. As in all things, Jesus Christ is our ultimate exemplar, “who for the joy that was set before him endured the cross.” Think of that! In order for Him to endure the most excruciating experience ever endured on earth, our Savior focused on joy!

And what was the joy that was set before Him? Surely it included the joy of cleansing, healing, and strengthening us; the joy of paying for the sins of all who would repent; the joy of making it possible for you and me to return home—clean and worthy—to live with our Heavenly Parents and families.

If we focus on the joy that will come to us, or to those we love, what can we endure that presently seems overwhelming, painful, scary, unfair, or simply impossible?

One father in a spiritually precarious situation focused on the joy of finally being clean and right with the Lord—the joy of being freed from guilt and shame—and the joy of having peace of mind. That focus gave him the courage to confess to his wife and bishop about his problem with pornography and his subsequent infidelity. He is now doing everything his bishop counsels him to do, striving with all his heart to regain the trust of his dear wife.

A young woman focused on the joy of staying sexually pure to help her endure the mocking of friends as she walked away from a popular and provocative, but spiritually dangerous, situation.

A man who frequently demeaned his wife and indulged in angry outbursts at his children focused on the joy of being worthy to have the Holy Ghost as his constant companion. That focus motivated him to put off the natural man, to which he had too often succumbed, and make needed changes.

A dear colleague recently told me of his past two decades of heavy trials. He said, “I have learned to suffer with joy. My suffering was swallowed up in the joy of Christ.”

What will you and I be able to endure as we focus on the joy that is “set before” us? What repenting will then be possible? What weakness will become a strength? What chastening will become a blessing? What disappointments, even tragedies, will turn to our good? And what challenging service to the Lord will we be able to give?

As we diligently focus on the Savior and then follow His pattern of focusing on joy, we need to avoid those things that can interrupt our joy. Remember Korihor, the anti-Christ? Spewing falsehoods about the Savior, Korihor went from place to place until he was brought before a high priest who asked him: “Why do ye go about perverting the ways of the Lord? Why do ye teach this people that there shall be no Christ, to interrupt their rejoicings?”

Anything that opposes Christ or His doctrine will interrupt our joy. That includes the philosophies of men, so abundant online and in the blogosphere, which do exactly what Korihor did.

If we look to the world and follow its formulas for happiness, we will never know joy. The unrighteous may experience any number of emotions and sensations, but they will never experience joy! Joy is a gift for the faithful. It is the gift that comes from intentionally trying to live a righteous life, as taught by Jesus Christ.

He taught us how to have joy. When we choose Heavenly Father to be our God and when we can feel the Savior’s Atonement working in our lives, we will be filled with joy. Every time we nurture our spouse and guide our children, every time we forgive someone or ask for forgiveness, we can feel joy.

Every day that you and I choose to live celestial laws, every day that we keep our covenants and help others to do the same, joy will be ours.

Heed these words of the Psalmist: “I have set the Lord always before me: because he is at my right hand, I shall not be moved. … In [His] presence is fulness of joy.” As this principle is embedded in our hearts, each and every day can be a day of joy and gladness. I so testify in the sacred name of Jesus Christ, amen.

\newpage
\settalk{Drawing the Power of Jesus Christ into Our Lives}
\setspeaker{Russell M. Nelson}
\settalkdate{April 2017}
\talkheading{Drawing the Power of Jesus Christ into Our Lives}{Russell M. Nelson}

My dear brothers and sisters, we live in a most difficult dispensation. Challenges, controversies, and complexities swirl around us. These turbulent times were foreseen by the Savior. He warned us that in our day the adversary would stir up anger in the hearts of men and lead them astray. Yet our Heavenly Father never intended that we would deal with the maze of personal problems and social issues on our own.

God so loved the world that He sent His Only Begotten Son to help us. And His Son, Jesus Christ, gave His life for us. All so that we could have access to godly power—power sufficient to deal with the burdens, obstacles, and temptations of our day. Today I would like to speak about how we can draw into our lives the power of our Lord and Master, Jesus Christ.

We begin by learning about Him. “It is impossible for [us] to be saved in ignorance.” The more we know about the Savior’s ministry and mission—the more we understand His doctrine and what He did for us—the more we know that He can provide the power that we need for our lives.

Earlier this year, I asked the young adults of the Church to consecrate a portion of their time each week to study everything Jesus said and did as recorded in the standard works. I invited them to let the scriptural citations about Jesus Christ in the Topical Guide become their personal core curriculum.

I gave that challenge because I had already accepted it myself. I read and underlined every verse cited about Jesus Christ, as listed under the main heading and the 57 subtitles in the Topical Guide. When I finished that exciting exercise, my wife asked me what impact it had on me. I told her, “I am a different man!”

I felt a renewed devotion to Him as I read again in the Book of Mormon the Savior’s own statement about His mission in mortality. He declared:

“I came into the world to do the will of my Father, because my Father sent me.

“And my Father sent me that I might be lifted up upon the cross.”

As Latter-day Saints, we refer to His mission as the Atonement of Jesus Christ, which made resurrection a reality for all and made eternal life possible for those who repent of their sins and receive and keep essential ordinances and covenants.

It is doctrinally incomplete to speak of the Lord’s atoning sacrifice by shortcut phrases, such as “the Atonement” or “the enabling power of the Atonement” or “applying the Atonement” or “being strengthened by the Atonement.” These expressions present a real risk of misdirecting faith by treating the event as if it had living existence and capabilities independent of our Heavenly Father and His Son, Jesus Christ.

Under the Father’s great eternal plan, it is the Savior who suffered. It is the Savior who broke the bands of death. It is the Savior who paid the price for our sins and transgressions and blots them out on condition of our repentance. It is the Savior who delivers us from physical and spiritual death.

There is no amorphous entity called “the Atonement” upon which we may call for succor, healing, forgiveness, or power. Jesus Christ is the source. Sacred terms such as Atonement and Resurrection describe what the Savior did, according to the Father’s plan, so that we may live with hope in this life and gain eternal life in the world to come. The Savior’s atoning sacrifice—the central act of all human history—is best understood and appreciated when we expressly and clearly connect it to Him.

The importance of the Savior’s mission was emphasized by the Prophet Joseph Smith, who declared emphatically that “the fundamental principles of our religion are the testimony of the Apostles and Prophets, concerning Jesus Christ, that He died, was buried, and rose again the third day, and ascended into heaven; and all other things which pertain to our religion are only appendages to it.”

It was this very statement of the Prophet that provided the incentive for 15 prophets, seers, and revelators to issue and sign their testimony to commemorate the 2,000th anniversary of the Lord’s birth. That historic testimony is titled “The Living Christ.” Many members have memorized its truths. Others barely know of its existence. As you seek to learn more about Jesus Christ, I urge you to study “The Living Christ.”

As we invest time in learning about the Savior and His atoning sacrifice, we are drawn to participate in another key element to accessing His power: we choose to have faith in Him and follow Him.

True disciples of Jesus Christ are willing to stand out, speak up, and be different from the people of the world. They are undaunted, devoted, and courageous. I learned of such disciples during a recent assignment in Mexico, where I met with government officials as well as leaders of other religious denominations. Each thanked me for our members’ heroic and successful efforts to protect and preserve strong marriages and families in their country.

There is nothing easy or automatic about becoming such powerful disciples. Our focus must be riveted on the Savior and His gospel. It is mentally rigorous to strive to look unto Him in every thought. But when we do, our doubts and fears flee.

Recently I learned of a fearless young Laurel. She was invited to participate in a statewide competition for her high school on the same evening she had committed to participate in a stake Relief Society meeting. When she realized the conflict and explained to competition officials that she would need to leave the competition early to attend an important meeting, she was told she would be disqualified if she did so.

What did this latter-day Laurel do? She kept her commitment to participate in the Relief Society meeting. As promised, she was disqualified from the statewide competition. When asked about her decision, she replied simply, “Well, the Church is more important, isn’t it?”

Faith in Jesus Christ propels us to do things we otherwise would not do. Faith that motivates us to action gives us more access to His power.

We also increase the Savior’s power in our lives when we make sacred covenants and keep those covenants with precision. Our covenants bind us to Him and give us godly power. As faithful disciples, we repent and follow Him into the waters of baptism. We walk along the covenant path to receive other essential ordinances. And gratefully, God’s plan provides for those blessings to be extended to ancestors who died without an opportunity to obtain them during their mortal lives.

Covenant-keeping men and women seek for ways to keep themselves unspotted from the world so there will be nothing blocking their access to the Savior’s power. One faithful wife and mother wrote this recently: “These are troubled and perilous times. How blessed we are to have the increased knowledge of the plan of salvation and the inspired guidance from loving prophets, apostles, and leaders to help us sail these stormy seas safely. We stopped our habit of turning on the radio in the morning. Instead, we now listen to a general conference talk on our mobile phones every morning as we prepare ourselves for another day.”

Another element in drawing the Savior’s power into our lives is to reach up to Him in faith. Such reaching requires diligent, focused effort.

Do you remember the biblical story of the woman who suffered for 12 years with a debilitating problem? She exercised great faith in the Savior, exclaiming, “If I may touch but his clothes, I shall be whole.”

This faithful, focused woman needed to stretch as far as she could to access His power. Her physical stretching was symbolic of her spiritual stretching.

Many of us have cried out from the depths of our hearts a variation of this woman’s words: “If I could spiritually stretch enough to draw the Savior’s power into my life, I would know how to handle my heart-wrenching situation. I would know what to do. And I would have the power to do it.”

When you reach up for the Lord’s power in your life with the same intensity that a drowning person has when grasping and gasping for air, power from Jesus Christ will be yours. When the Savior knows you truly want to reach up to Him—when He can feel that the greatest desire of your heart is to draw His power into your life—you will be led by the Holy Ghost to know exactly what you should do.

When you spiritually stretch beyond anything you have ever done before, then His power will flow into you. And then you will understand the deep meaning of words we sing in the hymn “The Spirit of God”:

\begin{verse}
The Lord is extending the Saints’ understanding. …\\
The knowledge and power of God are expanding;\\
The veil o’er the earth is beginning to burst.
\end{verse}

The gospel of Jesus Christ is filled with His power, which is available to every earnestly seeking daughter or son of God. It is my testimony that when we draw His power into our lives, both He and we will rejoice.

As one of His special witnesses, I declare that God lives! Jesus is the Christ! His Church has been restored to the earth! God’s prophet upon the earth today is President Thomas S. Monson, whom I sustain with all my heart. I so testify, with my expression of love and blessing for each of you, in the sacred name of Jesus Christ, amen.

\newpage
\settalk{The Book of Mormon: What Would Your Life Be Like without It?}
\setspeaker{Russell M. Nelson}
\settalkdate{October 2017}
\talkheading{The Book of Mormon: What Would Your Life Be Like without It?}{Russell M. Nelson}

In 1986, I was invited to give a special lecture at a university in Accra, Ghana. There I met a number of dignitaries, including an African tribal king. As we visited prior to the lecture, the king spoke to me only through his linguist, who then translated for me. I responded to the linguist, and the linguist then translated my responses to the king.

After my lecture, the king made his way directly to me, but this time without his linguist. To my surprise, he spoke in perfect English—the Queen’s English, I might add!

The king seemed puzzled. “Just who are you?” he asked.

I replied, “I am an ordained Apostle of Jesus Christ.”

The king asked, “What can you teach me about Jesus Christ?”

I responded with a question: “May I ask what you already know about Him?”

The king’s response revealed he was a serious student of the Bible and one who loved the Lord.

I then asked if he knew about the ministry of Jesus Christ to the people of ancient America.

As I expected, he did not.

I explained that after the Savior’s Crucifixion and Resurrection, He came to the people of ancient America, where He taught His gospel. He organized His Church and asked His disciples to keep a record of His ministry among them.

“That record,” I continued, “is what we know as the Book of Mormon. It is another testament of Jesus Christ. It is a companion scripture to the Holy Bible.”

At this point, the king became very interested. I turned to the mission president accompanying me and asked if he had an extra copy of the Book of Mormon with him. He pulled one from his briefcase.

I opened it to 3 Nephi chapter 11, and together the king and I read the Savior’s sermon to the Nephites. I then presented the copy of the Book of Mormon to him. His response lodged in my mind and heart forever: “You could have given me diamonds or rubies, but nothing is more precious to me than this additional knowledge about the Lord Jesus Christ.”

After experiencing the power of the Savior’s words in 3 Nephi, the king proclaimed, “If I am converted and join the Church, I will bring my whole tribe with me.”

“Oh, King,” I said, “it doesn’t work that way. Conversion is an individual matter. The Savior ministered to the Nephites one by one. Each individual receives a witness and testimony of the gospel of Jesus Christ.”

My brothers and sisters, how precious is the Book of Mormon to you? If you were offered diamonds or rubies or the Book of Mormon, which would you choose? Honestly, which is of greater worth to you?

Remember in the Sunday morning session of the April 2017 general conference, President Thomas S. Monson pleaded with “each of us to prayerfully study and ponder the Book of Mormon each day.” Many have responded to our prophet’s plea.

Let me say that neither I nor eight-year-old Riley knew anyone was taking our pictures. Notice that Riley is reading his Book of Mormon with the help of an “I Am a Child of God” bookmark.

Something powerful happens when a child of God seeks to know more about Him and His Beloved Son. Nowhere are those truths taught more clearly and powerfully than in the Book of Mormon.

Since President Monson’s challenge six months ago, I have tried to follow his counsel. Among other things, I’ve made lists of what the Book of Mormon is, what it affirms, what it refutes, what it fulfills, what it clarifies, and what it reveals. Looking at the Book of Mormon through these lenses has been an insightful and inspiring exercise! I recommend it to each of you.

During these six months, I have invited various groups—including my Brethren in the Quorum of the Twelve, missionaries in Chile, and mission presidents and their wives gathered in Argentina—to consider three related questions that I urge you to think about today:

First, what would your life be like without the Book of Mormon? Second, what would you not know? And third, what would you not have?

Enthusiastic answers from these groups came straight from their hearts. Here are just a few of their comments:

“Without the Book of Mormon, I would be confused about the conflicting teachings and opinions about so many things. I would be just like I was before I found the Church, when I was searching for knowledge, faith, and hope.”

Another said: “I would not know about the role the Holy Ghost can play in my life.”

Another: “I would not clearly understand my purpose here on earth!”

Another respondent said: “I would not know that there is continuing progress after this life. Because of the Book of Mormon, I know that there really is life after death. That is the ultimate goal for which we are working.”

This last comment made me reflect on my life decades ago as a young surgical resident. One of the sobering responsibilities a surgeon bears, on occasion, is to inform the family when a loved one passes away. In one hospital where I worked, a special room was built with padded walls where family members could receive such news. There, some people manifested their grief by banging their heads against those padded walls. How I longed to teach those individuals that death, though difficult for surviving loved ones, is a necessary part of our immortal existence. Death allows us to progress to the next world.

Another respondent to my question said: “I did not have a life until I read the Book of Mormon. Even though I had prayed and gone to my church all my life, the Book of Mormon helped me to really communicate with Heavenly Father for the first time.”

Another said: “Without the Book of Mormon, I wouldn’t understand that the Savior not only suffered for my sins, but He can heal my pains and sorrows.”

And yet another: “I would not know that we have prophets to lead us.”

Immersing ourselves regularly in the truths of the Book of Mormon can be a life-changing experience. One of our missionary granddaughters, Sister Olivia Nelson, promised an investigator that if he would read the Book of Mormon daily, his test scores on his university exams would improve. He did, and they did.

My dear brothers and sisters, I testify that the Book of Mormon is truly the word of God. It contains the answers to life’s most compelling questions. It teaches the doctrine of Christ. It expands and clarifies many of the “plain and precious” truths that were lost through centuries of time and numerous translations of the Bible.

The Book of Mormon provides the fullest and most authoritative understanding of the Atonement of Jesus Christ to be found anywhere. It teaches what it really means to be born again. From the Book of Mormon we learn about the gathering of scattered Israel. We know why we are here on earth. These and other truths are more powerfully and persuasively taught in the Book of Mormon than in any other book. The full power of the gospel of Jesus Christ is contained in the Book of Mormon. Period.

The Book of Mormon both illuminates the teachings of the Master and exposes the tactics of the adversary. The Book of Mormon teaches true doctrine to dispel false religious traditions—such as the erroneous practice of performing infant baptisms. The Book of Mormon gives purpose to life by urging us to ponder the potential of eternal life and “never-ending happiness.” The Book of Mormon shatters the false beliefs that happiness can be found in wickedness and that individual goodness is all that is required to return to the presence of God. It abolishes forever the false concepts that revelation ended with the Bible and that the heavens are sealed today.

When I think of the Book of Mormon, I think of the word power. The truths of the Book of Mormon have the power to heal, comfort, restore, succor, strengthen, console, and cheer our souls.

My dear brothers and sisters, I promise that as you prayerfully study the Book of Mormon every day, you will make better decisions—every day. I promise that as you ponder what you study, the windows of heaven will open, and you will receive answers to your own questions and direction for your own life. I promise that as you daily immerse yourself in the Book of Mormon, you can be immunized against the evils of the day, even the gripping plague of pornography and other mind-numbing addictions.

Whenever I hear anyone, including myself, say, “I know the Book of Mormon is true,” I want to exclaim, “That’s nice, but it is not enough!” We need to feel, deep in “the inmost part” of our hearts, that the Book of Mormon is unequivocally the word of God. We must feel it so deeply that we would never want to live even one day without it. I might paraphrase President Brigham Young in saying, “I wish I had the voice of seven thunders to wake up the people” to the truth and power of the Book of Mormon.

We need to be like this young missionary serving in Europe who felt so deeply about the truth of the Book of Mormon that he literally ran with a copy of this sacred record to the man in a park he and his companion had just found.

I testify that Joseph Smith was and is the prophet of this last dispensation. It was he who, through the gift and power of God, translated this holy book. This is the book that will help to prepare the world for the Second Coming of the Lord.

I testify that Jesus Christ is the literal and living Son of our living God. He is our Savior, our Redeemer, our great Exemplar, and our Advocate with the Father. He was the promised Messiah, the mortal Messiah, and will be the millennial Messiah. I testify with my whole soul that in a most miraculous and singular way, the Book of Mormon teaches us of Jesus Christ and His gospel.

I know that President Thomas S. Monson is the prophet of God on the earth today. I love him and sustain him with all of my heart. I so testify in the sacred name of Jesus Christ, amen.

\intalkheader{President Nelson’s Book of Mormon Lists}

The Book of Mormon is:

The Book of Mormon affirms:

The Book of Mormon refutes notions that:

The Book of Mormon fulfills biblical prophecies that:

The Book of Mormon clarifies understanding about:

The Book of Mormon reveals information previously unknown:

\newpage
\settalk{Introductory Remarks}
\setspeaker{Russell M. Nelson}
\settalkdate{April 2018}
\talkheading{Introductory Remarks}{Russell M. Nelson}

Thank you, Brother Holmes, for your important message.

Dear brethren, we deeply miss President Thomas S. Monson and Elder Robert D. Hales. Yet we “all press on in the work of the Lord.”

I am very grateful for each man who bears the holy priesthood. You are the hope of our Redeemer, who desires “that every man might speak in the name of God the Lord, even the Savior of the world.” He wants all of His ordained sons to represent Him, to speak for Him, to act for Him, and to bless the lives of God’s children throughout the world, to the end “that faith also might increase in [all] the earth.”

Some of you serve where the Church has been established for generations. Others serve where the Church is relatively new. For some, your wards are large. For others, your branches are small and distances are great. Regardless of your individual circumstances, each of you is a member of a priesthood quorum with a divine mandate to learn and to teach, to love and to serve others.

Tonight we announce a significant restructuring of our Melchizedek Priesthood quorums to accomplish the work of the Lord more effectively. In each ward, the high priests and the elders will now be combined into one elders quorum. This adjustment will greatly enhance the capacity and the ability of men who bear the priesthood to serve others. Prospective elders will be welcomed in and fellowshipped by that quorum. In each stake, the stake presidency will continue to preside over the stake high priests quorum. But the composition of that quorum will be based on current priesthood callings, as will be explained later.

Elder D. Todd Christofferson and Elder Ronald A. Rasband of the Quorum of the Twelve Apostles will now teach us more about these important adjustments.

These modifications have been under study for many months. We have felt a pressing need to improve the way we care for our members and report our contacts with them. To do that better, we need to strengthen our priesthood quorums to give greater direction to the ministering of love and support that the Lord intends for His Saints.

These adjustments are inspired of the Lord. As we implement them, we will be even more effective than we have ever been previously.

We are engaged in the work of Almighty God. Jesus is the Christ! We are His humble servants! God bless you, brethren, as we learn of and do our duty, I pray in the name of Jesus Christ, amen.

\newpage
\settalk{Let Us All Press On}
\setspeaker{Russell M. Nelson}
\settalkdate{April 2018}
\talkheading{Let Us All Press On}{Russell M. Nelson}

My beloved brothers and sisters, as we come to the close of this historic conference, I join with you in thanking the Lord for His direction and inspiring influence. The music has been beautiful and uplifting. Not only have the messages been edifying, but they have been life-changing!

In solemn assembly we sustained a new First Presidency. Two great men have been placed in the Quorum of the Twelve Apostles. And eight new General Authority Seventies have been called.

Now a favorite hymn summarizes our renewed resolve, our challenge, and our charge going forward:

\begin{verse}
Let us all press on in the work of the Lord,\\
That when life is o’er we may gain a reward;\\
In the fight for right let us wield a sword,\\
The mighty sword of truth.
\end{verse}

I exhort you to study the messages of this conference frequently—even repeatedly—during the next six months. Conscientiously look for ways to incorporate these messages in your family home evenings, your gospel teaching, your conversations with family and friends, and even your discussions with those not of our faith. Many good people will respond to the truths taught in this conference when offered in love. And your desire to obey will be enhanced as you remember and reflect upon what you have felt these past two days.

This general conference marks the beginning of a new era of ministering. The Lord has made important adjustments in the way we care for each other. Sisters and brothers—old and young—will serve one another in a new, holier way. Elders quorums will be strengthened to bless the lives of men, women, and children throughout the world. Relief Society sisters will continue to minister in their unique and loving way, extending opportunities to younger sisters to join them as appropriately assigned.

Our message to the world is simple and sincere: we invite all of God’s children on both sides of the veil to come unto their Savior, receive the blessings of the holy temple, have enduring joy, and qualify for eternal life.

Eventual exaltation requires our complete fidelity now to covenants we make and ordinances we receive in the house of the Lord. At this time, we have 159 functioning temples, and more are under construction. We want to bring temples closer to the expanding membership of the Church. So we are now pleased to announce plans to construct seven more temples. Those temples will be located in the following locations: Salta, Argentina; Bengaluru, India; Managua, Nicaragua; Cagayan de Oro, Philippines; Layton, Utah; Richmond, Virginia; and a major city yet to be determined in Russia.

My dear brothers and sisters, construction of these temples may not change your life, but your time in the temple surely will. In that spirit, I bless you to identify those things you can set aside so you can spend more time in the temple. I bless you with greater harmony and love in your homes and a deeper desire to care for your eternal family relationships. I bless you with increased faith in the Lord Jesus Christ and a greater ability to follow Him as His true disciples.

I bless you to raise your voice in testimony, as I do now, that we are engaged in the work of Almighty God! Jesus is the Christ. This is His Church, which He directs through His anointed servants. I so testify, with my expression of love for each of you, in the sacred name of Jesus Christ, amen.

\newpage
\settalk{Ministering}
\setspeaker{Russell M. Nelson}
\settalkdate{April 2018}
\talkheading{Ministering}{Russell M. Nelson}

Thank you, Elder Gong and Elder Soares, for your heartfelt expressions of faith. We are so very grateful for you and your dear companions.

Dear brothers and sisters, we constantly seek direction from the Lord on how we can help our members keep the commandments of God, especially those two great commandments to love God and our neighbors.

For months we have been seeking a better way to minister to the spiritual and temporal needs of our people in the Savior’s way.

We have made the decision to retire home teaching and visiting teaching as we have known them. Instead, we will implement a newer, holier approach to caring for and ministering to others. We will refer to these efforts simply as “ministering.”

Effective ministering efforts are enabled by the innate gifts of the sisters and by the incomparable power of the priesthood. We all need such protection from the cunning wiles of the adversary.

Elder Jeffrey R. Holland of the Quorum of the Twelve Apostles and Sister Jean B. Bingham, General President of the Relief Society, will explain how assigned brethren of the priesthood and assigned sisters of the Relief Society and Young Women will now function in serving and watching over members of the Church throughout the world.

The First Presidency and the Twelve are united in endorsing their messages. Gratefully and prayerfully we open this new chapter in the history of the Church. In the name of Jesus Christ, amen.

\newpage
\settalk{Ministering with the Power and Authority of God}
\setspeaker{Russell M. Nelson}
\settalkdate{April 2018}
\talkheading{Ministering with the Power and Authority of God}{Russell M. Nelson}

My beloved brethren, thank you for your devotion to the Lord and His holy work. It is truly a joy to be with you. As a new First Presidency, we thank you for your prayers and for your sustaining efforts. We are grateful for your lives and for your service to the Lord. Your devotion to duty and your selfless service are just as important in your callings as ours are in our callings. Through a lifetime of service in this Church, I have learned that it really doesn’t matter where one serves. What the Lord cares about is how one serves.

I express deep gratitude for President Thomas S. Monson, who was an example to me for more than 50 years. And for his counselors, President Henry B. Eyring and President Dieter F. Uchtdorf, I express profound admiration. I commend them for their service to the Lord and His prophets. Both of these devoted servants have received new assignments. They continue to serve with vigor and commitment. I honor and love them both.

It is a remarkable blessing to serve in the Lord’s true and living Church with His authority and power. The restoration of the priesthood of God, including the keys of the priesthood, opens to worthy Latter-day Saints the greatest of all spiritual blessings. We see those blessings flowing to women, men, and children throughout the world.

We see faithful women who understand the power inherent in their callings and in their endowment and other temple ordinances. These women know how to call upon the powers of heaven to protect and strengthen their husbands, their children, and others they love. These are spiritually strong women who lead, teach, and minister fearlessly in their callings with the power and authority of God! How thankful I am for them!

Likewise, we see faithful men who live up to their privileges as bearers of the priesthood. They lead and serve by sacrifice in the Lord’s way with love, kindness, and patience. They bless, guide, protect, and strengthen others by the power of the priesthood they hold. They bring miracles to those they serve while they keep their own marriages and families safe. They shun evil and are mighty elders in Israel. I am most thankful for them!

Now, may I voice a concern? It is this: Too many of our brothers and sisters do not fully understand the concept of priesthood power and authority. They act as though they would rather satisfy their own selfish desires and appetites than use the power of God to bless His children.

I fear that too many of our brothers and sisters do not grasp the privileges that could be theirs. Some of our brethren, for example, act like they do not understand what the priesthood is and what it enables them to do. Let me give you some specific examples.

Not long ago, I attended a sacrament meeting in which a new baby was to be given a name and a father’s blessing. The young father held his precious infant in his arms, gave her a name, and then offered a beautiful prayer. But he did not give that child a blessing. That sweet baby girl got a name but no blessing! That dear elder did not know the difference between a prayer and a priesthood blessing. With his priesthood authority and power, he could have blessed his infant, but he did not. I thought, “What a missed opportunity!”

Let me cite some other examples. We know of brethren who set sisters apart as Primary, Young Women, or Relief Society leaders and teachers but fail to bless them—to bless them with the power to fulfill their callings. They give only admonitions and instructions. We see a worthy father who fails to give his wife and his children priesthood blessings when that is exactly what they need. Priesthood power has been restored to this earth, and yet far too many brothers and sisters go through terrible trials in life without ever receiving a true priesthood blessing. What a tragedy! That’s a tragedy that we can eliminate.

Brethren, we hold the holy priesthood of God! We have His authority to bless His people. Just think of the remarkable assurance the Lord gave us when He said, “Whomsoever you bless I will bless.” It is our privilege to act in the name of Jesus Christ to bless God’s children according to His will for them. Stake presidents and bishops, please ensure that every member of the quorums within your stewardship understands how to give a priesthood blessing—including the personal worthiness and spiritual preparation required to call fully upon the power of God.

To all brethren holding the priesthood, I invite you to inspire members to keep their covenants, fast and pray, study the scriptures, worship in the temple, and serve with faith as men and women of God. We can help all to see with the eye of faith that obedience and righteousness will draw them closer to Jesus Christ, allow them to enjoy the companionship of the Holy Ghost, and experience joy in life!

A hallmark of the Lord’s true and living Church will always be an organized, directed effort to minister to individual children of God and their families. Because it is His Church, we as His servants will minister to the one, just as He did. We will minister in His name, with His power and authority, and with His loving-kindness.

An experience I had more than 60 years ago in Boston taught me just how powerful the privilege of ministering one-on-one can be. I was then a resident surgeon at the Massachusetts General Hospital—on duty every day, every other night, and every other weekend. I had limited time for my wife, our four children, and Church activity. Nonetheless, our branch president assigned me to visit the home of Wilbur and Leonora Cox with the hope that Brother Cox might come back into activity in the Church. He and Leonora had been sealed in the temple. Yet Wilbur had not participated for many years.

My companion and I went to their home. As we entered, Sister Cox welcomed us warmly, but Brother Cox abruptly walked into another room and closed the door.

I went to the closed door and knocked. After a moment, I heard a muffled “Come in.” I opened the door to find Brother Cox sitting beside an array of amateur radio equipment. In that small room, he lit up a cigar. Clearly, my visit was not all that welcome.

I gazed about the room with wonderment and said, “Brother Cox, I have always wanted to learn more about amateur radio work. Would you be willing to teach me about it? I’m sorry I can’t stay any longer tonight, but could I come back another time?”

He hesitated for a moment and then said yes. That was the beginning of what became a wonderful friendship. I returned and he taught me. I began to love and respect him. Through our subsequent visits, the greatness of this man emerged. We became very good friends, as did our dear eternal companions. Then, with the passage of time, our family moved away. Local leaders continued to nurture the Cox family.

About eight years after that first visit, the Boston Stake was created. Can you guess who its first stake president was? Yes! Brother Cox! During subsequent years, he also served as a mission president and a temple president.

Years later, I, as a member of the Quorum of the Twelve, was assigned to create a new stake in Sanpete County, Utah. During the usual interviews, I was pleasantly surprised to encounter again my dear friend Brother Cox! I felt impressed to call him as the new stake patriarch. After I ordained him, we embraced each other and wept. People in the room were wondering why these two grown men were crying. But we knew. And Sister Cox knew. Ours were tears of joy! We silently remembered the incredible journey of love and repentance that began more than 30 years ago, one night in their home.

The account doesn’t end there. Brother and Sister Cox’s family grew to include 3 children, 20 grandchildren, and 54 great-grandchildren. Add to that their impact on hundreds of missionaries, on thousands more in the temple, and on hundreds more who received patriarchal blessings at the hands of Wilbur Cox. His and Leonora’s influence will continue to ripple through many generations throughout the world.

Experiences such as this with Wilbur and Leonora Cox occur every week—hopefully, every day—within this Church. Dedicated servants of the Lord Jesus Christ carry out His work, with His power and authority.

Brethren, there are doors we can open, priesthood blessings we can give, hearts we can heal, burdens we can lift, testimonies we can strengthen, lives we can save, and joy we can bring into the homes of the Latter-day Saints—all because we hold the priesthood of God. We are the men who have been “called and prepared from the foundation of the world according to the foreknowledge of God, on account of [our] exceeding faith,” to do this work.

Tonight I invite you literally to rise up with me in our great eternal brotherhood. When I name your priesthood office, please stand and remain standing. Deacons, please arise! Teachers, arise! Priests! Bishops! Elders! High priests! Patriarchs! Seventies! Apostles!

Now, brethren, will you please remain standing and join with our chorus in singing all three verses of “Rise Up, O Men of God.” While you sing, think of your duty as God’s mighty army to help prepare the world for the Second Coming of the Lord. This is our charge. This is our privilege. I so testify in the name of Jesus Christ, amen.

\newpage
\settalk{Revelation for the Church, Revelation for Our Lives}
\setspeaker{Russell M. Nelson}
\settalkdate{April 2018}
\talkheading{Revelation for the Church, Revelation for Our Lives}{Russell M. Nelson}

What a glorious privilege it has been to celebrate Easter with you on this Sunday of general conference! Nothing could be more fitting than to commemorate the most important event that ever occurred on this earth by worshipping the most important being who ever walked this earth. In this, The Church of Jesus Christ of Latter-day Saints, we worship Him who commenced His infinite Atonement in the Garden of Gethsemane. He was willing to suffer for the sins and weaknesses of each of us, which suffering caused Him “to bleed at every pore.” He was crucified on Calvary’s cross and rose the third day as the first resurrected being of our Heavenly Father’s children. I love Him and testify that He lives! It is He who leads and guides His Church.

Without our Redeemer’s infinite Atonement, not one of us would have hope of ever returning to our Heavenly Father. Without His Resurrection, death would be the end. Our Savior’s Atonement made eternal life a possibility and immortality a reality for all.

It is because of His transcendent mission and the peace He grants His followers that my wife, Wendy, and I felt comfort late on January 2, 2018, when we were awakened by a phone call telling us that President Thomas S. Monson had stepped through the veil.

How we miss President Monson! We honor his life and his legacy. A spiritual giant, he left an indelible imprint upon all who knew him and upon the Church that he loved.

On Sunday, January 14, 2018, in the upper room of the Salt Lake Temple, the First Presidency was reorganized in the simple yet sacred pattern established by the Lord. Then, at yesterday morning’s solemn assembly, members of the Church throughout the world raised their hands to confirm the earlier action taken by the Apostles. I am humbly grateful for your sustaining support.

I am also grateful for those upon whose shoulders I stand. It has been my privilege to serve in the Quorum of the Twelve Apostles for 34 years and to know personally 10 of the 16 previous Presidents of the Church. I learned much from each of them.

I also owe much to my forebears. All eight of my great-grandparents were converts to the Church in Europe. Each of these stalwart souls sacrificed everything to come to Zion. During subsequent generations, however, not all my ancestors remained so committed. As a result, I was not raised in a gospel-centered home.

I adored my parents. They meant the world to me and taught me crucial lessons. I cannot thank them enough for the happy homelife they created for me and my siblings. And yet, even as a boy, I knew I was missing something. One day I jumped on the streetcar and went to an LDS bookstore to find a book about the Church. I loved learning about the gospel.

As I came to understand the Word of Wisdom, I wanted my parents to live that law. So, one day when I was very young, I went to our basement and smashed on the concrete floor every bottle of liquor! I expected my father to punish me, but he never said a word.

As I matured and began to understand the magnificence of Heavenly Father’s plan, I often said to myself, “I don’t want one more Christmas present! I just want to be sealed to my parents.” That longed-for event did not happen until my parents were past 80, and then it did happen. I cannot fully express the joy that I felt that day, and each day I feel that joy of their sealing and my being sealed to them.

In 1945, while I was in medical school, I married Dantzel White in the Salt Lake Temple. She and I were blessed with nine splendid daughters and one precious son. Today our ever-growing family is one of the greatest joys of my life.

In 2005, after nearly 60 years of marriage, my dear Dantzel was unexpectedly called home. For a season, my grief was almost immobilizing. But the message of Easter and the promise of resurrection sustained me.

Then the Lord brought Wendy Watson to my side. We were sealed in the Salt Lake Temple on April 6, 2006. How I love her! She is an extraordinary woman—a great blessing to me, to our family, and to the entire Church.

Each of these blessings has come as a result of seeking and heeding the promptings of the Holy Ghost. Said President Lorenzo Snow, “This is the grand privilege of every Latter-day Saint … that it is our right to have the manifestations of the Spirit every day of our lives.”

One of the things the Spirit has repeatedly impressed upon my mind since my new calling as President of the Church is how willing the Lord is to reveal His mind and will. The privilege of receiving revelation is one of the greatest gifts of God to His children.

Through the manifestations of the Holy Ghost, the Lord will assist us in all our righteous pursuits. I remember in an operating room, I have stood over a patient—unsure how to perform an unprecedented procedure—and experienced the Holy Ghost diagramming the technique in my mind.

To strengthen my proposal to Wendy, I said to her, “I know about revelation and how to receive it.” To her credit—and, as I have come to learn, typical of her—she had already sought and received her own revelation about us, which gave her the courage to say yes.

As a member of the Quorum of the Twelve Apostles, I prayed daily for revelation and gave thanks to the Lord every time He spoke to my heart and mind.

Imagine the miracle of it! Whatever our Church calling, we can pray to our Heavenly Father and receive guidance and direction, be warned about dangers and distractions, and be enabled to accomplish things we simply could not do on our own. If we will truly receive the Holy Ghost and learn to discern and understand His promptings, we will be guided in matters large and small.

When I recently faced the daunting task of choosing two counselors, I wondered how I could possibly choose just two from twelve men whom I love and respect.

Because I know that good inspiration is based upon good information, I prayerfully met one-on-one with each Apostle. I then sequestered myself in a private room in the temple and sought the Lord’s will. I testify that the Lord instructed me to select President Dallin H. Oaks and President Henry B. Eyring to serve as my counselors in the First Presidency.

In like manner, I testify that the Lord inspired the call of Elder Gerrit W. Gong and Elder Ulisses Soares to be ordained as His Apostles. I and we welcome them to this unique brotherhood of service.

When we convene as a Council of the First Presidency and Quorum of the Twelve, our meeting rooms become rooms of revelation. The Spirit is palpably present. As we wrestle with complex matters, a thrilling process unfolds as each Apostle freely expresses his thoughts and point of view. Though we may differ in our initial perspectives, the love we feel for each other is constant. Our unity helps us to discern the Lord’s will for His Church.

In our meetings, the majority never rules! We listen prayerfully to one another and talk with each other until we are united. Then when we have reached complete accord, the unifying influence of the Holy Ghost is spine-tingling! We experience what the Prophet Joseph Smith knew when he taught, “By union of feeling we obtain power with God.” No member of the First Presidency or Quorum of the Twelve would ever leave decisions for the Lord’s Church to his own best judgment!

Brothers and sisters, how can we become the men and women—the Christlike servants—the Lord needs us to be? How can we find answers to questions that perplex us? If Joseph Smith’s transcendent experience in the Sacred Grove teaches us anything, it is that the heavens are open and that God speaks to His children.

The Prophet Joseph Smith set a pattern for us to follow in resolving our questions. Drawn to the promise of James that if we lack wisdom we may ask of God, the boy Joseph took his question directly to Heavenly Father. He sought personal revelation, and his seeking opened this last dispensation.

In like manner, what will your seeking open for you? What wisdom do you lack? What do you feel an urgent need to know or understand? Follow the example of the Prophet Joseph. Find a quiet place where you can regularly go. Humble yourself before God. Pour out your heart to your Heavenly Father. Turn to Him for answers and for comfort.

Pray in the name of Jesus Christ about your concerns, your fears, your weaknesses—yes, the very longings of your heart. And then listen! Write the thoughts that come to your mind. Record your feelings and follow through with actions that you are prompted to take. As you repeat this process day after day, month after month, year after year, you will “grow into the principle of revelation.”

Does God really want to speak to you? Yes! “As well might man stretch forth his puny arm to stop the Missouri river in its decreed course … as to hinder the Almighty from pouring down knowledge from heaven upon the heads of the Latter-day Saints.”

You don’t have to wonder about what is true. You do not have to wonder whom you can safely trust. Through personal revelation, you can receive your own witness that the Book of Mormon is the word of God, that Joseph Smith is a prophet, and that this is the Lord’s Church. Regardless of what others may say or do, no one can ever take away a witness borne to your heart and mind about what is true.

I urge you to stretch beyond your current spiritual ability to receive personal revelation, for the Lord has promised that “if thou shalt [seek], thou shalt receive revelation upon revelation, knowledge upon knowledge, that thou mayest know the mysteries and peaceable things—that which bringeth joy, that which bringeth life eternal.”

Oh, there is so much more that your Father in Heaven wants you to know. As Elder Neal A. Maxwell taught, “To those who have eyes to see and ears to hear, it is clear that the Father and the Son are giving away the secrets of the universe!”

Nothing opens the heavens quite like the combination of increased purity, exact obedience, earnest seeking, daily feasting on the words of Christ in the Book of Mormon, and regular time committed to temple and family history work.

To be sure, there may be times when you feel as though the heavens are closed. But I promise that as you continue to be obedient, expressing gratitude for every blessing the Lord gives you, and as you patiently honor the Lord’s timetable, you will be given the knowledge and understanding you seek. Every blessing the Lord has for you—even miracles—will follow. That is what personal revelation will do for you.

I am optimistic about the future. It will be filled with opportunities for each of us to progress, contribute, and take the gospel to every corner of the earth. But I am also not naive about the days ahead. We live in a world that is complex and increasingly contentious. The constant availability of social media and a 24-hour news cycle bombard us with relentless messages. If we are to have any hope of sifting through the myriad of voices and the philosophies of men that attack truth, we must learn to receive revelation.

Our Savior and Redeemer, Jesus Christ, will perform some of His mightiest works between now and when He comes again. We will see miraculous indications that God the Father and His Son, Jesus Christ, preside over this Church in majesty and glory. But in coming days, it will not be possible to survive spiritually without the guiding, directing, comforting, and constant influence of the Holy Ghost.

My beloved brothers and sisters, I plead with you to increase your spiritual capacity to receive revelation. Let this Easter Sunday be a defining moment in your life. Choose to do the spiritual work required to enjoy the gift of the Holy Ghost and hear the voice of the Spirit more frequently and more clearly.

With Moroni, I exhort you on this Easter Sabbath to “come unto Christ, and lay hold upon every good gift,” beginning with the gift of the Holy Ghost, which gift can and will change your life.

We are followers of Jesus Christ. The most important truth the Holy Ghost will ever witness to you is that Jesus is the Christ, the Son of the living God. He lives! He is our Advocate with the Father, our Exemplar, and our Redeemer. On this Easter Sunday, we commemorate His atoning sacrifice, His literal Resurrection, and His divinity.

This is His Church, restored through the Prophet Joseph Smith. I so testify, with my expression of love for each of you, in the sacred name of Jesus Christ, amen.

\newpage
\settalk{Becoming Exemplary Latter-day Saints}
\setspeaker{Russell M. Nelson}
\settalkdate{October 2018}
\talkheading{Becoming Exemplary Latter-day Saints}{Russell M. Nelson}

This has been an inspirational and historic conference. We look to the future with enthusiasm. We have been motivated to do better and to be better. The marvelous messages delivered from this pulpit by our General Authorities and General Officers and the music have been sublime! I urge you to study these messages, commencing this week. They express the mind and the will of the Lord for His people, today.

The new home-centered, Church-supported integrated curriculum has the potential to unleash the power of families, as each family follows through conscientiously and carefully to transform their home into a sanctuary of faith. I promise that as you diligently work to remodel your home into a center of gospel learning, over time your Sabbath days will truly be a delight. Your children will be excited to learn and to live the Savior’s teachings, and the influence of the adversary in your life and in your home will decrease. Changes in your family will be dramatic and sustaining.

During this conference we have strengthened our resolve to execute the essential effort to honor the Lord Jesus Christ every time we refer to His Church. I promise you that our rigorous attention to use the correct name of the Savior’s Church and its members will lead to increased faith and access to greater spiritual power for members of His Church.

Now let’s turn to the topic of temples. We know that our time in the temple is crucial to our salvation and exaltation and to that of our families.

After we receive our own temple ordinances and make sacred covenants with God, each one of us needs the ongoing spiritual strengthening and tutoring that is possible only in the house of the Lord. And our ancestors need us to serve as proxy for them.

Consider the great mercy and fairness of God, who, before the foundation of the world, provided a way to give temple blessings to those who died without a knowledge of the gospel. These sacred temple rites are ancient. To me that antiquity is thrilling and another evidence of their authenticity.

My dear brothers and sisters, the assaults of the adversary are increasing exponentially, in intensity and in variety. Our need to be in the temple on a regular basis has never been greater. I plead with you to take a prayerful look at how you spend your time. Invest time in your future and in that of your family. If you have reasonable access to a temple, I urge you to find a way to make an appointment regularly with the Lord—to be in His holy house—then keep that appointment with exactness and joy. I promise you that the Lord will bring the miracles He knows you need as you make sacrifices to serve and worship in His temples.

Currently we have 159 dedicated temples. The proper care and maintenance of those temples is very important to us. With the passage of time, temples are inevitably in need of refreshing and renewal. To that end, plans are now being made to renovate and update the Salt Lake Temple and other pioneer-generation temples. Details on these projects will be shared as they are developed.

Today we are pleased to announce plans to construct 12 more temples. Those temples will be built in the following locations: Mendoza, Argentina; Salvador, Brazil; Yuba City, California; Phnom Penh, Cambodia; Praia, Cape Verde; Yigo, Guam; Puebla, Mexico; Auckland, New Zealand; Lagos, Nigeria; Davao, Philippines; San Juan, Puerto Rico; and Washington County, Utah.

Building and maintaining temples may not change your life, but spending your time in the temple surely will. To those who have long been absent from the temple, I encourage you to prepare and return as soon as possible. Then I invite you to worship in the temple and pray to feel deeply the Savior’s infinite love for you, that each of you may gain your own testimony that He directs this sacred and ageless work.

Brothers and sisters, I thank you for your faith and sustaining efforts. I leave my love and blessing upon you, that you may feast upon the word of the Lord and apply His teachings in your personal lives. I assure you that revelation continues in the Church and will continue until “the purposes of God shall be accomplished, and the Great Jehovah shall say the work is done.”

I bless you with increased faith in Him and in His holy work, with faith and patience to endure your personal challenges in life. I bless you to become exemplary Latter-day Saints. I so bless you and bear my testimony that God lives! Jesus is the Christ! This is His Church. We are His people, in the name of Jesus Christ, amen.

\newpage
\settalk{Opening Remarks}
\setspeaker{Russell M. Nelson}
\settalkdate{October 2018}
\talkheading{Opening Remarks}{Russell M. Nelson}

My dear brothers and sisters, we have looked forward to gathering with you again in this October general conference of the Church. We extend our heartfelt welcome to each of you. We are deeply grateful for your sustaining prayers. We can feel their effect. Thank you!

We are grateful for your tremendous efforts to follow through with counsel given at general conference six months ago. Stake presidencies throughout the world have sought revelation necessary to reorganize elders quorums. Men of those quorums along with our dedicated Relief Society sisters are working diligently to minister to our brothers and sisters in a higher, holier way. We are inspired by your goodness and extraordinary efforts to bring the Savior’s love to your families, neighbors, and friends and to minister to them as He would.

Since April conference, Sister Nelson and I have met with members on four continents and on islands of the sea. From Jerusalem to Harare, from Winnipeg to Bangkok, we have experienced your great faith and the strength of your testimonies.

We are overjoyed with the number of our youth who have joined the Lord’s youth battalion to help gather scattered Israel. We thank you! And as you continue to follow my invitations offered at our worldwide youth devotional, you are setting a standard for the rest of us to follow. And what a difference you youth are making!

In recent years, we in the presiding councils of the Church have wrestled with a fundamental question: how can we take the gospel in its simple purity and the ordinances with their eternal efficacy to all of God’s children?

As Latter-day Saints, we have become accustomed to thinking of “church” as something that happens in our meetinghouses, supported by what happens at home. We need an adjustment to this pattern. It is time for a home-centered Church, supported by what takes place inside our branch, ward, and stake buildings.

As the Church continues to expand throughout the world, many members live where we have no chapels—and may not for the foreseeable future. I remember a family who, because of such circumstances, was required to meet in their home. I asked the mother how she liked going to church in her own home. She replied, “I like it! My husband uses better language at home now, knowing that he will bless the sacrament here each Sunday.”

The long-standing objective of the Church is to assist all members to increase their faith in our Lord Jesus Christ and in His Atonement, to assist them in making and keeping their covenants with God, and to strengthen and seal their families. In this complex world today, this is not easy. The adversary is increasing his attacks on faith and upon us and our families at an exponential rate. To survive spiritually, we need counterstrategies and proactive plans. Accordingly, we now want to put in place organizational adjustments that will further fortify our members and their families.

For many years, Church leaders have been working on an integrated curriculum to strengthen families and individuals through a home-centered and Church-supported plan to learn doctrine, strengthen faith, and foster greater personal worship. Our efforts over these recent years to hallow the Sabbath—to make it a delight and a personal sign to God of our love for Him—will be augmented by the adjustments we will now introduce.

This morning we will announce a new balance and connection between gospel instruction in the home and in the Church. We are each responsible for our individual spiritual growth. And scriptures make it clear that parents have the primary responsibility to teach the doctrine to their children. It is the responsibility of the Church to assist each member in the divinely defined goal of increasing his or her gospel knowledge.

Elder Quentin L. Cook will now explain these important adjustments. All members of the Council of the First Presidency and Quorum of the Twelve Apostles are united in endorsing this message. We gratefully acknowledge the inspiration from the Lord that has influenced the development of the plans and procedures that Elder Cook will present.

My dear brothers and sisters, I know that God lives! Jesus is the Christ! This is His Church that He directs by prophecy and revelation to His humble servants. I so testify in the name of Jesus Christ, amen.

\newpage
\settalk{Sisters’ Participation in the Gathering of Israel}
\setspeaker{Russell M. Nelson}
\settalkdate{October 2018}
\talkheading{Sisters’ Participation in the Gathering of Israel}{Russell M. Nelson}

It is wonderful to be with you, my dear and precious sisters. Perhaps a recent experience will give you a glimpse into how I feel about you and the supernal abilities with which you are endowed.

One day while I was speaking to a congregation in South America, I became exceedingly excited about my topic, and at a pivotal moment, I said, “As the mother of 10 children, I can tell you that …” And then I went on to complete my message.

I did not realize that I had said the word mother. My translator, assuming I had misspoken, changed the word mother to father, so the congregation never knew that I had referred to myself as mother. But my wife Wendy heard it, and she was delighted with my Freudian slip.

In that moment, the deep longing of my heart to make a difference in the world—like only a mother does—bubbled up from my heart. Through the years, whenever I have been asked why I chose to become a medical doctor, my answer has always been the same: “Because I could not choose to be a mother.”

Please note that anytime I use the word mother, I am not talking only about women who have given birth or adopted children in this life. I am speaking about all of our Heavenly Parents’ adult daughters. Every woman is a mother by virtue of her eternal divine destiny.

So tonight, as the father of 10 children—nine daughters and one son—and as President of the Church, I pray that you will sense how deeply I feel about you—about who you are and all the good you can do. No one can do what a righteous woman can do. No one can duplicate the influence of a mother.

Men can and often do communicate the love of Heavenly Father and the Savior to others. But women have a special gift for it—a divine endowment. You have the capacity to sense what someone needs—and when he or she needs it. You can reach out, comfort, teach, and strengthen someone in his or her very moment of need.

Women see things differently than men do, and oh, how we need your perspective! Your nature leads you to think of others first, to consider the effect that any course of action will have on others.

As President Eyring pointed out, it was our glorious Mother Eve—with her far-reaching vision of our Heavenly Father’s plan—who initiated what we call ‘‘the Fall.” Her wise and courageous choice and Adam’s supporting decision moved God’s plan of happiness forward. They made it possible for each of us to come to earth, receive a body, and prove that we would choose to stand up for Jesus Christ now, just as we did premortally.

My dear sisters, you have special spiritual gifts and propensities. Tonight I urge you, with all the hope of my heart, to pray to understand your spiritual gifts—to cultivate, use, and expand them, even more than you ever have. You will change the world as you do so.

As women, you inspire others and set a standard worthy of emulation. Let me give you a little background on two of the major announcements made at our last general conference. You, my dear sisters, were key to each.

First, ministering. The supreme standard for ministering is that of our Savior, Jesus Christ. Generally, women are, and always have been, closer to that standard than men. When you are truly ministering, you follow your feelings to help someone else experience more of the Savior’s love. The inclination to minister is inherent in righteous women. I know women who pray every day, “Whom wouldst Thou have me help today?”

Before the April 2018 announcement about the higher and holier way of caring for others, the tendency of some men was to check off their home teaching assignment as “done” and move on to the next task.

But when you sensed that a sister you visit taught needed help, you responded immediately and then throughout the month. Thus, it was how you visit taught that inspired our upward shift to ministering.

Second, in the last general conference, we also restructured Melchizedek Priesthood quorums. When we wrestled with how to help the men of the Church be more effective in their responsibilities, we carefully considered the example of the Relief Society.

In Relief Society, women in various ages and stages of life meet together. Each decade of life brings unique challenges, and yet, there you were, week after week, mingling together, growing and teaching the gospel together, and making a real difference in the world.

Now, following your example, Melchizedek Priesthood bearers are members of the elders quorum. These men range in age from 18 to 98 (maybe more), with equally wide-ranging priesthood and Church experiences. These brethren can now create stronger fraternal linkages, learn together, and bless others more effectively.

You remember that last June, Sister Nelson and I spoke to the youth of the Church. We invited them to enlist in the Lord’s youth battalion to help gather Israel on both sides of the veil. This gathering is “the greatest challenge, the greatest cause, and the greatest work on earth today”!

It is a cause that desperately needs women, because women shape the future. So tonight I’m extending a prophetic plea to you, the women of the Church, to shape the future by helping to gather scattered Israel.

Where can you start?

May I offer four invitations:

First, I invite you to participate in a 10-day fast from social media and from any other media that bring negative and impure thoughts to your mind. Pray to know which influences to remove during your fast. The effect of your 10-day fast may surprise you. What do you notice after taking a break from perspectives of the world that have been wounding your spirit? Is there a change in where you now want to spend your time and energy? Have any of your priorities shifted—even just a little? I urge you to record and follow through with each impression.

Second, I invite you to read the Book of Mormon between now and the end of the year. As impossible as that may seem with all you are trying to manage in your life, if you will accept this invitation with full purpose of heart, the Lord will help you find a way to achieve it. And, as you prayerfully study, I promise that the heavens will open for you. The Lord will bless you with increased inspiration and revelation.

As you read, I would encourage you to mark each verse that speaks of or refers to the Savior. Then, be intentional about talking of Christ, rejoicing in Christ, and preaching of Christ with your families and friends. You and they will be drawn closer to the Savior through this process. And changes, even miracles, will begin to happen.

This morning the announcement was made regarding the new Sunday schedule and home-centered, Church-supported curriculum. You, my dear sisters, are a key to the success of this new, balanced, and coordinated gospel-teaching effort. Please teach those whom you love what you are learning from the scriptures. Teach them how to turn to the Savior for His healing and cleansing power when they sin. And teach them how to draw upon His strengthening power every day of their lives.

Third, establish a pattern of regular temple attendance. This may require a little more sacrifice in your life. More regular time in the temple will allow the Lord to teach you how to draw upon His priesthood power with which you have been endowed in His temple. For those of you who don’t live near a temple, I invite you to study prayerfully about temples in the scriptures and in the words of living prophets. Seek to know more, to understand more, to feel more about temples than you ever have before.

In our worldwide youth devotional last June, I spoke about a young man whose life changed when his parents exchanged his smartphone for a flip phone. This young man’s mother is a fearless woman of faith. She saw her son drifting toward choices that could prevent him from serving a mission. She took her pleadings to the temple to know how best to help her son. Then she followed through with every impression.

She said: “I felt the Spirit guiding me to check my son’s phone at specific times to catch specific things. I don’t know how to navigate these smartphones, but the Spirit guided me through all the social media that I don’t even use! I know the Spirit helps parents who are seeking guidance to protect their children. [At first] my son was furious with me. … But after only three days, he thanked me! He could feel the difference.”

Her son’s behavior and attitudes changed dramatically. He became more helpful at home, smiled more, and was more attentive at church. He loved serving for a time in a temple baptistry and preparing for his mission.

My fourth invitation, for you who are of age, is to participate fully in Relief Society. I urge you to study the current Relief Society purpose statement. It is inspiring. It may guide you in developing your own purpose statement for your own life. I also entreat you to savor the truths in the Relief Society declaration published almost 20 years ago. A framed copy of this declaration hangs on the wall in the office of the First Presidency. I am thrilled every time I read it. It describes who you are and who the Lord needs you to be at this precise time as you do your part to help gather scattered Israel.

My dear sisters, we need you! We “need your strength, your conversion, your conviction, your ability to lead, your wisdom, and your voices.” We simply cannot gather Israel without you.

I love you and thank you and now bless you with the ability to leave the world behind as you assist in this crucial and urgent work. Together we can do all that our Heavenly Father needs us to do to prepare the world for the Second Coming of His Beloved Son.

Jesus is the Christ. This is His Church. Of this I testify in the name of Jesus Christ, amen.

\newpage
\settalk{The Correct Name of the Church}
\setspeaker{Russell M. Nelson}
\settalkdate{October 2018}
\talkheading{The Correct Name of the Church}{Russell M. Nelson}

My beloved brothers and sisters, on this beautiful Sabbath day we rejoice together in our many blessings from the Lord. We are very grateful for your testimonies of the restored gospel of Jesus Christ, for the sacrifices you’ve made to stay on or return to His covenant path, and for your consecrated service in His Church.

Today I feel compelled to discuss with you a matter of great importance. Some weeks ago, I released a statement regarding a course correction for the name of the Church. I did this because the Lord impressed upon my mind the importance of the name He decreed for His Church, even The Church of Jesus Christ of Latter-day Saints.

As you would expect, responses to this statement and to the revised style guide have been mixed. Many members immediately corrected the name of the Church on their blogs and social media pages. Others wondered why, with all that’s going on in the world, it was necessary to emphasize something so “inconsequential.” And some said it couldn’t be done, so why even try? Let me explain why we care so deeply about this issue. But first let me state what this effort is not:

Instead, it is a correction. It is the command of the Lord. Joseph Smith did not name the Church restored through him; neither did Mormon. It was the Savior Himself who said, “For thus shall my church be called in the last days, even The Church of Jesus Christ of Latter-day Saints.”

Even earlier, in AD 34, our resurrected Lord gave similar instruction to members of His Church when He visited them in the Americas. At that time He said:

“Ye shall call the church in my name. …

“And how be it my church save it be called in my name? For if a church be called in Moses’ name then it be Moses’ church; or if it be called in the name of a man then it be the church of a man; but if it be called in my name then it is my church.”

Thus, the name of the Church is not negotiable. When the Savior clearly states what the name of His Church should be and even precedes His declaration with, “Thus shall my church be called,” He is serious. And if we allow nicknames to be used or adopt or even sponsor those nicknames ourselves, He is offended.

What’s in a name or, in this case, a nickname? When it comes to nicknames of the Church, such as the “LDS Church,” the “Mormon Church,” or the “Church of the Latter-day Saints,” the most important thing in those names is the absence of the Savior’s name. To remove the Lord’s name from the Lord’s Church is a major victory for Satan. When we discard the Savior’s name, we are subtly disregarding all that Jesus Christ did for us—even His Atonement.

Consider this from His perspective: Premortally, He was Jehovah, God of the Old Testament. Under the direction of His Father, He was the Creator of this and other worlds. He chose to submit to the will of His Father and do something for all of God’s children that no one else could do! Condescending to come to earth as the Only Begotten of the Father in the flesh, He was brutally reviled, mocked, spit upon, and scourged. In the Garden of Gethsemane, our Savior took upon Himself every pain, every sin, and all of the anguish and suffering ever experienced by you and me and by everyone who has ever lived or will ever live. Under the weight of that excruciating burden, He bled from every pore. All of this suffering was intensified as He was cruelly crucified on Calvary’s cross.

Through these excruciating experiences and His subsequent Resurrection—His infinite Atonement—He granted immortality to all and ransomed each one of us from the effects of sin on condition of our repentance.

Following the Savior’s Resurrection and the death of His Apostles, the world plunged into centuries of darkness. Then in the year 1820, God the Father and His Son, Jesus Christ, appeared to the Prophet Joseph Smith to initiate the Restoration of the Lord’s Church.

After all He had endured—and after all He had done for humankind—I realize with profound regret that we have unwittingly acquiesced in the Lord’s restored Church being called by other names, each of which expunges the sacred name of Jesus Christ!

Every Sunday as we worthily partake of the sacrament, we make anew our sacred promise to our Heavenly Father that we are willing to take upon us the name of His Son, Jesus Christ. We promise to follow Him, repent, keep His commandments, and always remember Him.

When we omit His name from His Church, we are inadvertently removing Him as the central focus of our lives.

Taking the Savior’s name upon us includes declaring and witnessing to others—through our actions and our words—that Jesus is the Christ. Have we been so afraid to offend someone who called us “Mormons” that we have failed to defend the Savior Himself, to stand up for Him even in the name by which His Church is called?

If we as a people and as individuals are to have access to the power of the Atonement of Jesus Christ—to cleanse and heal us, to strengthen and magnify us, and ultimately to exalt us—we must clearly acknowledge Him as the source of that power. We can begin by calling His Church by the name He decreed.

For much of the world, the Lord’s Church is presently disguised as the “Mormon Church.” But we as members of the Lord’s Church know who stands at its head: Jesus Christ Himself. Unfortunately, many who hear the term Mormon may think that we worship Mormon. Not so! We honor and respect that great ancient American prophet. But we are not Mormon’s disciples. We are the Lord’s disciples.

In the early days of the restored Church, terms such as Mormon Church and Mormons were often used as epithets—as cruel terms, abusive terms—designed to obliterate God’s hand in restoring the Church of Jesus Christ in these latter days.

Brothers and sisters, there are many worldly arguments against restoring the correct name of the Church. Because of the digital world in which we live and with search engine optimization that helps all of us find information we need almost instantly—including information about the Lord’s Church—critics say that a correction at this point is unwise. Others feel that because we are known so widely as “Mormons” and as the “Mormon Church,” we should make the best of it.

If this were a discussion about branding a man-made organization, those arguments might prevail. But in this crucial matter, we look to Him whose Church this is and acknowledge that the Lord’s ways are not, and never will be, man’s ways. If we will be patient and if we will do our part well, the Lord will lead us through this important task. After all, we know that the Lord helps those who seek to do His will, just as He helped Nephi accomplish the task of building a ship to cross the sea.

We will want to be courteous and patient in our efforts to correct these errors. Responsible media will be sympathetic in responding to our request.

In a previous general conference, Elder Benjamín De Hoyos spoke of such an event. He said:

“Some years ago while serving in the office of public affairs of the Church in Mexico, [a companion and I] were invited to participate in a radio talk show. … [One of the program directors] asked [us], ‘Why does the Church have such a long name? …’

“My companion and I smiled at such a magnificent question and then proceeded to explain that the name of the Church was not chosen by man. It was given by the Savior. … The program director immediately and respectfully responded, ‘We will thus repeat it with great pleasure.’”

That report provides a pattern. One by one, our best efforts as individuals will be required to correct errors that have crept in through the years. The rest of the world may or may not follow our lead in calling us by the correct name. But it is disingenuous for us to be frustrated if most of the world calls the Church and its members by the wrong names if we do the same.

Our revised style guide is helpful. It states: “In the first reference, the full name of the Church is preferred: ‘The Church of Jesus Christ of Latter-day Saints.’ When a shortened [second] reference is needed, the terms ‘the Church’ or the ‘Church of Jesus Christ’ are encouraged. The ‘restored Church of Jesus Christ’ is also accurate and encouraged.”

If someone should ask, “Are you a Mormon?” you could reply, “If you are asking if I am a member of The Church of Jesus Christ of Latter-day Saints, yes, I am!”

If someone asks, “Are you a Latter-day Saint?” you might respond, “Yes, I am. I believe in Jesus Christ and am a member of His restored Church.”

My dear brothers and sisters, I promise you that if we will do our best to restore the correct name of the Lord’s Church, He whose Church this is will pour down His power and blessings upon the heads of the Latter-day Saints, the likes of which we have never seen. We will have the knowledge and power of God to help us take the blessings of the restored gospel of Jesus Christ to every nation, kindred, tongue, and people and to prepare the world for the Second Coming of the Lord.

So, what’s in a name? When it comes to the name of the Lord’s Church, the answer is “Everything!” Jesus Christ directed us to call the Church by His name because it is His Church, filled with His power.

I know that God lives. Jesus is the Christ. He leads His Church today. I so testify in the sacred name of Jesus Christ, amen.

\newpage
\settalk{We Can Do Better and Be Better}
\setspeaker{Russell M. Nelson}
\settalkdate{April 2019}
\talkheading{We Can Do Better and Be Better}{Russell M. Nelson}

My dear brethren, it is inspiring to look out over this vast congregation of the Lord’s battalion of priesthood bearers. What a mighty force for good you are! We love you. We pray for you. And we are most grateful for you.

Recently I have found myself drawn to the Lord’s instruction given through the Prophet Joseph Smith: “Say nothing but repentance unto this generation.” This declaration is often repeated throughout scripture. It prompts an obvious question: “Does everyone need to repent?” The answer is yes.

Too many people consider repentance as punishment—something to be avoided except in the most serious circumstances. But this feeling of being penalized is engendered by Satan. He tries to block us from looking to Jesus Christ, who stands with open arms, hoping and willing to heal, forgive, cleanse, strengthen, purify, and sanctify us.

The word for repentance in the Greek New Testament is metanoeo. The prefix meta- means “change.” The suffix -noeo is related to Greek words that mean “mind,” “knowledge,” “spirit,” and “breath.”

Thus, when Jesus asks you and me to “repent,” He is inviting us to change our mind, our knowledge, our spirit—even the way we breathe. He is asking us to change the way we love, think, serve, spend our time, treat our wives, teach our children, and even care for our bodies.

Nothing is more liberating, more ennobling, or more crucial to our individual progression than is a regular, daily focus on repentance. Repentance is not an event; it is a process. It is the key to happiness and peace of mind. When coupled with faith, repentance opens our access to the power of the Atonement of Jesus Christ.

Whether you are diligently moving along the covenant path, have slipped or stepped from the covenant path, or can’t even see the path from where you are now, I plead with you to repent. Experience the strengthening power of daily repentance—of doing and being a little better each day.

When we choose to repent, we choose to change! We allow the Savior to transform us into the best version of ourselves. We choose to grow spiritually and receive joy—the joy of redemption in Him. When we choose to repent, we choose to become more like Jesus Christ!

Brethren, we need to do better and be better because we are in a battle. The battle with sin is real. The adversary is quadrupling his efforts to disrupt testimonies and impede the work of the Lord. He is arming his minions with potent weapons to keep us from partaking of the joy and love of the Lord.

Repentance is the key to avoiding misery inflicted by traps of the adversary. The Lord does not expect perfection from us at this point in our eternal progression. But He does expect us to become increasingly pure. Daily repentance is the pathway to purity, and purity brings power. Personal purity can make us powerful tools in the hands of God. Our repentance—our purity—will empower us to help in the gathering of Israel.

The Lord taught the Prophet Joseph Smith “that the rights of the priesthood are inseparably connected with the powers of heaven, and that the powers of heaven cannot be controlled nor handled only upon the principles of righteousness.”

We know what will give us greater access to the powers of heaven. We also know what will hinder our progress—what we need to stop doing to increase our access to the powers of heaven. Brethren, prayerfully seek to understand what stands in the way of your repentance. Identify what stops you from repenting. And then, change! Repent! All of us can do better and be better than ever before.

There are specific ways in which we can likely improve. One is in the way we treat our bodies. I stand in awe of the miracle of the human body. It is a magnificent creation, essential to our gradual ascent toward our ultimate divine potential. We cannot progress without it. In giving us the gift of a body, God has allowed us to take a vital step toward becoming more like Him.

Satan understands this. He chafes at the fact that his premortal apostasy permanently disqualifies him from this privilege, leaving him in a constant state of jealousy and resentment. Thus many, if not most, of the temptations he puts in our path cause us to abuse our bodies or the bodies of others. Because Satan is miserable without a body, he wants us to be miserable because of ours.

Your body is your personal temple, created to house your eternal spirit. Your care of that temple is important. Now, I ask you, brethren, are you more interested in dressing and grooming your body to appeal to the world than to please God? Your answer to this question sends a direct message to Him about your feelings regarding His transcendent gift to you. In this reverence for our bodies, brethren, I think we can do better and be better.

Another way we can also do better and be better is how we honor the women in our lives, beginning with our wives and daughters, our mothers and sisters.

Months ago, I received a heartbreaking letter from a dear sister. She wrote: “[My daughters and I] feel we are in fierce competition for our husbands’ and sons’ undivided attention, with 24/7 sports updates, video games, stock market updates, [and] endless analyzing and watching of games of every [conceivable] sport. It feels like we’re losing our front-row seats with our husbands and sons because of their permanent front-row seats with [sports and games].”

Brethren, your first and foremost duty as a bearer of the priesthood is to love and care for your wife. Become one with her. Be her partner. Make it easy for her to want to be yours. No other interest in life should take priority over building an eternal relationship with her. Nothing on TV, a mobile device, or a computer is more important than her well-being. Take an inventory of how you spend your time and where you devote your energy. That will tell you where your heart is. Pray to have your heart attuned to your wife’s heart. Seek to bring her joy. Seek her counsel, and listen. Her input will improve your output.

If you have a need to repent because of the way you have treated the women closest to you, begin now. And remember that it is your responsibility to help the women in your life receive the blessings that derive from living the Lord’s law of chastity. Never be the reason that a woman is unable to receive her temple blessings.

Brethren, we all need to repent. We need to get up off the couch, put down the remote, and wake up from our spiritual slumber. It is time to put on the full armor of God so we can engage in the most important work on earth. It is time to “thrust in [our] sickles, and reap with all [our] might, mind, and strength.” The forces of evil have never raged more forcefully than they do today. As servants of the Lord, we cannot be asleep while this battle rages.

Your family needs your leadership and love. Your quorum and those in your ward or branch need your strength. And all who meet you need to know what a true disciple of the Lord looks like and acts like.

My dear brethren, you were chosen by our Father to come to earth at this crucial time because of your premortal spiritual valor. You are among the finest, most valiant men who have ever come to the earth. Satan knows who you are and who you were premortally, and he understands the work that must be done before the Savior returns. And after millennia of practicing his cunning arts, the adversary is experienced and incorrigible.

Gratefully, the priesthood we hold is far stronger than are the wiles of the adversary. I plead with you to be the men and young men the Lord needs you to be. Make your focus on daily repentance so integral to your life that you can exercise the priesthood with greater power than ever before. This is the only way you will keep yourself and your family spiritually safe in the challenging days ahead.

The Lord needs selfless men who put the welfare of others ahead of their own. He needs men who intentionally work to hear the voice of the Spirit with clarity. He needs men of the covenant who keep their covenants with integrity. He needs men who are determined to keep themselves sexually pure—worthy men who can be called upon at a moment’s notice to give blessings with pure hearts, clean minds, and willing hands. The Lord needs men eager to repent—men with a zeal to serve and be part of the Lord’s battalion of worthy priesthood bearers.

I bless you to become those men. I bless you with the courage to repent daily and learn how to exercise full priesthood power. I bless you to communicate the love of the Savior to your wife and children and to all who know you. I bless you to do better and be better. And I bless you that as you make these efforts, you will experience miracles in your life.

We are engaged in the work of Almighty God. Jesus is the Christ. We are Their servants. I so testify in the name of Jesus Christ, amen.

\newpage
\settalk{“Come, Follow Me”}
\setspeaker{Russell M. Nelson}
\settalkdate{April 2019}
\talkheading{“Come, Follow Me”}{Russell M. Nelson}

My beloved brothers and sisters, my wife Wendy and I rejoice in being with you on this Sabbath morning. Much has happened since our last general conference. New temples have been dedicated in Concepción, Chile; Barranquilla, Colombia; and Rome, Italy. We experienced a rich outpouring of the Spirit at these sacred events.

I congratulate the many women (and men) who have recently read the Book of Mormon and discovered joy and hidden treasures. I am inspired by reports about miracles received.

I marvel at 11-year-old young men who, now as deacons, worthily pass the sacrament each Sunday. They go to the temple along with our 11-year-old young women, who are now eagerly learning and serving as Beehives. Both young men and young women are preaching gospel truths with clarity and conviction.

I rejoice with children and youth who are helping to teach the gospel in their homes as they work with their parents to follow the home-centered, Church-supported curriculum.

We received this photo of four-year-old Blake, who, early on a Saturday morning, grabbed a Church book and exclaimed, “I need to feed my spirit!”

Blake, we are thrilled with you and others who are choosing to feed their spirits by feasting on the truths of the restored gospel of Jesus Christ. And we delight in knowing that many are receiving the power of God in their lives as they worship and serve in the temple.

As many of you know, our family experienced a tender separation three months ago when our daughter Wendy departed from this mortal life. In the final days of her battle with cancer, I was blessed with the opportunity to have our farewell daddy-daughter conversation.

I held her hands and told her how much I loved her and how grateful I was to be her father. I said: “You married in the temple and faithfully honored your covenants. You and your husband welcomed seven children into your home and raised them to be devout disciples of Jesus Christ, valiant Church members, and contributing citizens. And they have chosen spouses of that same caliber. Your daddy is very, very proud of you. You have brought me much joy!”

She quietly responded, “Thank you, Daddy.”

It was a tender, tearful moment for us. During her 67 years, we worked together, sang together, and often skied together. But that evening, we talked of things that matter most, such as covenants, ordinances, obedience, faith, family, fidelity, love, and eternal life.

We miss our daughter greatly. However, because of the restored gospel of Jesus Christ, we do not worry about her. As we continue to honor our covenants with God, we live in anticipation of our being with her again. Meanwhile, we’re serving the Lord here and she is serving Him there—in paradise.

Actually, my wife and I visited Paradise earlier this year—Paradise, California, that is. As it happened, our scheduled visit there came less than 40 hours after our daughter departed from this world. We, along with Elder Kevin W. Pearson and his wife, June, were bolstered by the Saints of the Chico California Stake. We learned of their great faith, their ministering, and the miracles that occurred even amidst their devastating losses from the most destructive wildfire in the history of California.

While there, we spoke at length with a young police officer, John, who was one of many brave first responders. He recalled the thick darkness that descended upon Paradise on November 8, 2018, as flames and embers raced through the town, devouring property and possessions like a scourge and leaving nothing but piles of ash and stark brick chimneys.

For 15 hours, John drove through an impenetrable darkness that was streaked with javelins of threatening embers as he helped person after person, family after family escape to safety—all at the peril of his own life. Yet during that strenuous ordeal, what terrified John most was his all-consuming question: “Where is my family?” After many long, terrifying hours of anguish, he finally learned of their safe evacuation.

The account of John’s concern for his family has prompted me to speak today with those of you who may ask when approaching the end of your mortal life, “Where is my family?” In that coming day when you will complete your mortal probation and enter the spirit world, you will be brought face-to-face with that heart-wrenching question: “Where is my family?”

Jesus Christ teaches the way back to our eternal home. He understands our Heavenly Father’s plan of eternal progression better than any of us. After all, He is the keystone of it all. He is our Redeemer, our Healer, and our Savior.

Ever since Adam and Eve were expelled from the Garden of Eden, Jesus the Christ has offered His mighty arm to help all who choose to follow Him. Repeatedly, scriptures record that despite all kinds of sins from all kinds of people, His arms are outstretched still.

The spirit in each of us naturally yearns for family love to last forever. Love songs perpetuate a false hope that love is all you need if you want to be together forever. And some erroneously believe that the Resurrection of Jesus Christ provides a promise that all people will be with their loved ones after death.

In truth, the Savior Himself has made it abundantly clear that while His Resurrection assures that every person who ever lived will indeed be resurrected and live forever, much more is required if we want to have the high privilege of exaltation. Salvation is an individual matter, but exaltation is a family matter.

Listen to these words spoken by the Lord Jesus Christ to His prophet: “All covenants, contracts, bonds, obligations, oaths, vows, performances, connections, associations, or expectations, that are not made and entered into and sealed by the Holy Spirit of promise … are of no efficacy, virtue, or force in and after the resurrection from the dead; for all contracts that are not made unto this end have an end when men are dead.”

So, what is required for a family to be exalted forever? We qualify for that privilege by making covenants with God, keeping those covenants, and receiving essential ordinances.

This has been true since the beginning of time. Adam and Eve, Noah and his wife, Abraham and Sarah, Lehi and Sariah, and all other devoted disciples of Jesus Christ—since the world was created—have made the same covenants with God. They have received the same ordinances that we as members of the Lord’s restored Church today have made: those covenants that we receive at baptism and in the temple.

The Savior invites all to follow Him into the waters of baptism and, in time, to make additional covenants with God in the temple and receive and be faithful to those further essential ordinances. All these are required if we want to be exalted with our families and with God forever.

The anguish of my heart is that many people whom I love, whom I admire, and whom I respect decline His invitation. They ignore the pleadings of Jesus Christ when He beckons, “Come, follow me.”

I understand why God weeps. I also weep for such friends and relatives. They are wonderful men and women, devoted to their family and civic responsibilities. They give generously of their time, energy, and resources. And the world is better for their efforts. But they have chosen not to make covenants with God. They have not received the ordinances that will exalt them with their families and bind them together forever.

How I wish I could visit with them and invite them to consider seriously the enabling laws of the Lord. I’ve wondered what I could possibly say so they would feel how much the Savior loves them and know how much I love them and come to recognize how covenant-keeping women and men can receive a “fulness of joy.”

They need to understand that while there is a place for them hereafter—with wonderful men and women who also chose not to make covenants with God—that is not the place where families will be reunited and be given the privilege to live and progress forever. That is not the kingdom where they will experience the fulness of joy—of never-ending progression and happiness. Those consummate blessings can come only by living in an exalted celestial realm with God, our Eternal Father; His Son, Jesus Christ; and our wonderful, worthy, and qualified family members.

I feel to say to my reticent friends:

“In this life, you have never settled for second best in anything. Yet, as you resist fully embracing the restored gospel of Jesus Christ, you are choosing to settle for second best.

“The Savior said, ‘In my Father’s house are many mansions.’ However, as you choose not to make covenants with God, you are settling for a most meager roof over your head throughout all eternity.”

I would further entreat my reticent friends by saying:

“Pour out your heart to God. Ask Him if these things are true. Make time to study His words. Really study! If you truly love your family and if you desire to be exalted with them throughout eternity, pay the price now—through serious study and fervent prayer—to know these eternal truths and then to abide by them.

“If you are not sure you even believe in God, start there. Understand that in the absence of experiences with God, one can doubt the existence of God. So, put yourself in a position to begin having experiences with Him. Humble yourself. Pray to have eyes to see God’s hand in your life and in the world around you. Ask Him to tell you if He is really there—if He knows you. Ask Him how He feels about you. And then listen.”

One such dear friend of mine had limited experiences with God. But he longed to be with his departed wife. So he asked me to help him. I encouraged him to meet with our missionaries in order to understand the doctrine of Christ and learn of gospel covenants, ordinances, and blessings.

That he did. But he felt the course they advised would require him to make too many changes in his life. He said, “Those commandments and covenants are just too difficult for me. Also, I can’t possibly pay tithing, and I don’t have time to serve in the Church.” Then he asked me, “Once I die, please do the necessary temple work for my wife and me so that we can be together again.”

Thankfully, I am not this man’s judge. But I do question the efficacy of proxy temple work for a man who had the opportunity to be baptized in this life—to be ordained to the priesthood and receive temple blessings while here in mortality—but who made the conscious decision to reject that course.

My dear brothers and sisters, Jesus Christ invites us to take the covenant path back home to our Heavenly Parents and be with those we love. He invites us to “come, follow me.”

Now, as President of His Church, I plead with you who have distanced yourselves from the Church and with you who have not yet really sought to know that the Savior’s Church has been restored. Do the spiritual work to find out for yourselves, and please do it now. Time is running out.

I testify that God lives! Jesus is the Christ. His Church and the fulness of His gospel have been restored to bless our lives with joy, here and hereafter. I so testify in the name of Jesus Christ, amen.

\newpage
\settalk{Closing Remarks}
\setspeaker{Russell M. Nelson}
\settalkdate{April 2019}
\talkheading{Closing Remarks}{Russell M. Nelson}

My beloved brothers and sisters, as we come to the end of this historic conference, we thank the Lord for His inspiration and protection. The messages have instructed and edified us.

Topics were not assigned to the speakers. They each prayed for personal revelation in preparing their messages. To me, it is remarkable how those themes seem to fit so well with each other. As you study them, seek to learn what the Lord is trying to teach you through His servants.

The music has been glorious. We are deeply grateful to the many musicians who have blended their talents to bring the Spirit of the Lord into each session. And He has blessed the prayers and the congregations in every session. Truly, conference has again been a spiritual feast for us all.

We hope and pray that each member’s home will become a true sanctuary of faith, where the Spirit of the Lord may dwell. Despite contention all around us, one’s home can become a heavenly place, where study, prayer, and faith can be merged with love. We can truly become disciples of the Lord, standing up and speaking up for Him wherever we are.

God’s objective should be our objective. He wants His children to choose to return to Him, prepared, qualified, endowed, sealed, and faithful to covenants made in holy temples.

We now have 162 dedicated temples. The earliest ones stand as monuments to the faith and vision of our beloved pioneers. Each temple constructed by them resulted from their great personal sacrifice and effort. Each one stands as a stunning jewel in the crown of pioneer achievement.

Ours is a sacred responsibility to care for them. Therefore, these pioneer temples will soon undergo a period of renewal and refreshing and, for some, a major restoration. Efforts will be made to preserve the unique historicity of each temple wherever possible, preserving the inspiring beauty and unique craftsmanship of generations long-since passed.

Details for the St. George Utah Temple have already been released. Plans for the renovation of the Salt Lake Temple, Temple Square, and the adjoining plaza near the Church Office Building will be announced on Friday, April 19, 2019.

The Manti and Logan Utah Temples will also be renovated in coming years. When those plans are prepared, they too will be announced.

This work will require that each temple be closed for a period of time. Church members may continue to enjoy temple worship and service in other temples nearby. When each project is completed, each historic temple will be rededicated.

Brothers and sisters, we regard a temple as the most sacred structure in the Church. As we announce plans to construct a new temple, it becomes part of our sacred history. Now, please listen carefully and reverently. If I announce a temple in a place that is special to you, may I suggest that you simply bow your head with a silent prayer of gratitude in your heart. We do not want any verbal outbursts to detract from the sacred nature of this conference and the Lord’s holy temples.

Today we are pleased to announce plans to construct more temples, to be built in the following locations:

Pago Pago, American Samoa; Okinawa City, Okinawa; Neiafu, Tonga; Tooele Valley, Utah; Moses Lake, Washington; San Pedro Sula, Honduras; Antofagasta, Chile; Budapest, Hungary.

Thank you, my dear brothers and sisters.

As we speak of our temples old and new, may each of us signify by our actions that we are true disciples of the Lord Jesus Christ. May we renovate our lives through our faith and trust in Him. May we access the power of His Atonement by our repentance each day. And may we dedicate and rededicate our lives to serving God and His children—on both sides of the veil.

I leave my love and blessing with you, assuring you that revelation continues in this, the Lord’s Church. It will continue until “the purposes of God shall be accomplished, and the Great Jehovah shall say the work is done.”

I so bless you and bear my testimony that God lives! Jesus is the Christ! This is His Church. We are His people. In the sacred name of Jesus Christ, amen.

\newpage
\settalk{Witnesses, Aaronic Priesthood Quorums, and Young Women Classes}
\setspeaker{Russell M. Nelson}
\settalkdate{October 2019}
\talkheading{Witnesses, Aaronic Priesthood Quorums, and Young Women Classes}{Russell M. Nelson}

Dear brothers and sisters, it is wonderful to be with you again in general conference. Earlier this week, announcements were made to members of the Church about changes in policy regarding who may serve as witnesses to baptism and sealing ordinances. I would like to highlight those three points.

These policy adjustments are procedural. The underlying doctrine and covenants are unchanged. They are equally efficacious in all ordinances. These changes should greatly enhance family participation in these ordinances.

I also wanted to speak with you at this time to introduce adjustments that pertain to our youth and their leaders.

You will remember that I have invited the youth of The Church of Jesus Christ of Latter-day Saints to enlist in the Lord’s youth battalion to participate in the greatest cause on earth today—the gathering of Israel. I issued this invitation to our youth because they are unusually gifted in reaching out to others and sharing what they believe in a convincing fashion. The cause of the gathering is an essential part of helping to prepare the world and its people for the Second Coming of the Lord.

In each ward, the Lord’s youth battalion is led by a bishop, a dedicated servant of God. His first and foremost responsibility is to care for the young men and young women of his ward. The bishop and his counselors direct the work of the Aaronic Priesthood quorums and the Young Women classes in the ward.

The adjustments we will now announce are intended to help young men and young women develop their sacred personal potential. We also want to strengthen Aaronic Priesthood quorums and Young Women classes and provide support to bishops and other adult leaders as they serve this rising generation.

Elder Quentin L. Cook will now discuss the adjustments that relate to the young men. And tonight, at the general women’s session, Sister Bonnie H. Cordon, Young Women General President, will discuss the adjustments that relate to the young women.

The First Presidency and the Twelve are united in endorsing these efforts to strengthen our youth. Oh, how we love them and pray for them! They are the “hope of Israel, Zion’s army, children of the promised day.” We express our complete confidence in our youth and our gratitude for them. In the sacred name of Jesus Christ, amen.

\newpage
\settalk{Spiritual Treasures}
\setspeaker{Russell M. Nelson}
\settalkdate{October 2019}
\talkheading{Spiritual Treasures}{Russell M. Nelson}

Thank you for that beautiful music. As we all stood to sing that intermediate hymn, “We Thank Thee, O God, for a Prophet,” I had two overpowering thoughts come to me. One is about the Prophet Joseph Smith, the prophet of this dispensation. My love and admiration for him grows with every passing day. The second thought occurred as I looked at my wife, my daughters, granddaughters, and great-granddaughters. I felt like I’d like to claim every one of you as part of my family.

Several months ago, at the end of a temple endowment session, I said to my wife Wendy, “I hope the sisters understand the spiritual treasures that are theirs in the temple.” Sisters, I often find myself thinking about you, including two months ago when Wendy and I visited Harmony, Pennsylvania.

This was our second trip there. Both times we have been deeply moved as we walked on that sacred ground. It was near Harmony that John the Baptist appeared to Joseph Smith and restored the Aaronic Priesthood.

It was there that the Apostles Peter, James, and John appeared to restore the Melchizedek Priesthood.

It was in Harmony that Emma Hale Smith served as her husband’s first scribe while the Prophet translated the Book of Mormon.

It was also in Harmony that Joseph received a revelation manifesting the Lord’s will to Emma. The Lord instructed Emma to expound the scriptures, to exhort the Church, to receive the Holy Ghost, and to spend her time “learning much.” Emma was also counseled to “lay aside the things of this world and seek for the things of a better” and to hold fast to her covenants with God. The Lord concluded His instruction with these compelling words: “This is my voice unto all.”

Everything that happened in this area has profound implications for your lives. The restoration of the priesthood, along with the Lord’s counsel to Emma, can guide and bless each of you. How I yearn for you to understand that the restoration of the priesthood is just as relevant to you as a woman as it is to any man. Because the Melchizedek Priesthood has been restored, both covenant-keeping women and men have access to “all the spiritual blessings of the church” or, we might say, to all the spiritual treasures the Lord has for His children.

Every woman and every man who makes covenants with God and keeps those covenants, and who participates worthily in priesthood ordinances, has direct access to the power of God. Those who are endowed in the house of the Lord receive a gift of God’s priesthood power by virtue of their covenant, along with a gift of knowledge to know how to draw upon that power.

The heavens are just as open to women who are endowed with God’s power flowing from their priesthood covenants as they are to men who bear the priesthood. I pray that truth will register upon each of your hearts because I believe it will change your life. Sisters, you have the right to draw liberally upon the Savior’s power to help your family and others you love.

Now, you might be saying to yourself, “This sounds wonderful, but how do I do it? How do I draw the Savior’s power into my life?”

You won’t find this process spelled out in any manual. The Holy Ghost will be your personal tutor as you seek to understand what the Lord would have you know and do. This process is neither quick nor easy, but it is spiritually invigorating. What could possibly be more exciting than to labor with the Spirit to understand God’s power—priesthood power?

What I can tell you is that accessing the power of God in your life requires the same things that the Lord instructed Emma and each of you to do.

So, I invite you to study prayerfully section 25 of the Doctrine and Covenants and discover what the Holy Ghost will teach you. Your personal spiritual endeavor will bring you joy as you gain, understand, and use the power with which you have been endowed.

Part of this endeavor will require you to put aside many things of this world. Sometimes we speak almost casually about walking away from the world with its contention, pervasive temptations, and false philosophies. But truly doing so requires you to examine your life meticulously and regularly. As you do so, the Holy Ghost will prompt you about what is no longer needful, what is no longer worthy of your time and energy.

As you shift your focus away from worldly distractions, some things that seem important to you now will recede in priority. You will need to say no to some things, even though they may seem harmless. As you embark upon and continue this lifelong process of consecrating your life to the Lord, the changes in your perspective, feelings, and spiritual strength will amaze you!

Now a little word of warning. There are those who would undermine your ability to call upon the power of God. There are some who would have you doubt yourself and minimize your stellar spiritual capacity as a righteous woman.

Most certainly, the adversary does not want you to understand the covenant you made at baptism or the profound endowment of knowledge and power you have received or will receive in the temple—the house of the Lord. And Satan certainly does not want you to understand that every time you worthily serve and worship in the temple, you leave armed with God’s power and with His angels having “charge over” you.

Satan and his minions will constantly contrive roadblocks to prevent you from understanding the spiritual gifts with which you have been and can be blessed. Unfortunately, some roadblocks may be the result of another’s misbehavior. It grieves me to think that any of you have felt marginalized or have not been believed by a priesthood leader or have been abused or betrayed by a husband, father, or a supposed friend. I feel deep sorrow that any of you have felt sidelined, disrespected, or misjudged. Such offenses have no place in the kingdom of God.

Conversely, it thrills me when I learn of priesthood leaders who eagerly seek the participation of women in ward and stake councils. I am inspired by each husband who demonstrates that his most important priesthood responsibility is to care for his wife. I praise that man who deeply respects his wife’s ability to receive revelation and treasures her as an equal partner in their marriage.

When a man understands the majesty and power of a righteous, seeking, endowed Latter-day Saint woman, is it any wonder that he feels like standing when she enters the room?

From the dawning of time, women have been blessed with a unique moral compass—the ability to distinguish right from wrong. This gift is enhanced in those who make and keep covenants. And it diminishes in those who willfully ignore the commandments of God.

I hasten to add that I do not absolve men in any way from God’s requirement for them also to distinguish between right and wrong. But my dear sisters, your ability to discern truth from error, to be society’s guardians of morality, is crucial in these latter days. And we depend upon you to teach others to do likewise. Let me be very clear about this: if the world loses the moral rectitude of its women, the world will never recover.

We Latter-day Saints are not of the world; we are of covenant Israel. We are called to prepare a people for the Second Coming of the Lord.

Now, may I clarify several additional points with respect to women and priesthood. When you are set apart to serve in a calling under the direction of one who holds priesthood keys—such as your bishop or stake president—you are given priesthood authority to function in that calling.

Similarly, in the holy temple you are authorized to perform and officiate in priesthood ordinances every time you attend. Your temple endowment prepares you to do so.

If you are endowed but not currently married to a man who bears the priesthood and someone says to you, “I’m sorry you don’t have the priesthood in your home,” please understand that that statement is incorrect. You may not have a priesthood bearer in your home, but you have received and made sacred covenants with God in His temple. From those covenants flows an endowment of His priesthood power upon you. And remember, if your husband should die, you would preside in your home.

As a righteous, endowed Latter-day Saint woman, you speak and teach with power and authority from God. Whether by exhortation or conversation, we need your voice teaching the doctrine of Christ. We need your input in family, ward, and stake councils. Your participation is essential and never ornamental!

My dear sisters, your power will increase as you serve others. Your prayers, fasting, time in the scriptures, service in the temple, and family history work will open the heavens to you.

I entreat you to study prayerfully all the truths you can find about priesthood power. You might begin with Doctrine and Covenants sections 84 and 107. Those sections will lead you to other passages. The scriptures and teachings by modern prophets, seers, and revelators are filled with these truths. As your understanding increases and as you exercise faith in the Lord and His priesthood power, your ability to draw upon this spiritual treasure that the Lord has made available will increase. As you do so, you will find yourselves better able to help create eternal families that are united, sealed in the temple of the Lord, and full of love for our Heavenly Father and for Jesus Christ.

All our efforts to minister to each other, proclaim the gospel, perfect the Saints, and redeem the dead converge in the holy temple. We now have 166 temples throughout the world, and more are coming.

As you know, the Salt Lake Temple, Temple Square, and the adjoining plaza near the Church Office Building will be renewed in a project that will begin at the close of this year. This sacred temple must be preserved and prepared to inspire future generations, just as it has influenced us in this generation.

As the Church grows, more temples will be built so that more families can have access to that greatest of all blessings, that of eternal life. We regard a temple as the most sacred structure in the Church. Whenever plans are announced to construct a new temple, it becomes an important part of our history. As we have discussed here tonight, you sisters are vital to the work of the temple, and the temple is where you will receive your highest spiritual treasures.

Please listen carefully and reverently as I will now announce plans to build eight new temples. If one is announced in a place that is meaningful to you, I suggest that you simply bow your head prayerfully with gratitude in your heart. We are pleased to announce plans to construct temples in the following locations: Freetown, Sierra Leone; Orem, Utah; Port Moresby, Papua New Guinea; Bentonville, Arkansas; Bacolod, Philippines; McAllen, Texas; Cobán, Guatemala; and Taylorsville, Utah. Thank you, dear sisters. We deeply appreciate your receipt of these plans and your reverent response.

Now, in closing, I would like to leave a blessing upon you, that you may understand the priesthood power with which you have been endowed and that you will augment that power by exercising your faith in the Lord and in His power.

Dear sisters, with deep respect and gratitude, I express my love for you. Humbly, I declare that God lives! Jesus is the Christ. This is His Church. I so testify in the sacred name of Jesus Christ, amen.

\newpage
\settalk{The Second Great Commandment}
\setspeaker{Russell M. Nelson}
\settalkdate{October 2019}
\talkheading{The Second Great Commandment}{Russell M. Nelson}

My dear brothers and sisters, thank you for all you are doing to help gather Israel on both sides of the veil, to strengthen your families, and to bless the lives of those in need. Thank you for living as true followers of Jesus Christ. You know and love to obey His two great commandments, to love God and to love your neighbors.

During the last six months, Sister Nelson and I have met thousands of Saints as we have traveled to Central and South America, the islands of the Pacific, and various cities in the United States. As we travel, our hope is to build your faith. Yet we always return having had our faith strengthened by the members and friends we meet. May I share three meaningful moments from our recent experiences?

In May, Sister Nelson and I traveled with Elder Gerrit W. and Sister Susan Gong to the South Pacific. While in Auckland, New Zealand, we had the honor of meeting with imams from two mosques in Christchurch, New Zealand, where just two months earlier, innocent worshippers had been gunned down in an act of horrible violence.

We extended our sympathy to these brothers of another faith and reaffirmed our mutual commitment to religious freedom.

We also offered volunteer labor and modest financial assistance to rebuild their mosques. Our meeting with these Muslim leaders was filled with tender expressions of brotherhood.

In August, along with Elder Quentin L. and Sister Mary Cook, Sister Nelson and I met individuals in Buenos Aires, Argentina—most of them not of our faith—whose lives have been changed by wheelchairs provided to them through our Latter-day Saint Charities. We were inspired as they expressed joy-filled gratitude for their newfound mobility.

A third precious moment occurred just a few weeks ago here in Salt Lake City. It came from a unique letter I received on my birthday from a young woman I will call Mary—age 14.

Mary wrote about things she and I had in common: “You have 10 kids. We have 10 kids. You speak Mandarin. Seven of the kids in my family, including me, were adopted from China, so Mandarin is our first language. You are a heart surgeon. My sister has had two open-heart [operations]. You like two-hour church. We like two-hour church. You have perfect pitch. My brother has perfect pitch too. He is blind like me.”

Mary’s words touched me deeply, revealing not only her great spirit but also the consecration of her mother and father.

Latter-day Saints, as with other followers of Jesus Christ, are always looking for ways to help, to lift, and to love others. They who are willing to be called the Lord’s people “are willing to bear one another’s burdens, … to mourn with those that mourn; … and [to] comfort those that stand in need of comfort.”

They truly seek to live the first and second great commandments. When we love God with all our hearts, He turns our hearts to the well-being of others in a beautiful, virtuous cycle.

It would be impossible to calculate the amount of service that Latter-day Saints render around the globe every day of every year, but it is possible to calculate the good the Church as an organization does to bless men and women—boys and girls—who are in need of a helping hand.

The Church’s humanitarian outreach was launched in 1984. Then a Churchwide fast was held to raise funds to assist those afflicted by a devastating drought in eastern Africa. Church members donated \$6.4 million on that single fast day.

Then-Elder M. Russell Ballard and Brother Glenn L. Pace were dispatched to Ethiopia to assess how those consecrated funds could best be used. This effort proved to be the beginning of what would later be known as Latter-day Saint Charities.

Since that time, Latter-day Saint Charities has provided more than two billion dollars in aid to assist those in need throughout the world. This assistance is offered to recipients regardless of their church affiliation, nationality, race, sexual orientation, gender, or political persuasion.

That is not all. To assist members of the Lord’s Church in distress, we love and live the ancient law of the fast. We go hungry to help others who are hungry. One day each month, we go without food and donate the cost of that food (and more) to help those in need.

I will never forget my first visit to West Africa in 1986. The Saints came to our meetings in great numbers. Though they had little in terms of material possessions, most came dressed in spotless white clothing.

I asked the stake president how he cared for members who had so little. He replied that their bishops knew their people well. If members could afford two meals a day, no help was needed. But if they could afford only one meal or less—even with family help—bishops provided food, financed from fast offerings. Then he added this remarkable fact: their fast-offering contributions usually exceeded their expenses. Surplus fast offerings were then sent to people elsewhere whose needs exceeded theirs. Those stalwart African Saints taught me a great lesson about the power of the law and the spirit of the fast.

As members of the Church, we feel a kinship to those who suffer in any way. As sons and daughters of God, we are all brothers and sisters. We heed an Old Testament admonition: “Thou shalt open thine hand wide unto thy brother, to thy poor, and to thy needy.”

We also strive to live the teachings of the Lord Jesus Christ as recorded in Matthew 25:

“For I was an hungred, and ye gave me meat: I was thirsty, and ye gave me drink: I was a stranger, and ye took me in:

“Naked, and ye clothed me: I was sick, and ye visited me. …

“… Inasmuch as ye have done it unto one of the least of these my brethren, ye have done it unto me.”

Let me cite just a few examples of how the Church follows these teachings of the Savior.

To help relieve hunger, the Church operates 124 bishops’ storehouses throughout the world. Through them, approximately 400,000 food orders are given each year to individuals in need. In locations where no storehouse exists, bishops and branch presidents draw from fast-offering funds of the Church to provide food and supplies for their needy members.

However, the challenge of hunger goes far beyond the boundaries of the Church. It is increasing throughout the world. A recent United Nations report indicated that the number of undernourished people in the world now exceeds 820 million—or almost one in nine of the earth’s inhabitants.

What a sobering statistic! How grateful we are for your contributions. Thanks to your heartfelt generosity, millions throughout the world will receive much-needed food, clothing, temporary shelter, wheelchairs, medicines, clean water, and more.

Much sickness throughout the world is caused because of unclean water. To date, the Church’s humanitarian initiative has helped provide clean water in hundreds of communities in 76 countries.

A project in Luputa, in the Democratic Republic of the Congo, is a great example. With a population exceeding 100,000, the town had no running water. Citizens had to walk long distances for sources of safe water. A mountain spring was discovered 18 miles (29 km) away, but townspeople could not access that water on a regular basis.

When our humanitarian missionaries learned about this challenge, they worked with the leaders of Luputa by supplying materials and training to pipe the water to the city. The people of Luputa spent three years digging a one-meter-deep trench through rock and jungle. By working together, the joyful day finally arrived when fresh, clean water was available to all in that village.

The Church also helps refugees, whether from civil strife, the ravages of nature, or religious persecution. More than 70 million people are now displaced from their homes.

In the year 2018 alone, the Church provided emergency supplies to refugees in 56 countries. In addition, many Church members volunteer their time to help refugees integrate into new communities. We thank every one of you who reach out to help those who are trying to establish new homes.

Through generous donations to Deseret Industries outlets in the United States, millions of pounds of clothing are collected and sorted each year. While local bishops use this vast inventory to help members in need, the greatest portion is donated to other charitable organizations who distribute the items worldwide.

And just last year, the Church provided vision care for more than 300,000 people in 35 countries, newborn care for thousands of mothers and infants in 39 countries, and wheelchairs for more than 50,000 people living in dozens of countries.

The Church is well known for being among the first responders when tragedy strikes. Even before a hurricane hits, Church leaders and staff in the affected locations are mapping out plans for how they will deliver relief supplies and volunteer assistance to those who will be impacted.

Last year alone, the Church carried out more than 100 disaster-relief projects around the world, helping victims of hurricanes, fires, floods, earthquakes, and other calamities. Whenever possible, our Church members in yellow Helping Hands vests mobilize in great numbers to help those afflicted by the disaster. This kind of service, rendered by so many of you, is the very essence of ministering.

My dear brothers and sisters, the activities I have described are merely a small part of the growing welfare and humanitarian outreach of The Church of Jesus Christ of Latter-day Saints. And you are the ones who make all this possible. Because of your exemplary lives, your generous hearts, and your helping hands, it is no wonder that many communities and government leaders are praising your efforts.

Since becoming President of the Church, I have been amazed at how many presidents, prime ministers, and ambassadors have sincerely thanked me for our humanitarian aid to their people. And they have also expressed gratitude for the strength that our faithful members bring to their country as loyal, contributing citizens.

I have also marveled as world leaders have visited the First Presidency expressing their hope for the Church to be established in their lands. Why? Because they know that Latter-day Saints will help to build strong families and communities, making life better for others wherever they live.

Regardless of where we call home, members of the Church feel passionately about the fatherhood of God and the brotherhood of man. Thus, our greatest joy comes as we help our brothers and sisters, no matter where we live in this wonderful world.

Giving help to others—making a conscientious effort to care about others as much as or more than we care about ourselves—is our joy. Especially, I might add, when it is not convenient and when it takes us out of our comfort zone. Living that second great commandment is the key to becoming a true disciple of Jesus Christ.

My dear brothers and sisters, you are living exemplars of the fruits that come from following the teachings of Jesus Christ. I thank you! I love you!

I know that God lives. Jesus is the Christ. His Church has been restored in these latter days to fulfill its divine purposes. I so testify in the name of Jesus Christ, amen.

\newpage
\settalk{Closing Remarks}
\setspeaker{Russell M. Nelson}
\settalkdate{October 2019}
\talkheading{Closing Remarks}{Russell M. Nelson}

My beloved brothers and sisters, as we come to the end of this historic conference, we thank the Lord for inspiring the messages and the music that have edified us. We have truly enjoyed a spiritual feast.

We know the restored gospel of Jesus Christ will bring hope and joy to people who will hear and heed His doctrine. We also know that each home can become a true sanctuary of faith, where peace, love, and the Spirit of the Lord may dwell.

Of course, the crowning jewel of the Restoration is the holy temple. Its sacred ordinances and covenants are pivotal to preparing a people who are ready to welcome the Savior at His Second Coming. Presently we have 166 dedicated temples, and more are coming.

An open house will be held prior to the dedication of each new and renovated temple. Many friends not of our faith will participate in tours of those temples and will learn something about temple blessings. And some of those visitors will be moved upon to know more. Some will sincerely ask how they might qualify for the blessings of the temple.

As members of the Church, we need to be prepared to answer their questions. We can explain that the blessings of the temple are available to any and all people who will prepare themselves. But before they can enter a dedicated temple, they need to qualify. The Lord wants all His children to partake of the eternal blessings available in His temple. He has directed what each person must do to qualify to enter His holy house.

A good place for us to begin such a teaching opportunity is to call attention to the words etched on the temple’s exterior: “Holiness to the Lord: The House of the Lord.” President Henry B. Eyring’s message today and many others have inspired us to become more holy. Each temple is a holy place; each temple patron strives to become more holy.

All requirements to enter the temple relate to personal holiness. To assess that readiness, each person who wants to enjoy the blessings of the temple will have two interviews: first with a bishop, bishopric counselor, or branch president; second with a stake or mission president or one of his counselors. In those interviews, several questions will be asked.

Some of those questions have recently been edited for clarity. I would like to review them for you now:

Tomorrow, these revised temple recommend questions will be distributed to Church leaders throughout the world.

In addition to their answering those questions honestly, it is understood that each adult temple patron will wear the sacred garment of the priesthood under their regular clothing. This is symbolic of an inner commitment to strive each day to become more like the Lord. It also reminds us to remain faithful each day to covenants made and to walk on the covenant path each day in a higher and holier way.

Now, for just a moment, I would like to speak to our youth. We encourage you to qualify for limited-use temple recommends. You will be asked only those questions applicable to you in your preparation for the ordinances of proxy baptism and confirmation. We are very grateful for your worthiness and willingness to participate in that sacred temple work. We thank you!

Individual worthiness to enter the Lord’s house requires much individual spiritual preparation. But with the Lord’s help, nothing is impossible. In some respects, it is easier to build a temple than it is to build a people prepared for a temple. Individual worthiness requires a total conversion of mind and heart to be more like the Lord, to be an honest citizen, to be a better example, and to be a holier person.

I testify that such preparatory work brings innumerable blessings in this life and inconceivable blessings for the life to come, including the perpetuation of your family unit throughout all eternity “in a state of never-ending happiness.”

Now I would like to turn to another topic: plans for the coming year. In the springtime of the year 2020, it will be exactly 200 years since Joseph Smith experienced the theophany that we know as the First Vision. God the Father and His Beloved Son, Jesus Christ, appeared to Joseph, a 14-year-old youth. That event marked the onset of the Restoration of the gospel of Jesus Christ in its fulness, precisely as foretold in the Holy Bible.

Then came a succession of visits from heavenly messengers, including Moroni, John the Baptist, and the early Apostles Peter, James, and John. Others followed, including Moses, Elias, and Elijah. Each brought divine authority to bless God’s children on the earth once again.

Miraculously, we have also received the Book of Mormon: Another Testament of Jesus Christ, a companion scripture to the Holy Bible. The revelations published in the Doctrine and Covenants and the Pearl of Great Price have also greatly enriched our understanding of God’s commandments and eternal truth.

The keys and offices of the priesthood have been restored, including the offices of Apostle, Seventy, patriarch, high priest, elder, bishop, priest, teacher, and deacon. And women who love the Lord serve valiantly in the Relief Society, Primary, Young Women, Sunday School, and other Church callings—all vital parts of the Restoration of the gospel of Jesus Christ in its fulness.

Thus, the year 2020 will be designated as a bicentennial year. General conference next April will be different from any previous conference. In the next six months, I hope that every member and every family will prepare for a unique conference that will commemorate the very foundations of the restored gospel.

You may wish to begin your preparation by reading afresh Joseph Smith’s account of the First Vision as recorded in the Pearl of Great Price. Our course of study for next year in Come, Follow Me is the Book of Mormon. You may wish to ponder important questions such as, “How would my life be different if my knowledge gained from the Book of Mormon were suddenly taken away?” or “How have the events that followed the First Vision made a difference for me and my loved ones?” Also, with the Book of Mormon videos now becoming available, you may wish to incorporate them in your individual and family study.

Select your own questions. Design your own plan. Immerse yourself in the glorious light of the Restoration. As you do, general conference next April will be not only memorable; it will be unforgettable.

Now in closing, I leave with you my love and my blessing that each of you may become happier and holier with each passing day. Meanwhile, please be assured that revelation continues in the Church and will continue under the Lord’s direction until “the purposes of God shall be accomplished, and the Great Jehovah shall say the work is done.”

I so bless you, reaffirming my love for you, with my testimony that God lives! Jesus is the Christ! This is His Church and we are His people. In the sacred name of Jesus Christ, amen.

\newpage
\settalk{Opening Message}
\setspeaker{Russell M. Nelson}
\settalkdate{April 2020}
\talkheading{Opening Message}{Russell M. Nelson}

My beloved brothers and sisters, as we welcome you to this historic April 2020 general conference of The Church of Jesus Christ of Latter-day Saints, for reasons you know, I stand before you in an empty auditorium!

Little did I know, when I promised you at the October 2019 general conference that this April conference would be “memorable” and “unforgettable,” that speaking to a visible congregation of fewer than 10 people would make this conference so memorable and unforgettable for me! Yet the knowledge that you are participating by electronic transmission, and the choir’s beautiful rendition of “It Is Well with My Soul,” bring great comfort to my soul.

As you know, attendance at this general conference has been strictly limited as part of our efforts to be good global citizens and do all we can to limit the spread of COVID-19. This virus has had a major impact throughout the world. It has also altered our Church meetings, missionary service, and temple work for a while.

Though today’s restrictions relate to a virulent virus, life’s personal trials stretch far beyond this pandemic. Future trials could result from an accident, a natural disaster, or an unexpected personal heartache.

How can we endure such trials? The Lord has told us that “if ye are prepared ye shall not fear.” Of course, we can store our own reserves of food, water, and savings. But equally crucial is our need to fill our personal spiritual storehouses with faith, truth, and testimony.

Our ultimate quest in life is to prepare to meet our Maker. We do this by striving daily to become more like our Savior, Jesus Christ. And we do that as we repent daily and receive His cleansing, healing, and strengthening power. Then we can feel enduring peace and joy, even during turbulent times. This is exactly why the Lord has implored us to stand in holy places and “be not moved.”

This year, we commemorate the 200th anniversary of one of the most significant events in the history of the world—namely, the appearance of God the Father and His Beloved Son, Jesus Christ, to Joseph Smith. During that singular vision, God the Father pointed to Jesus Christ and said: “This is My Beloved Son. Hear Him!”

That admonition given to Joseph is for each of us. We are to seek, in every way we can, to hear Jesus Christ, who speaks to us through the power and ministering of the Holy Ghost.

The purpose of this and every general conference is to help us to hear Him. We have prayed, and invite you to pray, that the Spirit of the Lord will be with us in such rich abundance that you can hear the messages that the Savior has especially for you—messages that will bring peace to your soul. Messages that will heal your broken heart. Messages that will illuminate your mind. Messages that will help you know what to do as you move ahead through times of turmoil and trial.

We pray that this conference will be memorable and unforgettable because of the messages you will hear, the unique announcements which will be made, and the experiences in which you will be invited to participate.

For example, at the conclusion of the Sunday morning session, we will convene a worldwide solemn assembly when I will lead you in the sacred Hosanna Shout. We pray that this will be a spiritual highlight for you as we express in global unison our profound gratitude to God the Father and His Beloved Son by praising Them in this unique way.

For this sacred experience, we use clean white handkerchiefs. But if you do not have one, you may simply wave your hand. At the conclusion of the Hosanna Shout, the congregation will join with the choir in singing “The Spirit of God.”

My dear brothers and sisters, this conference will be magnificent. This year will be extraordinary as we focus intently on the Savior and His restored gospel. The most important lasting effects of this historic conference will be as our hearts change and we commence a lifelong quest to hear Him.

Welcome to April 2020 general conference! I know that God, our Heavenly Father, and His Son, Jesus Christ, are mindful of us. They will be with us throughout the proceedings of these two glorious days as we seek to draw closer to Them and honor Them. In the sacred name of Jesus Christ, amen.

\newpage
\settalk{Opening the Heavens for Help}
\setspeaker{Russell M. Nelson}
\settalkdate{April 2020}
\talkheading{Opening the Heavens for Help}{Russell M. Nelson}

What a unique and wonderful session this has been! Thank you, dear Laudy and Enzo. You represented the magnificent young women and young men of the Church so very well.

My dear brothers and sisters, we have heard much today about the Restoration of the Church—the very Church that our Savior, Jesus Christ, established during His earthly ministry. That Restoration began 200 years ago this spring when God the Father and His Son, Jesus Christ, appeared to the young Joseph Smith.

Ten years after this transcendent vision, the Prophet Joseph Smith and five others were called as founding members of the Lord’s restored Church.

From that small group assembled on April 6, 1830, has come a global organization of more than 16 million members. The good this Church accomplishes around the world to alleviate human suffering and provide uplift for humankind is widely known. But its prime purpose is to help men, women, and children follow the Lord Jesus Christ, keep His commandments, and qualify for the greatest of all blessings—that of eternal life with God and their loved ones.

As we commemorate the event that was launched in 1820, it is important to remember that while we revere Joseph Smith as a prophet of God, this is not the church of Joseph Smith, nor is it the church of Mormon. This is the Church of Jesus Christ. He decreed exactly what His Church should be called: “For thus shall my church be called in the last days, even The Church of Jesus Christ of Latter-day Saints.”

I have spoken previously about a needed course correction in the way we refer to the name of the Church. Since that time, much has been done to accomplish this correction. I am very grateful to President M. Russell Ballard and the entire Quorum of the Twelve Apostles, who have done so much to lead these efforts as well as those related to another initiative that I will announce this evening.

Church leaders and departments, related entities, and millions of members—and others—now use the correct name of the Church. The Church’s official style guide has been adjusted. The Church’s principal website is now ChurchofJesusChrist.org. Addresses for email, domain names, and social media channels have been updated. Our beloved choir is now the Tabernacle Choir at Temple Square.

We have gone to these extraordinary efforts because when we remove the Lord’s name from the name of His Church, we inadvertently remove Him as the central focus of our worship and our lives. When we take the Savior’s name upon us at baptism, we commit to witness, by our words, thoughts, and actions, that Jesus is the Christ.

Previously, I promised that if we would “do our best to restore the correct name of the Lord’s Church,” He would “pour down His power and blessings upon the heads of the Latter-day Saints, the likes of which we have never seen.” I renew that promise today.

To help us remember Him and to identify The Church of Jesus Christ of Latter-day Saints as the Lord’s Church, we are pleased to introduce a symbol that will signify the central place of Jesus Christ in His Church.

This symbol includes the name of the Church contained within a cornerstone. Jesus Christ is the chief cornerstone.

At the center of the symbol is a representation of Thorvaldsen’s marble statue the Christus. It portrays the resurrected, living Lord reaching out to embrace all who will come unto Him.

Symbolically, Jesus Christ is standing under an arch. The arch reminds us of the resurrected Savior emerging from the tomb on the third day following His Crucifixion.

This symbol should feel familiar to many, as we have long identified the restored gospel with the living, resurrected Christ.

The symbol will now be used as a visual identifier for official literature, news, and events of the Church. It will remind all that this is the Savior’s Church and that all we do as members of His Church centers on Jesus Christ and His gospel.

Now, my dear brothers and sisters, tomorrow is Palm Sunday, as Elder Gong has so eloquently taught. Then we enter the special week that culminates with Easter. As followers of Jesus Christ, living in a day when the COVID-19 pandemic has put the whole world in commotion, let us not just talk of Christ or preach of Christ or employ a symbol representing Christ.

Let us put our faith in the Lord Jesus Christ into action!

As you know, members of the Church observe the law of the fast one day each month.

The doctrine of fasting is ancient. It has been practiced by biblical heroes from the earliest days. Moses, David, Ezra, Nehemiah, Esther, Isaiah, Daniel, Joel, and many others fasted and preached of fasting. Through Isaiah’s writings, the Lord said: “Is not this the fast that I have chosen? to loose the bands of wickedness, to undo the heavy burdens, and to let the oppressed go free?”

The Apostle Paul admonished Saints in Corinth to “give yourselves to fasting and prayer.” The Savior Himself declared that certain things go “not out but by prayer and fasting.”

I said recently in a social media video that “as a physician and surgeon, I have great admiration for medical professionals, scientists, and all who are working around the clock to curb the spread of COVID-19.”

Now, as President of The Church of Jesus Christ of Latter-day Saints and an Apostle of Jesus Christ, I know that God “has all power, all wisdom, and all understanding; he comprehendeth all things, and he is a merciful Being, even unto salvation, to those who will repent and believe on his name.”

So, during times of deep distress, as when illness reaches pandemic proportions, the most natural thing for us to do is to call upon our Heavenly Father and His Son—the Master Healer—to show forth Their marvelous power to bless the people of the earth.

In my video message, I invited all to join in fasting on Sunday, March 29, 2020. Many of you may have seen the video and joined in the fast. Some may have not. Now we still need help from heaven.

So tonight, my dear brothers and sisters, in the spirit of the sons of Mosiah, who gave themselves to much fasting and prayer, and as part of our April 2020 general conference, I am calling for another worldwide fast. For all whose health may permit, let us fast, pray, and unite our faith once again. Let us prayerfully plead for relief from this global pandemic.

I invite all, including those not of our faith, to fast and pray on Good Friday, April 10, that the present pandemic may be controlled, caregivers protected, the economy strengthened, and life normalized.

How do we fast? Two meals or a period of 24 hours is customary. But you decide what would constitute a sacrifice for you, as you remember the supreme sacrifice the Savior made for you. Let us unite in pleading for healing throughout the world.

Good Friday would be the perfect day to have our Heavenly Father and His Son hear us!

Dear brothers and sisters, I express my deep love for you, along with my testimony of the divinity of the work in which we are engaged. This is The Church of Jesus Christ of Latter-day Saints. He stands at its head and directs all that we do. I know that He will respond to the pleadings of His people. I so testify in the sacred name of Jesus Christ, amen.

\newpage
\settalk{Hear Him}
\setspeaker{Russell M. Nelson}
\settalkdate{April 2020}
\talkheading{Hear Him}{Russell M. Nelson}

My dear brothers and sisters, how thankful I am that through the use of technology we have been able to meet together and worship on this Sunday morning. How blessed we are to know that the gospel of Jesus Christ has been restored to the earth!

In the past several weeks, most of us have experienced disruptions in our personal lives. Earthquakes, fires, floods, plagues, and their aftermaths have disrupted routines and caused shortages of food, staples, and savings.

Amidst all of this, we commend you and thank you for choosing to hear the word of the Lord during this time of turmoil by joining with us for general conference. The increasing darkness that accompanies tribulation makes the light of Jesus Christ shine ever brighter. Just think of the good each of us can do during this time of global upheaval. Your love of and faith in the Savior may very well be the catalyst for someone to discover the Restoration of the fulness of the gospel of Jesus Christ.

In the past two years, Sister Nelson and I have met with thousands of you around the world. We have convened with you in outdoor arenas and in hotel ballrooms. In each location, I have felt that I was in the presence of the Lord’s elect and that I was seeing the gathering of Israel occur before my eyes.

We live in the day that “our forefathers have awaited with anxious expectation.” We have front-row seats to witness live what the prophet Nephi saw only in vision, that “the power of the Lamb of God” would descend “upon the covenant people of the Lord, who were scattered upon all the face of the earth; and they were armed with righteousness and with the power of God in great glory.”

You, my brothers and sisters, are among those men, women, and children whom Nephi saw. Think of that!

Regardless of where you live or what your circumstances are, the Lord Jesus Christ is your Savior, and God’s prophet Joseph Smith is your prophet. He was foreordained before the foundation of the earth to be the prophet of this last dispensation, when “nothing shall be withheld” from the Saints. Revelation continues to flow from the Lord during this ongoing process of restoration.

What does it mean for you that the gospel of Jesus Christ has been restored to the earth?

It means that you and your family can be sealed together forever! It means that because you have been baptized by one who has authority from Jesus Christ and have been confirmed a member of His Church, you can enjoy the constant companionship of the Holy Ghost. He will guide and protect you. It means you will never be left comfortless or without access to the power of God to help you. It means that priesthood power can bless you as you receive essential ordinances and make covenants with God and keep them. What an anchor to our souls are these truths, especially during these times when the tempest is raging.

The Book of Mormon chronicles the classic rise and fall of two major civilizations. Their history demonstrates how easy it is for a majority of the people to forget God, reject warnings of the Lord’s prophets, and seek power, popularity, and pleasures of the flesh. Repeatedly, past prophets have declared “great and marvelous things unto the people, which they did not believe.”

It is no different in our day. Through the years, great and marvelous things have been heard from dedicated pulpits across the earth. Yet most people do not embrace these truths—either because they do not know where to look for them or because they are listening to those who do not have the whole truth or because they have rejected truth in favor of worldly pursuits.

The adversary is clever. For millennia he has been making good look evil and evil look good. His messages tend to be loud, bold, and boastful.

However, messages from our Heavenly Father are strikingly different. He communicates simply, quietly, and with such stunning plainness that we cannot misunderstand Him.

For example, whenever He has introduced His Only Begotten Son to mortals upon the earth, He has done so with remarkably few words. On the Mount of Transfiguration to Peter, James, and John, God said, “This is my beloved Son: hear him.” His words to the Nephites in ancient Bountiful were “Behold my Beloved Son, in whom I am well pleased, in whom I have glorified my name—hear ye him.” And to Joseph Smith, in that profound declaration that opened this dispensation, God simply said, “This is My Beloved Son. Hear Him!”

Now, my dear brothers and sisters, consider the fact that in these three instances just mentioned, just before the Father introduced the Son, the people involved were in a state of fear and, to some degree, desperation.

The Apostles were afraid when they saw Jesus Christ encircled by a cloud on the Mount of Transfiguration.

The Nephites were afraid because they had been through destruction and darkness for several days.

Joseph Smith was in the grips of a force of darkness just before the heavens opened.

Our Father knows that when we are surrounded by uncertainty and fear, what will help us the very most is to hear His Son.

Because when we seek to hear—truly hear—His Son, we will be guided to know what to do in any circumstance.

The very first word in the Doctrine and Covenants is hearken. It means “to listen with the intent to obey.” To hearken means to “hear Him”—to hear what the Savior says and then to heed His counsel. In those two words—“Hear Him”—God gives us the pattern for success, happiness, and joy in this life. We are to hear the words of the Lord, hearken to them, and heed what He has told us!

As we seek to be disciples of Jesus Christ, our efforts to hear Him need to be ever more intentional. It takes conscious and consistent effort to fill our daily lives with His words, His teachings, His truths.

We simply cannot rely upon information we bump into on social media. With billions of words online and in a marketing-saturated world constantly infiltrated by noisy, nefarious efforts of the adversary, where can we go to hear Him?

We can go to the scriptures. They teach us about Jesus Christ and His gospel, the magnitude of His Atonement, and our Father’s great plan of happiness and redemption. Daily immersion in the word of God is crucial for spiritual survival, especially in these days of increasing upheaval. As we feast on the words of Christ daily, the words of Christ will tell us how to respond to difficulties we never thought we would face.

We can also hear Him in the temple. The house of the Lord is a house of learning. There the Lord teaches in His own way. There each ordinance teaches about the Savior. There we learn how to part the veil and communicate more clearly with heaven. There we learn how to rebuke the adversary and draw upon the Lord’s priesthood power to strengthen us and those we love. How eager each of us should be to seek refuge there.

When these temporary COVID-19 restrictions are lifted, please schedule regular time to worship and serve in the temple. Every minute of that time will bless you and your family in ways nothing else can. Take time to ponder what you hear and feel when you are there. Ask the Lord to teach you how to open the heavens to bless your life and the lives of those you love and serve.

While worshipping in the temple is presently not possible, I invite you to increase your participation in family history, including family history research and indexing. I promise that as you increase your time in temple and family history work, you will increase and improve your ability to hear Him.

We also hear Him more clearly as we refine our ability to recognize the whisperings of the Holy Ghost. It has never been more imperative to know how the Spirit speaks to you than right now. In the Godhead, the Holy Ghost is the messenger. He will bring thoughts to your mind which the Father and Son want you to receive. He is the Comforter. He will bring a feeling of peace to your heart. He testifies of truth and will confirm what is true as you hear and read the word of the Lord.

I renew my plea for you to do whatever it takes to increase your spiritual capacity to receive personal revelation.

Doing so will help you know how to move ahead with your life, what to do during times of crisis, and how to discern and avoid the temptations and the deceptions of the adversary.

And, finally, we hear Him as we heed the words of prophets, seers, and revelators. Ordained Apostles of Jesus Christ always testify of Him. They point the way as we make our way through the heart-wrenching maze of our mortal experiences.

What will happen as you more intentionally hear, hearken, and heed what the Savior has said and what He is saying now through His prophets? I promise that you will be blessed with additional power to deal with temptation, struggles, and weakness. I promise miracles in your marriage, family relationships, and daily work. And I promise that your capacity to feel joy will increase even if turbulence increases in your life.

This April 2020 general conference is our time to commemorate an event that changed the world. As we anticipated this 200th anniversary of Joseph Smith’s First Vision, the First Presidency and Council of the Twelve Apostles wondered what we might do to commemorate appropriately this singular event.

That theophany initiated the Restoration of the fulness of the gospel of Jesus Christ and ushered in the dispensation of the fulness of times.

We wondered if a monument should be erected. But as we considered the unique historic and international impact of that First Vision, we felt impressed to create a monument not of granite or stone but of words—words of solemn and sacred proclamation—written, not to be carved in “tables of stone” but rather to be etched in the “fleshy tables” of our hearts.

Since the Church was organized, only five proclamations have been issued, with the last being “The Family: A Proclamation to the World,” presented by President Gordon B. Hinckley in 1995.

Now as we contemplate this significant time in the history of the world and the Lord’s charge to gather scattered Israel in preparation for the Second Coming of Jesus Christ, we, the First Presidency and Council of the Twelve Apostles, issue the following proclamation. Its title is “The Restoration of the Fulness of the Gospel of Jesus Christ: A Bicentennial Proclamation to the World.” It is authored by the First Presidency and the Council of the Twelve Apostles of The Church of Jesus Christ of Latter-day Saints. It is dated April 2020. To prepare for today, I previously recorded this proclamation in the Sacred Grove, where Joseph Smith first saw the Father and the Son.

“We solemnly proclaim that God loves His children in every nation of the world. God the Father has given us the divine birth, the incomparable life, and the infinite atoning sacrifice of His Beloved Son, Jesus Christ. By the power of the Father, Jesus rose again and gained the victory over death. He is our Savior, our Exemplar, and our Redeemer.

“Two hundred years ago, on a beautiful spring morning in 1820, young Joseph Smith, seeking to know which church to join, went into the woods to pray near his home in upstate New York, USA. He had questions regarding the salvation of his soul and trusted that God would direct him.

“In humility, we declare that in answer to his prayer, God the Father and His Son, Jesus Christ, appeared to Joseph and inaugurated the ‘restitution of all things’ (Acts 3:21) as foretold in the Bible. In this vision, he learned that following the death of the original Apostles, Christ’s New Testament Church was lost from the earth. Joseph would be instrumental in its return.

“We affirm that under the direction of the Father and the Son, heavenly messengers came to instruct Joseph and reestablish the Church of Jesus Christ. The resurrected John the Baptist restored the authority to baptize by immersion for the remission of sins. Three of the original twelve Apostles—Peter, James, and John—restored the apostleship and keys of priesthood authority. Others came as well, including Elijah, who restored the authority to join families together forever in eternal relationships that transcend death.

“We further witness that Joseph Smith was given the gift and power of God to translate an ancient record: the Book of Mormon—Another Testament of Jesus Christ. Pages of this sacred text include an account of the personal ministry of Jesus Christ among people in the Western Hemisphere soon after His Resurrection. It teaches of life’s purpose and explains the doctrine of Christ, which is central to that purpose. As a companion scripture to the Bible, the Book of Mormon testifies that all human beings are sons and daughters of a loving Father in Heaven, that He has a divine plan for our lives, and that His Son, Jesus Christ, speaks today as well as in days of old.

“We declare that The Church of Jesus Christ of Latter-day Saints, organized on April 6, 1830, is Christ’s New Testament Church restored. This Church is anchored in the perfect life of its chief cornerstone, Jesus Christ, and in His infinite Atonement and literal Resurrection. Jesus Christ has once again called Apostles and has given them priesthood authority. He invites all of us to come unto Him and His Church, to receive the Holy Ghost, the ordinances of salvation, and to gain enduring joy.

“Two hundred years have now elapsed since this Restoration was initiated by God the Father and His Beloved Son, Jesus Christ. Millions throughout the world have embraced a knowledge of these prophesied events.

“We gladly declare that the promised Restoration goes forward through continuing revelation. The earth will never again be the same, as God will ‘gather together in one all things in Christ’ (Ephesians 1:10).

“With reverence and gratitude, we as His Apostles invite all to know—as we do—that the heavens are open. We affirm that God is making known His will for His beloved sons and daughters. We testify that those who prayerfully study the message of the Restoration and act in faith will be blessed to gain their own witness of its divinity and of its purpose to prepare the world for the promised Second Coming of our Lord and Savior, Jesus Christ.”

Beloved brothers and sisters, that is our bicentennial proclamation to the world regarding the Restoration of the gospel of Jesus Christ in its fulness. It has been translated into 12 languages. Other languages will soon follow. It will be available immediately on the Church website, from which you may obtain a copy. Study it privately and with your family members and friends. Ponder the truths and think of the impact those truths will have on your life if you will hear them, hearken to them, and heed the commandments and covenants that accompany them.

I know that Joseph Smith is the foreordained prophet whom the Lord chose to open this last dispensation. Through him the Lord’s Church was restored to the earth. Joseph sealed his testimony with his blood. How I love and honor him!

God lives! Jesus is the Christ! His Church has been restored! He and His Father, our Heavenly Father, are watching over us. I so testify in the sacred name of Jesus Christ, amen.

\newpage
\settalk{Hosanna Shout}
\setspeaker{Presented by President Russell M. Nelson}
\settalkdate{April 2020}
\talkheading{Hosanna Shout}{Presented by President Russell M. Nelson}

Now, my dear brothers and sisters, as we commemorate Joseph Smith’s First Vision of the Father and the Son, we felt that it would be appropriate to rejoice together by participating in the Hosanna Shout.

This sacred shout was first given in this dispensation at the dedication of the Kirtland Temple on March 27, 1836. It is now given at the dedication of each temple. It is a sacred tribute to the Father and the Son, symbolizing the reaction of the multitude when the Savior made His triumphal entry into Jerusalem. It also reaffirms what young Joseph experienced that day in the Sacred Grove—namely, that the Father and the Son are two glorified Beings, whom we worship and praise.

I will now demonstrate how the Hosanna Shout is given. As I do, I invite our colleagues in the media to treat this very sacred observance with dignity and respect.

Each one participating takes a clean white handkerchief, holding it by one corner, and waves it while saying in unison, “Hosanna, Hosanna, Hosanna to God and the Lamb,” repeated three times, followed by “Amen, Amen, and Amen.” If you do not have a white handkerchief, you may simply wave your hand.

Brothers and sisters, I now invite you to stand and participate in the Hosanna Shout, following which the “Hosanna Anthem” and “The Spirit of God” will be sung.

Upon a signal from the conductor, please join in singing “The Spirit of God.”

\begin{verse}
Hosanna, Hosanna, Hosanna to God and the Lamb.\\
Hosanna, Hosanna, Hosanna to God and the Lamb.\\
Hosanna, Hosanna, Hosanna to God and the Lamb.\\
Amen, Amen, and Amen.
\end{verse}

\newpage
\settalk{Go Forward in Faith}
\setspeaker{Russell M. Nelson}
\settalkdate{April 2020}
\talkheading{Go Forward in Faith}{Russell M. Nelson}

My beloved brothers and sisters, as we come to the end of this historic conference, we express our gratitude to the Lord. The music has been sublime and the messages inspiring.

During this conference, we have experienced many highlights. On this bicentennial anniversary, we have introduced a proclamation to the world declaring the reality of the Restoration of the gospel of Jesus Christ in its fulness.

We commemorated the Restoration with the Hosanna Shout.

We unveiled a new symbol signifying our faith in the Lord Jesus Christ and for visual recognition of official Church information and materials.

We have called for a global day of fasting and prayer, that the present pandemic may be controlled, caregivers protected, the economy strengthened, and life normalized. This fast will be held on Good Friday, April 10. What a great Friday that will be!

Next Sunday is Easter Sunday, when we will again commemorate the Atonement and Resurrection of our Lord Jesus Christ. Because of His Atonement, His gift of resurrection will come to all who have ever lived. And His gift of eternal life will come to all who qualify by fidelity to ordinances and covenants made in His holy temples.

The many inspiring components of this April 2020 general conference—and the sacred week that we now begin—can be summarized by two divinely decreed words: “Hear Him.” We pray that your focus on Heavenly Father, who spoke those words, and on His Beloved Son, Jesus Christ, will loom largest in your memories of all that has transpired. We pray that you will begin anew truly to hear, hearken to, and heed the words of the Savior. I promise that decreased fear and increased faith will follow.

Thank you for your desire to make your homes true sanctuaries of faith, where the Spirit of the Lord may dwell. Our gospel study curriculum, Come, Follow Me, will continue to bless your lives. Your consistent efforts in this endeavor—even during those moments when you feel that you are not being particularly successful—will change your life, that of your family, and the world. We will be strengthened as we become even more valiant disciples of the Lord, standing up and speaking up for Him, wherever we are.

Now, let’s talk about temples. We have 168 dedicated temples across the world. Others are at various stages of planning and construction. When plans are announced to erect a new temple, it becomes part of our sacred history.

It may seem odd to announce new temples when all our temples are closed for a while.

More than a century ago, President Wilford Woodruff foresaw conditions such as ours today, as recorded in his dedicatory prayer of the Salt Lake Temple, given in 1893. Some of you may have recently seen excerpts from this remarkable prayer on social media.

Hear these pleadings from a mighty prophet of God: “When Thy people shall not have the opportunity of entering this holy house … and they are oppressed and in trouble, surrounded by difficulties … and shall turn their faces towards this Thy holy house and ask Thee for deliverance, for help, for Thy power to be extended in their behalf, we beseech Thee, to look down from Thy holy habitation in mercy … and listen to their cries. Or when the children of Thy people, in years to come, shall be separated, through any cause, from this place, … and they shall cry unto Thee from the depths of their affliction and sorrow to extend relief and deliverance to them, we humbly entreat Thee to … hearken to their cries, and grant unto them the blessings for which they ask.”

Brothers and sisters, during times of our distress when temples are closed, you can still draw upon the power of your temple covenants and endowment as you honor your covenants. Please use this time when temples are closed to continue to live a temple-worthy life or to become temple worthy.

Talk about the temple with your family and friends. Because Jesus Christ is at the center of everything we do in the temple, as you think more about the temple you will be thinking more about Him. Study and pray to learn more about the power and knowledge with which you have been endowed—or with which you will yet be endowed.

Today we are pleased to announce plans to construct eight new temples in the following locations: Bahía Blanca, Argentina; Tallahassee, Florida; Lubumbashi, Democratic Republic of the Congo; Pittsburgh, Pennsylvania; Benin City, Nigeria; Syracuse, Utah; Dubai, United Arab Emirates; and Shanghai, People’s Republic of China.

In all eight locations, Church architects will work with local officials so that the temple will harmonize with and be a beautiful addition to each community.

The plan for a temple in Dubai comes in response to their gracious invitation, which we gratefully acknowledge.

Context for the plan for Shanghai is very important. For more than two decades, temple-worthy members in the People’s Republic of China have attended the Hong Kong China Temple. But in July 2019, that temple was closed for long-planned and much-needed renovation.

In Shanghai, a modest multipurpose meeting place will provide a way for Chinese members to continue to participate in ordinances of the temple—in the People’s Republic of China—for them and their ancestors.

In every country, this Church teaches its members to honor, obey, and sustain the law. We teach the importance of the family, of being good parents and exemplary citizens. Because we respect the laws and regulations of the People’s Republic of China, the Church does not send proselyting missionaries there; nor will we do so now.

Expatriate and Chinese congregations will continue to meet separately. The Church’s legal status there remains unchanged. In an initial phase of facility use, entry will be by appointment only. The house of the Lord in Shanghai will not be a destination place for tourists from other countries.

These eight new temples will bless the lives of many people on both sides of the veil of death. Temples are a crowning part of the Restoration of the fulness of the gospel of Jesus Christ. In God’s goodness and generosity, He is bringing the blessings of the temple closer to His children everywhere.

As the Restoration continues, I know that God will continue to reveal many great and important things pertaining to His kingdom here on earth. That kingdom is The Church of Jesus Christ of Latter-day Saints.

Dear brothers and sisters, I express my love for you. During this time of tension and uncertainty, and invoking the authority vested in me, I would like to confer upon you an apostolic blessing.

I bless you with peace and increasing faith in the Lord.

I bless you with a desire to repent and become a little more like Him each day.

I bless you to know that the Prophet Joseph Smith is the prophet of the Restoration of the gospel of Jesus Christ in its fulness.

Should there be illness among you or your loved ones, I leave a blessing of healing, consistent with the will of the Lord.

I so bless you, adding once more my expression of love for each of you, in the sacred name of Jesus Christ, amen.

\newpage
\settalk{Moving Forward}
\setspeaker{Russell M. Nelson}
\settalkdate{October 2020}
\talkheading{Moving Forward}{Russell M. Nelson}

My dear brothers and sisters, what a joy it is to be with you as we begin the 190th Semiannual General Conference of The Church of Jesus Christ of Latter-day Saints. I love joining with you in your homes or wherever you are to listen together to the messages of prophets, seers, and revelators and other Church leaders.

How grateful we are for the technology that allows us to be connected as one great worldwide gathering of disciples of Jesus Christ. General conference last April was viewed by more people than any preceding it, and we have every expectation that will happen again.

During the past few months, a global pandemic, raging wildfires, and other natural disasters have turned our world upside down. I grieve with each of you who has lost a loved one during this time. And I pray for all who are currently suffering.

Meanwhile, the work of the Lord is steadily moving forward. Amid social distancing, face masks, and Zoom meetings, we have learned to do some things differently and some even more effectively. Unusual times can bring unusual rewards.

Our missionaries and mission leaders have been resourceful, resilient, and truly remarkable. Although most missionaries have had to find new, creative ways to do their work, many missions have reported doing more teaching than ever.

We had to close temples for a time, and some construction projects were briefly delayed, but now they are all moving forward. In the calendar year 2020, we will have broken ground for 20 new temples!

Family history work has increased exponentially. Many new wards and stakes have been created. And we are gratified to report that the Church has provided pandemic humanitarian aid for 895 projects in 150 countries.

Increased gospel study in many homes is resulting in stronger testimonies and family relationships. One mother wrote: “We feel much closer to our children and grandchildren now that we gather on Zoom every Sunday. Each takes a turn giving their thoughts on Come, Follow Me. Prayers for our family members have changed because we better understand what they need.”

I pray that we as a people are using this unique time to grow spiritually. We are here on earth to be tested, to see if we will choose to follow Jesus Christ, to repent regularly, to learn, and to progress. Our spirits long to progress. And we do that best by staying firmly on the covenant path.

Through it all, our Heavenly Father and His Son, Jesus Christ, love us! They care for us! They and Their holy angels are watching over us. I know that is true.

As we gather to hear the words the Lord has inspired His servants to deliver, I invite you to ponder a promise the Lord made. He declared that “whosoever will may lay hold upon the word of God, which is quick and powerful, which shall divide asunder all the cunning … and the wiles of the devil, and lead the [disciple] of Christ in a strait and narrow course.”

I pray that you will choose to lay hold upon the word of God as it is declared during this general conference. And I pray that you may feel the Lord’s perfect love for you, in the sacred name of Jesus Christ, amen.

\newpage
\settalk{Embrace the Future with Faith}
\setspeaker{Russell M. Nelson}
\settalkdate{October 2020}
\talkheading{Embrace the Future with Faith}{Russell M. Nelson}

This has been an unforgettable evening. My dear sisters, I am honored to be with you. You have been on my mind so often during these past few months. You are more than eight million strong. You have not only the numbers but the spiritual power to change the world. I have watched you doing just that during this pandemic.

Some of you suddenly found yourselves searching for scarce supplies or a new job. Many tutored children and checked on neighbors. Some welcomed missionaries home earlier than expected, while others transformed your homes into missionary training centers. You have used technology to connect with family and friends, to minister to those who have felt isolated, and to study Come, Follow Me with others. You have found new ways to make the Sabbath a delight. And you have made protective masks—millions of them!

With heartfelt compassion and love, my heart goes out to the many women around the world whose loved ones have died. We weep with you. And we pray for you. We praise and pray for all who work tirelessly to safeguard the health of others.

You young women have also been remarkable. Though social media has been flooded with contention, many of you have found ways to encourage others and share our Savior’s light.

Sisters, you have all been absolutely heroic! I marvel at your strength and your faith. You have shown that in difficult circumstances, you bravely carry on. I love you, and I assure you that the Lord loves you and sees the great work you are performing. Thank you! Once again, you have proven that you are literally the hope of Israel!

You embody the hopes that President Gordon B. Hinckley had for you when he introduced “The Family: A Proclamation to the World” 25 years ago in the September 1995 general Relief Society meeting. It is significant that he chose to introduce this important proclamation to the sisters of the Church. By doing so, President Hinckley underscored the irreplaceable influence of women in the Lord’s plan.

Now, I would love to know what you have learned this year. Have you grown closer to the Lord, or do you feel further away from Him? And how have current events made you feel about the future?

Admittedly, the Lord has spoken of our day in sobering terms. He warned that in our day “men’s hearts [would fail] them” and that even the very elect would be at risk of being deceived. He told the Prophet Joseph Smith that “peace [would] be taken from the earth” and calamities would befall mankind.

Yet the Lord has also provided a vision of how remarkable this dispensation is. He inspired the Prophet Joseph Smith to declare that “the work of … these last days, is one of vast magnitude. … Its glories are past description, and its grandeur unsurpassable.”

Now, grandeur may not be the word you would choose to describe these past few months! How are we to deal with both the somber prophecies and the glorious pronouncements about our day? The Lord told us how with simple, but stunning, reassurance: “If ye are prepared ye shall not fear.”

What a promise! It is one that can literally change the way we see our future. I recently heard a woman of deep testimony admit that the pandemic, combined with an earthquake in the Salt Lake Valley, had helped her realize she was not as prepared as she thought she was. When I asked whether she was referring to her food storage or her testimony, she smiled and said, “Yes!”

If preparation is our key to embracing this dispensation and our future with faith, how can we best prepare?

For decades, the Lord’s prophets have urged us to store food, water, and financial reserves for a time of need. The current pandemic has reinforced the wisdom of that counsel. I urge you to take steps to be temporally prepared. But I am even more concerned about your spiritual and emotional preparation.

In that regard, we can learn a lot from Captain Moroni. As commander of the Nephite armies, he faced opposing forces that were stronger, greater in number, and meaner. So, Moroni prepared his people in three essential ways.

First, he helped them create areas where they would be safe—“places of security” he called them. Second, he prepared “the minds of the people to be faithful unto the Lord.” And third, he never stopped preparing his people—physically or spiritually. Let us consider these three principles.

\intalkheader{Principle Number One: Create Places of Security}

Moroni fortified every Nephite city with embankments, forts, and walls. When the Lamanites came against them, they “were astonished exceedingly, because of the wisdom of the Nephites in preparing their places of security.”

Similarly, as turmoil rages around us, we need to create places where we are safe, both physically and spiritually. When your home becomes a personal sanctuary of faith—where the Spirit resides—your home becomes the first line of defense.

Likewise, the stakes of Zion are “a refuge from the storm” because they are led by those who hold priesthood keys and exercise priesthood authority. As you continue to follow the counsel of those whom the Lord has authorized to guide you, you will feel greater safety.

The temple—the house of the Lord—is a place of security unlike any other. There, you sisters are endowed with priesthood power through the sacred priesthood covenants you make. There, your families are sealed for eternity. Even this year, when access to our temples has been seriously limited, your endowment has given you constant access to God’s power as you have honored your covenants with Him.

Simply said, a place of security is anywhere you can feel the presence of the Holy Ghost and be guided by Him. When the Holy Ghost is with you, you can teach truth, even when it runs counter to prevailing opinions. And you can ponder sincere questions about the gospel in an environment of revelation.

I invite you, my dear sisters, to create a home that is a place of security. And I renew my invitation for you to increase your understanding of priesthood power and of temple covenants and blessings. Having places of security to which you can retreat will help you embrace the future with faith.

\intalkheader{Principle Number Two: Prepare Your Mind to Be Faithful to God}

We have undertaken a major project to extend the life and capacity of the Salt Lake Temple.

Some questioned the need for taking such extraordinary measures. However, when the Salt Lake Valley suffered a 5.7-magnitude earthquake earlier this year, this venerable temple shook hard enough that the trumpet on the statue of the angel Moroni fell!

Just as the physical foundation of the Salt Lake Temple must be strong enough to withstand natural disasters, our spiritual foundations must be solid. Then, when metaphorical earthquakes rock our lives, we can stand “steadfast and immovable” because of our faith.

The Lord taught us how to increase our faith by seeking “learning, even by study and also by faith.” We strengthen our faith in Jesus Christ as we strive to keep His commandments and “always remember him.” Further, our faith increases every time we exercise our faith in Him. That is what learning by faith means.

For example, each time we have the faith to be obedient to God’s laws—even when popular opinions belittle us—or each time we resist entertainment or ideologies that celebrate covenant-breaking, we are exercising our faith, which in turn increases our faith.

Further, few things build faith more than does regular immersion in the Book of Mormon. No other book testifies of Jesus Christ with such power and clarity. Its prophets, as inspired by the Lord, saw our day and selected the doctrine and truths that would help us most. The Book of Mormon is our latter-day survival guide.

Of course, our ultimate security comes as we yoke ourselves to Heavenly Father and Jesus Christ! Life without God is a life filled with fear. Life with God is a life filled with peace. This is because spiritual blessings come to the faithful. Receiving personal revelation is one of the greatest of those blessings.

The Lord has promised that if we will ask, we may receive “revelation upon revelation.” I promise that as you increase your capacity to receive revelation, the Lord will bless you with increased direction for your life and with boundless gifts of the Spirit.

\intalkheader{Principle Number Three: Never Stop Preparing}

Even when things went well, Captain Moroni continued to prepare his people. He never stopped. He never became complacent.

The adversary never stops attacking. So, we can never stop preparing! The more self-reliant we are—temporally, emotionally, and spiritually—the more prepared we are to thwart Satan’s relentless assaults.

Dear sisters, you are adept at creating places of security for yourselves and those you love. Further, you have a divine endowment that enables you to build faith in others in compelling ways. And you never stop. You have demonstrated that once again this year.

Please, keep going! Your vigilance in safeguarding your homes and instilling faith in the hearts of your loved ones will reap rewards for generations to come.

My dear sisters, we have so much to look forward to! The Lord placed you here now because He knew you had the capacity to negotiate the complexities of the latter part of these latter days. He knew you would grasp the grandeur of His work and be eager to help bring it to pass.

I am not saying that the days ahead will be easy, but I promise you that the future will be glorious for those who are prepared and who continue to prepare to be instruments in the Lord’s hands.

My dear sisters, let us not just endure this current season. Let us embrace the future with faith! Turbulent times are opportunities for us to thrive spiritually. They are times when our influence can be much more penetrating than in calmer times.

I promise that as we create places of security, prepare our minds to be faithful to God, and never stop preparing, God will bless us. He will “deliver us; yea, insomuch that he [will] speak peace to our souls, and [will] grant unto us great faith, and … cause us that we [can] hope for our deliverance in him.”

As you prepare to embrace the future with faith, these promises will be yours! I so testify, with my expression of love for you and my confidence in you, in the sacred name of Jesus Christ, amen.

\newpage
\settalk{Let God Prevail}
\setspeaker{Russell M. Nelson}
\settalkdate{October 2020}
\talkheading{Let God Prevail}{Russell M. Nelson}

My dear brothers and sisters, how grateful I am for the marvelous messages of this conference and for my privilege to speak with you now.

For the more than 36 years I’ve been an Apostle, the doctrine of the gathering of Israel has captured my attention. Everything about it has intrigued me, including the ministries and names of Abraham, Isaac, and Jacob; their lives and their wives; the covenant God made with them and extended through their lineage; the dispersion of the twelve tribes; and the numerous prophecies about the gathering in our day.

I have studied the gathering, prayed about it, feasted upon every related scripture, and asked the Lord to increase my understanding.

So imagine my delight when I was led recently to a new insight. With the help of two Hebrew scholars, I learned that one of the Hebraic meanings of the word Israel is “let God prevail.” Thus the very name of Israel refers to a person who is willing to let God prevail in his or her life. That concept stirs my soul!

The word willing is crucial to this interpretation of Israel. We all have our agency. We can choose to be of Israel, or not. We can choose to let God prevail in our lives, or not. We can choose to let God be the most powerful influence in our lives, or not.

For a moment, let us recall a crucial turning point in the life of Jacob, the grandson of Abraham. At the place Jacob named Peniel (which means “the face of God”), Jacob wrestled with a serious challenge. His agency was tested. Through this wrestle, Jacob proved what was most important to him. He demonstrated that he was willing to let God prevail in his life. In response, God changed Jacob’s name to Israel, meaning “let God prevail.” God then promised Israel that all the blessings that had been pronounced upon Abraham’s head would also be his.

Sadly, Israel’s posterity broke their covenants with God. They stoned the prophets and were not willing to let God prevail in their lives. Subsequently, God scattered them to the four corners of the earth. Mercifully, He later promised to gather them, as reported by Isaiah: “For a small moment have I forsaken thee [Israel]; but with great mercies will I gather thee.”

With the Hebraic definition of Israel in mind, we find that the gathering of Israel takes on added meaning. The Lord is gathering those who are willing to let God prevail in their lives. The Lord is gathering those who will choose to let God be the most important influence in their lives.

For centuries, prophets have foretold this gathering, and it is happening right now! As an essential prelude to the Second Coming of the Lord, it is the most important work in the world!

This premillennial gathering is an individual saga of expanding faith and spiritual courage for millions of people. And as members of The Church of Jesus Christ of Latter-day Saints, or “latter-day covenant Israel,” we have been charged to assist the Lord with this pivotal work.

When we speak of gathering Israel on both sides of the veil, we are referring, of course, to missionary, temple, and family history work. We are also referring to building faith and testimony in the hearts of those with whom we live, work, and serve. Anytime we do anything that helps anyone—on either side of the veil—to make and keep their covenants with God, we are helping to gather Israel.

Not long ago, the wife of one of our grandsons was struggling spiritually. I will call her “Jill.” Despite fasting, prayer, and priesthood blessings, Jill’s father was dying. She was gripped with fear that she would lose both her dad and her testimony.

Late one evening, my wife, Sister Wendy Nelson, told me of Jill’s situation. The next morning Wendy felt impressed to share with Jill that my response to her spiritual wrestle was one word! The word was myopic.

Jill later admitted to Wendy that initially she was devastated by my response. She said, “I was hoping for Grandfather to promise me a miracle for my dad. I kept wondering why the word myopic was the one he felt compelled to say.”

After Jill’s father passed on, the word myopic kept coming to her mind. She opened her heart to understand even more deeply that myopic meant “nearsighted.” And her thinking began to shift. Jill then said, “Myopic caused me to stop, think, and heal. That word now fills me with peace. It reminds me to expand my perspective and seek the eternal. It reminds me that there is a divine plan and that my dad still lives and loves and looks out for me. Myopic has led me to God.”

I am very proud of our precious granddaughter-in-law. During this heart-wrenching time in her life, dear Jill is learning to embrace God’s will for her dad, with an eternal perspective for her own life. By choosing to let God prevail, she is finding peace.

If we will allow it, there are many ways this Hebraic interpretation of Israel can help us. Imagine how our prayers for our missionaries—and for our own efforts to gather Israel—could change with this concept in mind. We often pray that we and the missionaries will be led to those who are prepared to receive the truths of the restored gospel of Jesus Christ. I wonder, to whom will we be led when we plead to find those who are willing to let God prevail in their lives?

We may be led to some who have never believed in God or Jesus Christ but who are now yearning to learn about Them and Their plan of happiness. Others may have been “born in the covenant” but have since wandered away from the covenant path. They may now be ready to repent, return, and let God prevail. We can assist them by welcoming them with open arms and hearts. And some to whom we may be led may have always felt there was something missing in their lives. They too are longing for the wholeness and joy that come to those who are willing to let God prevail in their lives.

The gospel net to gather scattered Israel is expansive. There is room for each person who will fully embrace the gospel of Jesus Christ. Each convert becomes one of God’s covenant children, whether by birth or by adoption. Each becomes a full heir to all that God has promised the faithful children of Israel!

Each of us has a divine potential because each is a child of God. Each is equal in His eyes. The implications of this truth are profound. Brothers and sisters, please listen carefully to what I am about to say. God does not love one race more than another. His doctrine on this matter is clear. He invites all to come unto Him, “black and white, bond and free, male and female.”

I assure you that your standing before God is not determined by the color of your skin. Favor or disfavor with God is dependent upon your devotion to God and His commandments and not the color of your skin.

I grieve that our Black brothers and sisters the world over are enduring the pains of racism and prejudice. Today I call upon our members everywhere to lead out in abandoning attitudes and actions of prejudice. I plead with you to promote respect for all of God’s children.

The question for each of us, regardless of race, is the same. Are you willing to let God prevail in your life? Are you willing to let God be the most important influence in your life? Will you allow His words, His commandments, and His covenants to influence what you do each day? Will you allow His voice to take priority over any other? Are you willing to let whatever He needs you to do take precedence over every other ambition? Are you willing to have your will swallowed up in His?

Consider how such willingness could bless you. If you are unmarried and seeking an eternal companion, your desire to be “of Israel” will help you decide whom to date and how.

If you are married to a companion who has broken his or her covenants, your willingness to let God prevail in your life will allow your covenants with God to remain intact. The Savior will heal your broken heart. The heavens will open as you seek to know how to move forward. You do not need to wander or wonder.

If you have sincere questions about the gospel or the Church, as you choose to let God prevail, you will be led to find and understand the absolute, eternal truths that will guide your life and help you stay firmly on the covenant path.

When you are faced with temptation—even if the temptation comes when you are exhausted or feeling alone or misunderstood—imagine the courage you can muster as you choose to let God prevail in your life and as you plead with Him to strengthen you.

When your greatest desire is to let God prevail, to be part of Israel, so many decisions become easier. So many issues become nonissues! You know how best to groom yourself. You know what to watch and read, where to spend your time, and with whom to associate. You know what you want to accomplish. You know the kind of person you really want to become.

Now, my dear brothers and sisters, it takes both faith and courage to let God prevail. It takes persistent, rigorous spiritual work to repent and to put off the natural man through the Atonement of Jesus Christ. It takes consistent, daily effort to develop personal habits to study the gospel, to learn more about Heavenly Father and Jesus Christ, and to seek and respond to personal revelation.

During these perilous times of which the Apostle Paul prophesied, Satan is no longer even trying to hide his attacks on God’s plan. Emboldened evil abounds. Therefore, the only way to survive spiritually is to be determined to let God prevail in our lives, to learn to hear His voice, and to use our energy to help gather Israel.

Now, how does the Lord feel about people who will let God prevail? Nephi summed it up well: “[The Lord] loveth those who will have him to be their God. Behold, he loved our fathers, and he covenanted with them, yea, even Abraham, Isaac, and Jacob; and he remember[s] the covenants which he [has] made.”

And what is the Lord willing to do for Israel? The Lord has pledged that He will “fight [our] battles, and [our] children’s battles, and our children’s children’s [battles] … to the third and fourth generation”!

As you study your scriptures during the next six months, I encourage you to make a list of all that the Lord has promised He will do for covenant Israel. I think you will be astounded! Ponder these promises. Talk about them with your family and friends. Then live and watch for these promises to be fulfilled in your own life.

My dear brothers and sisters, as you choose to let God prevail in your lives, you will experience for yourselves that our God is “a God of miracles.” As a people, we are His covenant children, and we will be called by His name. Of this I testify in the sacred name of Jesus Christ, amen.

\newpage
\settalk{A New Normal}
\setspeaker{Russell M. Nelson}
\settalkdate{October 2020}
\talkheading{A New Normal}{Russell M. Nelson}

My dear brothers and sisters, these two days of general conference have been glorious! I agree with Elder Jeffrey R. Holland. As he mentioned, the messages, the prayers, and the music have all been inspired by the Lord. I am grateful to all who have participated in any way.

Throughout the proceedings, I have pictured you in my mind listening to conference. I have asked the Lord to help me understand what you are feeling, worrying about, or trying to resolve. I have wondered what I might say to conclude this conference that would send you forth with the optimism about the future that I know the Lord wants you to feel.

We live in a glorious age, foreseen by prophets for centuries. This is the dispensation when no spiritual blessing will be withheld from the righteous. Despite the world’s commotion, the Lord would have us look forward to the future “with joyful anticipation.” Let us not spin our wheels in the memories of yesterday. The gathering of Israel moves forward. The Lord Jesus Christ directs the affairs of His Church, and it will achieve its divine objectives.

The challenge for you and me is to make certain that each of us will achieve his or her divine potential. Today we often hear about “a new normal.” If you really want to embrace a new normal, I invite you to turn your heart, mind, and soul increasingly to our Heavenly Father and His Son, Jesus Christ. Let that be your new normal.

Embrace your new normal by repenting daily. Seek to be increasingly pure in thought, word, and deed. Minister to others. Keep an eternal perspective. Magnify your callings. And whatever your challenges, my dear brothers and sisters, live each day so that you are more prepared to meet your Maker.

That is why we have temples. The Lord’s ordinances and covenants prepare us for eternal life, the greatest of all of God’s blessings. As you know, the COVID pandemic required a temporary closure of our temples. Then we commenced a carefully coordinated, phased reopening. With phase 2 now in place in many temples, thousands of couples have been sealed and thousands have received their own endowments just in the past few months. We look forward to the day when all worthy members of the Church can again serve their ancestors and worship in a holy temple.

Now I am pleased to announce plans for the construction of six new temples to be built in the following locations: Tarawa, Kiribati; Port Vila, Vanuatu; Lindon, Utah; Greater Guatemala City, Guatemala; São Paulo East, Brazil; and Santa Cruz, Bolivia.

As we build and maintain these temples, we pray that each of you will build and maintain yourself so you can be worthy to enter the holy temple.

Now, my dear brothers and sisters, I bless you to be filled with the peace of the Lord Jesus Christ. His peace is beyond all mortal understanding. I bless you with an increased desire and ability to obey the laws of God. I promise that as you do, you will be showered with blessings, including greater courage, increased personal revelation, sweeter harmony in your homes, and joy even amid uncertainty.

May we go forward together to fulfill our divine mandate—that of preparing ourselves and the world for the Second Coming of the Lord. I so pray, with my expression of love for you, in the sacred name of Jesus Christ, amen.

\newpage
\settalk{Welcome Message}
\setspeaker{Russell M. Nelson}
\settalkdate{April 2021}
\talkheading{Welcome Message}{Russell M. Nelson}

My dear brothers and sisters and friends throughout the world, I offer my personal welcome to this general conference. We gather as a great global family desiring to worship our Lord and Savior, Jesus Christ. Thank you for joining with us.

This last year has been one for the record books. No doubt we have each learned things we did not know previously. Some lessons that I knew before have been written on my heart in new and instructive ways.

For example, I know for sure that the Lord directs the affairs of His Church. He said, “I will show [you] that I am able to do mine own work.”

Often, my counselors and I have watched through tear-brimmed eyes as He has interceded in extremely challenging circumstances after we had done our best and could do no more. We do indeed stand all amazed.

I also understand better now what He meant when He said, “Behold, I will hasten my work in its time.” Over and over again I have rejoiced as He has directed and executed the hastening of His work—even during a global pandemic.

My dear brothers and sisters, the strength of the Church lies in the efforts and ever-growing testimonies of its members. Testimonies are best cultivated in the home. During this past year, many of you have dramatically increased the study of the gospel in your homes. I thank you, and your children will thank you.

The huge project to renovate the Salt Lake Temple continues. From my office I have a front-row seat to watch the work taking place on the temple plaza.

As I have watched workers dig out old tree roots, plumbing, wiring, and a leaky fountain, I have thought about the need for each of us to remove, with the Savior’s help, the old debris in our lives.

The gospel of Jesus Christ is a gospel of repentance. Because of the Savior’s Atonement, His gospel provides an invitation to keep changing, growing, and becoming more pure. It is a gospel of hope, of healing, and of progress. Thus, the gospel is a message of joy! Our spirits rejoice with every small step forward we take.

Part of the gathering of Israel, and a very important part, is the charge for us as a people to be worthy and willing to help prepare the world for the Second Coming of the Lord.

As we listen to the messages that have been carefully prepared by our leaders under the direction of the Holy Ghost, I invite you to pray to identify the debris you should remove from your life so you can become more worthy.

I love you, my dear brothers and sisters, and testify that our Heavenly Father and His Beloved Son know and love you individually. They stand ready to assist you in every step forward you take. Welcome to general conference and to the privilege of hearing the voice of the Lord. In the name of Jesus Christ, amen.

\newpage
\settalk{What We Are Learning and Will Never Forget}
\setspeaker{Russell M. Nelson}
\settalkdate{April 2021}
\talkheading{What We Are Learning and Will Never Forget}{Russell M. Nelson}

My dear brethren, I have looked forward to this virtual meeting with you. The last time we held a priesthood session of general conference was in April 2019. Much has happened in the past two years! Some of you have lost loved ones. Others have lost jobs, livelihood, or health. Still others have lost a sense of peace or hope for the future. My heart goes out to each one of you who has suffered these or other losses. I pray constantly that the Lord will comfort you. As you continue to let God prevail in your life, I know that He is just as optimistic about your future as He has ever been.

Amid the losses we have experienced, there are also some things we have found. Some have found deeper faith in our Heavenly Father and His Son, Jesus Christ. Many have found a fresh perspective on life—even an eternal perspective. You may have found stronger relationships with your loved ones and with the Lord. I hope you have found an increased ability to hear Him and receive personal revelation. Difficult trials often provide opportunities to grow that would not have come in any other way.

Think back on the past two years. How have you grown? What have you learned? You might initially wish you could go back to 2019 and stay there! But if you look at your life prayerfully, I believe you will see many ways in which the Lord has been guiding you through this time of hardship, helping you to become a more devoted, more converted man—a true man of God.

I know the Lord has great and marvelous plans for us—individually and collectively. With compassion and patience, He says:

“Ye are little children, and ye have not as yet understood how great blessings the Father hath … prepared for you;

“And ye cannot bear all things now; nevertheless, be of good cheer, for I will lead you along.”

My dear brothers, I testify that He has been, and is, indeed leading us along, as we seek to hear Him. He wants us to grow and to learn, even through—perhaps especially through—adversity.

Adversity is a great teacher. What have you learned in the past two years that you always want to remember? Your answers will be unique to you, but may I suggest four lessons I hope we have all learned and will never forget.

\intalkheader{Lesson 1: The Home Is the Center of Faith and Worship}

Often when the Lord warns us about the perils of the last days, He counsels thus: “Stand ye in holy places, and be not moved.” These “holy places” certainly include the Lord’s temples and meetinghouses. But as our ability to gather in these places has been restricted in varying degrees, we have learned that one of the holiest of places on earth is the home—yes, even your home.

Brethren, you bear the priesthood of God. “The rights of the priesthood are inseparably connected with the powers of heaven.” You and your family have received priesthood ordinances. It is “in the ordinances [of the priesthood that] the power of godliness is manifest.” That power is available to you and your family in your own home as you keep the covenants you have made.

Just 185 years ago, this very day, April 3, 1836, Elijah restored the keys of the priesthood that allow our families to be sealed together forever. That is why it felt so good to administer the sacrament in your home. How do you think it affected your family members to see you—their father, grandfather, husband, son, or brother—administer this holy ordinance? What will you do to retain that sacred feeling in your family?

You may feel that there is still more you need to do to make your home truly a sanctuary of faith. If so, please do it! If you are married, counsel with your wife as your equal partner in this crucial work. There are few pursuits more important than this. Between now and the time the Lord comes again, we all need our homes to be places of serenity and security.

Attitudes and actions that invite the Spirit will increase the holiness of your home. Equally certain is the fact that holiness will vanish if there is anything in your behavior or environment that offends the Holy Spirit, for then “the heavens withdraw themselves.”

Have you ever wondered why the Lord wants us to make our homes the center of gospel learning and gospel living? It is not just to prepare us for, and help us through, a pandemic. Present restrictions on gathering will eventually end. However, your commitment to make your home your primary sanctuary of faith should never end. As faith and holiness decrease in this fallen world, your need for holy places will increase. I urge you to continue to make your home a truly holy place “and be not moved” from that essential goal.

\intalkheader{Lesson 2: We Need Each Other}

God wants us to work together and help each other. That is why He sends us to earth in families and organizes us into wards and stakes. That is why He asks us to serve and minister to each other. That is why He asks us to live in the world but not be of the world. We can accomplish so much more together than we can alone. God’s plan of happiness would be frustrated if His children remained isolated one from another.

The recent pandemic has been unique in that it has affected everyone in the world at essentially the same time. While some have suffered more than others, we have all been challenged in some way. Because of this, our common trial has the potential to help unite God’s children as never before. So, I ask, has this shared trial drawn you closer to your neighbors—to your brothers and sisters across the street and around the world?

In this regard, the two great commandments can guide us: first, to love God and, second, to love our neighbor. We show our love by serving.

If you know of anyone who is alone, reach out—even if you feel alone too! You do not need to have a reason or a message or business to transact. Just say hello and show your love. Technology can help you. Pandemic or not, each precious child of God needs to know that he or she is not alone!

\intalkheader{Lesson 3: Your Priesthood Quorum Is Meant for More Than Just a Meeting}

During the pandemic, Sunday quorum meetings were canceled for a time. Some quorums are now able to meet virtually. Nevertheless, the work that the Lord has given to priesthood quorums was never meant to be confined to a meeting. Meetings are only a small part of what a quorum is meant for and what it can do.

My brethren of the Aaronic Priesthood and elders quorums, expand your vision of why we have quorums. How does the Lord wish you would use your quorum to accomplish His work—now? Seek revelation from the Lord. Humble yourself! Ask! Listen! If you have been called to lead, counsel as a presidency and with quorum members. Whatever your priesthood office or calling, let God prevail in your commitment as a member of your quorum and in your service. Experience with joy the righteousness you will bring to pass as you are “anxiously engaged in a good cause.” Quorums are in a unique position to accelerate the gathering of Israel on both sides of the veil.

\intalkheader{Lesson 4: We Hear Jesus Christ Better When We Are Still}

We live in a time prophesied long ago, when “all things shall be in commotion; and surely, men’s hearts shall fail them; for fear shall come upon all people.” That was true before the pandemic, and it will be true after. Commotion in the world will continue to increase. In contrast, the voice of the Lord is not “a voice of a great tumultuous noise, but … it [is] a still voice of perfect mildness, [like] a whisper, and it [pierces] even to the very soul.” In order to hear this still voice, you too must be still!

For a time, the pandemic has canceled activities that would normally fill our lives. Soon we may be able to choose to fill that time again with the noise and commotion of the world. Or we can use our time to hear the voice of the Lord whispering His guidance, comfort, and peace. Quiet time is sacred time—time that will facilitate personal revelation and instill peace.

Discipline yourself to have time alone and with your loved ones. Open your heart to God in prayer. Take time to immerse yourself in the scriptures and worship in the temple.

My dear brethren, there are many things the Lord wants us to learn from our experiences during this pandemic. I have listed only four. I invite you to make your own list, consider it carefully, and share it with those you love.

The future is bright for God’s covenant-keeping people. The Lord will increasingly call upon His servants who worthily hold the priesthood to bless, comfort, and strengthen mankind and to help prepare the world and its people for His Second Coming. It behooves each of us to measure up to the sacred ordination we have received. We can do this! I so testify, with my expression of love for each of you, my beloved brethren, in the sacred name of Jesus Christ, amen.

\newpage
\settalk{Christ Is Risen; Faith in Him Will Move Mountains}
\setspeaker{Russell M. Nelson}
\settalkdate{April 2021}
\talkheading{Christ Is Risen; Faith in Him Will Move Mountains}{Russell M. Nelson}

My dear brothers and sisters, I am grateful for the privilege of speaking with you on this Easter Sunday. The atoning sacrifice and Resurrection of Jesus Christ changed each of our lives forever. We love Him and gratefully worship Him and our Heavenly Father.

During the past six months, we have continued to grapple with a global pandemic. I marvel at your resilience and spiritual strength in the face of illness, loss, and isolation. I pray constantly that, through it all, you will feel the Lord’s unfailing love for you. If you have responded to your trials with a stronger discipleship, this past year will not have been in vain.

This morning, we have heard from Church leaders who come from every populated continent on earth. Truly, the blessings of the gospel are for every race, language, and people. The Church of Jesus Christ is a global church. Jesus Christ is our leader.

Thankfully, even a pandemic has not been able to slow the onward march of His truth. The gospel of Jesus Christ is exactly what is needed in this confused, contentious, and weary world.

Each of God’s children deserves the opportunity to hear and accept the healing, redeeming message of Jesus Christ. No other message is more vital to our happiness—now and forever. No other message is more filled with hope. No other message can eliminate contention in our society.

Faith in Jesus Christ is the foundation of all belief and the conduit of divine power. According to the Apostle Paul, “Without faith it is impossible to please [God]: for he that cometh to God must believe that he is, and that he is a rewarder of them that diligently seek him.”

Everything good in life—every potential blessing of eternal significance—begins with faith. Allowing God to prevail in our lives begins with faith that He is willing to guide us. True repentance begins with faith that Jesus Christ has the power to cleanse, heal, and strengthen us.

“Deny not the power of God,” the prophet Moroni declared, “for he worketh by power, according to the faith of the children of men.” It is our faith that unlocks the power of God in our lives.

And yet, exercising faith can seem overwhelming. At times we may wonder if we can possibly muster enough faith to receive the blessings that we so desperately need. However, the Lord put those fears to rest through the words of the Book of Mormon prophet Alma.

Alma asks us simply to experiment upon the word and “exercise a particle of faith, yea, even if [we] can no more than desire to believe.” The phrase “particle of faith” reminds me of the Lord’s biblical promise that if we “have faith as a grain of mustard seed,” we shall be able to “say unto this mountain, Remove hence to yonder place; and it shall remove; and nothing shall be impossible unto [us].”

The Lord understands our mortal weakness. We all falter at times. But He also knows of our great potential. The mustard seed starts small but grows into a tree large enough for birds to nest in its branches. The mustard seed represents a small but growing faith.

The Lord does not require perfect faith for us to have access to His perfect power. But He does ask us to believe.

My dear brothers and sisters, my call to you this Easter morning is to start today to increase your faith. Through your faith, Jesus Christ will increase your ability to move the mountains in your life, even though your personal challenges may loom as large as Mount Everest.

Your mountains may be loneliness, doubt, illness, or other personal problems. Your mountains will vary, and yet the answer to each of your challenges is to increase your faith. That takes work. Lazy learners and lax disciples will always struggle to muster even a particle of faith.

To do anything well requires effort. Becoming a true disciple of Jesus Christ is no exception. Increasing your faith and trust in Him takes effort. May I offer five suggestions to help you develop that faith and trust.

First, study. Become an engaged learner. Immerse yourself in the scriptures to understand better Christ’s mission and ministry. Know the doctrine of Christ so that you understand its power for your life. Internalize the truth that the Atonement of Jesus Christ applies to you. He took upon Himself your misery, your mistakes, your weakness, and your sins. He paid the compensatory price and provided the power for you to move every mountain you will ever face. You obtain that power with your faith, trust, and willingness to follow Him.

Moving your mountains may require a miracle. Learn about miracles. Miracles come according to your faith in the Lord. Central to that faith is trusting His will and timetable—how and when He will bless you with the miraculous help you desire. Only your unbelief will keep God from blessing you with miracles to move the mountains in your life.

The more you learn about the Savior, the easier it will be to trust in His mercy, His infinite love, and His strengthening, healing, and redeeming power. The Savior is never closer to you than when you are facing or climbing a mountain with faith.

Second, choose to believe in Jesus Christ. If you have doubts about God the Father and His Beloved Son or the validity of the Restoration or the veracity of Joseph Smith’s divine calling as a prophet, choose to believe and stay faithful. Take your questions to the Lord and to other faithful sources. Study with the desire to believe rather than with the hope that you can find a flaw in the fabric of a prophet’s life or a discrepancy in the scriptures. Stop increasing your doubts by rehearsing them with other doubters. Allow the Lord to lead you on your journey of spiritual discovery.

Third, act in faith. What would you do if you had more faith? Think about it. Write about it. Then receive more faith by doing something that requires more faith.

Fourth, partake of sacred ordinances worthily. Ordinances unlock the power of God for your life.

And fifth, ask your Heavenly Father, in the name of Jesus Christ, for help.

Faith takes work. Receiving revelation takes work. But “every one that asketh receiveth; and he that seeketh findeth; and to him that knocketh it shall be opened.” God knows what will help your faith grow. Ask, and then ask again.

A nonbeliever might say that faith is for the weak. But this assertion overlooks the power of faith. Would the Savior’s Apostles have continued to teach His doctrine after His death, at the peril of their lives, if they had doubted Him? Would Joseph and Hyrum Smith have suffered martyrs’ deaths defending the Restoration of the Lord’s Church unless they had a sure witness that it was true? Would nearly 2,000 Saints have died along the pioneer trail if they did not have faith that the gospel of Jesus Christ had been restored? Truly, faith is the power that enables the unlikely to accomplish the impossible.

Do not minimize the faith you already have. It takes faith to join the Church and remain faithful. It takes faith to follow prophets rather than pundits and popular opinion. It takes faith to serve a mission during a pandemic. It takes faith to live a chaste life when the world shouts that God’s law of chastity is now outmoded. It takes faith to teach the gospel to children in a secular world. It takes faith to plead for the life of a loved one and even more faith to accept a disappointing answer.

Two years ago, Sister Nelson and I visited Samoa, Tonga, Fiji, and Tahiti. Each of those island nations had experienced heavy rains for days. Members had fasted and prayed that their outdoor meetings would be protected from the rain.

In Samoa, Fiji, and Tahiti, just as the meetings began, the rain stopped. But in Tonga, the rain did not stop. Yet 13,000 faithful Saints came hours early to get a seat, waited patiently through a steady downpour, and then sat through a very wet two-hour meeting.

We saw vibrant faith at work among each of those islanders—faith sufficient to stop the rain and faith to persevere when the rain did not stop.

The mountains in our lives do not always move how or when we would like. But our faith will always propel us forward. Faith always increases our access to godly power.

Please know this: if everything and everyone else in the world whom you trust should fail, Jesus Christ and His Church will never fail you. The Lord never slumbers, nor does He sleep. He “is the same yesterday, today, and [tomorrow].” He will not forsake His covenants, His promises, or His love for His people. He works miracles today, and He will work miracles tomorrow.

Faith in Jesus Christ is the greatest power available to us in this life. All things are possible to them that believe.

Your growing faith in Him will move mountains—not the mountains of rock that beautify the earth but the mountains of misery in your lives. Your flourishing faith will help you turn challenges into unparalleled growth and opportunity.

On this Easter Sunday, with my deep feelings of love and gratitude, I declare my witness that Jesus Christ is indeed risen. He is risen to lead His Church. He is risen to bless the lives of all of God’s children, wherever they live. With faith in Him, we can move the mountains in our lives. I so testify in the sacred name of Jesus Christ, amen.

\newpage
\settalk{COVID-19 and Temples}
\setspeaker{Russell M. Nelson}
\settalkdate{April 2021}
\talkheading{COVID-19 and Temples}{Russell M. Nelson}

My beloved brothers and sisters, we have truly had a spiritual feast. How grateful I am for the prayers, messages, and music of the entire conference. Thanks to each of you for joining with us, wherever you are.

Early last year, because of the COVID-19 pandemic and our desire to be good global citizens, we made the difficult decision to close all temples temporarily. During the ensuing months, we have felt inspired to reopen temples gradually through a very cautious approach. Temples are now being opened in four phases, adhering strictly to local government regulations and safety protocols.

For temples in phase 1, qualified couples who have previously received their own endowments can be sealed as husband and wife.

For temples in phase 2, all ordinances for the living are performed, including one’s own endowment, the sealing of husband and wife, and the sealing of children to parents. We have recently amended provisions of phase 2 and allow our youth, new members, and others with a limited-use recommend to participate in proxy baptisms for their ancestors.

For temples in phase 3, those with scheduled appointments may participate not only in ordinances for the living but also in all proxy ordinances for deceased ancestors.

Phase 4 is a return to full, regular temple activity.

We are grateful for your patience and devoted service during this changing and challenging period. I pray that your desire to worship and serve in the temple burns more brightly now than ever.

You may be wondering when you will be able to return to the temple. Answer: Your temple will be open when local government regulations allow it. When the incidence of COVID-19 in your area is within safe limits, your temple will be reopened. Do all you can to bring COVID numbers down in your area so that your temple opportunities can increase.

Meanwhile, keep your temple covenants and blessings foremost in your minds and hearts. Stay true to the covenants you have made.

We are building now for the future! Forty-one temples are presently under construction or renovation. Just last year, despite the pandemic, ground was broken for 21 new temples!

We want to bring the house of the Lord even closer to our members, that they may have the sacred privilege of attending the temple as often as their circumstances allow.

As I announce our plans to construct 20 more temples, I ponder and praise pioneers—past and present—whose consecrated lives have helped to make this history today. A new temple will be built in each of the following locations: Oslo, Norway; Brussels, Belgium; Vienna, Austria; Kumasi, Ghana; Beira, Mozambique; Cape Town, South Africa; Singapore, Republic of Singapore; Belo Horizonte, Brazil; Cali, Colombia; Querétaro, Mexico; Torreón, Mexico; Helena, Montana; Casper, Wyoming; Grand Junction, Colorado; Farmington, New Mexico; Burley, Idaho; Eugene, Oregon; Elko, Nevada; Yorba Linda, California; and Smithfield, Utah.

Temples are a vital part of the Restoration of the gospel of Jesus Christ in its fulness. Ordinances of the temple fill our lives with power and strength available in no other way. We thank God for those blessings.

As we close this conference, we again express our love for you. We pray that God will shower His blessings and watchcare upon each of you. Together we are engaged in His sacred service. With courage, let us all press on in the glorious work of the Lord! For this I pray in the sacred name of Jesus Christ, amen.

\newpage
\settalk{Pure Truth, Pure Doctrine, and Pure Revelation}
\setspeaker{Russell M. Nelson}
\settalkdate{October 2021}
\talkheading{Pure Truth, Pure Doctrine, and Pure Revelation}{Russell M. Nelson}

My dear brothers and sisters, welcome to general conference! What a joy it is to be with you! You have been on my mind almost constantly during the past six months. I have prayed about you and for you. During recent weeks I have prayed intently that this conference would be a time of revelation and reflection for all who seek those blessings.

We are delighted to be speaking to you from the Conference Center once again. Most of the seats remain empty, but the presence of some members of the Tabernacle Choir is a wonderful step forward. We welcome you all to this largely virtual conference, wherever you are.

We are still dealing with the ravages of COVID-19 and its variants. We thank you for following our counsel and the advice of medical experts and government officials in your own communities.

We convene each general conference as directed by the Lord. The format has varied over the years. When I was much younger, conference lasted three or four days. Later, conference was reduced to two days. Every message—then and now—is the result of earnest prayer and much spiritual preparation.

The General Authorities and General Officers of the Church who speak will focus their messages on our Savior, Jesus Christ, His mercy, and His infinite redeeming power. There has never been a time in the history of the world when knowledge of our Savior is more personally vital and relevant to every human soul. Imagine how quickly the devastating conflicts throughout the world—and those in our individual lives—would be resolved if we all chose to follow Jesus Christ and heed His teachings.

In that spirit, I invite you to listen for three things during this conference: pure truth, the pure doctrine of Christ, and pure revelation. Contrary to the doubts of some, there really is such a thing as right and wrong. There really is absolute truth—eternal truth. One of the plagues of our day is that too few people know where to turn for truth. I can assure you that what you will hear today and tomorrow constitutes pure truth.

The pure doctrine of Christ is powerful. It changes the life of everyone who understands it and seeks to implement it in his or her life. The doctrine of Christ helps us find and stay on the covenant path. Staying on that narrow but well-defined path will ultimately qualify us to receive all that God has. Nothing could be worth more than all our Father has!

Finally, pure revelation for the questions in your heart will make this conference rewarding and unforgettable. If you have not yet sought for the ministering of the Holy Ghost to help you hear what the Lord would have you hear during these two days, I invite you to do so now. Please make this conference a time of feasting on messages from the Lord through His servants. Learn how to apply them in your life.

This is The Church of Jesus Christ of Latter-day Saints. We are His covenant people. The Lord declared that He would hasten His work in its time, and He is doing so at an ever-increasing pace. We are privileged to participate in His holy work.

I invoke a blessing upon all who are seeking greater light, knowledge, and truth. I express my love for each of you, in the sacred name of Jesus Christ, amen.

\newpage
\settalk{The Temple and Your Spiritual Foundation}
\setspeaker{Russell M. Nelson}
\settalkdate{October 2021}
\talkheading{The Temple and Your Spiritual Foundation}{Russell M. Nelson}

My dear brothers and sisters, I am grateful to be with you this morning to share the feelings of my heart.

As you know, we are performing major renovations on the historic Salt Lake Temple. This complex project includes major reinforcement of its original foundation, which has served well for more than a century. But this temple must stand much longer. In late May, I inspected the progress on this massive project. I thought you would appreciate seeing what my wife Wendy and I saw. I think you’ll see why the hymn “How Firm a Foundation” has come to have new meaning for us.

Video from the site of the Salt Lake Temple renovation: “We are looking at the original foundation of the Salt Lake Temple. I am standing in an area beneath what was the Garden Room. As I examine the craftsmanship of this entire building, I marvel at what the pioneers accomplished. I am totally in awe when I consider that they built this magnificent temple with only tools and techniques available to them more than a century ago.

“These many decades later, however, if we examine the foundation closely, we can see the effects of erosion, gaps in the original stonework, and varying stages of stability in the masonry.

“Now as I witness what modern engineers, architects, and construction experts can do to reinforce that original foundation, I am absolutely amazed. Their work is astonishing!

“The foundation of any building, particularly one as large as this one, must be strong and resilient enough to withstand earthquakes, corrosion, high winds, and the inevitable settling that affects all buildings. The complex task of strengthening now underway will reinforce this sacred temple with the foundation that can and will stand the test of time.”

We are sparing no effort to give this venerable temple, which had become increasingly vulnerable, a foundation that will withstand the forces of nature into the Millennium. In like manner, it is now time that we each implement extraordinary measures—perhaps measures we have never taken before—to strengthen our personal spiritual foundations. Unprecedented times call for unprecedented measures.

My dear brothers and sisters, these are the latter days. If you and I are to withstand the forthcoming perils and pressures, it is imperative that we each have a firm spiritual foundation built upon the rock of our Redeemer, Jesus Christ.

So I ask each of you, how firm is your foundation? And what reinforcements to your testimony and understanding of the gospel are needed?

The temple lies at the center of strengthening our faith and spiritual fortitude because the Savior and His doctrine are the very heart of the temple. Everything taught in the temple, through instruction and through the Spirit, increases our understanding of Jesus Christ. His essential ordinances bind us to Him through sacred priesthood covenants. Then, as we keep our covenants, He endows us with His healing, strengthening power. And oh, how we will need His power in the days ahead.

We have been promised that “if [we] are prepared [we] shall not fear.” This assurance has profound implications today. The Lord has declared that despite today’s unprecedented challenges, those who build their foundations upon Jesus Christ, and have learned how to draw upon His power, need not succumb to the unique anxieties of this era.

Temple ordinances and covenants are ancient. The Lord instructed Adam and Eve to pray, make covenants, and offer sacrifices. Indeed, “whenever the Lord has had a people on the earth who will obey His word, they have been commanded to build temples.” The standard works are replete with references to temple teachings, clothing, language, and more. Everything we believe and every promise God has made to His covenant people come together in the temple. In every age, the temple has underscored the precious truth that those who make covenants with God and keep them are children of the covenant.

Thus, in the house of the Lord, we can make the same covenants with God that Abraham, Isaac, and Jacob made. And we can receive the same blessings!

Temples have been part of this dispensation from its earliest days. Elijah committed the keys of sealing authority to Joseph Smith in the Kirtland Temple. The fulness of the priesthood was restored in the Nauvoo Temple.

Until his martyrdom, Joseph Smith continued to receive revelations that furthered the restoration of the endowment and sealing ordinances. He recognized, however, that further refinement was needed. After administering the endowment to Brigham Young in May 1842, Joseph told Brigham, “This is not arranged right, but we have done the best we could under the circumstances in which we are placed, and I wish you to take this matter in hand and organize and systematize all these ceremonies.”

Following the Prophet’s death, President Young oversaw the completion of the Nauvoo Temple and later built temples in the Utah Territory. At the dedication of the lower stories of the St. George Temple, Brigham Young vigorously declared the urgency of vicarious temple work when he said, “When I think upon this subject, I want the tongues of seven thunders to wake up the people.”

From that time forward, temple ordinances were gradually refined. President Harold B. Lee explained why procedures, policies, and even the administration of temple ordinances continue to change within the Savior’s restored Church. President Lee said: “The principles of the gospel of Jesus Christ are divine. Nobody changes the principles and [doctrine] of the Church except the Lord by revelation. But methods change as the inspired direction comes to those who preside at a given time.”

Consider how administering the sacrament has changed over the years. In earlier days, the water of the sacrament was offered to the congregation in one large vessel. Everyone drank from it. Now we use individual disposable cups. The procedure changed, but the covenants remain the same.

Ponder these three truths:

The First Presidency and Quorum of the Twelve Apostles have often asked the Lord if there are better ways to take the blessings of the temple to His faithful children. We seek guidance regularly on how to ensure worldwide accuracy and consistency of temple instruction, covenants, and ordinances despite differences in language and culture.

Under the Lord’s direction and in answer to our prayers, recent procedural adjustments have been made. He is the One who wants you to understand with great clarity exactly what you are making covenants to do. He is the One who wants you to experience fully His sacred ordinances. He wants you to comprehend your privileges, promises, and responsibilities. He wants you to have spiritual insights and awakenings you’ve never had before. This He desires for all temple patrons, no matter where they live.

Current adjustments in temple procedures, and others that will follow, are continuing evidence that the Lord is actively directing His Church. He is providing opportunities for each of us to bolster our spiritual foundations more effectively by centering our lives on Him and on the ordinances and covenants of His temple. When you bring your temple recommend, a contrite heart, and a seeking mind to the Lord’s house of learning, He will teach you.

Should distance, health challenges, or other constraints prohibit your temple attendance for a season, I invite you to set a regular time to rehearse in your mind the covenants you have made.

If you don’t yet love to attend the temple, go more often—not less. Let the Lord, through His Spirit, teach and inspire you there. I promise you that over time, the temple will become a place of safety, solace, and revelation.

If it were possible for me to speak one-on-one with every young adult, I would plead with you to seek a companion with whom you can be sealed in the temple. You may wonder what difference this will make in your life. I promise it will make all the difference! As you marry in the temple and return repeatedly, you will be strengthened and guided in your decisions.

If I could speak with each husband and wife who have still not been sealed in the temple, I would plead with you to take the necessary steps to receive that crowning, life-changing ordinance. Will it make a difference? Only if you want to progress forever and be together forever. Wishing to be together forever will not make it so. No other ceremony or contract will make it so.

If I could speak to each man or woman who longs for marriage but has not yet found his or her eternal companion, I would urge you not to wait until marriage to be endowed in the house of the Lord. Begin now to learn and experience what it means to be armed with priesthood power.

And to each of you who has made temple covenants, I plead with you to seek—prayerfully and consistently—to understand temple covenants and ordinances. Spiritual doors will open. You will learn how to part the veil between heaven and earth, how to ask for God’s angels to attend you, and how better to receive direction from heaven. Your diligent efforts to do so will reinforce and strengthen your spiritual foundation.

My dear brothers and sisters, when renovations on the Salt Lake Temple are completed, there will be no safer place during an earthquake in the Salt Lake Valley than inside that temple.

Likewise, whenever any kind of upheaval occurs in your life, the safest place to be spiritually is living inside your temple covenants!

Please believe me when I say that when your spiritual foundation is built solidly upon Jesus Christ, you have no need to fear. As you are true to your covenants made in the temple, you will be strengthened by His power. Then, when spiritual earthquakes occur, you will be able to stand strong because your spiritual foundation is solid and immovable.

I love you, dear brothers and sisters. These truths I know: God, our Heavenly Father, wants you to choose to come home to Him. His plan of eternal progression is not complicated, and it honors your agency. You are free to choose who you will be—and with whom you will be—in the world to come!

God lives! Jesus is the Christ! This is His Church, restored to help you fulfill your divine destiny. I so testify in the sacred name of Jesus Christ, amen.

\newpage
\settalk{Make Time for the Lord}
\setspeaker{Russell M. Nelson}
\settalkdate{October 2021}
\talkheading{Make Time for the Lord}{Russell M. Nelson}

My dear brothers and sisters, for two days we have been well taught by servants of the Lord who have sought diligently to know what He would have them say.

We have been given our charge for the next six months. Now the question is, how will we be different because of what we have heard and felt?

The pandemic has demonstrated how quickly life can change, at times from circumstances beyond our control. However, there are many things we can control. We set our own priorities and determine how we use our energy, time, and means. We decide how we will treat each other. We choose those to whom we will turn for truth and guidance.

The voices and pressures of the world are engaging and numerous. But too many voices are deceptive, seductive, and can pull us off the covenant path. To avoid the inevitable heartbreak that follows, I plead with you today to counter the lure of the world by making time for the Lord in your life—each and every day.

If most of the information you get comes from social or other media, your ability to hear the whisperings of the Spirit will be diminished. If you are not also seeking the Lord through daily prayer and gospel study, you leave yourself vulnerable to philosophies that may be intriguing but are not true. Even Saints who are otherwise faithful can be derailed by the steady beat of Babylon’s band.

My brothers and sisters, I plead with you to make time for the Lord! Make your own spiritual foundation firm and able to stand the test of time by doing those things that allow the Holy Ghost to be with you always.

Never underestimate the profound truth that “the Spirit speaketh … of things as they really are, and of things as they really will be.” “It will show unto you all things what ye should do.”

Nothing invites the Spirit more than fixing your focus on Jesus Christ. Talk of Christ, rejoice in Christ, feast upon the words of Christ, and press forward with steadfastness in Christ. Make your Sabbath a delight as you worship Him, partake of the sacrament, and keep His day holy.

As I emphasized this morning, please make time for the Lord in His holy house. Nothing will strengthen your spiritual foundation like temple service and temple worship.

We thank all who are working on our new temples. They are being built all over the world. Today I am pleased to announce our plans to build more temples at or near the following locations: Kaohsiung, Taiwan; Tacloban, Philippines; Monrovia, Liberia; Kananga, Democratic Republic of the Congo; Antananarivo, Madagascar; Culiacán, Mexico; Vitória, Brazil; La Paz, Bolivia; Santiago West, Chile; Fort Worth, Texas; Cody, Wyoming; Rexburg North, Idaho; Heber Valley, Utah; and reconstruction of the Provo Utah Temple after the Orem Utah Temple is dedicated.

I love you, dear brothers and sisters. The Lord knows you and loves you. He is your Savior and your Redeemer. He leads and guides His Church. He will lead and guide you in your personal life if you will make time for Him in your life—each and every day.

May God be with you until we meet again, I pray in the sacred name of Jesus Christ, amen.

\newpage
\settalk{Preaching the Gospel of Peace}
\setspeaker{Russell M. Nelson}
\settalkdate{April 2022}
\talkheading{Preaching the Gospel of Peace}{Russell M. Nelson}

My dear brothers and sisters, welcome to general conference! I have looked forward to this day with great anticipation. I pray for you every day. I have also prayed that this conference will be a time of spiritual rejuvenation for each one of you.

Since last conference, difficulties in the world have continued. The global pandemic still affects our lives. And now the world has been rocked by a conflict that is raining terror on millions of innocent men, women, and children.

Prophets have foreseen our day, when there would be wars and rumors of wars and when the whole earth would be in commotion. As followers of Jesus Christ, we plead with leaders of nations to find peaceful resolutions to their differences. We call upon people everywhere to pray for those in need, to do what they can to help the distressed, and to seek the Lord’s help in ending any major conflicts.

Brothers and sisters, the gospel of Jesus Christ has never been needed more than it is today. Contention violates everything the Savior stood for and taught. I love the Lord Jesus Christ and testify that His gospel is the only enduring solution for peace. His gospel is a gospel of peace.

His gospel is the only answer when many in the world are stunned with fear. This underscores the urgent need for us to follow the Lord’s instruction to His disciples to “go … into all the world, and preach the gospel to every creature.” We have the sacred responsibility to share the power and peace of Jesus Christ with all who will listen and who will let God prevail in their lives.

Every person who has made covenants with God has promised to care about others and serve those in need. We can demonstrate faith in God and always be ready to respond to those who ask about “the hope that is in [us].” Each of us has a role to play in the gathering of Israel.

Today I reaffirm strongly that the Lord has asked every worthy, able young man to prepare for and serve a mission. For Latter-day Saint young men, missionary service is a priesthood responsibility. You young men have been reserved for this time when the promised gathering of Israel is taking place. As you serve missions, you play a pivotal role in this unprecedented event!

For you young and able sisters, a mission is also a powerful, but optional, opportunity. We love sister missionaries and welcome them wholeheartedly. What you contribute to this work is magnificent! Pray to know if the Lord would have you serve a mission, and the Holy Ghost will respond to your heart and mind.

Dear young friends, you are each vital to the Lord. He has held you in reserve until now to help gather Israel. Your decision to serve a mission, whether a proselyting or a service mission, will bless you and many others. We also welcome senior couples to serve when their circumstances permit. Their efforts are simply irreplaceable.

All missionaries teach and testify of the Savior. The spiritual darkness in the world makes the light of Jesus Christ needed more than ever. Everyone deserves the chance to know about the restored gospel of Jesus Christ. Every person deserves to know where they can find the hope and peace that “[pass] all understanding.”

May this conference be a time of peace and spiritual feasting for you. May you seek and receive personal revelation during these sessions, I pray in the sacred name of Jesus Christ, amen.

\newpage
\settalk{The Power of Spiritual Momentum}
\setspeaker{Russell M. Nelson}
\settalkdate{April 2022}
\talkheading{The Power of Spiritual Momentum}{Russell M. Nelson}

My dear brothers and sisters, I love you. I cherish this opportunity to speak with you today. I pray daily that you will be protected from the fierce attacks of the adversary and have the strength to push forward through whatever challenges you face.

Some trials are deeply private burdens no one else can see. Others are played out on the world stage. The armed conflict in eastern Europe is one of these. I have been to Ukraine and Russia many times. I love those lands, the people, and their languages. I weep and pray for all who are affected by this conflict. As a church we’re doing all we can to help those who are suffering and struggling to survive. We invite everyone to continue to fast and pray for all the people being hurt by this calamity. Any war is a horrifying violation of everything the Lord Jesus Christ stands for and teaches.

None of us can control nations or the actions of others or even members of our own families. But we can control ourselves. My call today, dear brothers and sisters, is to end conflicts that are raging in your heart, your home, and your life. Bury any and all inclinations to hurt others—whether those inclinations be a temper, a sharp tongue, or a resentment for someone who has hurt you. The Savior commanded us to turn the other cheek, to love our enemies, and to pray for those who despitefully use us.

It can be painfully difficult to let go of anger that feels so justified. It can seem impossible to forgive those whose destructive actions have hurt the innocent. And yet, the Savior admonished us to “forgive all men.”

We are followers of the Prince of Peace. Now more than ever, we need the peace only He can bring. How can we expect peace to exist in the world when we are not individually seeking peace and harmony? Brothers and sisters, I know what I’m suggesting is not easy. But followers of Jesus Christ should set the example for all the world to follow. I plead with you to do all you can to end personal conflicts that are currently raging in your hearts and in your lives.

May I underscore this call to action by discussing a concept I was reminded of recently while watching a basketball game.

In that game, the first half was a seesaw battle, back and forth. Then, during the last five seconds of the first half, a guard on one team made a beautiful three-point shot. With only one second left, his teammate stole the inbound pass and made another basket at the buzzer! So that team went into the locker room four points ahead with a palpable surge of momentum. They were able to carry that momentum into the second half and win the game.

Momentum is a powerful concept. We all have experienced it in one form or another—for example, in a vehicle that picks up speed or with a disagreement that suddenly turns into an argument.

So I ask, what can ignite spiritual momentum? We have seen examples of both positive and negative momentum. We know followers of Jesus Christ who became converted and grew in their faith. But we also know of once-committed believers who fell away. Momentum can swing either way.

We have never needed positive spiritual momentum more than we do now, to counteract the speed with which evil and the darker signs of the times are intensifying. Positive spiritual momentum will keep us moving forward amid the fear and uncertainty created by pandemics, tsunamis, volcanic eruptions, and armed hostilities. Spiritual momentum can help us withstand the relentless, wicked attacks of the adversary and thwart his efforts to erode our personal spiritual foundation.

Many actions can ignite positive spiritual momentum. Obedience, love, humility, service, and gratitude are but a few.

Today I would like to suggest five specific actions we can take to help us maintain positive spiritual momentum.

First: Get on the covenant path and stay there.

Not long ago, I had a vivid dream in which I met a large group of people. They asked me many questions, the most frequent of which was about the covenant path and why it is so important.

In my dream, I explained that we enter the covenant path by being baptized and making our first covenant with God. Each time we partake of the sacrament, we promise again to take the name of the Savior upon us, to remember Him, and to keep His commandments. In return, God assures us that we may always have the Spirit of the Lord to be with us.

Later we make additional covenants in the temple, where we receive even greater promises. Ordinances and covenants give us access to godly power. The covenant path is the only path that leads to exaltation and eternal life.

In my dream, a woman then asked how someone who has broken his or her covenants can get back on that path. My answer to her question leads to my second suggestion:

Discover the joy of daily repentance.

How important is repentance? Alma taught that we should “preach nothing save it were repentance and faith on the Lord.” Repentance is required of every accountable person who desires eternal glory. There are no exceptions. In a revelation to the Prophet Joseph Smith, the Lord chastised early Church leaders for not teaching the gospel to their children. Repenting is the key to progress. Pure faith keeps us moving forward on the covenant path.

Please do not fear or delay repenting. Satan delights in your misery. Cut it short. Cast his influence out of your life! Start today to experience the joy of putting off the natural man. The Savior loves us always but especially when we repent. He promised that though “the mountains shall depart, and the hills be removed … my kindness shall not depart from thee.”

If you feel you have strayed off the covenant path too far or too long and have no way to return, that simply is not true. Please contact your bishop or branch president. He is the Lord’s agent and will help you experience the joy and relief of repenting.

Now, a caution: Returning to the covenant path does not mean that life will be easy. This path is rigorous and at times will feel like a steep climb. This ascent, however, is designed to test and teach us, refine our natures, and help us to become saints. It is the only path that leads to exaltation. One prophet described the “blessed and happy state of those that keep the commandments of God. For behold, they are blessed in all things, both temporal and spiritual; and if they hold out faithful to the end they are received into heaven … [and] dwell with God in a state of never-ending happiness.”

Walking the covenant path, coupled with daily repentance, fuels positive spiritual momentum.

My third suggestion: Learn about God and how He works.

One of our greatest challenges today is distinguishing between the truths of God and the counterfeits of Satan. That is why the Lord warned us to “pray always, … that [we] may conquer Satan, and … escape the hands of the servants of Satan that do uphold his work.”

Moses provided an example of how to discern between God and Satan. When Satan came tempting Moses, he detected the deception because he had just had a face-to-face interaction with God. Moses quickly realized who Satan was and commanded him to depart. When Satan persisted, Moses knew how to call upon God for more help. Moses received divine strength and rebuked the evil one again, saying, “Depart from me, Satan, for this one God only will I worship.”

We should follow that example. Cast Satan’s influence out of your life! Please do not follow him down to his “gulf of misery and endless wo.”

With frightening speed, a testimony that is not nourished daily “by the good word of God” can crumble. Thus, the antidote to Satan’s scheme is clear: we need daily experiences worshipping the Lord and studying His gospel. I plead with you to let God prevail in your life. Give Him a fair share of your time. As you do, notice what happens to your positive spiritual momentum.

Suggestion number 4: Seek and expect miracles.

Moroni assured us that “God has not ceased to be a God of miracles.” Every book of scripture demonstrates how willing the Lord is to intervene in the lives of those who believe in Him. He parted the Red Sea for Moses, helped Nephi retrieve the brass plates, and restored His Church through the Prophet Joseph Smith. Each of these miracles took time and may not have been exactly what those individuals originally requested from the Lord.

In the same way, the Lord will bless you with miracles if you believe in Him, “doubting nothing.” Do the spiritual work to seek miracles. Prayerfully ask God to help you exercise that kind of faith. I promise that you can experience for yourself that Jesus Christ “giveth power to the faint; and to them that have no might he increaseth strength.” Few things will accelerate your spiritual momentum more than realizing the Lord is helping you to move a mountain in your life.

Suggestion number 5: End conflict in your personal life.

I repeat my call to end the conflicts in your life. Exercise the humility, courage, and strength required both to forgive and to seek forgiveness. The Savior has promised that “if [we] forgive men their trespasses, [our] heavenly Father will also forgive [us].”

Two weeks from today we celebrate Easter. Between now and then, I invite you to seek an end to a personal conflict that has weighed you down. Could there be a more fitting act of gratitude to Jesus Christ for His Atonement? If forgiveness presently seems impossible, plead for power through the atoning blood of Jesus Christ to help you. As you do so, I promise personal peace and a burst of spiritual momentum.

When the Savior atoned for all mankind, He opened a way that those who follow Him can have access to His healing, strengthening, and redeeming power. These spiritual privileges are available to all who seek to hear Him and follow Him.

My dear brothers and sisters, with all the pleadings of my heart, I urge you to get on the covenant path and stay there. Experience the joy of repenting daily. Learn about God and how He works. Seek and expect miracles. Strive to end conflict in your life.

As you act on these pursuits, I promise you the ability to move forward on the covenant path with increased momentum, despite whatever obstacles you face. And I promise you greater strength to resist temptation, more peace of mind, freedom from fear, and greater unity in your families.

God lives! Jesus is the Christ! He lives! He loves us and will help us. Of this I testify in the sacred name of our Redeemer, Jesus Christ, amen.

\newpage
\settalk{Now Is the Time}
\setspeaker{Russell M. Nelson}
\settalkdate{April 2022}
\talkheading{Now Is the Time}{Russell M. Nelson}

My dear brothers and sisters, this conference has been historic in many ways. We have been blessed by the prayers, messages, and music. We have been inspired by servants of the Lord.

We have received important direction for the future. My prayer is that the Spirit has spoken to you directly about things the Lord would have you do.

The future is always uncertain. Weather changes. Economic cycles are unpredictable. Disasters, accidents, and illness can change life quickly. These actions are largely beyond our control. But there are some things we can control, including how we spend our time each day.

I like this poem by Henry Van Dyke, posted on a sundial at Wells College in New York. It reads:

\begin{verse}
The shadow by my finger cast\\
Divides the future from the past:\\
Before it, sleeps the unborn hour\\
In darkness, and beyond thy power:\\
Behind its unreturning line,\\
The vanished hour, no longer thine:\\
One hour alone is in thy hands,—\\
The NOW on which the shadow stands.
\end{verse}

Yes, we should learn from the past, and yes, we should prepare for the future. But only now can we do. Now is the time we can learn. Now is the time we can repent. Now is the time we can bless others and “lift up the hands which hang down.” As Mormon counseled his son Moroni, “Let us labor diligently; … for we have a labor to perform [while] in this tabernacle of clay, that we may conquer the enemy of all righteousness, and rest our souls in the kingdom of God.”

The adversary never sleeps. There will always be opposition to the truth. I repeat my urging from this morning to do those things that will increase your positive spiritual momentum, that lift Elder Dieter F. Uchtdorf was talking about, that will keep you moving forward through whatever challenges and opportunities come.

Positive spiritual momentum increases as we worship in the temple and grow in our understanding of the magnificent breadth and depth of the blessings we receive there. I plead with you to counter worldly ways by focusing on the eternal blessings of the temple. Your time there brings blessings for eternity.

As the Church grows, we strive to keep pace by building more temples. Forty-four new temples are presently under construction. More are being renewed. I pray for the skilled people who work on those projects across the world.

In the spirit of prayerful gratitude, I am pleased to announce our plans to build a new temple in each of the following locations: Wellington, New Zealand; Brazzaville, Republic of the Congo; Barcelona, Spain; Birmingham, United Kingdom; Cusco, Peru; Maceió, Brazil; Santos, Brazil; San Luis Potosí, Mexico; Mexico City Benemérito, Mexico; Tampa, Florida; Knoxville, Tennessee; Cleveland, Ohio; Wichita, Kansas; Austin, Texas; Missoula, Montana; Montpelier, Idaho; and Modesto, California.

These 17 temples will bless countless lives on both sides of the veil. I love you, my dear brothers and sisters. More important, the Lord loves you. He is your Savior and your Redeemer. He leads and guides His Church. May we be a people worthy of the Lord, who said, “Ye shall be my people, and I will be your God.”

For this I pray in the sacred name of Jesus Christ, amen.

\newpage
\settalk{What Is True?}
\setspeaker{Russell M. Nelson}
\settalkdate{October 2022}
\talkheading{What Is True?}{Russell M. Nelson}

My beloved brothers and sisters, thanks to all for this inspiring session! Since our conference last April, we have witnessed many world events, ranging from the heartbreaking to the sublime.

We are delighted with reports of large youth conferences being held throughout the world. At these conferences, our noble youth are learning that no matter what happens in their lives, their greatest strength comes from the Lord.

We rejoice that more temples are being built across the world. With the dedication of each new temple, additional godly power comes into the world to strengthen us and counteracts the intensifying efforts of the adversary.

Abuse constitutes the influence of the adversary. It is a grievous sin. As President of the Church, I affirm the teachings of the Lord Jesus Christ on this issue. Let me be perfectly clear: any kind of abuse of women, children, or anyone is an abomination to the Lord. He grieves and I grieve whenever anyone is harmed. He mourns and we all mourn for each person who has fallen victim to abuse of any kind. Those who perpetrate these hideous acts are not only accountable to the laws of man but will also face the wrath of Almighty God.

For decades now, the Church has taken extensive measures to protect—in particular—children from abuse. There are many aids on the Church website. I invite you to study them. These guidelines are in place to protect the innocent. I urge each of us to be alert to anyone who might be in danger of being abused and to act promptly to protect them. The Savior will not tolerate abuse, and as His disciples, neither can we.

The adversary has other disturbing tactics. Among them are his efforts to blur the line between what is true and what is not true. The flood of information available at our fingertips, ironically, makes it increasingly difficult to determine what is true.

This challenge reminds me of an experience Sister Nelson and I had when we visited a dignitary in a country where relatively few people have heard of Jesus Christ. This dear aging friend had recently been quite ill. He told us that during his many days in bed, he often stared at the ceiling and asked, “What is true?”

Many on earth today “are only kept from the truth because they know not where to find it.” Some would have us believe that truth is relative—that each person should determine for himself or herself what is true. Such a belief is but wishful thinking for those who mistakenly think they will not also be accountable to God.

Dear brothers and sisters, God is the source of all truth. The Church of Jesus Christ of Latter-day Saints embraces all truth that God conveys to His children, whether learned in a scientific laboratory or received by direct revelation from Him.

From this pulpit today and tomorrow, you will continue to hear truth. Please make notes of thoughts that catch your attention and those that come into your mind and stay in your heart. Prayerfully ask the Lord to confirm that what you have heard is true.

I love you, my dear brothers and sisters. I pray that this conference will provide the spiritual feast you are seeking. In the sacred name of Jesus Christ, amen.

\newpage
\settalk{Overcome the World and Find Rest}
\setspeaker{Russell M. Nelson}
\settalkdate{October 2022}
\talkheading{Overcome the World and Find Rest}{Russell M. Nelson}

My dear brothers and sisters, I am grateful to greet you on this glorious Sabbath morning. You are constantly on my mind. I marvel at the way you spring into action whenever you see others in need. I stand amazed at the faith and testimony you demonstrate again and again. I weep over your heartaches, disappointments, and worries. I love you. I assure you that our Heavenly Father and His Beloved Son, Jesus Christ, love you. They are intimately aware of your circumstances, your goodness, your needs, and your prayers for help. Again and again, I pray for you to feel Their love for you.

Experiencing Their love is vital, as it seems that we are accosted daily by an onslaught of sobering news. You may have had days when you wished you could don your pajamas, curl up in a ball, and ask someone to awaken you when the turmoil is over.

But, my dear brothers and sisters, so many wonderful things are ahead. In coming days, we will see the greatest manifestations of the Savior’s power that the world has ever seen. Between now and the time He returns “with power and great glory,” He will bestow countless privileges, blessings, and miracles upon the faithful.

Nonetheless, we are presently living in what surely is a most complicated time in the history of the world. The complexities and challenges leave many people feeling overwhelmed and exhausted. However, consider a recent experience that might shed light on how you and I can find rest.

During the recent open house of the Washington D.C. Temple, a member of the open house committee witnessed an insightful interchange as he escorted several prominent journalists through the temple. Somehow a young family became attached to this media tour. One reporter kept asking about the “journey” of a temple patron as he or she moves through the temple. He wanted to know if the temple journey is symbolic of the challenges in a person’s journey through life.

A young boy in the family picked up on the conversation. When the tour group entered an endowment room, the boy pointed to the altar, where people kneel to make covenants with God, and said, “Oh, that’s nice. Here is a place for people to rest on their temple journey.”

I doubt that the boy knew just how profound his observation was. He likely had no idea about the direct connection between making a covenant with God in the temple and the Savior’s stunning promise:

“Come unto me, all ye that labour and are heavy laden, and I will give you rest.

“Take my yoke upon you, and learn of me; … and ye shall find rest unto your souls.

“For my yoke is easy, and my burden is light.”

Dear brothers and sisters, I grieve for those who leave the Church because they feel membership requires too much of them. They have not yet discovered that making and keeping covenants actually makes life easier! Each person who makes covenants in baptismal fonts and in temples—and keeps them—has increased access to the power of Jesus Christ. Please ponder that stunning truth!

The reward for keeping covenants with God is heavenly power—power that strengthens us to withstand our trials, temptations, and heartaches better. This power eases our way. Those who live the higher laws of Jesus Christ have access to His higher power. Thus, covenant keepers are entitled to a special kind of rest that comes to them through their covenantal relationship with God.

Before the Savior submitted Himself to the agony of Gethsemane and Calvary, He declared to His Apostles, “In the world ye shall have tribulation: but be of good cheer; I have overcome the world.” Subsequently, Jesus entreated each of us to do the same when He said, “I will that ye should overcome the world.”

Dear brothers and sisters, my message to you today is that because Jesus Christ overcame this fallen world, and because He atoned for each of us, you too can overcome this sin-saturated, self-centered, and often exhausting world.

Because the Savior, through His infinite Atonement, redeemed each of us from weakness, mistakes, and sin, and because He experienced every pain, worry, and burden you have ever had, then as you truly repent and seek His help, you can rise above this present precarious world.

You can overcome the spiritually and emotionally exhausting plagues of the world, including arrogance, pride, anger, immorality, hatred, greed, jealousy, and fear. Despite the distractions and distortions that swirl around us, you can find true rest—meaning relief and peace—even amid your most vexing problems.

This important truth prompts three fundamental questions:

First, what does it mean to overcome the world?

Second, how do we do it?

And third, how does overcoming the world bless our lives?

What does it mean to overcome the world? It means overcoming the temptation to care more about the things of this world than the things of God. It means trusting the doctrine of Christ more than the philosophies of men. It means delighting in truth, denouncing deception, and becoming “humble followers of Christ.” It means choosing to refrain from anything that drives the Spirit away. It means being willing to “give away” even our favorite sins.

Now, overcoming the world certainly does not mean becoming perfect in this life, nor does it mean that your problems will magically evaporate—because they won’t. And it does not mean that you won’t still make mistakes. But overcoming the world does mean that your resistance to sin will increase. Your heart will soften as your faith in Jesus Christ increases. Overcoming the world means growing to love God and His Beloved Son more than you love anyone or anything else.

How, then, do we overcome the world? King Benjamin taught us how. He said that “the natural man is an enemy to God” and remains so forever “unless he yields to the enticings of the Holy Spirit, and putteth off the natural man and becometh a saint through the atonement of Christ the Lord.” Each time you seek for and follow the promptings of the Spirit, each time you do anything good—things that “the natural man” would not do—you are overcoming the world.

Overcoming the world is not an event that happens in a day or two. It happens over a lifetime as we repeatedly embrace the doctrine of Christ. We cultivate faith in Jesus Christ by repenting daily and keeping covenants that endow us with power. We stay on the covenant path and are blessed with spiritual strength, personal revelation, increasing faith, and the ministering of angels. Living the doctrine of Christ can produce the most powerful virtuous cycle, creating spiritual momentum in our lives.

As we strive to live the higher laws of Jesus Christ, our hearts and our very natures begin to change. The Savior lifts us above the pull of this fallen world by blessing us with greater charity, humility, generosity, kindness, self-discipline, peace, and rest.

Now, you may be thinking this sounds more like hard spiritual work than rest. But here is the grand truth: while the world insists that power, possessions, popularity, and pleasures of the flesh bring happiness, they do not! They cannot! What they do produce is nothing but a hollow substitute for “the blessed and happy state of those [who] keep the commandments of God.”

The truth is that it is much more exhausting to seek happiness where you can never find it! However, when you yoke yourself to Jesus Christ and do the spiritual work required to overcome the world, He, and He alone, does have the power to lift you above the pull of this world.

Now, how does overcoming the world bless our lives? The answer is clear: entering into a covenant relationship with God binds us to Him in a way that makes everything about life easier. Please do not misunderstand me: I did not say that making covenants makes life easy. In fact, expect opposition, because the adversary does not want you to discover the power of Jesus Christ. But yoking yourself with the Savior means you have access to His strength and redeeming power.

I reaffirm a profound teaching of President Ezra Taft Benson: “Men and women who turn their lives over to God will discover that He can make a lot more out of their lives than they can. He will deepen their joys, expand their vision, quicken their minds, … lift their spirits, multiply their blessings, increase their opportunities, comfort their souls, raise up friends, and pour out peace.”

These incomparable privileges follow those who seek the support of heaven to help them overcome this world. To this end, I extend to members of the entire Church the same charge I gave to our young adults last May. I urged them then—and I plead with you now—to take charge of your own testimony of Jesus Christ and His gospel. Work for it. Nurture it so that it will grow. Feed it truth. Don’t pollute it with false philosophies of unbelieving men and women. As you make the continual strengthening of your testimony of Jesus Christ your highest priority, watch for miracles to happen in your life.

My plea to you this morning is to find rest from the intensity, uncertainty, and anguish of this world by overcoming the world through your covenants with God. Let Him know through your prayers and your actions that you are serious about overcoming the world. Ask Him to enlighten your mind and send the help you need. Each day, record the thoughts that come to you as you pray; then follow through diligently. Spend more time in the temple, and seek to understand how the temple teaches you to rise above this fallen world.

As I have stated before, the gathering of Israel is the most important work taking place on earth today. One crucial element of this gathering is preparing a people who are able, ready, and worthy to receive the Lord when He comes again, a people who have already chosen Jesus Christ over this fallen world, a people who rejoice in their agency to live the higher, holier laws of Jesus Christ.

I call upon you, my dear brothers and sisters, to become this righteous people. Cherish and honor your covenants above all other commitments. As you let God prevail in your life, I promise you greater peace, confidence, joy, and yes, rest.

With the power of the holy apostleship vested in me, I bless you in your quest to overcome this world. I bless you to increase your faith in Jesus Christ and learn better how to draw upon His power. I bless you to be able to discern truth from error. I bless you to care more about the things of God than the things of this world. I bless you to see the needs of those around you and strengthen those you love. Because Jesus Christ overcame this world, you can too. I so testify in the sacred name of Jesus Christ, amen.

\newpage
\settalk{Focus on the Temple}
\setspeaker{Russell M. Nelson}
\settalkdate{October 2022}
\talkheading{Focus on the Temple}{Russell M. Nelson}

Dear brothers and sisters, during these five magnificent sessions of general conference, we have once again experienced that the heavens are open! I pray that you have recorded your impressions and will follow through with them. Our Heavenly Father and His Beloved Son, Jesus Christ, stand ready to help you. I urge you to increase your efforts to seek Their help.

Recently, Sister Nelson and I had the opportunity to preview the new season 4 of the Book of Mormon Videos series. We were inspired by them! May I show you a brief excerpt from the scene depicting the Savior’s appearance to the Nephites.

It is significant that the Savior chose to appear to the people at the temple. It is His house. It is filled with His power. Let us never lose sight of what the Lord is doing for us now. He is making His temples more accessible. He is accelerating the pace at which we are building temples. He is increasing our ability to help gather Israel. He is also making it easier for each of us to become spiritually refined. I promise that increased time in the temple will bless your life in ways nothing else can.

We currently have 168 operating temples and 53 new temples under construction and another 54 in the preconstruction design phase! I am pleased to announce our plans to build a new temple in each of the following locations: Busan, Korea; Naga, Philippines; Santiago, Philippines; Eket, Nigeria; Chiclayo, Peru; Buenos Aires City Center, Argentina; Londrina, Brazil; Ribeirão Prêto, Brazil; Huehuetenango, Guatemala; Jacksonville, Florida; Grand Rapids, Michigan; Prosper, Texas; Lone Mountain, Nevada; and Tacoma, Washington.

We are also planning to build multiple temples in selected large metropolitan areas where travel time to an existing temple is a major challenge. Therefore, I am pleased to announce four additional locations near Mexico City where new temples will be built in Cuernavaca, Pachuca, Toluca, and Tula.

My dear brothers and sisters, may you focus on the temple in ways you never have before. I bless you to grow closer to God and Jesus Christ every day. I love you. May God be with you until we meet again, I pray in the sacred name of Jesus Christ, amen.

\newpage
\settalk{Peacemakers Needed}
\setspeaker{Russell M. Nelson}
\settalkdate{April 2023}
\talkheading{Peacemakers Needed}{Russell M. Nelson}

My dear brothers and sisters, it is a joy to be with you. During these past six months, you have been constantly on my mind and in my prayers. I pray that the Holy Ghost will communicate what the Lord wants you to hear as I speak to you now.

During my surgical internship many years ago, I assisted a surgeon who was amputating a leg filled with highly infectious gangrene. The operation was difficult. Then, to add to the tension, one of the team performed a task poorly, and the surgeon erupted in anger. In the middle of his tantrum, he threw his scalpel loaded with germs. It landed in my forearm!

Everyone in the operating room—except the out-of-control surgeon—was horrified by this dangerous breach of surgical practice. Gratefully, I did not become infected. But this experience left a lasting impression on me. In that very hour, I promised myself that whatever happened in my operating room, I would never lose control of my emotions. I also vowed that day never to throw anything in anger—whether it be scalpels or words.

Even now, decades later, I find myself wondering if the contaminated scalpel that landed in my arm was any more toxic than the venomous contention that infects our civic dialogue and too many personal relationships today. Civility and decency seem to have disappeared during this era of polarization and passionate disagreements.

Vulgarity, faultfinding, and evil speaking of others are all too common. Too many pundits, politicians, entertainers, and other influencers throw insults constantly. I am greatly concerned that so many people seem to believe that it is completely acceptable to condemn, malign, and vilify anyone who does not agree with them. Many seem eager to damage another’s reputation with pathetic and pithy barbs!

Anger never persuades. Hostility builds no one. Contention never leads to inspired solutions. Regrettably, we sometimes see contentious behavior even within our own ranks. We hear of those who belittle their spouses and children, of those who use angry outbursts to control others, and of those who punish family members with the “silent treatment.” We hear of youth and children who bully and of employees who defame their colleagues.

My dear brothers and sisters, this should not be. As disciples of Jesus Christ, we are to be examples of how to interact with others—especially when we have differences of opinion. One of the easiest ways to identify a true follower of Jesus Christ is how compassionately that person treats other people.

The Savior made this clear in His sermons to followers in both hemispheres. “Blessed are the peacemakers,” He said. “Whosoever shall smite thee on thy right cheek, turn to him the other also.” And then, of course, He gave the admonition that challenges each of us: “Love your enemies, bless them that curse you, do good to them that hate you, and pray for them which despitefully use you, and persecute you.”

Before His death, the Savior commanded His Twelve Apostles to love one another as He had loved them. And then He added, “By this shall all men know that ye are my disciples, if ye have love one to another.”

The Savior’s message is clear: His true disciples build, lift, encourage, persuade, and inspire—no matter how difficult the situation. True disciples of Jesus Christ are peacemakers.

Today is Palm Sunday. We are preparing to commemorate the most important and transcendent event ever recorded on earth, which is the Atonement and Resurrection of the Lord Jesus Christ. One of the best ways we can honor the Savior is to become a peacemaker.

The Savior’s Atonement made it possible for us to overcome all evil—including contention. Make no mistake about it: contention is evil! Jesus Christ declared that those who have “the spirit of contention” are not of Him but are “of the devil, who is the father of contention, and [the devil] stirreth up the hearts of men to contend with anger, one with another.” Those who foster contention are taking a page out of Satan’s playbook, whether they realize it or not. “No man can serve two masters.” We cannot support Satan with our verbal assaults and then think that we can still serve God.

My dear brothers and sisters, how we treat each other really matters! How we speak to and about others at home, at church, at work, and online really matters. Today, I am asking us to interact with others in a higher, holier way. Please listen carefully. “If there is anything virtuous, lovely, or of good report or praiseworthy” that we can say about another person—whether to his face or behind her back—that should be our standard of communication.

If a couple in your ward gets divorced, or a young missionary returns home early, or a teenager doubts his testimony, they do not need your judgment. They need to experience the pure love of Jesus Christ reflected in your words and actions.

If a friend on social media has strong political or social views that violate everything you believe in, an angry, cutting retort by you will not help. Building bridges of understanding will require much more of you, but that is exactly what your friend needs.

Contention drives away the Spirit—every time. Contention reinforces the false notion that confrontation is the way to resolve differences; but it never is. Contention is a choice. Peacemaking is a choice. You have your agency to choose contention or reconciliation. I urge you to choose to be a peacemaker, now and always.

Brothers and sisters, we can literally change the world—one person and one interaction at a time. How? By modeling how to manage honest differences of opinion with mutual respect and dignified dialogue.

Differences of opinion are part of life. I work every day with dedicated servants of the Lord who do not always see an issue the same way. They know I want to hear their ideas and honest feelings about everything we discuss—especially sensitive issues.

My two noble counselors, President Dallin H. Oaks and President Henry B. Eyring, are exemplary in the way they express their feelings—especially when they may differ. They do so with pure love for each other. Neither suggests that he knows best and therefore must rigorously defend his position. Neither evidences the need to compete with the other. Because each is filled with charity, “the pure love of Christ,” our deliberations can be guided by the Spirit of the Lord. How I love and honor these two great men!

Charity is the antidote to contention. Charity is the spiritual gift that helps us to cast off the natural man, who is selfish, defensive, prideful, and jealous. Charity is the principal characteristic of a true follower of Jesus Christ. Charity defines a peacemaker.

When we humble ourselves before God and pray with all the energy of our hearts, God will grant us charity.

Those blessed with this supernal gift are long-suffering and kind. They do not envy others and are not caught up in their own importance. They are not easily provoked and do not think evil of others.

Brothers and sisters, the pure love of Christ is the answer to the contention that ails us today. Charity propels us “to bear one another’s burdens” rather than heap burdens upon each other. The pure love of Christ allows us “to stand as witnesses of God at all times and in all things”—especially in tense situations. Charity allows us to demonstrate how men and women of Christ speak and act—especially when under fire.

Now, I am not talking about “peace at any price.” I am talking about treating others in ways that are consistent with keeping the covenant you make when you partake of the sacrament. You covenant to always remember the Savior. In situations that are highly charged and filled with contention, I invite you to remember Jesus Christ. Pray to have the courage and wisdom to say or do what He would. As we follow the Prince of Peace, we will become His peacemakers.

At this point you may be thinking that this message would really help someone you know. Perhaps you are hoping that it will help him or her to be nicer to you. I hope it will! But I also hope that you will look deeply into your heart to see if there are shards of pride or jealousy that prevent you from becoming a peacemaker.

If you are serious about helping to gather Israel and about building relationships that will last throughout the eternities, now is the time to lay aside bitterness. Now is the time to cease insisting that it is your way or no way. Now is the time to stop doing things that make others walk on eggshells for fear of upsetting you. Now is the time to bury your weapons of war. If your verbal arsenal is filled with insults and accusations, now is the time to put them away. You will arise as a spiritually strong man or woman of Christ.

The temple can help us in our quest. There we are endowed with God’s power, giving us the ability to overcome Satan, the instigator of all contention. Cast him out of your relationships! Note that we also rebuke the adversary every time we heal a misunderstanding or refuse to take offense. Instead, we can show the tender mercy that is characteristic of true disciples of Jesus Christ. Peacemakers thwart the adversary.

Let us as a people become a true light on the hill—a light that “cannot be hid.” Let us show that there is a peaceful, respectful way to resolve complex issues and an enlightened way to work out disagreements. As you demonstrate the charity that true followers of Jesus Christ manifest, the Lord will magnify your efforts beyond your loftiest imagination.

The gospel net is the largest net in the world. God has invited all to come unto Him, “black and white, bond and free, male and female.” There is room for everyone. However, there is no room for prejudice, condemnation, or contention of any kind.

My dear brothers and sisters, the best is yet to come for those who spend their lives building up others. Today I invite you to examine your discipleship within the context of the way you treat others. I bless you to make any adjustments that may be needed so that your behavior is ennobling, respectful, and representative of a true follower of Jesus Christ.

I bless you to replace belligerence with beseeching, animosity with understanding, and contention with peace.

God lives! Jesus is the Christ. He stands at the head of this Church. We are His servants. He will help us to become His peacemakers. I so testify in the sacred name of Jesus Christ, amen.

\newpage
\settalk{The Answer Is Always Jesus Christ}
\setspeaker{Russell M. Nelson}
\settalkdate{April 2023}
\talkheading{The Answer Is Always Jesus Christ}{Russell M. Nelson}

My dear brothers and sisters, we have been spiritually fed during these last two days. The choir music was superb. Those who have spoken have been instruments for the Lord. I pray that you will seek for the Holy Ghost to guide you in your study as you contemplate the truths taught from this pulpit. They have truly been heaven sent.

One week from today is Easter Sunday. It is the most important religious observance for followers of Jesus Christ. The main reason we celebrate Christmas is because of Easter. The Come, Follow Me lesson this week will prompt you to study the Savior’s triumphal entry into Jerusalem, His cleansing of the temple, His suffering in the Garden of Gethsemane, His Crucifixion, His glorious Resurrection, and His subsequent appearances to His followers.

Savor these sacred verses and find every way you can to thank our Heavenly Father for sending us His Only Begotten Son. Because of Jesus Christ, we can repent and be forgiven of our sins. Because of Him, each of us will be resurrected.

I also invite you to study again the account of the Savior’s appearance to the Nephites in the Americas, as recorded in 3 Nephi. Not long before that appearance, His voice was heard among the people, including these words of supplication:

“Will ye not now return unto me, and repent of your sins, and be converted, that I may heal you?

“… Behold, mine arm of mercy is extended towards you, and whosoever will come, him will I receive.”

Dear brothers and sisters, Jesus Christ extends that same invitation to you today. I plead with you to come unto Him so that He can heal you! He will heal you from sin as you repent. He will heal you from sadness and fear. He will heal you from the wounds of this world.

Whatever questions or problems you have, the answer is always found in the life and teachings of Jesus Christ. Learn more about His Atonement, His love, His mercy, His doctrine, and His restored gospel of healing and progression. Turn to Him! Follow Him!

Jesus Christ is the reason we build temples. Each is His holy house. Making covenants and receiving essential ordinances in the temple, as well as seeking to draw closer to Him there, will bless your life in ways no other kind of worship can. For this reason, we are doing all within our power to make the blessings of the temple more accessible to our members around the world. Today, I am grateful to announce our plans to build a new temple in each of the following locations:

My dear brothers and sisters, I bear witness that Jesus Christ directs the affairs of His Church. I testify that following Him is the only way to enduring happiness. I know that His power is descending upon His covenant-keeping people, who are “armed with righteousness and with the power of God in great glory.” I so testify, with my love and blessing for each of you, in the sacred name of Jesus Christ, amen.

\newpage
\settalk{Think Celestial!}
\setspeaker{Russell M. Nelson}
\settalkdate{October 2023}
\talkheading{Think Celestial!}{Russell M. Nelson}

My dear brothers and sisters, I am deeply grateful to speak with you today. At my age, each new day brings wonderful as well as challenging surprises. Three weeks ago, I injured the muscles of my back. So, while I have delivered more than 100 general conference addresses standing, today I thought I would do so sitting. I pray that the Spirit will carry my message into your hearts today.

I recently celebrated my 99th birthday and thus commenced my 100th year of living. I am often asked the secret to living so long. A better question would be “What have I learned in nearly a century of living?”

Time today does not allow me to answer that question fully, but may I share one of the most crucial lessons I have learned.

I have learned that Heavenly Father’s plan for us is fabulous, that what we do in this life really matters, and that the Savior’s Atonement is what makes our Father’s plan possible.

As I have wrestled with the intense pain caused by my recent injury, I have felt even deeper appreciation for Jesus Christ and the incomprehensible gift of His Atonement. Think of it! The Savior suffered “pains and afflictions and temptations of every kind” so that He can comfort us, heal us, rescue us in times of need. Jesus Christ described His experience in Gethsemane and on Calvary: “Which suffering caused myself, even God, the greatest of all, to tremble because of pain, and to bleed at every pore.” My injury has caused me to reflect again and again on “the greatness of the Holy One of Israel.” During my healing, the Lord has manifested His divine power in peaceful and unmistakable ways.

Because of Jesus Christ’s infinite Atonement, our Heavenly Father’s plan is a perfect plan! An understanding of God’s fabulous plan takes the mystery out of life and the uncertainty out of our future. It allows each of us to choose how we will live here on earth and where we will live forever. The baseless notion that we should “eat, drink, and be merry, for tomorrow we die; and it shall be well with us” is one of the most absurd lies in the universe.

Here is the great news of God’s plan: the very things that will make your mortal life the best it can be are exactly the same things that will make your life throughout all eternity the best it can be! Today, to assist you to qualify for the rich blessings Heavenly Father has for you, I invite you to adopt the practice of “thinking celestial”! Thinking celestial means being spiritually minded. We learn from the Book of Mormon prophet Jacob that “to be spiritually-minded is life eternal.”

Mortality is a master class in learning to choose the things of greatest eternal import. Far too many people live as though this life is all there is. However, your choices today will determine three things: where you will live throughout all eternity, the kind of body with which you will be resurrected, and those with whom you will live forever. So, think celestial.

In my first message as President of the Church, I encouraged you to begin with the end in mind. This means making the celestial kingdom your eternal goal and then carefully considering where each of your decisions while here on earth will place you in the next world.

The Lord has clearly taught that only men and women who are sealed as husband and wife in the temple, and who keep their covenants, will be together throughout the eternities. He said, “All covenants, contracts, bonds, obligations, oaths, vows, performances, connections, associations, or expectations, that are not made and entered into and sealed by the Holy Spirit of promise … have an end when men are dead.”

Thus, if we unwisely choose to live telestial laws now, we are choosing to be resurrected with a telestial body. We are choosing not to live with our families forever.

So, my dear brothers and sisters, how and where and with whom do you want to live forever? You get to choose.

When you make choices, I invite you to take the long view—an eternal view. Put Jesus Christ first because your eternal life is dependent upon your faith in Him and in His Atonement. It is also dependent upon your obedience to His laws. Obedience paves the way for a joyful life for you today and a grand, eternal reward tomorrow.

When you are confronted with a dilemma, think celestial! When tested by temptation, think celestial! When life or loved ones let you down, think celestial! When someone dies prematurely, think celestial. When someone lingers with a devastating illness, think celestial. When the pressures of life crowd in upon you, think celestial! As you recover from an accident or injury, as I am doing now, think celestial!

As you focus on thinking celestial, expect to encounter opposition. Decades ago, a professional colleague criticized me for having “too much temple” in me, and more than one supervisor penalized me because of my faith. I am convinced, however, that thinking celestial enhanced my career.

As you think celestial, your heart will gradually change. You will want to pray more often and more sincerely. Please don’t let your prayers sound like a shopping list. The Lord’s perspective transcends your mortal wisdom. His response to your prayers may surprise you and will help you to think celestial.

Consider the Lord’s response to Joseph Smith when he pleaded for relief in Liberty Jail. The Lord taught the Prophet that his inhumane treatment would give him experience and be for his good. “If thou endure it well,” the Lord promised, “God shall exalt thee on high.” The Lord was teaching Joseph to think celestial and to envision an eternal reward rather than focus on the excruciating difficulties of the day. Our prayers can be—and should be—living discussions with our Heavenly Father.

As you think celestial, you will find yourself avoiding anything that robs you of your agency. Any addiction—be it gaming, gambling, debt, drugs, alcohol, anger, pornography, sex, or even food—offends God. Why? Because your obsession becomes your god. You look to it rather than to Him for solace. If you struggle with an addiction, seek the spiritual and professional help you need. Please do not let an obsession rob you of your freedom to follow God’s fabulous plan.

Thinking celestial will also help you obey the law of chastity. Few things will complicate your life more quickly than violating this divine law. For those who have made covenants with God, immorality is one of the quickest ways to lose your testimony.

Many of the adversary’s most relentless temptations involve violations of moral purity. The power to create life is the one privilege of godhood that Heavenly Father allows His mortal children to exercise. Thus, God set clear guidelines for the use of this living, divine power. Physical intimacy is only for a man and a woman who are married to each other.

Much of the world does not believe this, but public opinion is not the arbiter of truth. The Lord has declared that no unchaste person will attain the celestial kingdom. So when you make decisions regarding morality, please think celestial. And if you have been unchaste, I plead with you to repent. Come unto Christ and receive His promise of complete forgiveness as you fully repent of your sins.

As you think celestial, you will view trials and opposition in a new light. When someone you love attacks truth, think celestial, and don’t question your testimony. The Apostle Paul prophesied that “in the latter times some shall depart from the faith, giving heed to seducing spirits, and doctrines of devils.”

There is no end to the adversary’s deceptions. Please be prepared. Never take counsel from those who do not believe. Seek guidance from voices you can trust—from prophets, seers, and revelators and from the whisperings of the Holy Ghost, who “will show unto you all things what ye should do.” Please do the spiritual work to increase your capacity to receive personal revelation.

As you think celestial, your faith will increase. When I was a young intern, my income was \$15 a month. One night, my wife Dantzel asked if I was paying tithing on that meager stipend. I was not. I quickly repented and began paying the additional \$1.50 in monthly tithing.

Was the Church any different because we increased our tithing? Of course not. However, becoming a full-tithe payer changed me. That is when I learned that paying tithing is all about faith, not money. As I became a full-tithe payer, the windows of heaven began to open for me. I attribute several subsequent professional opportunities to our faithful payment of tithes.

Paying tithing requires faith, and it also builds faith in God and His Beloved Son.

Choosing to live a virtuous life in a sexualized, politicized world builds faith.

Spending more time in the temple builds faith. And your service and worship in the temple will help you to think celestial. The temple is a place of revelation. There you are shown how to progress toward a celestial life. There you are drawn closer to the Savior and given greater access to His power. There you are guided in solving the problems in your life, even your most perplexing problems.

The ordinances and covenants of the temple are of eternal significance. We continue to build more temples to make these sacred possibilities become a reality in each of your lives. We are grateful to announce our plans to build a temple in each of the following 20 locations:

The Lord is directing us to build these temples to help us think celestial. God lives. Jesus is the Christ. His Church has been restored to bless all of God’s children. I so testify in the sacred name of Jesus Christ, amen.

\newpage
\settalk{Rejoice in the Gift of Priesthood Keys}
\setspeaker{Russell M. Nelson}
\settalkdate{April 2024}
\talkheading{Rejoice in the Gift of Priesthood Keys}{Russell M. Nelson}

My dear brothers and sisters, today is an historic day for President Dallin H. Oaks and me. It was 40 years ago, on April 7, 1984, when we were sustained to the Quorum of the Twelve Apostles. We have rejoiced in each and every general conference since then, including this one. We have once again been blessed with a sacred outpouring of the Spirit. I hope you will repeatedly study the messages of this conference throughout the coming months.

When I was born, there were six functioning temples in the Church—one each in St. George, Logan, Manti, and Salt Lake City, Utah; as well as in Cardston, Alberta, Canada; and Laie, Hawaii. Two earlier temples had functioned briefly in Kirtland, Ohio, and Nauvoo, Illinois. As the body of the Church moved west, the Saints were forced to leave those two temples behind.

The Nauvoo Temple was destroyed by an arsonist’s fire. It was rebuilt and then dedicated by President Gordon B. Hinckley. The Kirtland Temple was desecrated by enemies of the Church. Later the Kirtland Temple was acquired by Community of Christ, which has owned it for many years.

Last month we announced that The Church of Jesus Christ of Latter-day Saints has purchased the Kirtland Temple, along with several significant historic sites in Nauvoo. We greatly appreciate the cordial and mutually beneficial discussions we had with leaders from Community of Christ that led to this agreement.

The Kirtland Temple has unusual significance in the Restoration of the gospel of Jesus Christ. Several events that took place there had been prophesied for millennia and were essential for the Lord’s restored Church to fulfill its latter-day mission.

The most important of these events occurred on Easter Sunday, April 3, 1836. On that day, Joseph Smith and Oliver Cowdery experienced a series of remarkable visitations. First, the Lord Jesus Christ appeared. The Prophet recorded that the Savior’s “eyes were as a flame of fire; the hair of his head was white like the pure snow; his countenance shone above the brightness of the sun; and his voice was as the sound of the rushing of great waters.”

During this visitation, the Lord affirmed His identity. He said, “I am the first and the last; I am he who liveth, I am he who was slain; I am your advocate with the Father.”

Jesus Christ then declared that He had accepted the temple as His house and made this stunning promise: “I will manifest myself to my people in mercy in this house.”

This significant promise applies to every dedicated temple today. I invite you to ponder what the Lord’s promise means for you personally.

Following the Savior’s visitation, Moses appeared. Moses conferred upon Joseph Smith the keys for the gathering of Israel and the return of the ten tribes.

When this vision closed, “Elias appeared, and committed the dispensation of the gospel of Abraham” to Joseph.

Then Elijah the prophet appeared. His appearance fulfilled Malachi’s promise that before the Second Coming, the Lord would send Elijah to “turn the heart of the fathers to the children, and the heart of the children to their fathers.” Elijah conferred the keys of the sealing power upon Joseph Smith.

The significance of these keys being returned to the earth by three heavenly messengers under the direction of the Lord cannot be overstated. Priesthood keys constitute the authority and power of presidency. Priesthood keys govern how the priesthood of God may be used to bring about the Lord’s purposes and bless all who accept the restored gospel of Jesus Christ.

It is important to note that prior to the organization of the Church, heavenly messengers had conferred the Aaronic and Melchizedek Priesthoods upon the Prophet Joseph and had given him keys of both priesthoods. These keys gave Joseph Smith authority to organize the Church in 1830.

Then in the Kirtland Temple in 1836, the conferral of these three additional priesthood keys—namely, keys of the gathering of Israel, keys of the gospel of Abraham, and keys of the sealing power—was essential. These keys authorized Joseph Smith—and all succeeding Presidents of the Lord’s Church—to gather Israel on both sides of the veil, to bless all covenant children with the blessings of Abraham, to place a ratifying seal on priesthood ordinances and covenants, and to seal families eternally. The power of these priesthood keys is infinite and breathtaking.

Consider how your life would be different if priesthood keys had not been restored to the earth. Without priesthood keys, you could not be endowed with the power of God. Without priesthood keys, the Church could serve only as a significant teaching and humanitarian organization but not much more. Without priesthood keys, none of us would have access to essential ordinances and covenants that bind us to our loved ones eternally and allow us eventually to live with God.

Priesthood keys distinguish The Church of Jesus Christ of Latter-day Saints from any other organization on earth. Many other organizations can and do make your life better here in mortality. But no other organization can and will influence your life after death.

Priesthood keys give us the authority to extend all of the blessings promised to Abraham to every covenant-keeping man and woman. Temple work makes these exquisite blessings available to all of God’s children, regardless of where or when they lived or now live. Let us rejoice that priesthood keys are once again on the earth!

I invite you to consider carefully the following three statements:

Priesthood keys restored through the Prophet Joseph Smith make it possible for every covenant-keeping man and woman to enjoy incredible personal spiritual privileges. Here again, there is much we can learn from the sacred history of the Kirtland Temple.

Joseph Smith’s dedicatory prayer of the Kirtland Temple is a tutorial about how the temple spiritually empowers you and me to meet the challenges of life in these last days. I encourage you to study that prayer, recorded in Doctrine and Covenants section 109. That dedicatory prayer, which was received by revelation, teaches that the temple is “a house of prayer, a house of fasting, a house of faith, a house of learning, a house of glory, a house of order, a house of God.”

This list of attributes is much more than a description of a temple. It is a promise about what will happen to those who serve and worship in the house of the Lord. They can expect to receive answers to prayer, personal revelation, greater faith, strength, comfort, increased knowledge, and increased power.

Time in the temple will help you to think celestial and to catch a vision of who you really are, who you can become, and the kind of life you can have forever. Regular temple worship will enhance the way you see yourself and how you fit into God’s magnificent plan. I promise you that.

We are also promised that in the temple we may “receive a fulness of the Holy Ghost.” Imagine what that promise means in terms of having the heavens open for each earnest seeker of eternal truth.

We are instructed that all who worship in the temple will have the power of God and angels having “charge over them.” How much does it increase your confidence to know that, as an endowed woman or man armed with the power of God, you do not have to face life alone? What courage does it give you to know that angels really will help you?

Finally, we are promised that “no combination of wickedness” will prevail over those who worship in the house of the Lord.

Understanding the spiritual privileges made possible in the temple is vital to each of us today.

My dear brothers and sisters, here is my promise. Nothing will help you more to hold fast to the iron rod than worshipping in the temple as regularly as your circumstances permit. Nothing will protect you more as you encounter the world’s mists of darkness. Nothing will bolster your testimony of the Lord Jesus Christ and His Atonement or help you understand God’s magnificent plan more. Nothing will soothe your spirit more during times of pain. Nothing will open the heavens more. Nothing!

The temple is the gateway to the greatest blessings God has in store for each of us, for the temple is the only place on earth where we may receive all of the blessings promised to Abraham. That is why we are doing all within our power, under the direction of the Lord, to make the temple blessings more accessible to members of the Church. Thus, we are pleased to announce that we plan to build a new temple in each of the following 15 locations:

My dear brothers and sisters, I testify that this is The Church of Jesus Christ of Latter-day Saints. He stands at its head. We are His disciples.

Let us rejoice in the restoration of priesthood keys, which make it possible for you and me to enjoy every spiritual blessing we are willing and worthy to receive. I so testify in the sacred name of Jesus Christ, amen.

\newpage
\settalk{The Lord Jesus Christ Will Come Again}
\setspeaker{Russell M. Nelson}
\settalkdate{October 2024}
\talkheading{The Lord Jesus Christ Will Come Again}{Russell M. Nelson}

My dear brothers and sisters, I am grateful that the Lord has blessed me to speak with you.

In this conference, the Lord has spoken to us through His servants. I urge you to study their messages. Use them as a litmus test of what is true and what is not during the next six months.

The preservation and renovation of the Salt Lake Temple and other areas on Temple Square have been underway for nearly five years. Present projections indicate that this work will be completed by the end of 2026. We are grateful for all who are working on this massive project.

During the last six months, we have dedicated or rededicated nine temples in five countries. Between now and the end of the year, we will dedicate five more.

Today we are pleased to announce plans to construct 17 more temples. Please listen reverently as I announce the locations.

My dear brothers and sisters, do you see what is happening right before our eyes? I pray that we will not miss the majesty of this moment! The Lord is indeed hastening His work.

Why are we building temples at such an unprecedented pace? Why? Because the Lord has instructed us to do so. The blessings of the temple help to gather Israel on both sides of the veil. These blessings also help to prepare a people who will help prepare the world for the Second Coming of the Lord!

As the prophet Isaiah prophesied, and as memorialized in Handel’s Messiah, when Jesus Christ returns, “the glory of the Lord shall be revealed, and all flesh shall see it together.”

In that day “the government shall be upon his shoulder: and his name shall be called Wonderful, Counsellor, The mighty God, The everlasting Father, The Prince of Peace.”

Jesus Christ will govern from both old Jerusalem and the New Jerusalem “built upon the American continent.” From these two centers, He will direct the affairs of His Church.

In that day the Lord will be known as “King of Kings, and Lord of Lords.” Those who are with Him will be “called, and chosen, and faithful.”

Brothers and sisters, now is the time for you and for me to prepare for the Second Coming of our Lord and Savior, Jesus the Christ. Now is the time for us to make our discipleship our highest priority. In a world filled with dizzying distractions, how can we do this?

Regular worship in the temple will help us. In the house of the Lord, we focus on Jesus Christ. We learn of Him. We make covenants to follow Him. We come to know Him. As we keep our temple covenants, we gain greater access to the Lord’s strengthening power. In the temple, we receive protection from the buffetings of the world. We experience the pure love of Jesus Christ and our Heavenly Father in great abundance! We feel peace and spiritual reassurance, in contrast to the turbulence of the world.

Here is my promise to you: Every sincere seeker of Jesus Christ will find Him in the temple. You will feel His mercy. You will find answers to your most vexing questions. You will better comprehend the joy of His gospel.

I have learned that the most crucial question we each must answer is this: To whom or to what will I give my life?

My decision to follow Jesus Christ is the most important decision I have ever made. During medical school, I gained a testimony of the divinity of God the Father and His Son, Jesus Christ. Since then our Savior has been the rock upon which I have built my life. That choice has made all the difference! That decision has made so many other decisions easier. That decision has given me purpose and direction. It has also helped me weather the storms of life. Let me share two examples:

First, when my wife Dantzel unexpectedly passed away, I could not reach any of our children. There I was, alone, devastated, and crying out for help. Gratefully, through His Spirit, the Lord taught me why my dear Dantzel had been taken home. With that understanding, I was comforted. Over time, I was better able to cope with my grief. Later, I married my beloved wife Wendy. She was a central part of my second example.

When Wendy and I were on assignment in a distant land, armed robbers put a gun to my head and pulled the trigger. But the gun did not fire. Throughout that experience, both of our lives were threatened. Yet Wendy and I felt an undeniable peace. It was the peace “which passeth all understanding.”

Brothers and sisters, the Lord will comfort you too! He will strengthen you. He will bless you with peace, even amidst chaos.

Please listen to this promise of Jesus Christ to you: “I will be on your right hand and on your left, and my Spirit shall be in your hearts, and mine angels round about you, to bear you up.”

There is no limit to the Savior’s capacity to help you. His incomprehensible suffering in Gethsemane and on Calvary was for you! His infinite Atonement is for you!

I urge you to devote time each week—for the rest of your life—to increase your understanding of the Atonement of Jesus Christ. My heart aches for those who are mired in sin and don’t know how to get out. I weep for those who struggle spiritually or who carry heavy burdens alone because they do not understand what Jesus Christ did for them.

Jesus Christ took upon Himself your sins, your pains, your heartaches, and your infirmities. You do not have to bear them alone! He will forgive you as you repent. He will bless you with what you need. He will heal your wounded soul. As you yoke yourself to Him, your burdens will feel lighter. If you will make and keep covenants to follow Jesus Christ, you will find that the painful moments of your life are temporary. Your afflictions will be “swallowed up in the joy of Christ.”

It is neither too early nor too late for you to become a devout disciple of Jesus Christ. Then you will experience fully the blessings of His Atonement. You will also be more effective in helping to gather Israel.

My dear brothers and sisters, in a coming day, Jesus Christ will return to the earth as the millennial Messiah. So today I call upon you to rededicate your lives to Jesus Christ. I call upon you to help gather scattered Israel and to prepare the world for the Second Coming of the Lord. I call upon you to talk of Christ, testify of Christ, have faith in Christ, and rejoice in Christ!

Come unto Christ and “offer your whole [soul]” to Him. This is the secret to a life of joy!

The best is yet to come, my dear brothers and sisters, because the Savior is coming again! The best is yet to come because the Lord is hastening His work. The best is yet to come as we fully turn our hearts and our lives to Jesus Christ.

I bear my solemn witness that Jesus Christ is the Son of God. I am His disciple. I am honored to be His servant. At His Second Coming, “the glory of the Lord shall be revealed, and all flesh shall see it together.” That day will be filled with joy for the righteous!

Through the power of the sacred priesthood keys I hold, I declare this truth to you and to all the world! In the name of Jesus Christ, amen.

\newpage
\settalk{Confidence in the Presence of God}
\setspeaker{Russell M. Nelson}
\settalkdate{April 2025}
\talkheading{Confidence in the Presence of God}{Russell M. Nelson}

My dear brothers and sisters, I am grateful to speak to you today in this important general conference. My eyes continue to get older. Thank you for your understanding as I deliver my message.

We live in a day of great momentum in The Church of Jesus Christ of Latter-day Saints. I am filled with joy as I see progress in so many areas.

I am especially inspired by our youth. They are giving service in large numbers. They are seeking their ancestors and performing ordinances in the temple. Our young men and women are submitting applications for missionary service in record numbers. The rising generation is rising up as stalwart followers of Jesus Christ.

Recently I met a new great-granddaughter. When I contemplate the challenges she will experience in her life, I feel a great desire to help her build faith in Jesus Christ. Living His gospel is essential for her future happiness.

She, like each of us, will face challenges. We all will experience illness, disappointment, temptation, and loss. These challenges can knock our self-confidence. However, disciples of Jesus Christ have access to a different kind of confidence.

When we make and keep covenants with God, we can have confidence that is born of the Spirit. The Lord told the Prophet Joseph Smith that our confidence can “wax strong in the presence of God.” Imagine the comfort of having confidence in the presence of God!

When I speak of having confidence before God, I am referring to having confidence in approaching God right now! I am referring to praying with confidence that Heavenly Father hears us, that He understands our needs better than we do. I am referring to having confidence that He loves us more than we can comprehend, that He sends angels to be with us and with those we love. I am referring to having confidence that He yearns to help each of us reach our highest potential.

Now, how do we gain such confidence? The Lord answers this question with these words: “Let thy bowels … be full of charity towards all men, … and let virtue garnish thy thoughts unceasingly; then shall thy confidence wax strong in the presence of God.”

There’s the key! In the Lord’s own words, charity and virtue open the way to having confidence before God! Brothers and sisters, we can do this! Our confidence can truly wax strong in the presence of God, right now!

Let’s consider both charity and virtue.

First, charity. Two years ago, I called upon us, as covenant followers of Jesus Christ, to be peacemakers. I repeat what I said then: “Anger never persuades. Hostility builds no one. Contention never leads to inspired solutions.”

True charity towards all men is the hallmark of peacemakers! It is imperative that we have charity in our discourse, both public and private. I thank those of you who took my previous counsel to heart. But we can still do better.

The present hostility in public dialogue and on social media is alarming. Hateful words are deadly weapons. Contention prevents the Holy Ghost from being our constant companion.

As followers of Jesus Christ, we should lead the way as peacemakers. As charity becomes part of our nature, we will lose the impulse to demean others. We will stop judging others. We will have charity for those from all walks of life. Charity towards all men is essential to our progress. Charity is the foundation of a godly character.

Let us plead with our Heavenly Father to fill our hearts with greater charity—especially for those who are difficult to love—for charity is a gift from our Heavenly Father for true followers of Jesus Christ. The Savior is the Prince of Peace. We are to be His instruments for peace.

Now, let us speak of virtue. The Lord tells us to garnish our thoughts unceasingly with virtue. Imagine the boost you will receive to any positive thought when you enhance it with virtue. Virtue makes everything better and happier! On the other hand, imagine what will happen when you add virtue to an impure thought, a cruel thought, or a depressing thought. Virtue will drive away those thoughts. Virtue will free you from anxious, troublesome thoughts.

Brothers and sisters, as the world grows more wicked, we need to grow increasingly pure. Our thoughts, words, and actions need to be unfailingly virtuous and filled with the pure love of Jesus Christ towards all men. The great opportunity before us is to become the people God needs us to be.

Regular worship in the house of the Lord increases our capacity for both virtue and charity. Thus, time in the temple increases our confidence before the Lord. Increased time in the temple will help us prepare for the Second Coming of our Savior, Jesus Christ. We do not know the day or the hour of His coming. But I do know that the Lord is prompting me to urge us to get ready for that “great and dreadful day.”

As we diligently seek to have charity and virtue fill our lives, our confidence in approaching God will increase. I invite you to take intentional steps to grow in your confidence before the Lord. Then, as we go to our Heavenly Father with increasing confidence, we will be filled with more joy, and your faith in Jesus Christ will increase. We will begin to experience spiritual power that exceeds our greatest hopes.

We are grateful to the Lord for the acceleration of temple-building in recent years. Under His direction, today we announce plans to construct a temple in each of these 15 locations:

I testify that Jesus Christ, the Redeemer of Israel, leads this, His Church. He is preparing to come again. May we likewise prepare to receive Him. For this I pray in the name of Jesus Christ, amen.

\end{document}
